% Options for packages loaded elsewhere
\PassOptionsToPackage{unicode}{hyperref}
\PassOptionsToPackage{hyphens}{url}
\PassOptionsToPackage{dvipsnames,svgnames,x11names}{xcolor}
\documentclass[
  dutch,
  a4paper,
]{article}
\usepackage{xcolor}
\usepackage[margin=2.5cm]{geometry}
\usepackage{amsmath,amssymb}
\setcounter{secnumdepth}{-\maxdimen} % remove section numbering
\usepackage{iftex}
\ifPDFTeX
  \usepackage[T1]{fontenc}
  \usepackage[utf8]{inputenc}
  \usepackage{textcomp} % provide euro and other symbols
\else % if luatex or xetex
  \usepackage{unicode-math} % this also loads fontspec
  \defaultfontfeatures{Scale=MatchLowercase}
  \defaultfontfeatures[\rmfamily]{Ligatures=TeX,Scale=1}
\fi
\usepackage{lmodern}
\ifPDFTeX\else
  % xetex/luatex font selection
  \setmainfont[]{Liberation Serif}
  \setmonofont[]{Liberation Mono}
\fi
% Use upquote if available, for straight quotes in verbatim environments
\IfFileExists{upquote.sty}{\usepackage{upquote}}{}
\IfFileExists{microtype.sty}{% use microtype if available
  \usepackage[]{microtype}
  \UseMicrotypeSet[protrusion]{basicmath} % disable protrusion for tt fonts
}{}
\makeatletter
\@ifundefined{KOMAClassName}{% if non-KOMA class
  \IfFileExists{parskip.sty}{%
    \usepackage{parskip}
  }{% else
    \setlength{\parindent}{0pt}
    \setlength{\parskip}{6pt plus 2pt minus 1pt}}
}{% if KOMA class
  \KOMAoptions{parskip=half}}
\makeatother
\usepackage{longtable,booktabs,array}
\usepackage{caption}
\captionsetup[table]{skip=6pt}
\newcounter{none} % for unnumbered tables
\usepackage{calc} % for calculating minipage widths
% Correct order of tables after \paragraph or \subparagraph
\usepackage{etoolbox}
\makeatletter
\patchcmd\longtable{\par}{\if@noskipsec\mbox{}\fi\par}{}{}
\makeatother
% Allow footnotes in longtable head/foot
\IfFileExists{footnotehyper.sty}{\usepackage{footnotehyper}}{\usepackage{footnote}}
\makesavenoteenv{longtable}
\ifLuaTeX
\usepackage[bidi=basic,shorthands=off]{babel}
\else
\usepackage[bidi=default,shorthands=off]{babel}
\fi
\ifPDFTeX
\else
\babelfont{rm}[]{Liberation Serif}
\fi
\ifLuaTeX
  \usepackage{selnolig} % disable illegal ligatures
\fi
\setlength{\emergencystretch}{3em} % prevent overfull lines
\providecommand{\tightlist}{%
  \setlength{\itemsep}{0pt}\setlength{\parskip}{0pt}}
\usepackage{xcolor}
\usepackage{booktabs}
\usepackage{fancyhdr}
\usepackage{etoolbox}
\setlength{\headheight}{14pt}
\pagestyle{fancy}
\fancyhf{}
\fancyhead[L]{\leftmark}
\fancyhead[R]{\rightmark}
\fancyfoot[C]{\thepage}
\renewcommand{\headrulewidth}{0.4pt}
\makeatletter
\patchcmd{\tableofcontents}
  {\@mkboth{\MakeUppercase\contentsname}{\MakeUppercase\contentsname}}
  {\@mkboth{\MakeUppercase\contentsname}{}}
  {}{\PackageWarning{validation-report}{Could not patch TOC page marks}}
\makeatother
\usepackage{bookmark}
\IfFileExists{xurl.sty}{\usepackage{xurl}}{} % add URL line breaks if available
\urlstyle{same}
% fallback for those not using the hyperref driver hyperxmp:
\makeatletter
\@ifundefined{xmpquote}{\newcommand{\xmpquote}[1]{#1}}{}
\makeatother
\hypersetup{
  pdftitle={Kwalitatieve data Neuro-symbolic model: Gold \& Model output},
  pdflang={nl},
  colorlinks=true,
  linkcolor={Maroon},
  filecolor={Maroon},
  citecolor={Blue},
  urlcolor={Blue},
  pdfcreator={LaTeX via pandoc}}

\title{Kwalitatieve data Neuro-symbolic model: Gold \& Model output}
\author{}
\date{}

\begin{document}
\maketitle

\renewcommand*\contentsname{Inhoudsopgave}
{
\setcounter{tocdepth}{3}
\tableofcontents
}
\newpage

\markboth{JOR 2014/78}{Gold standard}

\section{JOR 2014/78}\label{jor-201478}

\texttt{ECLI:NL:GHARL:2013:822}, uitspraakdatum 12-02-2013

\subsection{Gold}\label{gold}

\begin{center}
\small
\begin{tabular}{@{}p{8cm}p{2.6cm}@{}}
\toprule
Variabele & Waarde \\
\midrule
Is rechtshandeling & \texttt{true} \\
Datum rechtshandeling & \texttt{2005-12-30} \\
Verricht door de schuldenaar & \texttt{true} \\
Om niet & \texttt{true} \\
Datum faillietverklaring & \texttt{2009-06-23} \\
Onverplicht & \texttt{true} \\
Benadeling schuldeisers & \texttt{true} \\
Wetenschap schuldenaar aangetoond & \texttt{true} \\
Wetenschap wederpartij aangetoond & \texttt{null} \\
Vermoeden wetenschap schuldenaar tegenbewijs & \texttt{null} \\
Vermoeden wetenschap wederpartij tegenbewijs & \texttt{null} \\
Tegen onderwaarde & \texttt{null} \\
Voldoening of zekerheid niet opeisbare schuld & \texttt{null} \\
Gelieerde partij & \texttt{null} \\
Reeds verplicht voor aanvang termijn & \texttt{null} \\
\textbf{Vernietigbaar} & \textbf{\texttt{true}} \\
\bottomrule
\end{tabular}
\end{center}

\subsubsection{Uitspraak}\label{uitspraak}

(\ldots; red.)

\paragraph{3. De vaststaande feiten}\label{de-vaststaande-feiten}

3.1. Tussen partijen staan, als enerzijds gesteld en anderzijds niet althans onvoldoende gemotiveerd betwist, de volgende feiten vast.

3.2. Van Doorn Beheer was aandeelhouder van de Vennootschap. Op 9 respectievelijk 28 december 2005 heeft Van Doorn Beheer de aandelen in de Vennootschap verkocht en geleverd aan Vewa Beheer B.V. (hierna: Vewa).

3.3. In de koopovereenkomst tussen Van Doorn Beheer en Vewa is onder meer het volgende bepaald:

``(\ldots)

overwegende dat,

\begin{enumerate}
\def\labelenumi{\alph{enumi}.}
\item
  Verkoper houder is van alle geplaatste aandelen in het kapitaal van F. van Doorn B.V. (\ldots) (\ldots{} hierna te noemen `de Vennootschap');
\item
  Verkoper een vordering in rekening-courantheeft op de Vennootschap;
\item
  (\ldots)
\item
  Partijen (\ldots) tot overeenstemming zijn gekomen over de verkrijging door Koper van de door Verkoper gehouden aandelen in het kapitaal van de Vennootschap en van de vordering van Verkoper op de vennootschap;
\item
  de algemene vergadering van aandeelhouders van Verkoper heeft ingestemd met deze verkoop, en de directie van de Vennootschap zich akkoord heeft verklaard met de overdracht aan Koper van haar schuld aan Verkoper, (\ldots),
\end{enumerate}

verklaren te zijn overeengekomen als volgt:

Artikel 1 -- Object van koop en verkoop

Koper koopt van Verkoper:

\begin{enumerate}
\def\labelenumi{\arabic{enumi}.}
\item
  de volle en onbezwaarde eigendom van de aandelen (\ldots) in het aandelenkapitaal van de Vennootschap (\ldots).
\item
  de vordering in rekening-courant van de Verkoper op de Vennootschap.
\end{enumerate}

Artikel 2 -- Koopprijs

De koopprijs van de door Koper overgenomen aandelen bedraagt € 1.

De koopprijs van de door Koper overgenomen vordering in rekening-courant op de Vennootschap bedraagt € 349.999. (\ldots)

(\ldots)

Artikel 6 -- Betaling van de koopprijs

Koper zal uiterlijk op de Leveringsdatum de koopprijs van de aandelen in contanten voldoen (\ldots).

Koper zal uiterlijk op de Leveringsdatum de koopprijs van de vordering in rekening-courant (onder aftrek van het hierna genoemde door Verkoper aan Koper te lenen bedrag) voldoen (\ldots).

Een bedrag ter grootte van € 150.000 zal door Verkoper aan Koper worden geleend. (\ldots) De belangrijkste bepalingen in de leningsovereenkomst zijn:

(\ldots)

-- zekerheid in de vorm van een tweede pandrecht op de debiteuren en de inventaris van de Vennootschap. (\ldots)''

3.4. In de overeenkomst van geldlening tussen Van Doorn Beheer als schuldeiser en Vewa als schuldenaar d.d. 30 december 2005 is onder meer het volgende bepaald:

``(\ldots)

In aanmerking nemende dat:

-- De schuldenaar van schuldeiser de aandelen in F. van Doorn BV en de vordering op F. van Doorn BV heeft gekocht voor een bedrag van € 349.999

-- De schuldenaar en schuldeiser een deel van de koopsom groot € 150.000 schuldig is gebleven, welk bedrag bij deze door schuldeiser aan schuldenaar als geldlening zal worden verstrekt

(\ldots)

Verklaren vast te leggen en te zijn overeengekomen als volgt:

(\ldots)

\begin{enumerate}
\def\labelenumi{\arabic{enumi}.}
\setcounter{enumi}{7}
\tightlist
\item
  De schuldenaar verpandt de volgende zekerheden aan schuldeiser die schuldeiser bij deze als pand aanvaardt
\end{enumerate}

\begin{enumerate}
\def\labelenumi{\alph{enumi})}
\item
  De gehele bedrijfsinventaris behorende tot de onderneming (\ldots)
\item
  Alle vorderingen (\ldots) van de schuldenaar die thans reeds bestaan of zullen voortvloeien uit een thans bestaande rechtsverhouding ter zake van de uitoefening door schuldenaar van zijn bedrijf (\ldots)''
\end{enumerate}

3.5. In de overeenkomst van verpanding tussen Van Doorn Beheer als pandhouder en de Vennootschap als pandgever d.d. 30 december 2005 is onder meer het volgende bepaald:

``(\ldots)

in aanmerking nemende

dat F. van Doorn Beheer BV aan Vewa Beheer BV een geldbedrag heeft geleend groot € 150.000.

dat F. van Doorn BV zich voor dit bedrag borg heeft gesteld ten behoeve van Vewa Beheer BV

dat F. van Doorn BV in dit kader in pand aan F. van Doorn Beheer BV verstrekt haar debiteuren en inventaris onder de navolgende voorwaarden.

verklaren het volgende te zijn overeengekomen:

De pandgever verpandt de volgende zekerheden aan pandhouder die pandhouder bij deze in pand aanvaardt

\begin{enumerate}
\def\labelenumi{\alph{enumi})}
\item
  De gehele bedrijfsinventaris behorende tot de onderneming (\ldots)
\item
  Alle vorderingen (\ldots) van de pandhouder die thans reeds bestaan of zullen voortvloeien uit een thans bestaande rechtsverhouding ter zake van de uitoefening door pandgever van zijn bedrijf (\ldots)''
\end{enumerate}

3.6. In de periode oktober 2007 tot en met januari 2009 heeft de Vennootschap, naast huurpenningen, € 49.500 aan Van Doorn Beheer betaald.

3.7. De Vennootschap is op 23 juni 2009 in staat van faillissement verklaard.

3.8. Na het faillissement van de Vennootschap heeft de curator de bedrijfsinventaris verkocht en is Van Doorn Beheer tot inning van de openstaande vorderingen van de Vennootschap overgegaan.

\paragraph{4. De motivering van de beslissing in hoger beroep}\label{de-motivering-van-de-beslissing-in-hoger-beroep}

4.1. De curator betwist dat de activa van de Vennootschap rechtsgeldig aan Van Doorn Beheer zijn verpand en dat door de Vennootschap rechtsgeldig een borgtocht is verstrekt. Hij heeft in dat verband -- voor zover in hoger beroep nog van belang en zakelijk samengevat -- verklaringen van recht gevorderd dat de Vennootschap geen verplichtingen heeft uit hoofde van een borgtocht voor de verplichtingen van Vewa jegens Van Doorn Beheer, althans dat die borgtocht (ver)nietig(baar) is, en dat Van Doorn Beheer geen pandrecht toekomt op inventaris of vorderingen van de Vennootschap, althans dat deze pandrechten (ver)nietig(baar) zijn. Tevens heeft de curator gevorderd dat Van Doorn Beheer zal worden veroordeeld tot het doen van rekening en verantwoording over de door haar geïnde debiteuren van de Vennootschap en tot afdracht van het daaruit blijkende saldo. Voorts heeft de curator zich op het standpunt gesteld dat de in rechtsoverweging 3.6 bedoelde betalingen onverschuldigd zijn gedaan en heeft hij terugbetaling daarvan gevorderd.

4.2. Van Doorn Beheer heeft zich daartegenover op het standpunt gesteld dat de borgtocht en pandrechten rechtsgeldig zijn verstrekt en dat de curator, door zich op een ander standpunt te stellen, onrechtmatig jegens Van Doorn Beheer handelt. Van Doorn Beheer heeft in dat kader in reconventie -- zakelijk samengevat -- verklaringen van recht gevorderd alsmede veroordeling van de curator tot betaling van een bedrag van € 160.000, verhoogd met rente, uit hoofde van de borgstelling en tot schadevergoeding nader op te maken bij staat en te vereffenen volgens de wet.

4.3. De rechtbank heeft geoordeeld dat de door de Vennootschap verstrekte borgtocht en pandrechten nietig zijn wegens het bepaalde in artikel 2:207c van het Burgerlijk Wetboek (BW), de vorderingen van de curator in conventie toegewezen en de vorderingen van Van Doorn Beheer in reconventie afgewezen. Tegen die beslissingen (met uitzondering van de reeds gedane rekening en verantwoording) richten zich de grieven in het hoger beroep. Evenals de rechtbank, zal het hof bij de beoordeling er veronderstellenderwijze van uit gaan dat de borgtocht is aangegaan en dat de pandrechten zijn verstrekt.

4.4. Met haar eerste grief richt Van Doorn Beheer zich tegen het oordeel van de rechtbank dat de borgtocht en pandrechten nietig zijn wegens het bepaalde in artikel 2:207c BW. Artikel 2:207c BW is met de inwerkingtreding van de Wet vereenvoudiging en flexibilisering bv-recht (Stb. 2012, 299 en 301) per 1 oktober 2012 vervallen. Dit brengt op grond van het in artikel V.1 lid 1 Invoeringswet vereenvoudiging en flexibilisering bv-recht (Stb. 2012, 300 en 301) van overeenkomstige toepassing verklaarde artikel 81 van de Overgangswet nieuw Burgerlijk Wetboek met zich dat rechtshandelingen die op grond van artikel 2:207c BW nietig zouden zijn in beginsel met terugwerkende kracht geldig en onaantastbaar worden. Dat is evenwel anders indien, zoals in het onderhavige geval, een onmiddellijk belanghebbende zich reeds op de nietigheid van de rechtshandeling had beroepen. Toch komt het hof niet aan een inhoudelijke beoordeling van de grief toe. Mocht deze grief namelijk slagen, dan komt alsnog het beroep van de curator op vernietiging van de borgtocht en de pandrechten op grond van artikel 42 van de Faillissementwet (Fw) aan de orde. Het hof ziet aanleiding dat beroep eerst te behandelen.

4.5. Voor een succesvol beroep op artikel 42 Fw dient aan drie voorwaarden te zijn voldaan. Er dient -- kort gezegd -- sprake te zijn van een onverplichte rechtshandeling (i), een daardoor veroorzaakte benadeling van schuldeisers (ii) en wetenschap van benadeling (iii).

4.6. Zonder nadere redengeving, die ontbreekt, valt niet in te zien op welke grond de Vennootschap gehouden was om jegens Van Doorn Beheer zekerheden te verstrekken voor een schuld van Vewa aan Van Doorn Beheer. De omstandigheid dat, zoals Van Doorn Beheer in dit verband heeft aangevoerd, de deal tussen Van Doorn Beheer en Vewa geen doorgang zou vinden zonder dat Van Doorn Beheer zekerheid zou krijgen voor de door haar aan Vewa te verstrekken lening, levert nog geen verplichting op zijdens de Vennootschap. Dat geldt ook voor de in dit verband door Vewa jegens Van Doorn Beheer aangegane verplichting tot het verstrekken van zekerheden. Die verplichting regardeert de Vennootschap niet. Overigens is niets gesteld of gebleken waaruit zou kunnen volgen dat het verstrekken van de borgtocht en de pandrechten desalniettemin berustten op een uit de wet of overeenkomst voortvloeiende verplichting van de Vennootschap. Daarmee is aan het eerste vereiste van artikel 42 Fw voldaan.

4.7. De vraag of de verstrekking van de borgtocht en pandrechten door de Vennootschap benadeling van haar schuldeisers tot gevolg heeft, dient te worden beoordeeld naar het moment waarop op het door de curator gedane beroep op de vernietiging wordt beslist en moet worden beantwoord door de hypothetische situatie waarin de schuldeisers zouden hebben verkeerd zonder de gewraakte rechtshandelingen te vergelijken met de situatie waarin zij feitelijk verkeren als die rechtshandelingen onaangetast blijven. Van een dergelijke benadeling is sprake. Door de borgtocht en de daarmee samenhangende pandrechten heeft Van Doorn Beheer zich een sterkere verhaalspositie verschaft jegens de Vennootschap, terwijl de vordering van Van Doorn Beheer op de Vennootschap, al is zij gecedeerd, is blijven bestaan. Dat deze ``feitelijk'' zou zijn kwijtgescholden, zoals Van Doorn Beheer heeft gesteld, laat onverlet dat de vordering op de Vennootschap niet is verminderd. Er is geen verandering ontstaan in de totale schuldenlast van de Vennootschap terwijl door de borgstelling en de verpanding wel minder verhaalsobjecten voor de overige schuldeisers beschikbaar zijn. Daarmee is ook aan het tweede vereiste van artikel 42 Fw voldaan.

4.8. Voor de beoordeling van het derde vereiste van artikel 42 Fw dient eerst te worden beoordeeld of sprake is van rechtshandelingen om niet; in dat geval is geen wetenschap van benadeling aan de zijde van Van Doorn Beheer vereist. Tegen de stelling van de curator dat de verstrekking door de Vennootschap van de borgtocht en pandrechten om niet plaatsvond, heeft Van Doorn Beheer geen verweer gevoerd. Voor zover haar verweer dat de Vennootschap heeft geprofiteerd doordat de vordering van Van Doorn Beheer feitelijk deels is kwijtgescholden mede in dit verband moet worden betrokken, faalt het. Uit de omstandigheid dat Van Doorn Beheer haar vordering op de vennootschap (deels) boekhoudkundig afboekte, kan niet worden geconcludeerd dat de rechtshandelingen anders dan om niet hebben plaatsgevonden. De schuld van de Vennootschap is immers geheel blijven bestaan. Dat de Vennootschap wist of behoorde te weten dat door de verstrekking van de borgtocht en de pandrechten haar schuldeisers zouden worden benadeeld, is door Van Doorn Beheer onvoldoende gemotiveerd betwist, zodat dit vaststaat. Daarmee staat tevens vast dat ook aan het derde vereiste van artikel 42 Fw is voldaan.

4.9. Dit betekent dat de curator zich op goede gronden op vernietiging van de borgtocht en de pandrechten heeft beroepen. Waar de curator reeds zelf vernietigingsverklaringen heeft uitgebracht, is daarmee op grond van het bepaalde in artikel 3:53 lid 1 BW de nietigheid van de rechtshandelingen reeds ingetreden (een rechtshandeling wordt immers door de vernietiging met terugwerkende kracht tot een nietige rechtshandeling) en kunnen de daartoe strekkende verklaringen van recht worden bekrachtigd.

4.10. De borgtocht en de pandrechten strekten ertoe Van Doorn Beheer een bepaalde mate van zekerheid te verschaffen ten aanzien van de door haar aan Vewa te verstrekken lening. Het hof vermag niet, ook niet met hetgeen door Van Doorn Beheer ter zake van de toepassing van artikel 3:42 BW is opgemerkt, in te zien van welke, niet-paulianeuze, rechtshandelingen de strekking in voldoende mate beantwoordt aan die van de borgtocht en de pandrechten dat kan worden aangenomen dat die rechtshandelingen zouden zijn verricht indien van de borgtocht en pandrechten wegens hun vernietigbaarheid was afgezien. Voor conversie op grond van de wet is immers vereist dat de vervangende geldige rechtshandeling naar haar strekking met de nietige rechtshandeling vergelijkbaar dient te zijn en partijen dient te verschaffen hetgeen zij blijkens hun nietige rechtshandeling (in bedrijfseconomisch opzicht) wilden bereiken. Een toelaatbaar alternatief voor naar haar strekking met artikel 42 Fw strijdige borgtocht en pandrechten laat zich in dit geval bezwaarlijk indenken. Aan de vereisten voor conversie is dan ook niet voldaan. Daarmee behoeft de eerste grief verder geen behandeling meer.

4.11. Van Doorn Beheer heeft, daarop ziet haar tweede grief, nog aangevoerd dat de door de Vennootschap aan Van Doorn Beheer gedane betalingen zijn verricht namens Vewa en niet onverschuldigd zijn gedaan. Door de curator is evenwel betwist dat tussen de Vennootschap en Vewa sprake is van een rechtsverhouding op grond waarvan kan worden aangenomen dat de door de Vennootschap gedane betalingen niet in haar hoedanigheid van borg maar rechtstreeks namens (of: ten behoeve van) Vewa zijn verricht. De enkele stelling dat de betalingen zijn geboekt in rekening-courant tussen Vewa en de Vennootschap is, tegenover deze betwisting, onvoldoende om te kunnen aannemen dat de Vennootschap namens Vewa heeft betaald. Ook indien de Vennootschap uit hoofde van de (vernietigde) borgtocht heeft betaald, had zij immers regres op Vewa en laat de boeking in rekening-courant zich evenzeer daardoor verklaren. Van Doorn Beheer heeft zich onder II.1 van haar memorie van grieven ook zelf primair op het standpunt gesteld dat ``de vennootschap als borg ervan uit mocht gaan dat zij op dat moment werd uitgewonnen''. Voor zover Van Doorn Beheer desalniettemin bedoeld heeft te stellen dat de Vennootschap rechtstreeks namens Vewa heeft betaald, en mitsdien niet onverschuldigd, heeft zij dat verweer onvoldoende onderbouwd. Dit verweer wordt daarom gepasseerd.

4.12. Uit het voorgaande volgt dat de pandrechten en borgtocht niet geldig zijn. De rechtbank heeft de vorderingen in reconventie op die grond afgewezen. Onder II.2 heeft Van Doorn Beheer hiertegen gegriefd, maar uitsluitend door te stellen dat zij zich wel op rechtsgeldige pandrechten en borgstelling kan beroepen. Die grief deelt in het lot van de eerste grief.

4.13. De derde grief richt zich tegen het passeren van het bewijsaanbod door de rechtbank. Die grief faalt reeds op de grond dat, indien Van Doorn Beheer in hoger beroep alsnog tot bewijslevering toegelaten had willen worden, zij daartoe een voldoende concreet en specifiek aanbod had moeten doen. Dat heeft zij evenwel nagelaten. Overigens geldt ten aanzien van het in eerste aanleg gedane bewijsaanbod dat dit, voor zover geconcretiseerd en gespecificeerd, geen betrekking heeft op feiten die, indien bewezen, tot andere conclusies aanleiding zouden kunnen geven. Het bewijsaanbod is dus onvoldoende specifiek en terzake dienend en daarom terecht gepasseerd. Ook op die grond mist de grief doel.

4.14. De vierde grief mist zelfstandig belang en deelt in het lot van de andere grieven.

4.15. Uit het voorgaande volgt dat het hoger beroep faalt. De bestreden vonnissen zullen worden bekrachtigd. Van Doorn Beheer zal als de in het ongelijk gestelde partij worden veroordeeld in de kosten van het hoger beroep.

\paragraph{5. De beslissing}\label{de-beslissing}

Het hof, recht doende in hoger beroep:

bekrachtigt de tussen partijen gewezen vonnissen van de rechtbank Arnhem van 1 juni 2011 en 14 december 2011;

veroordeelt Van Doorn Beheer in de kosten van het hoger beroep (\ldots; red.).

\newpage

\markboth{JOR 2014/78}{mistral/mistral-medium-3.5-128b:fp8 (k = 0)}

\subsection{k = 0}\label{k-0}

\subsubsection{Mistral}\label{mistral}

\texttt{mistral/mistral-medium-3.5-128b:fp8}

\begin{center}
\small
\begin{tabular}{@{}p{7.4cm}p{2.6cm}p{2.6cm}@{}}
\toprule
Variabele & Waarde (model) & Waarde (gold) \\
\midrule
Is rechtshandeling & \texttt{true} & \texttt{true} \\
Datum rechtshandeling & \texttt{2005-12-30} & \texttt{2005-12-30} \\
Verricht door de schuldenaar & \texttt{true} & \texttt{true} \\
Om niet & \textcolor{gray}{\texttt{false}} & \texttt{true} \\
Datum faillietverklaring & \texttt{2009-06-23} & \texttt{2009-06-23} \\
Onverplicht & \texttt{true} & \texttt{true} \\
Benadeling schuldeisers & \texttt{true} & \texttt{true} \\
Wetenschap schuldenaar aangetoond & \textcolor{gray}{\texttt{false}} & \texttt{true} \\
Wetenschap wederpartij aangetoond & \texttt{false} & \texttt{null} \\
Vermoeden wetenschap schuldenaar tegenbewijs & \texttt{null} & \texttt{null} \\
Vermoeden wetenschap wederpartij tegenbewijs & \texttt{null} & \texttt{null} \\
Tegen onderwaarde & \texttt{null} & \texttt{null} \\
Voldoening of zekerheid niet opeisbare schuld & \texttt{null} & \texttt{null} \\
Gelieerde partij & \textcolor{gray}{\texttt{true}} & \texttt{null} \\
Reeds verplicht voor aanvang termijn & \texttt{null} & \texttt{null} \\
\textbf{Vernietigbaar} & \textbf{\textcolor{gray}{\texttt{false}}} & \textbf{\texttt{true}} \\
\bottomrule
\end{tabular}
\end{center}

\paragraph{Is rechtshandeling - overwegingen}\label{is-rechtshandeling---overwegingen}

De verpanding van de bedrijfsinventaris en vorderingen door de Vennootschap aan Van Doorn Beheer (3.5) is een goederenrechtelijke handeling, gericht op een rechtsgevolg (vestiging van pandrecht).

\paragraph{Verricht door de schuldenaar - overwegingen}\label{verricht-door-de-schuldenaar---overwegingen}

De Vennootschap (de latere gefailleerde) heeft zelf de verpanding verricht (3.5).

\paragraph{Om niet - overwegingen}\label{om-niet---overwegingen}

De verpanding is verricht in het kader van een lening van € 150.000 door Van Doorn Beheer aan Vewa (3.4 en 3.5), waarbij de Vennootschap als pandgever optrad. Er is sprake van een wederprestatie (lening vs.~pandrecht), dus de handeling is `anders dan om niet'.

\paragraph{Onverplicht - overwegingen}\label{onverplicht---overwegingen}

De verpanding van de bedrijfsinventaris en vorderingen door de Vennootschap aan Van Doorn Beheer is onverplicht verricht in de zin van artikel 42 Fw. Uit de casus blijkt dat de Vennootschap zich borg heeft gesteld voor de lening van € 150.000 die Van Doorn Beheer aan Vewa Beheer heeft verstrekt. Deze borgtocht en de daarmee samenhangende verpanding zijn niet gebaseerd op een bestaande rechtsplicht van de Vennootschap. Er was geen wettelijke of contractuele verplichting voor de Vennootschap om deze zekerheid te stellen. De borgtocht en verpanding zijn voortgevloeid uit de koopovereenkomst tussen Van Doorn Beheer en Vewa, waarbij de Vennootschap geen partij was. De Vennootschap had derhalve geen rechtsplicht om deze zekerheid te verlenen. Het feit dat de directie van de Vennootschap akkoord is gegaan met de overdracht van de schuld aan Vewa, betekent niet dat er een rechtsplicht bestond om zekerheid te stellen voor de lening aan Vewa. De verpanding is daarom onverplicht verricht.

\paragraph{Benadeling schuldeisers - overwegingen}\label{benadeling-schuldeisers---overwegingen}

De rechtshandeling betreft de verkoop van aandelen en een vordering in rekening-courant door Van Doorn Beheer aan Vewa, inclusief de daarmee samenhangende lening van € 150.000 en de verpanding van de bedrijfsinventaris en debiteuren van de Vennootschap als zekerheid voor deze lening.

\begin{enumerate}
\def\labelenumi{\arabic{enumi}.}
\tightlist
\item
  \textbf{Vermogensvergelijking}:

  \begin{itemize}
  \tightlist
  \item
    \textbf{Hypothetische situatie zonder de rechtshandeling}: De Vennootschap had geen verpanding van haar activa (inventaris en debiteuren) aan Van Doorn Beheer en geen borgtocht voor Vewa's schuld. De activa van de Vennootschap (inclusief de debiteuren en inventaris) zouden volledig ten bate van alle schuldeisers in het faillissement zijn gekomen.
  \item
    \textbf{Feitelijke situatie met de rechtshandeling}: De activa van de Vennootschap zijn belast met een pandrecht ten gunste van Van Doorn Beheer, en de Vennootschap heeft zich borg gesteld voor Vewa's schuld. Hierdoor kunnen de schuldeisers zich niet of beperkt verhalen op deze activa, terwijl Van Doorn Beheer als pandhouder voorrang heeft. Bovendien heeft Van Doorn Beheer na het faillissement de debiteuren geïnd, wat ten koste gaat van de boedel.
  \end{itemize}
\item
  \textbf{Benadeling}:

  \begin{itemize}
  \tightlist
  \item
    De verpanding en borgtocht leiden tot een \textbf{vermindering van het actief} voor de overige schuldeisers, omdat de activa van de Vennootschap niet meer vrij beschikbaar zijn voor verhaal.
  \item
    Er is sprake van \textbf{verstoring van de rangorde}: Van Doorn Beheer verkrijgt door het pandrecht een preferente positie ten opzichte van andere schuldeisers, terwijl deze handeling niet in het kader van een normale zakelijke transactie tegen marktconforme voorwaarden lijkt te zijn gesloten (de koopprijs van de aandelen was € 1, en de vordering werd verkocht voor € 349.999, terwijl direct een lening van € 150.000 werd verstrekt met zekerheid op de activa van de Vennootschap).
  \item
    De betalingen van € 49.500 door de Vennootschap aan Van Doorn Beheer (3.6) zijn eveneens benadelend, omdat deze gelden niet ten bate van de boedel zijn gekomen.
  \end{itemize}
\item
  \textbf{Kennis van benadeling}:

  \begin{itemize}
  \tightlist
  \item
    Van Doorn Beheer en Vewa wisten of behoorden te weten dat de verpanding en borgtocht tot benadeling van de schuldeisers van de Vennootschap zouden leiden, gelet op de constructie (verkoop van aandelen en vordering tegen een symbolische prijs, directe lening met zekerheid op de activa van de Vennootschap) en de latere betalingen door de Vennootschap aan Van Doorn Beheer.
  \end{itemize}
\item
  \textbf{Conclusie}: De schuldeisers zijn benadeeld doordat de activa van de Vennootschap zijn belast met een pandrecht ten gunste van Van Doorn Beheer, de rangorde is verstoord, en gelden zijn afgevloeid naar Van Doorn Beheer die anders ten bate van de boedel zouden zijn gekomen. Dit voldoet aan de voorwaarden van artikel 42 Fw.
\end{enumerate}

\paragraph{Wetenschap - overwegingen}\label{wetenschap---overwegingen}

De Vennootschap (schuldenaar) heeft op 30 december 2005 een pandrecht verstrekt aan Van Doorn Beheer (wederpartij) op haar inventaris en vorderingen. Van Doorn Beheer was aandeelhouder van de Vennootschap en valt onder art. 43 lid 1 sub 4° Fw (gelieerde partij: aandeelhouder). De rechtshandeling (verpanding) is verricht binnen een jaar voor de faillietverklaring (23 juni 2009) en de Vennootschap had zich niet reeds voor de aanvang van die termijn daartoe verplicht. Derhalve wordt de wetenschap van benadeling aan beide zijden (schuldenaar en wederpartij) vermoed (art. 43 lid 1 Fw). Er is geen sprake van tegenbewijs in de casus, dus de vermoedens blijven staan. De curator hoeft derhalve niet rechtstreeks de wetenschap aan te tonen.

\newpage

\markboth{JOR 2014/78}{gpt-5.5 (k = 0)}

\subsubsection{GPT}\label{gpt}

\texttt{gpt-5.5}

\begin{center}
\small
\begin{tabular}{@{}p{7.4cm}p{2.6cm}p{2.6cm}@{}}
\toprule
Variabele & Waarde (model) & Waarde (gold) \\
\midrule
Is rechtshandeling & \texttt{true} & \texttt{true} \\
Datum rechtshandeling & \texttt{2005-12-30} & \texttt{2005-12-30} \\
Verricht door de schuldenaar & \texttt{true} & \texttt{true} \\
Om niet & \texttt{true} & \texttt{true} \\
Datum faillietverklaring & \texttt{2009-06-23} & \texttt{2009-06-23} \\
Onverplicht & \texttt{true} & \texttt{true} \\
Benadeling schuldeisers & \texttt{true} & \texttt{true} \\
Wetenschap schuldenaar aangetoond & \textcolor{gray}{\texttt{false}} & \texttt{true} \\
Wetenschap wederpartij aangetoond & \texttt{false} & \texttt{null} \\
Vermoeden wetenschap schuldenaar tegenbewijs & \texttt{null} & \texttt{null} \\
Vermoeden wetenschap wederpartij tegenbewijs & \texttt{null} & \texttt{null} \\
Tegen onderwaarde & \texttt{null} & \texttt{null} \\
Voldoening of zekerheid niet opeisbare schuld & \texttt{null} & \texttt{null} \\
Gelieerde partij & \texttt{null} & \texttt{null} \\
Reeds verplicht voor aanvang termijn & \texttt{null} & \texttt{null} \\
\textbf{Vernietigbaar} & \textbf{\textcolor{gray}{\texttt{false}}} & \textbf{\texttt{true}} \\
\bottomrule
\end{tabular}
\end{center}

\paragraph{Is rechtshandeling - overwegingen}\label{is-rechtshandeling---overwegingen-1}

De gewraakte handelingen betreffen het stellen van borgtocht door de Vennootschap voor de schuld van Vewa en het vestigen van pandrechten op inventaris en vorderingen ten gunste van Van Doorn Beheer. Dit zijn op rechtsgevolg gerichte verklaringen, waaronder goederenrechtelijke zekerheidshandelingen.

\paragraph{Verricht door de schuldenaar - overwegingen}\label{verricht-door-de-schuldenaar---overwegingen-1}

De failliete Vennootschap, F. van Doorn B.V., heeft zich volgens de feiten borg gesteld en trad in de verpandingsovereenkomst op als pandgever tegenover Van Doorn Beheer.

\paragraph{Om niet - overwegingen}\label{om-niet---overwegingen-1}

Uit de feiten blijkt niet dat de Vennootschap voor het stellen van de borgtocht en het verstrekken van pandrechten een reële wederprestatie heeft bedongen. De zekerheid werd verstrekt voor een lening van Van Doorn Beheer aan Vewa.

\paragraph{Onverplicht - overwegingen}\label{onverplicht---overwegingen-1}

Voor de door de Vennootschap verstrekte borgtocht en verpanding van haar inventaris en vorderingen bestond geen in rechte afdwingbare rechtsplicht uit wet of overeenkomst. De koopovereenkomst en de geldleningsovereenkomst zijn gesloten tussen Van Doorn Beheer en Vewa; de Vennootschap was daarbij niet als schuldenaar van de zekerheidsverplichting partij. Dat in de koopovereenkomst is vermeld dat zekerheid zou worden verstrekt in de vorm van een tweede pandrecht op debiteuren en inventaris van de Vennootschap, schept op zichzelf geen afdwingbare verplichting voor de Vennootschap, nu zij die contractuele verplichting niet als partij is aangegaan. Ook de omstandigheid dat de directie van de Vennootschap akkoord was met de overdracht van de schuld/vordering aan Vewa levert geen verplichting op om zich borg te stellen of pandrechten te vestigen voor de schuld van Vewa aan Van Doorn Beheer. De borgtocht en pandrechten dienden ter zekerheid van een schuld van Vewa, niet van een eigen reeds afdwingbare schuld van de Vennootschap. De rechtshandeling is daarom onverplicht verricht in de zin van artikel 42 Fw.

\paragraph{Benadeling schuldeisers - overwegingen}\label{benadeling-schuldeisers---overwegingen-1}

De aangevochten rechtshandelingen bestaan erin dat de Vennootschap zich ten behoeve van de schuld van Vewa aan Van Doorn Beheer borg heeft gesteld en haar inventaris en vorderingen aan Van Doorn Beheer heeft verpand. Bij de vereiste vermogensvergelijking geldt dat, zonder deze borgtocht en pandrechten, Van Doorn Beheer geen separatistenpositie en geen verhaalsrecht op de activa van de Vennootschap zou hebben gehad voor de schuld van Vewa. De inventaris en debiteurenopbrengsten zouden dan in beginsel beschikbaar zijn geweest voor de gezamenlijke schuldeisers van de Vennootschap. Als de rechtshandelingen in stand blijven, kan Van Doorn Beheer zich daarentegen met voorrang verhalen op die activa en heeft zij na faillissement ook debiteuren van de Vennootschap geïnd. Daarnaast heeft de Vennootschap vóór faillissement € 49.500 aan Van Doorn Beheer betaald, hetgeen het actief eveneens heeft verminderd voor zover die betalingen samenhangen met de gewraakte rechtsverhouding. De rechtshandelingen leiden daarmee tot vermindering van het voor de gezamenlijke schuldeisers beschikbare actief en in elk geval tot verstoring van de rangorde doordat Van Doorn Beheer een zekerheidspositie verkrijgt waar zij zonder de rechtshandelingen geen aanspraak op had. Uit de casus blijkt niet dat de betrokken activa ook zonder deze pandrechten volledig buiten het verhaal van de schuldeisers zouden zijn gebleven, bijvoorbeeld door uitwinning door een hoger gerangschikte zekerheidsgerechtigde. Eventuele voordelen voor de Vennootschap zijn evenmin concreet: de lening werd verstrekt aan Vewa en de aandelenoverdracht kwam niet als zodanig aan het vermogen van de Vennootschap ten goede. Daarom zijn de schuldeisers in de zin van artikel 42 Fw benadeeld.

\paragraph{Wetenschap - overwegingen}\label{wetenschap---overwegingen-1}

Uit de casus volgt onvoldoende dat de Vennootschap ten tijde van het verstrekken van de borgtocht en pandrechten wist of behoorde te weten dat benadeling van schuldeisers het gevolg zou zijn. Er zijn geen concrete feiten gesteld waaruit blijkt dat op dat moment het faillissement en een tekort daarin met een redelijke mate van waarschijnlijkheid voorzienbaar waren; het faillissement is pas op 23 juni 2009 uitgesproken. Ook ten aanzien van Van Doorn Beheer is die wetenschap niet rechtstreeks uit de feiten af te leiden. De bewijsvermoedens van art. 43 Fw zijn op basis van de gegeven feiten niet toepasbaar gemaakt: er is geen duidelijke grond gesteld zoals aanmerkelijke onderwaarde, voldoening of zekerheid voor een niet-opeisbare schuld binnen de relevante termijn, of een kwalificerende gelieerdheid in de zin van art. 43 Fw. Art. 45 Fw biedt evenmin aanknopingspunten voor een vermoeden dat niet is ontzenuwd, omdat uit de casus geen binnen één jaar vóór faillissement verrichte rechtshandeling om niet blijkt waarop dit vermoeden hier concreet kan worden gebaseerd.

\newpage

\markboth{JOR 2014/78}{gemini-3.5-flash (k = 0)}

\subsubsection{Gemini}\label{gemini}

\texttt{gemini-3.5-flash}

\begin{center}
\small
\begin{tabular}{@{}p{7.4cm}p{2.6cm}p{2.6cm}@{}}
\toprule
Variabele & Waarde (model) & Waarde (gold) \\
\midrule
Is rechtshandeling & \texttt{true} & \texttt{true} \\
Datum rechtshandeling & \texttt{2005-12-30} & \texttt{2005-12-30} \\
Verricht door de schuldenaar & \texttt{true} & \texttt{true} \\
Om niet & \texttt{true} & \texttt{true} \\
Datum faillietverklaring & \texttt{2009-06-23} & \texttt{2009-06-23} \\
Onverplicht & \texttt{true} & \texttt{true} \\
Benadeling schuldeisers & \texttt{true} & \texttt{true} \\
Wetenschap schuldenaar aangetoond & \textcolor{gray}{\texttt{false}} & \texttt{true} \\
Wetenschap wederpartij aangetoond & \texttt{false} & \texttt{null} \\
Vermoeden wetenschap schuldenaar tegenbewijs & \texttt{null} & \texttt{null} \\
Vermoeden wetenschap wederpartij tegenbewijs & \texttt{null} & \texttt{null} \\
Tegen onderwaarde & \texttt{null} & \texttt{null} \\
Voldoening of zekerheid niet opeisbare schuld & \texttt{null} & \texttt{null} \\
Gelieerde partij & \texttt{null} & \texttt{null} \\
Reeds verplicht voor aanvang termijn & \texttt{null} & \texttt{null} \\
\textbf{Vernietigbaar} & \textbf{\textcolor{gray}{\texttt{false}}} & \textbf{\texttt{true}} \\
\bottomrule
\end{tabular}
\end{center}

\paragraph{Is rechtshandeling - overwegingen}\label{is-rechtshandeling---overwegingen-2}

De overeenkomst van verpanding en de daarin opgenomen borgtocht van 30 december 2005 zijn rechtshandelingen, omdat zij gericht zijn op het scheppen van rechtsgevolgen, te weten het vestigen van een pandrecht en het aangaan van een borgtochtverplichting.

\paragraph{Verricht door de schuldenaar - overwegingen}\label{verricht-door-de-schuldenaar---overwegingen-2}

De overeenkomst van verpanding en de borgtocht zijn aangegaan en verricht door de Vennootschap zelf (als pandgever en borg), welke partij later failliet is verklaard.

\paragraph{Om niet - overwegingen}\label{om-niet---overwegingen-2}

De Vennootschap heeft zich borg gesteld en haar activa verpand voor een schuld van een derde (Vewa Beheer B.V.) aan Van Doorn Beheer. Nu de Vennootschap hiervoor zelf geen tegenprestatie heeft bedongen of ontvangen, is de rechtshandeling om niet verricht.

\paragraph{Onverplicht - overwegingen}\label{onverplicht---overwegingen-2}

De rechtshandelingen waarbij de Vennootschap zich borg heeft gesteld en haar activa (inventaris en debiteuren) heeft verpand aan Van Doorn Beheer voor de schuld van Vewa, zijn onverplicht verricht. De verplichting tot het stellen van zekerheid voor de geldlening van € 150.000 was overeengekomen tussen Van Doorn Beheer (als verkoper/schuldeiser) en Vewa (als koper/schuldenaar) in hun koop- en leningsovereenkomst. De Vennootschap was bij deze overeenkomsten geen partij. Er bestond voor de Vennootschap dan ook geen eigen, rechtstreekse wettelijke of contractuele rechtsplicht jegens Van Doorn Beheer om zich borg te stellen of haar activa te verpanden voor de schuld van haar nieuwe moedermaatschappij. Het verstrekken van zekerheid en borgtocht voor een schuld van een derde is in een dergelijk geval onverplicht.

\paragraph{Benadeling schuldeisers - overwegingen}\label{benadeling-schuldeisers---overwegingen-2}

De rechtshandelingen van de Vennootschap bestaan uit het aangaan van een borgtocht voor een schuld van haar nieuwe moedermaatschappij (Vewa) aan Van Doorn Beheer en het verstrekken van pandrechten op haar eigen inventaris en debiteuren tot zekerheid van die borgtocht. Dit zijn onverplichte rechtshandelingen. Door deze verpanding zijn de belangrijkste activa van de Vennootschap bezwaard geraakt ten behoeve van een derde, zonder dat daar voor de Vennootschap zelf een reële tegenprestatie tegenover stond. In de hypothetische situatie zonder deze rechtshandelingen zouden de opbrengsten van de inventaris en de debiteuren volledig in de boedel zijn gevallen ten behoeve van de gezamenlijke schuldeisers. Als gevolg van de verpanding kan Van Doorn Beheer zich nu bij voorrang op deze activa verhalen, waardoor het voor de overige schuldeisers beschikbare verhaalsactief aanzienlijk is verminderd. Er is derhalve sprake van benadeling van de schuldeisers in de zin van artikel 42 Fw.

\paragraph{Wetenschap - overwegingen}\label{wetenschap---overwegingen-2}

De rechtshandelingen (borgtocht en verpanding) zijn verricht op 30 december 2005, terwijl het faillissement op 23 juni 2009 is uitgesproken. Omdat de rechtshandelingen meer dan een jaar voor de faillietverklaring zijn verricht, zijn de bewijsvermoedens van art. 43 en 45 Fw niet van toepassing. De curator moet de wetenschap van benadeling aan beide zijden rechtstreeks aantonen conform de hoofdregel van art. 42 Fw. Er zijn in de casus geen feiten gesteld waaruit blijkt dat ten tijde van de rechtshandelingen in 2005 een faillissement en een tekort daarin met een redelijke mate van waarschijnlijkheid waren te voorzien, waardoor de wetenschap van benadeling aan beide zijden niet is aangetoond.

\newpage

\markboth{JOR 2014/78}{mistral/mistral-medium-3.5-128b:fp8 (k = 5)}

\subsection{k = 5}\label{k-5}

\subsubsection{Mistral}\label{mistral-1}

\texttt{mistral/mistral-medium-3.5-128b:fp8}

\begin{center}
\small
\begin{tabular}{@{}p{7.4cm}p{2.6cm}p{2.6cm}@{}}
\toprule
Variabele & Waarde (model) & Waarde (gold) \\
\midrule
Is rechtshandeling & \texttt{true} & \texttt{true} \\
Datum rechtshandeling & \texttt{2005-12-30} & \texttt{2005-12-30} \\
Verricht door de schuldenaar & \texttt{true} & \texttt{true} \\
Om niet & \textcolor{gray}{\texttt{false}} & \texttt{true} \\
Datum faillietverklaring & \texttt{2009-06-23} & \texttt{2009-06-23} \\
Onverplicht & \texttt{true} & \texttt{true} \\
Benadeling schuldeisers & \texttt{true} & \texttt{true} \\
Wetenschap schuldenaar aangetoond & \textcolor{gray}{\texttt{false}} & \texttt{true} \\
Wetenschap wederpartij aangetoond & \texttt{false} & \texttt{null} \\
Vermoeden wetenschap schuldenaar tegenbewijs & \texttt{null} & \texttt{null} \\
Vermoeden wetenschap wederpartij tegenbewijs & \texttt{null} & \texttt{null} \\
Tegen onderwaarde & \texttt{null} & \texttt{null} \\
Voldoening of zekerheid niet opeisbare schuld & \textcolor{gray}{\texttt{true}} & \texttt{null} \\
Gelieerde partij & \textcolor{gray}{\texttt{true}} & \texttt{null} \\
Reeds verplicht voor aanvang termijn & \texttt{null} & \texttt{null} \\
\textbf{Vernietigbaar} & \textbf{\textcolor{gray}{\texttt{false}}} & \textbf{\texttt{true}} \\
\bottomrule
\end{tabular}
\end{center}

\paragraph{Is rechtshandeling - overwegingen}\label{is-rechtshandeling---overwegingen-3}

De verpanding van de bedrijfsinventaris en vorderingen door de Vennootschap aan Van Doorn Beheer (3.5) is een goederenrechtelijke handeling, gericht op een rechtsgevolg (totstandkoming pandrecht).

\paragraph{Verricht door de schuldenaar - overwegingen}\label{verricht-door-de-schuldenaar---overwegingen-3}

De Vennootschap (de latere gefailleerde) heeft zelf als pandgever de verpanding verricht (3.5).

\paragraph{Om niet - overwegingen}\label{om-niet---overwegingen-3}

De verpanding is verricht in het kader van een lening van € 150.000 door Van Doorn Beheer aan Vewa (3.4 en 3.5), waarbij de Vennootschap als pandgever optrad. Er is sprake van een tegenprestatie (de lening), zodat de handeling niet `om niet' is.

\paragraph{Onverplicht - overwegingen}\label{onverplicht---overwegingen-3}

De verpanding van de activa van de Vennootschap aan Van Doorn Beheer (art. 3.5) is een onverplichte rechtshandeling. Uit de casus blijkt dat de Vennootschap zich als pandgever heeft verbonden tot verpanding van haar debiteuren en inventaris ten behoeve van Van Doorn Beheer, ter zekerheid van een lening die Van Doorn Beheer aan Vewa Beheer B.V. had verstrekt. Er bestond geen op de wet of een overeenkomst berustende verplichting voor de Vennootschap om deze zekerheid te verstrekken. De verpanding was immers niet gericht op het zekerstellen van een eigen schuld van de Vennootschap, maar van een schuld van Vewa Beheer B.V. aan Van Doorn Beheer. Bovendien was de Vennootschap niet contractueel verplicht om zekerheid te stellen voor deze lening. De stelplicht en bewijslast van het onverplichte karakter rusten op de curator, en uit de feiten blijkt dat er geen verplichting bestond om deze zekerheid te verstrekken. Zie ook ECLI:NL:RBOVE:2018:1717, rov. 4.1.5, waar wordt overwogen dat het vestigen van een pandrecht onverplicht is wanneer daartoe geen op een overeenkomst gebaseerde verplichting bestond.

\paragraph{Benadeling schuldeisers - overwegingen}\label{benadeling-schuldeisers---overwegingen-3}

In de onderhavige casus heeft Van Doorn Beheer als aandeelhouder van de Vennootschap op 9 en 28 december 2005 de aandelen in de Vennootschap verkocht en geleverd aan Vewa Beheer B.V. (Vewa), inclusief een vordering in rekening-courant van € 349.999. De koopprijs van de aandelen bedroeg € 1, terwijl de koopprijs van de vordering € 349.999 bedroeg. Daarnaast is een lening van € 150.000 verstrekt door Van Doorn Beheer aan Vewa, met als zekerheid een tweede pandrecht op de debiteuren en inventaris van de Vennootschap.

De Vennootschap is op 23 juni 2009 in staat van faillissement verklaard. De curator betwist de rechtsgeldigheid van de pandrechten en de borgtocht en vordert vernietiging van de rechtshandelingen op grond van artikel 42 Fw.

\textbf{Benadeling van schuldeisers:}
De transactie tussen Van Doorn Beheer en Vewa behelst een onverplichte rechtshandeling, waarbij de vordering van Van Doorn Beheer op de Vennootschap (€ 349.999) is overgedragen aan Vewa. Deze vordering maakte deel uit van het actief van Van Doorn Beheer, maar door de overdracht is deze uit het vermogen van de Vennootschap verdwenen. De Vennootschap heeft hierdoor een actief verloren zonder dat daar een gelijkwaardige tegenprestatie tegenover stond.

Daarnaast is er sprake van een pandrecht op de debiteuren en inventaris van de Vennootschap ten gunste van Van Doorn Beheer. Dit pandrecht zorgt ervoor dat Van Doorn Beheer voorrang heeft ten opzichte van andere schuldeisers bij de vereffening van de boedel. Hierdoor worden de overige schuldeisers benadeeld, omdat de opbrengsten uit de verkoop van de verpande goederen niet aan de boedel (en dus de gezamenlijke schuldeisers) ten goede komen, maar aan Van Doorn Beheer.

De benadeling uit zich in een \textbf{vermindering van het actief} (verlies van de vordering en de verpande goederen voor de boedel) en een \textbf{verstoring van de rangorde tussen de schuldeisers} (Van Doorn Beheer verkrijgt voorrang ten opzichte van andere schuldeisers).

\textbf{Vergelijking hypothetische en feitelijke situatie:}
- \textbf{Hypothetische situatie (zonder de gewraakte rechtshandeling):} De vordering van € 349.999 en de debiteuren/inventaris zouden deel uitmaken van het actief van de Vennootschap en beschikbaar zijn voor alle schuldeisers, inclusief Van Doorn Beheer als gewone schuldeiser.
- \textbf{Feitelijke situatie (met de gewraakte rechtshandeling):} De vordering is overgedragen aan Vewa, en de debiteuren/inventaris zijn verpand aan Van Doorn Beheer. Hierdoor heeft Van Doorn Beheer een bevoorrechte positie verkregen ten opzichte van de andere schuldeisers, terwijl de overige schuldeisers minder verhaalsmogelijkheden hebben.

De schuldeisers zijn derhalve benadeeld in de zin van artikel 42 Fw, omdat de hypothetische situatie voor hen gunstiger was dan de feitelijke situatie. De benadeling is aanwezig op het moment waarop de curator zijn rechten doet gelden (leer van de middellijke benadeling).

\paragraph{Wetenschap - overwegingen}\label{wetenschap---overwegingen-3}

De Vennootschap heeft op 30 december 2005 een pandrecht gevestigd op haar debiteuren en inventaris ten behoeve van Van Doorn Beheer (art. 3.5). Deze rechtshandeling is verricht binnen een jaar voor de faillietverklaring (23 juni 2009) en de Vennootschap had zich niet reeds voor de aanvang van die termijn daartoe verplicht. De pandrechtvestiging is een rechtshandeling ter zekerheidstelling voor een niet-opeisbare schuld (de lening van € 150.000, art. 3.4), waardoor art. 43 lid 1 sub 2° Fw van toepassing is. Dit betekent dat de wetenschap van benadeling aan beide zijden (Vennootschap en Van Doorn Beheer) wordt vermoed. Van Doorn Beheer is een gelieerde partij (art. 43 lid 1 sub 3° Fw), omdat het een aandeelhouder was van de Vennootschap (art. 3.1). Dit versterkt het vermoeden. Er is geen sprake van tegenbewijs, omdat niet is gebleken dat het faillissement en een tekort daarin niet met een redelijke mate van waarschijnlijkheid waren te voorzien.

\newpage

\markboth{JOR 2014/78}{gpt-5.5 (k = 5)}

\subsubsection{GPT}\label{gpt-1}

\texttt{gpt-5.5}

\begin{center}
\small
\begin{tabular}{@{}p{7.4cm}p{2.6cm}p{2.6cm}@{}}
\toprule
Variabele & Waarde (model) & Waarde (gold) \\
\midrule
Is rechtshandeling & \texttt{true} & \texttt{true} \\
Datum rechtshandeling & \texttt{2005-12-30} & \texttt{2005-12-30} \\
Verricht door de schuldenaar & \texttt{true} & \texttt{true} \\
Om niet & \texttt{true} & \texttt{true} \\
Datum faillietverklaring & \texttt{2009-06-23} & \texttt{2009-06-23} \\
Onverplicht & \texttt{true} & \texttt{true} \\
Benadeling schuldeisers & \texttt{true} & \texttt{true} \\
Wetenschap schuldenaar aangetoond & \textcolor{gray}{\texttt{false}} & \texttt{true} \\
Wetenschap wederpartij aangetoond & \texttt{false} & \texttt{null} \\
Vermoeden wetenschap schuldenaar tegenbewijs & \texttt{null} & \texttt{null} \\
Vermoeden wetenschap wederpartij tegenbewijs & \texttt{null} & \texttt{null} \\
Tegen onderwaarde & \texttt{null} & \texttt{null} \\
Voldoening of zekerheid niet opeisbare schuld & \textcolor{gray}{\texttt{true}} & \texttt{null} \\
Gelieerde partij & \texttt{null} & \texttt{null} \\
Reeds verplicht voor aanvang termijn & \texttt{null} & \texttt{null} \\
\textbf{Vernietigbaar} & \textbf{\textcolor{gray}{\texttt{false}}} & \textbf{\texttt{true}} \\
\bottomrule
\end{tabular}
\end{center}

\paragraph{Is rechtshandeling - overwegingen}\label{is-rechtshandeling---overwegingen-4}

De gewraakte handelingen betreffen het zich borg stellen door de Vennootschap voor de schuld van Vewa aan Van Doorn Beheer en het verstrekken van pandrechten op inventaris en vorderingen. Dit zijn op rechtsgevolg gerichte verklaringen/overeenkomsten en daarmee rechtshandelingen.

\paragraph{Verricht door de schuldenaar - overwegingen}\label{verricht-door-de-schuldenaar---overwegingen-4}

De failliete Vennootschap, F. van Doorn B.V., wordt in de overeenkomst van verpanding aangeduid als pandgever en heeft volgens de tekst haar debiteuren en inventaris aan Van Doorn Beheer in pand verstrekt. Ook staat vermeld dat de Vennootschap zich borg heeft gesteld ten behoeve van Vewa.

\paragraph{Om niet - overwegingen}\label{om-niet---overwegingen-4}

Uit de feiten blijkt niet dat de Vennootschap voor het stellen van de borgtocht of het verstrekken van de pandrechten een reële wederprestatie heeft bedongen of ontvangen. De lening van € 150.000 is verstrekt aan Vewa, niet aan de Vennootschap.

\paragraph{Onverplicht - overwegingen}\label{onverplicht---overwegingen-4}

Voor de door de Vennootschap verrichte rechtshandelingen --- het stellen van borgtocht voor de schuld van Vewa aan Van Doorn Beheer en het vestigen van pandrechten op haar debiteuren en inventaris --- blijkt uit de casus niet van een reeds bestaande, in rechte afdwingbare verplichting uit wet of overeenkomst. De lening van € 150.000 is verstrekt door Van Doorn Beheer aan Vewa; de Vennootschap was daarbij niet de schuldenaar. Dat in de koopovereenkomst tussen Van Doorn Beheer en Vewa was voorzien in zekerheid op activa van de Vennootschap, bindt de Vennootschap niet zonder meer, nu zij geen partij bij die overeenkomst was en uit de geciteerde tekst slechts volgt dat haar directie akkoord was met de overdracht van haar schuld, niet dat de Vennootschap zich reeds rechtens had verplicht tot borgstelling of verpanding. Ook een eventuele feitelijke of economische druk om de aandelenoverdracht en financiering mogelijk te maken schept geen rechtsplicht. De latere pandovereenkomst zelf kan niet als vooraf bestaande verplichting gelden; zij vormt juist de aangevochten rechtshandeling. Dit sluit aan bij de lijn dat zekerheidstelling voor een bestaande of andermans schuld onverplicht is als daartoe niet al eerder een afdwingbare verplichting bestond, vgl. onder meer ECLI:NL:RBOVE:2018:1717, terwijl ECLI:NL:GHSHE:2016:136 juist ziet op de tegengestelde situatie waarin wél een eerdere contractuele zekerheidsverplichting bestond.

\paragraph{Benadeling schuldeisers - overwegingen}\label{benadeling-schuldeisers---overwegingen-4}

Ja. Bij de vereiste vermogensvergelijking moet de feitelijke situatie waarin de borgtocht en pandrechten ten gunste van Van Doorn Beheer in stand blijven worden vergeleken met de hypothetische situatie zonder die rechtshandeling. Zonder de borgstelling en verpanding zou Van Doorn Beheer voor de aan Vewa verstrekte lening van € 150.000 in beginsel geen verhaalsrecht op het vermogen van de Vennootschap hebben gehad. De inventaris en vorderingen van de Vennootschap zouden dan, behoudens eventueel reeds bestaande oudere zekerheden, beschikbaar zijn gebleven voor de gezamenlijke schuldeisers. Door de aangevochten rechtshandeling is daarentegen een verhaalspositie van Van Doorn Beheer op activa van de Vennootschap gecreëerd ter zekerheid van een schuld van Vewa, terwijl geen gelijkwaardige tegenprestatie aan de Vennootschap blijkt. Dat leidt tot vermeerdering van de passiva van de Vennootschap door de borgtocht en tot verstoring van de rangorde doordat Van Doorn Beheer zich als pandhouder bij voorrang kan verhalen op debiteuren en inventaris. Dat nadeel is ook concreet geworden: vóór faillissement heeft de Vennootschap naast huurpenningen € 49.500 aan Van Doorn Beheer betaald en na faillissement is Van Doorn Beheer tot inning van de openstaande vorderingen overgegaan. Daarmee zijn activa of opbrengsten aan de boedel en de gezamenlijke schuldeisers onttrokken althans aan één schuldeiser met voorrang toegekomen. Er is niet gebleken dat de gewone schuldeisers zich zonder de rechtshandeling toch niet op deze goederen of opbrengsten hadden kunnen verhalen, zoals bijvoorbeeld het geval kan zijn bij een reeds volledig dekkend ouder zekerheidsrecht. De situatie verschilt daarom van gevallen waarin een separatist ook zonder de aangevochten handeling dezelfde opbrengst zou hebben gekregen. De benadeling bestaat hier zowel uit vermindering van verhaal als uit rangordeverstoring, vergelijkbaar met de lijn dat het onverplicht verschaffen van zekerheid of het overdragen van activa met verrekening ten koste van de overige schuldeisers benadelend kan zijn.

\paragraph{Wetenschap - overwegingen}\label{wetenschap---overwegingen-4}

De te beoordelen borgtocht/verpanding door de Vennootschap dateert van december 2005, terwijl het faillissement pas op 23 juni 2009 is uitgesproken. Voor het wetenschapsvereiste geldt daarom de hoofdregel van art. 42 Fw: de curator moet rechtstreeks aantonen dat ten tijde van de rechtshandeling het faillissement én een tekort daarin met een redelijke mate van waarschijnlijkheid waren te voorzien. De verstrekte feiten wijzen wel op een voor de Vennootschap nadelige zekerheidsconstructie ten behoeve van de koopsomfinanciering van Vewa, maar bevatten onvoldoende concrete aanwijzingen dat de Vennootschap eind 2005 al zodanig financieel in zwaar weer verkeerde dat faillissement met tekort voorzienbaar was. Ook ten aanzien van Van Doorn Beheer is die wetenschap niet rechtstreeks aangetoond. Een bewijsvermoeden van art. 43 Fw helpt de curator hier niet, omdat de gewraakte borgtocht/verpanding ruim buiten de eenjaarstermijn ligt; voor zover de verpanding kwalificeert als zekerheid voor een niet-opeisbare schuld, komt dat vermoeden daardoor niet tot toepassing. Art. 45 Fw is evenmin aan de orde voor een vermoeden binnen één jaar. Van succesvol of onsuccesvol tegenbewijs tegen een toegepast vermoeden is daarom geen sprake.

\newpage

\markboth{JOR 2014/78}{gemini-3.5-flash (k = 5)}

\subsubsection{Gemini}\label{gemini-1}

\texttt{gemini-3.5-flash}

\begin{center}
\small
\begin{tabular}{@{}p{7.4cm}p{2.6cm}p{2.6cm}@{}}
\toprule
Variabele & Waarde (model) & Waarde (gold) \\
\midrule
Is rechtshandeling & \texttt{true} & \texttt{true} \\
Datum rechtshandeling & \texttt{2005-12-30} & \texttt{2005-12-30} \\
Verricht door de schuldenaar & \texttt{true} & \texttt{true} \\
Om niet & \texttt{true} & \texttt{true} \\
Datum faillietverklaring & \texttt{2009-06-23} & \texttt{2009-06-23} \\
Onverplicht & \texttt{true} & \texttt{true} \\
Benadeling schuldeisers & \texttt{true} & \texttt{true} \\
Wetenschap schuldenaar aangetoond & \textcolor{gray}{\texttt{false}} & \texttt{true} \\
Wetenschap wederpartij aangetoond & \texttt{false} & \texttt{null} \\
Vermoeden wetenschap schuldenaar tegenbewijs & \texttt{null} & \texttt{null} \\
Vermoeden wetenschap wederpartij tegenbewijs & \texttt{null} & \texttt{null} \\
Tegen onderwaarde & \texttt{null} & \texttt{null} \\
Voldoening of zekerheid niet opeisbare schuld & \texttt{null} & \texttt{null} \\
Gelieerde partij & \texttt{null} & \texttt{null} \\
Reeds verplicht voor aanvang termijn & \texttt{null} & \texttt{null} \\
\textbf{Vernietigbaar} & \textbf{\textcolor{gray}{\texttt{false}}} & \textbf{\texttt{true}} \\
\bottomrule
\end{tabular}
\end{center}

\paragraph{Is rechtshandeling - overwegingen}\label{is-rechtshandeling---overwegingen-5}

De overeenkomst van verpanding en de daarin opgenomen borgtocht van 30 december 2005 zijn rechtshandelingen, omdat zij gericht zijn op het scheppen van rechtsgevolgen, te weten het vestigen van een pandrecht en het aangaan van een borgtochtverplichting.

\paragraph{Verricht door de schuldenaar - overwegingen}\label{verricht-door-de-schuldenaar---overwegingen-5}

De overeenkomst van verpanding en de borgtocht zijn aangegaan en verricht door de Vennootschap zelf (als pandgever en borg), welke partij later failliet is verklaard.

\paragraph{Om niet - overwegingen}\label{om-niet---overwegingen-5}

De Vennootschap heeft zich borg gesteld en haar activa verpand voor een schuld van een derde (Vewa Beheer B.V.) aan Van Doorn Beheer. Nu de Vennootschap hiervoor zelf geen tegenprestatie heeft bedongen of ontvangen, is de rechtshandeling om niet verricht.

\paragraph{Onverplicht - overwegingen}\label{onverplicht---overwegingen-5}

De rechtshandelingen van de Vennootschap (het aangaan van de borgtocht en de verpanding van haar activa ten behoeve van de schuld van Vewa aan Van Doorn Beheer) zijn onverplicht verricht. Een rechtshandeling is onverplicht wanneer daartoe geen op de wet of een overeenkomst berustende rechtsplicht bestond die in rechte kon worden afgedwongen. In deze casus was de Vennootschap zelf geen partij bij de koopovereenkomst of de geldleningsovereenkomst tussen Van Doorn Beheer en Vewa, waarin de verplichting tot het stellen van zekerheid was overeengekomen. De Vennootschap had derhalve geen voorafgaande, in rechte afdwingbare rechtsplicht jegens Van Doorn Beheer om zich borg te stellen voor de schuld van Vewa of om haar eigen activa (debiteuren en inventaris) aan Van Doorn Beheer te verpanden. Het feit dat de directie van de Vennootschap zich in de considerans van de koopovereenkomst akkoord had verklaard met de overdracht van de schuld, schept geen rechtsplicht tot het verlenen van derdenzekerheid of borgtocht.

\paragraph{Benadeling schuldeisers - overwegingen}\label{benadeling-schuldeisers---overwegingen-5}

In deze casus heeft de Vennootschap zich borg gesteld en haar activa (inventaris en debiteuren) verpand aan Van Doorn Beheer tot zekerheid van een geldlening van € 150.000 die haar nieuwe aandeelhouder, Vewa, was aangegaan voor de overname van de aandelen. Deze rechtshandelingen zijn onverplicht verricht. Om te beoordelen of de schuldeisers zijn benadeeld, dient een vermogensvergelijking te worden gemaakt tussen de feitelijke situatie na de rechtshandelingen en de hypothetische situatie waarin deze achterwege waren gebleven. Door de verpanding en borgtocht zijn de activa van de Vennootschap bezwaard voor een schuld van een derde (Vewa), zonder dat de Vennootschap hiervoor een reële tegenprestatie ontving. Na het faillissement heeft Van Doorn Beheer deze zekerheden uitgewonnen door de inventaris te verkopen en debiteuren te innen. Zonder de verpanding en borgtocht zouden deze activa onbezwaard in de boedel zijn gevallen en beschikbaar zijn geweest voor verhaal door de gezamenlijke schuldeisers. De rechtshandelingen hebben aldus geleid tot een rechtstreekse vermindering van het actieve boedelvolume en een verstoring van de rangorde ten nadele van de gezamenlijke schuldeisers. Er is derhalve sprake van benadeling in de zin van artikel 42 Fw.

\paragraph{Wetenschap - overwegingen}\label{wetenschap---overwegingen-5}

De omstreden rechtshandelingen (de borgtocht en de verpanding) zijn verricht op 30 december 2005, terwijl het faillissement van de Vennootschap pas op 23 juni 2009 is uitgesproken. Omdat deze handelingen ruim buiten de jaartermijn van art. 43 en 45 Fw vallen, zijn de wettelijke bewijsvermoedens niet van toepassing. De curator draagt de reguliere bewijslast op grond van art. 42 Fw. Ten tijde van de rechtshandelingen in december 2005 (circa 3,5 jaar vóór het faillissement) waren het faillissement en een tekort daarin niet met een redelijke mate van waarschijnlijkheid te voorzien. De wetenschap van benadeling is dan ook aan geen van beide zijden aangetoond. De latere betalingen tussen 2007 en 2009 zijn verricht ter voldoening van een (alsdan) opeisbare verbintenis en zijn derhalve niet onverplicht.

\newpage

\markboth{JOR 2014/78}{mistral/mistral-medium-3.5-128b:fp8 (k = 10)}

\subsection{k = 10}\label{k-10}

\subsubsection{Mistral}\label{mistral-2}

\texttt{mistral/mistral-medium-3.5-128b:fp8}

\begin{center}
\small
\begin{tabular}{@{}p{7.4cm}p{2.6cm}p{2.6cm}@{}}
\toprule
Variabele & Waarde (model) & Waarde (gold) \\
\midrule
Is rechtshandeling & \texttt{true} & \texttt{true} \\
Datum rechtshandeling & \texttt{2005-12-30} & \texttt{2005-12-30} \\
Verricht door de schuldenaar & \texttt{true} & \texttt{true} \\
Om niet & \textcolor{gray}{\texttt{false}} & \texttt{true} \\
Datum faillietverklaring & \texttt{2009-06-23} & \texttt{2009-06-23} \\
Onverplicht & \texttt{true} & \texttt{true} \\
Benadeling schuldeisers & \texttt{true} & \texttt{true} \\
Wetenschap schuldenaar aangetoond & \textcolor{gray}{\texttt{false}} & \texttt{true} \\
Wetenschap wederpartij aangetoond & \texttt{false} & \texttt{null} \\
Vermoeden wetenschap schuldenaar tegenbewijs & \texttt{null} & \texttt{null} \\
Vermoeden wetenschap wederpartij tegenbewijs & \texttt{null} & \texttt{null} \\
Tegen onderwaarde & \texttt{null} & \texttt{null} \\
Voldoening of zekerheid niet opeisbare schuld & \texttt{null} & \texttt{null} \\
Gelieerde partij & \textcolor{gray}{\texttt{true}} & \texttt{null} \\
Reeds verplicht voor aanvang termijn & \texttt{false} & \texttt{null} \\
\textbf{Vernietigbaar} & \textbf{\textcolor{gray}{\texttt{false}}} & \textbf{\texttt{true}} \\
\bottomrule
\end{tabular}
\end{center}

\paragraph{Is rechtshandeling - overwegingen}\label{is-rechtshandeling---overwegingen-6}

De verpanding van de bedrijfsinventaris en vorderingen door de Vennootschap aan Van Doorn Beheer (3.5) is een goederenrechtelijke handeling, gericht op een rechtsgevolg (vestiging van pandrecht).

\paragraph{Verricht door de schuldenaar - overwegingen}\label{verricht-door-de-schuldenaar---overwegingen-6}

De Vennootschap (de latere gefailleerde) heeft zelf als pandgever de verpanding verricht (3.5).

\paragraph{Om niet - overwegingen}\label{om-niet---overwegingen-6}

De verpanding is verricht in het kader van een lening van € 150.000 door Van Doorn Beheer aan Vewa (3.4 en 3.5), waarbij de Vennootschap als pandgever optrad. Er is sprake van een wederprestatie (lening tegen zekerheid), zodat de handeling niet `om niet' is.

\paragraph{Onverplicht - overwegingen}\label{onverplicht---overwegingen-6}

De verpanding van de activa van de Vennootschap aan Van Doorn Beheer is onverplicht verricht. Uit de casus blijkt dat de Vennootschap een vordering in rekening-courant had op Van Doorn Beheer, maar er was geen sprake van een bestaande verplichting om zekerheid te stellen voor deze vordering. De pandovereenkomst van 30 december 2005 (art. 3.4) vermeldt weliswaar dat de Vennootschap zich als borg heeft gesteld voor de lening van Vewa aan Van Doorn Beheer, maar dit impliceert niet dat de Vennootschap verplicht was om haar eigen activa te verpanden. De verpanding is derhalve een onverplichte rechtshandeling, omdat er geen op de wet of een overeenkomst berustende rechtsplicht bestond om deze zekerheid te verstrekken. Morele, feitelijke of economische druk (zoals de relatie tussen Van Doorn Beheer en de Vennootschap) schept geen rechtsplicht (vgl. ECLI:NL:RBOVE:2018:1717, rov. 4.1.5).

\paragraph{Benadeling schuldeisers - overwegingen}\label{benadeling-schuldeisers---overwegingen-6}

In de onderhavige casus heeft Van Doorn Beheer op 9 en 28 december 2005 de aandelen in de Vennootschap verkocht en geleverd aan Vewa Beheer B.V. (Vewa), inclusief de vordering in rekening-courant van Van Doorn Beheer op de Vennootschap. De koopprijs voor de aandelen bedroeg € 1, terwijl de koopprijs voor de vordering € 349.999 bedroeg. Daarnaast is een lening van € 150.000 verstrekt door Van Doorn Beheer aan Vewa, met als zekerheid een tweede pandrecht op de debiteuren en inventaris van de Vennootschap.

Voor de beoordeling of de schuldeisers zijn benadeeld in de zin van artikel 42 Fw, dient een vermogensvergelijking te worden gemaakt tussen de hypothetische situatie zonder de gewraakte rechtshandeling en de feitelijke situatie met de rechtshandeling.

\begin{enumerate}
\def\labelenumi{\arabic{enumi}.}
\tightlist
\item
  \textbf{Vermogensvergelijking}:

  \begin{itemize}
  \tightlist
  \item
    \textbf{Hypothetische situatie (zonder de rechtshandeling)}: De vordering van Van Doorn Beheer op de Vennootschap (€ 349.999) zou in het vermogen van Van Doorn Beheer zijn gebleven. De Vennootschap zou de vordering niet hoeven te voldoen, en de activa (debiteuren en inventaris) van de Vennootschap zouden onbelast blijven met het pandrecht ten gunste van Van Doorn Beheer.
  \item
    \textbf{Feitelijke situatie (met de rechtshandeling)}: De vordering is overgedragen aan Vewa, en de Vennootschap heeft een pandrecht op haar activa ten gunste van Van Doorn Beheer (via Vewa). Dit betekent dat de activa van de Vennootschap nu belast zijn met een pandrecht, wat de verhaalsmogelijkheden voor andere schuldeisers beperkt.
  \end{itemize}
\item
  \textbf{Benadeling}:

  \begin{itemize}
  \tightlist
  \item
    Door de overdracht van de vordering en de vestiging van het pandrecht zijn de activa van de Vennootschap (debiteuren en inventaris) belast met een zekerheidsrecht ten gunste van Van Doorn Beheer. Dit betekent dat andere schuldeisers van de Vennootschap zich niet meer onbeperkt op deze activa kunnen verhalen. De verhaalsmogelijkheden voor de overige schuldeisers zijn hierdoor verminderd, wat neerkomt op benadeling in de zin van artikel 42 Fw.
  \item
    De benadeling is aanwezig op het moment waarop de curator zijn rechten doet gelden (leer van de middellijke benadeling), omdat de activa van de Vennootschap nu belast zijn met een pandrecht dat de verhaalsmogelijkheden voor andere schuldeisers beperkt.
  \end{itemize}
\item
  \textbf{Conclusie}:
  De schuldeisers zijn benadeeld doordat de activa van de Vennootschap zijn belast met een pandrecht ten gunste van Van Doorn Beheer, wat hun verhaalsmogelijkheden beperkt. Dit voldoet aan de eis van benadeling in de zin van artikel 42 Fw.
\end{enumerate}

\paragraph{Wetenschap - overwegingen}\label{wetenschap---overwegingen-6}

De rechtshandeling betreft de verkoop en levering van aandelen en een vordering in rekening-courant door Van Doorn Beheer aan Vewa Beheer B.V. (Vewa) op 9 en 28 december 2005, alsmede de vestiging van een pandrecht op de activa van de Vennootschap ten gunste van Van Doorn Beheer op 30 december 2005. De Vennootschap is op 23 juni 2009 failliet verklaard. De curator vernietigt deze rechtshandelingen op grond van art. 42 Fw.

De rechtshandeling is onverplicht, omdat er geen verplichting bestond tot het vestigen van het pandrecht of de verkoop van de aandelen en vordering. De benadeling van schuldeisers is aannemelijk, omdat de activa van de Vennootschap door het pandrecht aan Van Doorn Beheer zijn verpand, waardoor de overige schuldeisers in hun verhaalsmogelijkheden zijn beperkt.

Wat betreft de wetenschap van benadeling: de rechtshandelingen zijn verricht binnen een jaar voor de faillietverklaring (december 2005 vs.~juni 2009). Van Doorn Beheer was aandeelhouder van de Vennootschap, en de rechtshandeling is verricht met een gelieerde partij in de zin van art. 43 lid 1 sub 3° Fw (aandeelhouder). Dit betekent dat het wettelijk vermoeden van art. 43 lid 1 Fw van toepassing is: de wetenschap van benadeling wordt vermoed aan beide zijden (schuldenaar en wederpartij) te bestaan, tenzij tegenbewijs wordt geleverd. Er is geen sprake van dat de schuldenaar zich reeds voor de aanvang van de jaartermijn (d.w.z. voor december 2005) tot de rechtshandeling had verplicht.

Er is geen tegenbewijs geleverd dat het faillissement en een tekort daarin niet met een redelijke mate van waarschijnlijkheid waren te voorzien. De omstandigheden (financiële problemen van de Vennootschap, de aard van de transactie) rechtvaardigen het vermoeden dat Van Doorn Beheer en de Vennootschap wisten of behoorden te weten dat benadeling van schuldeisers het gevolg zou zijn. Derhalve is de wetenschap van benadeling aan beide zijden aannemelijk.

\newpage

\markboth{JOR 2014/78}{gpt-5.5 (k = 10)}

\subsubsection{GPT}\label{gpt-2}

\texttt{gpt-5.5}

\begin{center}
\small
\begin{tabular}{@{}p{7.4cm}p{2.6cm}p{2.6cm}@{}}
\toprule
Variabele & Waarde (model) & Waarde (gold) \\
\midrule
Is rechtshandeling & \texttt{true} & \texttt{true} \\
Datum rechtshandeling & \texttt{2005-12-30} & \texttt{2005-12-30} \\
Verricht door de schuldenaar & \texttt{true} & \texttt{true} \\
Om niet & \texttt{true} & \texttt{true} \\
Datum faillietverklaring & \texttt{2009-06-23} & \texttt{2009-06-23} \\
Onverplicht & \texttt{true} & \texttt{true} \\
Benadeling schuldeisers & \texttt{true} & \texttt{true} \\
Wetenschap schuldenaar aangetoond & \textcolor{gray}{\texttt{false}} & \texttt{true} \\
Wetenschap wederpartij aangetoond & \texttt{false} & \texttt{null} \\
Vermoeden wetenschap schuldenaar tegenbewijs & \texttt{null} & \texttt{null} \\
Vermoeden wetenschap wederpartij tegenbewijs & \texttt{null} & \texttt{null} \\
Tegen onderwaarde & \texttt{null} & \texttt{null} \\
Voldoening of zekerheid niet opeisbare schuld & \texttt{null} & \texttt{null} \\
Gelieerde partij & \texttt{null} & \texttt{null} \\
Reeds verplicht voor aanvang termijn & \texttt{null} & \texttt{null} \\
\textbf{Vernietigbaar} & \textbf{\textcolor{gray}{\texttt{false}}} & \textbf{\texttt{true}} \\
\bottomrule
\end{tabular}
\end{center}

\paragraph{Is rechtshandeling - overwegingen}\label{is-rechtshandeling---overwegingen-7}

De gewraakte handelingen betreffen het zich borg stellen door de Vennootschap voor de schuld van Vewa aan Van Doorn Beheer en het verstrekken van pandrechten op haar inventaris en vorderingen. Dit zijn op rechtsgevolg gerichte rechtshandelingen, vastgelegd in de overeenkomst van verpanding van 30 december 2005.

\paragraph{Verricht door de schuldenaar - overwegingen}\label{verricht-door-de-schuldenaar---overwegingen-7}

De failliete Vennootschap, F. van Doorn B.V., heeft zich volgens de feiten borg gesteld en trad in de overeenkomst van verpanding op als pandgever. De handelingen zijn dus door de schuldenaar zelf verricht.

\paragraph{Om niet - overwegingen}\label{om-niet---overwegingen-7}

Uit de feiten blijkt dat de borgtocht en pandrechten door de Vennootschap zijn verstrekt ter zekerheid van een lening van Van Doorn Beheer aan Vewa. Er is geen wederprestatie aan de Vennootschap vermeld of bedongen. Op basis van de aangeleverde feiten is de rechtshandeling daarom om niet.

\paragraph{Onverplicht - overwegingen}\label{onverplicht---overwegingen-7}

De aangevochten rechtshandeling betreft, voor zover hier relevant, het door de Vennootschap stellen van borgtocht voor de schuld van Vewa aan Van Doorn Beheer en het vestigen van pandrechten op haar debiteuren en inventaris tot zekerheid van die schuld. Voor een verplicht karakter is beslissend of Van Doorn Beheer deze prestatie van de Vennootschap in rechte kon afdwingen op grond van wet of overeenkomst. Uit de casus blijkt niet dat de Vennootschap partij was bij de koopovereenkomst tussen Van Doorn Beheer en Vewa of zich daarin jegens Van Doorn Beheer reeds afdwingbaar had verbonden tot borgstelling of verpanding. De omstandigheid dat in die koopovereenkomst is voorzien in zekerheid in de vorm van een tweede pandrecht op activa van de Vennootschap, schept zonder meer geen rechtsplicht voor de Vennootschap zelf. Ook het feit dat de directie van de Vennootschap akkoord was met de overdracht van haar schuld aan Vewa ziet niet op een afdwingbare verplichting tot het stellen van borgtocht of pandrechten. Dat de zekerheden economisch of feitelijk onderdeel waren van de aandelenoverdracht en financieringsconstructie is onvoldoende: feitelijke of commerciële druk levert geen rechtsplicht op. Vergelijkbaar volgt uit onder meer ECLI:NL:RBOVE:2018:1717 dat zekerheidstelling voor een bestaande schuld onverplicht is wanneer geen eerdere contractuele verplichting tot zekerheidstelling bestond, terwijl uit ECLI:NL:GHSHE:2016:136 juist volgt dat dit anders kan zijn indien een eerdere geldleningsovereenkomst wél een afdwingbare zekerheidsverplichting bevatte. Een dergelijke voorafgaande afdwingbare verplichting van de Vennootschap blijkt hier niet.

\paragraph{Benadeling schuldeisers - overwegingen}\label{benadeling-schuldeisers---overwegingen-7}

De benadeling moet worden beoordeeld door vergelijking van de feitelijke situatie waarin de borgtocht/pandrechten en de daarop gebaseerde betalingen/inningen in stand blijven met de hypothetische situatie zonder die rechtshandelingen. In de feitelijke situatie heeft de Vennootschap zich borg gesteld voor een schuld van Vewa aan Van Doorn Beheer en heeft zij haar debiteuren en inventaris aan Van Doorn Beheer verpand. Daardoor is Van Doorn Beheer, voor een schuld die primair niet van de Vennootschap maar van Vewa was, in een bevoorrechte verhaalspositie gebracht. Na faillissement heeft Van Doorn Beheer bovendien openstaande vorderingen van de Vennootschap geïnd en aanspraak gemaakt op de opbrengst van activa, terwijl de Vennootschap vóór faillissement ook al € 49.500 aan Van Doorn Beheer heeft betaald. Zonder de borgtocht en verpanding zouden de debiteuren, inventaris/opbrengst en betaalde bedragen in beginsel voor verhaal door de gezamenlijke schuldeisers beschikbaar zijn gebleven. De rechtshandelingen leiden dus tot vermindering van het actief en in elk geval tot verstoring van de rangorde tussen schuldeisers, doordat Van Doorn Beheer ten opzichte van de overige schuldeisers wordt bevoordeeld. Dat de verpanding als tweede pandrecht is aangeduid, sluit benadeling niet uit, nu niet blijkt dat een hoger gerangschikte pandhouder de betrokken activa volledig zou hebben uitgewonnen zodat de overige schuldeisers zich daarop hoe dan ook niet hadden kunnen verhalen. Evenmin blijkt van een gelijkwaardig voordeel voor de Vennootschap dat deze nadelige gevolgen compenseert; de financiering strekte primair tot betaling van de door Vewa verschuldigde koopprijs. Dit sluit aan bij de lijn dat zekerheidsverschaffing voor een bestaande of derde schuld benadeling kan opleveren door vermindering van verhaalsmogelijkheden of rangverstoring, vgl. onder meer ECLI:NL:RBOVE:2018:1717 en ECLI:NL:GHARL:2023:8459.

\paragraph{Wetenschap - overwegingen}\label{wetenschap---overwegingen-7}

Voor de borgtocht en verpanding door de Vennootschap ten gunste van Van Doorn Beheer blijkt uit de gegeven feiten onvoldoende dat op het moment van die rechtshandelingen het faillissement van de Vennootschap en een tekort daarin met een redelijke mate van waarschijnlijkheid voorzienbaar waren. Een oogmerk tot benadeling is niet vereist, maar de enkele omstandigheid dat de Vennootschap zekerheid stelde voor een schuld van Vewa aan Van Doorn Beheer, in het kader van de aandelentransactie, en dat zij enkele jaren later failliet is gegaan, is onvoldoende voor rechtstreekse wetenschap van benadeling. Voor zover de rechtshandeling als anders dan om niet moet worden beoordeeld, is evenmin rechtstreeks aangetoond dat Van Doorn Beheer deze wetenschap had. De bewijsvermoedens van art. 43 Fw worden hier niet dragend toegepast: de relevante rechtshandelingen dateren uit december 2005, terwijl het faillissement pas op 23 juni 2009 is uitgesproken. Art. 45 Fw biedt om dezelfde reden geen steun voor een vermoeden. Tegenbewijs tegen een toegepast vermoeden is daarom niet aan de orde.

\newpage

\markboth{JOR 2014/78}{gemini-3.5-flash (k = 10)}

\subsubsection{Gemini}\label{gemini-2}

\texttt{gemini-3.5-flash}

\begin{center}
\small
\begin{tabular}{@{}p{7.4cm}p{2.6cm}p{2.6cm}@{}}
\toprule
Variabele & Waarde (model) & Waarde (gold) \\
\midrule
Is rechtshandeling & \texttt{true} & \texttt{true} \\
Datum rechtshandeling & \texttt{2005-12-30} & \texttt{2005-12-30} \\
Verricht door de schuldenaar & \texttt{true} & \texttt{true} \\
Om niet & \texttt{true} & \texttt{true} \\
Datum faillietverklaring & \texttt{2009-06-23} & \texttt{2009-06-23} \\
Onverplicht & \texttt{true} & \texttt{true} \\
Benadeling schuldeisers & \texttt{true} & \texttt{true} \\
Wetenschap schuldenaar aangetoond & \textcolor{gray}{\texttt{false}} & \texttt{true} \\
Wetenschap wederpartij aangetoond & \texttt{false} & \texttt{null} \\
Vermoeden wetenschap schuldenaar tegenbewijs & \texttt{null} & \texttt{null} \\
Vermoeden wetenschap wederpartij tegenbewijs & \texttt{null} & \texttt{null} \\
Tegen onderwaarde & \texttt{null} & \texttt{null} \\
Voldoening of zekerheid niet opeisbare schuld & \texttt{null} & \texttt{null} \\
Gelieerde partij & \texttt{null} & \texttt{null} \\
Reeds verplicht voor aanvang termijn & \texttt{null} & \texttt{null} \\
\textbf{Vernietigbaar} & \textbf{\textcolor{gray}{\texttt{false}}} & \textbf{\texttt{true}} \\
\bottomrule
\end{tabular}
\end{center}

\paragraph{Is rechtshandeling - overwegingen}\label{is-rechtshandeling---overwegingen-8}

De overeenkomst van verpanding en de daarin opgenomen borgtocht d.d. 30 december 2005 zijn rechtshandelingen, aangezien zij gericht zijn op het in het leven roepen van goederenrechtelijke en obligatoire rechtsgevolgen (het vestigen van een pandrecht en het aangaan van een borgtocht).

\paragraph{Verricht door de schuldenaar - overwegingen}\label{verricht-door-de-schuldenaar---overwegingen-8}

De overeenkomst van verpanding en de borgtocht zijn aangegaan en verricht door de Vennootschap (de latere failliet) zelf als pandgever en borg.

\paragraph{Om niet - overwegingen}\label{om-niet---overwegingen-8}

De Vennootschap heeft zich borg gesteld en haar activa verpand ten behoeve van een schuld van een derde (Vewa Beheer B.V.) aan Van Doorn Beheer, zonder dat de Vennootschap daar zelf een reële tegenprestatie voor ontving. De rechtshandeling is derhalve om niet verricht.

\paragraph{Onverplicht - overwegingen}\label{onverplicht---overwegingen-8}

De rechtshandelingen van de Vennootschap (F. van Doorn B.V.) bestaande uit het verstrekken van een borgtocht en het vestigen van pandrechten op haar debiteuren en inventaris ten behoeve van Van Doorn Beheer B.V. zijn onverplicht verricht. Er bestond voor de Vennootschap geen eerdere, op de wet of een overeenkomst berustende rechtsplicht om zich jegens Van Doorn Beheer te verbinden voor de schuld van de koper (Vewa Beheer B.V.). De verplichting om deze zekerheden te stellen was weliswaar overeengekomen in de koopovereenkomst en de leningsovereenkomst, maar deze overeenkomsten waren gesloten tussen Van Doorn Beheer en Vewa. De Vennootschap was hierbij geen partij en kon derhalve niet in rechte worden gedwongen om deze zekerheden te verstrekken. Het verstrekken van zekerheid voor de schuld van een derde zonder eerdere eigen rechtsplicht is onverplicht (vergelijk ECLI:NL:RBOVE:2018:1717 en ECLI:NL:RBROT:2009:BH6448).

\paragraph{Benadeling schuldeisers - overwegingen}\label{benadeling-schuldeisers---overwegingen-8}

De aangevochten rechtshandelingen betreffen de door de Vennootschap verstrekte borgtocht en de verpanding van haar inventaris en debiteuren aan Van Doorn Beheer, tot zekerheid van een geldlening van € 150.000,- die Van Doorn Beheer heeft verstrekt aan de koper van de aandelen (Vewa). Dit is een onverplichte rechtshandeling waarbij de Vennootschap haar eigen activa bezwaart voor een schuld van een derde (haar nieuwe moedermaatschappij), zonder dat daar voor de Vennootschap zelf een reële tegenprestatie of bedrijfseconomisch voordeel tegenover staat. Door deze verpanding zijn de belangrijkste verhaalscomponenten van de Vennootschap (inventaris en debiteuren) bezwaard met een pandrecht. Bij een vergelijking tussen de hypothetische situatie zonder deze rechtshandelingen (waarin deze activa onbezwaard in de boedel zouden vallen) en de feitelijke situatie (waarin de activa bij voorrang toekomen aan de pandhouder), staat vast dat de verhaalsmogelijkheden van de gezamenlijke schuldeisers zijn verminderd. Er is derhalve sprake van daadwerkelijke benadeling van de schuldeisers in de zin van artikel 42 Fw.

\paragraph{Wetenschap - overwegingen}\label{wetenschap---overwegingen-8}

De rechtshandelingen (de borgtocht en verpanding) zijn verricht op 30 december 2005, terwijl het faillissement van de Vennootschap pas op 23 juni 2009 is uitgesproken. Omdat deze handelingen ruim buiten de jaartermijn van art. 43 en 45 Fw vallen, zijn de wettelijke bewijsvermoedens niet van toepassing. De curator moet de wetenschap van benadeling derhalve rechtstreeks aantonen op grond van art. 42 Fw. Hiervoor is vereist dat ten tijde van de rechtshandelingen het faillissement en een tekort daarin met een redelijke mate van waarschijnlijkheid voorzienbaar waren. Gelet op het tijdsverloop van circa 3,5 jaar en het feit dat de Vennootschap pas veel later insolvent is geraakt, is deze wetenschap van benadeling aan beide zijden niet aangetoond.

\newpage

\markboth{JOR 2026/137}{Gold standard}

\section{JOR 2026/137}\label{jor-2026137}

\texttt{ECLI:NL:GHARL:2025:4787}, uitspraakdatum 08-07-2025

\subsection{Gold}\label{gold-1}

\begin{center}
\small
\begin{tabular}{@{}p{8cm}p{2.6cm}@{}}
\toprule
Variabele & Waarde \\
\midrule
Is rechtshandeling & \texttt{true} \\
Datum rechtshandeling & \texttt{2022-02-04} \\
Verricht door de schuldenaar & \texttt{true} \\
Om niet & \texttt{false} \\
Datum faillietverklaring & \texttt{2022-05-02} \\
Onverplicht & \texttt{false} \\
Benadeling schuldeisers & \texttt{false} \\
Wetenschap schuldenaar aangetoond & \texttt{false} \\
Wetenschap wederpartij aangetoond & \texttt{null} \\
Vermoeden wetenschap schuldenaar tegenbewijs & \texttt{null} \\
Vermoeden wetenschap wederpartij tegenbewijs & \texttt{null} \\
Tegen onderwaarde & \texttt{null} \\
Voldoening of zekerheid niet opeisbare schuld & \texttt{null} \\
Gelieerde partij & \texttt{null} \\
Reeds verplicht voor aanvang termijn & \texttt{null} \\
\textbf{Vernietigbaar} & \textbf{\texttt{false}} \\
\bottomrule
\end{tabular}
\end{center}

\subsubsection{Uitspraak}\label{uitspraak-1}

(\ldots; red.)

\paragraph{2. De kern van de zaak}\label{de-kern-van-de-zaak}

2.1. Het gaat in deze zaak om de vraag of boekingen die Durisol enkele maanden voor haar faillissement op de rekening van een van haar groepsvennootschappen heeft gedaan door de curator kunnen worden vernietigd, of dat de curator althans een schadevordering op Mevriet c.s. heeft die op deze betalingen is gebaseerd. Dat geschil heeft de volgende feitelijke achtergrond.

2.2. {[}geïntimeerde6{]} en {[}geïntimeerde5{]} zijn bestuurders en enig aandeelhouders van respectievelijk B2R en {[}geïntimeerde4{]}. Deze vennootschappen zijn of waren op hun beurt bestuurders en aandeelhouders van Mevriet, Mevriet Vastgoed en (voorheen) van Durisol.

2.3. Rabobank Salland U.A. heeft op enig moment een kredietfaciliteit met een maximale debetstand van € 350.000.00 aan Durisol verleend. Tot zekerheid van betaling van onder meer deze vordering is op 26 juli 2016 een recht van hypotheek gevestigd met als onderpand het bedrijfsterrein met opstallen van Mevriet Vastgoed B.V.

2.4. Op 11 januari 2022 hebben onder meer Durisol en Mevriet Vastgoed de roerende en onroerende activa van hun ondernemingen verkocht. Daaronder vielen met name de met hypotheek bezwaarde bedrijfsterreinen en opstallen van Mevriet Vastgoed. Op dat moment beliep de debetstand van de rekening van Durisol € 56.837,93. Rabobank wilde dat bedrag na royement van de hypotheek op 1 februari 2022 aan Durisol overmaken en het restant aan Mevriet Vastgoed, maar maakte het gehele surplus van € 349.981.74 aan Durisol over (€ 350.000, verminderd met € 18,26 aan rente). {[}naam1{]} van Rabo heeft hierover het volgende geschreven: ``Dit bedrag had Rabobank moet{[}en{]} doen toekomen aan Mevriet: zij was immers eigenaar, hypotheekgever en daardoor rechthebbende op de koopsom nadat de schuld van Durisol aan Rabobank was voldaan''.

2.5. Enkele dagen later, Op 4 februari 2022, boekte Durisol € 275.000 op de rekening van Mevriet Vastgoed, en op 9 februari 2022 volgde een laatste overboeking van € 18.143,81 -- steeds met de vermelding ``Surplus afgelost bankkrediet Rabobank''. In totaal heeft Durisol dus binnen ongeveer een week tijd aan Mevriet Vastgoed een bedrag betaald dat gelijk is aan het onverschuldigd aan haarzelf door Rabo betaalde en voor Mevriet Vastgoed bestemde bedrag. De laatste betaling is gedaan na overleg met Rabo, die geen bezwaar maakte tegen deze ``rechtstreekse overboekingen''. Rabo heeft ook achteraf ingestemd met het standpunt van Durisol dat met de betalingen haar vordering uit onverschuldigde betaling op Durisol was afgelost. Enkele dagen na de laatste betaling heeft Rabo namelijk aan {[}geïntimeerde5{]} geschreven dat het klopt dat door verwerking van het royement het creditsaldo op de zakelijke rekening van Durisol aan Mevriet Vastgoed toebehoorde, en dat de overboekingen correct waren verwerkt.

2.6. Op 14 april 2022 heeft Durisol besloten eigen aangifte van faillissement te doen. Het faillissement is daarna op 2 mei 2022 uitgesproken.

2.7. De curator vordert nu van alle gedaagden (hoofdelijk) betaling van het door Rabo aan Durisol betaalde, en door die vennootschap aan Mevriet Vastgoed doorbetaalde bedrag van € 293.143,81, vermeerderd met rente en kosten. De rechtbank heeft deze vorderingen afgewezen. De bedoeling van het hoger beroep is dat alsnog toewijzing volgt.

\paragraph{3. Het oordeel van het hof}\label{het-oordeel-van-het-hof}

\subparagraph{Inleiding}\label{inleiding}

3.1. Het hof zal oordelen dat de bestreden beslissing in stand blijft. Dat wordt hierna uitgelegd. De bezwaren (grieven) zullen daarbij thematisch worden behandeld.

\subparagraph{Is sprake van een onverplichte betaling aan Mevriet Vastgoed als bedoeld in artikel 42 Fw of van een onverschuldigde betaling (artikel 6:203 BW)?}\label{is-sprake-van-een-onverplichte-betaling-aan-mevriet-vastgoed-als-bedoeld-in-artikel-42-fw-of-van-een-onverschuldigde-betaling-artikel-6203-bw}

3.2. Artikel 42 Fw biedt de curator de mogelijkheid rechtshandelingen te vernietigen als de gefailleerde schuldenaar die onverplicht heeft verricht, en de schuldenaar bij het verrichten van de rechtshandeling had moeten weten of behoren te weten dat benadeling van schuldeisers daarvan het gevolg zou zijn. De curator beroept zich op deze bepaling: Durisol had wel een schuld uit onverschuldigde betaling aan Rabo, maar was niet verplicht aan Mevriet Vastgoed te betalen. De overboekingen werden gedaan om het geld uit het vermogen van Durisol te trekken (``veilig te stellen''), teneinde te voorkomen dat een schuldeiser er beslag op zou leggen. Deze betalingen aan Mevriet Vastgoed waren onverschuldigd, en Durisol had moeten weten dat haar schuldeisers daardoor werden benadeeld.

3.3. Voor zover de curator doelt op benadeling van schuldeisers die door de transacties in hun beslagmogelijkheden werden gefrustreerd, is dat onterecht. Als de transacties niet hadden plaatsgehad en wel beslag was gelegd, dan zou dat beslag door het faillissement immers vervallen en zouden deze schuldeisers zich ook in de rij bij de andere schuldeisers moeten voegen. Het is bovendien niet de taak van de curator de belangen van een individuele schuldeiser te vertegenwoordigen, maar die van de gezamenlijke failissementscrediteuren.

3.4. Voor zover wordt gedoeld op de frustratie van de belangen van die gezamenlijke faillissementscrediteuren, volgt het hof de curator evenmin. Rabo had immers ten belope van de door Durisol verrichte betalingen een direct opeisbare vordering uit onverschuldigde betaling op Durisol. Daarnaast had Mevriet Vastgoed voor datzelfde bedrag een vordering op Rabo. Dat was duidelijk voor de indirecte bestuurders van beide vennootschappen. Hun lezing van de gang van zaken is, dat zij de fout van Rabo hebben willen rechtzetten door het surplus naar de rechthebbende partij over te hevelen1. Die uitleg komt het hof aannemelijk voor, omdat deze in overeenstemming is met de omschrijving die bij de betalingen is gegeven (``Surplus afgelost bankkrediet Rabobank'') en het overleg met Rabo dat aan de laatste tranche is voorafgegaan. Zo beschouwd, heeft Durisol gekozen voor een praktische oplossing van het door Rabo gecreëerde probleem: het gelijktijdig wegstrepen van twee vorderingen. Juridisch moet dat worden uitgelegd als de gelijktijdige betaling van de schuld van Rabo aan Mevriet Vastgoed en (vooral) voldoening van de even grote eigen, opeisbare schuld van Durisol aan Rabo. Hoewel Rabo hier ten tijde van de eerste betalingen van in totaal € 250.000 nog niet mee had ingestemd, heeft de bank de betaling van de schuld van Durisol aan haar wel meteen bekrachtigd. Dat blijkt uit het op 9 februari gevoerde overleg, de mededeling dat met deze gang van zaken de overboekingen correct waren verwerkt en het feit dat Rabo nadien geen verificatievordering in het faillissement van Durisol heeft ingediend. Bovendien is niet gesteld of gebleken dat Mevriet Vastgoed nog pretendeert voor het genoemde bedrag een vordering op Rabo te hebben.

3.5. Durisol heeft de betalingen dus niet onverplicht verricht. Aan die conclusie staat niet in de weg dat de adviseur van Mevriet c.s., Bredewoud, heeft opgemerkt dat het ontstane surplus op deze manier bij Mevriet vastgoed werd ``veiliggesteld'' om te voorkomen dat het door beslaglegging niet mogelijk zou zijn de voorgenomen correctie van de door Rabo gemaakte fout uit te voeren. Met die formulering kan in de door Mevriet c.s. gegeven lezing immers slechts zijn bedoeld dat de voorgenomen oplossing niet door beslaglegging kon worden gefrustreerd -- en niet dat Durisol onverplicht een betaling aan haar zustervennootschap deed. Voor dat laatste biedt het dossier (in weerwil van al het voorgaande) ook geen enkel aanknopingspunt.

3.6. Een en ander betekent dat het standpunt dat Durisol zonder rechtsgrond (lees ook: onverschuldigd) aan Mevriet Vastgoed heeft betaald, geen steun vindt in de feiten of in wat verder is aangevoerd: de curator heeft geen feiten of omstandigheden naar voren gebracht die -- indien ze vast zouden komen te staan -- tot een andere conclusie zouden kunnen leiden. Voor bewijsvoering ziet het hof om die reden geen ruimte.

\subparagraph{Is sprake van samenspanning (artikel 47 Fw) of onrechtmatig handelen (artikel 6:162 BW?}\label{is-sprake-van-samenspanning-artikel-47-fw-of-onrechtmatig-handelen-artikel-6162-bw}

3.7. Op grond van artikel 47 Fw kan de voldoening door de schuldenaar aan een opeisbare schuld onder meer worden vernietigd wanneer wordt aangetoond dat de betaling het gevolg was van overleg tussen de schuldenaar en de schuldeiser dat ten doel had de schuldeiser door die betaling boven andere schuldeisers te begunstigen2. De vraag of van benadeling sprake is, moet worden beantwoord door de hypothetische situatie waarin de schuldeisers zouden hebben verkeerd zonder de gewraakte rechtshandeling te vergelijken met de situatie waarin zij feitelijk verkeren als die handeling onaangetast blijft3. De enkele omstandigheid dat bij beide partijen wetenschap bestond van benadeling van overige schuldeisers is in dit laatste geval niet voldoende. Aan de gestelde voorwaarden is alleen in geval van samenspanning voldaan.

3.8. De curator doet een beroep op deze regeling voor het geval het hof van oordeel zou zijn dat met de betalingen de vordering van Rabo op Mevriet Vastgoed is voldaan. Hij gaat daarbij uit van samenspanning tussen Durisol en Mevriet Vastgoed, en voert aan dat (1) alle activiteiten van Durisol ten tijde van de betalingen al waren gestaakt en de activa waren verkocht, (2) de verkoopopbrengst niet volstond voor de betaling van alle crediteuren (de voorlopig erkende crediteurenlast beloopt € 2.52,452,08, bij een boedelsaldo van € 15.241,22), (3) het faillissement voorzienbaar was en (4) Durisol deze betalingen heeft verricht om te voorkomen dat beslag zou worden gelegd op de tegoeden die door de onverschuldigde betaling waren ontstaan.

3.9. Het hof kan niet de conclusie trekken dat onder de door de curator aangevoerde omstandigheden is voldaan aan de hiervoor genoemde vereisten voor paulianeuze samenspanning. Mevriet Vastgoed was immers geen schuldeiser van Durisol. Dat was Rabo. Daarop strandt het betoog van de curator al.~Bovendien is op geen enkele manier onderbouwd dat het de bedoeling van Durisol was om Mevriet Vastgoed met de betalingen te begunstigen ten opzichte van de schuldeisers van Durisol. Zowel de bedoeling als het gevolg van de betalingen was immers slechts dat een opeisbare vordering van een van die schuldeisers erdoor werd voldaan, en onder de geschetste omstandigheden stond het Durisol vrij om die schuld af te lossen. Het hof wijst er wat dat aangaat op dat zelfs een onmiskenbaar onverschuldigde betaling na faillissement in beginsel moet worden gerestorneerd.

3.10. Gelet op alles wat hiervoor is overwogen, kan er evenmin van worden uitgegaan dat Mevriet c.s. onrechtmatig hebben gehandeld, zoals de curator ook nog heeft aangevoerd.

\subparagraph{De conclusie}\label{de-conclusie}

3.11. Het hoger beroep slaagt niet. Omdat de curator in het ongelijk zal worden gesteld, zal het hof hem tot betaling van de proceskosten in hoger beroep veroordelen. Onder die kosten vallen ook de nakosten die nodig zijn voor de betekening van de uitspraak en de wettelijke rente daarover. De rente is verschuldigd vanaf veertien dagen na die betekening.

\paragraph{4. De beslissing}\label{de-beslissing-1}

Het hof:

bekrachtigt het vonnis van de rechtbank Overijssel, locatie Zwolle van 3 april 2024;

veroordeelt de curator tot betaling van de volgende proceskosten van Mevriet c.s.:

(\ldots; red.)

bepaalt dat al deze kosten moeten worden betaald binnen 14 dagen na vandaag. Als niet op tijd wordt betaald, dan worden die kosten verhoogd met de wettelijke rente;

verklaart de proceskostenveroordeling uitvoerbaar bij voorraad;

wijst af wat verder is gevorderd.

\newpage

\markboth{JOR 2026/137}{mistral/mistral-medium-3.5-128b:fp8 (k = 0)}

\subsection{k = 0}\label{k-0-1}

\subsubsection{Mistral}\label{mistral-3}

\texttt{mistral/mistral-medium-3.5-128b:fp8}

\begin{center}
\small
\begin{tabular}{@{}p{7.4cm}p{2.6cm}p{2.6cm}@{}}
\toprule
Variabele & Waarde (model) & Waarde (gold) \\
\midrule
Is rechtshandeling & \texttt{true} & \texttt{true} \\
Datum rechtshandeling & \texttt{2022-02-04} & \texttt{2022-02-04} \\
Verricht door de schuldenaar & \texttt{true} & \texttt{true} \\
Om niet & \texttt{false} & \texttt{false} \\
Datum faillietverklaring & \texttt{2022-05-02} & \texttt{2022-05-02} \\
Onverplicht & \textcolor{gray}{\texttt{true}} & \texttt{false} \\
Benadeling schuldeisers & \textcolor{gray}{\texttt{true}} & \texttt{false} \\
Wetenschap schuldenaar aangetoond & \texttt{false} & \texttt{false} \\
Wetenschap wederpartij aangetoond & \texttt{false} & \texttt{null} \\
Vermoeden wetenschap schuldenaar tegenbewijs & \texttt{null} & \texttt{null} \\
Vermoeden wetenschap wederpartij tegenbewijs & \texttt{null} & \texttt{null} \\
Tegen onderwaarde & \texttt{null} & \texttt{null} \\
Voldoening of zekerheid niet opeisbare schuld & \texttt{null} & \texttt{null} \\
Gelieerde partij & \textcolor{gray}{\texttt{true}} & \texttt{null} \\
Reeds verplicht voor aanvang termijn & \texttt{null} & \texttt{null} \\
\textbf{Vernietigbaar} & \textbf{\textcolor{gray}{\texttt{true}}} & \textbf{\texttt{false}} \\
\bottomrule
\end{tabular}
\end{center}

\paragraph{Is rechtshandeling - overwegingen}\label{is-rechtshandeling---overwegingen-9}

De overboekingen van Durisol naar Mevriet Vastgoed zijn op een rechtsgevolg gerichte wilsverklaringen (betalingen) en vallen daarmee onder het begrip rechtshandeling.

\paragraph{Verricht door de schuldenaar - overwegingen}\label{verricht-door-de-schuldenaar---overwegingen-9}

De overboekingen zijn verricht door Durisol, de latere gefailleerde, en dus door de schuldenaar zelf.

\paragraph{Om niet - overwegingen}\label{om-niet---overwegingen-9}

De betalingen zijn gedaan ter afwikkeling van een onverschuldigde betaling door Rabobank aan Durisol, waarbij Durisol het surplus (dat eigenlijk aan Mevriet Vastgoed toekwam) doorbetaalde. Dit impliceert een verband met een tegenprestatie (afwikkeling van de onverschuldigde betaling), waardoor de handeling niet `om niet' is.

\paragraph{Onverplicht - overwegingen}\label{onverplicht---overwegingen-9}

De betalingen van Durisol aan Mevriet Vastgoed (€ 275.000 op 4 februari 2022 en € 18.143,81 op 9 februari 2022) zijn onverplicht in de zin van artikel 42 Fw.

Allereerst was Durisol niet verplicht om het surplusbedrag (€ 293.143,81) aan Mevriet Vastgoed over te maken. Het surplus was door Rabobank ten onrechte aan Durisol betaald, terwijl het rechtens toekwam aan Mevriet Vastgoed als eigenaar en hypotheekgever van het verkochte onroerend goed. Durisol had geen contractuele of wettelijke verplichting om dit bedrag aan Mevriet Vastgoed te betalen. De verplichting tot betaling berustte niet op een overeenkomst of wet, maar was het gevolg van een onverschuldigde betaling door Rabobank aan Durisol.

Daarnaast was er geen sprake van een bestaande schuld van Durisol aan Mevriet Vastgoed die afgelost moest worden. De betalingen waren derhalve niet verplicht in de zin van artikel 42 Fw, omdat er geen rechtsplicht bestond om deze te verrichten. De curator heeft dan ook voldoende aannemelijk gemaakt dat de betalingen onverplicht zijn verricht.

\paragraph{Benadeling schuldeisers - overwegingen}\label{benadeling-schuldeisers---overwegingen-9}

In deze casus heeft Durisol, kort voor haar faillissement, een bedrag van € 293.143,81 (€ 275.000 + € 18.143,81) overgemaakt naar Mevriet Vastgoed. Dit bedrag betrof het surplus van de verkoopopbrengst van de met hypotheek bezwaarde activa van Mevriet Vastgoed, dat onverschuldigd door Rabobank aan Durisol was betaald. De overboekingen zijn gedaan met de vermelding `Surplus afgelost bankkrediet Rabobank', wat aangeeft dat het geld eigenlijk toebehoorde aan Mevriet Vastgoed.

Bij de beoordeling of de schuldeisers zijn benadeeld, dient een vermogensvergelijking te worden gemaakt. Zonder de overboekingen zou het bedrag van € 293.143,81 in de boedel van Durisol zijn gebleven, waar de schuldeisers zich op hadden kunnen verhalen. Door de overboekingen is dit bedrag echter uit de boedel van Durisol verdwenen en naar Mevriet Vastgoed gegaan. Dit leidt tot een vermindering van het actief van Durisol, wat de schuldeisers benadeelt.

Daarnaast is er sprake van een verstoring van de rangorde tussen de schuldeisers, omdat het bedrag specifiek ten goede is gekomen aan Mevriet Vastgoed (via de groepsvennootschap), terwijl andere schuldeisers hierdoor minder of niets ontvangen.

De benadeling is aanwezig op het moment waarop de curator zijn rechten doet gelden, omdat het bedrag niet meer in de boedel aanwezig is en de schuldeisers hierdoor een nadeel ondervinden.

Tot slot is er geen sprake van een situatie waarin de schuldeisers zich toch niet op het bedrag hadden kunnen verhalen, omdat het bedrag niet was belast met zekerheidsrechten ten gunste van andere schuldeisers. Het bedrag was immers het surplus na aflossing van de hypotheek van Rabobank, en had dus vrij in de boedel van Durisol moeten blijven.

\paragraph{Wetenschap - overwegingen}\label{wetenschap---overwegingen-9}

De betalingen van Durisol aan Mevriet Vastgoed (€ 275.000 op 4 februari 2022 en € 18.143,81 op 9 februari 2022) vonden plaats binnen een jaar voor de faillietverklaring (2 mei 2022). Mevriet Vastgoed is een groepsvennootschap van Durisol (art. 43 lid 1 sub 6° Fw), waardoor het bewijsvermoeden van art. 43 lid 1 Fw van toepassing is. Dit vermoeden geldt voor zowel Durisol als Mevriet Vastgoed, tenzij tegenbewijs wordt geleverd. Er is geen sprake van een eerdere verplichting (reeds\_verplicht\_voor\_aanvang\_termijn is niet van toepassing). De betalingen zijn niet om niet, dus art. 45 Fw is niet relevant. De curator hoeft derhalve niet rechtstreeks de wetenschap aan te tonen, tenzij Mevriet Vastgoed succesvol tegenbewijs levert.

\newpage

\markboth{JOR 2026/137}{gpt-5.5 (k = 0)}

\subsubsection{GPT}\label{gpt-3}

\texttt{gpt-5.5}

\begin{center}
\small
\begin{tabular}{@{}p{7.4cm}p{2.6cm}p{2.6cm}@{}}
\toprule
Variabele & Waarde (model) & Waarde (gold) \\
\midrule
Is rechtshandeling & \texttt{true} & \texttt{true} \\
Datum rechtshandeling & \texttt{2022-02-04} & \texttt{2022-02-04} \\
Verricht door de schuldenaar & \texttt{true} & \texttt{true} \\
Om niet & \texttt{false} & \texttt{false} \\
Datum faillietverklaring & \texttt{2022-05-02} & \texttt{2022-05-02} \\
Onverplicht & \texttt{false} & \texttt{false} \\
Benadeling schuldeisers & \texttt{false} & \texttt{false} \\
Wetenschap schuldenaar aangetoond & \texttt{false} & \texttt{false} \\
Wetenschap wederpartij aangetoond & \texttt{false} & \texttt{null} \\
Vermoeden wetenschap schuldenaar tegenbewijs & \texttt{false} & \texttt{null} \\
Vermoeden wetenschap wederpartij tegenbewijs & \texttt{false} & \texttt{null} \\
Tegen onderwaarde & \texttt{false} & \texttt{null} \\
Voldoening of zekerheid niet opeisbare schuld & \texttt{false} & \texttt{null} \\
Gelieerde partij & \textcolor{gray}{\texttt{true}} & \texttt{null} \\
Reeds verplicht voor aanvang termijn & \texttt{false} & \texttt{null} \\
\textbf{Vernietigbaar} & \textbf{\texttt{false}} & \textbf{\texttt{false}} \\
\bottomrule
\end{tabular}
\end{center}

\paragraph{Is rechtshandeling - overwegingen}\label{is-rechtshandeling---overwegingen-10}

De gewraakte handelingen bestaan uit overboekingen/betalingen door Durisol aan Mevriet Vastgoed. Daarmee werd beoogd een rechtsgevolg teweeg te brengen, namelijk betaling/doorbetaling van het door Rabobank aan Durisol overgemaakte surplus.

\paragraph{Verricht door de schuldenaar - overwegingen}\label{verricht-door-de-schuldenaar---overwegingen-10}

Durisol heeft zelf op 4 februari 2022 € 275.000 en op 9 februari 2022 € 18.143,81 overgeboekt naar Mevriet Vastgoed. Durisol is de later failliet verklaarde schuldenaar.

\paragraph{Om niet - overwegingen}\label{om-niet---overwegingen-10}

De betalingen waren volgens de feiten bedoeld als doorbetaling van het surplus dat aan Mevriet Vastgoed toebehoorde. Rabobank heeft achteraf ingestemd met het standpunt dat hiermee haar vordering uit onverschuldigde betaling op Durisol was afgelost. Daarmee stond tegenover de betalingen een daarmee verband houdend rechtsgevolg, zodat de handeling niet als om niet wordt aangemerkt.

\paragraph{Onverplicht - overwegingen}\label{onverplicht---overwegingen-10}

Voor de overboekingen bestond een in rechte afdwingbare rechtsplicht. Rabobank had het surplus na royement van de hypotheek abusievelijk aan Durisol overgemaakt, terwijl dit bedrag aan Mevriet Vastgoed toekwam. Daardoor had Rabobank jegens Durisol een vordering uit onverschuldigde betaling, althans moest Durisol het niet aan haar toekomende bedrag doorgeleiden. De betalingen aan Mevriet Vastgoed zijn gedaan ter correctie van die fout, met de omschrijving ``Surplus afgelost bankkrediet Rabobank'', na overleg met Rabobank en met haar instemming. Rabobank heeft ook bevestigd dat de overboekingen correct waren verwerkt en dat daarmee haar vordering uit onverschuldigde betaling op Durisol was afgelost. Dat Durisol niet rechtstreeks aan Rabobank terugbetaalde maar, met instemming van Rabobank, aan Mevriet Vastgoed betaalde, maakt de uitvoering van de bestaande verplichting niet onverplicht. Een bestaande verplichting kan immers op een niet-voorgeschreven wijze worden uitgevoerd zonder dat de uitvoering daardoor onverplicht wordt.

\paragraph{Benadeling schuldeisers - overwegingen}\label{benadeling-schuldeisers---overwegingen-10}

Bij de vermogensvergelijking moet de feitelijke situatie na de overboekingen worden vergeleken met de reële hypothetische situatie zonder die overboekingen. Het bedrag van € 293.143,81 was niet afkomstig uit eigen vermogen van Durisol, maar betrof het restant van de verkoopopbrengst van aan Mevriet Vastgoed toebehorende, met hypotheek bezwaarde activa. Rabobank heeft dit bedrag per vergissing aan Durisol overgemaakt, terwijl het na aflossing van de debetstand van Durisol aan Mevriet Vastgoed had moeten worden uitgekeerd. Door de doorbetaling aan Mevriet Vastgoed is die fout feitelijk hersteld en is een daarmee corresponderende restitutie-/correctieverplichting van Durisol jegens Rabobank dan wel Mevriet Vastgoed weggevallen. Zonder de betalingen zou niet aannemelijk zijn dat dit bedrag vrij voor verhaal door de gezamenlijke schuldeisers van Durisol beschikbaar zou zijn gebleven; Rabobank had het onverschuldigd betaalde bedrag kunnen terugvorderen of verrekenen/corrigeren. De schuldeisers van Durisol zijn daardoor niet daadwerkelijk in hun verhaalspositie verslechterd, maar missen hoogstens een voordeel dat zou zijn ontstaan door de vergissing van Rabobank. Ook is geen relevante rangordeverstoring aannemelijk, omdat het niet gaat om betaling uit vrij beschikbaar vermogen van Durisol aan een gewone schuldeiser ten koste van de overige schuldeisers. Daarom is geen sprake van benadeling van schuldeisers in de zin van artikel 42 Fw.

\paragraph{Wetenschap - overwegingen}\label{wetenschap---overwegingen-10}

Uit de gegeven feiten volgt niet rechtstreeks dat Durisol op 4 en 9 februari 2022 wist of behoorde te weten dat faillissement en een tekort daarin met een redelijke mate van waarschijnlijkheid te voorzien waren. Dat het faillissement enkele maanden later is aangevraagd en uitgesproken, is daarvoor op zichzelf onvoldoende. Ook aan de zijde van Mevriet Vastgoed is rechtstreekse wetenschap niet concreet aangetoond. Wel is sprake van een betaling aan een gelieerde/groepsvennootschap: Durisol en Mevriet Vastgoed behoren tot dezelfde vennootschappelijke structuur met dezelfde achterliggende bestuurders/aandeelhouders. Daarmee is, indien de overige voorwaarden voor art. 43 Fw zijn vervuld, het bewijsvermoeden van wetenschap aan beide zijden van toepassing. Er is geen aanmerkelijke onderwaarde en evenmin blijkt van voldoening of zekerheidstelling voor een niet-opeisbare schuld. Uit de casus blijkt ook niet dat Durisol zich al vóór de aanvang van de jaartermijn tot deze concrete overboekingen had verplicht. De omstandigheid dat Rabobank de doorboeking als correctie van een onjuiste betaling beschouwde, levert in deze feitenopstelling onvoldoende tegenbewijs op tegen de veronderstelde voorzienbaarheid van faillissement en tekort.

\newpage

\markboth{JOR 2026/137}{gemini-3.5-flash (k = 0)}

\subsubsection{Gemini}\label{gemini-3}

\texttt{gemini-3.5-flash}

\begin{center}
\small
\begin{tabular}{@{}p{7.4cm}p{2.6cm}p{2.6cm}@{}}
\toprule
Variabele & Waarde (model) & Waarde (gold) \\
\midrule
Is rechtshandeling & \texttt{true} & \texttt{true} \\
Datum rechtshandeling & \texttt{2022-02-04} & \texttt{2022-02-04} \\
Verricht door de schuldenaar & \texttt{true} & \texttt{true} \\
Om niet & \texttt{false} & \texttt{false} \\
Datum faillietverklaring & \texttt{2022-05-02} & \texttt{2022-05-02} \\
Onverplicht & \textcolor{gray}{\texttt{true}} & \texttt{false} \\
Benadeling schuldeisers & \texttt{false} & \texttt{false} \\
Wetenschap schuldenaar aangetoond & \texttt{false} & \texttt{false} \\
Wetenschap wederpartij aangetoond & \texttt{false} & \texttt{null} \\
Vermoeden wetenschap schuldenaar tegenbewijs & \textcolor{gray}{\texttt{true}} & \texttt{null} \\
Vermoeden wetenschap wederpartij tegenbewijs & \textcolor{gray}{\texttt{true}} & \texttt{null} \\
Tegen onderwaarde & \texttt{null} & \texttt{null} \\
Voldoening of zekerheid niet opeisbare schuld & \texttt{null} & \texttt{null} \\
Gelieerde partij & \textcolor{gray}{\texttt{true}} & \texttt{null} \\
Reeds verplicht voor aanvang termijn & \texttt{false} & \texttt{null} \\
\textbf{Vernietigbaar} & \textbf{\texttt{false}} & \textbf{\texttt{false}} \\
\bottomrule
\end{tabular}
\end{center}

\paragraph{Is rechtshandeling - overwegingen}\label{is-rechtshandeling---overwegingen-11}

De overboekingen van in totaal € 293.143,81 door Durisol aan Mevriet Vastgoed zijn rechtshandelingen, aangezien deze betalingen gericht zijn op het teweegbrengen van een rechtsgevolg, namelijk het overdragen van vermogensbestanddelen ter voldoening van een vermeende schuld.

\paragraph{Verricht door de schuldenaar - overwegingen}\label{verricht-door-de-schuldenaar---overwegingen-11}

De overboekingen zijn verricht door Durisol, de schuldenaar die naderhand failliet is verklaard.

\paragraph{Om niet - overwegingen}\label{om-niet---overwegingen-11}

De rechtshandeling is anders dan om niet (om baat). De betalingen vonden plaats ter aflossing van een vordering uit onverschuldigde betaling c.q. ter doorstorting van gelden die Rabobank per abuis aan Durisol had overgemaakt maar die toebehoorden aan Mevriet Vastgoed. Er stond derhalve een prestatie (bevrijding van een schuld/verplichting) tegenover.

\paragraph{Onverplicht - overwegingen}\label{onverplicht---overwegingen-11}

De betalingen door Durisol aan Mevriet Vastgoed zijn onverplicht verricht. Een rechtshandeling is onverplicht wanneer daartoe geen op de wet of een overeenkomst rustende rechtsplicht bestond. In dit geval bestond er geen directe rechtsverhouding tussen Durisol en Mevriet Vastgoed die Durisol verplichtte tot betaling van het surplus aan Mevriet Vastgoed. De rechtsplicht tot terugbetaling van het onverschuldigd ontvangen surplus rustte op Durisol jegens Rabobank (op grond van artikel 6:203 BW), terwijl Mevriet Vastgoed een vordering had op Rabobank. Hoewel Rabobank achteraf heeft ingestemd met de rechtstreekse betalingen ter afwikkeling van de respectievelijke vorderingen, bestond er op het moment van de betalingen geen in rechte afdwingbare verplichting voor Durisol om aan Mevriet Vastgoed te betalen. De achteraf verleende instemming of bekrachtiging door de bank maakt de betalingen niet achteraf alsnog verplicht op het moment van verrichten.

\paragraph{Benadeling schuldeisers - overwegingen}\label{benadeling-schuldeisers---overwegingen-11}

Voor de vraag of er sprake is van benadeling van schuldeisers in de zin van artikel 42 Fw, moet een vergelijking worden gemaakt tussen de feitelijke situatie waarin de schuldeisers verkeren na de betalingen en de hypothetische situatie waarin zij zouden hebben verkeerd als de betalingen achterwege waren gebleven. In het hypothetische scenario waarin de foutieve overboeking door de Rabobank niet had plaatsgevonden, zou de Rabobank het surplus van de verkoopopbrengst rechtstreeks aan Mevriet Vastgoed (als eigenaar en hypotheekgever van het verkochte pand) hebben overgemaakt. Durisol zou dit bedrag dan nooit hebben ontvangen, en de schuldeisers van Durisol zouden zich er niet op hebben kunnen verhalen. Door de foutieve overboeking te herstellen en het surplus alsnog door te boeken naar de werkelijke rechthebbende (Mevriet Vastgoed), is de feitelijke situatie materieel gelijk aan de hypothetische situatie waarin de fout achterwege was gebleven. Het vermogen van Durisol is per saldo niet verminderd ten opzichte van de situatie waarin de rechtshandeling niet was verricht, aangezien Durisol geen recht had op dit surplus. Er is derhalve geen sprake van benadeling van de schuldeisers.

\paragraph{Wetenschap - overwegingen}\label{wetenschap---overwegingen-11}

De betalingen door Durisol aan Mevriet Vastgoed vonden plaats binnen een jaar voor het faillissement en tussen gelieerde partijen (groepsmaatschappijen), waardoor het bewijsvermoeden van art. 43 lid 1 sub 5°/6° Fw in beginsel van toepassing is. Echter, de betalingen betroffen de correctie van een foutieve overboeking door de Rabobank van een surplus dat rechtens en economisch aan Mevriet Vastgoed toebehoorde. Hierdoor is het vermoeden van wetenschap van benadeling aan beide zijden met succes ontzenuwd; ten tijde van de correctieboeking was een faillissement en een tekort daarin niet met een redelijke mate van waarschijnlijkheid te voorzien als gevolg van deze handeling, aangezien het vermogen van Durisol per saldo niet nadelig werd beïnvloed (er stond een gelijke schuld uit onverschuldigde betaling tegenover). Rechtstreekse wetenschap van benadeling is evenmin aangetoond.

\newpage

\markboth{JOR 2026/137}{mistral/mistral-medium-3.5-128b:fp8 (k = 5)}

\subsection{k = 5}\label{k-5-1}

\subsubsection{Mistral}\label{mistral-4}

\texttt{mistral/mistral-medium-3.5-128b:fp8}

\begin{center}
\small
\begin{tabular}{@{}p{7.4cm}p{2.6cm}p{2.6cm}@{}}
\toprule
Variabele & Waarde (model) & Waarde (gold) \\
\midrule
Is rechtshandeling & \texttt{true} & \texttt{true} \\
Datum rechtshandeling & \texttt{2022-02-04} & \texttt{2022-02-04} \\
Verricht door de schuldenaar & \texttt{true} & \texttt{true} \\
Om niet & \texttt{false} & \texttt{false} \\
Datum faillietverklaring & \texttt{2022-05-02} & \texttt{2022-05-02} \\
Onverplicht & \textcolor{gray}{\texttt{true}} & \texttt{false} \\
Benadeling schuldeisers & \textcolor{gray}{\texttt{true}} & \texttt{false} \\
Wetenschap schuldenaar aangetoond & \texttt{false} & \texttt{false} \\
Wetenschap wederpartij aangetoond & \texttt{false} & \texttt{null} \\
Vermoeden wetenschap schuldenaar tegenbewijs & \texttt{null} & \texttt{null} \\
Vermoeden wetenschap wederpartij tegenbewijs & \texttt{null} & \texttt{null} \\
Tegen onderwaarde & \texttt{null} & \texttt{null} \\
Voldoening of zekerheid niet opeisbare schuld & \texttt{null} & \texttt{null} \\
Gelieerde partij & \textcolor{gray}{\texttt{true}} & \texttt{null} \\
Reeds verplicht voor aanvang termijn & \texttt{null} & \texttt{null} \\
\textbf{Vernietigbaar} & \textbf{\textcolor{gray}{\texttt{true}}} & \textbf{\texttt{false}} \\
\bottomrule
\end{tabular}
\end{center}

\paragraph{Is rechtshandeling - overwegingen}\label{is-rechtshandeling---overwegingen-12}

De overboekingen van Durisol naar Mevriet Vastgoed zijn op een rekening gedaan, wat een rechtshandeling is (goederenrechtelijke handeling).

\paragraph{Verricht door de schuldenaar - overwegingen}\label{verricht-door-de-schuldenaar---overwegingen-12}

De overboekingen zijn verricht door Durisol, de latere gefailleerde, zelf.

\paragraph{Om niet - overwegingen}\label{om-niet---overwegingen-12}

De overboekingen zijn gedaan ter afwikkeling van een onverschuldigde betaling door Rabobank aan Durisol, wat een reële tegenprestatie (afwikkeling schuld) inhoudt. Derhalve is de handeling `anders dan om niet'.

\paragraph{Onverplicht - overwegingen}\label{onverplicht---overwegingen-12}

De betalingen door Durisol aan Mevriet Vastgoed (€ 275.000 op 4 februari 2022 en € 18.143,81 op 9 februari 2022) zijn onverplicht verricht in de zin van artikel 42 lid 1 Fw.

Allereerst was er geen rechtsplicht voor Durisol om deze bedragen aan Mevriet Vastgoed te betalen. Het surplus van € 349.981,74 was door Rabobank ten onrechte aan Durisol overgemaakt, terwijl dit bedrag toebehoorde aan Mevriet Vastgoed als eigenaar en hypotheekgever van het bedrijfsterrein (zie 2.4). Rabobank heeft achteraf bevestigd dat het creditsaldo op de rekening van Durisol aan Mevriet Vastgoed toebehoorde (2.5). Dit betekent dat Durisol geen contractuele of wettelijke verplichting had om dit geld aan Mevriet Vastgoed over te maken; het was immers niet haar eigen geld.

Daarnaast was er geen sprake van een bestaande schuld of verplichting tussen Durisol en Mevriet Vastgoed die deze betalingen rechtvaardigde. De betalingen waren derhalve niet gebaseerd op een op de wet of een overeenkomst berustende rechtsplicht. Morele, feitelijke of economische druk (zoals de relatie binnen de groepsvennootschappen) schept geen rechtsplicht (vgl. ECLI:NL:GHSHE:2017:4528, rov. 3.5.4).

Tot slot was de uitvoering van de betalingen (doorboeken van het surplus) geen verplichte nakoming van een bestaande verbintenis, maar een handeling die voortvloeide uit de onjuiste toewijzing van het geld door Rabobank. De curator heeft derhalve voldoende aannemelijk gemaakt dat de betalingen onverplicht zijn verricht.

\paragraph{Benadeling schuldeisers - overwegingen}\label{benadeling-schuldeisers---overwegingen-12}

In de onderhavige casus heeft Durisol op 4 en 9 februari 2022 in totaal € 293.143,81 overgemaakt naar Mevriet Vastgoed, met de vermelding `Surplus afgelost bankkrediet Rabobank'. Dit bedrag betrof het onverschuldigd door Rabobank aan Durisol betaalde surplus van € 349.981,74, dat eigenlijk aan Mevriet Vastgoed toekwam (ECLI:NL:GHSHE:2014:4568, rov. 3.7).

De benadeling van de schuldeisers moet worden beoordeeld door een vermogensvergelijking: de hypothetische situatie zonder de gewraakte rechtshandeling versus de feitelijke situatie. Zonder de overboekingen zou het bedrag van € 293.143,81 in de boedel van Durisol zijn gebleven en ter beschikking hebben gestaan voor de gezamenlijke schuldeisers. Door de overboekingen is dit bedrag aan Mevriet Vastgoed toegekomen, waardoor het actief van Durisol is verminderd en de schuldeisers in het algemeen zijn benadeeld.

De overboekingen waren onverplicht in de zin van artikel 42 lid 1 Fw, omdat Durisol geen verplichting had om dit bedrag aan Mevriet Vastgoed te betalen. Het surplus was immers onverschuldigd aan Durisol betaald door Rabobank en had derhalve in de boedel moeten blijven. Bovendien was Durisol op het moment van de overboekingen (februari 2022) al in financiële moeilijkheden, zoals blijkt uit het feit dat Durisol op 14 april 2022 eigen aangifte van faillissement heeft gedaan en het faillissement op 2 mei 2022 is uitgesproken. Dit duidt erop dat Durisol wist of behoorde te weten dat de overboekingen benadeling van de schuldeisers tot gevolg zouden hebben (ECLI:NL:RBAMS:2023:897, rov. 1.12).

Daarnaast is Mevriet Vastgoed, als ontvanger van de betalingen, een groepsvennootschap van Durisol, en is het aannemelijk dat ook Mevriet Vastgoed wist of behoorde te weten dat de overboekingen benadelend waren voor de schuldeisers van Durisol (artikel 42 lid 2 Fw).

Er is dus sprake van een vermindering van het actief van Durisol ten nadele van de gezamenlijke schuldeisers, wat voldoende is om benadeling aan te nemen (ECLI:NL:GHSHE:2017:4528, rov. 3.7.4).

\paragraph{Wetenschap - overwegingen}\label{wetenschap---overwegingen-12}

De betalingen van Durisol aan Mevriet Vastgoed (€ 275.000 op 4 februari 2022 en € 18.143,81 op 9 februari 2022) vonden plaats binnen een jaar voor de faillietverklaring (2 mei 2022) en betroffen een groepsmaatschappij (art. 43 lid 1 sub 6° Fw). Hierdoor wordt de wetenschap van benadeling aan beide zijden (Durisol en Mevriet Vastgoed) vermoed. Er is geen sprake van een reeds bestaande verplichting vóór de jaartermijn (de betalingen waren ad hoc en niet gebaseerd op een eerdere verbintenis). Het vermoeden van art. 43 Fw is dus van toepassing. Er is geen tegenbewijs geleverd dat het faillissement en een tekort daarin niet met een redelijke mate van waarschijnlijkheid waren te voorzien. De betalingen waren onverplicht (geen contractuele of wettelijke grondslag) en leidden tot benadeling van de schuldeisers, omdat het geld anders in de boedel zou zijn gevallen. De curator hoeft derhalve de wetenschap niet rechtstreeks aan te tonen.

\newpage

\markboth{JOR 2026/137}{gpt-5.5 (k = 5)}

\subsubsection{GPT}\label{gpt-4}

\texttt{gpt-5.5}

\begin{center}
\small
\begin{tabular}{@{}p{7.4cm}p{2.6cm}p{2.6cm}@{}}
\toprule
Variabele & Waarde (model) & Waarde (gold) \\
\midrule
Is rechtshandeling & \texttt{true} & \texttt{true} \\
Datum rechtshandeling & \texttt{2022-02-04} & \texttt{2022-02-04} \\
Verricht door de schuldenaar & \texttt{true} & \texttt{true} \\
Om niet & \texttt{false} & \texttt{false} \\
Datum faillietverklaring & \texttt{2022-05-02} & \texttt{2022-05-02} \\
Onverplicht & \texttt{false} & \texttt{false} \\
Benadeling schuldeisers & \texttt{false} & \texttt{false} \\
Wetenschap schuldenaar aangetoond & \texttt{false} & \texttt{false} \\
Wetenschap wederpartij aangetoond & \texttt{false} & \texttt{null} \\
Vermoeden wetenschap schuldenaar tegenbewijs & \texttt{false} & \texttt{null} \\
Vermoeden wetenschap wederpartij tegenbewijs & \texttt{false} & \texttt{null} \\
Tegen onderwaarde & \texttt{null} & \texttt{null} \\
Voldoening of zekerheid niet opeisbare schuld & \texttt{null} & \texttt{null} \\
Gelieerde partij & \textcolor{gray}{\texttt{true}} & \texttt{null} \\
Reeds verplicht voor aanvang termijn & \texttt{false} & \texttt{null} \\
\textbf{Vernietigbaar} & \textbf{\texttt{false}} & \textbf{\texttt{false}} \\
\bottomrule
\end{tabular}
\end{center}

\paragraph{Is rechtshandeling - overwegingen}\label{is-rechtshandeling---overwegingen-13}

De gewraakte handelingen zijn de overboekingen door Durisol van € 275.000 op 4 februari 2022 en € 18.143,81 op 9 februari 2022 aan Mevriet Vastgoed, met de vermelding ``Surplus afgelost bankkrediet Rabobank''. Deze betalingen waren gericht op rechtsgevolg, namelijk betaling/doorbetaling van het door Rabobank aan Durisol betaalde surplus en volgens Rabobank aflossing van de vordering uit onverschuldigde betaling op Durisol.

\paragraph{Verricht door de schuldenaar - overwegingen}\label{verricht-door-de-schuldenaar---overwegingen-13}

De overboekingen zijn gedaan door Durisol vanaf haar rekening. Durisol is de schuldenaar die later failliet is verklaard.

\paragraph{Om niet - overwegingen}\label{om-niet---overwegingen-13}

De betalingen zijn niet zonder enige daarmee verband houdende tegenprestatie gedaan. Uit de feiten volgt dat het bedrag door Rabobank ten onrechte aan Durisol was betaald en voor Mevriet Vastgoed bestemd was, en dat Rabobank achteraf heeft ingestemd met het standpunt dat door de betalingen haar vordering uit onverschuldigde betaling op Durisol was afgelost. De rechtshandeling is daarom anders dan om niet.

\paragraph{Onverplicht - overwegingen}\label{onverplicht---overwegingen-13}

De betalingen van Durisol aan Mevriet Vastgoed moeten worden beoordeeld tegen de achtergrond dat Rabobank het netto-surplus van de verkoopopbrengst, na voldoening van de schuld van Durisol aan Rabobank, per vergissing aan Durisol heeft overgemaakt terwijl dit bedrag aan Mevriet Vastgoed toekwam als eigenaar/hypotheekgever/rechthebbende op de koopsom. Durisol beschikte daarmee niet over een vrij aan haar toekomend actief, maar over een bedrag dat zij zonder rechtsgrond had ontvangen en dat moest worden teruggeleid naar de rechthebbende dan wel, bezien vanuit Rabobank, moest worden terugbetaald wegens onverschuldigde betaling. Die verplichting was in rechte afdwingbaar. Dat Durisol het bedrag niet aan Rabobank terugbetaalde maar rechtstreeks aan Mevriet Vastgoed overmaakte, maakt de uitvoering niet onverplicht: een bestaande verplichting kan op een niet oorspronkelijk voorgeschreven wijze worden uitgevoerd zonder dat de betaling daardoor onverplicht wordt, temeer nu Rabobank met deze rechtstreekse overboekingen instemde en bevestigde dat daarmee haar vordering uit onverschuldigde betaling op Durisol was afgelost. De situatie vertoont overeenkomst met de doorbetalingszaken waarin ontvangen bedragen die materieel aan een ander toekwamen op grond van een rechtsverhouding moesten worden doorgestort, zoals in ECLI:NL:GHLEE:2007:BB3627. De curator kan daarom niet volstaan met de stelling dat de betaling aan een groepsvennootschap plaatsvond; uit de feiten volgt juist een rechtsgrond en een afdwingbare verplichting tot afdracht van het surplus.

\paragraph{Benadeling schuldeisers - overwegingen}\label{benadeling-schuldeisers---overwegingen-13}

Bij de vermogensvergelijking moet niet alleen worden gekeken naar het feit dat € 293.143,81 van de rekening van Durisol is afgeboekt, maar ook naar de herkomst en bestemming van dat bedrag en naar wat zonder die betalingen reëel zou zijn gebeurd. Het bedrag betrof het door Rabobank per vergissing aan Durisol uitgekeerde surplus dat na voldoening van Durisols bankschuld toekwam aan Mevriet Vastgoed als eigenaar/hypotheekgever van de verkochte onroerende zaken. Door de doorbetaling aan Mevriet Vastgoed is weliswaar het banksaldo van Durisol verminderd, maar daar stond tegenover dat Rabobank haar vordering uit onverschuldigde betaling op Durisol als voldaan beschouwde. Gelet op de feiten --- Rabobank had het bedrag oorspronkelijk aan Mevriet Vastgoed willen uitkeren, heeft de vergissing onderkend, geen bezwaar gemaakt tegen de rechtstreekse overboekingen en deze achteraf als correcte verwerking bevestigd --- is het reële hypothetische scenario dat Rabobank het bedrag zonder de doorbetalingen zou hebben teruggehaald/gecorrigeerd of verrekend en alsnog aan Mevriet Vastgoed zou hebben doen toekomen. De schuldeisers van Durisol hadden zich dus niet daadwerkelijk op dit surplus kunnen verhalen. Dit verschilt van gevallen zoals ECLI:NL:GHSHE:2014:4568, waarin een vóór faillissement ontvangen onverschuldigde betaling door doorbetaling aan een derde leidde tot vermindering van het actief tegenover slechts het wegvallen van een concurrente vordering. Hier is geen reële vermindering van verhaalsmogelijkheden of verstoring van de rangorde aannemelijk.

\paragraph{Wetenschap - overwegingen}\label{wetenschap---overwegingen-13}

Rechtstreekse wetenschap van benadeling is op de gegeven feiten niet zonder meer aangetoond: ten tijde van de betalingen op 4 en 9 februari 2022 was het faillissement nog niet aangevraagd en het besluit tot eigen aangifte volgde pas op 14 april 2022; concrete gegevens waaruit blijkt dat op de betaaldatums een faillissement met tekort reeds met redelijke mate van waarschijnlijkheid voorzienbaar was, ontbreken. Dit wijkt af van ECLI:NL:GHSHE:2014:4568, waar het faillissement al was aangevraagd. Wel gaat het om betalingen binnen de jaartermijn aan een groeps-/gelieerde vennootschap, zodat bij toepassing van art. 43 lid 1 sub 5/6 Fw de wetenschap van benadeling aan beide zijden wordt vermoed, voor zover de rechtshandeling benadelend en onverplicht is. De verplichting tot doorbetaling van het door Rabobank abusievelijk aan Durisol betaalde surplus is in elk geval niet al vóór aanvang van de jaartermijn ontstaan. Uit de casus volgt geen geslaagd tegenbewijs dat faillissement en tekort op dat moment niet voorzienbaar waren; de omstandigheid dat de betalingen de bankvergissing corrigeerden ziet vooral op rechtsgrond/benadeling, niet op het ontzenuwen van het voorzienbaarheidsvermoeden.

\newpage

\markboth{JOR 2026/137}{gemini-3.5-flash (k = 5)}

\subsubsection{Gemini}\label{gemini-4}

\texttt{gemini-3.5-flash}

\begin{center}
\small
\begin{tabular}{@{}p{7.4cm}p{2.6cm}p{2.6cm}@{}}
\toprule
Variabele & Waarde (model) & Waarde (gold) \\
\midrule
Is rechtshandeling & \texttt{true} & \texttt{true} \\
Datum rechtshandeling & \texttt{2022-02-04} & \texttt{2022-02-04} \\
Verricht door de schuldenaar & \texttt{true} & \texttt{true} \\
Om niet & \texttt{false} & \texttt{false} \\
Datum faillietverklaring & \texttt{2022-05-02} & \texttt{2022-05-02} \\
Onverplicht & \texttt{false} & \texttt{false} \\
Benadeling schuldeisers & \textcolor{gray}{\texttt{true}} & \texttt{false} \\
Wetenschap schuldenaar aangetoond & \textcolor{gray}{\texttt{true}} & \texttt{false} \\
Wetenschap wederpartij aangetoond & \textcolor{gray}{\texttt{true}} & \texttt{null} \\
Vermoeden wetenschap schuldenaar tegenbewijs & \texttt{false} & \texttt{null} \\
Vermoeden wetenschap wederpartij tegenbewijs & \texttt{false} & \texttt{null} \\
Tegen onderwaarde & \texttt{null} & \texttt{null} \\
Voldoening of zekerheid niet opeisbare schuld & \texttt{null} & \texttt{null} \\
Gelieerde partij & \textcolor{gray}{\texttt{true}} & \texttt{null} \\
Reeds verplicht voor aanvang termijn & \texttt{false} & \texttt{null} \\
\textbf{Vernietigbaar} & \textbf{\texttt{false}} & \textbf{\texttt{false}} \\
\bottomrule
\end{tabular}
\end{center}

\paragraph{Is rechtshandeling - overwegingen}\label{is-rechtshandeling---overwegingen-14}

De overboekingen van in totaal € 293.143,81 door Durisol aan Mevriet Vastgoed zijn rechtshandelingen, aangezien deze betalingen gericht zijn op het teweegbrengen van een rechtsgevolg, namelijk het overdragen van vermogensbestanddelen ter voldoening van een vermeende schuld.

\paragraph{Verricht door de schuldenaar - overwegingen}\label{verricht-door-de-schuldenaar---overwegingen-14}

De overboekingen zijn verricht door Durisol, de schuldenaar die naderhand failliet is verklaard.

\paragraph{Om niet - overwegingen}\label{om-niet---overwegingen-14}

De rechtshandeling is anders dan om niet (om baat). De betalingen vonden plaats ter aflossing van een vordering uit onverschuldigde betaling c.q. ter doorstorting van gelden die Rabobank per abuis aan Durisol had overgemaakt maar die toebehoorden aan Mevriet Vastgoed. Er stond derhalve een prestatie (bevrijding van een schuld/verplichting) tegenover.

\paragraph{Onverplicht - overwegingen}\label{onverplicht---overwegingen-14}

Durisol ontving per abuis het surplus van de verkoopopbrengst van het bedrijfsterrein van Mevriet Vastgoed van de Rabobank. Hierdoor ontstond voor Durisol een in rechte afdwingbare, wettelijke verplichting (uit onverschuldigde betaling, art. 6:203 BW) om dit bedrag aan de Rabobank terug te betalen. Rabobank was op haar beurt verplicht dit surplus aan Mevriet Vastgoed te betalen. Door het surplus rechtstreeks aan Mevriet Vastgoed over te boeken, heeft Durisol haar schuld aan Rabobank voldaan, waarmee Rabobank ook expliciet heeft ingestemd. Hoewel de betaling aan een derde (Mevriet Vastgoed) is verricht, betreft dit de nakoming van een bestaande, opeisbare rechtsplicht van Durisol jegens Rabobank. Het feit dat een bestaande verplichting op een niet-voorgeschreven wijze wordt uitgevoerd (namelijk door rechtstreekse betaling aan de uiteindelijke rechthebbende in plaats van aan de schuldeiser zelf), maakt de rechtshandeling niet onverplicht.

\paragraph{Benadeling schuldeisers - overwegingen}\label{benadeling-schuldeisers---overwegingen-14}

In deze casus heeft Rabobank per vergissing een surplus van € 349.981,74, dat bestemd was voor Mevriet Vastgoed, overgemaakt naar de rekening van Durisol. Durisol heeft vervolgens in totaal € 293.143,81 doorgeboekt naar Mevriet Vastgoed. Indien deze betalingen achterwege waren gebleven, had dit bedrag zich ten tijde van het faillissement nog in het vermogen van Durisol bevonden en had het tot het actieve boedelvolume behoord waarop de gezamenlijke schuldeisers zich konden verhalen. De vordering die Mevriet Vastgoed (of Rabobank) op Durisol zou hebben wegens de onverschuldigde betaling, betreft een vordering die vóór de faillietverklaring is ontstaan. Overeenkomstig de jurisprudentie (zoals ECLI:NL:GHSHE:2014:4568) levert een onverschuldigde betaling vóór faillissement slechts een concurrente faillissementsvordering op. Door de feitelijke doorbetaling is het actief van de boedel verminderd en is Mevriet Vastgoed volledig voldaan, terwijl zij anders als concurrente schuldeiser zou hebben meegedeeld in het faillissement. Het tenietgaan van deze concurrente schuld levert de gezamenlijke schuldeisers geen voordeel op. Er is derhalve sprake van een vermindering van het actief en een verstoring van de rangorde (paritas creditorum), hetgeen benadeling van de schuldeisers in de zin van artikel 42 Fw oplevert.

\paragraph{Wetenschap - overwegingen}\label{wetenschap---overwegingen-14}

De betalingen van in totaal € 293.143,81 door Durisol aan haar groepsvennootschap Mevriet Vastgoed vonden plaats in februari 2022, binnen een jaar voor de faillietverklaring op 2 mei 2022. Omdat de betalingen zijn verricht aan een groepsmaatschappij, is het wettelijk bewijsvermoeden van artikel 43 lid 1 sub 6° Fw van toepassing. Dit vermoeden houdt in dat de wetenschap van benadeling aan beide zijden wordt vermoed te bestaan. De rechtshandeling was onverplicht; hoewel Rabobank de foutieve betaling aan Durisol deed, creëerde dit slechts een concurrente vordering uit onverschuldigde betaling (analoog aan ECLI:NL:GHSHE:2014:4568), waardoor de directe doorbetaling aan Mevriet Vastgoed de overige crediteuren benadeelde. Het tegenbewijs om dit vermoeden te ontzenuwen slaagt niet: ten tijde van de betalingen in februari 2022 was het faillissement en een tekort daarin met een redelijke mate van waarschijnlijkheid te voorzien, aangezien de activa reeds in januari 2022 waren verkocht en het eigen faillissement in april 2022 werd aangevraagd. Bovendien is de wetenschap ook rechtstreeks aangetoond, nu de betrokken bestuurders via de holdingstructuur volledige wetenschap hadden van de penibele financiële situatie van Durisol.

\newpage

\markboth{JOR 2026/137}{mistral/mistral-medium-3.5-128b:fp8 (k = 10)}

\subsection{k = 10}\label{k-10-1}

\subsubsection{Mistral}\label{mistral-5}

\texttt{mistral/mistral-medium-3.5-128b:fp8}

\begin{center}
\small
\begin{tabular}{@{}p{7.4cm}p{2.6cm}p{2.6cm}@{}}
\toprule
Variabele & Waarde (model) & Waarde (gold) \\
\midrule
Is rechtshandeling & \texttt{true} & \texttt{true} \\
Datum rechtshandeling & \texttt{2022-02-04} & \texttt{2022-02-04} \\
Verricht door de schuldenaar & \texttt{true} & \texttt{true} \\
Om niet & \texttt{false} & \texttt{false} \\
Datum faillietverklaring & \texttt{2022-05-02} & \texttt{2022-05-02} \\
Onverplicht & \textcolor{gray}{\texttt{true}} & \texttt{false} \\
Benadeling schuldeisers & \textcolor{gray}{\texttt{true}} & \texttt{false} \\
Wetenschap schuldenaar aangetoond & \texttt{false} & \texttt{false} \\
Wetenschap wederpartij aangetoond & \texttt{false} & \texttt{null} \\
Vermoeden wetenschap schuldenaar tegenbewijs & \texttt{null} & \texttt{null} \\
Vermoeden wetenschap wederpartij tegenbewijs & \texttt{null} & \texttt{null} \\
Tegen onderwaarde & \texttt{null} & \texttt{null} \\
Voldoening of zekerheid niet opeisbare schuld & \texttt{null} & \texttt{null} \\
Gelieerde partij & \textcolor{gray}{\texttt{true}} & \texttt{null} \\
Reeds verplicht voor aanvang termijn & \texttt{null} & \texttt{null} \\
\textbf{Vernietigbaar} & \textbf{\textcolor{gray}{\texttt{true}}} & \textbf{\texttt{false}} \\
\bottomrule
\end{tabular}
\end{center}

\paragraph{Is rechtshandeling - overwegingen}\label{is-rechtshandeling---overwegingen-15}

De overboekingen van Durisol naar Mevriet Vastgoed zijn op een bankrekening gedaan, wat een rechtshandeling is (goederenrechtelijke handeling).

\paragraph{Verricht door de schuldenaar - overwegingen}\label{verricht-door-de-schuldenaar---overwegingen-15}

De overboekingen zijn verricht door Durisol, de latere gefailleerde, zelf.

\paragraph{Om niet - overwegingen}\label{om-niet---overwegingen-15}

De overboekingen zijn gedaan ter afwikkeling van een onverschuldigde betaling door Rabobank aan Durisol, waarbij Durisol het surplus (dat eigenlijk aan Mevriet Vastgoed toekwam) doorbetaalde. Dit is geen om niet-handeling, omdat er een verband is met de afwikkeling van een financiële verplichting (onverschuldigde betaling).

\paragraph{Onverplicht - overwegingen}\label{onverplicht---overwegingen-15}

De betalingen door Durisol aan Mevriet Vastgoed (€ 275.000 op 4 februari 2022 en € 18.143,81 op 9 februari 2022) zijn onverplicht verricht in de zin van artikel 42 lid 1 Fw.

Allereerst is vastgesteld dat Rabobank per abuis het gehele surplus van € 349.981,74 (na royement van de hypotheek) aan Durisol heeft overgemaakt, terwijl dit bedrag (min de schuld van Durisol aan Rabobank) toekwam aan Mevriet Vastgoed als eigenaar van het onderpand (art. 2.4). Durisol had dus geen recht op dit surplus en was derhalve niet verplicht om dit bedrag aan Mevriet Vastgoed door te betalen.

Daarnaast was er geen sprake van een bestaande verplichting van Durisol om Mevriet Vastgoed te betalen. De betalingen zijn gedaan om een onverschuldigd ontvangen bedrag door te storten naar de rechtmatige rechthebbende (Mevriet Vastgoed), maar dit was geen verplichting die op Durisol rustte. Er bestond geen contractuele of wettelijke grondslag voor deze betalingen.

Morele, feitelijke of economische druk (zoals de wens om de situatie recht te zetten) schept geen rechtsplicht (vgl. ECLI:NL:GHSHE:2017:4528, r.o. 3.5.4). De stelplicht en bewijslast van het onverplichte karakter rusten op de curator, en in deze casus is voldoende aannemelijk gemaakt dat Durisol niet verplicht was tot deze betalingen.

\paragraph{Benadeling schuldeisers - overwegingen}\label{benadeling-schuldeisers---overwegingen-15}

In de onderhavige casus heeft Durisol op 4 en 9 februari 2022 een totaalbedrag van € 293.143,81 (€ 275.000 + € 18.143,81) overgemaakt aan Mevriet Vastgoed, met de vermelding `Surplus afgelost bankkrediet Rabobank'. Dit bedrag was eerder door Rabobank ten onrechte aan Durisol overgemaakt, terwijl het eigenlijk toekwam aan Mevriet Vastgoed als rechthebbende op de koopsom na royement van de hypotheek.

De benadeling van de schuldeisers moet worden beoordeeld door een vermogensvergelijking: de hypothetische situatie waarin de schuldeisers zouden hebben verkeerd zonder de gewraakte rechtshandeling (de overboekingen door Durisol aan Mevriet Vastgoed) wordt vergeleken met de situatie waarin zij feitelijk verkeren als die handeling in stand blijft.

In de hypothetische situatie zonder de overboekingen zou Durisol het bedrag van € 293.143,81 hebben behouden als onderdeel van haar vermogen. Dit bedrag had dan beschikbaar kunnen zijn voor de gezamenlijke schuldeisers van Durisol, waaronder Rabobank, om hun vorderingen te verhalen. In de feitelijke situatie is dit bedrag echter overgemaakt aan Mevriet Vastgoed, waardoor het niet langer deel uitmaakt van het vermogen van Durisol. Hierdoor is het actief van Durisol verminderd, wat direct ten nadele komt van de schuldeisers, die nu minder vermogen hebben waarop zij zich kunnen verhalen.

Er is dus sprake van een vermindering van het actief van Durisol, wat een vorm van benadeling van de schuldeisers is in de zin van artikel 42 Fw. De benadeling is aanwezig op het moment waarop de curator zijn rechten doet gelden, nu het bedrag niet meer in de boedel aanwezig is en de schuldeisers hierdoor daadwerkelijk nadeel ondervinden in hun verhaalsmogelijkheden.

Daarnaast is er geen sprake van een gunstig gevolg voor de boedel dat de benadeling zou compenseren. Het bedrag is immers definitief uit de boedel verdwenen en niet teruggekeerd naar Durisol of haar schuldeisers.

Gelet op het voorgaande zijn de schuldeisers door de overboekingen benadeeld in de zin van artikel 42 Fw.

\paragraph{Wetenschap - overwegingen}\label{wetenschap---overwegingen-15}

De betalingen van Durisol aan Mevriet Vastgoed (€ 275.000 op 4 februari 2022 en € 18.143,81 op 9 februari 2022) zijn verricht binnen een jaar voor de faillietverklaring van Durisol (2 mei 2022). Mevriet Vastgoed is een groepsvennootschap van Durisol, wat valt onder art. 43 lid 1 sub 6° Fw. Hierdoor wordt de wetenschap van benadeling aan beide zijden (Durisol en Mevriet Vastgoed) vermoed. Er is geen sprake van een verplichting die vóór de jaartermijn was aangegaan, dus het vermoeden van art. 43 Fw is van toepassing. Er is geen tegenbewijs geleverd dat het vermoeden ontzenuwt.

\newpage

\markboth{JOR 2026/137}{gpt-5.5 (k = 10)}

\subsubsection{GPT}\label{gpt-5}

\texttt{gpt-5.5}

\begin{center}
\small
\begin{tabular}{@{}p{7.4cm}p{2.6cm}p{2.6cm}@{}}
\toprule
Variabele & Waarde (model) & Waarde (gold) \\
\midrule
Is rechtshandeling & \texttt{true} & \texttt{true} \\
Datum rechtshandeling & \texttt{2022-02-04} & \texttt{2022-02-04} \\
Verricht door de schuldenaar & \texttt{true} & \texttt{true} \\
Om niet & \texttt{false} & \texttt{false} \\
Datum faillietverklaring & \texttt{2022-05-02} & \texttt{2022-05-02} \\
Onverplicht & \texttt{false} & \texttt{false} \\
Benadeling schuldeisers & \texttt{false} & \texttt{false} \\
Wetenschap schuldenaar aangetoond & \texttt{false} & \texttt{false} \\
Wetenschap wederpartij aangetoond & \texttt{false} & \texttt{null} \\
Vermoeden wetenschap schuldenaar tegenbewijs & \texttt{false} & \texttt{null} \\
Vermoeden wetenschap wederpartij tegenbewijs & \texttt{false} & \texttt{null} \\
Tegen onderwaarde & \texttt{null} & \texttt{null} \\
Voldoening of zekerheid niet opeisbare schuld & \texttt{null} & \texttt{null} \\
Gelieerde partij & \textcolor{gray}{\texttt{true}} & \texttt{null} \\
Reeds verplicht voor aanvang termijn & \texttt{false} & \texttt{null} \\
\textbf{Vernietigbaar} & \textbf{\texttt{false}} & \textbf{\texttt{false}} \\
\bottomrule
\end{tabular}
\end{center}

\paragraph{Is rechtshandeling - overwegingen}\label{is-rechtshandeling---overwegingen-16}

De gewraakte handelingen zijn de overboekingen door Durisol aan Mevriet Vastgoed van € 275.000 op 4 februari 2022 en € 18.143,81 op 9 februari 2022. Deze betalingen/overboekingen waren gericht op rechtsgevolg, namelijk het doorbetalen van het door Rabobank aan Durisol betaalde surplus en volgens de feiten het aflossen van de vordering uit onverschuldigde betaling.

\paragraph{Verricht door de schuldenaar - overwegingen}\label{verricht-door-de-schuldenaar---overwegingen-16}

De betalingen zijn verricht door Durisol. Durisol is de vennootschap waarvan op 2 mei 2022 het faillissement is uitgesproken.

\paragraph{Om niet - overwegingen}\label{om-niet---overwegingen-16}

De betalingen zijn niet zonder meer zonder wederprestatie verricht. Volgens de feiten betrof het de doorbetaling van een aan Durisol onverschuldigd betaald en voor Mevriet Vastgoed bestemd bedrag, waarbij Rabobank ermee instemde dat door deze betalingen haar vordering uit onverschuldigde betaling op Durisol was afgelost. Daarmee stond tegenover de betalingen een daarmee verband houdend vermogensrechtelijk gevolg voor Durisol.

\paragraph{Onverplicht - overwegingen}\label{onverplicht---overwegingen-16}

De overboekingen door Durisol aan Mevriet Vastgoed betroffen het deel van de door Rabobank abusievelijk aan Durisol betaalde opbrengst dat niet aan Durisol, maar aan Mevriet Vastgoed als eigenaar/hypotheekgever toekwam. Rabobank heeft dit ook zo bevestigd: na voldoening van de schuld van Durisol aan Rabobank was Mevriet Vastgoed rechthebbende op het restant van de koopsom. Durisol had daarom geen vrije beschikkingsruimte over dit bedrag. Zij moest het onverschuldigd aan haar betaalde, voor Mevriet Vastgoed bestemde bedrag terugleiden naar de rechthebbende dan wel daarmee de restitutieverplichting jegens Rabobank uitvoeren. Dat de betaling rechtstreeks aan Mevriet Vastgoed is gedaan, maakt die uitvoering niet onverplicht; een bestaande verplichting kan immers ook op een niet uitdrukkelijk voorgeschreven wijze worden uitgevoerd zonder dat de uitvoering daardoor onverplicht wordt. Dat geldt temeer nu Rabobank geen bezwaar had tegen de rechtstreekse overboekingen en achteraf heeft bevestigd dat het creditsaldo aan Mevriet Vastgoed toebehoorde en dat de overboekingen correct waren verwerkt. De situatie is vergelijkbaar met de doorbetalingsjurisprudentie waarin betalingen die de schuldenaar voor een ander ontvangt en vervolgens doorbetaalt, niet als onverplicht worden aangemerkt indien een rechtsgrond voor doorbetaling bestaat, zoals in ECLI:NL:GHLEE:2007:BB3627. De curator draagt de stelplicht en bewijslast van het onverplichte karakter, maar uit de casus volgt juist dat voor Durisol een rechtens relevante verplichting bestond om het surplus niet te behouden maar door te betalen.

\paragraph{Benadeling schuldeisers - overwegingen}\label{benadeling-schuldeisers---overwegingen-16}

De benadeling moet worden beoordeeld door vergelijking van de feitelijke situatie met de reële hypothetische situatie zonder de overboekingen van Durisol aan Mevriet Vastgoed. Feitelijk heeft Durisol € 293.143,81 doorbetaald aan Mevriet Vastgoed; daarmee is de door de foutieve betaling van Rabobank ontstane terugbetalings-/correctieverhouding afgewikkeld. Zonder deze doorbetaling zou Durisol weliswaar formeel een creditsaldo hebben behouden, maar dat saldo was het gevolg van een evidente vergissing van Rabobank bij de afwikkeling van de verkoopopbrengst van onroerend goed van Mevriet Vastgoed. Rabobank heeft uitdrukkelijk verklaard dat dit bedrag aan Mevriet Vastgoed toekwam als eigenaar/hypotheekgever en dat het creditsaldo op de rekening van Durisol aan Mevriet Vastgoed toebehoorde. Ook heeft Rabobank de rechtstreekse overboekingen aanvaard als correcte verwerking/aflossing van haar vordering uit onverschuldigde betaling op Durisol. In de hypothetische situatie is daarom niet reëel dat dit bedrag als vrij verhaalsactief voor de gezamenlijke schuldeisers van Durisol beschikbaar zou zijn gebleven; Rabobank zou de fout hebben gecorrigeerd dan wel haar terugbetalingsvordering hebben verrekend/uitgewonnen, waarna het surplus alsnog bij Mevriet Vastgoed zou zijn terechtgekomen. Anders dan in ECLI:NL:GHSHE:2014:4568 gaat het hier dus niet enkel om betaling van een concurrente schuldeiser uit vrij actief, maar om correctie van een bankfout bij de verdeling van een opbrengst waarop de schuldeisers van Durisol zich in werkelijkheid niet mochten verhalen. De doorbetaling heeft daarom niet geleid tot minder verhaal of tot een relevante rangordeverstoring; hoogstens is de schuldeisers een toevallig voordeel ontgaan.

\paragraph{Wetenschap - overwegingen}\label{wetenschap---overwegingen-16}

Rechtstreekse wetenschap van benadeling is op de gegeven feiten niet zelfstandig aangetoond: uit de betalingen kort vóór faillissement en de latere eigen aangifte volgt op zichzelf nog niet dat op 4 en 9 februari 2022 het faillissement en een tekort daarin met redelijke mate van waarschijnlijkheid voorzienbaar waren voor Durisol en Mevriet Vastgoed. Wel geldt een bewijsvermoeden van art. 43 Fw: de betalingen zijn binnen één jaar vóór faillissement verricht, een verplichting van vóór aanvang van die jaartermijn blijkt niet, en Durisol en Mevriet Vastgoed zijn gelieerde/groepsvennootschappen. Daardoor wordt de wetenschap aan beide zijden vermoed. De omstandigheid dat Rabobank het surplus per vergissing aan Durisol had betaald en achteraf instemde met doorbetaling aan Mevriet Vastgoed kan relevant zijn voor rechtsgrond/benadeling, maar ontzenuwt op zichzelf niet het art. 43-vermoeden dat faillissement en tekort niet voorzienbaar waren. Van succesvol tegenbewijs blijkt daarom niet. Art. 45 Fw is niet doorslaggevend, omdat de doorbetaling kennelijk strekte tot aflossing/verrekening van de onverschuldigde betaling en dus niet zonder meer als rechtshandeling om niet moet worden aangemerkt.

\newpage

\markboth{JOR 2026/137}{gemini-3.5-flash (k = 10)}

\subsubsection{Gemini}\label{gemini-5}

\texttt{gemini-3.5-flash}

\begin{center}
\small
\begin{tabular}{@{}p{7.4cm}p{2.6cm}p{2.6cm}@{}}
\toprule
Variabele & Waarde (model) & Waarde (gold) \\
\midrule
Is rechtshandeling & \texttt{true} & \texttt{true} \\
Datum rechtshandeling & \texttt{2022-02-04} & \texttt{2022-02-04} \\
Verricht door de schuldenaar & \texttt{true} & \texttt{true} \\
Om niet & \texttt{false} & \texttt{false} \\
Datum faillietverklaring & \texttt{2022-05-02} & \texttt{2022-05-02} \\
Onverplicht & \textcolor{gray}{\texttt{true}} & \texttt{false} \\
Benadeling schuldeisers & \textcolor{gray}{\texttt{true}} & \texttt{false} \\
Wetenschap schuldenaar aangetoond & \textcolor{gray}{\texttt{true}} & \texttt{false} \\
Wetenschap wederpartij aangetoond & \textcolor{gray}{\texttt{true}} & \texttt{null} \\
Vermoeden wetenschap schuldenaar tegenbewijs & \texttt{false} & \texttt{null} \\
Vermoeden wetenschap wederpartij tegenbewijs & \texttt{false} & \texttt{null} \\
Tegen onderwaarde & \texttt{null} & \texttt{null} \\
Voldoening of zekerheid niet opeisbare schuld & \texttt{null} & \texttt{null} \\
Gelieerde partij & \textcolor{gray}{\texttt{true}} & \texttt{null} \\
Reeds verplicht voor aanvang termijn & \texttt{false} & \texttt{null} \\
\textbf{Vernietigbaar} & \textbf{\textcolor{gray}{\texttt{true}}} & \textbf{\texttt{false}} \\
\bottomrule
\end{tabular}
\end{center}

\paragraph{Is rechtshandeling - overwegingen}\label{is-rechtshandeling---overwegingen-17}

De overboekingen van in totaal € 293.143,81 door Durisol aan Mevriet Vastgoed zijn rechtshandelingen, aangezien deze betalingen gericht zijn op het teweegbrengen van een rechtsgevolg, namelijk het overdragen van vermogensbestanddelen ter voldoening van een vermeende schuld.

\paragraph{Verricht door de schuldenaar - overwegingen}\label{verricht-door-de-schuldenaar---overwegingen-17}

De overboekingen zijn verricht door Durisol, de schuldenaar die naderhand failliet is verklaard.

\paragraph{Om niet - overwegingen}\label{om-niet---overwegingen-17}

De rechtshandeling is anders dan om niet (om baat). De betalingen vonden plaats ter aflossing van een vordering uit onverschuldigde betaling c.q. ter doorstorting van gelden die Rabobank per abuis aan Durisol had overgemaakt maar die toebehoorden aan Mevriet Vastgoed. Er stond derhalve een prestatie (bevrijding van een schuld/verplichting) tegenover.

\paragraph{Onverplicht - overwegingen}\label{onverplicht---overwegingen-17}

De betalingen door Durisol aan Mevriet Vastgoed B.V. zijn onverplicht verricht. Een rechtshandeling is onverplicht wanneer daartoe geen op de wet of een overeenkomst berustende rechtsplicht bestond die in rechte kon worden afgedwongen. Hoewel Rabobank per abuis het surplus van de verkoopopbrengst (dat toebehoorde aan Mevriet Vastgoed) aan Durisol had overgemaakt, creëerde dit enkel een verbintenis uit onverschuldigde betaling tussen Rabobank en Durisol. Er bestond geen directe contractuele of wettelijke rechtsbetrekking tussen Durisol en Mevriet Vastgoed op grond waarvan Mevriet Vastgoed betaling van Durisol kon afdwingen. De omstandigheid dat Rabobank achteraf heeft ingestemd met de rechtstreekse overboekingen ter aflossing van haar eigen schuld aan Mevriet Vastgoed en haar vordering op Durisol, maakt de betalingen jegens Mevriet Vastgoed niet verplicht. Nu Mevriet Vastgoed de betalingen niet in rechte van Durisol kon afdwingen, zijn de overboekingen onverplicht verricht.

\paragraph{Benadeling schuldeisers - overwegingen}\label{benadeling-schuldeisers---overwegingen-17}

Om te bepalen of er sprake is van benadeling van schuldeisers in de zin van artikel 42 Fw, moet een vergelijking worden gemaakt tussen de feitelijke situatie waarin de schuldeisers verkeren en de hypothetische situatie zonder de betalingen van in totaal € 293.143,81 aan Mevriet Vastgoed. Indien de betalingen niet waren verricht, zou dit bedrag (dat door Rabobank per abuis aan Durisol was overgemaakt) in het vermogen van Durisol zijn achtergebleven. Mevriet Vastgoed zou wegens deze foutieve betaling een vordering uit onverschuldigde betaling of ongerechtvaardigde verrijking hebben op Durisol. Omdat deze onverschuldigde betaling voorafgaand aan het faillissement heeft plaatsgevonden, zou de vordering van Mevriet Vastgoed in het faillissement van Durisol slechts een concurrente vordering zijn (conform ECLI:NL:GHSHE:2014:4568). Door de betalingen is de concurrente vordering van Mevriet Vastgoed volledig voldaan, waardoor het actieve boedelvermogen met € 293.143,81 is verminderd. Hierdoor is de verhaalsruimte voor de overige schuldeisers verkleind, wat leidt tot een feitelijke benadeling van de schuldeisers.

\paragraph{Wetenschap - overwegingen}\label{wetenschap---overwegingen-17}

De betalingen van Durisol aan haar groepsvennootschap Mevriet Vastgoed ter correctie van de foutieve overboeking door Rabobank zijn onverplicht. Conform de lijn uit ECLI:NL:GHSHE:2014:4568 (Studio '83) leidt het terugstorten van een vóór faillissement abusievelijk ontvangen bedrag tot benadeling van de schuldeisers, omdat de werkelijke rechthebbende (Mevriet Vastgoed) na faillissement slechts een concurrente vordering op de boedel zou hebben. Omdat de rechtshandeling binnen een jaar voor de faillietverklaring is verricht met een groepsmaatschappij, geldt het wettelijk bewijsvermoeden van art. 43 lid 1 sub 6° Fw. Dit vermoeden is niet ontzenuwd. Daarnaast is de wetenschap van benadeling aan beide zijden ook rechtstreeks aangetoond: de vennootschappen zijn nauw verweven via gemeenschappelijke bestuurders en het faillissement werd kort daarna (binnen twee maanden) door Durisol zelf aangevraagd, waardoor het faillissement en een tekort daarin destijds met een redelijke mate van waarschijnlijkheid voorzienbaar waren.

\newpage

\markboth{ECLI:NL:RBAMS:2026:3787}{Gold standard}

\section{ECLI:NL:RBAMS:2026:3787}\label{eclinlrbams20263787}

uitspraakdatum 15-04-2026

\subsection{Gold}\label{gold-2}

\begin{center}
\small
\begin{tabular}{@{}p{8cm}p{2.6cm}@{}}
\toprule
Variabele & Waarde \\
\midrule
Is rechtshandeling & \texttt{true} \\
Datum rechtshandeling & \texttt{2021-01-21} \\
Verricht door de schuldenaar & \texttt{true} \\
Om niet & \texttt{true} \\
Datum faillietverklaring & \texttt{2021-07-13} \\
Onverplicht & \texttt{true} \\
Benadeling schuldeisers & \texttt{true} \\
Wetenschap schuldenaar aangetoond & \texttt{false} \\
Wetenschap wederpartij aangetoond & \texttt{null} \\
Vermoeden wetenschap schuldenaar tegenbewijs & \texttt{false} \\
Vermoeden wetenschap wederpartij tegenbewijs & \texttt{false} \\
Tegen onderwaarde & \texttt{true} \\
Voldoening of zekerheid niet opeisbare schuld & \texttt{null} \\
Gelieerde partij & \texttt{true} \\
Reeds verplicht voor aanvang termijn & \texttt{false} \\
\textbf{Vernietigbaar} & \textbf{\texttt{true}} \\
\bottomrule
\end{tabular}
\end{center}

Opmerking: deze casus wordt in het kader van de faillissementspauliana en art. 42 Fw beperkt tot de vraag of het samenstel van rechtshandelingen waarbij vermogensbestanddelen van {[}eiser 1{]} en {[}eiser 2{]} zijn overgedragen aan {[}gedaagde 3{]} paulianeus is.
Gebruik voor de datum van faillietverklaring de datum van {[}eiser 1{]}; voer 1 datum in.

\subsubsection{Uitspraak}\label{uitspraak-2}

\paragraph{RECHTBANK Amsterdam}\label{rechtbank-amsterdam}

Civiel recht

Zaaknummers (gevoegd): C/13/746164 / HA ZA 24-124 en

C/13/759431 / HA ZA 24-1246

\textbf{Vonnis van 15 april 2026}

\emph{in de zaak met nummer C/13/746164 / HA ZA 24-124 van}

\textbf{{[}curator{]}}

\textbf{in zijn hoedanigheid van curator in het faillissement van}

\textbf{{[}eiser 1{]} B.V.},

kantoorhoudende in {[}plaats{]} ,

eisende partij,

hierna te noemen: de curator,

advocaat: mr. F.A.E. Diderich,

tegen

\textbf{{[}gedaagde 2{]}}
,

wonende in {[}woonplaats 2{]} ,

gedaagde partij,

advocaat: mr. E. Hoekstra,

\emph{in de zaak met nummer C/13/759431 / HA ZA 24-1246 van}

\paragraph{1 {[}curator{]}}\label{curator}

\textbf{in hoedanigheid van curator in het faillissement van}

\textbf{{[}eiser 1{]} B.V.},

\begin{enumerate}
\def\labelenumi{\arabic{enumi}.}
\setcounter{enumi}{1}
\tightlist
\item
  \textbf{{[}curator{]}}
\end{enumerate}

\textbf{in hoedanigheid van curator in het faillissement van {[}eiser 2{]} B.V.},

kantoorhoudende in {[}plaats{]} ,

eisende partijen,

hierna samen te noemen: de curator,

advocaat: mr. F.A.E. Diderich,

tegen

\paragraph{1 {[}gedaagde 1{]} ,}\label{gedaagde-1}

wonende in {[}woonplaats 1{]} ,

hierna te noemen: {[}gedaagde 1{]} ,

\begin{enumerate}
\def\labelenumi{\arabic{enumi}.}
\setcounter{enumi}{1}
\tightlist
\item
  \textbf{{[}gedaagde 2{]}},
\end{enumerate}

wonende in {[}woonplaats 2{]} ,

hierna te noemen: {[}gedaagde 2{]} ,

\paragraph{3 {[}gedaagde 3{]} B.V.,}\label{gedaagde-3-b.v.}

gevestigd in {[}vestigingsplaats{]} ,

hierna te noemen: {[}gedaagde 3{]} ,

\begin{enumerate}
\def\labelenumi{\arabic{enumi}.}
\setcounter{enumi}{3}
\tightlist
\item
  \textbf{{[}gedaagde 4{]} B.V.},
\end{enumerate}

gevestigd in {[}vestigingsplaats{]} ,

hierna te noemen: {[}gedaagde 4{]} ,

\begin{enumerate}
\def\labelenumi{\arabic{enumi}.}
\setcounter{enumi}{4}
\tightlist
\item
  \textbf{{[}gedaagde 5{]} C.V.},
\end{enumerate}

gevestigd in {[}vestigingsplaats{]} ,

hierna te noemen: {[}gedaagde 5{]} CV,

\begin{enumerate}
\def\labelenumi{\arabic{enumi}.}
\setcounter{enumi}{5}
\tightlist
\item
  \textbf{{[}gedaagde 6{]}},
\end{enumerate}

wonende in {[}woonplaats 2{]} ,

hierna te noemen: {[}gedaagde 6{]} ,

gedaagde partijen,

advocaat gedaagde 1: mr. K.R. Stephan,

advocaat gedaagden 2 t/m 6: mr. E. Hoekstra.

{[}eiser 2{]} B.V. en {[}eiser 1{]} B.V. worden hierna ` {[}eiser 2{]} ' en ` {[}eiser 1{]} ' genoemd.

\textbf{Waar gaat de zaak over?}

Deze zaak gaat vooral over bestuurdersaansprakelijkheid. {[}gedaagde 1{]} en {[}gedaagde 2{]} zijn door de curator aangesproken voor het boedeltekort in de faillissementen van {[}eiser 2{]} en {[}eiser 1{]} van in totaal ruim 3 miljoen. Volgens de curator is sprake van kennelijk onbehoorlijk bestuur omdat de boekhoudplicht is geschonden, maar ook door aanzienlijke privé-onttrekkingen, onzakelijke transacties en paulianeus handelen. {[}gedaagde 1{]} en {[}gedaagde 2{]} hebben ruzie gekregen en wijzen voor de onregelmatigheden vooral naar elkaar.

De rechtbank oordeelt dat sprake is van kennelijk onbehoorlijk bestuur en dat {[}gedaagde 1{]} en {[}gedaagde 2{]} beide aansprakelijk zijn voor het boedeltekort.

\paragraph{1 De procedures in beide zaken}\label{de-procedures-in-beide-zaken}

1.1. Het verloop van de procedure in de zaak C/13/746164 / HA ZA 24-124 blijkt uit:

\begin{itemize}
\item
  de dagvaarding van 16 november 2023 met producties,
\item
  de conclusie van antwoord van {[}gedaagde 2{]} , met producties
\item
  het tussenvonnis van 22 oktober 2025 waarin een mondelinge behandeling is bepaald,
\item
  het proces-verbaal van de mondelinge behandeling van 4 maart 2026 en de daarin genoemde stukken.
\end{itemize}

1.2. Het verloop van de procedure in de zaak C/13/759431 / HA ZA 24-1246 blijkt uit:

\begin{itemize}
\item
  de dagvaarding van 30 oktober 2024 met een voegingsverzoek, en producties,
\item
  de conclusie van antwoord van {[}gedaagde 1{]} , met producties,
\item
  het vonnis van 8 januari 2025 waarin de zaak is gevoegd,
\item
  de conclusie van antwoord van {[}gedaagde 2{]} c.s.,
\item
  het proces-verbaal van de mondelinge behandeling van 4 maart 2026 en de daarin genoemde stukken.
\end{itemize}

1.3. Daarna is vonnis bepaald in beide zaken.

\paragraph{2 De feiten in beide zaken}\label{de-feiten-in-beide-zaken}

2.1. Betrokken partijen

\emph{{[}eiser 2{]} , {[}eiser 1{]} (en haar bestuurders {[}gedaagde 1{]} en {[}gedaagde 2{]} )}

2.1.1. {[}eiser 2{]} is in 1988 opgericht door {[}gedaagde 2{]} . Sinds 2014 was {[}gedaagde 1{]} samen met {[}gedaagde 2{]} bestuurder van {[}eiser 2{]} . Eind 2014 hebben {[}gedaagde 1{]} en {[}gedaagde 2{]} {[}eiser 1{]} opgericht. Zij waren vanaf de oprichting van die vennootschap beide bestuurder. {[}eiser 2{]} en {[}eiser 1{]} hielden zich bezig met de handel in gebruikte motorboten in het luxe segment.

2.1.2. De expertise van {[}eiser 2{]} en {[}eiser 1{]} lag vooral bij het opkopen van schepen in het middellandse zeegebied die beschadigd waren of waarvan het liggeld niet was betaald. Deze schepen werden door {[}eiser 2{]} en {[}eiser 1{]} gekocht, naar Nederland getransporteerd, opgeknapt en daarna weer te koop aangeboden.

2.1.3. Voor de exploitatie maakten {[}eiser 2{]} en {[}eiser 1{]} vooral gebruik van particuliere financiering, waarbij zij soms leningen aangingen met een rente van 15\% of 20\% per jaar. Tussen februari 2015 en december 2019 hebben {[}eiser 2{]} en {[}eiser 1{]} in totaal € 3.155.000,00 van negen particuliere financiers ontvangen.

\emph{{[}gedaagde 5{]} CV, {[}gedaagde 6{]} , {[}gedaagde 4{]} en {[}gedaagde 3{]}}

2.1.4. In 1987 is {[}gedaagde 5{]} CV opgericht. {[}gedaagde 6{]} , de echtgenote van {[}gedaagde 2{]} , is ingeschreven als beherend vennoot van {[}gedaagde 5{]} CV en {[}gedaagde 2{]} zelf is commanditair vennoot. Eind 2015 hebben {[}gedaagde 2{]} en {[}gedaagde 1{]} met een derde {[}gedaagde 4{]} opgericht. {[}gedaagde 4{]} hield zich bezig met het foliewerk voor de schepen die door {[}eiser 1{]} en {[}eiser 2{]} werden opgeknapt en verkocht. {[}gedaagde 3{]} is door {[}gedaagde 2{]} opgericht op 22 december 2020 en houdt zich bezig met de verkoop en bemiddeling van schepen.

2.2. {[}naam boot 1{]} en Moody

\emph{Verkoop helft {[}naam boot 1{]} door {[}gedaagde 1{]}}

2.2.1. Op 6 oktober 2017 heeft ` {[}gedaagde 5{]} ' een koopovereenkomst gesloten met {[}naam 1{]} (hierna: {[}naam 1{]} ) voor een schip, type Pershing 65 met de naam {[}naam boot 1{]} (hierna: {[}naam boot 1{]} ). {[}eiser 2{]} heeft op 19 oktober 2017 de in de overeenkomst genoemde koopsom van € 135.000,00 betaald.

2.2.2. Op 7 november 2017 heeft {[}gedaagde 1{]} aan {[}naam 2{]} (hierna: {[}naam 2{]} ) voorgesteld de onverdeelde helft van de {[}naam boot 1{]} aan hem te verkopen en de operationele kosten ieder voor de helft te dragen, waarbij partijen zijn uitgegaan van een koopsom van € 300.000,00. {[}naam 2{]} heeft {[}gedaagde 1{]} gevraagd om de koopovereenkomst met {[}naam 1{]} .

2.2.3. Op 11 december 2017 heeft {[}gedaagde 1{]} via zijn eigen e-mailadres een koopovereenkomst van 6 oktober 2017 tussen ` {[}gedaagde 5{]} ' en {[}naam 1{]} voor de {[}naam boot 1{]} gestuurd naar het e-mailadres van {[}eiser 2{]} en {[}eiser 1{]} . Daarin is in plaats van € 135.000,00 een koopsom van € 275.000,00 opgenomen en de ondertekening van {[}naam 1{]} wijkt af van die in de in 2.2.1 genoemde koopovereenkomst. {[}gedaagde 1{]} heeft op 11 december 2017 de afwijkende versie vanaf het e-mailadres van {[}eiser 2{]} en {[}eiser 1{]} naar {[}naam 2{]} doorgestuurd en die dag de onverdeelde helft van de {[}naam boot 1{]} voor € 150.000,00 aan {[}naam 2{]} verkocht. De {[}naam boot 1{]} is op 7 december 2017 bij het RDW op naam van {[}gedaagde 1{]} geregistreerd.

\emph{Rentebetalingen {[}eiser 1{]} voor Moody van {[}gedaagde 2{]}}

2.2.4. Op 23 januari 2018 heeft {[}gedaagde 2{]} een Moody 54 met de naam {[}naam boot 2{]} (hierna: Moody) bij het Kadaster als zijn privé-eigendom teboekgesteld. Op 13 april 2018 heeft {[}gedaagde 2{]} ten gunste van {[}naam 3{]} (hierna: {[}naam 3{]} ) een hypotheekrecht op de Moody gevestigd. Tussen 18 juni 2018 en 14 april 2021 heeft {[}eiser 1{]} door betalingen van steeds € 1.500,00 in totaal € 48.000,00 aan {[}naam 3{]} betaald onder vermelding van `rente'.

2.3. Vakantiewoning {[}adres 1{]}

2.3.1. Op 24 juli 2018 heeft {[}eiser 2{]} een vakantiewoning aan de {[}adres 1{]} (hierna: de vakantiewoning) gekocht en betaald met twee schepen. De vakantiewoning is in juli 2017 getaxeerd op een marktwaarde van € 460.000,00.

2.3.2. Op 3 september 2018 hebben {[}gedaagde 1{]} en {[}gedaagde 2{]} {[}bedrijf 2{]} B.V. (hierna: {[}bedrijf 2{]} ) opgericht. {[}eiser 2{]} heeft op 28 september 2018 de vakantiewoning voor € 300.000,00 aan {[}bedrijf 2{]} verkocht, waarna {[}bedrijf 2{]} de vakantiewoning op 14 december 2018 voor € 397.500,00 heeft verkocht aan een derde.

2.4. Administratievoering {[}eiser 2{]} en {[}eiser 1{]}

2.4.1. In de loop van 2018 is {[}naam 4{]} (hierna: {[}naam 4{]} ) als extern boekhouder bij {[}eiser 2{]} en {[}eiser 1{]} betrokken. Hij heeft op 17 oktober 2022 bij de rechter-commissaris verklaard:

``\emph{Er waren vaak boekingen van de ene onderneming naar de andere. Als de ene onderneming te weinig geld had voor de financiering van een boot, dan kwam er geld vanuit de andere onderneming. Er was een soort rekening courant tussen de twee ondernemingen. In de loop van 2020 kwam hier een derde onderneming bij, CV. Voor 2020 gebeurde er bijna niets in de CV.}''

2.4.2. In de periode vanaf 2018 tot en met april 2021 zijn van en naar {[}eiser 2{]} en {[}eiser 1{]} , {[}gedaagde 1{]} , {[}gedaagde 2{]} , {[}gedaagde 4{]} , {[}gedaagde 5{]} CV en {[}gedaagde 3{]} vele betalingen gedaan met de omschrijving `lening'. De curator heeft voor deze betalingen in de administratie van {[}eiser 2{]} en {[}eiser 1{]} geen leningsovereenkomsten gevonden.

\emph{Aanpassingen administratie}

2.4.3. Eind 2019 heeft {[}naam 4{]} {[}gedaagde 1{]} en {[}gedaagde 2{]} op de hoogte gebracht van het verlies van {[}eiser 2{]} van € 298.000,00 over 2018. In reactie daarop heeft {[}gedaagde 1{]} {[}naam 4{]} verzocht de administratie zo aan te passen dat {[}eiser 2{]} over 2018 een winst van € 30.000,00 heeft behaald. Uiteindelijk is het resultaat over 2018, na discussie tussen {[}gedaagde 1{]} en {[}naam 4{]} over mogelijke onzakelijke transacties, aangepast naar een winst van € 21.397,00.

2.4.4. Op 24 juni 2020 heeft {[}naam 4{]} de voorlopige cijfers van 2019 van {[}eiser 2{]} en {[}eiser 1{]} aan {[}gedaagde 1{]} en {[}gedaagde 2{]} gepresenteerd. Voor {[}eiser 2{]} was dit een verlies van € 839.510,00 en voor {[}eiser 1{]} een winst van € 127.638,00. Na overleg met {[}gedaagde 1{]} en {[}gedaagde 2{]} is dit in de voorlopige cijfers van 2019 aangepast naar een winst voor {[}eiser 2{]} van € 50.487,00 en een winst voor {[}eiser 1{]} van € 139.887,00.

2.4.5. Op 25 september 2020 hebben {[}gedaagde 1{]} en {[}gedaagde 2{]} overleg gehad met {[}naam 4{]} over het verlies over 2019 van {[}eiser 2{]} (zie ook hierna in 2.4.6). Na het overleg is in de voorlopige jaarrekening opgenomen dat {[}eiser 2{]} een winst van € 13.142,00 heeft behaald en {[}eiser 1{]} een verlies van € 522.000,00.

\emph{Meldingen boekhouder over slechte financiële situatie}

2.4.6. Op 13 januari 2020 heeft {[}naam 4{]} {[}gedaagde 1{]} en {[}gedaagde 2{]} laten weten dat de continuïteit van {[}eiser 2{]} en {[}eiser 1{]} mogelijk gevaar loopt. {[}naam 4{]} heeft op 30 april 2020 de Belastingdienst verzocht het negatieve jaarsaldo van {[}eiser 1{]} voor 2019 vast te stellen op € 895.153,00. Volgens opgave van 8 juli 2020 van {[}naam 4{]} was de openstaande belastingschuld van {[}eiser 1{]} € 22.922,00 en van {[}eiser 2{]} € 111.46,001. Voorafgaand aan de geplande bespreking van 25 september 2020 met {[}gedaagde 1{]} en {[}gedaagde 2{]} heeft {[}naam 4{]} hen laten weten dat de continuïteit van {[}eiser 2{]} niet meer gegarandeerd is en dat een faillissement van {[}eiser 2{]} niet ondenkbaar is.

2.5. Ruzie {[}gedaagde 1{]} en {[}gedaagde 2{]}

2.5.1. Op 15 oktober 2020 hebben {[}gedaagde 1{]} en {[}gedaagde 2{]} , na de relatiebreuk tussen {[}gedaagde 1{]} en de dochter van {[}gedaagde 2{]} , een gesprek gehad op het kantoor van {[}eiser 1{]} en {[}eiser 2{]} . Dit gesprek is geëindigd in ruzie, waarna {[}gedaagde 2{]} voor {[}gedaagde 1{]} de toegang heeft geblokkeerd tot e-mailaccounts, bankrekeningen, de software en het bedrijfspand. Na de breuk heeft {[}gedaagde 2{]} aanzienlijke bedragen overgemaakt van {[}eiser 2{]} en {[}eiser 1{]} naar zichzelf en naar {[}gedaagde 1{]} . {[}gedaagde 2{]} heeft later ook weer bedragen teruggestort. Er zijn na 15 oktober 2020 meerdere procedures gestart tussen {[}gedaagde 1{]} aan de ene kant en {[}gedaagde 2{]} en {[}eiser 2{]} aan de andere kant waarbij over en weer beslagen zijn gelegd. Op 7 december 2020 is {[}gedaagde 1{]} ontslagen als statutair bestuurder van {[}eiser 2{]} en op 1 maart 2021 als statutair bestuurder van {[}eiser 1{]} .

2.6. Diverse gebeurtenissen na de ruzie tussen {[}gedaagde 1{]} en {[}gedaagde 2{]}

\emph{{[}gedaagde 5{]} CV}

2.6.1. Vanaf januari 2021 heeft {[}gedaagde 5{]} CV bedragen ontvangen voor de verkoop van schepen van {[}eiser 1{]} :

\begin{itemize}
\item
  op 12 januari 2021 heeft {[}gedaagde 5{]} CV bij de verkoop van de Targa 2000 van {[}eiser 1{]} de koopsom van € 164.000,00 ontvangen,
\item
  op 4 maart 2021 heeft {[}gedaagde 5{]} CV € 12.665,40 ontvangen van Bates Wharf voor annulering van de aankoop door {[}eiser 1{]} van de Azimut 46 in oktober 2020. Daarna is op 16 maart 2021 de rekeningcourantschuld van {[}eiser 1{]} aan {[}gedaagde 5{]} CV met € 12.215,00 naar beneden bijgesteld,
\item
  op 16 maart 2021 heeft {[}gedaagde 5{]} CV bij de verkoop van de Astondoa 45 € 130.000,00 van een derde ontvangen,
\item
  op 9 april 2021 heeft {[}gedaagde 5{]} CV de Princess 65 van {[}eiser 1{]} verkocht aan een derde en daarvoor de aanbetaling van € 15.000,00 ontvangen,
\item
  op 30 april heeft {[}gedaagde 5{]} CV de Sea Ray van {[}eiser 1{]} verkocht aan een derde en daarvoor € 18.000,00 en een Fairline Corsa ter waarde van € 69.500,00 ontvangen,
\item
  op 20 mei 2021 heeft {[}eiser 1{]} haar rekening-courant schuld aan {[}gedaagde 5{]} CV verrekend met de verkoop van twee Zodiacs van € 3.500,00 en € 2.250,00 door {[}gedaagde 5{]} CV aan een derde voor € 5.750,00.
\end{itemize}

2.6.2. In de administratie heeft de curator een factuur gevonden van 17 januari 2021 voor de aankoop door {[}gedaagde 5{]} CV van de Boston Whaler van {[}eiser 1{]} voor € 60.000,00. In de administratie van {[}eiser 2{]} is op 1 maart 2021 verwerkt dat per 25 januari 2021 door de verkoop van de Boston Whaler de rekeningcourantschuld van \emph{{[}eiser 2{]}} aan {[}gedaagde 5{]} CV met € 39.500,00 is afgenomen. {[}gedaagde 5{]} CV heeft op 24 maart 2021 de RDW gevraagd de Boston Whaler op haar naam te registreren.

2.6.3. Op 12 maart 2021 heeft {[}gedaagde 2{]} gevraagd de verzekering voor de handelsvoorraad van {[}eiser 1{]} over te schrijven naar {[}gedaagde 5{]} CV. De verzekeraar heeft daarop een creditnota van € 8.470,00 aan {[}eiser 1{]} gestuurd. Daarna is € 8.470,00 verrekend met de verzekeringspremie die {[}gedaagde 5{]} CV voor de voorraad heeft betaald.

\emph{{[}gedaagde 3{]}}

2.6.4. Op 21 januari 2021 heeft {[}eiser 2{]} via SIDN al haar domeinnamen, waaronder {[}domeinnaam 1{]} en {[}domeinnaam 2{]} , aan {[}gedaagde 3{]} overgedragen. Daarna hebben {[}eiser 1{]} en {[}eiser 2{]} op 10 februari 2021 de rechten uit de huurovereenkomsten per 1 januari 2021 overgedragen aan {[}gedaagde 3{]} .

2.6.5. Op 22 juni 2021 heeft {[}naam 4{]} aan de salarisadministrateur van {[}eiser 2{]} en {[}eiser 1{]} laten weten dat het personeel van {[}eiser 2{]} per 1 juli 2021 overgaat naar {[}gedaagde 3{]} . In augustus 2021 heeft {[}gedaagde 3{]} € 13.249,95 aan salarissen van het personeel betaald en dit bedrag via de rekening-courant ten laste van {[}eiser 2{]} gebracht.

\emph{Beslaglegging particuliere financiers}

2.6.6. Op 8 februari 2021 hebben enkele particuliere financiers beslag gelegd onder {[}eiser 2{]} , {[}eiser 1{]} en {[}gedaagde 2{]} . {[}gedaagde 2{]} heeft op 10 februari 2021 een factuur aangemaakt met datum 2 januari 2021 van {[}gedaagde 4{]} aan {[}eiser 1{]} ten bedrage van € 20.500,00 aan werkzaamheden. Deze werd in rekeningcourant geboekt, waardoor er geen vordering van {[}eiser 1{]} op {[}gedaagde 4{]} resteerde.

2.6.7. Op 16 maart 2021 heeft {[}eiser 1{]} per 8 februari 2021 volgens opgave van {[}naam 4{]} van derdebeslagenen te vorderen:

\begin{itemize}
\item
  vordering {[}eiser 1{]} op {[}gedaagde 2{]} : € 26.681,00
\item
  vordering {[}eiser 1{]} op {[}gedaagde 5{]} CV: € 250.200,00
\item
  vordering {[}eiser 1{]} op {[}eiser 2{]} : € 648.934,00
\end{itemize}

2.6.8. Op 17 maart 2021 heeft {[}gedaagde 2{]} aan {[}naam 4{]} gevraagd:

``\emph{{[}naam 4{]} SVP ALLE SCHULD WEG BOEKEN}

\emph{SVP {[}naam 5{]} OOK BOEK OP PRIVE IS TOCH NIET MEER TE CONTROLEREN}''

Daarna zijn de rekening-courantvorderingen op derdebeslagenen per 8 februari 2021 als volgt aangepast:

\begin{itemize}
\item
  schuld {[}eiser 1{]} aan {[}gedaagde 2{]} : € 20.819,00
\item
  schuld {[}eiser 1{]} aan {[}gedaagde 5{]} CV: € 46.285,00
\item
  vordering {[}eiser 1{]} op {[}eiser 2{]} : € 816.434,38
\end{itemize}

2.7. Faillissementen van {[}eiser 2{]} en {[}eiser 1{]}

2.7.1. Op 13 juli 2021 heeft de rechtbank, op aanvraag van {[}gedaagde 1{]} en de particuliere financiers vermeld in 2.6.6, het faillissement van {[}eiser 1{]} uitgesproken. In het faillissement van\\
{[}eiser 1{]} heeft de curator aan preferente vorderingen € 786.738,00 voorlopig erkend en € 101.002,79 aan concurrente vorderingen. Per 3 maart 2026 is volgens opgave van de curator het boedeltekort € 886.548,01 + p.m. De rechtbank heeft op 28 september 2021 het faillissement van {[}eiser 2{]} uitgesproken. Per 7 oktober 2024 heeft de curator € 456.949,84 aan preferente vorderingen voorlopig erkend en € 1.648.915,76 aan concurrente vorderingen. Het boedeltekort per 3 maart 2026 is volgens opgave van de curator € 2.208.533,29 + p.m.

\paragraph{3 De geschillen in beide zaken}\label{de-geschillen-in-beide-zaken}

3.1. De curator van {[}eiser 1{]} vordert in de zaak met nummer C/13/746164 / HA ZA 24-124 dat de rechtbank, uitvoerbaar bij voorraad:

I. {[}gedaagde 2{]} veroordeelt om aan de curator € 48.000,00 te betalen, of een door de rechtbank te bepalen bedrag, met wettelijke rente,

\begin{enumerate}
\def\labelenumi{\Roman{enumi}.}
\setcounter{enumi}{1}
\item
  {[}gedaagde 2{]} veroordeelt om aan de curator € 35.000,00 (rekening-courantvordering) te betalen, of een door de rechtbank te bepalen bedrag, met wettelijke rente,
\item
  {[}gedaagde 2{]} veroordeelt om aan de curator € 3.318,86 te betalen aan beslagkosten,
\end{enumerate}

of een door de rechtbank te bepalen bedrag, met wettelijke rente,

\begin{enumerate}
\def\labelenumi{\Roman{enumi}.}
\setcounter{enumi}{3}
\tightlist
\item
  {[}gedaagde 2{]} veroordeelt in de kosten van de procedure, met wettelijke rente.
\end{enumerate}

3.2. In de zaak met nummer C/13/759431 / HA ZA 24-1246, vordert de curator dat de rechtbank, uitvoerbaar bij voorraad:

\emph{I. {[}gedaagde 4{]} (rekening-courantvorderingen)}

\begin{enumerate}
\def\labelenumi{\arabic{enumi}.}
\item
  {[}gedaagde 4{]} veroordeelt om aan de curator van {[}eiser 1{]} € 20.500,00 te betalen,
\item
  {[}gedaagde 4{]} veroordeelt om aan de curator van {[}eiser 2{]} € 12.486,48 te betalen,
\end{enumerate}

\emph{II. {[}gedaagde 3{]} (oneigenlijke doorstart en terugbetaling lening)}

\begin{enumerate}
\def\labelenumi{\arabic{enumi}.}
\setcounter{enumi}{2}
\item
  primair voor recht verklaart dat de curator van {[}eiser 1{]} en {[}eiser 2{]} rechtsgeldig de buitengerechtelijke vernietiging van het samenstel van handelingen dat samen de oneigenlijke doorstart vormt heeft ingeroepen (of gerechtelijke vernietiging), en subsidiair voor recht verklaart dat {[}gedaagde 3{]} door het samenstel van rechtshandelingen dat heeft geleid tot de oneigenlijke doorstart ten koste van {[}eiser 1{]} en {[}eiser 2{]} ongerechtvaardigd is verrijkt,
\item
  {[}gedaagde 3{]} veroordeelt om aan de curator van {[}eiser 1{]} € 32.284,34 te betalen, vermeerderd met de waarde van de softe activa van {[}eiser 1{]} , te bepalen op € 200.000,00, of een door de rechtbank te bepalen bedrag, of nader op te maken bij staat en te vereffenen volgens de wet onder betaling van een voorschot van € 100.000,00,
\item
  {[}gedaagde 3{]} veroordeelt om aan de curator van {[}eiser 2{]} € 317.285,70 te betalen, vermeerderd met de waarde van de softe activa van {[}eiser 2{]} , te bepalen op € 200.000,00, of een door de rechtbank te bepalen bedrag, of nader op te maken bij staat en te vereffenen volgens de wet onder betaling van een voorschot van € 100.000,00,
\end{enumerate}

\emph{III. {[}gedaagde 5{]} CV, {[}gedaagde 6{]} en {[}gedaagde 2{]} (handelsvoorraad {[}eiser 1{]} en terugbetaling lening)}

\begin{enumerate}
\def\labelenumi{\arabic{enumi}.}
\setcounter{enumi}{5}
\item
  {[}gedaagde 5{]} CV, {[}gedaagde 6{]} en {[}gedaagde 2{]} hoofdelijk of elk apart veroordeelt om aan de curator van {[}eiser 1{]} € 448.985,00 te betalen,
\item
  primair {[}gedaagde 5{]} CV, {[}gedaagde 6{]} en {[}gedaagde 2{]} hoofdelijk of elk apart veroordeelt tot afgifte van de Boston Whaler aan de curator van {[}eiser 1{]} of betaling van een waardevergoeding van € 60.000,00, en subsidiair tot afgifte van de Boston Whaler aan de curator van {[}eiser 2{]} of betaling van een waardevergoeding van € 60.000,00,
\end{enumerate}

\emph{IV. {[}gedaagde 2{]} (privé-onttrekkingen en paulianeus handelen)}

\begin{enumerate}
\def\labelenumi{\arabic{enumi}.}
\setcounter{enumi}{7}
\item
  {[}gedaagde 2{]} veroordeelt om aan de curator van {[}eiser 1{]} € 290.000,00 te betalen,
\item
  primair {[}gedaagde 2{]} veroordeelt om aan de curator van {[}eiser 2{]} € 231.690,00 te betalen en subsidiair de betaling van € 20.000,00 van {[}eiser 2{]} aan {[}gedaagde 2{]} te vernietigen, en {[}gedaagde 2{]} veroordeelt om € 166.990,00 aan de curator van {[}eiser 2{]} te betalen,
\end{enumerate}

\emph{V. {[}gedaagde 1{]} (privé-onttrekkingen en paulianeus handelen)}

\begin{enumerate}
\def\labelenumi{\arabic{enumi}.}
\setcounter{enumi}{9}
\item
  {[}gedaagde 1{]} veroordeelt om € 117.584,61 aan de curator van {[}eiser 1{]} te betalen,
\item
  {[}gedaagde 1{]} veroordeelt om aan € 415.185,00 aan de curator van {[}eiser 2{]} te betalen,
\end{enumerate}

\emph{VI. Bestuurdersaansprakelijkheid {[}gedaagde 2{]} en {[}gedaagde 1{]}}

Primair, aansprakelijk voor het boedeltekort (artikel 2:248 BW)

\begin{enumerate}
\def\labelenumi{\arabic{enumi}.}
\setcounter{enumi}{11}
\item
  voor recht verklaart dat {[}gedaagde 2{]} en {[}gedaagde 1{]} hoofdelijk of elk apart aansprakelijk zijn voor het tekort in het faillissement van {[}eiser 1{]} , nader op te maken bij staat en te vereffenen volgens de wet,
\item
  voor recht verklaart dat {[}gedaagde 2{]} en {[}gedaagde 1{]} hoofdelijk of elk apart aansprakelijk zijn voor het tekort in het faillissement van {[}eiser 2{]} , nader op te maken bij staat en te vereffenen volgens de wet,
\item
  {[}gedaagde 2{]} en {[}gedaagde 1{]} hoofdelijk of elk apart veroordeelt om aan de curator van {[}eiser 1{]} te betalen een voorschot op het boedeltekort van {[}eiser 1{]} van € 800.000,00, of een door de rechtbank te bepalen bedrag,
\item
  {[}gedaagde 2{]} en {[}gedaagde 1{]} hoofdelijk of elk apart veroordeelt om aan de curator van {[}eiser 2{]} te betalen een voorschot op het boedeltekort van {[}eiser 2{]} van € 2.000.000,00, of een door de rechtbank te bepalen bedrag,
\end{enumerate}

Subsidiair, aansprakelijk voor individuele schadeposten (artikel 2:9 BW)

\begin{enumerate}
\def\labelenumi{\arabic{enumi}.}
\setcounter{enumi}{15}
\item
  {[}gedaagde 2{]} en {[}gedaagde 1{]} hoofdelijk veroordeelt om aan de curator van {[}eiser 1{]} € 2.424.075,00 te betalen, of {[}gedaagde 2{]} veroordeelt om aan de curator van {[}eiser 1{]} € 2.424.075,00 te betalen en {[}gedaagde 1{]} € 1.468.997,00,
\item
  {[}gedaagde 2{]} en {[}gedaagde 1{]} hoofdelijk veroordeelt om aan de curator van {[}eiser 2{]} € 2.558.118,96 te betalen, of {[}gedaagde 2{]} veroordeelt om aan de curator van {[}eiser 2{]} € 2.558.118,96 te betalen en {[}gedaagde 1{]} € 1.278.239,54,
\end{enumerate}

\emph{VII. Vernietiging onverplichte vermindering verhaalsmogelijkheden (artikel 2:248 lid 9 BW)}

\begin{enumerate}
\def\labelenumi{\arabic{enumi}.}
\setcounter{enumi}{17}
\item
  de vestiging van het hypotheekrecht op 9 maart 2021 door {[}gedaagde 2{]} en {[}gedaagde 6{]} ten gunste van {[}naam 6{]} te vernietigen,
\item
  {[}gedaagde 2{]} en {[}gedaagde 6{]} veroordeelt om binnen 7 dagen na de betekening van dit vonnis de inschrijving van {[}naam 6{]} op de onroerende zaak, staande en gelegen aan het adres
\end{enumerate}

{[}adres 2{]} (kadastraal bekend {[}kadastrale aanduidingen{]} ) zorg te dragen voor opheffing dan wel doorhaling van de hypothecaire

inschrijving, op straffe van een dwangsom van € 100.000,00 ineens en € 5.000,00 voor iedere dag dat {[}gedaagde 2{]} en {[}gedaagde 6{]} in gebreke blijven aan deze veroordeling te voldoen, met een maximum van € 200.000,00,

\begin{enumerate}
\def\labelenumi{\arabic{enumi}.}
\setcounter{enumi}{19}
\tightlist
\item
  te bepalen dat, als {[}gedaagde 2{]} en {[}gedaagde 6{]} in gebreke blijven met voldoening, het
\end{enumerate}

vonnis in de plaats treedt van de medewerking van {[}gedaagde 2{]} en {[}gedaagde 6{]} , en dat dus de handtekeningen van {[}gedaagde 2{]} en {[}gedaagde 6{]} onder alle documenten die getekend moeten worden om te komen tot doorhaling van het recht van hypotheek dat ten name van {[}gedaagde 2{]} en {[}gedaagde 6{]} is gevestigd en ingeschreven op de data en tijdstippen als vermeld in

het kadaster volgens het als productie 120 in de dagvaarding bedoelde kadastraal

hypothecaire uittreksel, met machtiging aan de curator een afschrift van dit vonnis te doen

inschrijven in de openbare registers,

\begin{enumerate}
\def\labelenumi{\arabic{enumi}.}
\setcounter{enumi}{20}
\item
  de overdracht, althans overschrijving bij het RDW van de BMW met kenteken {[}kenteken{]} door {[}gedaagde 2{]} aan {[}gedaagde 3{]} te vernietigen,
\item
  de overdracht, althans overschrijving bij het RDW van de BMW met kenteken\\
  {[}kenteken{]} door {[}gedaagde 3{]} aan {[}naam 7{]} te vernietigen,
\end{enumerate}

vorderingen I tot en met VI vermeerderd met de buitengerechtelijke incassokosten en de

wettelijke rente, met hoofdelijke of afzonderlijke veroordeling in de proceskosten.

3.3. {[}gedaagde 1{]} en {[}gedaagde 2{]} c.s. voeren verweer en willen dat de vorderingen van de curator worden afgewezen, met veroordeling van de curator in de proceskosten, voor {[}gedaagde 2{]} uitvoerbaar bij voorraad.

3.4. Op de stellingen van partijen wordt hierna, voor zover nodig, verder ingegaan.

\paragraph{4 De beoordeling in beide zaken}\label{de-beoordeling-in-beide-zaken}

4.1. Aansprakelijkheid {[}gedaagde 1{]} en {[}gedaagde 2{]} op grond van artikel 2:248 BW

4.1.1. De rechtbank beoordeelt als eerste of sprake is van (kennelijk) onbehoorlijk bestuur van {[}gedaagde 1{]} en {[}gedaagde 2{]} op grond van artikel 2:248 BW. Dat bepaalt dat in geval van faillissement van de vennootschap iedere bestuurder tegenover de boedel hoofdelijk aansprakelijk is voor het faillissementstekort als het bestuur zijn taak kennelijk onbehoorlijk heeft vervuld en aannemelijk is dat dit een belangrijke oorzaak is van het faillissement. Van kennelijk onbehoorlijke taakvervulling in de zin van dit artikel kan alleen worden gesproken als geen redelijk denkend bestuurder onder dezelfde omstandigheden aldus gehandeld zou hebben.\footnote{Hoge Raad 8 juni 2001, ECLI:NL:HR:2001:AB2021 (Panmo).} Artikel 2:248 lid 2 BW bepaalt dat als het bestuur niet heeft voldaan aan zijn verplichtingen uit artikel 2:10 BW (de administratieplicht) of 2:394 BW (publicatieplicht), het zijn taak onbehoorlijk heeft vervuld en wordt vermoed dat onbehoorlijke taakvervulling een belangrijke oorzaak is van het faillissement. Dit wettelijk vermoeden kan worden ontzenuwd door de bestuurder door aannemelijk te maken dat (ook) andere feiten of omstandigheden dan zijn onbehoorlijke taakvervulling een belangrijke oorzaak van het faillissement zijn geweest.

\emph{De administratie voldoet niet aan de vereisten}

4.1.2. De curator stelt allereerst dat sprake is van kennelijk onbehoorlijk bestuur omdat niet is voldaan aan de administratieplicht. De administratie van {[}eiser 2{]} en {[}eiser 1{]} voldoet volgens de curator niet aan de vereisten van artikel 2:10 lid 1 BW, onder andere omdat:

\begin{itemize}
\item
  in de administratie nauwelijks onderscheid is gemaakt tussen {[}eiser 2{]} en {[}eiser 1{]} ,
\item
  een projectadministratie per schip ontbreekt, zodat niet inzichtelijk is wat de kosten en de werkelijke verkoopopbrengst van een schip zijn,
\item
  door de curator geen debiteuren- en crediteurenlijsten zijn aangetroffen,
\item
  privékosten als zakelijke kosten werden geboekt en niet van zakelijke kosten te onderscheiden zijn,
\item
  er geen leningsovereenkomsten met bestuurders of gelieerde entiteiten zijn aangetroffen, terwijl er wel vaak leningen werden verstrekt,
\item
  facturen en overeenkomsten geantedateerd werden,
\item
  jaarrekeningen komen niet overeen met de gevoerde administratie en sluiten niet aan op eerdere boekjaren.
\end{itemize}

4.1.3. Volgens artikel 2:10 lid 1 BW is het bestuur verplicht om van de vermogenstoestand en van alles over de werkzaamheden van een vennootschap, naar de eisen die uit deze werkzaamheden voortvloeien, zo'n administratie te voeren en te bewaren, dat daarin altijd de rechten en verplichtingen van de vennootschap kunnen worden gekend.

\emph{Projectadministratie ontbreekt}

4.1.4. De rechtbank oordeelt dat de administratie van zowel {[}eiser 2{]} als {[}eiser 1{]} niet voldoet aan de eisen van artikel 2:10 lid 1 BW. Uit de door de curator aangetroffen administratie kunnen namelijk niet de rechten en verplichtingen van {[}eiser 2{]} en {[}eiser 1{]} worden gekend. Ten eerste niet omdat een projectadministratie per schip ontbreekt. Daardoor was het niet mogelijk om een realistische en winstgevende verkoopprijs voor een schip te bepalen waarmee de kosten die voor een schip waren gemaakt (bijvoorbeeld voor een reparatie) terugverdiend konden worden. Weliswaar was, zoals door {[}gedaagde 1{]} naar voren gebracht, een projectadministratie door de globalisatieregeling niet vereist, maar het voeren van zo'n administratie per schip was wel nodig met het oog op de activiteiten van {[}eiser 2{]} en {[}eiser 1{]} : het opkopen en daarna repareren en verkopen van schepen. Daarbij komt dat door het ontbreken van een projectadministratie {[}gedaagde 1{]} en {[}gedaagde 2{]} ook geen inzicht hadden in de waarde van de activa van {[}eiser 2{]} en {[}eiser 1{]} . Dat er volgens {[}gedaagde 1{]} vóór de breuk in oktober 2020 wel een projectadministratie was, is uit de door hem overgelegde foto's niet af te leiden en is verder ook niet gebleken.

4.1.5. Dat er geen inzicht was in de kosten van een schip wordt ook bevestigd door {[}naam 4{]} in zijn verklaring bij de rechter-commissaris op 17 oktober 2022:

``\emph{{[}gedaagde 2{]} had geen idee van hoeveel een boot moest kosten. Ik heb dit in kaart gebracht. De winstmarge was ten opzichte van de omzet bijzonder laag. {[}gedaagde 2{]} kocht bijvoorbeeld een boot in Frankrijk of Zuid Spanje voor 2 ton, deze moest opgeknapt worden. Hij had geen goed idee hoeveel alles kostte om het op te knappen, dat kon dan bijvoorbeeld ook weer 2 ton kosten. Daarbij moest er dan bijvoorbeeld ook een motor in van 60.000 euro. Hij maakte geen calculatie. Zo'n boot werd daarna dan voor bijvoorbeeld 3 ton verkocht en dan dacht hij dat hij een mooie verkoop had gedaan. Ik heb gezegd dat hij moest bijhouden hoeveel geld hij in een boot stopte. Ze wisten eigenlijk niet waar ze mee bezig waren}.''

4.1.6. Het verlies van {[}eiser 2{]} is, zo volgt uit de verklaring van {[}naam 4{]} in de concept jaarrekening 2019, het gevolg van de vaststelling van onjuiste verkoopprijzen:

``\emph{Het verlies over 2019 wordt veroorzaakt doordat onjuiste verkoopprijzen voor de boten zijn vastgesteld. De brutowinst bedraagt slechts 1,5 \% van de netto-omzet, hetgeen onvoldoende is om alle bedrijfskosten van {[}eiser 2{]} BV te dekken. Door de verliezen van {[}eiser 2{]} is de continuïteit niet meer gegarandeerd en is een faillissement niet ondenkbaar. Doordat geen marktconforme/kostendekkende verkoopprijs wordt gehanteerd (met name voor leveringen aan {[}eiser 1{]}} {[} {[}eiser 1{]} , rb{]}\emph{) wordt {[}eiser 2{]} benadeeld en {[}eiser 1{]} bevoordeeld. Aangezien beide ondernemingen hetzelfde bestuur hebben is sprake van onzakelijk handelen/besmette transacties. In geval van faillissement kan de curator de directie van {[}eiser 2{]} privé aansprakelijk stellen i.v.m. onbehoorlijk bestuur en bestuurdersaansprakelijkheid inroepen.}''

\emph{Geen onderscheid {[}eiser 2{]} en {[}eiser 1{]} en aanpassingen zonder rechtvaardiging}

4.1.7. Ten tweede kunnen uit de administratie de rechten en verplichtingen van {[}eiser 2{]} en {[}eiser 1{]} niet worden gekend door het vele betalingsverkeer tussen de vennootschappen en het onderling verstrekken van leningen, zonder dat daarvoor een verantwoording is gevonden (zie 2.4.1 en 2.4.2). Feitelijk werd alles op één hoop gegooid. {[}gedaagde 1{]} en {[}gedaagde 2{]} hebben dit ook niet weersproken. Zij hebben aangegeven dat {[}eiser 2{]} en {[}eiser 1{]} `communicerende vaten' waren, maar dat rechtvaardigt deze werkwijze van administratie voeren niet. Daarnaast hebben ook de vele aanpassingen in de administratie (zie 2.4.4 en 2.4.5) ervoor gezorgd dat deze niet voldoet aan de eisen van artikel 2:10 lid 1 BW. {[}gedaagde 1{]} en {[}gedaagde 2{]} hebben ook niet betwist dat die aanpassingen na overleg met hen hebben plaatsgevonden. {[}gedaagde 1{]} wordt niet gevolgd in zijn standpunt dat deze zijn gedaan uit boekhoudkundige en fiscale overwegingen. Zonder verdere toelichting is dit namelijk geen rechtvaardiging voor de aanpassingen. Daarbij komt dat aannemelijk is dat de administratie is aangepast om een positief beeld van {[}eiser 2{]} en {[}eiser 1{]} te geven, bijvoorbeeld om financiering te krijgen. Zo meldt {[}naam 4{]} op 28 september 2020 aan {[}gedaagde 1{]} en {[}gedaagde 2{]} , na aanpassing van de cijfers: ``\emph{Vanwege de hypotheekaanvraag is het van belang dat het resultaat in {[}eiser 2{]} BV positief is}''. Bovendien hebben de aanpassingen ervoor gezorgd dat de administratie niet betrouwbaar is en niet overeenkomt met de werkelijkheid. Ook dit maakt dat daaruit niet de rechten en verplichtingen van de {[}eiser 2{]} en {[}eiser 1{]} kunnen worden gekend.

\emph{Privé-onttrekkingen, vermenging van zakelijk en privé en onzakelijke transacties}

4.1.8. Ook privé-onttrekkingen, vermenging van zakelijk en privé en onzakelijke transacties hebben ervoor gezorgd dat uit de administratie niet de rechten en verplichtingen van {[}eiser 2{]} en {[}eiser 1{]} kunnen worden gekend. Illustratief daarvoor is naar het oordeel van de rechtbank de werkwijze van {[}gedaagde 1{]} en {[}gedaagde 2{]} met betrekking de {[}naam boot 1{]} , Moody en de vakantiewoning (zie 2.2 en 2.3). Daarvoor geldt het volgende.

4.1.9. De rechtbank volgt de curator in zijn standpunt dat {[}eiser 2{]} door de overdracht van de {[}naam boot 1{]} een direct verlies van € 135.000,00 heeft geleden en een winst van € 150.000,00 is misgelopen. {[}gedaagde 1{]} heeft hiertegen ingebracht dat {[}eiser 2{]} niet is benadeeld omdat {[}eiser 2{]} de {[}naam boot 1{]} heeft ingekocht voor € 135.000,00 en voor € 150.000,00 heeft verkocht aan {[}gedaagde 1{]} , met als resultaat een winst van € 15.000,00 voor {[}eiser 2{]} . Anders dan {[}gedaagde 1{]} naar voren heeft gebracht, is echter niet gebleken dat hij {[}eiser 2{]} € 150.000,00 voor de {[}naam boot 1{]} heeft betaald. Dit volgt niet uit het door {[}gedaagde 1{]} overgelegde bankafschrift. De curator heeft voldoende onderbouwd dat {[}gedaagde 1{]} zich de door {[}eiser 2{]} voor € 135.000,00 gekochte {[}naam boot 1{]} heeft toegeëigend en daarna de helft van de {[}naam boot 1{]} voor € 150.000,00 aan een derde heeft verkocht. Het gaat hier dus om een vermenging van zakelijk en privé, waarbij {[}gedaagde 1{]} (in privé) de {[}naam boot 1{]} aan het vermogen van {[}eiser 2{]} heeft onttrokken.

4.1.10. Ook voor de rentebetalingen van € 48.000,00 door {[}eiser 1{]} aan {[}naam 3{]} voor de Moody, die eigendom is van {[}gedaagde 2{]} in privé, geldt dat sprake is van vermenging van zakelijk en privé en onttrekking van vermogen aan {[}eiser 1{]} . De Moody is volgens het Kadaster een schip van {[}gedaagde 2{]} , waarop een hypotheekrecht is gevestigd ten gunste van {[}naam 3{]} . Dat de Moody volgens {[}gedaagde 2{]} eigendom is van {[}eiser 2{]} en dat geen zekerheidsrecht door hem is gevestigd vindt geen steun in de feiten.

4.1.11. Daarnaast is met de transactie van de vakantiewoning, waarbij {[}eiser 2{]} een verlies van € 169.200,00 heeft geleden en {[}bedrijf 2{]} een winst van € 97.500,00 heeft genoten, sprake geweest van een onzakelijke transactie. De curator wordt ook hier gevolgd in zijn standpunt. {[}gedaagde 1{]} betwist dat {[}eiser 2{]} een boekverlies heeft geleden. Volgens hem waren de twee boten die voor de aankoop van de vakantiewoning zijn ingeruild maar € 210.000,00 waard en is de zakelijke waarde van de vakantiewoning vastgesteld op € 300.000,00, zodat {[}eiser 2{]} een winst van € 90.000,00 heeft behaald. Zonder verdere onderbouwing oordeelt de rechtbank dat {[}gedaagde 1{]} hiermee het standpunt van de curator onvoldoende gemotiveerd heeft betwist. Hiervoor is ook van belang de e-mail van 17 september 2018 van {[}naam 4{]} waarin hij {[}gedaagde 2{]} en {[}gedaagde 1{]} erop heeft gewezen dat de doorverkoop van de vakantiewoning door {[}eiser 2{]} aan {[}bedrijf 2{]} voor een bedrag van € 300.000 als een niet zakelijke transactie gezien kan worden vanwege de ``\emph{lage}'' verkoopprijs ``\emph{temeer daar de vakantiewoning enige weken hiervoor is verkocht voor € 460.000}''. Ook de stelling van {[}gedaagde 2{]} dat de curator moet aantonen welke waarde reëel was en welke zakelijke redenen er destijds speelden is onvoldoende onderbouwd. Daarmee miskent {[}gedaagde 2{]} dat het aan hem is om het standpunt van de curator over het onzakelijke karakter van de transactie van de vakantiewoning gemotiveerd te betwisten.

\emph{Kennelijk onbehoorlijk bestuur door schending administratieplicht}

4.1.12. De conclusie is dat de administratieplicht van artikel 2:10 BW is geschonden. Daarmee staat op grond van artikel 2:248 lid 2 BW vast dat het bestuur zijn taak onbehoorlijk heeft vervuld en wordt (weerlegbaar) vermoed dat dit een belangrijke oorzaak is van het faillissement. Door schending van de administratieplicht is dus formeel sprake van kennelijk onbehoorlijk bestuur op grond van artikel 2:248 BW. De hiervoor genoemde aanpassingen van de administratie, privé-onttrekkingen en onzakelijke transacties kwalificeren op zichzelf en in samenhang bezien ook materieel als kennelijk onbehoorlijk bestuur in de zin van artikel 2:248 lid 1 BW. {[}gedaagde 2{]} en {[}gedaagde 1{]} hebben daarmee voorzienbaar ten nadele van de schuldeisers van {[}eiser 2{]} en {[}eiser 1{]} gehandeld en vrijwel uitsluitend henzelf bevoordeeld. De rechtbank volgt de curator in zijn standpunt dat {[}gedaagde 2{]} en {[}gedaagde 1{]} met deze handelwijze zodanig roekeloos zijn omgegaan met de belangen van {[}eiser 2{]} en {[}eiser 1{]} en hun schuldeisers dat geen redelijk denkend bestuurder onder dezelfde omstandigheden aldus gehandeld zou hebben. Anders dan {[}gedaagde 2{]} tijdens de zitting naar voren heeft gebracht, zijn deze handelingen voldoende voor dit oordeel.

\emph{Schending administratieplicht is belangrijke oorzaak van het faillissement}

4.1.13. Op grond van artikel 2:248 lid 2 BW wordt de schending van de administratieplicht reeds vermoed een belangrijke oorzaak van de faillissementen van {[}eiser 2{]} en {[}eiser 1{]} te zijn. De curator heeft ter nadere onderbouwing gesteld dat het hiervoor genoemde handelen van {[}gedaagde 2{]} en {[}gedaagde 1{]} afzonderlijk en in samenhang bezien een belangrijke oorzaak van de faillissementen is geweest, omdat als:

\begin{itemize}
\tightlist
\item
  de administratie van {[}eiser 2{]} en {[}eiser 1{]} daadwerkelijk inzage had gegeven in hun rechten
\end{itemize}

en verplichtingen, zij geen nieuwe leningen hadden kunnen aangaan en de schuldenlast niet zo hoog had kunnen oplopen,

\begin{itemize}
\tightlist
\item
  het bestuur geen dusdanig hoge privéonttrekkingen had gedaan of als deze
\end{itemize}

privéonttrekkingen tijdig waren terugbetaald, {[}eiser 2{]} en {[}eiser 1{]} (bijna) al hun

schuldeisers hadden kunnen voldoen,

\begin{itemize}
\tightlist
\item
  het bestuur geen onzakelijke transacties had verricht waarbij de kosten ten laste van
\end{itemize}

{[}eiser 2{]} en {[}eiser 1{]} waren gekomen en de opbrengsten ten gunste van het bestuur of gelieerde vennootschappen, de opbrengsten aan {[}eiser 2{]} en {[}eiser 1{]} waren toegekomen en {[}eiser 2{]} en {[}eiser 1{]} over voldoende liquide middelen hadden beschikt en hun schuldeisers kunnen voldoen.

\emph{Bewijsvermoeden is niet ontzenuwd}

4.1.14. {[}gedaagde 1{]} en {[}gedaagde 2{]} hebben ter weerlegging van het bewijsvermoeden aangevoerd dat niet de schending van de administratieplicht, maar de breuk tussen hen, medio oktober 2020, de oorzaak is geweest van de faillissementen. De rechtbank volgt hen hierin niet. Aan {[}gedaagde 1{]} kan worden toegegeven dat na de breuk met {[}gedaagde 2{]} veel onregelmatigheden hebben plaatsgevonden, maar daar staat tegenover dat {[}gedaagde 1{]} vóór medio oktober 2020 ook in aanzienlijke mate zelf bij onregelmatigheden betrokken was. Voor zover de breuk tussen {[}gedaagde 1{]} en {[}gedaagde 2{]} wel als een oorzaak van het faillissement aangemerkt zou kunnen worden, geldt dat dit niet een externe van buitenaf komende oorzaak is en dat de breuk en verdere escalatie door {[}gedaagde 1{]} en {[}gedaagde 2{]} zelf zijn veroorzaakt en de gevolgen daarvan voor hun rekening komen. Niet gebleken is bovendien dat, zoals {[}gedaagde 1{]} naar voren heeft gebracht, de zaken goed liepen vóór de breuk medio oktober 2020. {[}naam 4{]} had in januari 2020 al gewaarschuwd dat de continuïteit van {[}eiser 2{]} en {[}eiser 1{]} mogelijk gevaar liep (zie 2.4.6). In juni 2020 rapporteerde hij voor {[}eiser 2{]} een verlies van € 839.510,00 en voor {[}eiser 1{]} een winst van € 127.638,00. Ook daarna heeft {[}naam 4{]} in september 2020 {[}gedaagde 1{]} en {[}gedaagde 2{]} gewezen op de schulden, onder meer die aan de Belastingdienst, en dat een faillissement van {[}eiser 2{]} niet ondenkbaar was.

4.1.15. Tijdens de zitting heeft {[}gedaagde 2{]} voor de oorzaak van de faillissementen nog genoemd dat het bij {[}eiser 2{]} en {[}eiser 1{]} ging om ondernemingen met lage marges, hoge inkoopbedragen, hoge opknapkosten, financieringsdruk waarbij gelden noodgedwongen werden geleend tegen excessieve rentevergoedingen, afhankelijkheid van consignatie en de coronapandemie. De genoemde omstandigheden zijn zeer algemeen van aard en betreffen het bedrijfsmodel van {[}eiser 2{]} en {[}eiser 1{]} zoals dat al jaren bestond. Daarmee heeft {[}gedaagde 2{]} dus niet aannemelijk gemaakt dat andere omstandigheden dan de kennelijk onbehoorlijke taakvervulling een belangrijke oorzaak van het faillissement zijn geweest. Dat zou mogelijk anders kunnen zijn voor de coronapandemie, maar op dat punt ontbreekt een voldoende concrete toelichting.

4.1.16. De conclusie van het voorgaande is dat {[}gedaagde 2{]} en {[}gedaagde 1{]} vanwege kennelijk onbehoorlijk bestuur op grond van artikel 2:248 BW (in beginsel) aansprakelijk zijn voor het boedeltekort in de faillissementen van {[}eiser 2{]} en {[}eiser 1{]} .

\emph{Beroep van {[}gedaagde 1{]} en {[}gedaagde 2{]} op disculpatie slaagt niet}

4.1.17. De rechtbank begrijpt dat {[}gedaagde 1{]} en {[}gedaagde 2{]} zich beide beroepen op disculpatie op grond van artikel 2:248 lid 3 BW. Dat artikel bepaalt dat een bestuurder niet aansprakelijk is als hij bewijst dat het kennelijk onbehoorlijk bestuur niet aan hem is te wijten en dat hij niet nalatig is geweest in het treffen van maatregelen om de gevolgen daarvan af te wentelen. Het bestuur is in beginsel collectief verantwoordelijk voor een behoorlijke taakvervulling. Een geslaagd beroep op disculpatie is - zo blijkt uit de rechtspraak - slechts bij hoge uitzondering mogelijk.

4.1.18. Volgens {[}gedaagde 1{]} hebben de onregelmatigheden plaatsgevonden na de ruzie medio oktober 2020. {[}gedaagde 1{]} kon hier niets tegen doen omdat hij geen toegang meer had tot de administratie. {[}gedaagde 2{]} heeft zich op het standpunt gesteld dat alleen {[}gedaagde 1{]} zich bezighield met de administratie. Ook is het {[}gedaagde 1{]} geweest die privé-onttrekkingen heeft gedaan, zodat dit {[}gedaagde 2{]} niet kan worden verweten. Bovendien heeft {[}gedaagde 1{]} nog een laptop met daarop de administratie van {[}eiser 2{]} en {[}eiser 1{]} die hij kon aanpassen.

4.1.19. De rechtbank oordeelt dat {[}gedaagde 1{]} en {[}gedaagde 2{]} zich niet kunnen disculperen. De schending van de administratieplicht valt binnen de collectieve verantwoordelijkheid van het bestuur. Niet is gebleken dat alleen {[}gedaagde 1{]} zich bezighield met de administratie of dat {[}gedaagde 2{]} en {[}gedaagde 1{]} intern een taakverdeling hadden afgesproken. Het personeel van {[}eiser 2{]} en {[}eiser 1{]} heeft bij de rechter-commissaris op 17 oktober 2022 ook verklaard dat zowel {[}gedaagde 1{]} als {[}gedaagde 2{]} hen opdrachten gaven voor facturatie en dat zij hierover beide beslissingen namen. Daarbij komt dat {[}gedaagde 1{]} en {[}gedaagde 2{]} beide steeds toegang hadden tot de administratie. {[}gedaagde 2{]} , anders dan hij zelf naar voren heeft gebracht zonder meer, en {[}gedaagde 1{]} in ieder geval vóór de breuk medio oktober 2020. De hierboven besproken tekortkomingen in de administratie waren toen al aanwezig. Zowel {[}gedaagde 2{]} als {[}gedaagde 1{]} konden dus beide maatregelen treffen om de gevolgen van de schending van de administratieplicht af te wenden. Niet gebleken is dat zij dit hebben gedaan. Het beroep op disculpatie van {[}gedaagde 1{]} en {[}gedaagde 2{]} slaagt daarom niet.

\emph{Geen aanleiding voor matiging van bedrag aansprakelijkheid {[}gedaagde 2{]}}

4.1.20. heeft daarnaast een beroep gedaan op matiging van het bedrag van de aansprakelijkheid voor het boedeltekort. Op grond van artikel 2:248 lid 4 BW kan de rechter het bedrag waarvoor de bestuurders aansprakelijk zijn verminderen als hem dit bovenmatig voorkomt, i) gelet op de aard en de ernst van de onbehoorlijke taakvervulling, ii) de andere oorzaken van het faillissement, iii) de wijze waarop dit is afgewikkeld en iv) gezien de tijd die de bestuurder in functie is geweest in de periode waarin de onbehoorlijke taakvervulling plaatsvond. De in artikel 2:248 lid 4 BW genoemde matigingsgronden zijn limitatief.

4.1.21. {[}gedaagde 2{]} heeft in de conclusie van antwoord uiteengezet waarom het bedrag waarvoor hij aansprakelijk is moet worden gematigd tot nihil. Zonder verdere toelichting valt niet in te zien waarom de door hem aangevoerde redenen, afzonderlijk of in samenhang, zouden moeten leiden tot matiging, laat staat tot matiging (al dan niet tot nihil) op de gronden zoals limitatief genoemd in artikel 2:248 lid 4 BW. Zo zijn leeftijd, geringe verdiencapaciteit of afwezigheid van privaat voordeel geen gronden voor matiging. De rechtbank ziet ook geen aanleiding om ambtshalve tot matiging over te gaan.

\emph{Gevorderde verklaringen voor recht en betalingen voorschotten toewijsbaar}

4.1.22. Dit betekent dat de door de curator gevorderde verklaringen voor recht dat {[}gedaagde 1{]} en {[}gedaagde 2{]} hoofdelijk aansprakelijk zijn voor het tekort in de faillissementen van {[}eiser 1{]} en {[}eiser 2{]} , nader op te maken in de schadestaatprocedure, worden toegewezen. Ook de vorderingen van de curator tot betaling van voorschotten van € 800.000,00 in het faillissement van {[}eiser 1{]} en € 2.000.000,00 in het faillissement van {[}eiser 2{]} zijn toewijsbaar, omdat gezien de stand van de boedels zoals vermeld onder 2.7.1 de boedeltekorten in de beide faillissementen tenminste die bedragen zullen belopen. De gevorderde wettelijke rente over de boedeltekorten en te betalen voorschotten wordt niet toegewezen. Artikel 2:248 BW biedt hiervoor geen grondslag. Omdat {[}gedaagde 1{]} en {[}gedaagde 2{]} aansprakelijk zijn voor het boedeltekort, zijn de gevorderde buitengerechtelijke kosten en de proceskosten, inclusief de van {[}gedaagde 2{]} gevorderde beslagkosten, niet toewijsbaar. Deze kosten worden (in het salaris en de verschotten van de curator) al meegeteld bij de berekening van het boedeltekort.

4.1.23. Voor zover {[}gedaagde 1{]} en {[}gedaagde 2{]} een beroep doen op verrekening van hun eventuele vorderingen op {[}eiser 2{]} en {[}eiser 1{]} met de vordering van de curator tot betaling van het boedeltekort gaat dit niet op. Verrekening is niet mogelijk op grond van artikel 2:248 lid 6 BW.

4.1.24. Met veroordeling van {[}gedaagde 1{]} en {[}gedaagde 2{]} tot betaling van het boedeltekort in de faillissementen wordt niet toegekomen aan beoordeling van de andere afzonderlijke individuele vorderingen van de curator op {[}gedaagde 1{]} en {[}gedaagde 2{]} zoals die uit hoofde van rekening-courant verhoudingen, paulianeus handelen of onbehoorlijke taakvervulling op grond van artikel 2:9 BW. Dit zijn de vorderingen genoemd in 3.1 en de vorderingen in 3.2 met nummers 8, 9, 10, 11, 16 en 17. Deze worden dan ook afgewezen bij gebrek aan belang.

4.2. {[}gedaagde 4{]}

\emph{{[}gedaagde 4{]} moet € 12.486,48 aan {[}eiser 2{]} betalen, met rente en kosten}

4.2.1. Volgens de curator volgt uit de administratie van {[}eiser 2{]} een rekening-courantvordering van {[}eiser 2{]} op {[}gedaagde 4{]} van € 12.486,48 op datum faillissement. Dit komt nagenoeg overeen met de kolommenbalansen van zowel {[}gedaagde 4{]} als {[}eiser 2{]} van 26 mei 2021 waaruit volgt dat {[}eiser 2{]} uit hoofde van rekening-courant € 12.586,48 van {[}gedaagde 4{]} te vorderen had. De curator verwijst hiervoor naar de gegevens in de door hem veiliggestelde administratie. De rechtbank oordeelt dat, bij gebrek aan gemotiveerde betwisting van {[}gedaagde 4{]} , voldoende is komen vast te staan dat {[}eiser 2{]} € 12.486,48 van {[}gedaagde 4{]} te vorderen heeft. Deze vordering van de curator is daarom toewijsbaar. Voldoende is gebleken dat de curator hiervoor buitengerechtelijke werkzaamheden heeft verricht. De curator heeft {[}gedaagde 4{]} op 3 augustus 2023 verzocht om betaling binnen zeven dagen. Daarmee zijn de gevorderde buitengerechtelijke incassokosten toewijsbaar, en wel tot € 899,86. De door de curator gevorderde wettelijke rente vanaf 11 augustus 2023 wordt eveneens toegewezen.

\emph{{[}gedaagde 4{]} moet € 20.500,00 aan {[}eiser 1{]} betalen, met rente en kosten}

4.2.2. De curator vordert van {[}gedaagde 4{]} € 20.500,00 uit hoofde van de rekening-courantverhouding met {[}eiser 1{]} . Volgens de administratie van {[}eiser 1{]} had zij eind december 2020 een rekening-courantvordering van € 20.500,00 op {[}gedaagde 4{]} . Na de beslaglegging door de particuliere financiers is deze vordering met het opmaken van een (geantedateerde) factuur van {[}gedaagde 4{]} aan {[}eiser 1{]} teruggebracht naar nihil (zie 2.6.6). De curator stelt dat niet is gebleken dat de werkzaamheden op die factuur daadwerkelijk hebben plaatsgevonden. {[}gedaagde 4{]} heeft dit niet gemotiveerd weersproken. Bij deze stand van zaken is de rechtbank van oordeel dat voldoende is komen vast te staan dat {[}eiser 1{]} een rekening-courantvordering van € 20.500,00 op {[}gedaagde 4{]} heeft. Deze vordering van de curator wordt dan ook toegewezen. Dat geldt ook voor de gevorderde buitengerechtelijke incassokosten en wettelijke rente vanaf 11 augustus 2023 (zie 4.2.1). Bij toewijzing van de hoofdsom van € 20.500,00 zijn de buitengerechtelijke incassokosten toewijsbaar tot € 980,00.

4.3. {[}gedaagde 3{]}

\emph{Oneigenlijke doorstart}

4.3.1. De curator stelt dat een oneigenlijke doorstart van de onderneming van {[}eiser 2{]} en {[}eiser 1{]} heeft plaatsgevonden door {[}gedaagde 3{]} zonder dat {[}gedaagde 3{]} voor de verkrijging van vermogensbestanddelen van {[}eiser 2{]} en {[}eiser 1{]} een tegenprestatie heeft verricht. Bij brief van 28 maart 2024 heeft de curator het samenstel van rechtshandelingen die hebben geleid tot de oneigenlijke doorstart (buitengerechtelijk) vernietigd met een beroep op de actio Pauliana op grond van artikel 42 Fw. De curator vordert een verklaring voor recht dat hij rechtsgeldig tot vernietiging is overgegaan en betaling van een waardevergoeding, althans verwijzing naar de schadestaat onder voldoening van een voorschot.

4.3.2. De rechtbank oordeelt dat de curator bij brief van 28 maart 2024 de rechtshandelingen waarbij om niet vermogensbestanddelen van {[}eiser 2{]} en {[}eiser 1{]} aan {[}gedaagde 3{]} zijn overgedragen en die hebben geleid tot de oneigenlijke doorstart van de activiteiten van {[}eiser 2{]} en {[}eiser 1{]} via {[}gedaagde 3{]} rechtsgeldig heeft vernietigd. Daarvoor geldt het volgende. Op grond van artikel 42 lid 1 Fw kan de curator ten behoeve van de boedel elke onverplichte rechtshandeling die de schuldenaar vóór de faillietverklaring heeft verricht en waarvan hij wist of behoorde te weten dat daarmee schuldeisers zouden worden benadeeld, buitengerechtelijk vernietigen. Hierbij geldt volgens artikel 42 lid 2 Fw dat rechtshandelingen, anders dan om niet, die meerzijdig zijn of eenzijdig gericht op bepaalde personen, alleen vernietigd kunnen worden als ook de wederpartij wetenschap had of behoorde te hebben van de benadeling. Een rechtshandeling is onverplicht als daarvoor geen wettelijke of contractuele verplichting bestaat. Benadeling van schuldeisers betekent dat hun verhaal beperkt wordt en moet aanwezig zijn op het moment dat de curator zijn rechten uitoefent. Dit wordt beoordeeld door de hypothetische situatie zonder de rechtshandeling te vergelijken met de feitelijke situatie met de rechtshandeling. Van wetenschap van benadeling is sprake als bij de rechtshandeling het faillissement en een tekort met redelijke zekerheid voorzienbaar waren voor zowel schuldenaar als wederpartij. Artikel 43 Fw bevat rechtsvermoedens van deze wetenschap, die gelden behalve als tegenbewijs wordt geleverd.

4.3.3. Volgens de curator ziet de oneigenlijke doorstart op de volgende rechtshandelingen:

\begin{itemize}
\item
  ter beschikking stellen van personeel aan {[}gedaagde 3{]} ,
\item
  overdragen van rechten uit de huurovereenkomst,
\item
  overdragen van de materiële vaste activa,
\item
  overdragen van rechten uit andere (duur)overeenkomsten,
\item
  ten laste van {[}eiser 2{]} en {[}eiser 1{]} en ten gunste van {[}gedaagde 3{]} laten komen van betalingen die betrekking hadden op {[}gedaagde 3{]} ,
\item
  het overdragen van de domeinnamen, handelsnamen, websites, logo's, klantenbestanden en onderhanden werk (goodwill),
\item
  het overdragen van de exploitatie van {[}website{]} .
\end{itemize}

4.3.4. Niet weersproken is dat {[}gedaagde 3{]} dezelfde activiteiten als {[}eiser 2{]} en {[}eiser 1{]} exploiteert, gebruikmaakt van hetzelfde personeel, bedrijfspand, logo's, briefpapier, klantenbestand en website. Ook staat vast dat op 30 september 2021 het woord ` {[}gedaagde 3{]} ' aan de website is toegevoegd en dat deze nu ongewijzigd door {[}gedaagde 3{]} wordt geëxploiteerd. Voldoende is gebleken dat {[}gedaagde 3{]} na haar oprichting de handelsactiviteiten van {[}eiser 2{]} en {[}eiser 1{]} heeft voortgezet. Een medewerker van {[}eiser 2{]} / {[}eiser 1{]} heeft hierover op 17 oktober 2022 bij de rechter-commissaris verklaard:

``\emph{De curator vraagt welke werkzaamheden in {[}gedaagde 3{]} werden verricht, Volgens mij niet veel. Op een gegeven moment werden facturen vanuit {[}gedaagde 3{]} verstuurd. Toen {[}eiser 2{]} op klappen stond, is {[}gedaagde 3{]} opgericht en kort daarna begon dit factureren. {[}gedaagde 2{]} heeft op een moment verteld dat hij een nieuwe B.V, namelijk {[}gedaagde 3{]} had opgericht zodat wij buiten de ellende zouden blijven staan. Dus ja, {[}gedaagde 3{]} is opgericht omdat er binnen {[}eiser 2{]} problemen waren. {[}gedaagde 2{]} wilde niet dat wij ons zorgen maakten. Er lag natuurlijk ook beslag op de rekening. Hij kwam op een gegeven moment naar mij toe, ergens in februari/maart 2021, dat {[}eiser 2{]} het niet zou gaan redden, op dat moment kregen wij zoveel brieven binnen dat het duidelijk was dat het niet goed ging met de onderneming. Daarom heeft hij {[}gedaagde 3{]} opgericht, dan had hij een B.V. erbij van waaruit hij kon factureren. Dit was volgens mij enkele maanden na oktober 2020 toen {[}gedaagde 1{]} wegging bij de onderneming, dus inderdaad rond februari/maart 2021.}''

En {[}gedaagde 2{]} zelf:

``\emph{De curator vraagt of de opbrengst van dit bemiddelingswerk, ten tijde van voor faillissement en na inbetalinggeving, bij mij privé, bij CV of bij {[}gedaagde 3{]} terecht is gekomen. Dat is in CV of {[}gedaagde 3{]} gekomen. Ik wilde na het gedoe met {[}gedaagde 1{]} de B.V. 's even met rust laten.}''

4.3.5. {[}gedaagde 3{]} heeft niet gemotiveerd betwist dat zij feitelijk de activiteiten en daarmee vermogensbestanddelen van {[}eiser 2{]} en {[}eiser 1{]} , waaronder (onder meer) domeinnamen, huurcontracten en goodwill, heeft overgenomen zonder dat zij daar een vergoeding voor heeft betaald.

4.3.6. Hierdoor zijn de schuldeisers in de faillissementen van {[}eiser 2{]} en {[}eiser 1{]} benadeeld. Als gevolg van de voortzetting van de activiteiten via {[}gedaagde 3{]} zijn de (nog resterende) activa van {[}eiser 2{]} en {[}eiser 1{]} naar {[}gedaagde 3{]} overgeheveld (zie 2.6.4) en zijn de opbrengsten van {[}eiser 2{]} en {[}eiser 1{]} via facturering aan {[}gedaagde 3{]} toegekomen (zie ook hiervoor in 4.3.4). Ook is zo een doorstart uit de faillissementen van {[}eiser 2{]} en {[}eiser 1{]} onmogelijk gemaakt. Het vermoeden dat zowel {[}eiser 2{]} en {[}eiser 1{]} aan de ene kant en {[}gedaagde 3{]} aan de andere kant, welke vennootschappen allemaal (feitelijk) bestuurd werden door {[}gedaagde 2{]} , op grond van artikel 45 en artikel 43 lid 1 aanhef en onder sub 1 en sub 5 onder b Fw wetenschap hadden van die benadeling, is niet weerlegd. {[}gedaagde 3{]} wordt niet gevolgd in haar standpunt dat bij de oprichting nog geen sprake was van een faillissement. {[}naam 4{]} had in 2020 al gewaarschuwd voor een mogelijk gevaar voor de continuïteit van {[}eiser 2{]} en een mogelijk faillissement, en herhaalde dit bij de concept-jaarrekening van 2019 (zie ook 4.1.6).

4.3.7. De conclusie is dat de curator bij brief van 28 maart 2024 de rechtshandelingen, waarbij vermogensbestanddelen van {[}eiser 2{]} en {[}eiser 1{]} om niet zijn overgedragen aan {[}gedaagde 3{]} die hebben geleid tot de oneigenlijke doorstart van de activiteiten van {[}eiser 2{]} en {[}eiser 1{]} via {[}gedaagde 3{]} , rechtsgeldig heeft vernietigd. De door de curator gevorderde verklaring voor recht wordt dan ook toegewezen.

\emph{Rechtsgevolg van de vernietiging: waardevergoeding}

4.3.8. Volgens de curator kunnen de reeds ingetreden gevolgen van de rechtshandelingen bezwaarlijk ongedaan worden gemaakt. Hij heeft daarom betaling van een waardevergoeding, althans verwijzing naar de schadestaatprocedure gevorderd onder voldoening van een voorschot. {[}gedaagde 3{]} heeft de waarde van de door de curator opgevoerde vermogensbestanddelen betwist. Omdat de curator de waarde van de overgehevelde activa heeft gespecificeerd, zal de rechtbank waar mogelijk de waarde zelf begroten en voor zover dit niet kan, de procedure verwijzen naar de schadestaatprocedure.

4.3.9. De curator van {[}eiser 1{]} vordert van {[}gedaagde 3{]} op grond van de vernietiging ex artikel 42 Fw € 232.284,34, bestaande uit:

\begin{enumerate}
\def\labelenumi{\alph{enumi}.}
\item
  salarisbetalingen voor personeel: € 5.937,34,
\item
  betalingen handelscrediteuren: € 26.347,00,
\item
  softe activa: € 200.000,00.
\end{enumerate}

4.3.10. De curator van {[}eiser 2{]} vordert in totaal € 510.585,70, bestaande uit:

\begin{enumerate}
\def\labelenumi{\alph{enumi}.}
\setcounter{enumi}{3}
\item
  salarisbetalingen voor personeel: € 98.191,62,
\item
  loonvordering UWV: € 36.662,44,
\item
  kosten verbouwing bedrijfspand inclusief afschrijving: € 79.468,44,
\item
  materiële vaste activa: € 12.798,72,
\item
  betalingen handelscrediteuren {[}gedaagde 3{]} : € 38.692,14,
\item
  betalingen inkoopkosten boten {[}gedaagde 3{]} : € 20.396,42,
\item
  botengids jan-sept 2021: € 2.390,00,
\item
  botengids oktober 2021 tot september 2024: € 21.985,92,
\item
  botengids vanaf oktober 2024: p.m.
\item
  softe activa: € 200.000,00.
\end{enumerate}

\emph{Verwijzing naar schadestaat voor softe activa, kosten verbouwing en botengids}

4.3.11. Volgens {[}gedaagde 3{]} is ten onrechte tweemaal softe activa (c en m) door de curator opgevoerd en is het bedrag van € 200.000,00 `uit de lucht gegrepen'. Het gaat hier om de goodwill, bestaande uit domeinnamen, de website, huisstijl, logo en handelsnamen, maar onduidelijk is nog waarom hieraan een waarde van tweemaal € 200.000,00 moet worden toegekend. Dit zal in de schadestaatprocedure moeten worden bepaald. Met betrekking tot de verbouwingskosten van het bedrijfspand (f) eind 2019, die volgens de curator € 99.335,55 hebben bedragen, is nog onduidelijk waarom {[}gedaagde 3{]} daarvan met de overname van de huurovereenkomst een bedrag van € 79.468,44 moet vergoeden, ook gelet op de overname door {[}gedaagde 3{]} van de openstaande huurschulden van {[}eiser 2{]} en {[}eiser 1{]} . Tot slot zal de waarde van de website `botengids' (j, k en l), waarvan niet weersproken is dat die na oktober 2021 door {[}gedaagde 3{]} wordt geëxploiteerd, ook in de schadestaatprocedure definitief moeten worden vastgesteld.

\emph{Posten betaling handelscrediteuren, inkoopkosten boten en loonbetaling niet toewijsbaar}

4.3.12. Voor de door de curator gespecificeerde posten loonbetalingen personeel, betaling handelscrediteuren en inkoopkosten boten (a, b, d, e, h en i) is niet duidelijk waarom deze voor vergoeding in aanmerking komen, althans heeft de curator dit onvoldoende onderbouwd gelet op de grondslag van zijn vordering, te weten: vernietiging ex artikel 42 Fw van de rechtshandelingen die tot de oneigenlijke doorstart hebben geleid. Deze gevorderde posten worden dan ook afgewezen.

\emph{{[}gedaagde 3{]} moet € 12.798,72 aan {[}eiser 2{]} betalen}

4.3.13. De rechtbank oordeelt dat de post materiële vaste activa (g) voldoende is komen vast te staan. Op de balans van {[}eiser 2{]} van 2020 staat na afschrijvingen een bedrag van € 4.020,72 aan kantoorinventaris, een bedrag van € 7.978,00 aan vervoer- en transportmiddelen en een bedrag van € 800,00 aan computers en printers geactiveerd. Dit maakt een totaal van € 12.798,72 aan vaste activa. {[}gedaagde 3{]} heeft onvoldoende onderbouwd waarom dit, naar haar eigen zeggen, nihil zou moeten zijn.

\emph{{[}gedaagde 3{]} moet lening van € 6.700,00 aan {[}eiser 2{]} terugbetalen}

4.3.14. Daarnaast vordert de curator nog terugbetaling van de door {[}eiser 2{]} aan {[}gedaagde 3{]} verstrekte lening voor een bedrag van € 6.700,00. De curator heeft dit onderbouwd met bankafschriften. Daaruit volgt dat {[}eiser 2{]} en Shipcaryachts over en weer tussen 5 februari 2021 en 2 augustus 2021 diverse leningen hebben verstrekt. {[}gedaagde 3{]} heeft daarbij in totaal € 17.700,00 met omschrijving `lening' op de bankrekening van {[}eiser 2{]} voldaan, terwijl {[}eiser 2{]} € 24.400,00 met omschrijving `lening' op de bankrekening van {[}gedaagde 3{]} heeft betaald. {[}gedaagde 3{]} heeft hier niets concreets tegen ingebracht. De rechtbank oordeelt daarom dat de curator voldoende heeft onderbouwd dat {[}gedaagde 3{]} € 6.700,00 aan {[}eiser 2{]} moet terugbetalen en wijst deze vordering toe.

\emph{Wettelijke rente en buitengerechtelijke incassokosten}

4.3.15. De wettelijke rente over het totaal van de toegewezen hoofdsom van € 19.498,72 (€ 12.798,72 + € 6.700,00) wordt toegewezen vanaf 12 april 2024, de datum waarop de in de brief van de curator van 28 maart 2024 gestelde betaaltermijn is verstreken. De buitengerechtelijke incassokosten zijn, met het oog op de toegewezen hoofdsom, toewijsbaar tot € 969,99.

4.4. {[}gedaagde 5{]} CV, {[}gedaagde 2{]} en {[}gedaagde 6{]}

\emph{Verschuiving handelsvoorraad van {[}eiser 1{]} naar {[}gedaagde 5{]} CV}

4.4.1. De curator van {[}eiser 1{]} vordert van {[}gedaagde 5{]} CV, {[}gedaagde 2{]} en {[}gedaagde 6{]} (als beherend vennoot, die daarmee hoofdelijk aansprakelijk is ex artikel 18 WvK voor de schulden van de CV) een bedrag van € 448.985,00 in verband met de verschuiving in de handelsvoorraad (zie 2.6.1) van {[}eiser 1{]} naar {[}gedaagde 5{]} CV. Naast betaling van dit bedrag vordert de curator ook teruggave van de Boston Whaler, primair aan {[}eiser 1{]} en subsidiair aan {[}eiser 2{]} . De curator stelt dat {[}gedaagde 5{]} CV de opbrengsten met betrekking tot schepen die eigendom waren van {[}eiser 1{]} zonder rechtsgrond heeft ontvangen. {[}gedaagde 5{]} CV is daarmee ongerechtvaardigd verrijkt. Voor zover de opbrengst is verrekend in rekening-courant, was {[}gedaagde 5{]} CV niet bevoegd tot verrekening op grond van artikel 54 Fw, aldus de curator.

4.4.2. Volgens {[}gedaagde 2{]} was het de gewoonte om geldstromen tussen de diverse {[}gedaagde 5{]} -entiteiten te salderen via rekening-courantverhoudingen en zijn de schepen zo verkocht en afgerekend met {[}gedaagde 5{]} CV. Ook ontving {[}gedaagde 5{]} CV commissie van {[}eiser 1{]} en heeft de koper van de Astondoa 45 € 25.000,00 aan de curator betaald, aldus {[}gedaagde 2{]} .

4.4.3. De rechtbank is van oordeel dat de curator voldoende heeft onderbouwd dat {[}gedaagde 5{]} CV de opbrengsten zoals genoemd in 2.6.1 voor schepen van {[}eiser 1{]} zonder rechtsgrond heeft ontvangen. {[}eiser 1{]} , die gerechtigd was tot de ontvangsten, is hierdoor verarmd en {[}gedaagde 5{]} CV is hiermee (ongerechtvaardigd) verrijkt op grond van artikel 6:212 BW. Daarom moet {[}gedaagde 5{]} CV € 414.915,00 aan de curator van {[}eiser 1{]} betalen.

\emph{{[}gedaagde 5{]} CV/ {[}gedaagde 2{]} / {[}gedaagde 6{]} moet Boston Whaler teruggeven aan {[}eiser 1{]}}

4.4.4. De curator vordert revindicatie van de Boston Whaler of een waardevergoeding van € 60.000,00 daarvoor, primair aan {[}eiser 1{]} en subsidiair aan {[}eiser 2{]} . Aannemelijk is dat de Boston Whaler aan {[}eiser 1{]} toebehoorde, omdat {[}eiser 1{]} ook de andere schepen genoemd in 2.6.1 in eigendom had en omdat vóór aanpassing van de rekening-courantverhouding met {[}eiser 2{]} de Boston Whaler (volgens de factuur van 17 januari 2021) van {[}eiser 1{]} was (zie 2.6.2). {[}gedaagde 5{]} CV heeft niet weersproken dat zij de Boston Whaler niet rechtsgeldig in eigendom heeft gekregen en nog houdt, zodat de vordering van de curator tot teruggave (revindicatie) op grond van artikel 5:2 BW toewijsbaar is. Dit betekent dat {[}gedaagde 5{]} CV wordt veroordeeld tot teruggave van de Boston Whaler aan de curator van {[}eiser 1{]} .

\emph{{[}gedaagde 5{]} CV moet € 8.470,00 aan restitutie van verzekeringspremies aan {[}eiser 1{]} betalen}

4.4.5. Ook vordert de curator teruggave van € 8.470,00 aan een door {[}gedaagde 5{]} CV ontvangen restitutie van verzekeringspremies (zie 2.6.3). De curator heeft voldoende onderbouwd dat {[}eiser 1{]} (en niet {[}gedaagde 5{]} CV) gerechtigd was tot ontvangst van dit bedrag en {[}gedaagde 5{]} CV heeft dit niet weersproken. {[}eiser 1{]} is dus als gerechtigde tot de restitutie verarmd en {[}gedaagde 5{]} CV is met de verrekening van € 8.470,00 ongerechtvaardigd verrijkt (artikel 6:212 BW). De vordering tot teruggave van dit bedrag aan de curator van {[}eiser 1{]} is dus toewijsbaar.

\emph{{[}gedaagde 5{]} CV moet de lening van € 30.600,00 aan {[}eiser 1{]} terugbetalen}

4.4.6. Daarnaast heeft de curator terugbetaling gevorderd van de door {[}eiser 1{]} aan {[}gedaagde 5{]} CV verstrekte lening voor € 30.600,00. Hij heeft dit onderbouwd met bankafschriften. Daaruit volgt dat {[}eiser 2{]} en Shipcaryachts over en weer tussen 5 april 2019 en 27 mei 2021 diverse leningen hebben verstrekt. {[}gedaagde 5{]} CV heeft daarbij in totaal € 41.900,00 met omschrijving `lening' op de bankrekening van {[}eiser 1{]} voldaan. {[}eiser 1{]} heeft € 72.500,00 met omschrijving `lening' op de bankrekening van {[}gedaagde 5{]} CV betaald. De rechtbank is van oordeel dat de curator, zonder (gemotiveerde) betwisting van {[}gedaagde 5{]} CV, voldoende heeft onderbouwd dat {[}gedaagde 5{]} CV het verschil van € 30.600,00 van {[}eiser 1{]} heeft geleend en aan {[}eiser 1{]} moet terugbetalen. Deze vordering van de curator is dus toewijsbaar.

\emph{Vordering van de curator van {[}eiser 1{]} op {[}gedaagde 5{]} CV wordt toegewezen tot € 448.985,00}

4.4.7. In totaal zou {[}gedaagde 5{]} CV dus aan de curator van {[}eiser 1{]} moeten betalen: € 414.915,00 + 8.470,00 + 30.600,00 = € 453.985,00. Echter, in het petitum is niet meer gevorderd dan € 448.985,00, zodat de vordering tot dat bedrag wordt toegewezen.

4.4.8. De wettelijke rente over het totaal van de toegewezen hoofdsom € 448.985,00 wordt, zoals gevorderd, toegewezen vanaf 13 juni 2024, de datum waarop de in de sommatiebrief van 29 mei 2024 door de curator gestelde betaaltermijn is verstreken. Met die brief is voldoende gebleken dat de curator buitengerechtelijke werkzaamheden heeft verricht. De buitengerechtelijke incassokosten zijn, met het oog op het totaal van de toegewezen hoofdsom van € 448.985,00 toewijsbaar tot € 4.019,93. Omdat {[}gedaagde 6{]} als beherend vennoot ex artikel 18 WvK hoofdelijk aansprakelijk is voor de schulden van de CV, wordt zij ook (hoofdelijk) veroordeeld tot betaling van € 448.985,00 met rente en kosten.

4.4.9. Ook {[}gedaagde 2{]} kan hoofdelijk aansprakelijk gehouden worden voor deze schuld met rente en kosten. Hij is weliswaar commanditaire vennoot, maar niet weersproken is dat hij feitelijk voor {[}gedaagde 5{]} CV heeft opgetreden bij de verkopen van de genoemde schepen, zodat hij op grond van artikel 21 WvK evenzeer aansprakelijk is voor alle schulden en verbintenissen van {[}gedaagde 5{]} CV.

4.5. Beroep op matiging {[}gedaagde 4{]} , {[}gedaagde 3{]} , {[}gedaagde 5{]} CV en {[}gedaagde 6{]}

4.5.1. Voor zover in de conclusie van antwoord een beroep op matiging is gedaan door {[}gedaagde 4{]} , {[}gedaagde 3{]} , {[}gedaagde 5{]} CV en {[}gedaagde 6{]} geldt dat dit niet slaagt. Matiging van een wettelijke verplichting tot schadevergoeding is geregeld in artikel 6:109 BW. Deze bepaling is niet van toepassing voor zover het gaat om nakoming van contractuele verplichtingen, zoals terugbetaling van leningen (zie 4.4.6). Wel is deze bepaling van toepassing op schadevergoeding wegens ongerechtvaardigde verrijking (zie 4.4.3 en 4.4.5). Op grond van artikel 6:109 BW is alleen dan voor matiging van schadevergoeding plaats, wanneer volledige schadevergoeding in de gegeven omstandigheden, waaronder de aard van de aansprakelijkheid, de tussen partijen bestaande rechtsverhouding en hun beider draagkracht tot kennelijk onaanvaardbare gevolgen zou leiden. De in die bepaling neergelegde maatstaf maakt dat de rechter met terughoudendheid gebruik maakt van zijn bevoegdheid om een wettelijke verplichting tot schadevergoeding te matigen.\footnote{Hoge Raad 28 mei 1999, NJ 1999/510, ECLI:NL:HR:1999:ZC2913.} De omstandigheden die {[}gedaagde 4{]} , {[}gedaagde 3{]} , {[}gedaagde 5{]} CV en {[}gedaagde 6{]} (zonder verdere toelichting) naar voren hebben gebracht, zijn onvoldoende voor matiging.

4.6. Vorderingen vernietiging van hypotheekrecht en overdracht BMW niet toewijsbaar

4.6.1. De curator heeft gerechtelijke vernietiging van rechtshandelingen gevorderd waarbij {[}naam 6{]} (de zus van {[}gedaagde 2{]} ), {[}gedaagde 3{]} en {[}naam 7{]} (de dochter van {[}gedaagde 2{]} ) betrokken zijn. Artikel 3:51 BW bepaalt dat de rechtsvordering tot vernietiging wordt ingesteld tegen hen die partij bij de rechtshandeling zijn (lid 2). De curator heeft {[}naam 6{]} en {[}naam 7{]} niet gedagvaard, zodat de vorderingen tot vernietiging (met nevenvorderingen) van rechtshandelingen waar zij bij zijn betrokken, al om die reden niet toewijsbaar zijn. Voor {[}gedaagde 3{]} is dit anders, maar omdat de vordering tot vernietiging richting {[}gedaagde 3{]} uiteindelijk strekt tot teruggave van de BMW van {[}naam 7{]} , ontbreekt het belang bij toewijzing daarvan.

4.7. {[}gedaagde 4{]} , {[}gedaagde 3{]} , {[}gedaagde 5{]} CV en {[}gedaagde 6{]} moeten wel proceskosten betalen, met rente

4.7.1. De vorderingen van de curator op {[}gedaagde 4{]} , {[}gedaagde 3{]} , {[}gedaagde 5{]} CV en {[}gedaagde 6{]} worden grotendeels toegewezen. {[}gedaagde 4{]} , {[}gedaagde 3{]} , {[}gedaagde 5{]} CV en {[}gedaagde 6{]} worden als de in het ongelijk gestelde partijen daarom hoofdelijk veroordeeld in de proceskosten van de curator (inclusief nakosten). Deze worden begroot op:

{\def\LTcaptype{none} % do not increment counter
\begin{longtable}[]{@{}
  >{\raggedright\arraybackslash}p{(\linewidth - 6\tabcolsep) * \real{0.2500}}
  >{\raggedright\arraybackslash}p{(\linewidth - 6\tabcolsep) * \real{0.2500}}
  >{\raggedright\arraybackslash}p{(\linewidth - 6\tabcolsep) * \real{0.2500}}
  >{\raggedright\arraybackslash}p{(\linewidth - 6\tabcolsep) * \real{0.2500}}@{}}
\toprule\noalign{}
\begin{minipage}[b]{\linewidth}\raggedright
- kosten van de dagvaarding
\end{minipage} & \begin{minipage}[b]{\linewidth}\raggedright
€
\end{minipage} & \begin{minipage}[b]{\linewidth}\raggedright
123,26
\end{minipage} & \begin{minipage}[b]{\linewidth}\raggedright
\end{minipage} \\
\midrule\noalign{}
\endhead
\bottomrule\noalign{}
\endlastfoot
- griffierecht & € & 2.626,00 & \\
- salaris advocaat & € & 9.262,00 & (2 punten × € 4.631,00) \\
- nakosten & € & 189,00 & (plus de verhoging zoals vermeld in de beslissing) \\
Totaal & € & 12.077,00 & \\
\end{longtable}
}

4.8. De wettelijke rente over de proceskosten wordt toegewezen als hierna vermeld in de beslissing.

\paragraph{5 De beslissing}\label{de-beslissing-2}

De rechtbank

\emph{{[}gedaagde 4{]}}

5.1. veroordeelt {[}gedaagde 4{]} om aan de curator van {[}eiser 2{]} € 12.486,48 te betalen, te vermeerderen met de wettelijke rente als bedoeld in artikel 6:119 BW vanaf 11 augustus 2023 tot de dag van volledige betaling,

5.2. veroordeelt {[}gedaagde 4{]} om aan de curator van {[}eiser 2{]} € 899,86 aan buitengerechtelijke incassokosten te betalen,

5.3. veroordeelt {[}gedaagde 4{]} om aan de curator van {[}eiser 1{]} € 20.500,00 te betalen, te vermeerderen met de wettelijke rente als bedoeld in artikel 6:119 BW vanaf 11 augustus 2023 tot de dag van volledige betaling,

5.4. veroordeelt {[}gedaagde 4{]} om aan de curator van {[}eiser 1{]} € 980,00 aan buitengerechtelijke incassokosten te betalen,

\emph{{[}gedaagde 3{]}}

5.5. verklaart voor recht dat de curator bij brief van 28 maart 2024 rechtsgeldig (buitengerechtelijk) de vernietiging van de rechtshandelingen, waarbij de in 4.3.11 en 4.3.13 genoemde vermogensbestanddelen van {[}eiser 2{]} en {[}eiser 1{]} om niet zijn overgedragen aan {[}gedaagde 3{]} en die hebben geleid tot de oneigenlijke doorstart van de activiteiten van {[}eiser 2{]} en {[}eiser 1{]} door {[}gedaagde 3{]} , heeft ingeroepen,

5.6. veroordeelt {[}gedaagde 3{]} tot betaling van € 19.498,72 aan de curator van {[}eiser 2{]} , te vermeerderen met de wettelijke rente als bedoeld in artikel 6:119 BW vanaf 12 april 2024 tot de dag van volledige betaling,

5.7. veroordeelt {[}gedaagde 5{]} {[}gedaagde 3{]} om aan de curator van {[}eiser 2{]} € 969,99 aan buitengerechtelijke incassokosten te betalen,

5.8. veroordeelt {[}gedaagde 3{]} tot betaling van een vergoeding aan de curator van {[}eiser 1{]} en {[}eiser 2{]} als gevolg van de in 5.5 genoemde vernietiging voor zover die nog niet is vastgesteld in dit vonnis, op te maken bij staat en te vereffenen volgens de wet,

\emph{{[}gedaagde 5{]} CV en {[}gedaagde 6{]}}

5.9. veroordeelt {[}gedaagde 5{]} CV, {[}gedaagde 2{]} en {[}gedaagde 6{]} hoofdelijk om aan de curator van {[}eiser 1{]} € 448.985,00 te betalen, te vermeerderen met de wettelijke rente als bedoeld in artikel 6:119 BW vanaf 12 april 2024 tot de dag van volledige betaling,

5.10. veroordeelt {[}gedaagde 5{]} CV, {[}gedaagde 2{]} en {[}gedaagde 6{]} hoofdelijk om aan de curator van {[}eiser 1{]} € 4.019,93 aan buitengerechtelijke incassokosten te betalen,

5.11. veroordeelt {[}gedaagde 5{]} CV, {[}gedaagde 2{]} en {[}gedaagde 6{]} hoofdelijk tot teruggave van de Boston Whaler aan de curator van {[}eiser 1{]} ,

\emph{{[}gedaagde 1{]} en {[}gedaagde 2{]} als bestuurders van {[}eiser 2{]} en {[}eiser 1{]}}

5.12. verklaart voor recht dat {[}gedaagde 2{]} en {[}gedaagde 1{]} hoofdelijk aansprakelijk zijn voor het bedrag van de schulden van {[}eiser 1{]} , voor zover deze niet door vereffening van de overige baten kunnen worden voldaan (het boedeltekort), nader op te maken bij staat en te vereffenen volgens de wet,

5.13. verklaart voor recht dat {[}gedaagde 2{]} en {[}gedaagde 1{]} hoofdelijk aansprakelijk zijn voor het bedrag van de schulden van {[}eiser 2{]} , voor zover deze niet door vereffening van de overige baten kunnen worden voldaan (het boedeltekort), nader op te maken bij staat en te vereffenen volgens de wet,

5.14. veroordeelt {[}gedaagde 2{]} en {[}gedaagde 1{]} hoofdelijk om aan de curator van {[}eiser 1{]} € 800.000,00 te betalen als voorschot op het boedeltekort,

5.15. veroordeelt {[}gedaagde 2{]} en {[}gedaagde 1{]} hoofdelijk om aan de curator van {[}eiser 2{]}

€ 2.000.000,00 te betalen als voorschot op het boedeltekort,

\emph{{[}gedaagde 4{]} , {[}gedaagde 3{]} , {[}gedaagde 5{]} CV en {[}gedaagde 6{]}}

5.16. veroordeelt {[}gedaagde 4{]} , {[}gedaagde 3{]} , {[}gedaagde 5{]} CV en {[}gedaagde 6{]} hoofdelijk tot betaling van € 12.077,00 aan proceskosten, te betalen binnen veertien dagen na aanschrijving daartoe, te vermeerderen met € 98,00 plus de kosten van betekening als {[}gedaagde 4{]} , {[}gedaagde 3{]} , {[}gedaagde 5{]} CV en/of {[}gedaagde 6{]} niet tijdig aan de veroordelingen voldoen en het vonnis daarna wordt betekend,

5.17. veroordeelt {[}gedaagde 4{]} , {[}gedaagde 3{]} , {[}gedaagde 5{]} CV en {[}gedaagde 6{]} hoofdelijk tot betaling van de wettelijke rente als bedoeld in artikel 6:119 BW over de proceskosten als deze niet binnen veertien dagen na aanschrijving zijn betaald,

5.18. verklaart dit vonnis tot zover uitvoerbaar bij voorraad,

5.19. wijst het meer of anders gevorderde af.

Dit vonnis is gewezen door mrs. R.H.C. Jongeneel, P. Vrugt en G.C. de Heer, rechters, bijgestaan door mr. N. Noordmans, griffier, in het openbaar uitgesproken op 15 april 2026 en ondertekend door de voorzitter en de griffier.

\begin{center}\rule{0.5\linewidth}{0.5pt}\end{center}

\newpage

\markboth{ECLI:NL:RBAMS:2026:3787}{mistral/mistral-medium-3.5-128b:fp8 (k = 0)}

\subsection{k = 0}\label{k-0-2}

\subsubsection{Mistral}\label{mistral-6}

\texttt{mistral/mistral-medium-3.5-128b:fp8}

\begin{center}
\small
\begin{tabular}{@{}p{7.4cm}p{2.6cm}p{2.6cm}@{}}
\toprule
Variabele & Waarde (model) & Waarde (gold) \\
\midrule
Is rechtshandeling & \texttt{true} & \texttt{true} \\
Datum rechtshandeling & \texttt{2021-01-21} & \texttt{2021-01-21} \\
Verricht door de schuldenaar & \texttt{true} & \texttt{true} \\
Om niet & \texttt{true} & \texttt{true} \\
Datum faillietverklaring & \texttt{2021-07-13} & \texttt{2021-07-13} \\
Onverplicht & \texttt{true} & \texttt{true} \\
Benadeling schuldeisers & \texttt{true} & \texttt{true} \\
Wetenschap schuldenaar aangetoond & \texttt{false} & \texttt{false} \\
Wetenschap wederpartij aangetoond & \texttt{false} & \texttt{null} \\
Vermoeden wetenschap schuldenaar tegenbewijs & \texttt{null} & \texttt{false} \\
Vermoeden wetenschap wederpartij tegenbewijs & \texttt{null} & \texttt{false} \\
Tegen onderwaarde & \textcolor{gray}{\texttt{null}} & \texttt{true} \\
Voldoening of zekerheid niet opeisbare schuld & \texttt{null} & \texttt{null} \\
Gelieerde partij & \texttt{true} & \texttt{true} \\
Reeds verplicht voor aanvang termijn & \texttt{null} & \texttt{false} \\
\textbf{Vernietigbaar} & \textbf{\texttt{true}} & \textbf{\texttt{true}} \\
\bottomrule
\end{tabular}
\end{center}

\paragraph{Is rechtshandeling - overwegingen}\label{is-rechtshandeling---overwegingen-18}

De overdracht van domeinnamen (2.6.4) is een goederenrechtelijke handeling en kwalificeert als rechtshandeling in de zin van art. 3:33 BW.

\paragraph{Verricht door de schuldenaar - overwegingen}\label{verricht-door-de-schuldenaar---overwegingen-18}

De overdracht van domeinnamen is verricht door {[}eiser 2{]} (2.6.4), die later failliet is verklaard. De overdracht van de huurovereenkomsten is verricht door {[}eiser 1{]} en {[}eiser 2{]} (2.6.4).

\paragraph{Om niet - overwegingen}\label{om-niet---overwegingen-18}

De curator stelt dat {[}gedaagde 3{]} geen tegenprestatie heeft verricht voor de verkrijging van vermogensbestanddelen (4.3.1). De overdracht van domeinnamen en huurovereenkomsten (2.6.4) is derhalve om niet.

\paragraph{Onverplicht - overwegingen}\label{onverplicht---overwegingen-18}

De overgedragen vermogensbestanddelen (domeinnamen, huurovereenkomsten en softe activa) van {[}eiser 1{]} en {[}eiser 2{]} aan {[}gedaagde 3{]} zijn zonder tegenprestatie overgedragen. Er bestond geen op de wet of een overeenkomst berustende rechtsplicht voor {[}eiser 1{]} en {[}eiser 2{]} om deze vermogensbestanddelen aan {[}gedaagde 3{]} over te dragen. De overdrachten waren derhalve onverplicht in de zin van art. 42 lid 1 Fw. De curator heeft voldoende aangeduid dat de overdrachten deel uitmaakten van een oneigenlijke doorstart, waarbij {[}gedaagde 3{]} zonder tegenprestatie vermogensbestanddelen heeft verkregen. Dit impliceert dat er geen afdwingbare verplichting bestond om deze overdrachten te verrichten, wat het onverplichte karakter bevestigt. De stelplicht en bewijslast van het onverplichte karakter rusten op de curator, en deze heeft voldoende feiten en omstandigheden aangedragen om aan te tonen dat de overdrachten onverplicht waren (bijv. gebrek aan leningsovereenkomsten, geen tegenprestatie, en de context van de oneigenlijke doorstart).

\paragraph{Benadeling schuldeisers - overwegingen}\label{benadeling-schuldeisers---overwegingen-18}

Uit de casus blijkt dat {[}eiser 1{]} en {[}eiser 2{]} in een slechte financiële situatie verkeerden, met aanzienlijke verliezen en een dreigend faillissement (zoals gemeld door de boekhouder {[}naam 4{]} in 2020). De overgedragen vermogensbestanddelen (domeinnamen, huurovereenkomsten en softe activa) hadden waarde voor de boedel van {[}eiser 1{]} en {[}eiser 2{]}, maar zijn zonder adequate tegenprestatie overgedragen aan {[}gedaagde 3{]}. Hierdoor is het actief van de failliete boedels verminderd, wat direct benadelend is voor de schuldeisers. De curator heeft een boedeltekort vastgesteld (€ 886.548,01 voor {[}eiser 1{]} en € 2.208.533,29 voor {[}eiser 2{]}), wat aantoont dat de schuldeisers niet volledig kunnen worden voldaan. Zonder de overgedragen activa zou de boedel groter zijn geweest, wat de schuldeisers ten goede zou zijn gekomen. De benadeling is dus aanwezig in de zin van artikel 42 Fw, omdat de rechtshandelingen (overdracht van activa) hebben geleid tot een vermindering van het actief ten nadele van de schuldeisers. Bovendien wisten of behoorden {[}gedaagde 1{]}, {[}gedaagde 2{]} en {[}gedaagde 3{]} te weten dat deze handelingen benadelend waren, gelet op de financiële toestand van {[}eiser 1{]} en {[}eiser 2{]} en de afwezigheid van een reële tegenprestatie. De leer van de middellijke benadeling is hier van toepassing, omdat de benadeling op het moment van de faillietverklaring (13 juli 2021 voor {[}eiser 1{]}) al aanwezig was en door de curator wordt ingeroepen. De overdrachten vonden plaats in januari en februari 2021, kort voor het faillissement, en droegen bij aan de onmogelijkheid om de schulden te voldoen. De benadeling is dus aantoonbaar en voldoet aan de voorwaarden van artikel 42 Fw, lid 1 en 2 (indien van toepassing op meerzijdige handelingen).

\paragraph{Wetenschap - overwegingen}\label{wetenschap---overwegingen-18}

De overdrachten van domeinnamen en huurovereenkomsten door {[}eiser 1{]} en {[}eiser 2{]} aan {[}gedaagde 3{]} op 21 januari 2021 en 10 februari 2021 vonden plaats binnen een jaar voor de faillietverklaring van {[}eiser 1{]} (13 juli 2021). {[}gedaagde 3{]} is opgericht door {[}gedaagde 2{]}, die bestuurder was van {[}eiser 1{]} en {[}eiser 2{]}. Hierdoor is {[}gedaagde 3{]} een gelieerde partij in de zin van art. 43 lid 1 sub 4° Fw (rechtshandeling door rechtspersoon-schuldenaar met een natuurlijk persoon die bestuurder is van de schuldenaar). Het bewijsvermoeden van art. 43 lid 1 Fw is dus van toepassing, tenzij {[}eiser 1{]} zich reeds voor 21 januari 2020 (een jaar voor de eerste overdracht) tot deze handelingen had verplicht, wat niet het geval lijkt. De wetenschap van benadeling wordt derhalve aan beide zijden (schuldenaar en wederpartij) vermoed. Er is geen aanwijzing dat dit vermoeden met succes is weerlegd (geen tegenbewijs aangedragen).

\newpage

\markboth{ECLI:NL:RBAMS:2026:3787}{gpt-5.5 (k = 0)}

\subsubsection{GPT}\label{gpt-6}

\texttt{gpt-5.5}

\begin{center}
\small
\begin{tabular}{@{}p{7.4cm}p{2.6cm}p{2.6cm}@{}}
\toprule
Variabele & Waarde (model) & Waarde (gold) \\
\midrule
Is rechtshandeling & \texttt{true} & \texttt{true} \\
Datum rechtshandeling & \texttt{2021-01-21} & \texttt{2021-01-21} \\
Verricht door de schuldenaar & \texttt{true} & \texttt{true} \\
Om niet & \texttt{true} & \texttt{true} \\
Datum faillietverklaring & \texttt{2021-07-13} & \texttt{2021-07-13} \\
Onverplicht & \texttt{true} & \texttt{true} \\
Benadeling schuldeisers & \texttt{true} & \texttt{true} \\
Wetenschap schuldenaar aangetoond & \textcolor{gray}{\texttt{true}} & \texttt{false} \\
Wetenschap wederpartij aangetoond & \textcolor{gray}{\texttt{true}} & \texttt{null} \\
Vermoeden wetenschap schuldenaar tegenbewijs & \texttt{false} & \texttt{false} \\
Vermoeden wetenschap wederpartij tegenbewijs & \texttt{false} & \texttt{false} \\
Tegen onderwaarde & \texttt{true} & \texttt{true} \\
Voldoening of zekerheid niet opeisbare schuld & \texttt{null} & \texttt{null} \\
Gelieerde partij & \texttt{true} & \texttt{true} \\
Reeds verplicht voor aanvang termijn & \texttt{false} & \texttt{false} \\
\textbf{Vernietigbaar} & \textbf{\texttt{true}} & \textbf{\texttt{true}} \\
\bottomrule
\end{tabular}
\end{center}

\paragraph{Is rechtshandeling - overwegingen}\label{is-rechtshandeling---overwegingen-19}

Het gaat om een samenstel van overdrachten van vermogensbestanddelen aan {[}gedaagde 3{]}, waaronder de overdracht van domeinnamen en de overdracht van rechten uit huurovereenkomsten. Deze overdrachten zijn gericht op rechtsgevolg en kwalificeren daarom als rechtshandelingen.

\paragraph{Verricht door de schuldenaar - overwegingen}\label{verricht-door-de-schuldenaar---overwegingen-19}

Uit de feiten volgt dat {[}eiser 2{]} haar domeinnamen aan {[}gedaagde 3{]} heeft overgedragen en dat {[}eiser 1{]} en {[}eiser 2{]} de rechten uit huurovereenkomsten aan {[}gedaagde 3{]} hebben overgedragen. De gewraakte rechtshandelingen zijn dus verricht door de schuldenaren zelf.

\paragraph{Om niet - overwegingen}\label{om-niet---overwegingen-19}

De curator stelt dat {[}gedaagde 3{]} voor de verkrijging van de vermogensbestanddelen van {[}eiser 1{]} en {[}eiser 2{]} geen tegenprestatie heeft verricht. Uit de aangeleverde feiten blijkt geen reële wederprestatie voor deze overdrachten.

\paragraph{Onverplicht - overwegingen}\label{onverplicht---overwegingen-19}

Uit de casus blijkt niet van een op wet of overeenkomst berustende, in rechte afdwingbare verplichting van {[}eiser 1{]} of {[}eiser 2{]} om de betrokken vermogensbestanddelen aan {[}gedaagde 3{]} over te dragen. Het gaat om overdracht van onder meer domeinnamen, rechten uit huurovereenkomsten en feitelijke voortzetting/overgang van bedrijfsactiviteiten naar een kort daarvoor opgerichte gelieerde vennootschap, zonder dat een tegenprestatie is gesteld. In de administratie zijn bovendien geen relevante overeenkomsten aangetroffen die een verplichting tot deze overdrachten onderbouwen. Voor zover sprake was van financiële druk, continuïteitsproblemen of de wens om de onderneming voort te zetten, levert dat geen rechtsplicht op in de zin van artikel 42 Fw. Ook het enkele feit dat personeel is overgegaan of dat {[}gedaagde 3{]} salarissen heeft betaald en via rekening-courant heeft doorbelast, maakt niet dat voor het samenstel van overdrachten een afdwingbare verplichting bestond.

\paragraph{Benadeling schuldeisers - overwegingen}\label{benadeling-schuldeisers---overwegingen-19}

De aangevochten handelingen moeten als samenstel worden beoordeeld: de overdracht van domeinnamen, huur-/gebruiksrechten en kennelijk ook de overgang van personeel en overige `softe' activa/ondernemingswaarde naar de kort daarvoor opgerichte {[}gedaagde 3{]}, in het kader van een oneigenlijke doorstart. Uit de casus volgt dat {[}gedaagde 3{]} daarvoor volgens de curator geen tegenprestatie heeft verricht. In de feitelijke situatie zijn vermogensbestanddelen en ondernemingswaarde uit {[}eiser 1{]} en {[}eiser 2{]} verdwenen zonder dat daar een corresponderende opbrengst voor de boedel tegenover stond. Dat betekent een vermindering van het actief. Daarnaast is vermeld dat {[}gedaagde 3{]} in augustus 2021 salarissen heeft betaald en dit bedrag via rekening-courant ten laste van {[}eiser 2{]} heeft gebracht, wat eerder wijst op een verdere verslechtering van de vermogenspositie van {[}eiser 2{]} dan op compensatie voor de overgedragen activa. In het hypothetische scenario zonder deze handelingen zouden de domeinnamen, contractuele rechten en overige activa/goingconcernwaarde bij {[}eiser 1{]} en {[}eiser 2{]} zijn gebleven en in beginsel voor verhaal door de gezamenlijke schuldeisers beschikbaar zijn geweest, dan wel door de curator te gelde kunnen zijn gemaakt. Ten tijde van het inroepen van de pauliana bestaan bovendien aanzienlijke boedeltekorten in beide faillissementen. Daardoor leidt het wegvallen van activa zonder reële tegenprestatie tot minder verhaal voor de schuldeisers in het algemeen. Er zijn geen feiten gesteld waaruit volgt dat de schuldeisers zich ook zonder de handelingen niet op deze vermogensbestanddelen hadden kunnen verhalen, bijvoorbeeld wegens uitwinning door een separatist, of dat alle schulden ook zonder die activa volledig voldaan hadden kunnen worden. De schuldeisers zijn daarom benadeeld in de zin van artikel 42 Fw.

\paragraph{Wetenschap - overwegingen}\label{wetenschap---overwegingen-19}

Ten tijde van de overdracht van domeinnamen, huurposities/ondernemingsactiviteiten en personeel aan de kort daarvoor door {[}gedaagde 2{]} opgerichte {[}gedaagde 3{]} was het faillissement met een tekort met een redelijke mate van waarschijnlijkheid voorzienbaar. De bestuurders van {[}eiser 1{]} en {[}eiser 2{]} waren al in 2020 door de boekhouder gewezen op ernstige continuïteitsproblemen, substantiële verliezen, belastingschulden en zelfs dat faillissement van {[}eiser 2{]} niet ondenkbaar was; daarnaast blijkt uit de aanpassingen van de administratie dat zij met de slechte financiële positie bekend waren. Die wetenschap wordt aan {[}eiser 1{]} en {[}eiser 2{]} toegerekend. Ook aan de zijde van {[}gedaagde 3{]} is wetenschap rechtstreeks aannemelijk, omdat {[}gedaagde 3{]} kort voor de doorstart door {[}gedaagde 2{]} is opgericht en de overdrachten plaatsvonden binnen dezelfde kring van betrokkenen die van de financiële problemen op de hoogte waren. Voor zover de handelingen niet als om niet maar als anders dan om niet worden gekwalificeerd, wijzen het ontbreken van een reële tegenprestatie en de gelieerdheid met {[}gedaagde 3{]} bovendien op toepasselijkheid van de bewijsvermoedens van art. 43 Fw; van een reeds vóór de jaartermijn bestaande verplichting blijkt niet. Voor zover sprake is van rechtshandelingen om niet, geldt het vermoeden van art. 45 Fw aan de zijde van de schuldenaren. Er zijn geen feiten gesteld die dit vermoeden ontzenuwen; integendeel bevestigen de waarschuwingen en financiële gegevens de voorzienbaarheid van faillissement en tekort.

\newpage

\markboth{ECLI:NL:RBAMS:2026:3787}{gemini-3.5-flash (k = 0)}

\subsubsection{Gemini}\label{gemini-6}

\texttt{gemini-3.5-flash}

\begin{center}
\small
\begin{tabular}{@{}p{7.4cm}p{2.6cm}p{2.6cm}@{}}
\toprule
Variabele & Waarde (model) & Waarde (gold) \\
\midrule
Is rechtshandeling & \texttt{true} & \texttt{true} \\
Datum rechtshandeling & \texttt{2021-01-21} & \texttt{2021-01-21} \\
Verricht door de schuldenaar & \texttt{true} & \texttt{true} \\
Om niet & \texttt{true} & \texttt{true} \\
Datum faillietverklaring & \texttt{2021-07-13} & \texttt{2021-07-13} \\
Onverplicht & \texttt{true} & \texttt{true} \\
Benadeling schuldeisers & \texttt{true} & \texttt{true} \\
Wetenschap schuldenaar aangetoond & \textcolor{gray}{\texttt{true}} & \texttt{false} \\
Wetenschap wederpartij aangetoond & \textcolor{gray}{\texttt{true}} & \texttt{null} \\
Vermoeden wetenschap schuldenaar tegenbewijs & \texttt{false} & \texttt{false} \\
Vermoeden wetenschap wederpartij tegenbewijs & \texttt{false} & \texttt{false} \\
Tegen onderwaarde & \textcolor{gray}{\texttt{null}} & \texttt{true} \\
Voldoening of zekerheid niet opeisbare schuld & \texttt{null} & \texttt{null} \\
Gelieerde partij & \texttt{true} & \texttt{true} \\
Reeds verplicht voor aanvang termijn & \texttt{false} & \texttt{false} \\
\textbf{Vernietigbaar} & \textbf{\texttt{true}} & \textbf{\texttt{true}} \\
\bottomrule
\end{tabular}
\end{center}

\paragraph{Is rechtshandeling - overwegingen}\label{is-rechtshandeling---overwegingen-20}

Het betreft een samenstel van handelingen, waaronder de overdracht van domeinnamen en huurrechten, die gericht zijn op een rechtsgevolg (de overgang van deze activa en rechten) en dus kwalificeren als rechtshandelingen.

\paragraph{Verricht door de schuldenaar - overwegingen}\label{verricht-door-de-schuldenaar---overwegingen-20}

De overdrachten van de domeinnamen en de huurrechten zijn rechtstreeks verricht door de schuldenaren zelf, te weten {[}eiser 1{]} en {[}eiser 2{]}.

\paragraph{Om niet - overwegingen}\label{om-niet---overwegingen-20}

De curator stelt dat de overdracht van de vermogensbestanddelen aan {[}gedaagde 3{]} heeft plaatsgevonden zonder dat daar een tegenprestatie tegenover stond, wat betekent dat de rechtshandelingen om niet zijn verricht.

\paragraph{Onverplicht - overwegingen}\label{onverplicht---overwegingen-20}

Het samenstel van rechtshandelingen waarbij {[}eiser 1{]} en {[}eiser 2{]} hun vermogensbestanddelen (zoals domeinnamen, huurrechten en personeel) hebben overgedragen aan {[}gedaagde 3{]} ten behoeve van een doorstart, is onverplicht verricht. Er is geen enkele wettelijke bepaling of voorafgaande contractuele overeenkomst gesteld of gebleken waaruit een in rechte afdwingbare verplichting voortvloeide voor {[}eiser 1{]} en {[}eiser 2{]} om deze activa aan {[}gedaagde 3{]} over te dragen.

\paragraph{Benadeling schuldeisers - overwegingen}\label{benadeling-schuldeisers---overwegingen-20}

Door het samenstel van rechtshandelingen zijn waardevolle activa van {[}eiser 1{]} en {[}eiser 2{]} (zoals domeinnamen, huurrechten en de operationele onderneming via de overgang van het personeel) overgedragen aan {[}gedaagde 3{]} ten behoeve van een oneigenlijke doorstart, zonder dat daar een reële tegenprestatie tegenover stond. Dit heeft geleid tot een rechtstreekse vermindering van het actieve vermogen van de schuldenaren. In de hypothetische situatie waarin deze overdrachten niet hadden plaatsgevonden, hadden deze activa (of de waarde daarvan) ten gunste van de boedel kunnen worden aangewend. Nu er sprake is van omvangrijke boedeltekorten in beide faillissementen, zijn de schuldeisers door deze onttrekking concreet in hun verhaalsmogelijkheden benadeeld.

\paragraph{Wetenschap - overwegingen}\label{wetenschap---overwegingen-20}

De curator vordert de vernietiging van het samenstel van rechtshandelingen die hebben geleid tot de oneigenlijke doorstart van {[}eiser 1{]} en {[}eiser 2{]} naar {[}gedaagde 3{]}. De rechtshandelingen (overdracht domeinnamen, huurrechten en personeel) hebben binnen een jaar voor de faillietverklaringen plaatsgevonden. Er is sprake van een transactie met een gelieerde partij in de zin van art. 43 lid 1 Fw, aangezien {[}gedaagde 3{]} is opgericht en wordt beheerst door {[}gedaagde 2{]}, die tevens bestuurder was van de failliete vennootschappen. Hierdoor geldt het wettelijk bewijsvermoeden van wetenschap van benadeling aan beide zijden. Dit vermoeden is niet ontzenuwd. Bovendien is de wetenschap van benadeling ook rechtstreeks aangetoond: de bestuurders waren volledig op de hoogte van de penibele financiële situatie, de waarschuwingen van de boekhouder over een dreigend faillissement en de gemanipuleerde jaarcijfers. Deze wetenschap aan de zijde van de schuldenaren kan volledig aan de wederpartij {[}gedaagde 3{]} worden toegerekend.

\newpage

\markboth{ECLI:NL:RBAMS:2026:3787}{mistral/mistral-medium-3.5-128b:fp8 (k = 5)}

\subsection{k = 5}\label{k-5-2}

\subsubsection{Mistral}\label{mistral-7}

\texttt{mistral/mistral-medium-3.5-128b:fp8}

\begin{center}
\small
\begin{tabular}{@{}p{7.4cm}p{2.6cm}p{2.6cm}@{}}
\toprule
Variabele & Waarde (model) & Waarde (gold) \\
\midrule
Is rechtshandeling & \texttt{true} & \texttt{true} \\
Datum rechtshandeling & \texttt{2021-01-21} & \texttt{2021-01-21} \\
Verricht door de schuldenaar & \texttt{true} & \texttt{true} \\
Om niet & \texttt{true} & \texttt{true} \\
Datum faillietverklaring & \texttt{2021-07-13} & \texttt{2021-07-13} \\
Onverplicht & \texttt{true} & \texttt{true} \\
Benadeling schuldeisers & \texttt{true} & \texttt{true} \\
Wetenschap schuldenaar aangetoond & \textcolor{gray}{\texttt{true}} & \texttt{false} \\
Wetenschap wederpartij aangetoond & \textcolor{gray}{\texttt{true}} & \texttt{null} \\
Vermoeden wetenschap schuldenaar tegenbewijs & \texttt{null} & \texttt{false} \\
Vermoeden wetenschap wederpartij tegenbewijs & \texttt{null} & \texttt{false} \\
Tegen onderwaarde & \textcolor{gray}{\texttt{null}} & \texttt{true} \\
Voldoening of zekerheid niet opeisbare schuld & \texttt{null} & \texttt{null} \\
Gelieerde partij & \texttt{true} & \texttt{true} \\
Reeds verplicht voor aanvang termijn & \texttt{null} & \texttt{false} \\
\textbf{Vernietigbaar} & \textbf{\texttt{true}} & \textbf{\texttt{true}} \\
\bottomrule
\end{tabular}
\end{center}

\paragraph{Is rechtshandeling - overwegingen}\label{is-rechtshandeling---overwegingen-21}

De overdracht van domeinnamen (2.6.4) is een goederenrechtelijke handeling en kwalificeert als rechtshandeling. Ook de overdracht van de rechten uit huurovereenkomsten (2.6.4) is een rechtshandeling.

\paragraph{Verricht door de schuldenaar - overwegingen}\label{verricht-door-de-schuldenaar---overwegingen-21}

De overdrachten zijn verricht door {[}eiser 2{]} (domeinnamen) en {[}eiser 1{]} en {[}eiser 2{]} (huurovereenkomsten), die de schuldenaren zijn in deze context.

\paragraph{Om niet - overwegingen}\label{om-niet---overwegingen-21}

De curator stelt dat {[}gedaagde 3{]} geen tegenprestatie heeft verricht voor de verkrijging van vermogensbestanddelen (4.3.1). De overdrachten van domeinnamen en huurovereenkomsten zijn derhalve om niet.

\paragraph{Onverplicht - overwegingen}\label{onverplicht---overwegingen-21}

De overdracht van de domeinnamen en de rechten uit de huurovereenkomsten door {[}eiser 1{]} en {[}eiser 2{]} aan {[}gedaagde 3{]} op respectievelijk 21 januari 2021 en 10 februari 2021 zijn onverplichte rechtshandelingen in de zin van artikel 42 lid 1 Fw.

Uit de casus blijkt dat er geen op de wet of een overeenkomst berustende rechtsplicht bestond voor {[}eiser 1{]} en {[}eiser 2{]} om deze vermogensbestanddelen aan {[}gedaagde 3{]} over te dragen. De overdrachten vonden plaats in het kader van een oneigenlijke doorstart, waarbij {[}gedaagde 3{]} de ondernemingsactiviteiten van {[}eiser 1{]} en {[}eiser 2{]} voortzette zonder dat daarvoor een tegenprestatie werd verricht.

De boekhouder {[}naam 4{]} heeft herhaaldelijk gewaarschuwd voor de slechte financiële situatie van {[}eiser 1{]} en {[}eiser 2{]} (o.a. op 13 januari 2020, 30 april 2020 en 25 september 2020), en de continuïteit van beide vennootschappen was niet meer gegarandeerd. Dit betekent dat de overdrachten plaatsvonden in een context waarin de schuldenaren wisten of behoorden te weten dat benadeling van schuldeisers het gevolg zou zijn.

Daarnaast was {[}gedaagde 3{]} op 22 december 2020 opgericht door {[}gedaagde 2{]}, die ook bestuurder was van {[}eiser 1{]} en {[}eiser 2{]}. Dit impliceert dat {[}gedaagde 3{]} op de hoogte was of behoorde te zijn van de financiële toestand van {[}eiser 1{]} en {[}eiser 2{]} en de gevolgen van de overdrachten voor de schuldeisers.

Kortom, de overdrachten waren niet verplicht en vonden plaats in een situatie waarin benadeling van schuldeisers voorzienbaar was, wat voldoet aan de voorwaarden van artikel 42 lid 1 Fw voor onverplichtheid. Zie ook ECLI:NL:RBZWB:2025:5504 (r.o. 4.4-4.5) en ECLI:NL:RBGEL:2015:6805 (r.o. 4.15-4.16).

\paragraph{Benadeling schuldeisers - overwegingen}\label{benadeling-schuldeisers---overwegingen-21}

Op basis van de casus en de toepasselijke jurisprudentie (met name ECLI:NL:RBZWB:2025:5504 en ECLI:NL:RBNNE:2021:3917) kan worden geconcludeerd dat de schuldeisers door de aangevochten rechtshandelingen zijn benadeeld in de zin van artikel 42 Fw.

Allereerst betreft het een \textbf{samenstel van rechtshandelingen} (overdracht domeinnamen, huurovereenkomsten, en overname personeel) die gezamenlijk de oneigenlijke doorstart van {[}eiser 1{]} en {[}eiser 2{]} via {[}gedaagde 3{]} mogelijk maakten. Deze handelingen zijn \textbf{onverplicht} verricht: er bestond geen contractuele of wettelijke verplichting voor {[}eiser 1{]} en {[}eiser 2{]} om hun activa (zoals domeinnamen, huurcontracten en personeel) aan {[}gedaagde 3{]} over te dragen.

Vervolgens is er sprake van \textbf{benadeling van de schuldeisers}. De vermogensvergelijking toont aan dat zonder deze handelingen de activa (domeinnamen, huurcontracten, en de waarde van het personeel/goodwill) in de boedel van {[}eiser 1{]} en {[}eiser 2{]} zouden zijn gebleven, waar de schuldeisers zich op hadden kunnen verhalen. Door de overdracht aan {[}gedaagde 3{]} zijn deze activa voor de boedel en dus voor de schuldeisers \textbf{verloren gegaan}. Dit geldt met name omdat:
- De domeinnamen en huurcontracten waardevolle activa waren die in de boedel hadden kunnen worden geliquideerd.
- De overname van personeel en activiteiten door {[}gedaagde 3{]} zonder tegenprestatie aan de boedel betekent dat de ondernemingswaarde (goodwill) aan de schuldeisers is onttrokken.
- De \textbf{verstoring van de paritas creditorum} is evident: {[}gedaagde 3{]} heeft een voorsprong verkregen ten opzichte van andere schuldeisers door de activa te verwerven zonder dat de opbrengst ten goede kwam aan de gezamenlijke schuldeisers. Dit is in lijn met de jurisprudentie (bijv. ECLI:NL:RBZWB:2025:5504, r.o. 4.11), waar benadeling wordt aangenomen wanneer activa aan één partij ten goede komen ten koste van de gezamenlijke schuldeisers, ook als per saldo geen vermogensvermindering optreedt.

Daarnaast was de \textbf{benadeling voorzienbaar}. Uit de casus blijkt dat {[}eiser 1{]} en {[}eiser 2{]} al geruime tijd in financiële moeilijkheden verkeerden (verliezen in 2018 en 2019, waarschuwingen van de boekhouder over continuïteitsrisico's, en een boedeltekort van respectievelijk € 886.548,01 en € 2.208.533,29). De boekhouder had {[}gedaagde 1{]} en {[}gedaagde 2{]} (bestuurders van zowel {[}eiser 1{]}, {[}eiser 2{]} als {[}gedaagde 3{]}) \textbf{herhaaldelijk gewaarschuwd} voor de slechte financiële situatie en het risico op faillissement (bijv. op 13 januari 2020, 30 april 2020, en 25 september 2020). Gelet op artikel 43 lid 1 onder 5° Fw (gelijke bestuurders bij {[}eiser 1{]}/{[}eiser 2{]} en {[}gedaagde 3{]}) en onder 6° Fw (behoren tot dezelfde groep), wordt \textbf{vermoed} dat zowel {[}eiser 1{]}/{[}eiser 2{]} als {[}gedaagde 3{]} wisten of behoorden te weten dat de schuldeisers door deze handelingen zouden worden benadeeld. Dit vermoeden is niet ontzenuwd: {[}gedaagde 3{]} is opgericht door {[}gedaagde 2{]} (die ook bestuurder was van {[}eiser 1{]} en {[}eiser 2{]}) en de overdrachten vonden plaats in een periode waarin het faillissement \textbf{redelijkerwijs voorzienbaar} was (ECLI:NL:RBZWB:2025:5504, r.o. 4.17-4.20).

Tot slot is de benadeling \textbf{aanwezig op het moment dat de curator zijn rechten doet gelden} (leer van de middellijke benadeling). De faillissementen van {[}eiser 1{]} (13 juli 2021) en {[}eiser 2{]} (28 september 2021) tonen aan dat de schuldeisers daadwerkelijk in hun verhaalsmogelijkheden zijn beperkt door het verlies van de activa.

Concluderend zijn alle vereisten voor benadeling in de zin van artikel 42 Fw vervuld: onverplichte rechtshandelingen, benadeling van schuldeisers (verlies van activa en verstoring paritas), en wetenschap/behoren te weten van benadeling bij zowel {[}eiser 1{]}/{[}eiser 2{]} als {[}gedaagde 3{]}.

\paragraph{Wetenschap - overwegingen}\label{wetenschap---overwegingen-21}

De overdracht van domeinnamen (21 januari 2021) en huurovereenkomsten (10 februari 2021) door {[}eiser 1{]} en {[}eiser 2{]} aan {[}gedaagde 3{]} valt binnen de termijn van één jaar voor de faillietverklaring van {[}eiser 1{]} (13 juli 2021). {[}gedaagde 3{]} is opgericht door {[}gedaagde 2{]}, die tevens bestuurder is van {[}eiser 1{]} en {[}eiser 2{]}. Dit betekent dat {[}eiser 1{]} (schuldenaar) en {[}gedaagde 3{]} (wederpartij) verbonden zijn via een gelieerde partij in de zin van art. 43 lid 1 sub 4° Fw (bestuurder van schuldenaar is betrokken bij de wederpartij). Hierdoor wordt de wetenschap van benadeling aan beide zijden vermoed (art. 43 lid 1 Fw). Er is geen sprake van een verplichting die voor de jaartermijn bestond, noch is er tegenbewijs geleverd dat het faillissement en tekort niet voorzienbaar waren. De boekhouder {[}naam 4{]} had al in januari 2020 en april 2020 gewaarschuwd voor continuïteitsrisico's, en in september 2020 werd het faillissement van {[}eiser 2{]} als niet ondenkbaar bestempeld. De wetenschap van benadeling is derhalve aannemelijk.

\newpage

\markboth{ECLI:NL:RBAMS:2026:3787}{gpt-5.5 (k = 5)}

\subsubsection{GPT}\label{gpt-7}

\texttt{gpt-5.5}

\begin{center}
\small
\begin{tabular}{@{}p{7.4cm}p{2.6cm}p{2.6cm}@{}}
\toprule
Variabele & Waarde (model) & Waarde (gold) \\
\midrule
Is rechtshandeling & \texttt{true} & \texttt{true} \\
Datum rechtshandeling & \texttt{2021-01-21} & \texttt{2021-01-21} \\
Verricht door de schuldenaar & \texttt{true} & \texttt{true} \\
Om niet & \texttt{true} & \texttt{true} \\
Datum faillietverklaring & \texttt{2021-07-13} & \texttt{2021-07-13} \\
Onverplicht & \texttt{true} & \texttt{true} \\
Benadeling schuldeisers & \texttt{true} & \texttt{true} \\
Wetenschap schuldenaar aangetoond & \textcolor{gray}{\texttt{true}} & \texttt{false} \\
Wetenschap wederpartij aangetoond & \textcolor{gray}{\texttt{true}} & \texttt{null} \\
Vermoeden wetenschap schuldenaar tegenbewijs & \texttt{false} & \texttt{false} \\
Vermoeden wetenschap wederpartij tegenbewijs & \texttt{false} & \texttt{false} \\
Tegen onderwaarde & \texttt{true} & \texttt{true} \\
Voldoening of zekerheid niet opeisbare schuld & \texttt{null} & \texttt{null} \\
Gelieerde partij & \texttt{true} & \texttt{true} \\
Reeds verplicht voor aanvang termijn & \texttt{false} & \texttt{false} \\
\textbf{Vernietigbaar} & \textbf{\texttt{true}} & \textbf{\texttt{true}} \\
\bottomrule
\end{tabular}
\end{center}

\paragraph{Is rechtshandeling - overwegingen}\label{is-rechtshandeling---overwegingen-22}

Het gewraakte samenstel betreft overdrachten van vermogensbestanddelen aan {[}gedaagde 3{]}, waaronder de overdracht van domeinnamen en rechten uit huurovereenkomsten. Deze handelingen zijn gericht op rechtsgevolg en kwalificeren op basis van de feiten als rechtshandelingen.

\paragraph{Verricht door de schuldenaar - overwegingen}\label{verricht-door-de-schuldenaar---overwegingen-22}

Uit de feiten volgt dat {[}eiser 2{]} haar domeinnamen aan {[}gedaagde 3{]} heeft overgedragen en dat {[}eiser 1{]} en {[}eiser 2{]} rechten uit huurovereenkomsten aan {[}gedaagde 3{]} hebben overgedragen. Het samenstel is dus door de schuldenaren zelf verricht.

\paragraph{Om niet - overwegingen}\label{om-niet---overwegingen-22}

De curator stelt dat {[}gedaagde 3{]} de vermogensbestanddelen van {[}eiser 1{]} en {[}eiser 2{]} heeft verkregen zonder daarvoor een tegenprestatie te verrichten. Op basis van deze feiten is de rechtshandeling om niet.

\paragraph{Onverplicht - overwegingen}\label{onverplicht---overwegingen-22}

Voor de beoordeelde rechtshandelingen --- de overdracht van domeinnamen, de overdracht van rechten uit huurovereenkomsten, de overgang/overheveling van personeel en meer in het algemeen het samenstel waarmee activa en bedrijfsactiviteiten van {[}eiser 1{]} en {[}eiser 2{]} bij {[}gedaagde 3{]} terechtkwamen --- blijkt uit de casus niet van een vooraf bestaande, in rechte afdwingbare verplichting uit wet of overeenkomst. Er zijn geen overeenkomsten of andere rechtsbronnen genoemd waaruit volgt dat {[}eiser 1{]} of {[}eiser 2{]} gehouden waren deze vermogensbestanddelen of bedrijfsactiviteiten aan de kort daarvoor opgerichte {[}gedaagde 3{]} over te dragen, laat staan zonder tegenprestatie. Dat een doorstart feitelijk, economisch of organisatorisch wenselijk of noodzakelijk werd geacht, schept geen rechtsplicht in de zin van artikel 42 Fw. Dit sluit aan bij de lijn in onder meer ECLI:NL:RBNNE:2021:3917 en ECLI:NL:RBGEL:2015:6805, waarin overdracht/overheveling van ondernemingsactiva in het kader van een oneigenlijke doorstart zonder afdwingbare verplichting als onverplicht werd aangemerkt. Ook de omstandigheid dat binnen de betrokken vennootschappen veel rekening-courant- en `lening'-boekingen voorkwamen, zonder aangetroffen leningsovereenkomsten, levert geen rechtsplicht tot deze overdrachten op.

\paragraph{Benadeling schuldeisers - overwegingen}\label{benadeling-schuldeisers---overwegingen-22}

Ja. Beoordeeld moet worden of de schuldeisers op het moment dat de curator zijn rechten doet gelden slechter af zijn dan in de hypothetische situatie zonder het samenstel van rechtshandelingen. Uit de casus volgt dat kort vóór het faillissement van {[}eiser 1{]} op 13 juli 2021 en later {[}eiser 2{]} vermogensbestanddelen en ondernemingsonderdelen aan de net opgerichte {[}gedaagde 3{]} zijn overgedragen: domeinnamen, rechten uit huurovereenkomsten, personeel en volgens de curator ook overige `softe activa'/goodwill die de doorstart mogelijk maakten. Daarvoor is geen, althans geen kenbare gelijkwaardige, tegenprestatie aan {[}eiser 1{]} of {[}eiser 2{]} voldaan. In de feitelijke situatie zijn deze activa en de daarmee verbonden going-concernwaarde dus aan de boedels onttrokken, terwijl de opbrengst niet beschikbaar is gekomen voor de gezamenlijke schuldeisers. In de hypothetische situatie zonder deze overdracht zouden de curator en de boedels deze activa zelf hebben aangetroffen en te gelde hebben kunnen maken, al dan niet als onderdeel van een verkoop van de onderneming of van afzonderlijke activa. Dat levert vermindering van verhaal op. Daarbij weegt mee dat in beide faillissementen aanzienlijke erkende vorderingen en boedeltekorten bestaan, zodat niet kan worden aangenomen dat de schuldeisers ook zonder de handelingen volledig zouden zijn voldaan. Evenmin blijkt dat schuldeisers zich toch niet op deze activa hadden kunnen verhalen wegens geldige zekerheidsrechten of uitwinning door derden. Voor zover de overdracht van personeel of huurposities mogelijk ook lasten kon verminderen, neemt dat de benadeling niet weg, omdat het samenstel in zijn geheel ertoe strekte dat de onderneming/going concern-waarde bij {[}gedaagde 3{]} terechtkwam zonder vergoeding aan de boedels; bovendien is een door {[}gedaagde 3{]} betaalde salarispost zelfs via rekening-courant ten laste van {[}eiser 2{]} gebracht. Dit sluit aan bij de lijn in onder meer ECLI:NL:RBNNE:2021:3917, waarin het zonder vergoeding voortzetten van een onderneming met gebruikmaking van activa van de failliet als benadelend werd aangemerkt, en bij ECLI:NL:RBZWB:2025:5504 dat ook benadeling kan bestaan door onttrekking van verhaalsobjecten en verstoring van de paritas creditorum, ook wanneer een transactie boekhoudkundig neutraal lijkt. De schuldeisers zijn daarom in hun verhaalsmogelijkheden benadeeld.

\paragraph{Wetenschap - overwegingen}\label{wetenschap---overwegingen-22}

De wetenschap van benadeling is rechtstreeks voldoende af te leiden uit de feiten. De bestuurders van {[}eiser 1{]} en {[}eiser 2{]} waren al in 2020 door de boekhouder gewezen op ernstige continuïteitsproblemen, grote verliezen, belastingschulden en de mogelijkheid van faillissement; bovendien zijn cijfers aangepast om verliezen te maskeren. Kort daarna is {[}gedaagde 3{]} door {[}gedaagde 2{]} opgericht en zijn domeinnamen, huurposities, personeel en overige (softe) activa zonder reële tegenprestatie naar deze vennootschap overgeheveld. Ten tijde van deze rechtshandelingen waren faillissementen en tekorten met een redelijke mate van waarschijnlijkheid voorzienbaar. Die wetenschap kan ook aan {[}gedaagde 3{]} worden toegerekend, omdat zij door {[}gedaagde 2{]}, bestuurder van {[}eiser 1{]} en {[}eiser 2{]}, was opgericht/gecontroleerd en de onderneming voortzette. Daarnaast geldt, voor zover de handelingen om niet zijn verricht, het vermoeden van art. 45 Fw aan de zijde van de schuldenaren; tegenbewijs blijkt niet. Voor zover art. 43 Fw wordt toegepast, zijn de handelingen binnen de relevante periode onverplicht verricht, tegen aanmerkelijke onderwaarde/zonder tegenprestatie en met een gelieerde partij, zodat wetenschap aan beide zijden wordt vermoed. Ook dat vermoeden is niet ontzenuwd. Vergelijkbaar is de benadering in o.a. ECLI:NL:RBNNE:2021:3917 en ECLI:NL:RBZWB:2025:5504.

\newpage

\markboth{ECLI:NL:RBAMS:2026:3787}{gemini-3.5-flash (k = 5)}

\subsubsection{Gemini}\label{gemini-7}

\texttt{gemini-3.5-flash}

\begin{center}
\small
\begin{tabular}{@{}p{7.4cm}p{2.6cm}p{2.6cm}@{}}
\toprule
Variabele & Waarde (model) & Waarde (gold) \\
\midrule
Is rechtshandeling & \texttt{true} & \texttt{true} \\
Datum rechtshandeling & \texttt{2021-01-21} & \texttt{2021-01-21} \\
Verricht door de schuldenaar & \texttt{true} & \texttt{true} \\
Om niet & \texttt{true} & \texttt{true} \\
Datum faillietverklaring & \texttt{2021-07-13} & \texttt{2021-07-13} \\
Onverplicht & \texttt{true} & \texttt{true} \\
Benadeling schuldeisers & \texttt{true} & \texttt{true} \\
Wetenschap schuldenaar aangetoond & \textcolor{gray}{\texttt{true}} & \texttt{false} \\
Wetenschap wederpartij aangetoond & \textcolor{gray}{\texttt{true}} & \texttt{null} \\
Vermoeden wetenschap schuldenaar tegenbewijs & \texttt{false} & \texttt{false} \\
Vermoeden wetenschap wederpartij tegenbewijs & \texttt{false} & \texttt{false} \\
Tegen onderwaarde & \texttt{true} & \texttt{true} \\
Voldoening of zekerheid niet opeisbare schuld & \texttt{null} & \texttt{null} \\
Gelieerde partij & \texttt{true} & \texttt{true} \\
Reeds verplicht voor aanvang termijn & \texttt{false} & \texttt{false} \\
\textbf{Vernietigbaar} & \textbf{\texttt{true}} & \textbf{\texttt{true}} \\
\bottomrule
\end{tabular}
\end{center}

\paragraph{Is rechtshandeling - overwegingen}\label{is-rechtshandeling---overwegingen-23}

Het betreft een samenstel van handelingen, waaronder de overdracht van domeinnamen en huurrechten, die gericht zijn op een rechtsgevolg (de overgang van deze activa en rechten) en dus kwalificeren als rechtshandelingen.

\paragraph{Verricht door de schuldenaar - overwegingen}\label{verricht-door-de-schuldenaar---overwegingen-23}

De overdrachten van de domeinnamen en de huurrechten zijn rechtstreeks verricht door de schuldenaren zelf, te weten {[}eiser 1{]} en {[}eiser 2{]}.

\paragraph{Om niet - overwegingen}\label{om-niet---overwegingen-23}

De curator stelt dat de overdracht van de vermogensbestanddelen aan {[}gedaagde 3{]} heeft plaatsgevonden zonder dat daar een tegenprestatie tegenover stond, wat betekent dat de rechtshandelingen om niet zijn verricht.

\paragraph{Onverplicht - overwegingen}\label{onverplicht---overwegingen-23}

Het samenstel van rechtshandelingen dat de oneigenlijke doorstart van {[}eiser 1{]} en {[}eiser 2{]} naar {[}gedaagde 3{]} vormt (waaronder de overdracht van domeinnamen, de overdracht van de rechten uit de huurovereenkomsten en de overgang van het personeel) is onverplicht verricht. Er bestond geen enkele op de wet of een voorafgaande overeenkomst berustende rechtsplicht voor {[}eiser 1{]} en {[}eiser 2{]} om deze vermogensbestanddelen en rechten aan de nieuw opgerichte vennootschap {[}gedaagde 3{]} over te dragen. De overdrachten konden dan ook niet in rechte worden afgedwongen.

\paragraph{Benadeling schuldeisers - overwegingen}\label{benadeling-schuldeisers---overwegingen-23}

In deze casus is sprake van een samenstel van rechtshandelingen (overdracht van domeinnamen, huurrechten en personeel) die hebben geleid tot een oneigenlijke doorstart van de ondernemingen van {[}eiser 1{]} en {[}eiser 2{]} door {[}gedaagde 3{]}. Deze overdracht van vermogensbestanddelen heeft plaatsgevonden zonder dat {[}gedaagde 3{]} daarvoor een reële tegenprestatie heeft verricht. Door deze transacties zijn waardevolle activa (de `softe activa' en contractuele rechten) aan het vermogen van de schuldenaren onttrokken, waardoor het voor verhaal beschikbare actief is verminderd. Indien deze handelingen worden weggedacht (de hypothetische situatie), hadden deze activa in de boedel gezeten en hadden de schuldeisers zich daarop kunnen verhalen. Nu deze activa om niet of zonder reële tegenprestatie zijn overgedragen, zijn de schuldeisers daadwerkelijk in hun verhaalsmogelijkheden benadeeld. Dit sluit aan bij vaste jurisprudentie waarin het overhevelen van activa ten behoeve van een doorstart zonder reële vergoeding als paulianeus en benadelend wordt gekwalificeerd (vgl. ECLI:NL:RBNNE:2021:3917 en ECLI:NL:RBGEL:2015:6805).

\paragraph{Wetenschap - overwegingen}\label{wetenschap---overwegingen-23}

De wetenschap van benadeling is bij zowel de schuldenaren (eiser 1 en eiser 2) als de wederpartij (gedaagde 3) rechtstreeks aangetoond. De bestuurders (gedaagde 1 en gedaagde 2) waren door hun boekhouder in 2020 herhaaldelijk en expliciet gewaarschuwd dat de continuïteit niet langer gegarandeerd was en een faillissement niet ondenkbaar was. Desondanks hebben zij begin 2021 de activa (domeinnamen, huurrechten en personeel) overgedragen aan de nieuw opgerichte vennootschap gedaagde 3. Omdat gedaagde 3 is opgericht en wordt beheerst door gedaagde 2, had de wederpartij exact dezelfde wetenschap. Daarnaast is het bewijsvermoeden van art. 43 lid 1 sub 5 Fw (gelieerde partijen) van toepassing, aangezien de transacties binnen een jaar voor faillissement zijn verricht tussen vennootschappen met dezelfde (indirecte) bestuurder. Dit vermoeden is niet ontzenuwd.

\newpage

\markboth{ECLI:NL:RBAMS:2026:3787}{mistral/mistral-medium-3.5-128b:fp8 (k = 10)}

\subsection{k = 10}\label{k-10-2}

\subsubsection{Mistral}\label{mistral-8}

\texttt{mistral/mistral-medium-3.5-128b:fp8}

\begin{center}
\small
\begin{tabular}{@{}p{7.4cm}p{2.6cm}p{2.6cm}@{}}
\toprule
Variabele & Waarde (model) & Waarde (gold) \\
\midrule
Is rechtshandeling & \texttt{true} & \texttt{true} \\
Datum rechtshandeling & \texttt{2021-01-21} & \texttt{2021-01-21} \\
Verricht door de schuldenaar & \texttt{true} & \texttt{true} \\
Om niet & \texttt{true} & \texttt{true} \\
Datum faillietverklaring & \texttt{2021-07-13} & \texttt{2021-07-13} \\
Onverplicht & \texttt{true} & \texttt{true} \\
Benadeling schuldeisers & \texttt{true} & \texttt{true} \\
Wetenschap schuldenaar aangetoond & \texttt{false} & \texttt{false} \\
Wetenschap wederpartij aangetoond & \texttt{false} & \texttt{null} \\
Vermoeden wetenschap schuldenaar tegenbewijs & \texttt{null} & \texttt{false} \\
Vermoeden wetenschap wederpartij tegenbewijs & \texttt{null} & \texttt{false} \\
Tegen onderwaarde & \textcolor{gray}{\texttt{null}} & \texttt{true} \\
Voldoening of zekerheid niet opeisbare schuld & \texttt{null} & \texttt{null} \\
Gelieerde partij & \texttt{true} & \texttt{true} \\
Reeds verplicht voor aanvang termijn & \texttt{null} & \texttt{false} \\
\textbf{Vernietigbaar} & \textbf{\texttt{true}} & \textbf{\texttt{true}} \\
\bottomrule
\end{tabular}
\end{center}

\paragraph{Is rechtshandeling - overwegingen}\label{is-rechtshandeling---overwegingen-24}

De overdracht van domeinnamen (2.6.4) is een goederenrechtelijke handeling en kwalificeert als rechtshandeling. De overdracht van de rechten uit de huurovereenkomsten (2.6.4) is eveneens een rechtshandeling (overdracht van vorderingsrechten).

\paragraph{Verricht door de schuldenaar - overwegingen}\label{verricht-door-de-schuldenaar---overwegingen-24}

De overdrachten zijn verricht door {[}eiser 2{]} (domeinnamen) en {[}eiser 1{]} en {[}eiser 2{]} (huurovereenkomsten), die de schuldenaren zijn.

\paragraph{Om niet - overwegingen}\label{om-niet---overwegingen-24}

De curator stelt dat {[}gedaagde 3{]} geen tegenprestatie heeft verricht voor de verkrijging van vermogensbestanddelen (4.3.1). Er is geen sprake van een reële wederprestatie.

\paragraph{Onverplicht - overwegingen}\label{onverplicht---overwegingen-24}

De overdracht van de domeinnamen en de huurovereenkomsten door {[}eiser 1{]} en {[}eiser 2{]} aan {[}gedaagde 3{]} op respectievelijk 21 januari 2021 en 10 februari 2021 moet worden beoordeeld op onverplichtheid in de zin van art. 42 lid 1 Fw.

Uit de casus blijkt dat {[}eiser 1{]} en {[}eiser 2{]} zich bezighielden met de handel in gebruikte motorboten en dat {[}gedaagde 3{]} op 22 december 2020 is opgericht door {[}gedaagde 2{]}, die ook bestuurder was van {[}eiser 1{]} en {[}eiser 2{]}. De overdracht van de domeinnamen en huurovereenkomsten vond plaats in een periode waarin de financiële situatie van {[}eiser 1{]} en {[}eiser 2{]} verslechterde, zoals blijkt uit de aanpassingen in de administratie (2.4.3 en 2.4.4) en de meldingen van de boekhouder over de slechte financiële situatie (2.4.6).

Er is geen sprake van een op de wet of een overeenkomst berustende rechtsplicht voor {[}eiser 1{]} en {[}eiser 2{]} om de domeinnamen en huurovereenkomsten aan {[}gedaagde 3{]} over te dragen. De overdracht was derhalve onverplicht. Morele, feitelijke of economische druk (zoals de wens om de onderneming voort te zetten) schept geen rechtsplicht. De stelplicht en bewijslast van het onverplichte karakter rusten op de curator, en uit de casus blijkt dat de curator deze stelplicht heeft vervuld door aan te tonen dat er geen verplichting bestond om de overdracht te verrichten.

Daarnaast is de overdracht van de huurovereenkomsten een eenzijdige rechtshandeling van {[}eiser 1{]} en {[}eiser 2{]} gericht tot {[}gedaagde 3{]}. Voor de toepassing van art. 42 lid 2 Fw is vereist dat ook {[}gedaagde 3{]} wist of behoorde te weten dat de overdracht benadeling van schuldeisers tot gevolg zou hebben. Gelet op de nauwe banden tussen {[}gedaagde 2{]} (bestuurder van {[}eiser 1{]}, {[}eiser 2{]} en oprichter van {[}gedaagde 3{]}) en de financiële problemen van {[}eiser 1{]} en {[}eiser 2{]} is aannemelijk dat {[}gedaagde 3{]} hiervan op de hoogte was of behoorde te zijn.

De overdracht van de domeinnamen en huurovereenkomsten is derhalve onverplicht verricht in de zin van art. 42 Fw.

\paragraph{Benadeling schuldeisers - overwegingen}\label{benadeling-schuldeisers---overwegingen-24}

De curator vordert vernietiging van het samenstel van rechtshandelingen waarbij vermogensbestanddelen van {[}eiser 1{]} en {[}eiser 2{]} zijn overgedragen aan {[}gedaagde 3{]}, omdat dit een oneigenlijke doorstart zou inhouden. Uit de casus blijkt dat {[}eiser 1{]} en {[}eiser 2{]} in financiële moeilijkheden verkeerden: de boekhouder {[}naam 4{]} heeft herhaaldelijk gewaarschuwd voor de slechte financiële situatie (2.4.6), en de continuïteit van beide vennootschappen was niet gegarandeerd. Op 13 januari 2020 en 30 april 2020 meldde {[}naam 4{]} dat de continuïteit in gevaar was en dat een faillissement van {[}eiser 2{]} niet ondenkbaar was. Bovendien was er sprake van een aanzienlijk boedeltekort in beide faillissementen ({[}eiser 1{]}: € 886.548,01 + p.m.; {[}eiser 2{]}: € 2.208.533,29 + p.m.) (2.7.1).

De overdracht van vermogensbestanddelen (domeinnamen, huurovereenkomsten, personeel) aan {[}gedaagde 3{]} (2.6.4 en 2.6.5) heeft geleid tot een situatie waarin deze activa niet meer beschikbaar waren voor verhaalsdoeleinden in de faillissementsboedel van {[}eiser 1{]} en {[}eiser 2{]}. Hierdoor zijn de schuldeisers in hun verhaalsmogelijkheden benadeeld. De hypothetische situatie zonder deze overdrachten zou zijn geweest dat deze activa in de boedel van {[}eiser 1{]} en {[}eiser 2{]} waren gebleven en voor de schuldeisers beschikbaar waren geweest. In de feitelijke situatie zijn deze activa aan {[}gedaagde 3{]} overgedragen, waardoor de schuldeisers hierop geen verhaal meer kunnen halen.

Daarnaast is er sprake van een verstoring van de rangorde tussen schuldeisers, omdat de activa selectief aan {[}gedaagde 3{]} zijn overgedragen, terwijl andere schuldeisers (zoals de preferente en concurrente schuldeisers) hierop geen aanspraak meer kunnen maken. Dit voldoet aan de definitie van benadeling in de zin van artikel 42 Fw, omdat de schuldeisers in het algemeen zijn benadeeld door de vermindering van het actief in de boedel en de verstoring van de paritas creditorum.

De benadeling was aanwezig op het moment dat de curator zijn rechten deed gelden (leer van de middellijke benadeling), omdat de activa op dat moment niet meer in de boedel aanwezig waren en de schuldeisers hierdoor daadwerkelijk in hun verhaalsmogelijkheden waren beperkt.

Gelet op het bovenstaande zijn de schuldeisers door de aangevochten rechtshandelingen benadeeld in de zin van artikel 42 Fw. Dit wordt bevestigd door de uitspraak ECLI:NL:RBZWB:2025:5504 (rov. 4.11), waarin wordt overwogen dat benadeling ook aanwezig is wanneer goederen aan één crediteur zijn verkocht tegen een reële koopsom, maar de opbrengst niet voor de gezamenlijke crediteuren beschikbaar is. In deze casus zijn de activa zonder adequate tegenprestatie of vergoeding aan {[}gedaagde 3{]} overgedragen, wat de benadeling nog duidelijker maakt.

Tot slot is er geen sprake van een situatie waarin de schuldeisers toch niet op de betrokken goederen hadden kunnen verhalen (bijvoorbeeld door zekerheidsrechten), omdat niet is gesteld of gebleken dat dergelijke rechten bestonden ten aanzien van de overgedragen activa (zoals domeinnamen en huurovereenkomsten).

\paragraph{Wetenschap - overwegingen}\label{wetenschap---overwegingen-24}

De overdracht van domeinnamen en huurovereenkomsten door {[}eiser 1{]} en {[}eiser 2{]} aan {[}gedaagde 3{]} op 21 januari 2021 en 10 februari 2021 valt binnen de termijn van één jaar voor de faillietverklaring van {[}eiser 1{]} (13 juli 2021). {[}gedaagde 3{]} is opgericht door {[}gedaagde 2{]}, die bestuurder was van zowel {[}eiser 1{]} als {[}eiser 2{]}. Dit betekent dat {[}gedaagde 3{]} een gelieerde partij is in de zin van art. 43 lid 1 sub 3° Fw (rechtspersoon opgericht door een bloedverwant van de schuldenaar, die tevens bestuurder was). De rechtshandelingen zijn onverplicht verricht, omdat er geen contractuele of wettelijke verplichting bestond voor {[}eiser 1{]} en {[}eiser 2{]} om de domeinnamen en huurovereenkomsten aan {[}gedaagde 3{]} over te dragen. Bovendien is sprake van benadeling van schuldeisers, omdat deze activa uit de boedel zijn onttrokken zonder adequate tegenprestatie. Gelet op art. 43 lid 1 Fw wordt de wetenschap van benadeling aan beide zijden (schuldenaar en wederpartij) vermoed, tenzij tegenbewijs wordt geleverd. In deze casus is geen sprake van tegenbewijs: de financiële situatie van {[}eiser 1{]} en {[}eiser 2{]} was al in 2020 problematisch (verliezen, aanpassingen in de administratie om winst te tonen, meldingen van de boekhouder over slechte financiële situatie en mogelijke faillissementen). Derhalve is het vermoeden van wetenschap niet ontzenuwd.

\newpage

\markboth{ECLI:NL:RBAMS:2026:3787}{gpt-5.5 (k = 10)}

\subsubsection{GPT}\label{gpt-8}

\texttt{gpt-5.5}

\begin{center}
\small
\begin{tabular}{@{}p{7.4cm}p{2.6cm}p{2.6cm}@{}}
\toprule
Variabele & Waarde (model) & Waarde (gold) \\
\midrule
Is rechtshandeling & \texttt{true} & \texttt{true} \\
Datum rechtshandeling & \texttt{2021-01-21} & \texttt{2021-01-21} \\
Verricht door de schuldenaar & \texttt{true} & \texttt{true} \\
Om niet & \texttt{true} & \texttt{true} \\
Datum faillietverklaring & \texttt{2021-07-13} & \texttt{2021-07-13} \\
Onverplicht & \texttt{true} & \texttt{true} \\
Benadeling schuldeisers & \texttt{true} & \texttt{true} \\
Wetenschap schuldenaar aangetoond & \textcolor{gray}{\texttt{true}} & \texttt{false} \\
Wetenschap wederpartij aangetoond & \textcolor{gray}{\texttt{true}} & \texttt{null} \\
Vermoeden wetenschap schuldenaar tegenbewijs & \texttt{false} & \texttt{false} \\
Vermoeden wetenschap wederpartij tegenbewijs & \texttt{false} & \texttt{false} \\
Tegen onderwaarde & \texttt{true} & \texttt{true} \\
Voldoening of zekerheid niet opeisbare schuld & \texttt{null} & \texttt{null} \\
Gelieerde partij & \texttt{true} & \texttt{true} \\
Reeds verplicht voor aanvang termijn & \texttt{false} & \texttt{false} \\
\textbf{Vernietigbaar} & \textbf{\texttt{true}} & \textbf{\texttt{true}} \\
\bottomrule
\end{tabular}
\end{center}

\paragraph{Is rechtshandeling - overwegingen}\label{is-rechtshandeling---overwegingen-25}

Het gaat om een samenstel van op rechtsgevolg gerichte handelingen, waaronder de overdracht van domeinnamen aan {[}gedaagde 3{]} op 21 januari 2021 en de overdracht van rechten uit huurovereenkomsten aan {[}gedaagde 3{]} op 10 februari 2021. Dit zijn rechtshandelingen en geen louter feitelijke handelingen.

\paragraph{Verricht door de schuldenaar - overwegingen}\label{verricht-door-de-schuldenaar---overwegingen-25}

De overgedragen vermogensbestanddelen waren afkomstig van {[}eiser 1{]} en/of {[}eiser 2{]}. Uit de feiten volgt dat {[}eiser 2{]} haar domeinnamen aan {[}gedaagde 3{]} heeft overgedragen en dat {[}eiser 1{]} en {[}eiser 2{]} hun rechten uit huurovereenkomsten aan {[}gedaagde 3{]} hebben overgedragen.

\paragraph{Om niet - overwegingen}\label{om-niet---overwegingen-25}

De curator stelt dat {[}gedaagde 3{]} voor de verkrijging van de vermogensbestanddelen van {[}eiser 1{]} en {[}eiser 2{]} geen tegenprestatie heeft verricht. Op basis van de aangeleverde feiten blijkt geen bedongen reële wederprestatie.

\paragraph{Onverplicht - overwegingen}\label{onverplicht---overwegingen-25}

Voor het samenstel van rechtshandelingen waarbij vermogensbestanddelen van {[}eiser 1{]} en {[}eiser 2{]} -- waaronder domeinnamen, rechten uit huurovereenkomsten, personeel/bedrijfsactiviteiten en overige softe activa/goodwill -- aan de kort daarvoor opgerichte {[}gedaagde 3{]} zijn overgedragen, blijkt uit de casus niet van een vooraf bestaande, in rechte afdwingbare verplichting uit wet of overeenkomst. Dat de overdrachten mogelijk pasten in een beoogde doorstart of feitelijk/economisch wenselijk werden geacht, schept geen rechtsplicht. Beslissend is of {[}eiser 1{]} en {[}eiser 2{]} rechtens konden worden gedwongen deze activa en bedrijfsactiviteiten aan {[}gedaagde 3{]} over te dragen; daarvoor ontbreken aanknopingspunten. Eventuele verplichtingen die pas uit de overdrachtsovereenkomsten zelf voortvloeien, zijn voor de onverplichtheid niet relevant. Dit sluit aan bij rechtspraak over oneigenlijke doorstarts en activaoverdrachten binnen gelieerde verhoudingen, waarin verkoop of overdracht van activa zonder vooraf bestaande overdrachtsplicht als onverplicht wordt aangemerkt, zoals ECLI:NL:RBNNE:2021:3917 en ECLI:NL:RBGEL:2015:6805.

\paragraph{Benadeling schuldeisers - overwegingen}\label{benadeling-schuldeisers---overwegingen-25}

Ja. Bij de vereiste vermogensvergelijking moet de feitelijke situatie waarin het samenstel van handelingen in stand blijft worden vergeleken met de hypothetische situatie zonder die handelingen. Feitelijk zijn vóór het faillissement vermogensbestanddelen en ondernemingsonderdelen van {[}eiser 1{]} en {[}eiser 2{]} naar de nieuw opgerichte {[}gedaagde 3{]} overgegaan: domeinnamen, rechten uit huurovereenkomsten, personeel en kennelijk ook de aan de onderneming verbonden `softe' activa/goodwill en exploitatiemogelijkheden. Uit de casus blijkt niet van een reële tegenprestatie door {[}gedaagde 3{]}. Daarmee zijn activa en verhaalsmogelijkheden uit de vermogens van {[}eiser 1{]} en {[}eiser 2{]} verdwenen zonder dat daarvoor een gelijkwaardig actief in de boedel is teruggekomen. Zonder deze overdrachten hadden de curator en schuldeisers zich kunnen verhalen op deze activa, althans had de curator de onderneming of onderdelen daarvan kunnen verkopen. Dat is vergelijkbaar met de benadering in onder meer ECLI:NL:RBNNE:2021:3917 en ECLI:NL:RBGEL:2015:6805, waarin een oneigenlijke doorstart/overheveling van activa zonder vergoeding als benadelend werd aangemerkt. Ook de betaling door {[}gedaagde 3{]} van salarissen die vervolgens via rekening-courant ten laste van {[}eiser 2{]} is gebracht, vormt geen compenserend voordeel voor de gezamenlijke schuldeisers; integendeel, zij vergroot of bevestigt juist een last bij {[}eiser 2{]} terwijl de onderneming bij {[}gedaagde 3{]} wordt voortgezet. Van omstandigheden die benadeling uitsluiten, zoals volledig verhaal voor alle schuldeisers of bestaande zekerheidsrechten waardoor de schuldeisers zich toch niet op de goederen hadden kunnen verhalen, blijkt niet. Integendeel: in beide faillissementen bestaan aanzienlijke erkende vorderingen en grote boedeltekorten. De schuldeisers zijn daarom in hun verhaalsmogelijkheden benadeeld.

\paragraph{Wetenschap - overwegingen}\label{wetenschap---overwegingen-25}

De overdrachten aan {[}gedaagde 3{]} vonden plaats in het kader van een oneigenlijke doorstart, kort na oprichting van {[}gedaagde 3{]} door {[}gedaagde 2{]} en binnen één jaar vóór het faillissement van {[}eiser 1{]}. Uit de meldingen van de boekhouder in 2020, de zeer slechte resultaten/negatieve saldi, oplopende belastingschulden en de aanpassingen van de administratie volgt dat de bestuurders van {[}eiser 1{]} en {[}eiser 2{]} wisten of behoorden te weten dat faillissement en een tekort daarin met een redelijke mate van waarschijnlijkheid voorzienbaar waren. Die wetenschap kan ook aan {[}gedaagde 3{]} worden toegerekend, nu zij door {[}gedaagde 2{]} is opgericht en de activa/activiteiten zonder reële tegenprestatie naar deze nieuwe vennootschap zijn overgebracht. Voor zover de handelingen als om niet moeten worden aangemerkt, wordt de wetenschap van de schuldenaar bovendien vermoed op grond van art. 45 Fw; er zijn geen feiten gesteld waaruit succesvol tegenbewijs volgt. Daarnaast ligt een bewijsvermoeden van art. 43 Fw voor de hand wegens de gelieerdheid tussen de rechtspersonen en het ontbreken van een reeds vóór de jaartermijn bestaande verplichting; ook dat vermoeden is niet ontzenuwd. De situatie is vergelijkbaar met de doorstart-/activa-overhevelingszaken waarin wetenschap werd aangenomen bij gelieerde partijen, zoals ECLI:NL:RBNNE:2021:3917 en ECLI:NL:RBGEL:2015:6805.

\newpage

\markboth{ECLI:NL:RBAMS:2026:3787}{gemini-3.5-flash (k = 10)}

\subsubsection{Gemini}\label{gemini-8}

\texttt{gemini-3.5-flash}

\begin{center}
\small
\begin{tabular}{@{}p{7.4cm}p{2.6cm}p{2.6cm}@{}}
\toprule
Variabele & Waarde (model) & Waarde (gold) \\
\midrule
Is rechtshandeling & \texttt{true} & \texttt{true} \\
Datum rechtshandeling & \texttt{2021-01-21} & \texttt{2021-01-21} \\
Verricht door de schuldenaar & \texttt{true} & \texttt{true} \\
Om niet & \texttt{true} & \texttt{true} \\
Datum faillietverklaring & \texttt{2021-07-13} & \texttt{2021-07-13} \\
Onverplicht & \texttt{true} & \texttt{true} \\
Benadeling schuldeisers & \texttt{true} & \texttt{true} \\
Wetenschap schuldenaar aangetoond & \textcolor{gray}{\texttt{true}} & \texttt{false} \\
Wetenschap wederpartij aangetoond & \textcolor{gray}{\texttt{true}} & \texttt{null} \\
Vermoeden wetenschap schuldenaar tegenbewijs & \texttt{false} & \texttt{false} \\
Vermoeden wetenschap wederpartij tegenbewijs & \texttt{false} & \texttt{false} \\
Tegen onderwaarde & \textcolor{gray}{\texttt{null}} & \texttt{true} \\
Voldoening of zekerheid niet opeisbare schuld & \texttt{null} & \texttt{null} \\
Gelieerde partij & \texttt{true} & \texttt{true} \\
Reeds verplicht voor aanvang termijn & \texttt{false} & \texttt{false} \\
\textbf{Vernietigbaar} & \textbf{\texttt{true}} & \textbf{\texttt{true}} \\
\bottomrule
\end{tabular}
\end{center}

\paragraph{Is rechtshandeling - overwegingen}\label{is-rechtshandeling---overwegingen-26}

De overdracht van domeinnamen, huurrechten en de overgang van personeel betreffen rechtshandelingen, aangezien deze gericht zijn op het bewerkstelligen van rechtsgevolgen (de overgang van eigendom en contractuele posities).

\paragraph{Verricht door de schuldenaar - overwegingen}\label{verricht-door-de-schuldenaar---overwegingen-26}

De overdrachten van de domeinnamen en de huurrechten zijn verricht door de schuldenaren {[}eiser 1{]} en {[}eiser 2{]} zelf.

\paragraph{Om niet - overwegingen}\label{om-niet---overwegingen-26}

De curator stelt dat de overdracht van de vermogensbestanddelen heeft plaatsgevonden zonder dat {[}gedaagde 3{]} daarvoor enige tegenprestatie heeft verricht, waardoor de rechtshandelingen om niet zijn verricht.

\paragraph{Onverplicht - overwegingen}\label{onverplicht---overwegingen-26}

Het samenstel van rechtshandelingen waarbij {[}eiser 1{]} en {[}eiser 2{]} vermogensbestanddelen (zoals domeinnamen en huurrechten) hebben overgedragen aan {[}gedaagde 3{]} ten behoeve van een doorstart, is onverplicht verricht. Er bestond geen op de wet of een voorafgaande overeenkomst berustende rechtsplicht die {[}eiser 1{]} en {[}eiser 2{]} verplichtte tot deze overdrachten, waardoor de prestaties niet in rechte konden worden afgedwongen.

\paragraph{Benadeling schuldeisers - overwegingen}\label{benadeling-schuldeisers---overwegingen-26}

Op basis van de casus is er sprake van benadeling van de schuldeisers in de zin van artikel 42 Faillissementswet (Fw). De vennootschappen {[}eiser 1{]} en {[}eiser 2{]} hebben waardevolle activa, waaronder domeinnamen en de rechten uit huurovereenkomsten, overgedragen aan {[}gedaagde 3{]}. Daarnaast is het personeel overgegaan naar {[}gedaagde 3{]}, waarbij de salariskosten via de rekening-courant ten laste van {[}eiser 2{]} zijn gebracht. De curator stelt dat voor deze overdrachten, die samenhangen met een oneigenlijke doorstart, geen enkele tegenprestatie is verricht. Door een vermogensvergelijking toe te passen tussen de feitelijke situatie en de hypothetische situatie waarin de overdrachten niet zouden hebben plaatsgevonden, blijkt dat het verhaalsactief van de schuldenaren substantieel is verminderd. Er zijn immers vermogensbestanddelen uit het vermogen van de schuldenaren verdwenen zonder dat daar actieve, gelijkwaardige middelen voor in de plaats zijn gekomen waarop schuldeisers zich hadden kunnen verhalen. Dit sluit aan bij de lijn in de rechtspraak (zoals ECLI:NL:RBNNE:2021:3917 en ECLI:NL:RBGEL:2015:6805), waarin wordt geoordeeld dat het overhevelen van activa of het bewerkstelligen van een doorstart zonder reële vergoeding leidt tot daadwerkelijke benadeling van de gezamenlijke schuldeisers.

\paragraph{Wetenschap - overwegingen}\label{wetenschap---overwegingen-26}

De wetenschap van benadeling aan de zijde van de schuldenaren (Eiser 1 en Eiser 2) is rechtstreeks aangetoond. De extern boekhouder had de bestuurders (Gedaagde 1 en Gedaagde 2) al in januari en september 2020 expliciet gewaarschuwd dat de continuïteit in gevaar was en een faillissement ondenkbaar was, waarna de administratie zelfs werd gemanipuleerd om verliezen te verhullen. De wetenschap van de wederpartij (Gedaagde 3) is eveneens rechtstreeks aangetoond, aangezien Gedaagde 3 is opgericht en wordt geleid door Gedaagde 2, die tevens bestuurder was van de schuldenaren; deze wetenschap wordt aan de rechtspersoon toegerekend. Daarnaast is sprake van een bewijsvermoeden op grond van art. 43 lid 1 sub 5 Fw (gelieerde partijen) omdat de rechtshandelingen binnen een jaar voor faillissement zijn verricht tussen vennootschappen met dezelfde (indirecte) bestuurder. Dit vermoeden is niet ontzenuwd.

\newpage

\markboth{ECLI:NL:RBOVE:2026:4561}{Gold standard}

\section{ECLI:NL:RBOVE:2026:4561}\label{eclinlrbove20264561}

uitspraakdatum 14-07-2026

\subsection{Gold}\label{gold-3}

\begin{center}
\small
\begin{tabular}{@{}p{8cm}p{2.6cm}@{}}
\toprule
Variabele & Waarde \\
\midrule
Is rechtshandeling & \texttt{true} \\
Datum rechtshandeling & \texttt{2023-10-23} \\
Verricht door de schuldenaar & \texttt{true} \\
Om niet & \texttt{true} \\
Datum faillietverklaring & \texttt{2024-08-20} \\
Onverplicht & \texttt{true} \\
Benadeling schuldeisers & \texttt{true} \\
Wetenschap schuldenaar aangetoond & \texttt{false} \\
Wetenschap wederpartij aangetoond & \texttt{null} \\
Vermoeden wetenschap schuldenaar tegenbewijs & \texttt{false} \\
Vermoeden wetenschap wederpartij tegenbewijs & \texttt{null} \\
Tegen onderwaarde & \texttt{null} \\
Voldoening of zekerheid niet opeisbare schuld & \texttt{null} \\
Gelieerde partij & \texttt{null} \\
Reeds verplicht voor aanvang termijn & \texttt{null} \\
\textbf{Vernietigbaar} & \textbf{\texttt{true}} \\
\bottomrule
\end{tabular}
\end{center}

\subsubsection{Uitspraak}\label{uitspraak-3}

\textbf{RECHTBANK}
\textbf{OVERIJSSEL}

Civiel recht

Kantonrechter

Zittingsplaats Zwolle

Zaaknummer: 12140745 ~CV EXPL 26-825

\textbf{Vonnis van 14 juli 2026}

in de zaak van

\textbf{MR. MARK LOEF},

in hoedanigheid van curator in het faillissement van \textbf{{[}naam 1{]}},

wonende in {[}woonplaats 1{]} ,

eisende partij,

hierna te noemen: de curator,

gemachtigde: mr. C. ter Braack

tegen

\textbf{{[}gedaagde{]}}
,

wonende in {[}woonplaats 2{]} ,

gedaagde partij,

hierna te noemen: {[}gedaagde{]} ,

gemachtigde: mr. R. Dufour.

\paragraph{1 De zaak in het kort}\label{de-zaak-in-het-kort}

1.1 {[}naam 1{]} , de vader van {[}gedaagde{]} , is failliet gegaan. In de boedel heeft de curator (onder meer) een woning aangetroffen, waarin op dat moment onder andere {[}gedaagde{]} woonde. De curator meent dat er sprake was van een huurovereenkomst tussen {[}gedaagde{]} en zijn vader, en vraagt in deze procedure om betaling van achterstallige huur. Verder vraagt de curator om betaling van een bedrag dat {[}gedaagde{]} van zijn vader heeft ontvangen voorafgaand aan het faillissement, waarbij het volgens de curator gaat om een vernietigbare schenking.

1.2 De kantonrechter is kort gezegd van oordeel dat niet kan worden vastgesteld dat {[}gedaagde{]} de woning van zijn vader huurde. Wel is sprake geweest van een schenking, die de curator met succes heeft vernietigd. Dat bedrag moet {[}gedaagde{]} aan de curator betalen.

\paragraph{2 De procedure}\label{de-procedure}

2.1 Het verloop van de procedure blijkt uit:

\begin{itemize}
\item
  de dagvaarding,\\
\item
  de conclusie van antwoord,\\
\item
  de brief waarin is meegedeeld dat een mondelinge behandeling is bepaald,
\item
  de nadere producties van de curator,
\item
  de nadere producties van {[}gedaagde{]} ,
\item
  de mondelinge behandeling van 16 juni 2026, waarvan door de griffier aantekeningen zijn gemaakt,
\item
  de spreekaantekeningen van beide partijen.
\end{itemize}

\paragraph{3 De feiten}\label{de-feiten}

3.1 De vader van {[}gedaagde{]} , de heer {[}naam 1{]} , is bij vonnis van 20 augustus 2024 van de rechtbank Overijssel failliet verklaard. Mr.~Loef is aangesteld als curator.

3.2 De curator heeft in de boedel onder andere de woning aan de {[}adres{]} (hierna: de woning) aangetroffen. In de woning woonden op dat moment onder anderen {[}gedaagde{]} , zijn jongere broer {[}naam 2{]} en de voormalige partner van {[}naam 1{]} , mevrouw {[}naam 3{]} (hierna: {[}naam 3{]} ).

3.3 {[}gedaagde{]} en zijn broer {[}naam 2{]} hebben de woning op 27 maart 2025 van de curator uit de faillissementsboedel gekocht en op {[}datum{]} 2025 geleverd gekregen.

3.4 {[}naam 1{]} heeft op 23 oktober 2023, voorafgaand aan zijn faillissement, een bedrag van € 3.050,00 naar {[}gedaagde{]} overgemaakt. Bij de omschrijving van deze overboeking staat `schenking'. Bij brief van 16 juli 2025 aan {[}gedaagde{]} heeft de curator verklaard de schenking te vernietigen en {[}gedaagde{]} aangemaand om dit bedrag aan de boedel terug te betalen.

\paragraph{4 Het geschil}\label{het-geschil}

4.1 De curator vordert -- samengevat -- veroordeling van {[}gedaagde{]} tot betaling van € 21.500,00, te vermeerderen met de wettelijke rente. Daarnaast vordert de curator een verklaring voor recht dat hij de schenking buitengerechtelijk heeft vernietigd en veroordeling van {[}gedaagde{]} tot betaling van € 3.050,00, te vermeerderen met de wettelijke rente. Tot slot vordert de curator betaling van een bedrag van € 1.020,50 aan buitengerechtelijke incassokosten, en veroordeling van {[}gedaagde{]} in de proceskosten.

4.2 {[}gedaagde{]} voert verweer.

4.3 Op de stellingen van partijen wordt hierna, voor zover nodig, nader ingegaan.

\paragraph{5 De beoordeling}\label{de-beoordeling}

\emph{De vordering tot betaling van € 21.500,00}

5.1 Volgens de curator is {[}gedaagde{]} € 21.500,00 aan achterstallige huur aan de boedel verschuldigd. De curator stelt dat {[}gedaagde{]} de woning van zijn vader heeft gehuurd voor een huursom van € 500,00 per maand. {[}gedaagde{]} is volgens de curator huur over een periode van 43 maanden verschuldigd, gerekend vanaf de achttiende verjaardag van {[}gedaagde{]} tot aan de levering van de woning aan {[}gedaagde{]} en zijn broer op {[}datum{]} 2025.

5.2 Naar het oordeel van de kantonrechter kan niet worden vastgesteld dat een huurovereenkomst tussen {[}gedaagde{]} en zijn vader heeft bestaan. De stelplicht en de bewijslast voor het bestaan van de huurovereenkomst rusten op de curator. Tegenover de gemotiveerde betwisting van {[}gedaagde{]} , heeft de curator zijn stelling dat {[}gedaagde{]} de woning van zijn vader huurde onvoldoende onderbouwd. Dit wordt als volgt toegelicht.

5.2.1 {[}gedaagde{]} is op 1 oktober 2018, toen hij veertien jaar oud was, samen met zijn vader en zijn broer verhuisd naar de woning, waarvan zijn vader eigenaar was. {[}gedaagde{]} en zijn broer wonen daar nu nog steeds. Het is in zijn algemeenheid niet gebruikelijk dat een ouder huur vraagt van een minderjarig inwonend kind, en ook niet, zo heeft {[}gedaagde{]} onweersproken aangevoerd, dat een ouder huur van zijn kind gaat verlangen om in zijn woning te mogen verblijven nadat het achttien jaar is geworden. Daarbij wordt opgemerkt dat vast staat dat {[}gedaagde{]} vader zelf ook in de woning woonde gedurende de minderjarigheid van {[}gedaagde{]} , en ook toen {[}gedaagde{]} achttien jaar werd (op {[}geboortedatum{]} 2021). Ook als zou worden uitgegaan van de lezing van de curator dat {[}gedaagde{]} vader uit de woning is vertrokken, gebeurde dat later pas, in mei 2022.\footnote{De curator stelt in nr. 7 van de dagvaarding dat {[}naam 1{]} vanaf 23 mei 2022 niet meer woonachtig is in de woning.}

5.2.2 Bovendien heeft {[}gedaagde{]} nooit enige huurbetaling aan zijn vader gedaan, zo heeft hij onweersproken aangevoerd.

5.2.3 In dit licht bezien kan wat de curator ter onderbouwing van zijn stelling heeft aangedragen, niet de conclusie rechtvaardigen dat er wel een huurovereenkomst tussen {[}gedaagde{]} en zijn vader bestond.

5.2.4 De curator beroept zich er in de eerste plaats op dat {[}gedaagde{]} zelf aan hem heeft meegedeeld dat hij de woning huurde. In dit verband is de volgende gang van zaken van belang. Op 8 oktober 2024 heeft de curator {[}gedaagde{]} en de andere bewoners bericht dat zij zonder recht of titel in de woning verbleven en dat zij voor dat verblijf samen een vergoeding van € 1.600,00 per maand verschuldigd waren. Daarbij heeft de curator hen ook aangemaand om de woning binnen vier weken te verlaten. Op 15 november 2024 berichtte {[}gedaagde{]} de curator dat hij interesse had in de koop van de woning, waarbij hij opmerkte dat hij de woning op dat moment huurde. Vervolgens schreef {[}gedaagde{]} op 15 december 2024 aan de curator: ``\emph{In september ontvingen we een brief over een achterstand van € 1.600,- per maand. Wij willen graag weten vanaf welke datum deze achterstand volgens jullie betaald moet worden. Wij kunnen samen met mijn broer € 500,- per maand betalen (€ 250,- ieder), met ingang van oktober. Eerder hebben we al veel bedragen verrekend''}. Verder berichtte {[}gedaagde{]} de curator op 20 december 2024: ``\emph{Wat betreft de huur: we zullen de huur vanaf oktober betalen en kunnen daarbij op dit moment samen met mijn broer een bedrag van maximaal € 500,- missen''.}

{[}gedaagde{]} heeft op de zitting verklaard dat zijn eerdere mededeling aan de curator dat hij huurder was, onwaar is. {[}gedaagde{]} was bang om door de curator uit de woning te worden gezet, en beweerde daarom dat hij huurder was. Zijn vader had hem geadviseerd dit te zeggen, aldus {[}gedaagde{]} . Hoewel het {[}gedaagde{]} valt aan te rekenen dat hij naar eigen zeggen opzettelijk onjuiste informatie aan de curator heeft verstrekt, acht de kantonrechter die verklaring plausibel. Er kan daarom niet zonder meer op basis van de eerdere mededelingen van {[}gedaagde{]} aan de curator van worden uitgegaan dat {[}gedaagde{]} de woning daadwerkelijk van zijn vader huurde. Hetzelfde geldt voor de stelling van de curator dat ook de vader van {[}gedaagde{]} tegen hem heeft gezegd dat {[}gedaagde{]} huurder was van de woning.

5.2.5 Daarnaast beroept de curator zich op een rapport van Hypotheek Visie van 30 oktober 2023 met een hypotheekadvies aan {[}gedaagde{]} en {[}naam 3{]} , die destijds van plan waren om samen de woning te kopen. Op pagina 4 van het rapport staat als huidige woonsituatie van {[}gedaagde{]} en {[}naam 3{]} vermeldt: ``\emph{Momenteel woont u in een huurwoning: {[}adres{]} , {[}adres{]}}''.

Uit het gegeven dat de woning hier als huurwoning wordt aangeduid, kan echter niet zonder meer worden afgeleid dat {[}gedaagde{]} de woning huurde. Er waren naast {[}gedaagde{]} immers ook andere bewoners die volgens de curator als huurder moeten worden aangemerkt. Bovendien staat op pagina 4 van het rapport bij de vermelde huidige woonlasten van {[}gedaagde{]} en {[}naam 3{]} , geen bedrag aan huur opgenomen.

5.2.6 Tot slot wijst de curator op een e-mail die {[}naam 1{]} op 6 augustus 2024 ontving van incassokantoor Flanderijn, ingeschakeld door hypotheekhouder ABN AMRO. Daarin staat dat het Nederlands Woning Waarde Instituut (NWWI) een taxatierapport van de woning niet wilde valideren omdat volgens het NWWI sprake was van kamerverhuur.

Ook hieruit kan echter nog niet worden opgemaakt dat {[}gedaagde{]} huurder was. Het is namelijk onduidelijk waarop de constatering van het NWWI is gebaseerd. Daar komt bij dat ook in de betreffende e-mail staat dat de taxateur het niet eens is met het NWWI, omdat de inrichting van het pand niet is gericht op kamerverhuur.

5.3 De kantonrechter begrijpt dat de curator subsidiair aan zijn vordering ten grondslag legt dat {[}gedaagde{]} een (gebruiks)vergoeding is verschuldigd op grond van ongerechtvaardigde verrijking (als bedoeld in artikel 6:212 BW). Dat {[}gedaagde{]} ongerechtvaardigd is verrijkt vanwege zijn kosteloze verblijf in de woning, is echter niet komen vast te staan. Voor zover al zou kunnen worden aangenomen dat {[}naam 1{]} of de boedel financieel nadeel hebben geleden als rechtstreeks gevolg van {[}gedaagde{]} kosteloze verblijf in de woning, geldt dat de curator onvoldoende heeft toegelicht dat geen redelijke grond aanwezig is voor het voordeel dat {[}gedaagde{]} daarvan heeft gehad. Nu niet is komen vast te staan dat {[}naam 1{]} een vergoeding van {[}gedaagde{]} heeft verlangd voor zijn verblijf in de woning, moet ervan uit worden gegaan dat {[}naam 1{]} vrijwillig (stilzwijgend) heeft toegestaan dat {[}gedaagde{]} kosteloos in de woning verbleef. Het stond {[}naam 1{]} (uiteraard) vrij om {[}gedaagde{]} , zijn zoon, gratis onderdak te verschaffen, nog los van de vraag of hij daartoe gehouden was op grond van zijn onderhoudsplicht. De curator heeft onvoldoende uitgelegd dat, en vanaf wanneer, het voordeel dat {[}gedaagde{]} heeft gehad van zijn kosteloze verblijf in de woning, als gevolg van het faillissement vanaf enig moment wel als ongerechtvaardigd moet worden aangemerkt.

5.4 Gelet op het voorgaande wordt de vordering tot betaling van € 21.500,00 afgewezen.

\emph{De vernietiging van de betaling van € 3.050,00}

5.5 De vader van {[}gedaagde{]} heeft voorafgaand aan zijn faillietverklaring, op 23 oktober 2023, een bedrag van € 3.050,00 aan {[}gedaagde{]} overgemaakt. De overboeking is omschreven als `schenking'. De curator heeft verklaard deze betaling buitengerechtelijk te vernietigen op grond van artikel 42 lid 1 Faillissementswet (Fw), de zogeheten actio pauliana.

5.6 Het beroep van de curator op vernietiging van de betaling slaagt. Als gevolg van de vernietiging kan de curator het bedrag van € 3.050,00 als onverschuldigd betaald van {[}gedaagde{]} terugvorderen.

5.6.1 Vast is komen te staan dat de betaling van € 3.050,00 van {[}naam 1{]} aan {[}gedaagde{]} een onverplichte rechtshandeling om niet (zonder tegenprestatie) betreft, in dit geval een schenking. {[}gedaagde{]} heeft dit onvoldoende betwist. {[}gedaagde{]} heeft wisselende en onderling tegenstrijdige verklaringen afgelegd over het doel van de door hem ontvangen betaling. {[}gedaagde{]} heeft in zijn antwoord op de dagvaarding aangevoerd dat de betaling bedoeld was voor {[}gedaagde{]} om in zijn levensonderhoud te voorzien, maar dat {[}gedaagde{]} , om zijn broer niet te benadelen, in tweede instantie met zijn vader heeft besproken om de betaling als lening aan te merken. {[}gedaagde{]} zou die lening vervolgens aan zijn vader hebben terugbetaald, en wel, op instructie van zijn vader, door twee overschrijvingen op 16 januari 2024 van € 1.000,00 en € 800,00 aan {[}bedrijf{]} B.V. (een vennootschap waarbij zijn vader was betrokken), en twee overschrijvingen op 17 juni 2024 en 5 juli 2024 van € 1.000,00 aan zijn vader. Op de zitting heeft {[}gedaagde{]} echter verklaard dat het bij de overboeking door zijn vader ging om een terugbetaling van een lening die {[}gedaagde{]} eerder aan zijn vader had verstrekt voor de bekostiging van `basisbenodigdheden'. Alleen al vanwege het gebrek aan consistentie van het betoog van {[}gedaagde{]} wordt daaraan voorbij gegaan. Daar komt bij, wat betreft de verklaring van {[}gedaagde{]} dat de overboeking door zijn vader een lening betrof die hij heeft terugbetaald, dat {[}gedaagde{]} niet steekhoudend heeft toegelicht waarom de hoogte van zijn vermeende terugbetalingen niet overeenkomt met de hoogte van de lening, en waarom zijn vader hem zou instrueren om de terugbetalingen deels aan de genoemde vennootschap te doen en niet aan hemzelf. Wat betreft de verklaring van {[}gedaagde{]} dat de overboeking door zijn vader een terugbetaling betrof van een lening die hij eerder aan zijn vader had verstrekt, komt daar bij dat {[}gedaagde{]} niet heeft uitgelegd waarom zijn vader de overboeking als schenking heeft omschreven.

5.6.2 De schuldeisers van {[}naam 1{]} zijn door de schenking benadeeld, zo stelt de kantonrechter vast. Door de overboeking aan {[}gedaagde{]} zijn de betreffende gelden niet beschikbaar voor verdeling onder de schuldeisers. Anders dan {[}gedaagde{]} betoogt, blijkt uit de stukken die de curator heeft ingebracht dat er naast de curator zelf ook andere schuldeisers in het faillissement zijn. Verder wordt op grond van artikel 45 Fw vermoed dat Qingjian {[}gedaagde{]} wist of behoorde te weten dat de schuldeisers als gevolg van de schenking zijn benadeeld. {[}gedaagde{]} heeft dat vermoeden niet ontkracht.

5.7 {[}gedaagde{]} beroept zich op verrekening van zijn schuld aan de boedel met een vordering die hij op de boedel stelt te hebben op grond van een huurovereenkomst die hij als verhuurder met zijn vader heeft gesloten na diens faillissement. {[}gedaagde{]} licht toe dat zijn vader de verschuldigde huurprijs tot dusver onbetaald heeft gelaten, en dat die huurprijs op grond van artikel 39 Fw een boedelschuld is.

Dit beroep wordt verworpen. Uit artikel 24 Fw volgt dat de boedel niet aansprakelijk is voor de huurverplichtingen die {[}naam 1{]} na zijn faillietverklaring met {[}gedaagde{]} is aangegaan, tenzij de boedel daarbij zou zijn gebaat. Dat laatste is niet het geval. Niet gesteld of gebleken is dat de huurovereenkomst waaraan {[}naam 1{]} zich als huurder heeft verbonden, voor de boedel een bate oplevert. Artikel 39 lid 1 Fw is niet van toepassing, anders dan {[}gedaagde{]} betoogt, aangezien die bepaling ziet op een huurovereenkomst die al bestond ten tijde van de faillietverklaring.

5.8 Gelet op het voorgaande is de gevorderde verklaring voor recht en de gevorderde betaling van € 3.050,00 toewijsbaar.

5.9 De curator vordert wettelijke rente over dit bedrag vanaf de dag dat {[}gedaagde{]} de schenking heeft ontvangen, namelijk 23 oktober 2023. De kantonrechter zal de gevorderde wettelijke rente pas toewijzen vanaf 31 juli 2025. Hoewel de vernietiging van de schenking terugwerkende kracht heeft, kan niet worden aangenomen dat {[}gedaagde{]} zonder ingebrekestelling in verzuim is geraakt. De curator heeft namelijk niet gesteld dat {[}gedaagde{]} het bedrag te kwader trouw heeft aangenomen (zie artikel 6:205 BW). De curator heeft {[}gedaagde{]} op 16 juli 2025 aangemaand om het bedrag van € 3.050,00 terug te betalen binnen veertien dagen na ontvangt van de brief. Deze aanmaning kan naar het oordeel van de kantonrechter worden aangemerkt als een ingebrekestelling, zodat {[}gedaagde{]} na het verloop van de gestelde veertiendagentermijn op 31 juli 2025 in verzuim is geraakt (artikel 6:82 BW).

\emph{Buitengerechtelijke incassokosten}

5.10 De curator vordert vergoeding van buitengerechtelijke incassokosten ten bedrage van € 1.020,50, te vermeerderen met rente. De gevorderde vergoeding zal worden toegewezen tot een bedrag van € 430,00. Gelet op de aard van de toegewezen vordering is het Besluit vergoeding voor buitengerechtelijke incassokosten niet van toepassing, zodat de vraag of buitengerechtelijke incassokosten zijn verschuldigd wordt getoetst aan de eisen die zijn geformuleerd in het Rapport BGK-integraal. De curator heeft toegelicht en onderbouwd dat hij meerdere brieven en e-mails aan {[}gedaagde{]} en zijn (voormalige) gemachtigde heeft geschreven om de vordering buitengerechtelijk te incasseren. Daarmee heeft de curator voldoende gesteld dat hij buitengerechtelijke kosten heeft gemaakt die voor vergoeding in aanmerking komen. {[}gedaagde{]} heeft de door de curator gestelde buitengerechtelijke werkzaamheden niet betwist. Gelet op de hoogte van de toegewezen vordering heeft de curator op basis van het toepasselijke tarief, aanspraak op vergoeding van € 430,00. De gevorderde wettelijke rente over deze vergoeding zal worden toegewezen vanaf de dag van de dagvaarding (16 februari 2026), nu niet is gesteld dat de buitengerechtelijke kosten op een eerdere datum zijn betaald.

\emph{De proceskosten}

5.11 Omdat beide partijen gedeeltelijk ongelijk krijgen, zullen de proceskosten tussen hen worden gecompenseerd, in die zin dat iedere partij de eigen kosten draagt.

\emph{Uitvoerbaar bij voorraad?}

5.12 {[}gedaagde{]} heeft zich verzet tegen de gevorderde uitvoerbaar bij voorraadverklaring van dit vonnis. Bij de beoordeling van dit verweer van {[}gedaagde{]} moeten de belangen van partijen worden afgewogen in het licht van de omstandigheden van het geval. Het belang van de curator bij uitvoerbaarverklaring bij voorraad is gegeven. Daar staat tegenover dat naar het oordeel van de kantonrechter {[}gedaagde{]} een reëel restitutierisico loopt. Als de vordering van de curator in een eventueel hoger beroep alsnog wordt afgewezen, zal de curator een mogelijk bij voorraad uitgewonnen geldsom als onverschuldigd moeten terugbetalen aan {[}gedaagde{]} . De vordering tot terugbetaling dient in een dergelijk geval te worden behandeld als een concurrente boedelvordering (zie het arrest van de Hoge Raad van 7 juni 2002, ECLI:NL:HR:2002:AE3796). In het algemeen krijgen concurrente boedelschuldeisers in een faillissement niet hun gehele vordering uitbetaald. Uit de stellingen van de curator kan niet worden afgeleid dat dit in dit geval wel zo zou zijn. De afweging van de belangen van de curator en {[}gedaagde{]} leidt er toe dat de kantonrechter de curator gevraagde uitvoerbaarverklaring bij voorraad wel zal uitspreken, maar onder voorwaarde van zekerheidstelling in de vorm van een bankgarantie. Het belang van {[}gedaagde{]} bij zekerheidstelling, gegeven het restitutierisico, weegt zwaarder dan het belang van de curator bij tenuitvoerlegging zonder de voorwaarde van zekerheidstelling.

\paragraph{6 De beslissing}\label{de-beslissing-3}

De kantonrechter

6.1 verklaart voor recht dat de betaling aan {[}gedaagde{]} van het bedrag van € 3.050,00 buitengerechtelijk door de curator is vernietigd,

6.2 veroordeelt {[}gedaagde{]} om het bedrag van € 3.050,00 aan de curator te betalen, te vermeerderen met de wettelijke rente als bedoeld in artikel 6:119 BW vanaf 31 juli 2025 tot de dag van volledige betaling,

6.3 veroordeelt {[}gedaagde{]} om € 430,00 aan de curator te betalen ter vergoeding van buitengerechtelijke kosten, te vermeerderen met de wettelijke rente als bedoeld in artikel 6:119 BW vanaf 16 februari 2026 tot de dag van volledige betaling,

6.4 verklaart deze veroordelingen uitvoerbaar bij voorraad,

6.5 verbindt aan de uitvoerbaarverklaring bij voorraad de voorwaarde dat de curator zekerheid stelt tot een bedrag van € 3.050,00 in de vorm van een bankgarantie,

6.6 compenseert de kosten van de procedure tussen partijen, in die zin dat iedere partij de eigen kosten draagt,

6.7 wijst het meer of anders gevorderde af.

Dit vonnis is gewezen door mr. H.J. Berends en in het openbaar uitgesproken op 14 juli 2026. (wv)

\begin{center}\rule{0.5\linewidth}{0.5pt}\end{center}

\newpage

\markboth{ECLI:NL:RBOVE:2026:4561}{mistral/mistral-medium-3.5-128b:fp8 (k = 0)}

\subsection{k = 0}\label{k-0-3}

\subsubsection{Mistral}\label{mistral-9}

\texttt{mistral/mistral-medium-3.5-128b:fp8}

\begin{center}
\small
\begin{tabular}{@{}p{7.4cm}p{2.6cm}p{2.6cm}@{}}
\toprule
Variabele & Waarde (model) & Waarde (gold) \\
\midrule
Is rechtshandeling & \texttt{true} & \texttt{true} \\
Datum rechtshandeling & \texttt{2023-10-23} & \texttt{2023-10-23} \\
Verricht door de schuldenaar & \texttt{true} & \texttt{true} \\
Om niet & \texttt{true} & \texttt{true} \\
Datum faillietverklaring & \texttt{2024-08-20} & \texttt{2024-08-20} \\
Onverplicht & \texttt{true} & \texttt{true} \\
Benadeling schuldeisers & \texttt{true} & \texttt{true} \\
Wetenschap schuldenaar aangetoond & \texttt{false} & \texttt{false} \\
Wetenschap wederpartij aangetoond & \texttt{false} & \texttt{null} \\
Vermoeden wetenschap schuldenaar tegenbewijs & \texttt{null} & \texttt{false} \\
Vermoeden wetenschap wederpartij tegenbewijs & \texttt{null} & \texttt{null} \\
Tegen onderwaarde & \texttt{null} & \texttt{null} \\
Voldoening of zekerheid niet opeisbare schuld & \texttt{null} & \texttt{null} \\
Gelieerde partij & \textcolor{gray}{\texttt{true}} & \texttt{null} \\
Reeds verplicht voor aanvang termijn & \texttt{null} & \texttt{null} \\
\textbf{Vernietigbaar} & \textbf{\texttt{true}} & \textbf{\texttt{true}} \\
\bottomrule
\end{tabular}
\end{center}

\paragraph{Is rechtshandeling - overwegingen}\label{is-rechtshandeling---overwegingen-27}

De overboeking van € 3.050,00 door {[}naam 1{]} aan {[}gedaagde{]} is een rechtshandeling, omdat het een op een rechtsgevolg gerichte wil (schenking) betreft die zich door een verklaring (overboeking met omschrijving `schenking') heeft geopenbaard.

\paragraph{Verricht door de schuldenaar - overwegingen}\label{verricht-door-de-schuldenaar---overwegingen-27}

De overboeking is verricht door {[}naam 1{]}, die later failliet is verklaard, en dus de schuldenaar is.

\paragraph{Om niet - overwegingen}\label{om-niet---overwegingen-27}

De overboeking is aangeduid als `schenking', wat per definitie een om niet-verrichting is (geen wederprestatie).

\paragraph{Onverplicht - overwegingen}\label{onverplicht---overwegingen-27}

De betaling van € 3.050,00 door {[}naam 1{]} aan {[}gedaagde{]} is aangeduid als `schenking' in de omschrijving van de overboeking. Een schenking is per definitie een onverplichte rechtshandeling, omdat er geen op de wet of een overeenkomst berustende rechtsplicht bestaat om een schenking te verrichten. De schuldenaar ({[}naam 1{]}) was niet verplicht om dit bedrag aan {[}gedaagde{]} te betalen, en de betaling kon derhalve niet in rechte worden afgedwongen. Morele, feitelijke of economische druk speelt hier geen rol, aangezien een schenking door haar aard vrijwillig en onverplicht is. De curator heeft derhalve voldoende aannemelijk gemaakt dat de betaling onverplicht is verricht in de zin van artikel 42 lid 1 Fw.

\paragraph{Benadeling schuldeisers - overwegingen}\label{benadeling-schuldeisers---overwegingen-27}

De schenking van € 3.050,00 door {[}naam 1{]} aan {[}gedaagde{]} op 23 oktober 2023 is een onverplichte rechtshandeling (art. 42 lid 1 Fw). Aangezien de schenking als zodanig is gekwalificeerd (omschrijving `schenking'), is er sprake van een vermogensvermindering van de boedel zonder tegenprestatie. Dit leidt tot een vermindering van het actief waar de schuldeisers zich op hadden kunnen verhalen. De benadeling is aanwezig op het moment dat de curator zijn rechten doet gelden (leer van de middellijke benadeling), omdat de boedel zonder deze schenking € 3.050,00 meer zou hebben bevatten. Er is geen aanwijzing dat de schuldeisers ook zonder deze handeling volledig zouden zijn voldaan of dat de schenking geen effect had op hun verhaalsmogelijkheden. Derhalve zijn de schuldeisers benadeeld in de zin van art. 42 Fw.

\paragraph{Wetenschap - overwegingen}\label{wetenschap---overwegingen-27}

De betaling van € 3.050,00 door {[}naam 1{]} aan {[}gedaagde{]} op 23 oktober 2023 is een schenking (rechtshandeling om niet) binnen één jaar voor de faillietverklaring (20 augustus 2024). Op grond van art. 45 Fw wordt vermoed dat {[}naam 1{]} wist of behoorde te weten dat benadeling van schuldeisers het gevolg zou zijn. Aangezien {[}gedaagde{]} de zoon is van {[}naam 1{]}, valt de schenking ook onder art. 43 lid 1 sub 3° Fw (bloedverwant tot in de derde graad), waardoor de wetenschap aan beide zijden wordt vermoed. Er is geen sprake van een verplichting vóór de jaartermijn, dus het vermoeden van art. 43 Fw is van toepassing. De curator hoeft derhalve geen rechtstreekse wetenschap aan te tonen, tenzij {[}gedaagde{]} succesvol tegenbewijs levert.

\newpage

\markboth{ECLI:NL:RBOVE:2026:4561}{gpt-5.5 (k = 0)}

\subsubsection{GPT}\label{gpt-9}

\texttt{gpt-5.5}

\begin{center}
\small
\begin{tabular}{@{}p{7.4cm}p{2.6cm}p{2.6cm}@{}}
\toprule
Variabele & Waarde (model) & Waarde (gold) \\
\midrule
Is rechtshandeling & \texttt{true} & \texttt{true} \\
Datum rechtshandeling & \texttt{2023-10-23} & \texttt{2023-10-23} \\
Verricht door de schuldenaar & \texttt{true} & \texttt{true} \\
Om niet & \texttt{true} & \texttt{true} \\
Datum faillietverklaring & \texttt{2024-08-20} & \texttt{2024-08-20} \\
Onverplicht & \texttt{true} & \texttt{true} \\
Benadeling schuldeisers & \texttt{true} & \texttt{true} \\
Wetenschap schuldenaar aangetoond & \texttt{false} & \texttt{false} \\
Wetenschap wederpartij aangetoond & \texttt{null} & \texttt{null} \\
Vermoeden wetenschap schuldenaar tegenbewijs & \texttt{false} & \texttt{false} \\
Vermoeden wetenschap wederpartij tegenbewijs & \texttt{false} & \texttt{null} \\
Tegen onderwaarde & \texttt{null} & \texttt{null} \\
Voldoening of zekerheid niet opeisbare schuld & \texttt{null} & \texttt{null} \\
Gelieerde partij & \textcolor{gray}{\texttt{true}} & \texttt{null} \\
Reeds verplicht voor aanvang termijn & \texttt{false} & \texttt{null} \\
\textbf{Vernietigbaar} & \textbf{\texttt{true}} & \textbf{\texttt{true}} \\
\bottomrule
\end{tabular}
\end{center}

\paragraph{Is rechtshandeling - overwegingen}\label{is-rechtshandeling---overwegingen-28}

De gewraakte handeling betreft de overmaking van € 3.050,00 door {[}naam 1{]} aan {[}gedaagde{]} met als omschrijving `schenking'. Dit duidt op een op rechtsgevolg gerichte vermogensverschuiving.

\paragraph{Verricht door de schuldenaar - overwegingen}\label{verricht-door-de-schuldenaar---overwegingen-28}

{[}naam 1{]}, de latere failliet verklaarde schuldenaar, heeft het bedrag zelf aan {[}gedaagde{]} overgemaakt.

\paragraph{Om niet - overwegingen}\label{om-niet---overwegingen-28}

Bij de overboeking staat als omschrijving `schenking' en uit de feiten blijkt niet van een bedongen wederprestatie door {[}gedaagde{]}.

\paragraph{Onverplicht - overwegingen}\label{onverplicht---overwegingen-28}

De aangevochten rechtshandeling betreft de overboeking van € 3.050,00 door {[}naam 1{]} aan {[}gedaagde{]} met als omschrijving `schenking'. Uit de casus blijkt niet van een wettelijke of contractuele verplichting van {[}naam 1{]} om dit bedrag aan {[}gedaagde{]} te betalen. Een schenking is naar haar aard een rechtshandeling om niet en berust, zonder nadere verplichtende rechtsgrond, niet op een in rechte afdwingbare betalingsplicht. Eventuele familiale, morele of feitelijke motieven zouden geen rechtsplicht opleveren. Op basis van de gegeven feiten was de betaling daarom onverplicht in de zin van artikel 42 Fw.

\paragraph{Benadeling schuldeisers - overwegingen}\label{benadeling-schuldeisers---overwegingen-28}

De betaling van € 3.050,00 door {[}naam 1{]} aan {[}gedaagde{]} is aangeduid als `schenking' en heeft ertoe geleid dat dit bedrag vóór faillissement uit het vermogen van {[}naam 1{]} is verdwenen. Bij de vermogensvergelijking moet de feitelijke situatie waarin de schenking in stand blijft worden vergeleken met de hypothetische situatie zonder die schenking. Zonder de betaling zou € 3.050,00 in beginsel tot het vermogen van {[}naam 1{]} hebben behoord en na faillissement beschikbaar zijn geweest voor de gezamenlijke schuldeisers. Als de schenking in stand blijft, ontbreekt dit bedrag in de boedel. Er zijn geen feiten gesteld waaruit volgt dat tegenover de betaling een boedelvoordeel stond, dat alle schuldeisers ook zonder dit bedrag volledig zouden worden voldaan, of dat schuldeisers zich niet op dit bedrag hadden kunnen verhalen. Ook de latere aankoop van de woning door {[}gedaagde{]} en zijn broer doet daaraan op basis van de gegeven feiten niet af. De rechtshandeling leidt daarom tot vermindering van het actief en daarmee tot benadeling van de schuldeisers in de zin van artikel 42 Fw.

\paragraph{Wetenschap - overwegingen}\label{wetenschap---overwegingen-28}

De betaling is aangeduid als `schenking' aan de zoon van de schuldenaar. Uit de gegeven feiten blijkt niet rechtstreeks dat {[}naam 1{]} ten tijde van de betaling wist of behoorde te weten dat faillissement en een tekort daarin met redelijke mate van waarschijnlijkheid waren te voorzien. Voor {[}gedaagde{]} geldt bij een rechtshandeling om niet bovendien geen zelfstandig wetenschapsvereiste van art. 42 lid 2 Fw. Wel speelt het bewijsvermoeden van art. 45 Fw voor de wetenschap van {[}naam 1{]} bij een schenking binnen één jaar vóór faillissement; op basis van de casus is geen geslaagd tegenbewijs gesteld of gebleken. Daarnaast is {[}gedaagde{]} als zoon een gelieerde partij in de zin van art. 43 lid 1 sub 3 Fw; ook daarvoor is geen eerdere verplichting of geslaagd tegenbewijs uit de feiten af te leiden.

\newpage

\markboth{ECLI:NL:RBOVE:2026:4561}{gemini-3.5-flash (k = 0)}

\subsubsection{Gemini}\label{gemini-9}

\texttt{gemini-3.5-flash}

\begin{center}
\small
\begin{tabular}{@{}p{7.4cm}p{2.6cm}p{2.6cm}@{}}
\toprule
Variabele & Waarde (model) & Waarde (gold) \\
\midrule
Is rechtshandeling & \texttt{true} & \texttt{true} \\
Datum rechtshandeling & \texttt{2023-10-23} & \texttt{2023-10-23} \\
Verricht door de schuldenaar & \texttt{true} & \texttt{true} \\
Om niet & \texttt{true} & \texttt{true} \\
Datum faillietverklaring & \texttt{2024-08-20} & \texttt{2024-08-20} \\
Onverplicht & \texttt{true} & \texttt{true} \\
Benadeling schuldeisers & \texttt{true} & \texttt{true} \\
Wetenschap schuldenaar aangetoond & \texttt{false} & \texttt{false} \\
Wetenschap wederpartij aangetoond & \texttt{null} & \texttt{null} \\
Vermoeden wetenschap schuldenaar tegenbewijs & \texttt{false} & \texttt{false} \\
Vermoeden wetenschap wederpartij tegenbewijs & \texttt{null} & \texttt{null} \\
Tegen onderwaarde & \texttt{null} & \texttt{null} \\
Voldoening of zekerheid niet opeisbare schuld & \texttt{null} & \texttt{null} \\
Gelieerde partij & \texttt{null} & \texttt{null} \\
Reeds verplicht voor aanvang termijn & \texttt{null} & \texttt{null} \\
\textbf{Vernietigbaar} & \textbf{\texttt{true}} & \textbf{\texttt{true}} \\
\bottomrule
\end{tabular}
\end{center}

\paragraph{Is rechtshandeling - overwegingen}\label{is-rechtshandeling---overwegingen-29}

De overboeking van een bedrag van € 3.050,00 met de omschrijving `schenking' betreft de uitvoering van een schenkingsovereenkomst. Dit is een rechtshandeling, aangezien deze is gericht op het teweegbrengen van een rechtsgevolg, namelijk de overdracht van het geldbedrag.

\paragraph{Verricht door de schuldenaar - overwegingen}\label{verricht-door-de-schuldenaar---overwegingen-29}

De betaling/overboeking is verricht door de schuldenaar zelf, de heer {[}naam 1{]}, aan zijn zoon {[}gedaagde{]}.

\paragraph{Om niet - overwegingen}\label{om-niet---overwegingen-29}

De overboeking is verricht met de omschrijving `schenking'. Er is geen sprake van een bedongen of geleverde tegenprestatie, waardoor de rechtshandeling om niet is verricht.

\paragraph{Onverplicht - overwegingen}\label{onverplicht---overwegingen-29}

De betaling van € 3.050,00 door de schuldenaar ({[}naam 1{]}) aan zijn zoon ({[}gedaagde{]}) is gedaan met de omschrijving `schenking'. Een schenking is een rechtshandeling om niet waarvoor geen voorafgaande, in rechte afdwingbare wettelijke of contractuele verplichting bestond. Nu er geen rechtsplicht was om dit bedrag aan de zoon over te maken, is de betaling/schenking onverplicht verricht in de zin van artikel 42 Fw.

\paragraph{Benadeling schuldeisers - overwegingen}\label{benadeling-schuldeisers---overwegingen-29}

Door de betaling van € 3.050,00 als schenking aan {[}gedaagde{]} is er een bedrag uit het vermogen van de schuldenaar verdwenen zonder dat daar enige tegenprestatie tegenover stond. Indien deze onverplichte rechtshandeling om niet achterwege was gebleven, zou het bedrag van € 3.050,00 in het vermogen van de schuldenaar zijn gebleven en na het faillissement in de faillissementsboedel zijn gevallen ten behoeve van de gezamenlijke schuldeisers. Door de schenking is het actief van de boedel rechtstreeks verminderd, waardoor de verhaalsmogelijkheden van de schuldeisers zijn beperkt. Er is derhalve sprake van benadeling van de schuldeisers in de zin van artikel 42 Fw.

\paragraph{Wetenschap - overwegingen}\label{wetenschap---overwegingen-29}

De betaling van € 3.050,00 door de schuldenaar aan zijn zoon is aangemerkt als een schenking, wat een rechtshandeling om niet betreft. Deze handeling vond plaats op 23 oktober 2023, binnen een jaar voor de faillietverklaring op 20 augustus 2024. Op grond van artikel 45 Fw wordt de wetenschap van benadeling aan de zijde van de schuldenaar (de vader) derhalve vermoed. Omdat het een rechtshandeling om niet betreft, is de wetenschap van de wederpartij (de zoon) niet vereist voor een geslaagd beroep op de pauliana (art. 42 lid 2 Fw). Het verweer van de gedaagde bevat geen feiten of omstandigheden die het wettelijke vermoeden van artikel 45 Fw ontzenuwen, waardoor het tegenbewijs niet slaagt en de wetenschap van de schuldenaar vaststaat.

\newpage

\markboth{ECLI:NL:RBOVE:2026:4561}{mistral/mistral-medium-3.5-128b:fp8 (k = 5)}

\subsection{k = 5}\label{k-5-3}

\subsubsection{Mistral}\label{mistral-10}

\texttt{mistral/mistral-medium-3.5-128b:fp8}

\begin{center}
\small
\begin{tabular}{@{}p{7.4cm}p{2.6cm}p{2.6cm}@{}}
\toprule
Variabele & Waarde (model) & Waarde (gold) \\
\midrule
Is rechtshandeling & \texttt{true} & \texttt{true} \\
Datum rechtshandeling & \texttt{2023-10-23} & \texttt{2023-10-23} \\
Verricht door de schuldenaar & \texttt{true} & \texttt{true} \\
Om niet & \texttt{true} & \texttt{true} \\
Datum faillietverklaring & \texttt{2024-08-20} & \texttt{2024-08-20} \\
Onverplicht & \texttt{true} & \texttt{true} \\
Benadeling schuldeisers & \texttt{true} & \texttt{true} \\
Wetenschap schuldenaar aangetoond & \textcolor{gray}{\texttt{true}} & \texttt{false} \\
Wetenschap wederpartij aangetoond & \textcolor{gray}{\texttt{true}} & \texttt{null} \\
Vermoeden wetenschap schuldenaar tegenbewijs & \texttt{null} & \texttt{false} \\
Vermoeden wetenschap wederpartij tegenbewijs & \texttt{null} & \texttt{null} \\
Tegen onderwaarde & \texttt{null} & \texttt{null} \\
Voldoening of zekerheid niet opeisbare schuld & \texttt{null} & \texttt{null} \\
Gelieerde partij & \textcolor{gray}{\texttt{true}} & \texttt{null} \\
Reeds verplicht voor aanvang termijn & \texttt{null} & \texttt{null} \\
\textbf{Vernietigbaar} & \textbf{\texttt{true}} & \textbf{\texttt{true}} \\
\bottomrule
\end{tabular}
\end{center}

\paragraph{Is rechtshandeling - overwegingen}\label{is-rechtshandeling---overwegingen-30}

De overboeking van € 3.050,00 door {[}naam 1{]} aan {[}gedaagde{]} is een rechtshandeling, omdat het een op een rechtsgevolg gerichte wil (schenking) betreft die zich door een verklaring (overboeking met omschrijving `schenking') heeft geopenbaard.

\paragraph{Verricht door de schuldenaar - overwegingen}\label{verricht-door-de-schuldenaar---overwegingen-30}

De overboeking is verricht door {[}naam 1{]}, die later failliet is verklaard, en dus de schuldenaar in deze context is.

\paragraph{Om niet - overwegingen}\label{om-niet---overwegingen-30}

De overboeking is aangeduid als `schenking', wat per definitie een om niet-verrichting is (geen wederprestatie).

\paragraph{Onverplicht - overwegingen}\label{onverplicht---overwegingen-30}

De betaling van € 3.050,00 door {[}naam 1{]} aan {[}gedaagde{]} is aangeduid als `schenking' in de overboekingsomschrijving. Een schenking is per definitie een rechtshandeling om niet, wat betekent dat er geen tegenprestatie of verplichting aan verbonden is. Volgens de jurisprudentie en de wettelijke definitie (art. 42 lid 1 Fw) is een schenking een onverplichte rechtshandeling, omdat er geen op de wet of een overeenkomst berustende rechtsplicht bestaat om deze te verrichten. De stelplicht en bewijslast van het onverplichte karakter rusten op de curator, maar in dit geval is de aard van de handeling (schenking) voldoende om aan te tonen dat deze onverplicht is. Zie ook ECLI:NL:GHSHE:2006:BC2959, waar het hof oordeelt dat een schenking als onverplichte rechtshandeling kan worden aangemerkt.

\paragraph{Benadeling schuldeisers - overwegingen}\label{benadeling-schuldeisers---overwegingen-30}

De schenking van € 3.050,00 door {[}naam 1{]} aan {[}gedaagde{]} op 23 oktober 2023 is een onverplichte rechtshandeling (art. 42 lid 1 Fw). Voor de beoordeling of de schuldeisers hierdoor zijn benadeeld, dient een vermogensvergelijking te worden gemaakt tussen de hypothetische situatie zonder de schenking en de feitelijke situatie met de schenking.

In de hypothetische situatie zou het bedrag van € 3.050,00 deel uitmaken van het vermogen van {[}naam 1{]} en derhalve beschikbaar zijn voor verhaalsmogelijkheden van de schuldeisers in het faillissement. In de feitelijke situatie is dit bedrag uit het vermogen van {[}naam 1{]} verdwenen en aan {[}gedaagde{]} toegekomen. Hierdoor is het actief van de boedel verminderd met € 3.050,00, wat direct leidt tot benadeling van de schuldeisers in het algemeen, aangezien zij zich op dit bedrag hadden kunnen verhalen.

De benadeling hoeft niet per se te leiden tot een tekort in de boedel; voldoende is dat de schuldeisers in een nadeliger positie verkeren dan zonder de rechtshandeling het geval zou zijn geweest (vgl. ECLI:NL:GHSHE:2006:BC2959, rov. 7.7 en ECLI:NL:GHAMS:2012:BW9644, rov. 3.4).

Daarnaast is de schenking een rechtshandeling \emph{om niet}, wat betekent dat voor de vernietiging op grond van art. 42 lid 1 Fw volstaat dat {[}naam 1{]} wist of behoorde te weten dat de schenking benadeling van schuldeisers tot gevolg zou hebben. De benadeling is hier evident, omdat het vermogen van {[}naam 1{]} direct is verminderd zonder tegenprestatie.

Gelet op het bovenstaande zijn de schuldeisers door de schenking benadeeld in de zin van art. 42 Fw.

\paragraph{Wetenschap - overwegingen}\label{wetenschap---overwegingen-30}

De schenking van € 3.050,00 door {[}naam 1{]} aan {[}gedaagde{]} (zijn kind) op 23 oktober 2023 is een rechtshandeling om niet, verricht binnen één jaar voor de faillietverklaring (20 augustus 2024). Op grond van art. 45 Fw wordt vermoed dat {[}naam 1{]} wist of behoorde te weten dat benadeling van schuldeisers het gevolg zou zijn. Dit vermoeden is van toepassing omdat de schenking binnen de jaartermijn valt en om niet is verricht. Aangezien {[}gedaagde{]} een bloedverwant in de eerste graad is van {[}naam 1{]}, valt de schenking ook onder art. 43 lid 1 sub 3° Fw, wat een bewijsvermoeden van wetenschap aan beide zijden meebrengt. Er is geen sprake van tegenbewijs (art. 43 lid 1 Fw) of een eerdere verplichting (art. 43 lid 1 Fw).

\newpage

\markboth{ECLI:NL:RBOVE:2026:4561}{gpt-5.5 (k = 5)}

\subsubsection{GPT}\label{gpt-10}

\texttt{gpt-5.5}

\begin{center}
\small
\begin{tabular}{@{}p{7.4cm}p{2.6cm}p{2.6cm}@{}}
\toprule
Variabele & Waarde (model) & Waarde (gold) \\
\midrule
Is rechtshandeling & \texttt{true} & \texttt{true} \\
Datum rechtshandeling & \texttt{2023-10-23} & \texttt{2023-10-23} \\
Verricht door de schuldenaar & \texttt{true} & \texttt{true} \\
Om niet & \texttt{true} & \texttt{true} \\
Datum faillietverklaring & \texttt{2024-08-20} & \texttt{2024-08-20} \\
Onverplicht & \texttt{true} & \texttt{true} \\
Benadeling schuldeisers & \texttt{true} & \texttt{true} \\
Wetenschap schuldenaar aangetoond & \texttt{false} & \texttt{false} \\
Wetenschap wederpartij aangetoond & \texttt{null} & \texttt{null} \\
Vermoeden wetenschap schuldenaar tegenbewijs & \texttt{false} & \texttt{false} \\
Vermoeden wetenschap wederpartij tegenbewijs & \texttt{false} & \texttt{null} \\
Tegen onderwaarde & \texttt{null} & \texttt{null} \\
Voldoening of zekerheid niet opeisbare schuld & \texttt{null} & \texttt{null} \\
Gelieerde partij & \textcolor{gray}{\texttt{true}} & \texttt{null} \\
Reeds verplicht voor aanvang termijn & \texttt{false} & \texttt{null} \\
\textbf{Vernietigbaar} & \textbf{\texttt{true}} & \textbf{\texttt{true}} \\
\bottomrule
\end{tabular}
\end{center}

\paragraph{Is rechtshandeling - overwegingen}\label{is-rechtshandeling---overwegingen-31}

De gewraakte handeling is de overboeking van € 3.050,00 met als omschrijving `schenking'. Dit duidt op een op rechtsgevolg gerichte vermogensverschuiving/schenking aan {[}gedaagde{]}.

\paragraph{Verricht door de schuldenaar - overwegingen}\label{verricht-door-de-schuldenaar---overwegingen-31}

De overboeking is verricht door {[}naam 1{]}, die later bij vonnis failliet is verklaard. De handeling is dus door de schuldenaar zelf verricht.

\paragraph{Om niet - overwegingen}\label{om-niet---overwegingen-31}

Bij de overboeking staat als omschrijving `schenking' en uit de feiten blijkt niet van een bedongen wederprestatie door {[}gedaagde{]}. De handeling is daarom om niet.

\paragraph{Onverplicht - overwegingen}\label{onverplicht---overwegingen-31}

De aangevochten rechtshandeling is de overboeking van € 3.050,00 door {[}naam 1{]} aan {[}gedaagde{]} op 23 oktober 2023. De betaling is in de omschrijving uitdrukkelijk aangeduid als `schenking'. Uit de casus blijkt niet van een wettelijke of contractuele verplichting van {[}naam 1{]} jegens {[}gedaagde{]} om dit bedrag te betalen, noch van een opeisbare schuld die met de betaling werd voldaan. Een schenking berust naar haar aard niet op een in rechte afdwingbare verplichting van de schenker. Ook als aan de betaling een familiaire, morele of feitelijke aanleiding ten grondslag zou hebben gelegen, levert dat geen rechtsplicht op in de zin van art. 42 Fw; vgl. ook ECLI:NL:GHSHE:2006:BC2959, waar voldoening aan een natuurlijke verbintenis het onverplichte karakter niet wegneemt. Op basis van de beschikbare feiten moet daarom worden aangenomen dat voor de betaling geen in rechte afdwingbare rechtsplicht bestond.

\paragraph{Benadeling schuldeisers - overwegingen}\label{benadeling-schuldeisers---overwegingen-31}

De betaling van € 3.050,00 is in de omschrijving aangeduid als `schenking' en uit de casus blijkt niet van enige tegenprestatie of andere vermogensvermeerdering voor {[}naam 1{]}. Bij de vermogensvergelijking leidt instandhouding van deze rechtshandeling ertoe dat € 3.050,00 aan liquide middelen buiten de boedel is geraakt, terwijl dit bedrag zonder de schenking in beginsel beschikbaar zou zijn geweest voor verhaal door de schuldeisers. Er zijn geen feiten gesteld waaruit volgt dat ook zonder deze betaling alle schuldeisers volledig zouden zijn voldaan, dat schuldeisers zich niet op dit bedrag hadden kunnen verhalen, of dat tegenover de betaling een boedelvoordeel stond. Daarmee is sprake van vermindering van het actief. Dit sluit aan bij de benadering in onder meer ECLI:NL:GHSHE:2006:BC2959, waarin een schenking benadeling oplevert doordat het vermogen van de schuldenaar met de waarde van het geschonkene is verminderd. De latere aankoop van de woning uit de boedel door {[}gedaagde{]} en zijn broer doet aan deze afzonderlijke vermogensonttrekking niet af.

\paragraph{Wetenschap - overwegingen}\label{wetenschap---overwegingen-31}

De overboeking is als `schenking' aangeduid en betreft daarmee een rechtshandeling om niet. Uit de casus volgen geen concrete feiten waaruit rechtstreeks blijkt dat {[}naam 1{]} ten tijde van de betaling wist of behoorde te weten dat faillissement en een tekort daarin met redelijke mate van waarschijnlijkheid te voorzien waren. Omdat de schenking binnen één jaar vóór faillietverklaring is verricht, geldt echter het bewijsvermoeden van art. 45 Fw aan de zijde van {[}naam 1{]}. Er zijn geen feiten gesteld waaruit volgt dat dit vermoeden met succes is ontzenuwd. Daarnaast is {[}gedaagde{]} de zoon van {[}naam 1{]}, zodat sprake is van een gelieerde partij in de zin van art. 43 lid 1 sub 3° Fw; er is geen aanwijzing dat {[}naam 1{]} zich al vóór de jaartermijn tot deze schenking had verplicht. Daarmee wordt wetenschap ook op grond van art. 43 Fw vermoed. Voor rechtstreekse wetenschap bij {[}gedaagde{]} zijn geen concrete feiten gegeven, maar bij een rechtshandeling om niet is die wetenschap voor vernietiging op grond van art. 42 lid 2 Fw niet constitutief vereist. Vergelijk de toepassing van het bewijsvermoeden bij schenkingen aan kinderen in ECLI:NL:GHSHE:2015:3242.

\newpage

\markboth{ECLI:NL:RBOVE:2026:4561}{gemini-3.5-flash (k = 5)}

\subsubsection{Gemini}\label{gemini-10}

\texttt{gemini-3.5-flash}

\begin{center}
\small
\begin{tabular}{@{}p{7.4cm}p{2.6cm}p{2.6cm}@{}}
\toprule
Variabele & Waarde (model) & Waarde (gold) \\
\midrule
Is rechtshandeling & \texttt{true} & \texttt{true} \\
Datum rechtshandeling & \texttt{2023-10-23} & \texttt{2023-10-23} \\
Verricht door de schuldenaar & \texttt{true} & \texttt{true} \\
Om niet & \texttt{true} & \texttt{true} \\
Datum faillietverklaring & \texttt{2024-08-20} & \texttt{2024-08-20} \\
Onverplicht & \texttt{true} & \texttt{true} \\
Benadeling schuldeisers & \texttt{true} & \texttt{true} \\
Wetenschap schuldenaar aangetoond & \texttt{false} & \texttt{false} \\
Wetenschap wederpartij aangetoond & \texttt{null} & \texttt{null} \\
Vermoeden wetenschap schuldenaar tegenbewijs & \texttt{false} & \texttt{false} \\
Vermoeden wetenschap wederpartij tegenbewijs & \texttt{null} & \texttt{null} \\
Tegen onderwaarde & \texttt{null} & \texttt{null} \\
Voldoening of zekerheid niet opeisbare schuld & \texttt{null} & \texttt{null} \\
Gelieerde partij & \texttt{null} & \texttt{null} \\
Reeds verplicht voor aanvang termijn & \texttt{null} & \texttt{null} \\
\textbf{Vernietigbaar} & \textbf{\texttt{true}} & \textbf{\texttt{true}} \\
\bottomrule
\end{tabular}
\end{center}

\paragraph{Is rechtshandeling - overwegingen}\label{is-rechtshandeling---overwegingen-32}

De overboeking van een bedrag van € 3.050,00 met de omschrijving `schenking' betreft de uitvoering van een schenkingsovereenkomst. Dit is een rechtshandeling, aangezien deze is gericht op het teweegbrengen van een rechtsgevolg, namelijk de overdracht van het geldbedrag.

\paragraph{Verricht door de schuldenaar - overwegingen}\label{verricht-door-de-schuldenaar---overwegingen-32}

De betaling/overboeking is verricht door de schuldenaar zelf, de heer {[}naam 1{]}, aan zijn zoon {[}gedaagde{]}.

\paragraph{Om niet - overwegingen}\label{om-niet---overwegingen-32}

De overboeking is verricht met de omschrijving `schenking'. Er is geen sprake van een bedongen of geleverde tegenprestatie, waardoor de rechtshandeling om niet is verricht.

\paragraph{Onverplicht - overwegingen}\label{onverplicht---overwegingen-32}

De rechtshandeling betreft een betaling van € 3.050,00 door de schuldenaar ({[}naam 1{]}) aan zijn zoon ({[}gedaagde{]}) met de omschrijving `schenking'. Een schenking is een rechtshandeling om niet. Er bestond voor de schuldenaar geen enkele op de wet of een voorafgaande overeenkomst berustende rechtsplicht om deze schenking te verrichten; de prestatie was derhalve niet in rechte afdwingbaar. Zelfs indien de schenking zou zijn verricht ter voldoening aan een morele plicht of natuurlijke verbintenis, ontneemt dit volgens vaste rechtspraak (zoals ECLI:NL:GHSHE:2006:BC2959) niet het onverplichte karakter aan de rechtshandeling. De betaling is dan ook onverplicht verricht.

\paragraph{Benadeling schuldeisers - overwegingen}\label{benadeling-schuldeisers---overwegingen-32}

De betaling van € 3.050,00 door {[}naam 1{]} aan {[}gedaagde{]} is verricht onder de omschrijving `schenking' en betreft derhalve een rechtshandeling om niet. Door deze schenking is het vermogen van de schuldenaar met dit bedrag verminderd zonder dat daar een tegenprestatie tegenover stond. Hierdoor is het voor verhaal beschikbare actief van de schuldenaar rechtstreeks verkleind. In de hypothetische situatie waarin de schenking niet zou zijn verricht, zou het bedrag van € 3.050,00 nog in het vermogen van de schuldenaar aanwezig zijn geweest, waardoor de schuldeisers zich daarop hadden kunnen verhalen. Overeenkomstig vaste rechtspraak (zoals ECLI:NL:GHSHE:2006:BC2959 en ECLI:NL:GHSHE:2015:3242) leidt een dergelijke onverplichte vermogensvermindering om niet tot benadeling van de schuldeisers in hun verhaalsmogelijkheden. Er is dan ook voldaan aan het vereiste van benadeling van schuldeisers in de zin van artikel 42 Faillissementswet.

\paragraph{Wetenschap - overwegingen}\label{wetenschap---overwegingen-32}

De rechtshandeling betreft een schenking van € 3.050,00 door de schuldenaar aan zijn zoon op 23 oktober 2023. Dit is een rechtshandeling om niet. De schuldenaar is op 20 augustus 2024 failliet verklaard, wat binnen één jaar na de schenking is. Op grond van artikel 45 Fw wordt daarom vermoed dat de schuldenaar wist of behoorde te weten dat benadeling van de schuldeisers het gevolg van de rechtshandeling zou zijn. Omdat het een rechtshandeling om niet betreft, is de wetenschap van de wederpartij (de zoon) niet vereist voor vernietiging op grond van artikel 42 lid 2 Fw. Er zijn geen feiten gesteld die het wettelijke vermoeden van artikel 45 Fw ontzenuwen, noch is de wetenschap van de schuldenaar rechtstreeks aangetoond.

\newpage

\markboth{ECLI:NL:RBOVE:2026:4561}{mistral/mistral-medium-3.5-128b:fp8 (k = 10)}

\subsection{k = 10}\label{k-10-3}

\subsubsection{Mistral}\label{mistral-11}

\texttt{mistral/mistral-medium-3.5-128b:fp8}

\begin{center}
\small
\begin{tabular}{@{}p{7.4cm}p{2.6cm}p{2.6cm}@{}}
\toprule
Variabele & Waarde (model) & Waarde (gold) \\
\midrule
Is rechtshandeling & \texttt{true} & \texttt{true} \\
Datum rechtshandeling & \texttt{2023-10-23} & \texttt{2023-10-23} \\
Verricht door de schuldenaar & \texttt{true} & \texttt{true} \\
Om niet & \texttt{true} & \texttt{true} \\
Datum faillietverklaring & \texttt{2024-08-20} & \texttt{2024-08-20} \\
Onverplicht & \texttt{true} & \texttt{true} \\
Benadeling schuldeisers & \texttt{true} & \texttt{true} \\
Wetenschap schuldenaar aangetoond & \textcolor{gray}{\texttt{true}} & \texttt{false} \\
Wetenschap wederpartij aangetoond & \textcolor{gray}{\texttt{true}} & \texttt{null} \\
Vermoeden wetenschap schuldenaar tegenbewijs & \texttt{null} & \texttt{false} \\
Vermoeden wetenschap wederpartij tegenbewijs & \texttt{null} & \texttt{null} \\
Tegen onderwaarde & \texttt{null} & \texttt{null} \\
Voldoening of zekerheid niet opeisbare schuld & \texttt{null} & \texttt{null} \\
Gelieerde partij & \textcolor{gray}{\texttt{true}} & \texttt{null} \\
Reeds verplicht voor aanvang termijn & \texttt{null} & \texttt{null} \\
\textbf{Vernietigbaar} & \textbf{\texttt{true}} & \textbf{\texttt{true}} \\
\bottomrule
\end{tabular}
\end{center}

\paragraph{Is rechtshandeling - overwegingen}\label{is-rechtshandeling---overwegingen-33}

De overboeking van € 3.050,00 door {[}naam 1{]} aan {[}gedaagde{]} is een rechtshandeling, omdat het een op een rechtsgevolg gerichte wil (schenking) betreft die zich door een verklaring (overboeking met omschrijving `schenking') heeft geopenbaard.

\paragraph{Verricht door de schuldenaar - overwegingen}\label{verricht-door-de-schuldenaar---overwegingen-33}

De overboeking is verricht door {[}naam 1{]}, die later failliet is verklaard, en dus de schuldenaar is.

\paragraph{Om niet - overwegingen}\label{om-niet---overwegingen-33}

De overboeking is aangeduid als `schenking', wat per definitie een om niet-verrichting is (geen wederprestatie).

\paragraph{Onverplicht - overwegingen}\label{onverplicht---overwegingen-33}

De betaling van € 3.050,00 door {[}naam 1{]} aan {[}gedaagde{]} is aangeduid als `schenking' in de overboekingsomschrijving. Een schenking is per definitie een rechtshandeling om niet, wat betekent dat er geen tegenprestatie tegenover staat. Omdat er geen wettelijke of contractuele verplichting bestond voor {[}naam 1{]} om deze betaling te verrichten, is de handeling onverplicht in de zin van artikel 42 lid 1 Fw. Morele, feitelijke of economische druk (zoals een eventuele intentie om de zoon te begunsten) schept geen rechtsplicht. De stelplicht en bewijslast van het onverplichte karakter rusten op de curator, en in dit geval is de aanduiding `schenking' voldoende om aan te tonen dat er geen verplichting bestond.

\paragraph{Benadeling schuldeisers - overwegingen}\label{benadeling-schuldeisers---overwegingen-33}

De schenking van € 3.050,00 door {[}naam 1{]} aan {[}gedaagde{]} op 23 oktober 2023 is een onverplichte rechtshandeling (art. 42 lid 1 Fw). Door deze schenking is het vermogen van {[}naam 1{]} verminderd met € 3.050,00, wat in de hypothetische situatie (zonder schenking) ten bate van de schuldeisers had kunnen komen. Aangezien {[}naam 1{]} op dat moment al in financiële moeilijkheden verkeerde (gelet op het latere faillissement op 20 augustus 2024), is aannemelijk dat de schuldeisers hierdoor zijn benadeeld. De benadeling is aanwezig op het moment dat de curator zijn rechten doet gelden (leer van de middellijke benadeling).

\paragraph{Wetenschap - overwegingen}\label{wetenschap---overwegingen-33}

De schenking van € 3.050,00 door {[}naam 1{]} aan {[}gedaagde{]} op 23 oktober 2023 is een rechtshandeling om niet. Omdat de schenking binnen één jaar voor de faillietverklaring (20 augustus 2024) is verricht, wordt op grond van art. 45 Fw vermoed dat {[}naam 1{]} wist of behoorde te weten dat benadeling van de schuldeisers het gevolg zou zijn. {[}gedaagde{]} is de zoon van {[}naam 1{]}, waardoor de schenking valt onder art. 43 lid 1 sub 3° Fw (gelieerde partij). Dit brengt mee dat ook de wetenschap van {[}gedaagde{]} wordt vermoed. Er is geen sprake van tegenbewijs, omdat niet is gesteld of gebleken dat {[}naam 1{]} of {[}gedaagde{]} aannemelijk hebben gemaakt dat het faillissement en een tekort daarin niet met een redelijke mate van waarschijnlijkheid waren te voorzien.

\newpage

\markboth{ECLI:NL:RBOVE:2026:4561}{gpt-5.5 (k = 10)}

\subsubsection{GPT}\label{gpt-11}

\texttt{gpt-5.5}

\begin{center}
\small
\begin{tabular}{@{}p{7.4cm}p{2.6cm}p{2.6cm}@{}}
\toprule
Variabele & Waarde (model) & Waarde (gold) \\
\midrule
Is rechtshandeling & \texttt{true} & \texttt{true} \\
Datum rechtshandeling & \texttt{2023-10-23} & \texttt{2023-10-23} \\
Verricht door de schuldenaar & \texttt{true} & \texttt{true} \\
Om niet & \texttt{true} & \texttt{true} \\
Datum faillietverklaring & \texttt{2024-08-20} & \texttt{2024-08-20} \\
Onverplicht & \texttt{true} & \texttt{true} \\
Benadeling schuldeisers & \texttt{true} & \texttt{true} \\
Wetenschap schuldenaar aangetoond & \texttt{false} & \texttt{false} \\
Wetenschap wederpartij aangetoond & \texttt{null} & \texttt{null} \\
Vermoeden wetenschap schuldenaar tegenbewijs & \texttt{false} & \texttt{false} \\
Vermoeden wetenschap wederpartij tegenbewijs & \texttt{false} & \texttt{null} \\
Tegen onderwaarde & \texttt{null} & \texttt{null} \\
Voldoening of zekerheid niet opeisbare schuld & \texttt{null} & \texttt{null} \\
Gelieerde partij & \textcolor{gray}{\texttt{true}} & \texttt{null} \\
Reeds verplicht voor aanvang termijn & \texttt{false} & \texttt{null} \\
\textbf{Vernietigbaar} & \textbf{\texttt{true}} & \textbf{\texttt{true}} \\
\bottomrule
\end{tabular}
\end{center}

\paragraph{Is rechtshandeling - overwegingen}\label{is-rechtshandeling---overwegingen-34}

De gewraakte handeling is de overboeking van € 3.050,00 door {[}naam 1{]} aan {[}gedaagde{]} met als omschrijving `schenking'. Een schenking/overboeking is gericht op rechtsgevolg, namelijk vermogensverschuiving ten gunste van {[}gedaagde{]}.

\paragraph{Verricht door de schuldenaar - overwegingen}\label{verricht-door-de-schuldenaar---overwegingen-34}

De betaling is verricht door {[}naam 1{]}, die later bij vonnis van 20 augustus 2024 failliet is verklaard. Daarmee is de handeling door de schuldenaar zelf verricht.

\paragraph{Om niet - overwegingen}\label{om-niet---overwegingen-34}

Bij de overboeking staat als omschrijving `schenking'. Uit de feiten blijkt niet van enige wederprestatie door {[}gedaagde{]}. De handeling moet daarom als om niet worden aangemerkt.

\paragraph{Onverplicht - overwegingen}\label{onverplicht---overwegingen-34}

De aangevochten rechtshandeling is de overboeking van € 3.050,00 door {[}naam 1{]} aan {[}gedaagde{]} op 23 oktober 2023, met als omschrijving `schenking'. Uit de casus blijkt niet van een op de wet of overeenkomst berustende, in rechte afdwingbare verplichting van {[}naam 1{]} om dit bedrag aan {[}gedaagde{]} te betalen. Een schenking is naar haar aard een bevoordeling om niet en is, behoudens een vooraf bestaande afdwingbare verplichting tot schenking of betaling, onverplicht. Van zo'n verplichting is niets gesteld of gebleken. Eventuele familiale, morele of feitelijke redenen om het bedrag aan een kind over te maken scheppen geen rechtsplicht in de zin van artikel 42 Fw. Dit sluit aan bij de lijn dat schenkingen aan kinderen in beginsel als onverplichte rechtshandelingen worden aangemerkt, vgl. onder meer ECLI:NL:GHSHE:2006:BC2959 en ECLI:NL:RBGEL:2018:1312.

\paragraph{Benadeling schuldeisers - overwegingen}\label{benadeling-schuldeisers---overwegingen-34}

De aangevochten rechtshandeling is de overboeking van € 3.050,00 door {[}naam 1{]} aan {[}gedaagde{]} met als omschrijving `schenking'. Bij de benadelingsbeoordeling moet de feitelijke situatie waarin de schenking in stand blijft worden vergeleken met de hypothetische situatie zonder die schenking, beoordeeld naar het moment waarop de curator zijn rechten doet gelden. Zonder de schenking zou € 3.050,00 in het vermogen van {[}naam 1{]} zijn gebleven en daarmee in beginsel beschikbaar zijn geweest voor verhaal door diens schuldeisers/voor de faillissementsboedel. Door de schenking is het actief van {[}naam 1{]} met dat bedrag verminderd, zonder dat daar een tegenprestatie tegenover stond. Er zijn geen feiten gesteld waaruit volgt dat de schuldeisers ook zonder dit bedrag volledig zouden zijn voldaan, dat zij zich niet op dit geld hadden kunnen verhalen, of dat de schenking per saldo een gelijkwaardig boedelvoordeel opleverde. Dat de curator ook de woning in de boedel heeft aangetroffen en deze later heeft verkocht, maakt dit niet anders zolang niet blijkt dat daardoor alle schuldeisers volledig kunnen worden voldaan. Vergelijkbaar is in de rechtspraak aangenomen dat een schenking van een vermogensbestanddeel tot benadeling leidt doordat het vermogen van de schuldenaar met de waarde daarvan vermindert (vgl. ECLI:NL:GHSHE:2006:BC2959 en ECLI:NL:GHSHE:2015:3242).

\paragraph{Wetenschap - overwegingen}\label{wetenschap---overwegingen-34}

De betaling is aangeduid als `schenking' aan de zoon van de schuldenaar. Op basis van de beperkte feiten is niet rechtstreeks aangetoond dat {[}naam 1{]} ten tijde van de betaling wist of behoorde te weten dat faillissement en een tekort daarin met redelijke mate van waarschijnlijkheid waren te voorzien; het latere faillissement op zichzelf volstaat daarvoor niet. Wel ligt bij een schenking binnen één jaar het bewijsvermoeden van art. 45 Fw voor de hand: wetenschap van benadeling aan de zijde van de schuldenaar wordt dan vermoed. Er zijn geen concrete feiten gesteld waaruit volgt dat dit vermoeden met succes is ontzenuwd. Daarnaast is sprake van een gelieerde partij in de zin van art. 43 lid 1 sub 3° Fw, omdat de schenking is verricht aan een kind; voor zover art. 43 Fw wordt toegepast, is geen tegenbewijs gebleken. Vergelijkbaar volgt uit onder meer ECLI:NL:RBGEL:2018:1312 en ECLI:NL:GHSHE:2015:3242 dat bij schenkingen aan kinderen binnen de jaartermijn het vermoeden zwaar weegt en concrete feiten nodig zijn om dit te ontzenuwen.

\newpage

\markboth{ECLI:NL:RBOVE:2026:4561}{gemini-3.5-flash (k = 10)}

\subsubsection{Gemini}\label{gemini-11}

\texttt{gemini-3.5-flash}

\begin{center}
\small
\begin{tabular}{@{}p{7.4cm}p{2.6cm}p{2.6cm}@{}}
\toprule
Variabele & Waarde (model) & Waarde (gold) \\
\midrule
Is rechtshandeling & \texttt{true} & \texttt{true} \\
Datum rechtshandeling & \texttt{2023-10-23} & \texttt{2023-10-23} \\
Verricht door de schuldenaar & \texttt{true} & \texttt{true} \\
Om niet & \texttt{true} & \texttt{true} \\
Datum faillietverklaring & \texttt{2024-08-20} & \texttt{2024-08-20} \\
Onverplicht & \texttt{true} & \texttt{true} \\
Benadeling schuldeisers & \texttt{true} & \texttt{true} \\
Wetenschap schuldenaar aangetoond & \texttt{false} & \texttt{false} \\
Wetenschap wederpartij aangetoond & \texttt{null} & \texttt{null} \\
Vermoeden wetenschap schuldenaar tegenbewijs & \texttt{false} & \texttt{false} \\
Vermoeden wetenschap wederpartij tegenbewijs & \texttt{null} & \texttt{null} \\
Tegen onderwaarde & \texttt{null} & \texttt{null} \\
Voldoening of zekerheid niet opeisbare schuld & \texttt{null} & \texttt{null} \\
Gelieerde partij & \texttt{null} & \texttt{null} \\
Reeds verplicht voor aanvang termijn & \texttt{null} & \texttt{null} \\
\textbf{Vernietigbaar} & \textbf{\texttt{true}} & \textbf{\texttt{true}} \\
\bottomrule
\end{tabular}
\end{center}

\paragraph{Is rechtshandeling - overwegingen}\label{is-rechtshandeling---overwegingen-35}

De overboeking van een geldbedrag van € 3.050,00 met de omschrijving `schenking' betreft een rechtshandeling (schenking en de goederenrechtelijke levering van het geld), aangezien deze gericht is op het rechtsgevolg van eigendomsoverdracht.

\paragraph{Verricht door de schuldenaar - overwegingen}\label{verricht-door-de-schuldenaar---overwegingen-35}

De overboeking is verricht door de schuldenaar zelf, te weten de heer {[}naam 1{]}, voorafgaand aan zijn faillissement.

\paragraph{Om niet - overwegingen}\label{om-niet---overwegingen-35}

De rechtshandeling is aangemerkt als een `schenking', waarbij geen sprake is van een tegenprestatie. De handeling is derhalve om niet verricht.

\paragraph{Onverplicht - overwegingen}\label{onverplicht---overwegingen-35}

De betaling van € 3.050,00 door de schuldenaar ({[}naam 1{]}) aan zijn zoon ({[}gedaagde{]}) is verricht onder de omschrijving `schenking'. Een schenking is een rechtshandeling om niet. Er bestond voor de schuldenaar geen enkele op de wet of een overeenkomst berustende rechtsplicht om dit bedrag aan gedaagde te betalen. De rechtshandeling kon derhalve niet in rechte worden afgedwongen. Zelfs indien er sprake zou zijn van het voldoen aan een morele plicht of natuurlijke verbintenis, neemt dit het onverplichte karakter van de rechtshandeling niet weg (overeenkomstig ECLI:NL:GHSHE:2006:BC2959). De schenking is dan ook een onverplichte rechtshandeling in de zin van artikel 42 Faillissementswet.

\paragraph{Benadeling schuldeisers - overwegingen}\label{benadeling-schuldeisers---overwegingen-35}

De overboeking van € 3.050,00 door {[}naam 1{]} aan {[}gedaagde{]} is verricht onder de omschrijving `schenking' en kwalificeert derhalve als een rechtshandeling om niet. Om te bepalen of er sprake is van benadeling van de schuldeisers in de zin van artikel 42 Fw, dient een vermogensvergelijking te worden gemaakt tussen de hypothetische situatie waarin de schuldeisers zouden hebben verkeerd zonder de schenking, en de feitelijke situatie waarin zij verkeren nu de schenking heeft plaatsgevonden. Door de schenking is het actieve vermogen van de schuldenaar met € 3.050,00 verminderd zonder dat daar enige tegenprestatie tegenover stond. Hierdoor is het voor verhaal beschikbare actief van de schuldenaar rechtstreeks afgenomen, wat leidt tot een verslechtering van de verhaalsmogelijkheden van de schuldeisers. In lijn met vaste rechtspraak (zoals ECLI:NL:GHSHE:2006:BC2959, r.o. 7.7 en ECLI:NL:RBGEL:2022:6533, r.o. 4.20) brengt een dergelijke onverplichte vermogensvermindering om niet benadeling van de schuldeisers met zich mee, aangezien het actieve boedelvolume hiermee is verkleind.

\paragraph{Wetenschap - overwegingen}\label{wetenschap---overwegingen-35}

De betaling van € 3.050,00 op 23 oktober 2023 is aangemerkt als een schenking, wat een rechtshandeling om niet is. Omdat deze rechtshandeling is verricht binnen één jaar vóór de faillietverklaring van de schuldenaar (20 augustus 2024), is het wettelijk bewijsvermoeden van artikel 45 Fw van toepassing. Hierdoor wordt vermoed dat de schuldenaar wist of behoorde te weten dat benadeling van de schuldeisers het gevolg zou zijn. Er zijn geen feiten gesteld die dit vermoeden ontzenuwen, waardoor het tegenbewijs niet slaagt. Omdat het een rechtshandeling om niet betreft, is de wetenschap van de wederpartij (de gedaagde) niet vereist voor vernietiging onder artikel 42 Fw. De bewijsvermoedens van artikel 43 Fw zijn hier niet aan de orde.

\newpage

\markboth{ECLI:NL:RBZWB:2026:4349}{Gold standard}

\section{ECLI:NL:RBZWB:2026:4349}\label{eclinlrbzwb20264349}

uitspraakdatum 20-05-2026

\subsection{Gold}\label{gold-4}

\begin{center}
\small
\begin{tabular}{@{}p{8cm}p{2.6cm}@{}}
\toprule
Variabele & Waarde \\
\midrule
Is rechtshandeling & \texttt{true} \\
Datum rechtshandeling & \texttt{2019-08-01} \\
Verricht door de schuldenaar & \texttt{true} \\
Om niet & \texttt{false} \\
Datum faillietverklaring & \texttt{2020-08-04} \\
Onverplicht & \texttt{true} \\
Benadeling schuldeisers & \texttt{false} \\
Wetenschap schuldenaar aangetoond & \texttt{false} \\
Wetenschap wederpartij aangetoond & \texttt{false} \\
Vermoeden wetenschap schuldenaar tegenbewijs & \texttt{null} \\
Vermoeden wetenschap wederpartij tegenbewijs & \texttt{null} \\
Tegen onderwaarde & \texttt{null} \\
Voldoening of zekerheid niet opeisbare schuld & \texttt{null} \\
Gelieerde partij & \texttt{null} \\
Reeds verplicht voor aanvang termijn & \texttt{null} \\
\textbf{Vernietigbaar} & \textbf{\texttt{false}} \\
\bottomrule
\end{tabular}
\end{center}

\subsubsection{Uitspraak}\label{uitspraak-4}

\paragraph{RECHTBANK Zeeland-West-Brabant}\label{rechtbank-zeeland-west-brabant}

Civiel recht

Zittingsplaats Breda

Zaaknummer: C/02/431468 / HA ZA 25-66

\textbf{Vonnis van 20 mei 2026}

in de zaak van

\textbf{MR. {[}curator{]} IN ZIJN HOEDANIGHEID VAN CURATOR VAN THINK BUILDING CONCEPTS B.V.},

te Tilburg,

eisende partij,

hierna te noemen: de curator,

advocaat: mr. {[}curator{]} ,

tegen

\textbf{PROFUNDO B.V.},

te Fijnaart,

gedaagde partij,

hierna te noemen: Profundo,

advocaat: mr. M. Littooij.

\paragraph{1 De zaak in het kort}\label{de-zaak-in-het-kort-1}

1.1. Concepts heeft op 1 augustus 2019 een betaling gedaan aan Profundo. Concepts is op 4 augustus 2020 failliet verklaard met de benoeming van de curator tot curator. De curator vordert in deze procedure terugbetaling van het betaalde bedrag op grond van de faillissementspauliana van artikel 42 Fw. Profundo voert verweer. Aan de vereisten voor een geslaagd beroep op de faillissementspauliana wordt volgens haar niet voldaan.

1.2. De rechtbank wijst de vorderingen van de curator af. Hierna wordt onder het kopje `de beoordeling' uitgelegd waarom. Eerst schetst de rechtbank het verloop van de procedure, de feiten en het geschil. Tot slot volgt de beslissing.

\paragraph{2 De procedure}\label{de-procedure-1}

2.1. Het verloop van de procedure blijkt uit:

\begin{itemize}
\item
  het tussenvonnis van 11 juni 2025 en de daarin genoemde processtukken,
\item
  de akte overlegging producties van de curator met daarbij producties 15 tot en met 31,
\item
  de aanvullende productie 9 van Profundo,
\item
  de mondelinge behandeling van 5 februari 2026, waarvan door de griffier aantekeningen zijn gemaakt en waarvan naderhand een proces-verbaal is opgemaakt,
\item
  de spreekaantekeningen van de curator en van mr. Littooij, zoals zij die tijdens de mondelinge behandeling hebben voorgelezen.
\end{itemize}

2.2. De curator heeft tijdens de mondelinge behandeling aangegeven een beroep te doen op het proces-verbaal van de mondelinge behandeling van 29 januari 2026 in de aanpalende procedure tussen Think Building \& Construction B.V. (hierna: `Construction') en Thetis Real Estate B.V. (hierna: `Thetis') met zaaknummer C/02/431465/HA ZA 25-65. De rechtbank overweegt dat een proces-verbaal van een andere procedure geen onderdeel uitmaakt van het procesdossier in deze zaak. Daarbij merkt de rechtbank op dat het betreffende proces-verbaal op het moment van de mondelinge behandeling in deze zaak ook nog niet was opgemaakt.

2.3. Vervolgens heeft de rechtbank besloten om in deze zaak vonnis te wijzen.

\paragraph{3 De feiten}\label{de-feiten-1}

3.1. Concepts is op 4 augustus 2020 op eigen verzoek failliet verklaard door de rechtbank Oost-Brabant, met benoeming van de curator tot curator.

3.2. Concepts behoorde tot een groep vennootschappen, waarvan de heer {[}naam 1{]} via zijn persoonlijke holding {[}B.V.{]} (hierna: ` {[}naam 1{]} Holding') grootaandeelhouder is. Statutair bestuurder van Concepts was de besloten vennootschap Marcha Beheer B.V., op haar beurt vertegenwoordigd door haar bestuurder de heer {[}naam 2{]} (hierna: ` {[}naam 2{]} '). Enig aandeelhouder van Concepts was de besloten vennootschap Think Holding B.V.

3.3. Profundo behoort tot dezelfde groep vennootschappen als Concepts. Ook Construction maakt onderdeel uit van deze groep. {[}naam 1{]} Holding is statutair bestuurder van Profundo.

3.4. Concepts hield zich bezig met het ontwikkelen en uitvoeren van bouwprojecten.

3.5. Concepts en Profundo hebben op 13 juli 2017 een overeenkomst van geldlening gesloten (hierna: `de overeenkomst'). In de considerans van de overeenkomst is overwogen dat tussen partijen een rekening-courantverhouding bestaat, op basis waarvan Concepts een bedrag van € 1.285.000,00 aan Profundo verschuldigd was. Partijen wilden deze rekening-courantverhouding omzetten in een langlopende lening. De overeenkomst legt de voorwaarden van de langlopende lening vast.

3.6. In artikel 2 van de overeenkomst is bepaald dat aan de looptijd van de lening geen termijn is verbonden. Profundo kan de overeenkomst met inachtneming van een opzegtermijn van 6 maanden opzeggen. De overeenkomst bevat geen aflossingstermijnen.

3.7. In artikel 4 van de overeenkomst is bepaald:

\emph{``De leningstand (inclusief de lopende rente) kan dadelijk worden opgeëist, zonder dat enige sommatie of ingebrekestelling of gerechtelijke tussenkomst zal zijn vereist, indien:}

\begin{enumerate}
\def\labelenumi{\alph{enumi}.}
\tightlist
\item
  \emph{een Schuldenaar surséance van betaling aanvraagt of failliet wordt verklaard;}
\item
  \emph{een Schuldenaar een akkoord wordt aangeboden buiten faillissement;}
\item
  \emph{een Schuldenaar, na ingebrekestelling, nalatig is enige verplichting, voortkomende uit deze overeenkomst, na te komen;}
\item
  \emph{enige maatregel van conservatoire of executoriale aard jegens een Schuldenaar genomen wordt;''}
\end{enumerate}

3.8. Uit de jaarrekening 2017 blijkt dat de langlopende schuld van Concepts aan Profundo uit hoofde van de overeenkomst per 31 december 2017 € 1.307.329,00 bedroeg.

3.9. Think Holding B.V. heeft in hoedanigheid van houder van de aandelen in Construction jegens Woningborg N.V. een dividendverklaring afgegeven, ondertekend door {[}naam 1{]} . In deze dividendverklaring is onder meer de garantie opgenomen dat geen uitkering van dividend of uitkering ten laste van de agioreserve door Construction aan de aandeelhouder zal worden gedaan indien het aansprakelijk vermogen van Construction door een dergelijke uitkering onder een bedrag van € 600.000,00 zou uitkomen.

3.10. Concepts heeft op 1 augustus 2019 een bedrag van € 552.500,00 betaald aan Profundo. Deze betaling heeft plaatsgevonden in het kader van een kasronde binnen de groep op die datum. De kasronde was als volgt:

\begin{enumerate}
\def\labelenumi{\arabic{enumi}.}
\tightlist
\item
  Profundo leent € 670.000,00 uit aan Think Holding B.V.;
\item
  Think Holding B.V. betaalt € 670.000,00 als aflossing van haar openstaande schuld aan Construction;
\item
  Construction betaalt Concepts € 552.500,00 voor aflossing van haar rekening-courantschuld;
\item
  Construction betaalt Thetis € 117.374,00 terug voor de in rekening-courant geboekte legeskosten die Thetis op 19 juli 2017 had betaald;
\item
  Concepts betaalt Profundo € 552.500,00 terug op de lening.
\end{enumerate}

3.11. De curator heeft de betaling van Concepts aan Profundo van € 552.500,00 op 26 juli 2024 met een beroep op de faillissementspauliana van artikel 42 Fw buitengerechtelijk vernietigd. De curator heeft Profundo gesommeerd om het bedrag uiterlijk 15 dagen na ontvangst van de brief terug te betalen.

3.12. Mr.~Froger, kantoorgenoot van mr. Littooij, reageerde hierop namens Profundo per e-mail van 5 augustus 2024 als volgt:

\emph{``Thetis en Profundo betwisten de vorderingen die uw cliënt stelt te hebben: de motivering daarvan is hem al lang bekend; daar reageert u niet op. Er is uw cliënt meermaals uitleg gegeven over de kasronde ter versteviging van het eigen vermogen en de achterliggende reden. Dit gebeurde allemaal op initiatief van de bestuurder, de heer {[}naam 2{]} . Een weerlegging hiervan heeft uw cliënt nooit verwoord. (\ldots)''}

3.13. Profundo heeft het bedrag van € 552.500,00 niet terugbetaald aan Concepts.

\paragraph{4 Het geschil}\label{het-geschil-1}

4.1. De curator vordert -- samengevat -- dat de rechtbank bij uitvoerbaar bij voorraad te verklaren vonnis:

I. voor recht verklaart dat de curator op juiste gronden tot vernietiging is overgegaan ex artikel 42 Fw van de betaling van € 552.500,00 en Profundo veroordeelt tot terugbetaling aan Concepts van dit bedrag wegens onverschuldigde betaling,

\begin{enumerate}
\def\labelenumi{\Roman{enumi}.}
\setcounter{enumi}{1}
\item
  Profundo veroordeelt tot betaling van de wettelijke rente over de gevorderde hoofdsom vanaf 11 augustus 2024 tot de dag van volledige betaling, te vermeerderen met incassokosten,
\item
  Profundo veroordeelt in de proceskosten.
\end{enumerate}

4.2. Profundo voert verweer. Profundo concludeert tot het ontzeggen dan wel afwijzen van de vorderingen van de curator, met uitvoerbaar bij voorraad te verklaren veroordeling van de curator in de kosten van deze procedure.

4.3. De rechtbank gaat hierna bij de beoordeling in op de relevante standpunten die partijen ter onderbouwing van de vorderingen en het verweer daartegen hebben aangevoerd.

\paragraph{5 De beoordeling}\label{de-beoordeling-1}

5.1. De curator doet ter onderbouwing van zijn vordering een beroep op artikel 42 Fw, de faillissementspauliana. Volgens Profundo wordt aan de voorwaarden voor een geslaagd beroep op de faillissementspauliana niet voldaan.

Maatstaf voor de beoordeling

\emph{Voorwaarden voor een geslaagd beroep op de faillissementspauliana}

5.2. Op grond van artikel 42 lid 1 jo lid 2 Fw kan een curator een rechtshandeling anders dan om niet die de failliet vóór de faillietverklaring onverplicht heeft verricht, vernietigen, indien zowel de failliet als degenen met of jegens wie de schuldenaar de rechtshandeling verrichtte, bij dit verrichten wisten of behoorden te weten dat daarvan benadeling van de schuldeisers het gevolg zou zijn. Uit dit artikel volgt dus dat voor een geslaagd beroep op faillissementspauliana vereist is dat (1) sprake is van een onverplichte rechtshandeling, (2) waardoor de schuldeisers zijn benadeeld en (3) de schuldenaar en de wederpartij wetenschap hadden van die benadeling. Uitgangspunt is dat het aan de curator is om voldoende te stellen en zo nodig te bewijzen dat aan de vereisten van faillissementspauliana is voldaan.

\emph{Samenhangend geheel van rechtshandelingen}

5.3. Artikel 42 Fw neemt de enkelvoudige rechtshandeling tot uitgangspunt. De omstandigheden van het geval kunnen echter meebrengen dat onder de reikwijdte van artikel 42 Fw niet alleen een enkelvoudige rechtshandeling kan worden beoordeeld, maar ook een samenhangend geheel van rechtshandelingen. Rechtshandelingen kunnen een zodanig samenhangend geheel vormen dat de (nadelige) gevolgen ervan in onderling verband moeten worden beoordeeld. Of sprake is van samenhangende rechtshandelingen hangt onder meer af van de inhoud van de betrokken rechtshandelingen, de onderlinge afstemming daarvan en de samenhang tussen die handelingen wat betreft het moment waarop zij tot stand zijn gekomen. Bepalend is of daaruit volgt dat zij, gelet op hun samenhang, als één geheel moeten of kunnen worden beschouwd. De bedoeling van alle betrokken partijen is hiervoor beslissend.\footnote{HR 19 december 2008, ECLI:NL:HR:2008:BG1117 (Air Holland).}

5.4. Bij de toepassing van artikel 42 Fw op een samenhangend geheel van rechtshandelingen gaat het erom of dit samenhangende geheel als zodanig schuldeisers heeft benadeeld. Daarmee strookt dat ook de beoordeling van de wetenschap van benadeling wordt betrokken op het samenhangende geheel van rechtshandelingen. Daarbij gaat het erom dat op enig moment bij het verrichten van een rechtshandeling die behoort tot het samenhangende geheel van rechtshandelingen, is voldaan aan de voorwaarde dat de schuldenaar en, in het geval van art. 42 lid 2 Fw, degene met of jegens wie de schuldenaar de rechtshandeling verrichtte, wist of behoorde te weten dat van dat samenhangende geheel van rechtshandelingen benadeling van schuldeisers het gevolg zou zijn. Met het voorgaande strookt eveneens dat ook de beoordeling van de onverplichtheid die artikel 42 Fw als voorwaarde stelt voor vernietiging, wordt betrokken op het samenhangende geheel van rechtshandelingen.\footnote{HR 20 juni 2025, ECLI:NL:HR:2025:975.}

Het oordeel van de rechtbank

5.5. De rechtbank komt tot het oordeel dat er sprake is van een samenhangend geheel van rechtshandelingen. De gewraakte rechtshandeling is onverplicht, maar de schuldeisers van Concepts zijn door deze rechtshandeling niet benadeeld. Ook is de rechtbank van oordeel dat niet voldaan is aan het vereiste van wetenschap van benadeling. Dit betekent dat de curator geen beroep kan doen op artikel 42 Fw en dat de vorderingen van de curator dus worden afgewezen. De rechtbank legt hierna uit hoe zij tot dit oordeel komt.

\emph{Samenhangend geheel van rechtshandelingen: de kasronde}

5.6. Profundo voert aan dat de gewraakte betaling plaatsvond in het kader van een kasronde, waarin sprake was van meerdere samenhangende rechtshandelingen. De kasronde heeft op dezelfde dag in onderlinge afstemming plaatsgevonden op verzoek van {[}naam 2{]} , ter opschoning van de rekening-courantverhoudingen na afronding van een project. Profundo voert aan dat het gebruikelijk was om na afronding van een project tot een dergelijke opschoning over te gaan.

5.7. De curator betoogt dat de kasronde is uitgevoerd ter versteviging van het eigen vermogen van Construction. De curator verwijst naar de e-mail van 5 augustus 2024 van mr. Froger (zie 3.12). De curator verwijst ook naar de dividendverklaring van 29 november 2017 (zie 3.9). Het op 1 augustus 2019 door Concepts vanuit Construction ontvangen bedrag van € 552.000,00 had niet in strijd met de hiervoor bedoelde dividendverklaring aan Concepts mogen worden doorbetaald. Het vermogen had gelet op de dividendverklaring in Construction moeten blijven. Volgens de curator gaat het niet alleen om de kasronde van 2019, maar moet ook deze verklaring jegens Woningborg uit 2017 betrokken worden in het gehele samenstel.

5.8. De rechtbank volgt Profundo in haar standpunt dat de kasronde bestaat uit met elkaar samenhangende rechtshandelingen. Vast staat dat de kasronde zoals gesteld door Profundo heeft plaatsgevonden, nu de curator dat heeft erkend. Het staat dus vast dat op dezelfde dag (1 augustus 2019) meerdere kasstromen zijn gaan vloeien, beginnend bij Profundo en eindigend bij Profundo, totaal belopend een bedrag van (afgerond)\\
€ 670.000,00. De betaling van € 552.500,00 door Concepts aan Profundo kon enkel plaatsvinden doordat Concepts hetzelfde bedrag in het kader van de kasronde had ontvangen van Construction. De curator heeft niet weersproken dat de kasronde op initiatief van {[}naam 2{]} heeft plaatsgevonden. Uit de kasronde zelf volgt al dat er onderlinge afstemming moet zijn geweest. Profundo heeft tijdens de zitting voldoende gemotiveerd weersproken dat de bedoeling van de kasronde was om het eigen vermogen van Construction te versterken of verband hield met een afspraak met Woningborg uit 2017. Met de kasronde is bedoeld de onderlinge verhoudingen op te schonen. Het gaat hier volgens Profundo om twee bedragen en twee aparte zaken. In het licht van deze betwisting heeft de curator zijn stellingen op dit punt onvoldoende onderbouwd. Er bestaat naar het oordeel van de rechtbank zodanige samenhang tussen de vijf rechtshandelingen die in het kader van de kasronde zijn verricht dat deze als één geheel moeten worden beschouwd.

\emph{Een onverplichte rechtshandeling}

5.9. De curator stelt zich op het standpunt dat de aflossing van € 552.500,00 op de geldlening is aan te merken als een onverplichte rechtshandeling, nu Concepts op grond van de overeenkomst van geldlening niet gehouden was om af te lossen. Profundo heeft zich wat betreft de onverplichtheid van de betaling gerefereerd aan het oordeel van de rechtbank.

5.10. Profundo heeft de stelling van de curator dat er sprake is van een onverplichte rechtshandeling daarmee niet betwist. De rechtbank gaat daarom uit van de juistheid van die stelling.

\emph{Geen benadeling van schuldeisers}

5.11. Vervolgens is de vraag of de schuldeisers van Concepts zijn benadeeld door deze onverplichte rechtshandeling.

5.12. De curator stelt zich op het standpunt dat de schuldeisers zijn benadeeld als gevolg van de onverplichte betaling van € 552.500,00 aan Profundo. Van benadeling van schuldeisers is volgens de curator sprake wanneer het (verhaalbare) vermogen van Concepts ``per saldo'' is afgenomen of de onderlinge rangorde tussen de schuldeisers als gevolg van de rechtshandeling is doorkruist. Het is noodzakelijk, maar ook voldoende dat de benadeling aanwezig is op het moment dat daarover door de rechter wordt beslist. De vraag of van benadeling sprake is, moet worden beantwoord door de hypothetische situatie waarin de schuldeisers zouden hebben verkeerd zonder de gewraakte rechtshandeling te vergelijken met de situatie waarin ze feitelijk verkeren als de rechtshandeling onaangetast blijft. De curator is van mening dat door de betaling van € 552.500,00 de paritas creditorum (gelijkheid van schuldeisers) is verstoord. Immers is Profundo betaald voor dit bedrag en hebben de (andere) schuldeisers in het faillissement van Concepts niets ontvangen. Het is niet relevant of de schulden voortkomen uit een rechtsverhouding die er op 1 augustus 2019 al was. De curator vindt de eerste twee stappen in de kasronde (de betaling van Profundo aan Think Holding B.V. en de daaropvolgende betaling van Think Holding B.V. aan Construction, zie 3.10) geoorloofd, maar daar had het moeten stoppen. Het geld had in Construction moeten blijven. De kasronde is echter vervolgd met de betaling van Construction aan Concepts. Aangezien Concepts na 1 augustus 2019 zelf geen inkomsten meer kreeg en er geen liquiditeit in de vennootschap was in juli 2019, betrof het ontvangen geld het laatste kapitaal van de vennootschap. De betaling aan Profundo door Concepts is in strijd met de uitkeringstoets van artikel 2:216 BW en heeft geleid tot benadeling van de schuldeisers, aldus de curator.

5.13. Profundo betwist dat er sprake is van benadeling. Bij beoordeling van het geheel van rechtshandelingen/betalingen in het kader van de kasronde blijkt volgens Profundo dat van benadeling geen sprake is. Als de kasronde niet zou hebben plaatsgevonden, dan zou de betaling door Concepts aan Profundo niet hebben plaatsgevonden, maar dan zou Concepts het bedrag van de betaling zelf ook niet ontvangen hebben. Door de betaling enerzijds en de ontvangst van een gelijk bedrag anderzijds is per saldo van verslechtering van de positie van de (overige) crediteuren geen sprake. Integendeel: zonder dat de liquide middelen afnamen, zijn de schulden wel afgenomen. De curator heeft ten onrechte ook geen vernietiging gevorderd van het samenstel van rechtshandelingen. Ook daarom zijn de op artikel 42 Fw gestoelde vorderingen niet toewijsbaar.

5.14. Profundo stelt verder dat er op 1 augustus 2019 geen bestaande schuldeisers met opeisbare of voorzienbare schulden zijn benadeeld. Benadeling kan ook toekomstige schuldeisers betreffen, als die door de bestreden betalingen in hun belangen zijn geschaad. Dat betekent dat er een causaal verband moet zijn tussen de bestreden handelingen op 1 augustus 2019 en de onbetaald gebleven facturen. Dat causale verband wordt niet aangetoond. De curator toont namelijk niet aan dat de betaling aan Profundo aangewend had kunnen worden voor de door de curator genoemde schuldeisers, nu Concepts na 1 augustus 2019 in totaal nog € 3,3 miljoen aan crediteuren heeft voldaan. Als de van Construction ontvangen € 552.500,00 niet zou zijn doorgestort aan Profundo, dan zou het bedrag zijn aangewend ter betaling van facturen die nu `gedekt' zijn in de totale betalingen van € 3,3 miljoen. Ook zonder de aflossing aan Profundo zouden de crediteuren die zich bij de curator hebben gemeld onbetaald zijn gebleven. Dat zij onbetaald zijn gebleven is niet veroorzaakt door de betalingen van 1 augustus 2019, maar door het feit dat {[}naam 1{]} in april 2020 besloot om geen aanvullende financiering meer te verstrekken aan Concepts.

5.15. De rechtbank overweegt als volgt.

5.16. Zoals hiervoor overwogen, beschouwt de rechtbank de kasronde als een samenstel van rechtshandelingen, zodat de rechtbank de (nadelige) gevolgen van deze rechtshandelingen in onderling verband moet beoordelen. Bij de toepassing van artikel 42 Fw op een samenhangend geheel van rechtshandelingen gaat het er kortom om of dit samenhangende geheel als zodanig schuldeisers heeft benadeeld. Een rechtshandeling die op zichzelf bezien niet benadelend is, kan door samenhang met andere rechtshandelingen wel benadelend zijn; andersom geldt dat een rechtshandeling die op zichzelf bezien benadelend is, in samenhang met andere rechtshandelingen niet benadelend kan zijn.

5.17. De curator neemt in zijn stellingen het geheel van rechtshandelingen niet tot uitgangspunt. Hij betoogt dat de eerste stappen van de kasronde geoorloofd zijn, maar dat het had moeten stoppen bij Construction. De ontvangen betaling had Construction niet mogen verlaten. Aan dit betoog gaat de rechtbank echter voorbij. De curator miskent daarmee namelijk dat de kasronde en de gevolgen daarvan in zijn geheel moet worden beschouwd. Profundo heeft daarnaast gemotiveerd weersproken dat de kasronde bedoeld was ter versteviging van het eigen vermogen van Construction of enig verband hield met een afspraak met Woningborg van twee jaar daarvoor. Afgezien daarvan ziet de rechtbank niet in waarom het voortzetten van de kasronde door Construction voor (de schuldeisers van) Concepts nadelige gevolgen heeft gehad. Stel dat de curator gevolgd zou worden in zijn gedachtenlijn, dan was het geld namelijk in Construction gebleven. Concepts had in dat geval dus ook niet over het bedrag van € 552.000,00 kunnen beschikken.

5.18. Verder is niet in geschil dat Concepts tot 1 augustus 2019 geen liquiditeit had om Profundo of andere schuldeisers te betalen. Alleen door deel te nemen aan de kasronde tussen verschillende groepsvennootschappen heeft Concepts op 1 augustus 2019 een bedrag van € 552.500,00 van Construction ontvangen. Met dit bedrag heeft Concepts diezelfde dag\\
€ 552.500,00 afgelost op de lening van Profundo. Deze aflossing was onderdeel van dezelfde kasronde en vond in feite plaats met geld dat afkomstig was van Profundo zelf, die met een lening van € 670.000,00 aan Think Holding B.V. de kasronde was gestart. Na de betaling aan Profundo als sluitstuk van de kasronde had Concepts minder schulden en (opnieuw) geen liquiditeit.

5.19. Om te bepalen of de kasronde, als samenhangend geheel, schuldeisers van Concepts heeft benadeeld, moet deze feitelijke situatie vergeleken worden met de situatie waarin schuldeisers zouden hebben verkeerd zonder de kasronde. Zonder de kasronde had Concepts op 1 augustus 2019 geen bedrag ontvangen van € 552.500,00 en dus geen liquiditeit gehad om schuldeisers te betalen. Dat Concepts op 1 augustus 2019 kortstondig over het bedrag van € 552.500,00 heeft beschikt, kan dan ook niet los gezien worden van de kasronde en de bedoeling daarvan. Daarmee zou in de ogen van de rechtbank niet in lijn zijn geweest dat Concepts wél deel zou nemen aan de kasronde door € 552.500,00 te ontvangen, maar vervolgens zelf geen verdere uitvoering zou geven aan de kasronde. Deze situatie laat zich naar het oordeel van de rechtbank ook niet vergelijken met een enkelvoudige betaling uit het laatste (eigen) kapitaal van een vennootschap.

5.20. Gezien al het voorgaande is de rechtbank van oordeel dat de kasronde (als samenhangend geheel) niet heeft geleid tot benadeling van schuldeisers van Concepts.

\emph{Ten overvloede: geen wetenschap van benadeling}

5.21. Ook als er wel sprake zou zijn geweest van benadeling van schuldeisers van Concepts, zou het beroep op artikel 42 Fw overigens niet slagen. De rechtbank volgt de curator namelijk ook niet in zijn stelling dat Concepts en Profundo op 1 augustus 2019 wetenschap van benadeling hadden. Van wetenschap van benadeling kan alleen sprake zijn geweest als zowel Concepts als Profundo op 1 augustus 2019 het faillissement van Concepts en een tekort daarin met een redelijke mate van waarschijnlijkheid konden voorzien.\footnote{Hoge Raad 22 december 2009, ECLI:NL:HR:2009:BI8493.} De curator stelt dat de wetenschap van benadeling bij beide partijen aanwezig moet worden geacht omdat {[}naam 1{]} 100\% eigenaar was van Profundo en over een bankvolmacht beschikte. Volgens de curator had {[}naam 1{]} moeten weten dat Concepts op 1 augustus 2019 geen (noemenswaardige) inkomsten meer kon verwachten en kunstmatig overeind moest worden gehouden. Ook wijst de curator op een procedure tussen Concepts en {[}bedrijf{]} die op 1 augustus 2019 aanhangig was.

5.22. Profundo heeft de wetenschap van benadeling gemotiveerd betwist. Profundo heeft daarbij onder meer gewezen op een onder ede afgelegde getuigenverklaring van {[}naam 1{]} , die daarin weerspreekt dat hij op de hoogte was van de financiële stand van zaken van Concepts. Ook heeft Profundo onweersproken aangevoerd dat Concepts in de procedure tegen {[}bedrijf{]} pas bij eindvonnis van 12 mei 2021 is veroordeeld tot een schadevergoeding op te maken bij staat wegens onrechtmatige beslaglegging. Dit was op 1 augustus 2019 nog niet bekend. Verder heeft Profundo aangevoerd dat Concepts ná 1 augustus 2019 nog ongeveer € 2,5 miljoen aan financiering heeft ontvangen en nog voor een bedrag van ongeveer € 3,3 miljoen euro aan schuldeisers heeft betaald. Pas in april 2020 heeft {[}naam 1{]} besloten geen aanvullende financiering meer aan Concepts te verstrekken, maar ook daarna is de onderneming van Concepts nog maanden voortgezet. Ook dit heeft de curator niet weersproken. In de dagvaarding schrijft de curator bovendien zelf dat Concepts en Profundo ``\emph{medio april 2020 wisten dat het faillissement aanstaande was en er sprake was van een onomkeerbaar groot tekort om alle crediteuren te betalen}'', omdat vanaf dat moment de geldkraan werd dichtgedraaid. Ook in het door de curator opgemaakte faillissementsverslag van Profundo van 13 augustus 2025 staat dat de aandeelhouder in het voorjaar 2020 de geldkraan naar Profundo heeft dichtgedraaid, ``\emph{als gevolg waarvan gefailleerde niet meer aan haar verplichtingen kon voldoen}''. Klaarblijkelijk is de curator dus ook zelf van mening dat het faillissement pas aanstaande was na het stopzetten van de financiering door de aandeelhouder. Dit gebeurde in april 2020. Voor die tijd werd Concepts als gezegd nog wel gefinancierd en heeft zij schuldeisers betaald.

5.23. Gelet op alle voornoemde omstandigheden valt niet in te zien dat Concepts en Profundo het faillissement en het tekort daarin al op 1 augustus 2019 met een redelijke mate van waarschijnlijkheid konden voorzien. Van wetenschap van benadeling kan op 1 augustus 2019 dus geen sprake zijn geweest.

\emph{Conclusie}

5.24. De rechtbank concludeert dat aan meerdere vereisten van artikel 42 Fw niet is voldaan. Dat betekent dat de rechtbank alle vorderingen van de curator zal afwijzen wegens gebrek aan grondslag.

Proceskosten

5.25. De curator is in het ongelijk gesteld en moet daarom de proceskosten (inclusief nakosten) betalen. De proceskosten van Profundo worden begroot op:

{\def\LTcaptype{none} % do not increment counter
\begin{longtable}[]{@{}
  >{\raggedright\arraybackslash}p{(\linewidth - 6\tabcolsep) * \real{0.2500}}
  >{\raggedright\arraybackslash}p{(\linewidth - 6\tabcolsep) * \real{0.2500}}
  >{\raggedright\arraybackslash}p{(\linewidth - 6\tabcolsep) * \real{0.2500}}
  >{\raggedright\arraybackslash}p{(\linewidth - 6\tabcolsep) * \real{0.2500}}@{}}
\toprule\noalign{}
\begin{minipage}[b]{\linewidth}\raggedright
- griffierecht
\end{minipage} & \begin{minipage}[b]{\linewidth}\raggedright
€
\end{minipage} & \begin{minipage}[b]{\linewidth}\raggedright
6.861,00
\end{minipage} & \begin{minipage}[b]{\linewidth}\raggedright
\end{minipage} \\
\midrule\noalign{}
\endhead
\bottomrule\noalign{}
\endlastfoot
- salaris advocaat & € & 7.446,00 & (2 punten × € 3.723,00) \\
- nakosten & € & 189,00 & (plus de verhoging zoals vermeld in de beslissing) \\
Totaal & € & 14.496,00 & \\
\end{longtable}
}

\paragraph{6 De beslissing}\label{de-beslissing-4}

De rechtbank

6.1. wijst de vorderingen van de curator af,

6.2. veroordeelt de curator in de proceskosten van € 14.496,00, te betalen binnen veertien dagen na aanschrijving daartoe, te vermeerderen met € 98,00 plus de kosten van betekening als de curator niet tijdig aan de veroordelingen voldoet en het vonnis daarna wordt betekend,

6.3. verklaart dit vonnis wat betreft de proceskostenveroordeling uitvoerbaar bij voorraad.

Dit vonnis is gewezen door mr. De Vlieger en in het openbaar uitgesproken door\\
mr. Van den Heuvel op 20 mei 2026.

\begin{center}\rule{0.5\linewidth}{0.5pt}\end{center}

\newpage

\markboth{ECLI:NL:RBZWB:2026:4349}{mistral/mistral-medium-3.5-128b:fp8 (k = 0)}

\subsection{k = 0}\label{k-0-4}

\subsubsection{Mistral}\label{mistral-12}

\texttt{mistral/mistral-medium-3.5-128b:fp8}

\begin{center}
\small
\begin{tabular}{@{}p{7.4cm}p{2.6cm}p{2.6cm}@{}}
\toprule
Variabele & Waarde (model) & Waarde (gold) \\
\midrule
Is rechtshandeling & \texttt{true} & \texttt{true} \\
Datum rechtshandeling & \texttt{2019-08-01} & \texttt{2019-08-01} \\
Verricht door de schuldenaar & \texttt{true} & \texttt{true} \\
Om niet & \texttt{false} & \texttt{false} \\
Datum faillietverklaring & \texttt{2020-08-04} & \texttt{2020-08-04} \\
Onverplicht & \texttt{true} & \texttt{true} \\
Benadeling schuldeisers & \textcolor{gray}{\texttt{true}} & \texttt{false} \\
Wetenschap schuldenaar aangetoond & \texttt{false} & \texttt{false} \\
Wetenschap wederpartij aangetoond & \texttt{false} & \texttt{false} \\
Vermoeden wetenschap schuldenaar tegenbewijs & \texttt{null} & \texttt{null} \\
Vermoeden wetenschap wederpartij tegenbewijs & \texttt{null} & \texttt{null} \\
Tegen onderwaarde & \texttt{null} & \texttt{null} \\
Voldoening of zekerheid niet opeisbare schuld & \texttt{null} & \texttt{null} \\
Gelieerde partij & \textcolor{gray}{\texttt{true}} & \texttt{null} \\
Reeds verplicht voor aanvang termijn & \textcolor{gray}{\texttt{true}} & \texttt{null} \\
\textbf{Vernietigbaar} & \textbf{\texttt{false}} & \textbf{\texttt{false}} \\
\bottomrule
\end{tabular}
\end{center}

\paragraph{Is rechtshandeling - overwegingen}\label{is-rechtshandeling---overwegingen-36}

De betaling van € 552.500,00 door Concepts aan Profundo is een rechtshandeling, omdat het een op een rechtsgevolg (betaling) gerichte wil betreft die zich door een verklaring (de daadwerkelijke betaling) heeft geopenbaard.

\paragraph{Verricht door de schuldenaar - overwegingen}\label{verricht-door-de-schuldenaar---overwegingen-36}

De betaling is verricht door Concepts, de latere gefailleerde, en dus door de schuldenaar zelf.

\paragraph{Om niet - overwegingen}\label{om-niet---overwegingen-36}

De betaling van € 552.500,00 door Concepts aan Profundo maakt deel uit van een kasronde binnen de groep, waarbij Concepts een tegenprestatie ontving (namelijk de aflossing van haar rekening-courantschuld door Construction). Er was dus sprake van een reële tegenprestatie, waardoor de handeling niet `om niet' is.

\paragraph{Onverplicht - overwegingen}\label{onverplicht---overwegingen-36}

De betaling van € 552.500,00 door Concepts aan Profundo op 1 augustus 2019 is onverplicht in de zin van artikel 42 lid 1 Fw. Uit de overeenkomst van geldlening (artikel 2 en 4) blijkt dat de lening geen vaste looptijd had en geen aflossingstermijnen kende. De lening was alleen dadelijk opeisbaar bij faillissement, surséance, akkoord buiten faillissement, nalatigheid na ingebrekestelling of conservatoire/executoriale maatregelen (artikel 4). Op 1 augustus 2019 was geen van deze voorwaarden ingetreden, zodat de schuld van Concepts aan Profundo nog niet opeisbaar was. Betaling van een niet-opeisbare schuld is onverplicht. De kasronde was een interne groepsregeling en schept geen rechtsplicht tot betaling. Morele, feitelijke of economische druk (zoals de groepsstructuur) schept geen rechtsplicht. De curator heeft derhalve aannemelijk gemaakt dat de betaling onverplicht was.

\paragraph{Benadeling schuldeisers - overwegingen}\label{benadeling-schuldeisers---overwegingen-36}

\begin{enumerate}
\def\labelenumi{\arabic{enumi}.}
\item
  \textbf{Vermogensvergelijking}: De betaling van € 552.500 door Concepts aan Profundo op 1 augustus 2019 heeft geleid tot een vermindering van het actief van Concepts. Zonder deze betaling zou dit bedrag in de boedel zijn gebleven en ten goede zijn gekomen aan alle schuldeisers (inclusief Profundo, maar dan in de rangorde van de faillissementswet).
\item
  \textbf{Hypothetische situatie}: In de hypothetische situatie zonder de betaling zou Concepts € 552.500 meer in kas hebben gehad. Gelet op de latere faillietverklaring (4 augustus 2020) en het feit dat Concepts na 1 augustus 2019 nog wel financiering ontving en schulden afbetaalde, maar uiteindelijk toch failliet ging, is aannemelijk dat dit bedrag had kunnen bijdragen aan de voldoening van schulden in het faillissement. De curator hoeft niet aan te tonen dat het bedrag \emph{volledig} ten goede zou zijn gekomen aan de schuldeisers; voldoende is dat de schuldeisers in het algemeen zijn benadeeld doordat het actief is verminderd.
\item
  \textbf{Middellijke benadeling}: De benadeling hoeft niet direct zichtbaar te zijn op het moment van de handeling (1 augustus 2019), maar moet aanwezig zijn op het moment waarop de curator zijn rechten doet gelden (26 juli 2024). Gelet op het faillissement van Concepts en de onvoldoende dekking van de schulden in de boedel, is er sprake van middellijke benadeling.
\item
  \textbf{Rangordeverstoring}: Hoewel de betaling mogelijk een schuld aan Profundo (als groepsvennootschap) afdeed, is niet gebleken dat Profundo een voorrechtelijke of geprefereerde vordering had. De betaling heeft derhalve de rangorde tussen schuldeisers verstoord, doordat Profundo (als groepsvennootschap) voorrang heeft gekregen boven andere schuldeisers. Dit is voldoende voor benadeling in de zin van artikel 42 Fw, ook als de lening reëel was.
\item
  \textbf{Geen tegenargumenten}: De door Profundo aangevoerde kasronde en versteviging van het eigen vermogen van de groep doen niet af aan het feit dat de betaling het actief van Concepts heeft verminderd. De latere financiering en betalingen door Concepts (€ 2,5 miljoen ontvangen, € 3,3 miljoen betaald) tonen aan dat de onderneming nog actief was, maar niet dat de betaling van € 552.500 geen benadeling heeft veroorzaakt. Immers, zonder deze betaling had Concepts meer middelen gehad om schulden af te lossen of de onderneming voort te zetten.
\item
  \textbf{Kennis van benadeling}: Voor de toepassing van artikel 42 lid 2 Fw (aangezien de betaling een meerzijdige rechtshandeling betreft) moet ook Profundo hebben geweten of behoord te weten dat de betaling benadeling van schuldeisers tot gevolg zou hebben. Gelet op de groepsstructuur, de betrokkenheid van {[}naam 1{]} en {[}naam 2{]} (die ook bestuurder was van Concepts), en de kasronde die duidelijk gericht was op interne groepsfinanciering, is aannemelijk dat Profundo bekend was met de financiële situatie van Concepts en de risico's van de betaling. Dit wordt versterkt door het feit dat de leningovereenkomst een clausule bevatte voor directe opeisbaarheid bij faillissement (artikel 4 onder a), wat wijst op bewustzijn van de risico's.
\end{enumerate}

\paragraph{Wetenschap - overwegingen}\label{wetenschap---overwegingen-36}

De betaling van € 552.500,- door Concepts aan Profundo op 1 augustus 2019 valt binnen de jaartermijn voor de faillietverklaring (4 augustus 2020). Echter, de schuld van Concepts aan Profundo bestond reeds uit de leningovereenkomst van 13 juli 2017, waardoor Concepts zich \emph{reeds voor de aanvang van de jaartermijn} (1 augustus 2018) tot de betaling had verplicht. Dit sluit het bewijsvermoeden van art. 43 Fw uit. Wel is sprake van een gelieerde partij (art. 43 lid 1 sub 6° Fw), maar dit vermoeden is niet van toepassing omdat de verplichting reeds bestond. De curator moet derhalve de wetenschap van benadeling (art. 42 lid 1 Fw) \emph{rechtstreeks} aantonen. Uit de casus blijkt niet dat de curator dit heeft gedaan. Profundo heeft tegenbewijs geleverd door aan te tonen dat Concepts na 1 augustus 2019 nog significante financiering ontving en betalingen deed, wat suggereert dat benadeling niet met redelijke waarschijnlijkheid voorzienbaar was. Derhalve is de wetenschap niet aangetoond.

\newpage

\markboth{ECLI:NL:RBZWB:2026:4349}{gpt-5.5 (k = 0)}

\subsubsection{GPT}\label{gpt-12}

\texttt{gpt-5.5}

\begin{center}
\small
\begin{tabular}{@{}p{7.4cm}p{2.6cm}p{2.6cm}@{}}
\toprule
Variabele & Waarde (model) & Waarde (gold) \\
\midrule
Is rechtshandeling & \texttt{true} & \texttt{true} \\
Datum rechtshandeling & \texttt{2019-08-01} & \texttt{2019-08-01} \\
Verricht door de schuldenaar & \texttt{true} & \texttt{true} \\
Om niet & \texttt{false} & \texttt{false} \\
Datum faillietverklaring & \texttt{2020-08-04} & \texttt{2020-08-04} \\
Onverplicht & \texttt{true} & \texttt{true} \\
Benadeling schuldeisers & \textcolor{gray}{\texttt{true}} & \texttt{false} \\
Wetenschap schuldenaar aangetoond & \texttt{false} & \texttt{false} \\
Wetenschap wederpartij aangetoond & \texttt{false} & \texttt{false} \\
Vermoeden wetenschap schuldenaar tegenbewijs & \textcolor{gray}{\texttt{true}} & \texttt{null} \\
Vermoeden wetenschap wederpartij tegenbewijs & \textcolor{gray}{\texttt{true}} & \texttt{null} \\
Tegen onderwaarde & \texttt{false} & \texttt{null} \\
Voldoening of zekerheid niet opeisbare schuld & \textcolor{gray}{\texttt{true}} & \texttt{null} \\
Gelieerde partij & \textcolor{gray}{\texttt{true}} & \texttt{null} \\
Reeds verplicht voor aanvang termijn & \texttt{false} & \texttt{null} \\
\textbf{Vernietigbaar} & \textbf{\texttt{false}} & \textbf{\texttt{false}} \\
\bottomrule
\end{tabular}
\end{center}

\paragraph{Is rechtshandeling - overwegingen}\label{is-rechtshandeling---overwegingen-37}

De gewraakte handeling betreft de betaling door Concepts van € 552.500,00 aan Profundo ter terugbetaling/aflossing op de lening. Een betaling en de daarmee samenhangende vermogensverschuiving zijn gericht op rechtsgevolg, namelijk voldoening van een schuld en vermindering van de vordering van Profundo op Concepts.

\paragraph{Verricht door de schuldenaar - overwegingen}\label{verricht-door-de-schuldenaar---overwegingen-37}

Uit de feiten volgt dat Concepts zelf op 1 augustus 2019 het bedrag van € 552.500,00 aan Profundo heeft betaald.

\paragraph{Om niet - overwegingen}\label{om-niet---overwegingen-37}

De betaling vond plaats als terugbetaling op de lening van Concepts aan Profundo. Daartegenover stond de vermindering van de schuld van Concepts aan Profundo; de handeling is daarom anders dan om niet.

\paragraph{Onverplicht - overwegingen}\label{onverplicht---overwegingen-37}

De aangevochten rechtshandeling betreft de betaling door Concepts aan Profundo van € 552.500 op 1 augustus 2019 ter aflossing op de geldlening. Wel bestond er een schuld van Concepts aan Profundo uit hoofde van de overeenkomst van geldlening, maar uit de overeenkomst volgt geen verplichting voor Concepts om op 1 augustus 2019 af te lossen. De lening had geen vaste looptijd en bevatte geen aflossingstermijnen. Profundo kon de lening slechts opzeggen met een opzegtermijn van zes maanden, en de lening was alleen dadelijk opeisbaar bij de in artikel 4 genoemde omstandigheden, zoals faillissement, surseance, een buitengerechtelijk akkoord, verzuim na ingebrekestelling of beslagmaatregelen. Uit de casus blijkt niet dat Profundo vóór 1 augustus 2019 had opgezegd of dat zich een contractuele opeisingsgrond had voorgedaan. De betaling vond plaats in het kader van een groepsinterne kasronde en kennelijk ter versteviging van het eigen vermogen binnen de groep, maar een dergelijke feitelijke, economische of groepsvennootschappelijke aanleiding schept geen in rechte afdwingbare rechtsplicht tot betaling. Omdat de schuld op dat moment niet opeisbaar was en Concepts niet contractueel of wettelijk verplicht was deze betaling toen te verrichten, is de betaling onverplicht verricht in de zin van artikel 42 Fw.

\paragraph{Benadeling schuldeisers - overwegingen}\label{benadeling-schuldeisers---overwegingen-37}

De benadeling moet worden beoordeeld door de feitelijke situatie, waarin de betaling van € 552.500 aan Profundo in stand blijft, te vergelijken met de hypothetische situatie zonder die betaling. De betaling leidde ertoe dat liquide middelen die Concepts kort daarvoor van Construction ontving, niet beschikbaar bleven voor de gezamenlijke schuldeisers maar werden aangewend ter voldoening van een groepsvordering van Profundo. Dat daartegenover de schuld aan Profundo met hetzelfde bedrag afnam, neemt de benadeling niet weg: bij faillissement is een selectieve betaling aan één concurrente schuldeiser benadelend wanneer daardoor de paritas creditorum wordt doorbroken en andere schuldeisers minder verhaal hebben dan zij zonder die betaling zouden hebben gehad. Profundo heeft door de betaling feitelijk volledige voldoening gekregen voor € 552.500, terwijl zij zonder betaling in het faillissement slechts naar rato had kunnen delen. De kasronde als geheel maakt dit niet anders. De ontvangst door Concepts van € 552.500 van Construction was wel een gunstig gevolg, maar die ontvangst betrof aflossing van een vordering van Concepts op Construction en werd direct doorgeschoven naar Profundo. Per saldo werd een verhaalsmogelijkheid voor de gezamenlijke schuldeisers aangewend ten gunste van één schuldeiser. Ook de door Profundo gestelde latere financiering van circa € 2,5 miljoen en latere betalingen aan schuldeisers van circa € 3,3 miljoen sluiten benadeling niet uit. Daaruit volgt niet dat zonder de betaling aan Profundo alle schuldeisers volledig zouden zijn voldaan, noch dat de gezamenlijke schuldeisers zich niet op het bedrag hadden kunnen verhalen. Op het moment dat de curator zijn rechten uitoefende, was het bedrag niet in de boedel aanwezig en was de rangorde tussen schuldeisers door de betaling verstoord.

\paragraph{Wetenschap - overwegingen}\label{wetenschap---overwegingen-37}

Rechtstreekse wetenschap van benadeling is niet aangetoond. Beslissend is of op 1 augustus 2019 het faillissement van Concepts en een tekort daarin met redelijke mate van waarschijnlijkheid voorzienbaar waren voor Concepts en, bij deze niet-om-niet verrichte betaling, ook voor Profundo. De beschikbare feiten wijzen daar onvoldoende op: onweersproken is dat Concepts na de betaling nog circa € 2,5 miljoen financiering ontving, nog circa € 3,3 miljoen aan schuldeisers betaalde, dat pas in april 2020 is besloten geen aanvullende financiering meer te verstrekken en dat de onderneming daarna nog maanden is voortgezet. Wel betreft de betaling inhoudelijk aflossing van een langlopende lening zonder aflossingstermijn/opeising en dus voldoening van een niet-opeisbare schuld aan een groepsmaatschappij/gelieerde partij; er bestond geen eerdere verplichting tot deze concrete aflossing. Voor zover daardoor een bewijsvermoeden van art. 43 Fw aan beide zijden zou gelden, wordt dat door de genoemde omstandigheden voldoende ontzenuwd. Art. 45 Fw ligt niet voor de hand, omdat de betaling strekte tot vermindering van een schuld en dus niet om niet was.

\newpage

\markboth{ECLI:NL:RBZWB:2026:4349}{gemini-3.5-flash (k = 0)}

\subsubsection{Gemini}\label{gemini-12}

\texttt{gemini-3.5-flash}

\begin{center}
\small
\begin{tabular}{@{}p{7.4cm}p{2.6cm}p{2.6cm}@{}}
\toprule
Variabele & Waarde (model) & Waarde (gold) \\
\midrule
Is rechtshandeling & \texttt{true} & \texttt{true} \\
Datum rechtshandeling & \texttt{2019-08-01} & \texttt{2019-08-01} \\
Verricht door de schuldenaar & \texttt{true} & \texttt{true} \\
Om niet & \texttt{false} & \texttt{false} \\
Datum faillietverklaring & \texttt{2020-08-04} & \texttt{2020-08-04} \\
Onverplicht & \texttt{true} & \texttt{true} \\
Benadeling schuldeisers & \textcolor{gray}{\texttt{true}} & \texttt{false} \\
Wetenschap schuldenaar aangetoond & \texttt{false} & \texttt{false} \\
Wetenschap wederpartij aangetoond & \texttt{false} & \texttt{false} \\
Vermoeden wetenschap schuldenaar tegenbewijs & \textcolor{gray}{\texttt{true}} & \texttt{null} \\
Vermoeden wetenschap wederpartij tegenbewijs & \textcolor{gray}{\texttt{true}} & \texttt{null} \\
Tegen onderwaarde & \texttt{null} & \texttt{null} \\
Voldoening of zekerheid niet opeisbare schuld & \textcolor{gray}{\texttt{true}} & \texttt{null} \\
Gelieerde partij & \textcolor{gray}{\texttt{true}} & \texttt{null} \\
Reeds verplicht voor aanvang termijn & \texttt{false} & \texttt{null} \\
\textbf{Vernietigbaar} & \textbf{\texttt{false}} & \textbf{\texttt{false}} \\
\bottomrule
\end{tabular}
\end{center}

\paragraph{Is rechtshandeling - overwegingen}\label{is-rechtshandeling---overwegingen-38}

De betaling van een bedrag van € 552.500,00 door Concepts aan Profundo ter aflossing van een lening is een rechtshandeling, aangezien deze gericht is op het rechtsgevolg van het tenietgaan van een verbintenis (nakoming).

\paragraph{Verricht door de schuldenaar - overwegingen}\label{verricht-door-de-schuldenaar---overwegingen-38}

De betaling is rechtstreeks verricht door de schuldenaar, Concepts, aan haar schuldeiser Profundo.

\paragraph{Om niet - overwegingen}\label{om-niet---overwegingen-38}

De betaling betreft de aflossing van een lening (terugbetaling op de lening), waardoor een schuld van Concepts aan Profundo tenietgaat. Er staat derhalve een prestatie tegenover (bevrijding uit de schuld), waardoor de handeling om baat (anders dan om niet) is.

\paragraph{Onverplicht - overwegingen}\label{onverplicht---overwegingen-38}

De betaling van € 552.500,00 door Concepts aan Profundo op 1 augustus 2019 vond plaats ter gedeeltelijke aflossing van een geldlening. Uit de overeenkomst van geldlening van 13 juli 2017 blijkt dat aan de lening geen looptijd of aflossingstermijnen waren verbonden, en dat Profundo de lening slechts kon opzeggen met een opzegtermijn van zes maanden. Gesteld noch gebleken is dat Profundo de lening voorafgaand aan de betaling had opgezegd of dat er sprake was van een directe opeisingsgrond als bedoeld in artikel 4 van de overeenkomst. De schuld was op 1 augustus 2019 dan ook niet opeisbaar. Het voldoen van een niet-opeisbare schuld is in beginsel onverplicht, omdat hiervoor op dat moment geen in rechte afdwingbare rechtsplicht bestaat. De omstandigheid dat de betaling deel uitmaakte van een concernbrede `kasronde' creëert slechts een feitelijke of economische beweegreden, maar geen juridische verplichting.

\paragraph{Benadeling schuldeisers - overwegingen}\label{benadeling-schuldeisers---overwegingen-38}

De betaling van € 552.500,00 door Concepts aan Profundo op 1 augustus 2019 betreft de aflossing van een op dat moment niet-opeisbare lening. Door deze betaling is een substantieel liquide bedrag aan het vermogen van Concepts onttrokken ten gunste van één specifieke, gelieerde schuldeiser (Profundo). Indien deze betaling achterwege was gebleven, had Concepts de van Construction ontvangen € 552.500,00 behouden. In de hypothetische situatie zonder deze rechtshandeling zou dit bedrag in de faillissementsboedel zijn gevallen en beschikbaar zijn geweest voor de gezamenlijke schuldeisers. Nu de betaling in stand blijft, is het actieve boedelbestanddeel met dit bedrag verminderd. Hoewel de schuld aan Profundo hiermee is afgenomen, leidt deze selectieve betaling tot een directe benadeling van de overige schuldeisers, aangezien Profundo als concurrente schuldeiser in het faillissement anders slechts een fractie van haar vordering voldaan zou hebben gekregen. Er is derhalve sprake van een vermindering van het actief en een verstoring van de paritas creditorum. De latere financieringen en betalingen veranderen niets aan het feit dat er op het moment dat de curator de vernietiging inriep sprake was van een nadeliger positie voor de gezamenlijke schuldeisers.

\paragraph{Wetenschap - overwegingen}\label{wetenschap---overwegingen-38}

De betaling van € 552.500,00 op 1 augustus 2019 vond plaats binnen één jaar voor de faillietverklaring (4 augustus 2020) en betrof de voldoening van een niet-opeisbare schuld aan een gelieerde partij binnen de groep. Hierdoor zijn de wettelijke bewijsvermoedens van art. 43 lid 1 sub 2° en sub 6° Fw van toepassing. Profundo heeft dit vermoeden echter met succes ontzenuwd door aan te tonen dat Concepts na de betaling nog circa € 2,5 miljoen aan financiering ontving en € 3,3 miljoen aan schuldeisers heeft voldaan, waarbij de onderneming nog maanden is voortgezet. Ten tijde van de betaling waren het faillissement en een tekort daarin dan ook niet met een redelijke mate van waarschijnlijkheid te voorzien. De curator heeft de wetenschap van benadeling bij partijen niet rechtstreeks kunnen aantonen.

\newpage

\markboth{ECLI:NL:RBZWB:2026:4349}{mistral/mistral-medium-3.5-128b:fp8 (k = 5)}

\subsection{k = 5}\label{k-5-4}

\subsubsection{Mistral}\label{mistral-13}

\texttt{mistral/mistral-medium-3.5-128b:fp8}

\begin{center}
\small
\begin{tabular}{@{}p{7.4cm}p{2.6cm}p{2.6cm}@{}}
\toprule
Variabele & Waarde (model) & Waarde (gold) \\
\midrule
Is rechtshandeling & \texttt{true} & \texttt{true} \\
Datum rechtshandeling & \texttt{2019-08-01} & \texttt{2019-08-01} \\
Verricht door de schuldenaar & \texttt{true} & \texttt{true} \\
Om niet & \texttt{false} & \texttt{false} \\
Datum faillietverklaring & \texttt{2020-08-04} & \texttt{2020-08-04} \\
Onverplicht & \texttt{true} & \texttt{true} \\
Benadeling schuldeisers & \textcolor{gray}{\texttt{true}} & \texttt{false} \\
Wetenschap schuldenaar aangetoond & \texttt{false} & \texttt{false} \\
Wetenschap wederpartij aangetoond & \texttt{false} & \texttt{false} \\
Vermoeden wetenschap schuldenaar tegenbewijs & \texttt{null} & \texttt{null} \\
Vermoeden wetenschap wederpartij tegenbewijs & \texttt{null} & \texttt{null} \\
Tegen onderwaarde & \texttt{null} & \texttt{null} \\
Voldoening of zekerheid niet opeisbare schuld & \texttt{null} & \texttt{null} \\
Gelieerde partij & \textcolor{gray}{\texttt{true}} & \texttt{null} \\
Reeds verplicht voor aanvang termijn & \texttt{false} & \texttt{null} \\
\textbf{Vernietigbaar} & \textbf{\texttt{false}} & \textbf{\texttt{false}} \\
\bottomrule
\end{tabular}
\end{center}

\paragraph{Is rechtshandeling - overwegingen}\label{is-rechtshandeling---overwegingen-39}

De betaling van € 552.500,00 door Concepts aan Profundo is een rechtshandeling, omdat het een op een rechtsgevolg (aflossing van een schuld) gerichte wil betrof die zich door een verklaring (de betaling) heeft geopenbaard.

\paragraph{Verricht door de schuldenaar - overwegingen}\label{verricht-door-de-schuldenaar---overwegingen-39}

De betaling is verricht door Concepts, de latere gefailleerde schuldenaar.

\paragraph{Om niet - overwegingen}\label{om-niet---overwegingen-39}

De betaling was geen om niet-handeling, omdat deze plaatsvond in het kader van een kasronde binnen de groep, waarbij Concepts een tegenprestatie ontving (de aflossing van haar rekening-courantschuld door Construction).

\paragraph{Onverplicht - overwegingen}\label{onverplicht---overwegingen-39}

De betaling van € 552.500,00 door Concepts aan Profundo op 1 augustus 2019 is onverplicht verricht in de zin van artikel 42 lid 1 Fw.

Uit de casus blijkt dat Concepts en Profundo op 13 juli 2017 een overeenkomst van geldlening hebben gesloten, waarbij een bestaande rekening-courantverhouding van € 1.285.000,00 werd omgezet in een langlopende lening (art. 3.5). In artikel 2 van deze overeenkomst is bepaald dat de lening geen vaste looptijd had en dat Profundo de lening met inachtneming van een opzegtermijn van 6 maanden kon opzeggen (art. 3.6). Artikel 4 van de overeenkomst voorziet in directe opeisbaarheid van de leningstand (inclusief rente) bij faillissement, surséance van betaling, akkoord buiten faillissement, of nalatigheid na ingebrekestelling (art. 3.7).

De betaling van € 552.500,00 op 1 augustus 2019 maakt deel uit van een kasronde binnen de groep, waarbij Profundo eerst € 670.000,00 uitleende aan Think Holding B.V., die dit bedrag gebruikte om haar schuld aan Construction af te lossen. Construction betaalde vervolgens € 552.500,00 aan Concepts voor aflossing van haar rekening-courantschuld, waarna Concepts dit bedrag doorbetaalde aan Profundo (art. 3.10).

Beslissend is of Concepts op grond van de leningsovereenkomst verplicht was om op 1 augustus 2019 € 552.500,00 aan Profundo te betalen. Uit de leningsovereenkomst volgt dat de lening \textbf{niet direct opeisbaar} was, tenzij een van de in artikel 4 genoemde omstandigheden zich voordeden (faillissement, surséance, akkoord, of nalatigheid na ingebrekestelling). Op 1 augustus 2019 was geen van deze omstandigheden aan de orde: Concepts was nog niet failliet (dat gebeurde pas op 4 augustus 2020), en er is geen sprake van een ingebrekestelling of andere opeisbaarheidsgrond. De lening was derhalve \textbf{niet opeisbaar} op het moment van betaling.

Daar de betaling niet voortvloeide uit een op de wet of een overeenkomst berustende rechtsplicht (de lening was immers niet opeisbaar), is de betaling \textbf{onverplicht} in de zin van artikel 42 lid 1 Fw. Dit wordt bevestigd door de uitspraak ECLI:NL:RBOVE:2021:73, waarin de rechtbank oordeelde dat betalingen op een langlopende lening zonder opeisbaarheidsgrond als onverplicht moeten worden aangemerkt. Ook in ECLI:NL:GHARL:2016:3413 werd benadrukt dat een betaling alleen verplicht is als de schuld opeisbaar is (in dat geval was de lening wel opeisbaar op de betalingsdatum).

Morele, feitelijke of economische druk (zoals de kasronde binnen de groep) schept geen rechtsplicht (vgl. de toelichting bij artikel 42 Fw). De stelplicht en bewijslast van het onverplichte karakter rusten op de curator, en in deze casus is voldoende aannemelijk gemaakt dat de betaling niet verplicht was.

\paragraph{Benadeling schuldeisers - overwegingen}\label{benadeling-schuldeisers---overwegingen-39}

De betaling van € 552.500,00 door Concepts aan Profundo op 1 augustus 2019 maakt onderdeel uit van een kasronde binnen de groep, waarbij Concepts een bedrag ontving van Construction (€ 552.500,00) en dit direct doorbetaalde aan Profundo. Hoewel de betaling op zichzelf mogelijk tegen een reële waarde heeft plaatsgevonden (er is geen indicatie dat de lening niet bestond of dat de aflossing onevenredig was), is er sprake van benadeling van de schuldeisers in de zin van artikel 42 Fw om de volgende redenen:

\begin{enumerate}
\def\labelenumi{\arabic{enumi}.}
\tightlist
\item
  \textbf{Vermogensvergelijking en middellijke benadeling}:

  \begin{itemize}
  \tightlist
  \item
    Zonder de kasronde zou Concepts het bedrag van € 552.500,00 in haar eigen vermogen hebben behouden, wat ten goede zou zijn gekomen aan \emph{alle} schuldeisers in de boedel. Door de betaling aan Profundo is dit bedrag uit de boedel verdwenen en ten goede gekomen aan \emph{één} schuldeiser (Profundo), terwijl andere schuldeisers hierdoor in hun verhaalsmogelijkheden zijn beperkt. Dit vormt een \textbf{verstoring van de rangorde tussen schuldeisers} (paritas creditorum), wat voldoende is voor benadeling in de zin van artikel 42 Fw, ook als per saldo geen vermogensvermindering optreedt (vgl. ECLI:NL:RBROT:2008:BC6115, r.o. 4.4).
  \item
    De curator heeft niet weersproken dat Concepts ná de betaling nog voor ca. € 2,5 miljoen aan financiering heeft ontvangen en ca. € 3,3 miljoen aan schuldeisers heeft betaald (3.11 en 5.22). Dit betekent echter niet dat de betaling aan Profundo \emph{niet} benadelend was: de benadeling moet worden beoordeeld op het moment van de rechtshandeling (leer van de middellijke benadeling). Op 1 augustus 2019 was de betaling aan Profundo een vermogensoverdracht die de boedel verminderde ten opzichte van de hypothetische situatie waarin het bedrag in de boedel was gebleven.
  \end{itemize}
\item
  \textbf{Geen compenserende voordelen voor de boedel}:

  \begin{itemize}
  \tightlist
  \item
    Er is geen aantoonbaar voordeel voor de boedel van Concepts als geheel. De kasronde was een interne groepstransactie die geen nieuwe waarde toevoegde aan Concepts, maar slechts een herverdeling van middelen binnen de groep bewerkstelligde. De betaling aan Profundo was derhalve \textbf{onverplicht} in de zin van artikel 42 lid 1 Fw, omdat Concepts niet verplicht was om de ontvangen gelden direct door te betalen aan Profundo in plaats van deze te behouden voor haar eigen schuldeisers.
  \end{itemize}
\item
  \textbf{Vergelijking met relevante jurisprudentie}:

  \begin{itemize}
  \tightlist
  \item
    In ECLI:NL:RBAMS:2023:897 (r.o. 1.12) oordeelde de rechtbank dat een wegboeking van een rekening-courantvordering paulianeus was omdat deze leidde tot benadeling van schuldeisers door verstoring van de paritas creditorum, ook al was er geen sprake van een directe vermogensvermindering. Dit is vergelijkbaar met de onderhavige casus, waar de betaling aan Profundo de boedel ten nadele van andere schuldeisers verminderde.
  \item
    In ECLI:NL:RBLEE:2004:AR1614 (r.o. 4.2-4.3) oordeelde de rechtbank dat betalingen die voortvloeiden uit een afspraak tussen partijen \emph{niet} paulianeus waren, omdat zonder die betalingen de activa en passiva van de schuldenaar gelijk zouden zijn gebleven. In de onderhavige casus is dit \emph{niet} het geval: zonder de betaling aan Profundo zou het bedrag van € 552.500,00 in de boedel zijn gebleven en ten goede zijn gekomen aan \emph{alle} schuldeisers, in plaats van slechts aan Profundo.
  \end{itemize}
\item
  \textbf{Subjectieve voorwaarden (wetenschap van benadeling)}:

  \begin{itemize}
  \tightlist
  \item
    Hoewel de benadeling objectief vaststaat, moet voor toepassing van artikel 42 lid 2 Fw (aangezien dit een meerzijdige rechtshandeling betreft) ook worden aangetoond dat Profundo wist of behoorde te weten dat de betaling benadeling van schuldeisers tot gevolg zou hebben. De curator heeft hiertoe echter geen concrete feiten gesteld (bijv. dat Profundo op de hoogte was van de slechte financiële positie van Concepts). Dit is echter \emph{niet} relevant voor de beoordeling of de schuldeisers \emph{objectief} zijn benadeeld -- dat is wel het geval. De subjectieve voorwaarden zijn van belang voor de \emph{vernietigbaarheid} van de handeling, maar niet voor de vaststelling of er sprake is van benadeling.
  \end{itemize}
\end{enumerate}

Conclusie: De schuldeisers zijn objectief benadeeld doordat de betaling aan Profundo de boedel heeft verminderd en de paritas creditorum heeft verstoord. Dit voldoet aan de definitie van benadeling in artikel 42 Fw.

\paragraph{Wetenschap - overwegingen}\label{wetenschap---overwegingen-39}

De betaling van € 552.500,00 door Concepts aan Profundo op 1 augustus 2019 valt binnen de termijn van één jaar voor de faillietverklaring (4 augustus 2020) en is niet reeds verplicht voor de aanvang van die termijn (de lening was niet opeisbaar op dat moment, art. 4 luidt dat de lening dadelijk opeisbaar is bij faillissement, maar niet daarvoor). De betaling is verricht tussen groepsmaatschappijen (Concepts en Profundo behoren tot dezelfde groep, art. 43 lid 1 sub 6 Fw), waardoor het bewijsvermoeden van wetenschap van benadeling aan beide zijden (art. 43 lid 1 Fw) van toepassing is. Profundo heeft geen succesvol tegenbewijs geleverd: de stelling dat de kasronde bedoeld was om het eigen vermogen te verstevigen en de onderneming voort te zetten, is onvoldoende om het vermoeden te ontzenuwen, gelet op de latere faillissementen en de omstandigheid dat Concepts na 1 augustus 2019 nog € 2,5 miljoen aan financiering ontving en € 3,3 miljoen aan schuldeisers betaalde, maar uiteindelijk alsnog failliet ging. De wetenschap van benadeling wordt dus vermoed en is niet weerlegd.

\newpage

\markboth{ECLI:NL:RBZWB:2026:4349}{gpt-5.5 (k = 5)}

\subsubsection{GPT}\label{gpt-13}

\texttt{gpt-5.5}

\begin{center}
\small
\begin{tabular}{@{}p{7.4cm}p{2.6cm}p{2.6cm}@{}}
\toprule
Variabele & Waarde (model) & Waarde (gold) \\
\midrule
Is rechtshandeling & \texttt{true} & \texttt{true} \\
Datum rechtshandeling & \texttt{2019-08-01} & \texttt{2019-08-01} \\
Verricht door de schuldenaar & \texttt{true} & \texttt{true} \\
Om niet & \texttt{false} & \texttt{false} \\
Datum faillietverklaring & \texttt{2020-08-04} & \texttt{2020-08-04} \\
Onverplicht & \texttt{true} & \texttt{true} \\
Benadeling schuldeisers & \textcolor{gray}{\texttt{true}} & \texttt{false} \\
Wetenschap schuldenaar aangetoond & \texttt{false} & \texttt{false} \\
Wetenschap wederpartij aangetoond & \texttt{false} & \texttt{false} \\
Vermoeden wetenschap schuldenaar tegenbewijs & \textcolor{gray}{\texttt{true}} & \texttt{null} \\
Vermoeden wetenschap wederpartij tegenbewijs & \textcolor{gray}{\texttt{true}} & \texttt{null} \\
Tegen onderwaarde & \texttt{null} & \texttt{null} \\
Voldoening of zekerheid niet opeisbare schuld & \textcolor{gray}{\texttt{true}} & \texttt{null} \\
Gelieerde partij & \textcolor{gray}{\texttt{true}} & \texttt{null} \\
Reeds verplicht voor aanvang termijn & \texttt{false} & \texttt{null} \\
\textbf{Vernietigbaar} & \textbf{\texttt{false}} & \textbf{\texttt{false}} \\
\bottomrule
\end{tabular}
\end{center}

\paragraph{Is rechtshandeling - overwegingen}\label{is-rechtshandeling---overwegingen-40}

De gewraakte handeling is de betaling door Concepts van € 552.500,00 aan Profundo ter terugbetaling op de lening. Een betaling ter voldoening of vermindering van een schuld is gericht op rechtsgevolg, namelijk nakoming en vermindering/tenietgaan van de schuld.

\paragraph{Verricht door de schuldenaar - overwegingen}\label{verricht-door-de-schuldenaar---overwegingen-40}

Uit de feiten volgt dat Concepts zelf het bedrag van € 552.500,00 aan Profundo heeft betaald. Concepts is de later failliet verklaarde schuldenaar.

\paragraph{Om niet - overwegingen}\label{om-niet---overwegingen-40}

De betaling vond plaats als terugbetaling op de lening van Concepts aan Profundo. Daartegenover stond de vermindering van de schuld van Concepts aan Profundo, zodat de handeling niet om niet maar om baat was.

\paragraph{Onverplicht - overwegingen}\label{onverplicht---overwegingen-40}

De aangevochten rechtshandeling is de betaling door Concepts aan Profundo van € 552.500 op 1 augustus 2019. Weliswaar bestond er sinds de overeenkomst van geldlening van 13 juli 2017 een schuld van Concepts aan Profundo, maar uit die overeenkomst volgt geen op dat moment in rechte afdwingbare verplichting tot deze aflossing. Aan de lening was geen looptijd verbonden en er waren geen aflossingstermijnen overeengekomen. Profundo kon de lening slechts opzeggen met een opzegtermijn van zes maanden, maar uit de casus blijkt niet dat vóór 1 augustus 2019 was opgezegd. Evenmin blijkt dat een in artikel 4 genoemde opeisingsgrond zich toen voordeed. De betaling vond plaats in het kader van een groepsinterne kasronde, maar een feitelijke, economische of groepsbelangmatige aanleiding schept geen rechtsplicht. Anders dan in gevallen waarin een opeisbare schuld wordt betaald, zoals in ECLI:NL:GHARL:2016:3413, of waarin betaling binnen een overeengekomen betalingstermijn plaatsvindt, zoals in ECLI:NL:GHSHE:2015:1681, ging het hier om vervroegde aflossing van een niet-opeisbare langlopende lening. Dat is in de zin van artikel 42 Fw onverplicht.

\paragraph{Benadeling schuldeisers - overwegingen}\label{benadeling-schuldeisers---overwegingen-40}

Uitgaande van de vermogensvergelijking is sprake van benadeling. Door de betaling van € 552.500 aan Profundo is weliswaar de schuld van Concepts aan Profundo met hetzelfde bedrag verminderd, zodat op balansniveau niet zonder meer een netto-vermogensvermindering ontstaat. Maar de betaling had tot gevolg dat Profundo vóór faillissement feitelijk voor € 552.500 werd voldaan, terwijl andere schuldeisers op de faillissementsboedel zijn aangewezen. Zonder de betaling zou Concepts het in de kasronde van Construction ontvangen bedrag hebben behouden, althans zou de daarmee corresponderende verhaalswaarde voor de gezamenlijke schuldeisers beschikbaar zijn gebleven, en zou Profundo voor dat bedrag als concurrente schuldeiser in de rangorde hebben moeten delen. Daarmee is de paritas creditorum verstoord en zijn de overige schuldeisers in hun verhaalsmogelijkheden benadeeld. Dat Concepts na 1 augustus 2019 nog financiering heeft ontvangen en schuldeisers heeft betaald, neemt deze benadeling niet weg: daaruit volgt niet dat uiteindelijk alle schuldeisers volledig konden worden voldaan of dat het betaalde bedrag zonder de rechtshandeling niet aan de boedel ten goede zou zijn gekomen. Er zijn ook geen feiten gesteld waaruit volgt dat de gezamenlijke schuldeisers zich hoe dan ook niet op dit bedrag of de daarmee corresponderende waarde hadden kunnen verhalen, bijvoorbeeld wegens zekerheidsrechten. Vergelijkbaar met de lijn in onder meer ECLI:NL:GHDHA:2013:150 kan een betaling aan een groepsvennootschap voorafgaand aan faillissement benadeling opleveren doordat één schuldeiser gunstiger wordt behandeld dan de overige schuldeisers.

\paragraph{Wetenschap - overwegingen}\label{wetenschap---overwegingen-40}

Rechtstreekse wetenschap van benadeling is niet aangetoond. De maatstaf is of ten tijde van de betaling het faillissement en een tekort daarin met redelijke mate van waarschijnlijkheid waren te voorzien. Profundo heeft onweersproken aangevoerd dat Concepts na 1 augustus 2019 nog circa € 2,5 miljoen financiering ontving, nog circa € 3,3 miljoen aan schuldeisers betaalde, dat pas in april 2020 werd besloten geen aanvullende financiering meer te verstrekken en dat de onderneming daarna nog maanden is voortgezet. Dat wijst erop dat faillissement met tekort op het moment van de kasronde niet redelijkerwijs voorzienbaar was. Wel zijn art. 43-gronden in beeld: de betaling strekte tot aflossing van een nog niet opeisbare langlopende lening aan een groepsmaatschappij/gelieerde partij, zonder eerdere verplichting tot deze aflossing. Voor zover daardoor wetenschap aan beide zijden wordt vermoed, is dat vermoeden door de genoemde onweersproken omstandigheden ontzenuwd. Art. 45 Fw is niet aan de orde bij deze betaling ter aflossing van een lening.

\newpage

\markboth{ECLI:NL:RBZWB:2026:4349}{gemini-3.5-flash (k = 5)}

\subsubsection{Gemini}\label{gemini-13}

\texttt{gemini-3.5-flash}

\begin{center}
\small
\begin{tabular}{@{}p{7.4cm}p{2.6cm}p{2.6cm}@{}}
\toprule
Variabele & Waarde (model) & Waarde (gold) \\
\midrule
Is rechtshandeling & \texttt{true} & \texttt{true} \\
Datum rechtshandeling & \texttt{2019-08-01} & \texttt{2019-08-01} \\
Verricht door de schuldenaar & \texttt{true} & \texttt{true} \\
Om niet & \texttt{false} & \texttt{false} \\
Datum faillietverklaring & \texttt{2020-08-04} & \texttt{2020-08-04} \\
Onverplicht & \texttt{true} & \texttt{true} \\
Benadeling schuldeisers & \textcolor{gray}{\texttt{true}} & \texttt{false} \\
Wetenschap schuldenaar aangetoond & \texttt{false} & \texttt{false} \\
Wetenschap wederpartij aangetoond & \texttt{false} & \texttt{false} \\
Vermoeden wetenschap schuldenaar tegenbewijs & \texttt{null} & \texttt{null} \\
Vermoeden wetenschap wederpartij tegenbewijs & \texttt{null} & \texttt{null} \\
Tegen onderwaarde & \texttt{null} & \texttt{null} \\
Voldoening of zekerheid niet opeisbare schuld & \texttt{null} & \texttt{null} \\
Gelieerde partij & \texttt{null} & \texttt{null} \\
Reeds verplicht voor aanvang termijn & \texttt{null} & \texttt{null} \\
\textbf{Vernietigbaar} & \textbf{\texttt{false}} & \textbf{\texttt{false}} \\
\bottomrule
\end{tabular}
\end{center}

\paragraph{Is rechtshandeling - overwegingen}\label{is-rechtshandeling---overwegingen-41}

De betaling van een bedrag van € 552.500,00 door Concepts aan Profundo ter aflossing van een lening is een rechtshandeling, aangezien deze handeling gericht is op het rechtsgevolg van het tenietgaan van een verbintenis (nakoming).

\paragraph{Verricht door de schuldenaar - overwegingen}\label{verricht-door-de-schuldenaar---overwegingen-41}

De betaling is rechtstreeks verricht door Concepts, de schuldenaar die naderhand failliet is verklaard.

\paragraph{Om niet - overwegingen}\label{om-niet---overwegingen-41}

De betaling betreft de aflossing van een bestaande schuld uit een overeenkomst van geldlening tussen Concepts en Profundo. Omdat de betaling strekt tot het verminderen van een schuldpositie, is er sprake van een rechtshandeling anders dan om niet (om baat).

\paragraph{Onverplicht - overwegingen}\label{onverplicht---overwegingen-41}

De betaling van € 552.500,00 door Concepts aan Profundo op 1 augustus 2019 is een onverplichte rechtshandeling. Uit de overeenkomst van geldlening van 13 juli 2017 volgt dat aan de looptijd van de lening geen termijn was verbonden en dat er geen aflossingstermijnen golden. Profundo kon de lening slechts opzeggen met inachtneming van een opzegtermijn van zes maanden (artikel 2). Gesteld noch gebleken is dat een dergelijke opzegging heeft plaatsgevonden voorafgaand aan de betaling. Tevens deed zich op 1 augustus 2019 geen van de in artikel 4 van de overeenkomst genoemde situaties voor die tot directe opeisbaarheid zouden leiden (het faillissement van Concepts werd pas op 4 augustus 2020 uitgesproken). De schuld was op het moment van betaling derhalve niet opeisbaar. Het voldoen van een niet-opeisbare schuld is onverplicht, aangezien de prestatie op dat moment niet in rechte kon worden afgedwongen (vergelijk ECLI:NL:RBOVE:2021:73 en in tegenstelling tot ECLI:NL:GHARL:2016:3413, waarin de lening wel reeds opeisbaar was).

\paragraph{Benadeling schuldeisers - overwegingen}\label{benadeling-schuldeisers---overwegingen-41}

Om te bepalen of de schuldeisers zijn benadeeld in de zin van artikel 42 Fw, moet een vermogensvergelijking worden gemaakt tussen de feitelijke situatie waarin de schuldeisers verkeren na de betaling van € 552.500,00 aan Profundo, en de hypothetische situatie waarin deze betaling achterwege zou zijn gebleven. Indien de betaling niet had plaatsgevonden, had Concepts de van Construction ontvangen € 552.500,00 (of de vordering op Construction, indien de gehele kasronde wordt weggedacht) behouden als actief. Dit bedrag zou dan beschikbaar zijn geweest voor de gezamenlijke schuldeisers. Door de betaling is dit actieve vermogen echter aan de boedel onttrokken ten gunste van één concurrente schuldeiser (Profundo), die daardoor voor dit deel voor 100\% is voldaan. Dit levert een directe vermindering van het voor verhaal beschikbare actief op en een verstoring van de paritas creditorum. De benadeling moet worden beoordeeld op het moment dat de curator de vernietiging inroept (2024). Op dat moment is er sprake van een faillissementstekort. Het verweer van Profundo dat er na de transactie nog € 2,5 miljoen aan financiering is ontvangen en € 3,3 miljoen aan schuldeisers is betaald, doet niet af aan het feit dat de boedel ten tijde van de vernietiging € 552.500,00 aan actief mist dat anders voor de gezamenlijke schuldeisers beschikbaar zou zijn geweest. Er is derhalve sprake van benadeling van de schuldeisers.

\paragraph{Wetenschap - overwegingen}\label{wetenschap---overwegingen-41}

De betaling van € 552.500,00 door Concepts aan Profundo vond plaats op 1 augustus 2019, terwijl het faillissement van Concepts op 4 augustus 2020 is uitgesproken. Omdat deze rechtshandeling meer dan een jaar vóór de faillietverklaring is verricht, zijn de wettelijke bewijsvermoedens van artikel 43 Fw en artikel 45 Fw niet van toepassing. De curator draagt derhalve de volledige en rechtstreekse bewijslast om de wetenschap van benadeling aan te tonen op grond van de hoofdregel van artikel 42 Fw. Uit de feiten blijkt dat Concepts na 1 augustus 2019 nog circa € 2,5 miljoen aan aanvullende financiering heeft ontvangen en voor € 3,3 miljoen aan schuldeisers heeft voldaan. Pas in april 2020 is besloten de financiering te staken. Ten tijde van de betaling op 1 augustus 2019 waren het faillissement en een tekort daarin dan ook niet met een redelijke mate van waarschijnlijkheid te voorzien. De wetenschap van benadeling is derhalve aan beide zijden niet aangetoond.

\newpage

\markboth{ECLI:NL:RBZWB:2026:4349}{mistral/mistral-medium-3.5-128b:fp8 (k = 10)}

\subsection{k = 10}\label{k-10-4}

\subsubsection{Mistral}\label{mistral-14}

\texttt{mistral/mistral-medium-3.5-128b:fp8}

\begin{center}
\small
\begin{tabular}{@{}p{7.4cm}p{2.6cm}p{2.6cm}@{}}
\toprule
Variabele & Waarde (model) & Waarde (gold) \\
\midrule
Is rechtshandeling & \texttt{true} & \texttt{true} \\
Datum rechtshandeling & \texttt{2019-08-01} & \texttt{2019-08-01} \\
Verricht door de schuldenaar & \texttt{true} & \texttt{true} \\
Om niet & \texttt{false} & \texttt{false} \\
Datum faillietverklaring & \texttt{2020-08-04} & \texttt{2020-08-04} \\
Onverplicht & \texttt{true} & \texttt{true} \\
Benadeling schuldeisers & \textcolor{gray}{\texttt{true}} & \texttt{false} \\
Wetenschap schuldenaar aangetoond & \texttt{false} & \texttt{false} \\
Wetenschap wederpartij aangetoond & \texttt{false} & \texttt{false} \\
Vermoeden wetenschap schuldenaar tegenbewijs & \texttt{false} & \texttt{null} \\
Vermoeden wetenschap wederpartij tegenbewijs & \texttt{false} & \texttt{null} \\
Tegen onderwaarde & \texttt{null} & \texttt{null} \\
Voldoening of zekerheid niet opeisbare schuld & \texttt{null} & \texttt{null} \\
Gelieerde partij & \textcolor{gray}{\texttt{true}} & \texttt{null} \\
Reeds verplicht voor aanvang termijn & \texttt{false} & \texttt{null} \\
\textbf{Vernietigbaar} & \textbf{\texttt{false}} & \textbf{\texttt{false}} \\
\bottomrule
\end{tabular}
\end{center}

\paragraph{Is rechtshandeling - overwegingen}\label{is-rechtshandeling---overwegingen-42}

De betaling van € 552.500,00 door Concepts aan Profundo is een rechtshandeling, omdat het een op een rechtsgevolg gerichte wil (betaling) betreft die zich door een verklaring (de daadwerkelijke betaling) heeft geopenbaard.

\paragraph{Verricht door de schuldenaar - overwegingen}\label{verricht-door-de-schuldenaar---overwegingen-42}

De betaling is verricht door Concepts, de latere gefailleerde, en dus door de schuldenaar zelf.

\paragraph{Om niet - overwegingen}\label{om-niet---overwegingen-42}

De betaling van € 552.500,00 door Concepts aan Profundo maakt deel uit van een kasronde binnen de groep, waarbij Concepts een tegenprestatie ontving (namelijk de aflossing van haar rekening-courantschuld door Construction). Er was dus sprake van een reële tegenprestatie, waardoor de handeling niet `om niet' is.

\paragraph{Onverplicht - overwegingen}\label{onverplicht---overwegingen-42}

De betaling van € 552.500,00 door Concepts aan Profundo op 1 augustus 2019 is onverplicht verricht in de zin van artikel 42 lid 1 Fw.

Uit de casus blijkt dat Concepts en Profundo op 13 juli 2017 een overeenkomst van geldlening hebben gesloten, waarbij een rekening-courantverhouding werd omgezet in een langlopende lening (3.5). In artikel 2 van deze overeenkomst is bepaald dat aan de looptijd van de lening geen termijn is verbonden en dat Profundo de lening met inachtneming van een opzegtermijn van 6 maanden kan opzeggen (3.6). Er zijn geen aflossingstermijnen overeengekomen.

In artikel 4 van de overeenkomst is wel bepaald dat de leningstand (inclusief lopende rente) dadelijk kan worden opgeëist in bepaalde gevallen, zoals bij faillissement of surséance van betaling van Concepts (3.7). Echter, op 1 augustus 2019 was geen van deze gevallen aan de orde. Derhalve bestond er op dat moment geen opeisbare verplichting voor Concepts om de lening af te lossen.

De betaling van € 552.500,00 was onderdeel van een kasronde binnen de groep (3.10), maar deze kasronde creëerde geen nieuwe verplichting voor Concepts om deze betaling te verrichten. De lening was immers niet opeisbaar op dat moment, en er bestond geen afspraak over aflossingstermijnen.

Gelet op hetgeen in de aandachtspunten is vermeld, is het betalen van een niet-opeisbare schuld onverplicht. Aangezien de lening op 1 augustus 2019 niet opeisbaar was, was de betaling derhalve onverplicht.

\paragraph{Benadeling schuldeisers - overwegingen}\label{benadeling-schuldeisers---overwegingen-42}

De betaling van € 552.500,00 door Concepts aan Profundo op 1 augustus 2019 in het kader van de kasronde binnen de groep moet worden beoordeeld op benadeling van de schuldeisers in de zin van artikel 42 Fw.

Allereerst wordt de hypothetische situatie vergeleken met de feitelijke situatie. Zonder de betaling van € 552.500,00 aan Profundo zou dit bedrag nog in de boedel van Concepts aanwezig zijn geweest. Dit bedrag had dan kunnen worden gebruikt om de schuldeisers van Concepts te voldoen. In de feitelijke situatie is dit bedrag echter aan Profundo betaald, waardoor het actief van Concepts is verminderd en de schuldeisers in hun verhaalsmogelijkheden zijn benadeeld.

Daarnaast is van belang dat de betaling deel uitmaakte van een kasronde binnen de groep, waarbij geldstromen tussen verschillende vennootschappen zijn verplaatst. Dit kan erop duiden dat de betaling niet in het belang van Concepts zelf was, maar veeleer in het belang van de groep als geheel of van andere vennootschappen binnen de groep. Dit versterkt de conclusie dat de schuldeisers van Concepts zijn benadeeld.

Tot slot is relevant dat Concepts na 1 augustus 2019 nog ongeveer € 2,5 miljoen aan financiering heeft ontvangen en nog voor een bedrag van ongeveer € 3,3 miljoen aan schuldeisers heeft betaald (ECLI:NL:RBZWB:2025:6049, r.o. 5.22). Dit betekent echter niet dat de betaling van € 552.500,00 aan Profundo geen benadeling heeft veroorzaakt. De benadeling moet immers worden beoordeeld op het moment waarop de curator zijn rechten doet gelden, en niet op het moment van de betaling zelf. Het feit dat Concepts later nog betalingen heeft verricht, doet niet af aan het feit dat de betaling aan Profundo het actief van Concepts op dat moment heeft verminderd.

Gelet op het voorgaande is er sprake van benadeling van de schuldeisers in de zin van artikel 42 Fw.

\paragraph{Wetenschap - overwegingen}\label{wetenschap---overwegingen-42}

De betaling van € 552.500,00 door Concepts aan Profundo op 1 augustus 2019 is verricht binnen een jaar voor de faillietverklaring van Concepts (4 augustus 2020) en Concepts had zich niet reeds voor de aanvang van die termijn tot deze betaling verplicht. Profundo en Concepts behoren tot dezelfde groep vennootschappen, waarbij {[}naam 1{]} Holding statutair bestuurder is van Profundo en {[}naam 1{]} via zijn persoonlijke holding grootaandeelhouder is van de groep. Concepts is een rechtspersoon en Profundo is een andere rechtspersoon binnen dezelfde groep. Op grond van art. 43 lid 1 sub 6 Fw wordt de wetenschap van benadeling aan beide zijden (Concepts en Profundo) vermoed. Profundo heeft geen succesvol tegenbewijs geleverd om dit vermoeden te ontzenuwen. De omstandigheid dat Concepts na 1 augustus 2019 nog financiering heeft ontvangen en schuldeisers heeft betaald, doet niet af aan het vermoeden, nu het faillissement en een tekort daarin met een redelijke mate van waarschijnlijkheid waren te voorzien (ECLI:NL:RBOVE:2021:73).

\newpage

\markboth{ECLI:NL:RBZWB:2026:4349}{gpt-5.5 (k = 10)}

\subsubsection{GPT}\label{gpt-14}

\texttt{gpt-5.5}

\begin{center}
\small
\begin{tabular}{@{}p{7.4cm}p{2.6cm}p{2.6cm}@{}}
\toprule
Variabele & Waarde (model) & Waarde (gold) \\
\midrule
Is rechtshandeling & \texttt{true} & \texttt{true} \\
Datum rechtshandeling & \texttt{2019-08-01} & \texttt{2019-08-01} \\
Verricht door de schuldenaar & \texttt{true} & \texttt{true} \\
Om niet & \texttt{false} & \texttt{false} \\
Datum faillietverklaring & \texttt{2020-08-04} & \texttt{2020-08-04} \\
Onverplicht & \texttt{true} & \texttt{true} \\
Benadeling schuldeisers & \textcolor{gray}{\texttt{true}} & \texttt{false} \\
Wetenschap schuldenaar aangetoond & \texttt{false} & \texttt{false} \\
Wetenschap wederpartij aangetoond & \texttt{false} & \texttt{false} \\
Vermoeden wetenschap schuldenaar tegenbewijs & \textcolor{gray}{\texttt{true}} & \texttt{null} \\
Vermoeden wetenschap wederpartij tegenbewijs & \textcolor{gray}{\texttt{true}} & \texttt{null} \\
Tegen onderwaarde & \texttt{null} & \texttt{null} \\
Voldoening of zekerheid niet opeisbare schuld & \textcolor{gray}{\texttt{true}} & \texttt{null} \\
Gelieerde partij & \textcolor{gray}{\texttt{true}} & \texttt{null} \\
Reeds verplicht voor aanvang termijn & \texttt{false} & \texttt{null} \\
\textbf{Vernietigbaar} & \textbf{\texttt{false}} & \textbf{\texttt{false}} \\
\bottomrule
\end{tabular}
\end{center}

\paragraph{Is rechtshandeling - overwegingen}\label{is-rechtshandeling---overwegingen-43}

De gewraakte handeling is de betaling door Concepts van € 552.500,00 aan Profundo ter terugbetaling op de lening. Een betaling strekt tot rechtsgevolg, namelijk vermindering of voldoening van de schuld, en is daarmee een rechtshandeling.

\paragraph{Verricht door de schuldenaar - overwegingen}\label{verricht-door-de-schuldenaar---overwegingen-43}

Uit de feiten volgt dat Concepts zelf het bedrag van € 552.500,00 aan Profundo heeft betaald. Concepts is de later failliet verklaarde schuldenaar.

\paragraph{Om niet - overwegingen}\label{om-niet---overwegingen-43}

De betaling vond plaats als terugbetaling op de lening van Concepts aan Profundo. Tegenover de betaling stond vermindering van de schuld van Concepts aan Profundo; de handeling is daarom anders dan om niet.

\paragraph{Onverplicht - overwegingen}\label{onverplicht---overwegingen-43}

De aangevochten rechtshandeling is de betaling door Concepts aan Profundo van € 552.500,00 op 1 augustus 2019 ter aflossing op de in 2017 omgezette rekening-courantschuld/langlopende lening. Voor de vraag of deze betaling onverplicht was, is beslissend of Profundo op dat moment in rechte betaling kon afdwingen. Uit de leningsovereenkomst volgt dat aan de looptijd geen termijn was verbonden, dat Profundo slechts met een opzegtermijn van zes maanden kon opzeggen en dat geen aflossingstermijnen waren overeengekomen. De onmiddellijke opeisingsgronden van artikel 4 van de overeenkomst zagen onder meer op faillissement/surseance, akkoord, tekortkoming na ingebrekestelling of beslag, maar uit de casus blijkt niet dat één van deze gronden op 1 augustus 2019 was ingetreden, noch dat Profundo de lening had opgezegd. De schuld bestond dus wel, maar was op het moment van betaling niet opeisbaar. Het betalen van een niet-opeisbare schuld is in de zin van artikel 42 Fw onverplicht. Dat de betaling plaatsvond in het kader van een groepsinterne kasronde en mogelijk werd ingegeven door bedrijfseconomische of groepsbelangen, schept geen rechtsplicht. Dit verschilt van gevallen waarin een lening reeds contractueel opeisbaar was, zoals in ECLI:NL:GHARL:2016:3413, of waarin binnen een overeengekomen betalingstermijn werd betaald, zoals in ECLI:NL:GHSHE:2015:1681.

\paragraph{Benadeling schuldeisers - overwegingen}\label{benadeling-schuldeisers---overwegingen-43}

De benadeling moet worden beoordeeld door de feitelijke situatie bij instandhouding van de betaling van € 552.500 aan Profundo te vergelijken met de hypothetische situatie zonder die betaling, per het moment waarop de curator de pauliana inroept. Concepts had uit de kasronde een betaling van Construction ontvangen ter aflossing van een rekening-courantvordering. Door het bedrag vervolgens aan Profundo door te betalen, is dit verhaalsactief aan Concepts onttrokken en is een groepscrediteur voldaan op een langlopende lening waarvoor geen aflossingstermijnen golden en die op dat moment niet opeisbaar lijkt te zijn. Daartegenover staat wel dat de schuld aan Profundo met hetzelfde bedrag is verminderd, zodat boekhoudkundig geen zuivere vermogensvermindering optreedt. Dat neemt de benadeling echter niet weg: de betaling leidt ertoe dat Profundo vóór faillissement feitelijk voor € 552.500 is voldaan, terwijl andere schuldeisers in het faillissement onbetaald blijven en slechts naar rang en evenredigheid verhaal kunnen nemen. Dit is een verstoring van de paritas creditorum/rangorde en levert benadeling op in de zin van artikel 42 Fw. De omstandigheid dat Concepts na 1 augustus 2019 nog financiering heeft ontvangen en aanzienlijke bedragen aan schuldeisers heeft betaald, maakt dit niet anders; daaruit volgt niet dat zonder de betaling aan Profundo alle schuldeisers zouden zijn voldaan of dat het bedrag onvermijdelijk op een voor de gezamenlijke schuldeisers gelijkwaardige wijze zou zijn aangewend. Er zijn ook geen aanwijzingen dat de schuldeisers zich niet op het betrokken bedrag hadden kunnen verhalen wegens zekerheidsrechten of dat slechts een voordeel is misgelopen. Vergelijkbaar met de benadering in onder meer ECLI:NL:RBROT:2008:BC6115 en ECLI:NL:GHDHA:2013:150 kan selectieve voldoening van een groepscrediteur vlak vóór faillissement benadeling opleveren doordat het actief voor de gezamenlijke schuldeisers wordt verminderd dan wel de onderlinge rangorde wordt verstoord.

\paragraph{Wetenschap - overwegingen}\label{wetenschap---overwegingen-43}

De betaling aan Profundo was een aflossing op een langlopende groepslening zonder aflossingstermijnen en zonder gebleken voorafgaande opzegging; daarmee gaat het materieel om voldoening van een niet-opeisbare schuld aan een gelieerde/groepsmaatschappij. Er was geen reeds bestaande verplichting tot deze vervroegde aflossing. Rechtstreekse wetenschap van benadeling is echter niet aangetoond: onweersproken is dat Concepts na 1 augustus 2019 nog circa € 2,5 miljoen financiering ontving, circa € 3,3 miljoen aan schuldeisers betaalde, dat pas in april 2020 werd besloten geen aanvullende financiering meer te verstrekken en dat de onderneming daarna nog maanden is voortgezet. Daarmee was ten tijde van de betaling niet met een redelijke mate van waarschijnlijkheid voorzienbaar dat een faillissement met tekort zou volgen. Voor zover een bewijsvermoeden van art. 43 Fw zou worden toegepast, is dat vermoeden aan beide zijden voldoende ontzenuwd. Art. 45 Fw is niet aan de orde, omdat de betaling niet als rechtshandeling om niet kwalificeert.

\newpage

\markboth{ECLI:NL:RBZWB:2026:4349}{gemini-3.5-flash (k = 10)}

\subsubsection{Gemini}\label{gemini-14}

\texttt{gemini-3.5-flash}

\begin{center}
\small
\begin{tabular}{@{}p{7.4cm}p{2.6cm}p{2.6cm}@{}}
\toprule
Variabele & Waarde (model) & Waarde (gold) \\
\midrule
Is rechtshandeling & \texttt{true} & \texttt{true} \\
Datum rechtshandeling & \texttt{2019-08-01} & \texttt{2019-08-01} \\
Verricht door de schuldenaar & \texttt{true} & \texttt{true} \\
Om niet & \texttt{false} & \texttt{false} \\
Datum faillietverklaring & \texttt{2020-08-04} & \texttt{2020-08-04} \\
Onverplicht & \texttt{true} & \texttt{true} \\
Benadeling schuldeisers & \textcolor{gray}{\texttt{true}} & \texttt{false} \\
Wetenschap schuldenaar aangetoond & \texttt{false} & \texttt{false} \\
Wetenschap wederpartij aangetoond & \texttt{false} & \texttt{false} \\
Vermoeden wetenschap schuldenaar tegenbewijs & \textcolor{gray}{\texttt{true}} & \texttt{null} \\
Vermoeden wetenschap wederpartij tegenbewijs & \textcolor{gray}{\texttt{true}} & \texttt{null} \\
Tegen onderwaarde & \texttt{null} & \texttt{null} \\
Voldoening of zekerheid niet opeisbare schuld & \textcolor{gray}{\texttt{true}} & \texttt{null} \\
Gelieerde partij & \textcolor{gray}{\texttt{true}} & \texttt{null} \\
Reeds verplicht voor aanvang termijn & \texttt{false} & \texttt{null} \\
\textbf{Vernietigbaar} & \textbf{\texttt{false}} & \textbf{\texttt{false}} \\
\bottomrule
\end{tabular}
\end{center}

\paragraph{Is rechtshandeling - overwegingen}\label{is-rechtshandeling---overwegingen-44}

De betaling van een bedrag van € 552.500,00 door Concepts aan Profundo ter aflossing van een lening is een rechtshandeling, aangezien deze gericht is op het rechtsgevolg van het tenietgaan van een verbintenis (nakoming).

\paragraph{Verricht door de schuldenaar - overwegingen}\label{verricht-door-de-schuldenaar---overwegingen-44}

De betaling is rechtstreeks verricht door de schuldenaar, Concepts, aan haar schuldeiser Profundo.

\paragraph{Om niet - overwegingen}\label{om-niet---overwegingen-44}

De betaling betreft de aflossing van een lening (terugbetaling op de lening), waardoor een schuld van Concepts aan Profundo tenietgaat. Er staat derhalve een prestatie tegenover (bevrijding uit de schuld), waardoor de handeling om baat (anders dan om niet) is.

\paragraph{Onverplicht - overwegingen}\label{onverplicht---overwegingen-44}

De betaling van € 552.500,00 door Concepts aan Profundo op 1 augustus 2019 is onverplicht verricht. Een rechtshandeling is onverplicht wanneer daartoe geen op de wet of een overeenkomst berustende rechtsplicht bestond die in rechte kon worden afgedwongen. Het betalen van een niet-opeisbare schuld is in beginsel onverplicht. Uit de overeenkomst van geldlening van 13 juli 2017 blijkt dat aan de lening geen looptijd of aflossingstermijnen waren verbonden. Profundo kon de lening weliswaar opzeggen, maar met inachtneming van een opzegtermijn van 6 maanden, en er is niet gesteld of gebleken dat een dergelijke opzegging heeft plaatsgevonden. Daarnaast deden de in artikel 4 van de overeenkomst genoemde situaties voor directe opeisbaarheid (zoals faillissement of beslag) zich op 1 augustus 2019 niet voor. De schuld was op het moment van de betaling dan ook niet opeisbaar. De omstandigheid dat de betaling plaatsvond in het kader van een kasronde binnen de groep creëert bovendien geen in rechte afdwingbare rechtsplicht tot betaling.

\paragraph{Benadeling schuldeisers - overwegingen}\label{benadeling-schuldeisers---overwegingen-44}

De beoordeling van de benadeling van schuldeisers in de zin van artikel 42 Faillissementswet (Fw) vindt plaats door middel van een vermogensvergelijking. Hierbij wordt de hypothetische situatie waarin de schuldeisers zouden hebben verkeerd zonder de gewraakte rechtshandeling vergeleken met de feitelijke situatie waarin zij verkeren nu de handeling in stand blijft. De rechtshandeling betreft hier de betaling van € 552.500,00 door Concepts aan Profundo op 1 augustus 2019, als sluitstuk van een kasronde binnen de groep. In de feitelijke situatie is het actief van Concepts verminderd met € 552.500,00 (althans is de vordering op Construction tenietgegaan door betaling, waarna de ontvangen middelen direct zijn doorbetaald aan Profundo) en is de schuld aan Profundo met een gelijke omvang verminderd. In de hypothetische situatie waarin de betaling aan Profundo achterwege was gebleven, zou Concepts de van Construction ontvangen € 552.500,00 hebben behouden (of zou de vordering op Construction nog hebben bestaan). Deze middelen zouden dan beschikbaar zijn geweest voor de gezamenlijke schuldeisers. Nu er sprake is van een faillissement met een aanzienlijk tekort, leidt de betaling aan één specifieke schuldeiser (Profundo) tot een directe vermindering van het voor de gezamenlijke schuldeisers beschikbare actieve boedelvolume en een verstoring van de paritas creditorum (de rangorde en gelijkheid van schuldeisers). De omstandigheid dat Concepts na de betaling nog maanden is voortgezet, € 2,5 miljoen aan nieuwe financiering heeft ontvangen en € 3,3 miljoen aan andere crediteuren heeft betaald, heft deze benadeling niet op. De benadeling dient immers te worden beoordeeld op het moment dat de curator de vernietiging inroept (juli 2024), op welk moment het tekort in de boedel onverminderd aanwezig is. Zonder de betaling zou de boedel wegens het ontbreken van die € 552.500,00 thans ongunstiger gepositioneerd zijn. Er is derhalve sprake van benadeling van de schuldeisers.

\paragraph{Wetenschap - overwegingen}\label{wetenschap---overwegingen-44}

De betaling van € 552.500,- door Concepts aan haar groepsmaatschappij Profundo op 1 augustus 2019 vond materieel plaats ter voldoening van een niet-opeisbare schuld (art. 43 lid 1 sub 2 Fw), aangezien de lening geen aflossingstermijnen kende en niet was opgezegd met inachtneming van de opzegtermijn van 6 maanden. Er bestond geen eerdere verplichting tot deze betaling. Hoewel sprake is van een groepsmaatschappij (art. 43 lid 1 sub 6 Fw), is de wetenschap van benadeling aan beide zijden niet rechtstreeks aangetoond en zijn de bewijsvermoedens succesvol ontzenuwd. Uit de onweersproken feiten blijkt namelijk dat Concepts na de betaling nog circa € 2,5 miljoen aan financiering heeft ontvangen en € 3,3 miljoen aan schuldeisers heeft betaald, en dat de financiering pas in april 2020 werd gestaakt. Ten tijde van de rechtshandeling in augustus 2019 waren het faillissement en een tekort daarin dan ook niet met een redelijke mate van waarschijnlijkheid te voorzien.

\newpage

\markboth{JOR 2009/46}{Gold standard}

\section{JOR 2009/46}\label{jor-200946}

uitspraakdatum 02-08-2007

\subsection{Gold}\label{gold-5}

\begin{center}
\small
\begin{tabular}{@{}p{8cm}p{2.6cm}@{}}
\toprule
Variabele & Waarde \\
\midrule
Is rechtshandeling & \texttt{true} \\
Datum rechtshandeling & \texttt{2004-01-08} \\
Verricht door de schuldenaar & \texttt{true} \\
Om niet & \texttt{true} \\
Datum faillietverklaring & \texttt{2004-02-04} \\
Onverplicht & \texttt{true} \\
Benadeling schuldeisers & \texttt{true} \\
Wetenschap schuldenaar aangetoond & \texttt{false} \\
Wetenschap wederpartij aangetoond & \texttt{null} \\
Vermoeden wetenschap schuldenaar tegenbewijs & \texttt{false} \\
Vermoeden wetenschap wederpartij tegenbewijs & \texttt{null} \\
Tegen onderwaarde & \texttt{null} \\
Voldoening of zekerheid niet opeisbare schuld & \texttt{null} \\
Gelieerde partij & \texttt{null} \\
Reeds verplicht voor aanvang termijn & \texttt{null} \\
\textbf{Vernietigbaar} & \textbf{\texttt{true}} \\
\bottomrule
\end{tabular}
\end{center}

\subsubsection{Uitspraak}\label{uitspraak-5}

(\ldots; red.)

\paragraph{3. Beoordeling}\label{beoordeling}

3.1. Het gaat in deze zaak kort gezegd om het navolgende.

3.1.1 ECPE is in november 2000 opgericht en exploiteerde een onderneming die zich bezighield met het ontwerpen, leveren, installeren en onderhouden van commerciële telefoonnetwerken. De aandelen van ECPE worden gehouden door ECP (65\%) en Ketelaars (35\%). Ketelaars was bestuurder van ECPE.

3.1.2 In de bijzondere algemene vergadering van aandeelhouders van ECPE van 8 december 2003 is besloten dat mr. drs. A.J. Fontijn (hierna te noemen: Fontijn) en A. Vergeer zouden optreden als interim managers van ECPE. ECP heeft Ketelaars bij brief van 12 december 2003 op non-actief gesteld als bestuurder van ECPE.

3.1.3 Op 8, 9, 20 en 23 januari 2004 zijn vier bedragen, in het totaal € 70.217,83, overgeboekt van de bankrekening van ECPE naar de bankrekening van ECP.

3.1.4 Bij vonnis van 4 februari 2004 is ECPE in staat van faillissement verklaard met aanstelling van de curator als zodanig.

3.1.5 Bij brief van 8 april 2004 aan ECP heeft de curator aanspraak gemaakt op betaling van het bedrag van € 70.217,83. De curator heeft zich daarbij primair op het standpunt gesteld dat sprake is van onverschuldigde betaling door ECPE. Subsidiair heeft de curator de betalingen als paulianeus aangemerkt en op die grond de vernietiging van de desbetreffende betalingen ingeroepen.

3.2. Voorzover in hoger beroep nog van belang, vordert de curator in dit geding veroordeling van ECP tot betaling aan de curator van bedrag van € 70.217,83, te vermeerderen met de wettelijke rente vanaf 23 januari 2004. Aan zijn vordering heeft de curator ten grondslag gelegd dat de betalingen van ECPE aan ECP onverschuldigd hebben plaatsgevonden en dat sprake is van faillissementspauliana. Subsidiair heeft de curator aangevoerd dat ECP onrechtmatig heeft gehandeld jegens de schuldeisers van ECPE. In reconventie vordert ECP de opheffing van de door de curator ten laste van ECP gelegde beslagen.

3.3. In het bestreden vonnis heeft de rechtbank in conventie de hierboven weergegeven vordering van de curator toegewezen en in reconventie de vordering van ECP afgewezen. De rechtbank heeft daartoe kort gezegd overwogen dat de hierboven onder 3.1.3 genoemde overboekingen moeten worden aangemerkt als onverschuldigde betalingen en dat, voorzover ECP vervolgens betalingen heeft verricht aan derden ten behoeve van ECPE, aan ECP geen beroep toekomt op verrekening omdat ECP niet te goeder trouw was.

3.4.1 De tweede grief is gericht tegen het oordeel van de rechtbank dat de bedoelde overboekingen aangemerkt moeten worden als onverschuldigde betalingen. ECP voert aan dat de overboekingen slechts plaatsvonden op praktische gronden, dat wil zeggen ten einde te voorkomen dat schuldeisers van ECPE conservatoir beslag zouden leggen. ECPE behield de beschikking over de overgemaakte gelden in die zin dat ECP de betreffende gelden, althans voorzover er een positief saldo resteerde op de rekeningcourant tussen ECP en ECPE, voor ECPE is gaan houden. Feitelijk functioneerde ECP aldus slechts als bank voor ECPE, aldus ECP.

3.4.2 De door ECP aangevoerde motieven voor de overboekingen vormen geen rechtsgrond voor betaling door een ECPE aan ECP. De stelling van ECP dat zij de overgemaakte gelden voor ECP is gaan houden voorzover er een positief saldo resteerde op de rekeningcourant tussen ECP en ECPE, bevestigt dat er voor deze overboekingen in zoverre geen rechtsgrond bestond. Anders dan ECP kennelijk meent, hebben de overboekingen ten gevolge gehad dat het in het totaal overgeboekte bedrag van het vermogen van ECPE naar het vermogen van ECP is overgegaan, zoals ook volgt uit het hieronder nader te bespreken standpunt van ECP dat zij betalingen aan schuldeisers van ECPE heeft gedaan en uit dien hoofde een vordering op ECPE heeft verkregen. Het hof onderschrijft dus het door de rechtbank onder 6.3 gegeven oordeel en de daaraan ten grondslag gelegde motivering.

3.4.3 Voorzover ECP zich op het standpunt stelt dat de betalingen van ECPE aan ECP strekten tot voldoening van een schuld van ECPE aan ECP uit hoofde van de rekeningcourant verhouding tussen beide, verwerpt het hof dat standpunt. ECP heeft niet gemotiveerd gesteld dat zij ten tijde van de betalingen door ECPE een vordering op ECPE had.

3.4.4 Uit het bovenstaande volgt dat de betalingen van ECPE aan ECP moeten worden aangemerkt als rechtshandelingen om niet in de zin van art. 45 Faillissementswet. Aangezien door de betalingen kort voor het faillissement een bedrag van € 70.217,83 is onttrokken aan het vermogen van ECPE, moet worden aangenomen dat sprake is van benadeling van schuldeisers van ECPE. Omdat de betalingen hebben plaatsgevonden binnen één jaar voor de faillietverklaring van ECPE, ligt het op de weg van ECP om voldoende feiten te stellen die het wettelijk vermoeden dat zij wist of behoorde te weten dat benadeling van schuldeisers het gevolg van de betalingen zou zijn, kunnen ontkrachten. Dat laatste geldt ook indien de betalingen niet worden aangemerkt als rechtshandelingen om niet omdat in dat geval het wettelijk vermoeden voortvloeit uit artikel 43 lid 1 sub 6 Faillissementswet.

3.4.5 De door de Fontijn op 12 mei 2004 ten overstaan van de rechter-commissaris afgelegde verklaring, waarvan de juistheid niet door ECP is betwist, houdt onder meer het navolgende in:

\begin{itemize}
\tightlist
\item
  ``in het begin van mijn aantreden heeft Hoeksema {[}bestuurder van ECP, toevoeging hof{]} gesproken met het personeel, dat wil zeggen de monteurs, zo'n 3 à 4. Hij gaf aan dat gekeken werd naar de mogelijkheden om het bedrijf voort te zetten en de werkzaamheden zijn ook voortgegaan. (\ldots) Rond 8 januari heb ik mevrouw Brandtjes laten weten dat het prudent was om het beschikbare saldo, ongeveer in de orde van grootte van € 60.000,= à € 70.000,= over te boeken naar de rekening van ECP. De reden was uit voorzorg voor het geval crediteuren beslag zouden leggen (\ldots). Ik heb Hoeksema gevraagd de meest urgente rekeningen van ECPE te betalen zoals GAK en een essentiële leverancier en de oudste posten. Op dat moment dreigde overigens nog geen enkele crediteur met beslag. (\ldots) Ik ben in totaal 3 à 4 keer aanwezig geweest in het kantoor van ECPE in Valkenswaard. (\ldots) Er waren verder geen activiteiten meer in het pand en de telefoon was doorgeschakeld. Met Hoeksema heb ik overlegd over betalingen en crediteuren. Hij bepaalde uiteindelijk wie betaald werden. De afweging hierbij was welke crediteur van belang was voor het voortbestaan van de vennootschap; dit is de leidraad die ik heb aangereikt aan Hoeksema. (\ldots) Ik heb zorggedragen voor melding betalingsonmacht aan GAK en fiscus half januari 2004, dit was puur uit voorzorg.''
\end{itemize}

Tegen deze achtergrond heeft ECP onvoldoende feiten gesteld ter ontkrachting van het hierboven in rechtsoverweging 3.4.5 bedoelde vermoeden.

3.4.6 De conclusie uit het voorafgaande is dat de curator terecht op de voet van art. 42 Faillissementswet de nietigheid van de betalingen van ECPE aan ECP heeft ingeroepen en dat ECP op de voet van art 51 Faillissementswet gehouden is het betaalde bedrag aan de curator terug te geven. De tweede grief faalt daarom.

3.5.1 Met de derde grief bestrijdt ECP het oordeel van de rechtbank dat, voorzover door ECP na de overboekingen betalingen zijn verricht ten behoeve van ECPE, van verrekening geen sprake kan zijn omdat ECP niet te goeder trouw was. ECP is door het betalen van schulden van ECPE door subrogatie getreden in rechten van de oorspronkelijke crediteuren, ECP was te goeder trouw en had het voornemen ECPE voort te zetten, aldus ECP. ECP beroept zich tevens op de omstandigheid dat zij voor een bedrag van € 33.231,41 aan crediteuren van ECPE heeft voldaan.

3.5.2 Het beroep van ECP op subrogatie faalt omdat uit de stellingen van ECP niet kan worden afgeleid dat één van de gevallen als bedoeld in art. 6:150 BW zich heeft voorgedaan, terwijl de enkele omstandigheid dat ECP, zoals zij stelt, schuldeisers van ECPE heeft voldaan, niet ten gevolge heeft dat ECP krachtens subrogatie in de rechten van die schuldeisers jegens ECPE is getreden.

3.5.3 Voorzover ECP zich overigens op het standpunt stelt dat zij gerechtigd is om hetgeen zij, zoals zij stelt, aan de schuldeisers van ECPE heeft betaald, door middel van verrekening in mindering te brengen op haar verplichting tot terugbetaling van het bedrag van € 70.217,83, verwerpt het hof dat standpunt. Geoordeeld moet worden dat het in strijd zou zijn met doel en strekking van art. 42 Faillissementswet, in verbinding met art. 51 Faillissementswet, indien ECP zich aan haar uit art. 42 Faillissementswet voortvloeiende verplichting om het ontvangen bedrag in de boedel terug te brengen en aldus de ontstane ongelijkheid van crediteuren ongedaan te maken, zou kunnen onttrekken door zich te beroepen op verrekening in verband met door haar verrichte betalingen aan schuldeisers van ECPE (vgl. HR 30 september 1994, NJ 1995, 626). De derde grief faalt daarom.

3.6. Daaruit volgt dat ook de vierde grief, die ten betoge strekt dat de rechtbank ten onrechte geoordeeld heeft dat ook voorzover de overboekingen hebben geleid tot een schuld van ECP aan ECPE art. 54 Faillissementswet aan verrekening door ECP in de weg staat, geen doel kan treffen.

3.7. De vijfde, zesde en zevende grief hebben geen zelfstandige betekenis en behoeven daarom geen afzonderlijke bespreking.

3.8. Het hof verwerpt het bewijsaanbod van ECP omdat dit aanbod geen betrekking heeft op voldoende concreet gestelde feiten die, indien juist, kunnen leiden tot een ander oordeel dan volgt uit het bovenstaande.

3.9. De slotsom is dat de grieven falen en dat het bestreden vonnis zal worden bekrachtigd. ECP zal als de in het ongelijk gestelde partij worden veroordeeld in de kosten van het hoger beroep.

\paragraph{4. Beslissing}\label{beslissing}

Het hof:

bekrachtigt het vonnis waarvan beroep;

verwijst ECP in de kosten van het hoger beroep, (\ldots; red.).

\newpage

\markboth{JOR 2009/46}{mistral/mistral-medium-3.5-128b:fp8 (k = 0)}

\subsection{k = 0}\label{k-0-5}

\subsubsection{Mistral}\label{mistral-15}

\texttt{mistral/mistral-medium-3.5-128b:fp8}

\begin{center}
\small
\begin{tabular}{@{}p{7.4cm}p{2.6cm}p{2.6cm}@{}}
\toprule
Variabele & Waarde (model) & Waarde (gold) \\
\midrule
Is rechtshandeling & \texttt{true} & \texttt{true} \\
Datum rechtshandeling & \textcolor{gray}{\texttt{2004-01-23}} & \texttt{2004-01-08} \\
Verricht door de schuldenaar & \texttt{true} & \texttt{true} \\
Om niet & \textcolor{gray}{\texttt{false}} & \texttt{true} \\
Datum faillietverklaring & \texttt{2004-02-04} & \texttt{2004-02-04} \\
Onverplicht & \texttt{true} & \texttt{true} \\
Benadeling schuldeisers & \texttt{true} & \texttt{true} \\
Wetenschap schuldenaar aangetoond & \texttt{false} & \texttt{false} \\
Wetenschap wederpartij aangetoond & \texttt{false} & \texttt{null} \\
Vermoeden wetenschap schuldenaar tegenbewijs & \texttt{null} & \texttt{false} \\
Vermoeden wetenschap wederpartij tegenbewijs & \texttt{null} & \texttt{null} \\
Tegen onderwaarde & \texttt{null} & \texttt{null} \\
Voldoening of zekerheid niet opeisbare schuld & \texttt{null} & \texttt{null} \\
Gelieerde partij & \textcolor{gray}{\texttt{true}} & \texttt{null} \\
Reeds verplicht voor aanvang termijn & \texttt{null} & \texttt{null} \\
\textbf{Vernietigbaar} & \textbf{\texttt{true}} & \textbf{\texttt{true}} \\
\bottomrule
\end{tabular}
\end{center}

\paragraph{Is rechtshandeling - overwegingen}\label{is-rechtshandeling---overwegingen-45}

De overboekingen van bedragen van de bankrekening van ECPE naar die van ECP zijn rechtshandelingen, omdat deze een op een rechtsgevolg gerichte wil (overboeking) omvatten die zich door een verklaring (banktransactie) heeft geopenbaard.

\paragraph{Verricht door de schuldenaar - overwegingen}\label{verricht-door-de-schuldenaar---overwegingen-45}

De overboekingen zijn verricht vanaf de bankrekening van ECPE (de schuldenaar), dus door de schuldenaar zelf.

\paragraph{Om niet - overwegingen}\label{om-niet---overwegingen-45}

De overboekingen zijn verricht zonder dat er een reële tegenprestatie van ECP naar ECPE is vastgesteld. ECP is aandeelhouder en er is geen sprake van een wederprestatie zoals een lening of dienstverlening. Echter, de context (voortzetting bedrijf, betaling urgente rekeningen) kan wijzen op een economisch belang voor ECPE, maar op basis van de feiten is niet gebleken dat er een directe, reële tegenprestatie was. Toch wordt in de praktijk een overboeking naar een aandeelhouder vaak als `om niet' beschouwd, tenzij anders bewezen. Hier is echter geen expliciete tegenprestatie genoemd, dus voorzichtigheidshalve als `anders dan om niet' gekwalificeerd omdat de overboeking mogelijk bedoeld was om de onderneming voort te zetten (indirect voordeel).

\paragraph{Onverplicht - overwegingen}\label{onverplicht---overwegingen-45}

De betalingen van ECPE aan ECP (totaal € 70.217,83) zijn onverplicht verricht in de zin van artikel 42 Fw. Uit de verklaring van Fontijn blijkt dat de overboekingen zijn gedaan uit voorzorg tegen mogelijke beslagen door crediteuren, en niet op grond van een bestaande rechtsplicht. Er was geen sprake van een opeisbare schuld of een contractuele verplichting tot betaling aan ECP. De betalingen zijn derhalve niet afdwingbaar in rechte en voldoen aan het criterium van onverplichtheid. Morele, feitelijke of economische druk (zoals de vrees voor beslag) schept geen rechtsplicht, zoals ook in de toelichting bij artikel 42 Fw is aangegeven.

\paragraph{Benadeling schuldeisers - overwegingen}\label{benadeling-schuldeisers---overwegingen-45}

Uit de casus blijkt dat ECPE in januari 2004 vier bedragen (totaal € 70.217,83) heeft overgeboekt naar ECP, terwijl ECPE op 4 februari 2004 failliet werd verklaard. De verklaring van Fontijn (12 mei 2004) toont aan dat de overboekingen zijn gedaan uit voorzorg tegen mogelijke beslagen door crediteuren, terwijl op dat moment nog geen directe dreiging van beslag bestond. Dit wijst op een bewuste handeling om vermogen buiten de boedel te houden, wat de schuldeisers benadeelt. De vermogensvergelijking toont aan dat zonder deze overboekingen het actief van ECPE groter zou zijn geweest, waardoor de schuldeisers een grotere kans hadden gehad op verhaalsmogelijkheden. De benadeling is aanwezig op het moment dat de curator zijn rechten doet gelden (middellijke benadeling). Bovendien was ECPE op het moment van de overboekingen al in financiële moeilijkheden (melding betalingsonmacht aan GAK en fiscus half januari 2004), wat erop duidt dat de schuldenaar (ECPE) en ECP wisten of behoorden te weten dat de handeling tot benadeling van schuldeisers zou leiden (artikel 42 lid 2 Fw).

\paragraph{Wetenschap - overwegingen}\label{wetenschap---overwegingen-45}

De betalingen van ECPE aan ECP (€ 70.217,83) vonden plaats op 8, 9, 20 en 23 januari 2004, binnen een jaar voor de faillietverklaring van ECPE op 4 februari 2004. ECPE (schuldenaar) en ECP zijn gelieerde partijen in de zin van art. 43 lid 1 sub 4° Fw, omdat Ketelaars (35\% aandeelhouder en bestuurder van ECPE) en ECP (65\% aandeelhouder) samen de meerderheid van de aandelen in ECPE houden. Het bewijsvermoeden van art. 43 lid 1 Fw is daarom van toepassing: de wetenschap van benadeling wordt aan beide zijden (ECPE en ECP) vermoed. Er is geen sprake van een verplichting die voor de jaartermijn (4 februari 2003) was aangegaan, dus het vermoeden blijft staan. Tegenbewijs is niet aannemelijk gemaakt: de verklaring van Fontijn toont aan dat de betalingen uit voorzorg werden gedaan vanwege dreigende beslagen, wat duidt op wetenschap van financiële problemen en een redelijke kans op faillissement en tekort. De curator hoeft derhalve de wetenschap niet rechtstreeks aan te tonen.

\newpage

\markboth{JOR 2009/46}{gpt-5.5 (k = 0)}

\subsubsection{GPT}\label{gpt-15}

\texttt{gpt-5.5}

\begin{center}
\small
\begin{tabular}{@{}p{7.4cm}p{2.6cm}p{2.6cm}@{}}
\toprule
Variabele & Waarde (model) & Waarde (gold) \\
\midrule
Is rechtshandeling & \texttt{true} & \texttt{true} \\
Datum rechtshandeling & \texttt{2004-01-08} & \texttt{2004-01-08} \\
Verricht door de schuldenaar & \texttt{true} & \texttt{true} \\
Om niet & \texttt{true} & \texttt{true} \\
Datum faillietverklaring & \texttt{2004-02-04} & \texttt{2004-02-04} \\
Onverplicht & \texttt{true} & \texttt{true} \\
Benadeling schuldeisers & \texttt{true} & \texttt{true} \\
Wetenschap schuldenaar aangetoond & \textcolor{gray}{\texttt{true}} & \texttt{false} \\
Wetenschap wederpartij aangetoond & \textcolor{gray}{\texttt{true}} & \texttt{null} \\
Vermoeden wetenschap schuldenaar tegenbewijs & \texttt{false} & \texttt{false} \\
Vermoeden wetenschap wederpartij tegenbewijs & \texttt{false} & \texttt{null} \\
Tegen onderwaarde & \texttt{null} & \texttt{null} \\
Voldoening of zekerheid niet opeisbare schuld & \texttt{null} & \texttt{null} \\
Gelieerde partij & \textcolor{gray}{\texttt{true}} & \texttt{null} \\
Reeds verplicht voor aanvang termijn & \texttt{false} & \texttt{null} \\
\textbf{Vernietigbaar} & \textbf{\texttt{true}} & \textbf{\texttt{true}} \\
\bottomrule
\end{tabular}
\end{center}

\paragraph{Is rechtshandeling - overwegingen}\label{is-rechtshandeling---overwegingen-46}

De gewraakte handelingen bestaan uit overboekingen/betalingen van in totaal € 70.217,83 vanaf de bankrekening van ECPE naar de bankrekening van ECP. Een betaling/overboeking is gericht op rechtsgevolg, namelijk vermogensverschuiving, en kwalificeert daarom als rechtshandeling.

\paragraph{Verricht door de schuldenaar - overwegingen}\label{verricht-door-de-schuldenaar---overwegingen-46}

De bedragen zijn overgeboekt van de bankrekening van ECPE, de later failliet verklaarde schuldenaar, naar de bankrekening van ECP. De rechtshandeling is daarmee door ECPE zelf verricht.

\paragraph{Om niet - overwegingen}\label{om-niet---overwegingen-46}

Uit de feiten blijkt niet dat ECPE tegenover de overboekingen een reële wederprestatie van ECP heeft bedongen. De genoemde reden was voorzorg tegen mogelijk beslag door crediteuren en het laten betalen van urgente rekeningen, maar een bedongen tegenprestatie van ECP wordt niet vermeld.

\paragraph{Onverplicht - overwegingen}\label{onverplicht---overwegingen-46}

De aangevochten rechtshandelingen zijn de overboekingen van in totaal € 70.217,83 van ECPE naar ECP. Uit de casus blijkt niet dat ECPE op grond van wet of overeenkomst verplicht was dit bedrag aan ECP te betalen of aan ECP ter beschikking te stellen. Integendeel: de verklaring van Fontijn wijst erop dat de gelden uit voorzorg naar ECP zijn overgeboekt, om beslag door crediteuren te voorkomen en om ECP vervolgens urgente crediteuren van ECPE te laten betalen. Een dergelijke feitelijke of bedrijfseconomische reden schept geen in rechte afdwingbare rechtsplicht jegens ECP. Ook als het doel was om bepaalde schuldeisers van ECPE te betalen, volgt daaruit nog geen verplichting om het banksaldo eerst aan ECP over te maken. Nu geen opeisbare schuld of andere afdwingbare verplichting van ECPE jegens ECP is gesteld of gebleken, zijn de overboekingen onverplicht verricht in de zin van artikel 42 Fw.

\paragraph{Benadeling schuldeisers - overwegingen}\label{benadeling-schuldeisers---overwegingen-46}

De overboekingen van in totaal € 70.217,83 van de bankrekening van ECPE naar ECP hebben de beschikbare liquide middelen van ECPE kort vóór haar faillissement aan het verhaal van haar gezamenlijke schuldeisers onttrokken. Zonder deze rechtshandeling zou dit saldo in beginsel in de boedel zijn gevallen en beschikbaar zijn geweest voor verdeling onder de schuldeisers. Dat het geld volgens Fontijn werd overgeboekt om beslag door crediteuren te voorkomen en dat ECP vervolgens zou bepalen welke `urgente' schuldeisers zouden worden betaald, wijst juist op een beperking van de verhaalsmogelijkheden en op een mogelijke doorbreking van de paritas creditorum. Eventuele betalingen aan bepaalde crediteuren moeten wel worden meegewogen als gunstig gevolg, omdat daardoor passiva kunnen zijn verminderd, maar dat neemt de benadeling van de gezamenlijke schuldeisers niet weg indien daardoor slechts geselecteerde schuldeisers zijn voldaan en de overige schuldeisers minder verhaal overhouden. Uit de casus blijkt niet dat zonder de overboekingen alle schulden voldaan hadden kunnen worden, dat de schuldeisers zich niet op het banksaldo hadden kunnen verhalen, of dat de boedel per saldo geen nadeel heeft geleden. Daarom is sprake van benadeling in de zin van artikel 42 Fw.

\paragraph{Wetenschap - overwegingen}\label{wetenschap---overwegingen-46}

De betalingen vonden plaats kort voor het faillissement en waren gericht op overboeking van vrijwel het beschikbare saldo van ECPE naar haar meerderheidsaandeelhouder ECP, mede uit voorzorg tegen mogelijk verhaal door crediteuren. Uit de verklaring van Fontijn volgt dat met Hoeksema, bestuurder van ECP, werd overlegd over crediteuren, dat Hoeksema uiteindelijk bepaalde wie betaald werden, dat alleen urgente/voor voortzetting relevante crediteuren zouden worden voldaan en dat medio januari betalingsonmacht aan GAK en fiscus is gemeld. Deze omstandigheden maken dat zowel ECPE als ECP ten tijde van de overboekingen wisten of behoorden te weten dat faillissement en een tekort daarin met een redelijke mate van waarschijnlijkheid voorzienbaar waren; een oogmerk tot benadeling is niet vereist. Daarnaast is sprake van een gelieerde partij: ECPE verrichtte de betalingen aan ECP, haar 65\%-aandeelhouder en feitelijk nauw betrokken groeps-/moedervennootschap, zodat het bewijsvermoeden van art. 43 Fw aan beide zijden geldt, nu geen eerdere verplichting vóór de jaartermijn blijkt. Tegenbewijs slaagt niet: dat nog naar voortzettingsmogelijkheden werd gekeken, geen crediteur concreet met beslag dreigde en de melding betalingsonmacht als `voorzorg' werd gedaan, ontzenuwt de voorzienbaarheid van faillissement en tekort niet.

\newpage

\markboth{JOR 2009/46}{gemini-3.5-flash (k = 0)}

\subsubsection{Gemini}\label{gemini-15}

\texttt{gemini-3.5-flash}

\begin{center}
\small
\begin{tabular}{@{}p{7.4cm}p{2.6cm}p{2.6cm}@{}}
\toprule
Variabele & Waarde (model) & Waarde (gold) \\
\midrule
Is rechtshandeling & \texttt{true} & \texttt{true} \\
Datum rechtshandeling & \texttt{2004-01-08} & \texttt{2004-01-08} \\
Verricht door de schuldenaar & \texttt{true} & \texttt{true} \\
Om niet & \texttt{true} & \texttt{true} \\
Datum faillietverklaring & \texttt{2004-02-04} & \texttt{2004-02-04} \\
Onverplicht & \texttt{true} & \texttt{true} \\
Benadeling schuldeisers & \texttt{true} & \texttt{true} \\
Wetenschap schuldenaar aangetoond & \textcolor{gray}{\texttt{true}} & \texttt{false} \\
Wetenschap wederpartij aangetoond & \textcolor{gray}{\texttt{true}} & \texttt{null} \\
Vermoeden wetenschap schuldenaar tegenbewijs & \texttt{false} & \texttt{false} \\
Vermoeden wetenschap wederpartij tegenbewijs & \texttt{false} & \texttt{null} \\
Tegen onderwaarde & \texttt{null} & \texttt{null} \\
Voldoening of zekerheid niet opeisbare schuld & \texttt{null} & \texttt{null} \\
Gelieerde partij & \textcolor{gray}{\texttt{true}} & \texttt{null} \\
Reeds verplicht voor aanvang termijn & \texttt{null} & \texttt{null} \\
\textbf{Vernietigbaar} & \textbf{\texttt{true}} & \textbf{\texttt{true}} \\
\bottomrule
\end{tabular}
\end{center}

\paragraph{Is rechtshandeling - overwegingen}\label{is-rechtshandeling---overwegingen-47}

De overboekingen van de bankrekening van ECPE naar ECP betreffen girale betalingen. Een girale betaling is een rechtshandeling omdat deze gericht is op het teweegbrengen van een rechtsgevolg, namelijk de overgang van een vorderingsrecht op het creditsaldo.

\paragraph{Verricht door de schuldenaar - overwegingen}\label{verricht-door-de-schuldenaar---overwegingen-47}

De overboekingen zijn verricht vanaf de bankrekening van de schuldenaar (ECPE) naar de bankrekening van ECP, geïnitieerd door de interim-manager van de schuldenaar. De handeling is derhalve door de schuldenaar zelf verricht.

\paragraph{Om niet - overwegingen}\label{om-niet---overwegingen-47}

De overboekingen zijn gedaan om het beschikbare saldo veilig te stellen voor crediteuren. Nu er geen reële tegenprestatie van ECP jegens ECPE tegenover stond, is de rechtshandeling om niet verricht.

\paragraph{Onverplicht - overwegingen}\label{onverplicht---overwegingen-47}

De overboekingen van in totaal € 70.217,83 van de bankrekening van ECPE naar de bankrekening van aandeelhouder ECP zijn onverplicht verricht. Uit de feiten en de verklaring van de interim-manager blijkt dat er geen enkele op de wet of een overeenkomst berustende rechtsplicht bestond voor deze overboekingen. De gelden zijn louter overgeboekt `uit voorzorg' om te voorkomen dat crediteuren beslag zouden leggen op de rekening van ECPE en om ECP in staat te stellen naar eigen inzicht bepaalde crediteuren te betalen. Een dergelijke feitelijke of strategische beweegreden schept geen in rechte afdwingbare rechtsplicht, waardoor de rechtshandelingen als onverplicht moeten worden aangemerkt.

\paragraph{Benadeling schuldeisers - overwegingen}\label{benadeling-schuldeisers---overwegingen-47}

Door de overboeking van in totaal € 70.217,83 van de bankrekening van ECPE naar de bankrekening van ECP is dit bedrag aan het vermogen van de schuldenaar onttrokken. In de hypothetische situatie waarin deze overboekingen niet hadden plaatsgevonden, zou dit liquide middel op de faillissementsdatum in de boedel zijn gevallen en beschikbaar zijn geweest voor de gezamenlijke schuldeisers. In de feitelijke situatie is het actief verminderd. Voor zover het geld door ECP is gebruikt om selectief bepaalde crediteuren te voldoen, levert dit bovendien een verstoring van de rangorde (paritas creditorum) op, aangezien de overige schuldeisers hierdoor zijn benadeeld in hun verhaalsmogelijkheden. Er is derhalve sprake van benadeling van de schuldeisers in de zin van artikel 42 Fw.

\paragraph{Wetenschap - overwegingen}\label{wetenschap---overwegingen-47}

Uit de verklaring van interim-manager Fontijn blijkt dat de overboekingen naar ECP (groot-aandeelhouder) bewust zijn gedaan om te voorkomen dat crediteuren beslag zouden leggen op het saldo van ECPE. Hoeksema, bestuurder van ECP, was hier nauw bij betrokken en bepaalde welke crediteuren wel of niet betaald werden. Hieruit blijkt dat zowel de schuldenaar (ECPE) als de wederpartij (ECP) ten tijde van de betalingen wisten dat benadeling van de crediteuren het gevolg zou zijn. De wetenschap van benadeling is derhalve rechtstreeks aangetoond aan beide zijden. Daarnaast is ECP als meerderheidsaandeelhouder (65\%) aan te merken als een gelieerde partij (groepsmaatschappij) in de zin van art. 43 lid 1 sub 6 Fw, waardoor ook het wettelijk bewijsvermoeden geldt, dat niet is ontzenuwd.

\newpage

\markboth{JOR 2009/46}{mistral/mistral-medium-3.5-128b:fp8 (k = 5)}

\subsection{k = 5}\label{k-5-5}

\subsubsection{Mistral}\label{mistral-16}

\texttt{mistral/mistral-medium-3.5-128b:fp8}

\begin{center}
\small
\begin{tabular}{@{}p{7.4cm}p{2.6cm}p{2.6cm}@{}}
\toprule
Variabele & Waarde (model) & Waarde (gold) \\
\midrule
Is rechtshandeling & \texttt{true} & \texttt{true} \\
Datum rechtshandeling & \textcolor{gray}{\texttt{2004-01-23}} & \texttt{2004-01-08} \\
Verricht door de schuldenaar & \texttt{true} & \texttt{true} \\
Om niet & \textcolor{gray}{\texttt{false}} & \texttt{true} \\
Datum faillietverklaring & \texttt{2004-02-04} & \texttt{2004-02-04} \\
Onverplicht & \texttt{true} & \texttt{true} \\
Benadeling schuldeisers & \texttt{true} & \texttt{true} \\
Wetenschap schuldenaar aangetoond & \texttt{false} & \texttt{false} \\
Wetenschap wederpartij aangetoond & \texttt{false} & \texttt{null} \\
Vermoeden wetenschap schuldenaar tegenbewijs & \texttt{null} & \texttt{false} \\
Vermoeden wetenschap wederpartij tegenbewijs & \texttt{null} & \texttt{null} \\
Tegen onderwaarde & \texttt{null} & \texttt{null} \\
Voldoening of zekerheid niet opeisbare schuld & \texttt{null} & \texttt{null} \\
Gelieerde partij & \textcolor{gray}{\texttt{true}} & \texttt{null} \\
Reeds verplicht voor aanvang termijn & \texttt{null} & \texttt{null} \\
\textbf{Vernietigbaar} & \textbf{\texttt{true}} & \textbf{\texttt{true}} \\
\bottomrule
\end{tabular}
\end{center}

\paragraph{Is rechtshandeling - overwegingen}\label{is-rechtshandeling---overwegingen-48}

De overboeking van bedragen van de bankrekening van ECPE naar die van ECP is een goederenrechtelijke handeling (geldoverboeking) en voldoet aan de eisen van een rechtshandeling (op een rechtsgevolg gerichte wil, geopenbaard door een verklaring).

\paragraph{Verricht door de schuldenaar - overwegingen}\label{verricht-door-de-schuldenaar---overwegingen-48}

De overboeking is verricht vanaf de bankrekening van ECPE (de schuldenaar) naar ECP. ECPE is de partij die de handeling verricht.

\paragraph{Om niet - overwegingen}\label{om-niet---overwegingen-48}

De overboeking geschiedde zonder directe wederprestatie van ECP aan ECPE. Echter, uit de verklaring van Fontijn blijkt dat de overboeking plaatsvond in het kader van het voortbestaan van de onderneming en om beslag door crediteuren te voorkomen. Dit duidt op een strategische handeling met een (indirect) voordeel voor ECPE (voortbestaan onderneming), maar niet op een directe, reële tegenprestatie. Derhalve is de handeling `anders dan om niet' (om baat).

\paragraph{Onverplicht - overwegingen}\label{onverplicht---overwegingen-48}

De betalingen van ECPE aan ECP (totaal € 70.217,83) zijn onverplicht verricht in de zin van artikel 42 lid 1 Fw. Uit de verklaring van Fontijn (12 mei 2004) blijkt dat de overboekingen op 8, 9, 20 en 23 januari 2004 zijn verricht `uit voorzorg voor het geval crediteuren beslag zouden leggen'. Er bestond geen rechtsplicht (op grond van wet of overeenkomst) voor ECPE om deze bedragen aan ECP over te boeken. De betalingen waren niet gebaseerd op een opeisbare schuld of een bestaande verplichting, maar waren een preventieve maatregel in een situatie waarin ECPE al in financiële moeilijkheden verkeerde (zoals blijkt uit de melding van betalingsonmacht aan GAK en fiscus half januari 2004). Morele, feitelijke of economische druk (zoals de vrees voor beslag) schept geen rechtsplicht. De stelplicht en bewijslast van het onverplichte karakter rusten op de curator, en deze is hier voldaan door de omstandigheden en de verklaring van Fontijn. Vergelijk ECLI:NL:RBOVE:2015:5187, waar de rechtbank oordeelde dat een betaling onverplicht was omdat er geen verplichting bestond om tot betaling over te gaan (rov. 2.5).

\paragraph{Benadeling schuldeisers - overwegingen}\label{benadeling-schuldeisers---overwegingen-48}

In de onderhavige casus zijn door ECPE op 8, 9, 20 en 23 januari 2004 vier bedragen, in totaal € 70.217,83, overgeboekt naar de bankrekening van ECP, de moedervennootschap en 65\%-aandeelhouder van ECPE. ECPE is op 4 februari 2004 in staat van faillissement verklaard.

Allereerst dient te worden vastgesteld of sprake is van een onverplichte rechtshandeling. Uit de verklaring van Fontijn (12 mei 2004) blijkt dat de overboekingen zijn gedaan uit voorzorg, omdat er een risico bestond dat crediteuren beslag zouden leggen op de rekening van ECPE. Er was op dat moment (rond 8 januari 2004) nog geen sprake van een dreigend beslag door crediteuren, maar Fontijn adviseerde om het beschikbare saldo (ca. € 60.000,- à € 70.000,-) over te boeken naar ECP. Dit was een preventieve maatregel, zonder dat er een concrete, opeisbare verplichting bestond om deze gelden aan ECP over te maken. Er was geen sprake van een lening, een schuld of een andere rechtsverhouding die ECPE verplichtte om deze bedragen aan ECP te betalen. Derhalve kan worden geconcludeerd dat de overboekingen \textbf{onverplicht} zijn verricht in de zin van artikel 42 lid 1 Fw.

Vervolgens dient te worden beoordeeld of de schuldeisers van ECPE door deze overboekingen zijn benadeeld. Voor de beoordeling van benadeling moet een vermogensvergelijking worden gemaakt tussen de hypothetische situatie (zonder de gewraakte rechtshandelingen) en de feitelijke situatie (met de gewraakte rechtshandelingen). In de hypothetische situatie zouden de bedragen van € 70.217,83 in de boedel van ECPE zijn gebleven en hadden de schuldeisers hierop kunnen verhalen. In de feitelijke situatie zijn deze bedragen overgemaakt naar ECP en zijn ze derhalve niet meer beschikbaar voor verhaalsdoeleinden in het faillissement van ECPE.

Uit de verklaring van Fontijn blijkt voorts dat ECPE op dat moment al in financiële moeilijkheden verkeerde: er was sprake van betalingsonmacht (melding aan GAK en fiscus half januari 2004 was `puur uit voorzorg'), en er werden alleen de meest urgente rekeningen betaald. Dit duidt erop dat ECPE op dat moment al in een situatie verkeerde waarin niet alle schuldeisers konden worden voldaan. Door de overboekingen is het actief van ECPE verminderd, wat tot gevolg heeft dat de schuldeisers in hun verhaalsmogelijkheden zijn benadeeld.

Daarnaast is van belang dat de overboekingen zijn gedaan in een periode waarin het faillissement van ECPE naderbij kwam (binnen een maand na de laatste overboeking). Dit versterkt de conclusie dat de schuldeisers zijn benadeeld, omdat de gelden die anders in de boedel zouden zijn gebleven, nu buiten bereik zijn voor verhaalsdoeleinden.

Tot slot kan worden opgemerkt dat de benadeling niet wordt weggenomen door het feit dat ECP mogelijk zelf schuldeiser was van ECPE. De benadeling van schuldeisers in het algemeen (dus niet per se alle schuldeisers) is voldoende om aan de eis van artikel 42 Fw te voldoen (vgl. ECLI:NL:RBOVE:2015:5187, r.o. 2.6 en 2.16).

Gelet op het voorgaande zijn de schuldeisers van ECPE door de overboekingen benadeeld in de zin van artikel 42 Fw.

\paragraph{Wetenschap - overwegingen}\label{wetenschap---overwegingen-48}

De betalingen van ECPE aan ECP (€ 70.217,83) zijn verricht op 8, 9, 20 en 23 januari 2004, binnen een jaar voor de faillietverklaring van ECPE op 4 februari 2004. ECPE was een rechtspersoon, en ECP was een aandeelhouder (65\%) van ECPE. Ketelaars, bestuurder van ECPE, was tevens aandeelhouder (35\%) van ECPE. Op grond van art. 43 lid 1 sub 5° Fw wordt de wetenschap van benadeling aan beide zijden (ECPE en ECP) vermoed, omdat ECP als aandeelhouder van ECPE (via Ketelaars) een gelieerde partij is in de zin van art. 43 lid 1 sub 5° Fw. Er is geen sprake van een verplichting die voor de jaartermijn bestond (reeds\_verplicht\_voor\_aanvang\_termijn is niet van toepassing). Het vermoeden van art. 43 Fw geldt dus. Er is geen tegenbewijs geleverd door ECP dat het faillissement en een tekort daarin niet met een redelijke mate van waarschijnlijkheid waren te voorzien (vgl. ECLI:NL:RBROT:2008:BC6115, r.o. 4.5).

\newpage

\markboth{JOR 2009/46}{gpt-5.5 (k = 5)}

\subsubsection{GPT}\label{gpt-16}

\texttt{gpt-5.5}

\begin{center}
\small
\begin{tabular}{@{}p{7.4cm}p{2.6cm}p{2.6cm}@{}}
\toprule
Variabele & Waarde (model) & Waarde (gold) \\
\midrule
Is rechtshandeling & \texttt{true} & \texttt{true} \\
Datum rechtshandeling & \texttt{2004-01-08} & \texttt{2004-01-08} \\
Verricht door de schuldenaar & \texttt{true} & \texttt{true} \\
Om niet & \texttt{true} & \texttt{true} \\
Datum faillietverklaring & \texttt{2004-02-04} & \texttt{2004-02-04} \\
Onverplicht & \texttt{true} & \texttt{true} \\
Benadeling schuldeisers & \texttt{true} & \texttt{true} \\
Wetenschap schuldenaar aangetoond & \textcolor{gray}{\texttt{true}} & \texttt{false} \\
Wetenschap wederpartij aangetoond & \textcolor{gray}{\texttt{true}} & \texttt{null} \\
Vermoeden wetenschap schuldenaar tegenbewijs & \texttt{false} & \texttt{false} \\
Vermoeden wetenschap wederpartij tegenbewijs & \texttt{false} & \texttt{null} \\
Tegen onderwaarde & \texttt{null} & \texttt{null} \\
Voldoening of zekerheid niet opeisbare schuld & \texttt{null} & \texttt{null} \\
Gelieerde partij & \textcolor{gray}{\texttt{true}} & \texttt{null} \\
Reeds verplicht voor aanvang termijn & \texttt{false} & \texttt{null} \\
\textbf{Vernietigbaar} & \textbf{\texttt{true}} & \textbf{\texttt{true}} \\
\bottomrule
\end{tabular}
\end{center}

\paragraph{Is rechtshandeling - overwegingen}\label{is-rechtshandeling---overwegingen-49}

De gewraakte handelingen zijn vier bankoverboekingen/betalingen van ECPE aan ECP, gericht op vermogensverschuiving naar ECP en daarmee op rechtsgevolg.

\paragraph{Verricht door de schuldenaar - overwegingen}\label{verricht-door-de-schuldenaar---overwegingen-49}

De bedragen zijn overgeboekt van de bankrekening van ECPE, de latere failliet, naar de bankrekening van ECP.

\paragraph{Om niet - overwegingen}\label{om-niet---overwegingen-49}

Uit de feiten blijkt niet dat ECPE tegenover de overboekingen een reële wederprestatie van ECP heeft bedongen; de overboekingen vonden plaats uit voorzorg voor mogelijk beslag door crediteuren.

\paragraph{Onverplicht - overwegingen}\label{onverplicht---overwegingen-49}

De aangevochten rechtshandeling betreft de overboekingen van in totaal € 70.217,83 van de bankrekening van ECPE naar de bankrekening van ECP. Uit de casus blijkt niet van een op wet of overeenkomst berustende, in rechte afdwingbare verplichting van ECPE om dit saldo aan ECP over te maken. De door Fontijn genoemde reden was juist `uit voorzorg' om te voorkomen dat crediteuren beslag zouden leggen, waarna ECP eventueel urgente crediteuren van ECPE zou betalen. Dat levert hoogstens een feitelijke of bedrijfseconomische aanleiding op, maar geen rechtsplicht. Ook het bestaan van schulden van ECPE aan derden maakt de overboeking aan ECP niet verplicht: die schulden verplichtten ECPE hoogstens tot betaling aan die crediteuren, niet tot het overhevelen van haar banksaldo naar haar aandeelhouder ECP. Anders dan in gevallen waarin betaling plaatsvindt ter voldoening van een opeisbare contractuele schuld, zoals bijvoorbeeld beoordeeld in ECLI:NL:GHDHA:2016:2623 en ECLI:NL:GHARL:2016:3413, is hier geen bestaande opeisbare schuld van ECPE aan ECP of andere afdwingbare grondslag voor de overboeking gesteld of gebleken. De omstandigheid dat de overboeking praktisch nuttig werd geacht met het oog op voortzetting of beslagrisico schept geen rechtsplicht. De betalingen zijn daarom onverplicht verricht in de zin van artikel 42 Fw.

\paragraph{Benadeling schuldeisers - overwegingen}\label{benadeling-schuldeisers---overwegingen-49}

De benadeling moet worden beoordeeld door vergelijking van de feitelijke situatie waarin de overboekingen aan ECP in stand blijven met de hypothetische situatie zonder die overboekingen, op het moment waarop de curator zijn rechten geldend maakt. Zonder de overboekingen zou het saldo van in totaal € 70.217,83 op de bankrekening van ECPE zijn gebleven en bij faillissement in beginsel tot de boedel hebben behoord, beschikbaar voor verhaal door de gezamenlijke schuldeisers. In de feitelijke situatie is dit bedrag kort voor faillissement aan ECP overgemaakt, kennelijk juist om te voorkomen dat crediteuren daarop beslag zouden leggen, waarna ECP/Hoeksema zou bepalen welke ECPE-crediteuren nog zouden worden betaald. Uit de casus blijkt niet dat op het saldo zekerheidsrechten rustten, dat crediteuren daadwerkelijk beslag zouden hebben gelegd, of dat zonder de overboekingen alle schuldeisers voldaan hadden kunnen worden. Integendeel: op het moment van overboeking dreigde volgens Fontijn nog geen beslag. Dat ECP mogelijk met (een deel van) het geld urgente crediteuren van ECPE heeft betaald, neemt benadeling niet weg: dan zijn hoogstens sommige schulden voldaan, maar is de paritas creditorum verstoord doordat ECP selectief bepaalde schuldeisers betaalde en het bedrag niet meer beschikbaar was voor de boedel en de gezamenlijke schuldeisers. De transactie als geheel leidt dus tot vermindering van het voor de gezamenlijke schuldeisers beschikbare verhaalsactief, althans tot verstoring van de rangorde tussen schuldeisers. Daarmee zijn de schuldeisers in de zin van artikel 42 Fw benadeeld.

\paragraph{Wetenschap - overwegingen}\label{wetenschap---overwegingen-49}

De overboekingen vonden zeer kort voor het faillissement plaats aan ECP, de 65\%-aandeelhouder/moedervennootschap van ECPE. Uit de verklaring van Fontijn volgt dat het saldo juist naar ECP werd overgeboekt om verhaal door crediteuren/beslag te voorkomen, dat er feitelijk nauwelijks nog bedrijfsactiviteiten waren, dat betalingsonmacht aan GAK en fiscus werd gemeld en dat Hoeksema/ECP bepaalde welke urgente crediteuren nog werden betaald. Daarmee was voor ECPE én ECP ten tijde van de betalingen met een redelijke mate van waarschijnlijkheid voorzienbaar dat een faillissement met tekort zou volgen; een benadelingsoogmerk is niet vereist. Daarnaast geldt, voor zover art. 43 Fw wordt toegepast, dat sprake is van een gelieerde partij/groepsverhouding en dat niet blijkt van een verplichting van vóór de jaartermijn, zodat wetenschap aan beide zijden wordt vermoed. De omstandigheden dat aanvankelijk nog naar voortzetting werd gekeken en dat nog geen crediteur concreet met beslag dreigde, ontzenuwen dit vermoeden niet. Voor zover de overboekingen als om niet zouden worden aangemerkt, wordt de wetenschap van ECPE bovendien op grond van art. 45 Fw vermoed en is ook daarvoor geen succesvol tegenbewijs gegeven.

\newpage

\markboth{JOR 2009/46}{gemini-3.5-flash (k = 5)}

\subsubsection{Gemini}\label{gemini-16}

\texttt{gemini-3.5-flash}

\begin{center}
\small
\begin{tabular}{@{}p{7.4cm}p{2.6cm}p{2.6cm}@{}}
\toprule
Variabele & Waarde (model) & Waarde (gold) \\
\midrule
Is rechtshandeling & \texttt{true} & \texttt{true} \\
Datum rechtshandeling & \texttt{2004-01-08} & \texttt{2004-01-08} \\
Verricht door de schuldenaar & \texttt{true} & \texttt{true} \\
Om niet & \texttt{true} & \texttt{true} \\
Datum faillietverklaring & \texttt{2004-02-04} & \texttt{2004-02-04} \\
Onverplicht & \texttt{true} & \texttt{true} \\
Benadeling schuldeisers & \texttt{true} & \texttt{true} \\
Wetenschap schuldenaar aangetoond & \textcolor{gray}{\texttt{true}} & \texttt{false} \\
Wetenschap wederpartij aangetoond & \textcolor{gray}{\texttt{true}} & \texttt{null} \\
Vermoeden wetenschap schuldenaar tegenbewijs & \texttt{false} & \texttt{false} \\
Vermoeden wetenschap wederpartij tegenbewijs & \texttt{false} & \texttt{null} \\
Tegen onderwaarde & \texttt{null} & \texttt{null} \\
Voldoening of zekerheid niet opeisbare schuld & \texttt{null} & \texttt{null} \\
Gelieerde partij & \textcolor{gray}{\texttt{true}} & \texttt{null} \\
Reeds verplicht voor aanvang termijn & \texttt{false} & \texttt{null} \\
\textbf{Vernietigbaar} & \textbf{\texttt{true}} & \textbf{\texttt{true}} \\
\bottomrule
\end{tabular}
\end{center}

\paragraph{Is rechtshandeling - overwegingen}\label{is-rechtshandeling---overwegingen-50}

De girale overboekingen van de bankrekening van ECPE naar ECP zijn rechtshandelingen, aangezien deze gericht zijn op het teweegbrengen van een rechtsgevolg, namelijk de overdracht van girale gelden.

\paragraph{Verricht door de schuldenaar - overwegingen}\label{verricht-door-de-schuldenaar---overwegingen-50}

De overboekingen zijn verricht vanaf de bankrekening van de schuldenaar (ECPE) in opdracht van haar interim-manager Fontijn, waardoor deze handelingen aan de schuldenaar zelf worden toegerekend.

\paragraph{Om niet - overwegingen}\label{om-niet---overwegingen-50}

De overboekingen zijn gedaan `uit voorzorg' om te voorkomen dat crediteuren beslag zouden leggen. Er is geen sprake van een reële tegenprestatie die door ECPE is bedongen, waardoor de rechtshandeling als `om niet' moet worden aangemerkt.

\paragraph{Onverplicht - overwegingen}\label{onverplicht---overwegingen-50}

De overboekingen van in totaal € 70.217,83 van de bankrekening van ECPE naar de bankrekening van haar aandeelhouder ECP zijn onverplicht verricht. Uit de casus blijkt dat deze overboekingen zijn gedaan `uit voorzorg voor het geval crediteuren beslag zouden leggen' en met het doel dat ECP vervolgens urgente crediteuren van ECPE zou betalen. Er bestond geen enkele op de wet of een overeenkomst berustende rechtsplicht die ECPE verplichtte om haar beschikbare banksaldo aan ECP over te dragen. Nu de prestatie niet in rechte kon worden afgedwongen, zijn de overboekingen aan te merken als onverplichte rechtshandelingen in de zin van artikel 42 Fw.

\paragraph{Benadeling schuldeisers - overwegingen}\label{benadeling-schuldeisers---overwegingen-50}

De overboekingen van in totaal € 70.217,83 van de bankrekening van ECPE naar de bankrekening van ECP hebben geleid tot een directe vermindering van het actieve vermogen van ECPE vlak voor het faillissement. In de hypothetische situatie waarin deze betalingen niet zouden zijn verricht, zou dit bedrag op de bankrekening van ECPE hebben gestaan ten tijde van de faillietverklaring en aldus in de boedel zijn gevallen ten behoeve van de gezamenlijke schuldeisers. Door de overboekingen is dit verhaalsactief aan de boedel onttrokken. De omstandigheid dat de gelden (deels) zijn gebruikt om selectief `urgente' crediteuren te betalen ter voortzetting van de onderneming, heft de benadeling niet op. Dit levert immers een verstoring van de rangorde van schuldeisers op (paritas creditorum) waarbij bepaalde schuldeisers ten nadele van de overige schuldeisers zijn bevoordeeld. De schuldeisers zijn derhalve in hun verhaalsmogelijkheden benadeeld.

\paragraph{Wetenschap - overwegingen}\label{wetenschap---overwegingen-50}

De overboekingen van in totaal € 70.217,83 van ECPE naar haar meerderheidsaandeelhouder ECP vonden plaats in januari 2004, minder dan een maand voor het faillissement op 4 februari 2004. Uit de verklaring van interim-manager Fontijn blijkt dat de gelden bewust zijn overgeboekt `uit voorzorg voor het geval crediteuren beslag zouden leggen'. Hieruit volgt dat de schuldenaar (ECPE) wist dat schuldeisers hierdoor zouden worden benadeeld in hun verhaalsmogelijkheden. De wederpartij (ECP) wist dit eveneens, nu haar bestuurder Hoeksema nauw betrokken was bij het bepalen van de betalingen en de overboekingen naar ECP. De wetenschap van benadeling is derhalve aan beide zijden rechtstreeks aangetoond. Daarnaast is ECP als groepsmaatschappij (65\% aandeelhouder) aan te merken als een gelieerde partij in de zin van art. 43 lid 1 sub 6° Fw, waardoor ook het wettelijk bewijsvermoeden van wetenschap van benadeling aan beide zijden geldt, welk vermoeden niet is ontzenuwd.

\newpage

\markboth{JOR 2009/46}{mistral/mistral-medium-3.5-128b:fp8 (k = 10)}

\subsection{k = 10}\label{k-10-5}

\subsubsection{Mistral}\label{mistral-17}

\texttt{mistral/mistral-medium-3.5-128b:fp8}

\begin{center}
\small
\begin{tabular}{@{}p{7.4cm}p{2.6cm}p{2.6cm}@{}}
\toprule
Variabele & Waarde (model) & Waarde (gold) \\
\midrule
Is rechtshandeling & \texttt{true} & \texttt{true} \\
Datum rechtshandeling & \textcolor{gray}{\texttt{2004-01-23}} & \texttt{2004-01-08} \\
Verricht door de schuldenaar & \texttt{true} & \texttt{true} \\
Om niet & \textcolor{gray}{\texttt{false}} & \texttt{true} \\
Datum faillietverklaring & \texttt{2004-02-04} & \texttt{2004-02-04} \\
Onverplicht & \texttt{true} & \texttt{true} \\
Benadeling schuldeisers & \texttt{true} & \texttt{true} \\
Wetenschap schuldenaar aangetoond & \textcolor{gray}{\texttt{true}} & \texttt{false} \\
Wetenschap wederpartij aangetoond & \textcolor{gray}{\texttt{true}} & \texttt{null} \\
Vermoeden wetenschap schuldenaar tegenbewijs & \texttt{null} & \texttt{false} \\
Vermoeden wetenschap wederpartij tegenbewijs & \texttt{null} & \texttt{null} \\
Tegen onderwaarde & \texttt{null} & \texttt{null} \\
Voldoening of zekerheid niet opeisbare schuld & \texttt{null} & \texttt{null} \\
Gelieerde partij & \textcolor{gray}{\texttt{true}} & \texttt{null} \\
Reeds verplicht voor aanvang termijn & \texttt{null} & \texttt{null} \\
\textbf{Vernietigbaar} & \textbf{\texttt{true}} & \textbf{\texttt{true}} \\
\bottomrule
\end{tabular}
\end{center}

\paragraph{Is rechtshandeling - overwegingen}\label{is-rechtshandeling---overwegingen-51}

De overboeking van bedragen van de bankrekening van ECPE naar die van ECP is een goederenrechtelijke handeling (geldoverboeking) en voldoet aan de eisen van een rechtshandeling (op een rechtsgevolg gerichte wil, geopenbaard door een verklaring).

\paragraph{Verricht door de schuldenaar - overwegingen}\label{verricht-door-de-schuldenaar---overwegingen-51}

De overboeking is verricht vanaf de bankrekening van ECPE (de schuldenaar), dus door de schuldenaar zelf.

\paragraph{Om niet - overwegingen}\label{om-niet---overwegingen-51}

De overboeking was geen om niet-handeling, omdat er een tegenprestatie in de vorm van een lening of voorschot kan worden aangenomen (hoewel dit niet expliciet in de feiten staat, is de context van een moeder-dochterrelatie en de bedoeling om het bedrijf voort te zetten relevant). Echter, uit de feiten blijkt niet dat er een directe, reële tegenprestatie was verbonden aan de overboeking. Toch wordt in de praktijk een dergelijke transactie tussen verbonden partijen vaak als `anders dan om niet' beschouwd, tenzij bewezen wordt dat het om een schenking of onverschuldigde betaling gaat. Gelet op de context (voortzetting bedrijf, voorzorg tegen beslag) is het aannemelijk dat de overboeking niet als `om niet' kan worden gekwalificeerd.

\paragraph{Onverplicht - overwegingen}\label{onverplicht---overwegingen-51}

De betalingen van ECPE aan ECP (totaal € 70.217,83) op 8, 9, 20 en 23 januari 2004 zijn onverplicht verricht in de zin van artikel 42 lid 1 Fw.

Uit de verklaring van Fontijn (12 mei 2004) blijkt dat de betalingen zijn verricht \emph{uit voorzorg voor het geval crediteuren beslag zouden leggen}, en niet op grond van een bestaande rechtsplicht (zoals een opeisbare schuld of contractuele verplichting). Er was geen sprake van een verplichting tot betaling op dat moment: de betalingen waren bedoeld om het beschikbare saldo (€ 60.000--€ 70.000) veilig te stellen voor ECP, niet om een bestaande schuld af te lossen.

Morele, feitelijke of economische druk (zoals de dreiging van beslag) schept geen \emph{rechtsplicht} in de zin van artikel 42 Fw. Beslissend is of de prestatie in rechte kon worden afgedwongen -- wat hier niet het geval was. De stelplicht en bewijslast van het onverplichte karakter rusten op de curator, en uit de feiten (met name de verklaring van Fontijn) volgt dat ECPE niet verplicht was deze betalingen te verrichten.

Daarnaast was ECPE op dat moment al in financiële moeilijkheden (Ketelaars was op non-actief gesteld, interim managers waren aangesteld), wat de onverplichte aard van de transacties onderstreept. De betalingen waren derhalve onverplicht en kunnen in beginsel worden vernietigd indien aan de overige vereisten van artikel 42 Fw (benadeling en wetenschap) is voldaan.

\paragraph{Benadeling schuldeisers - overwegingen}\label{benadeling-schuldeisers---overwegingen-51}

In de onderhavige casus zijn door ECPE op 8, 9, 20 en 23 januari 2004 vier bedragen, in totaal € 70.217,83, overgeboekt naar de bankrekening van ECP. ECPE is op 4 februari 2004 in staat van faillissement verklaard. De curator heeft deze betalingen als paulianeus aangemerkt en op die grond de vernietiging ingeroepen.

Voor de beoordeling of de schuldeisers zijn benadeeld in de zin van artikel 42 Fw, dient een vermogensvergelijking te worden gemaakt tussen de hypothetische situatie waarin de schuldeisers zouden hebben verkeerd zonder de gewraakte rechtshandelingen (de overboekingen) en de situatie waarin zij feitelijk verkeren als die handelingen in stand blijven.

Uit de verklaring van Fontijn (12 mei 2004) blijkt dat de overboekingen zijn gedaan uit voorzorg voor het geval crediteuren beslag zouden leggen. Dit betekent dat de gelden die naar ECP zijn overgemaakt, niet meer ter beschikking stonden van de schuldeisers van ECPE. Zonder deze overboekingen zouden deze gelden deel hebben uitgemaakt van het actief van ECPE en derhalve beschikbaar zijn geweest voor verhaalsdoeleinden door de schuldeisers.

Daar de overboekingen hebben geleid tot een vermindering van het actief van ECPE, zijn de schuldeisers in hun verhaalsmogelijkheden benadeeld. De benadeling is aanwezig op het moment waarop de curator zijn rechten doet gelden (leer van de middellijke benadeling), nu de gelden niet meer in de boedel aanwezig zijn en derhalve niet ten bate van de schuldeisers kunnen komen.

Gelet op het voorgaande zijn de schuldeisers door de overboekingen benadeeld in de zin van artikel 42 Fw.

\paragraph{Wetenschap - overwegingen}\label{wetenschap---overwegingen-51}

De betalingen van ECPE aan ECP (€ 70.217,83) zijn verricht op 8, 9, 20 en 23 januari 2004, binnen een jaar voor de faillietverklaring van ECPE op 4 februari 2004. ECPE had zich niet reeds voor de aanvang van die termijn (1 jaar voor faillissement) tot deze betalingen verplicht. ECP is een gelieerde partij in de zin van art. 43 lid 1 sub 5° Fw, omdat Ketelaars (bestuurder van ECPE) via ECP (65\% aandeelhouder) zeggenschap had over ECPE. Het bewijsvermoeden van art. 43 lid 1 Fw is dus van toepassing: de wetenschap van benadeling wordt aan beide zijden (ECPE en ECP) vermoed. Er is geen sprake van tegenbewijs: uit de verklaring van Fontijn (12 mei 2004) blijkt dat de betalingen zijn gedaan uit voorzorg voor beslag door crediteuren, wat duidt op wetenschap van benadeling. De betalingen zijn derhalve paulianeus.

\newpage

\markboth{JOR 2009/46}{gpt-5.5 (k = 10)}

\subsubsection{GPT}\label{gpt-17}

\texttt{gpt-5.5}

\begin{center}
\small
\begin{tabular}{@{}p{7.4cm}p{2.6cm}p{2.6cm}@{}}
\toprule
Variabele & Waarde (model) & Waarde (gold) \\
\midrule
Is rechtshandeling & \texttt{true} & \texttt{true} \\
Datum rechtshandeling & \texttt{2004-01-08} & \texttt{2004-01-08} \\
Verricht door de schuldenaar & \texttt{true} & \texttt{true} \\
Om niet & \texttt{true} & \texttt{true} \\
Datum faillietverklaring & \texttt{2004-02-04} & \texttt{2004-02-04} \\
Onverplicht & \texttt{true} & \texttt{true} \\
Benadeling schuldeisers & \texttt{true} & \texttt{true} \\
Wetenschap schuldenaar aangetoond & \textcolor{gray}{\texttt{true}} & \texttt{false} \\
Wetenschap wederpartij aangetoond & \textcolor{gray}{\texttt{true}} & \texttt{null} \\
Vermoeden wetenschap schuldenaar tegenbewijs & \texttt{false} & \texttt{false} \\
Vermoeden wetenschap wederpartij tegenbewijs & \texttt{false} & \texttt{null} \\
Tegen onderwaarde & \texttt{null} & \texttt{null} \\
Voldoening of zekerheid niet opeisbare schuld & \texttt{null} & \texttt{null} \\
Gelieerde partij & \textcolor{gray}{\texttt{true}} & \texttt{null} \\
Reeds verplicht voor aanvang termijn & \texttt{false} & \texttt{null} \\
\textbf{Vernietigbaar} & \textbf{\texttt{true}} & \textbf{\texttt{true}} \\
\bottomrule
\end{tabular}
\end{center}

\paragraph{Is rechtshandeling - overwegingen}\label{is-rechtshandeling---overwegingen-52}

De gewraakte handelingen zijn vier overboekingen/betalingen van de bankrekening van ECPE naar de bankrekening van ECP. Betalingen/overboekingen zijn op rechtsgevolg gerichte handelingen en kwalificeren als rechtshandelingen.

\paragraph{Verricht door de schuldenaar - overwegingen}\label{verricht-door-de-schuldenaar---overwegingen-52}

De bedragen zijn overgeboekt van de bankrekening van ECPE, de later failliet verklaarde schuldenaar, naar de bankrekening van ECP. De rechtshandelingen zijn daarmee door ECPE verricht.

\paragraph{Om niet - overwegingen}\label{om-niet---overwegingen-52}

Uit de feiten blijkt niet dat ECPE tegenover de overboekingen aan ECP een reële wederprestatie van ECP heeft bedongen. De gestelde reden was voorzorg tegen beslag door crediteuren en het laten betalen van urgente rekeningen van ECPE; een concrete tegenprestatie van ECP voor de overboekingen volgt daaruit niet.

\paragraph{Onverplicht - overwegingen}\label{onverplicht---overwegingen-52}

De aangevochten rechtshandelingen zijn de overboekingen van in totaal € 70.217,83 van ECPE aan ECP kort vóór faillissement. Uit de casus blijkt niet van een op wet of overeenkomst berustende, in rechte afdwingbare verplichting van ECPE om dit bedrag aan ECP over te maken. De door Fontijn genoemde reden was juist dat het beschikbare saldo `uit voorzorg' naar ECP werd overgeboekt om eventueel beslag door crediteuren te voorkomen, waarna ECP/Hoeksema zou bepalen welke urgente schuldeisers van ECPE zouden worden betaald. Dat duidt op een praktische of bedrijfseconomische keuze, niet op een rechtsplicht jegens ECP. Dat ECPE mogelijk verplichtingen had jegens eigen crediteuren, maakt de voorafgaande overboeking van haar saldo naar ECP niet verplicht: er is niet gesteld dat ECPE gehouden was die schulden via ECP te voldoen of zekerheid/gelden aan ECP te verschaffen. Morele, feitelijke of economische druk en de wens om beslag te vermijden scheppen geen rechtsplicht. De overboekingen zijn daarom onverplicht verricht in de zin van artikel 42 Fw.

\paragraph{Benadeling schuldeisers - overwegingen}\label{benadeling-schuldeisers---overwegingen-52}

De schuldeisers zijn door de overboekingen benadeeld. Zonder de aangevochten rechtshandeling zou het beschikbare banksaldo van ECPE van € 70.217,83 in beginsel op de rekening van ECPE zijn gebleven en bij faillissement beschikbaar zijn geweest voor de gezamenlijke schuldeisers. Uit de verklaring van Fontijn volgt bovendien dat het geld juist is overgemaakt naar ECP uit voorzorg tegen mogelijk beslag door crediteuren en dat ECP vervolgens zou bepalen welke urgente crediteuren zouden worden betaald. Op dat moment dreigde nog geen concreet beslag. In de feitelijke situatie is het bedrag dus aan de directe verhaalsmogelijkheden van de gezamenlijke schuldeisers onttrokken en onder controle van ECP gebracht. Dat met het bedrag mogelijk bepaalde urgente of voor voortzetting relevante schuldeisers zouden worden voldaan, neemt de benadeling niet weg: een dergelijke selectieve aanwending verstoort de paritas creditorum en komt niet aan de schuldeisers als collectief ten goede. Er zijn geen feiten gesteld waaruit volgt dat ook zonder de overboekingen alle schulden zouden zijn voldaan, dat schuldeisers zich niet op het banksaldo hadden kunnen verhalen, of dat een gelijkwaardig en daadwerkelijk verhaalbaar actief voor de boedel is teruggekomen. In de vermogensvergelijking verkeren de schuldeisers daarom slechter als de overboekingen in stand blijven.

\paragraph{Wetenschap - overwegingen}\label{wetenschap---overwegingen-52}

De betalingen vonden plaats aan ECP, meerderheidsaandeelhouder van ECPE, zodat sprake is van een gelieerde partij/groepsverhouding in de zin van art. 43 lid 1 sub 6 Fw; er blijkt niet van een reeds vóór de jaartermijn bestaande verplichting tot deze overboekingen. Daardoor wordt wetenschap van benadeling aan beide zijden vermoed en er zijn geen feiten die dit vermoeden succesvol ontzenuwen. Daarnaast volgt de wetenschap ook rechtstreeks uit de feiten: het beschikbare saldo werd uit voorzorg overgeboekt om beslag door crediteuren te voorkomen, er werden alleen nog urgente crediteuren geselecteerd betaald, in het pand waren geen activiteiten meer, de telefoon was doorgeschakeld en betalingsonmacht aan GAK/fiscus werd gemeld. Hoeksema, bestuurder van ECP, was hierbij betrokken en bepaalde uiteindelijk welke crediteuren werden betaald. Daarmee waren faillissement en een tekort daarin ten tijde van de betalingen met een redelijke mate van waarschijnlijkheid voorzienbaar voor zowel ECPE als ECP.

\newpage

\markboth{JOR 2009/46}{gemini-3.5-flash (k = 10)}

\subsubsection{Gemini}\label{gemini-17}

\texttt{gemini-3.5-flash}

\begin{center}
\small
\begin{tabular}{@{}p{7.4cm}p{2.6cm}p{2.6cm}@{}}
\toprule
Variabele & Waarde (model) & Waarde (gold) \\
\midrule
Is rechtshandeling & \texttt{true} & \texttt{true} \\
Datum rechtshandeling & \texttt{2004-01-08} & \texttt{2004-01-08} \\
Verricht door de schuldenaar & \texttt{true} & \texttt{true} \\
Om niet & \texttt{true} & \texttt{true} \\
Datum faillietverklaring & \texttt{2004-02-04} & \texttt{2004-02-04} \\
Onverplicht & \texttt{true} & \texttt{true} \\
Benadeling schuldeisers & \texttt{true} & \texttt{true} \\
Wetenschap schuldenaar aangetoond & \textcolor{gray}{\texttt{true}} & \texttt{false} \\
Wetenschap wederpartij aangetoond & \textcolor{gray}{\texttt{true}} & \texttt{null} \\
Vermoeden wetenschap schuldenaar tegenbewijs & \texttt{false} & \texttt{false} \\
Vermoeden wetenschap wederpartij tegenbewijs & \texttt{false} & \texttt{null} \\
Tegen onderwaarde & \texttt{null} & \texttt{null} \\
Voldoening of zekerheid niet opeisbare schuld & \texttt{null} & \texttt{null} \\
Gelieerde partij & \textcolor{gray}{\texttt{true}} & \texttt{null} \\
Reeds verplicht voor aanvang termijn & \texttt{false} & \texttt{null} \\
\textbf{Vernietigbaar} & \textbf{\texttt{true}} & \textbf{\texttt{true}} \\
\bottomrule
\end{tabular}
\end{center}

\paragraph{Is rechtshandeling - overwegingen}\label{is-rechtshandeling---overwegingen-53}

De girale overboekingen van de bankrekening van ECPE naar ECP zijn rechtshandelingen, aangezien deze gericht zijn op het teweegbrengen van een rechtsgevolg, namelijk de overgang van een vorderingsrecht op het creditsaldo.

\paragraph{Verricht door de schuldenaar - overwegingen}\label{verricht-door-de-schuldenaar---overwegingen-53}

De overboekingen zijn verricht vanaf de bankrekening van de schuldenaar (ECPE) en zijn geïnitieerd door haar interim-manager Fontijn. De handeling is dus door de schuldenaar zelf verricht.

\paragraph{Om niet - overwegingen}\label{om-niet---overwegingen-53}

De overboekingen zijn gedaan om het saldo veilig te stellen voor crediteuren en er is geen reële tegenprestatie van ECP jegens ECPE overeengekomen. De rechtshandeling is derhalve om niet verricht.

\paragraph{Onverplicht - overwegingen}\label{onverplicht---overwegingen-53}

De overboekingen van in totaal € 70.217,83 van de bankrekening van ECPE naar de bankrekening van haar aandeelhouder ECP zijn onverplicht verricht. Er bestond geen op de wet of een overeenkomst berustende rechtsplicht voor ECPE om haar liquide middelen aan ECP over te dragen. Uit de verklaring van de interim-manager blijkt dat de overboekingen louter zijn verricht `uit voorzorg' om te voorkomen dat crediteuren beslag zouden leggen op de bankrekening van ECPE. Omdat er geen in rechte afdwingbare verplichting bestond tot deze overboekingen, is de rechtshandeling onverplicht in de zin van artikel 42 Fw.

\paragraph{Benadeling schuldeisers - overwegingen}\label{benadeling-schuldeisers---overwegingen-53}

De overboekingen van in totaal € 70.217,83 van de bankrekening van ECPE naar de bankrekening van ECP in januari 2004, kort voor het faillissement op 4 februari 2004, hebben geleid tot benadeling van de schuldeisers. Door deze onverplichte overboekingen is het actieve vermogen van ECPE dat beschikbaar was voor verhaal door de gezamenlijke schuldeisers rechtstreeks verminderd. De overboekingen werden bovendien bewust verricht om te voorkomen dat crediteuren beslag zouden leggen op het saldo. Zelfs indien het geld naderhand is gebruikt om selectief bepaalde `urgente' crediteuren te voldoen, leidt dit tot een verstoring van de rangorde van schuldeisers en daarmee tot benadeling van de overige crediteuren die achtergesteld zijn. In de hypothetische situatie zonder deze rechtshandeling zou het volledige bedrag in de boedel zijn gevallen, wat gunstiger is voor de gezamenlijke schuldeisers dan de feitelijke situatie waarin dit actief is onttrokken.

\paragraph{Wetenschap - overwegingen}\label{wetenschap---overwegingen-53}

In deze casus is sprake van onverplichte betalingen door ECPE aan haar meerderheidsaandeelhouder ECP (65\%) binnen een maand voor het faillissement. Omdat ECP een gelieerde partij is (groepsmaatschappij in de zin van art. 43 lid 1 sub 6 Fw), geldt het wettelijk bewijsvermoeden van art. 43 lid 1 Fw. Dit vermoeden is niet ontzenuwd. Bovendien is de wetenschap van benadeling aan beide zijden rechtstreeks aangetoond. Uit de verklaring van interim-manager Fontijn blijkt immers dat de gelden bewust zijn overgeboekt naar ECP `uit voorzorg voor het geval crediteuren beslag zouden leggen'. Zowel de schuldenaar (ECPE via Fontijn) als de wederpartij (ECP via haar bestuurder Hoeksema, die de betalingen coördineerde) wisten derhalve dat het faillissement en een tekort daarin met een redelijke mate van waarschijnlijkheid te voorzien waren en dat crediteuren hierdoor zouden worden benadeeld.

\newpage

\markboth{JOR 2010/210}{Gold standard}

\section{JOR 2010/210}\label{jor-2010210}

uitspraakdatum 03-11-2009

\subsection{Gold}\label{gold-6}

\begin{center}
\small
\begin{tabular}{@{}p{8cm}p{2.6cm}@{}}
\toprule
Variabele & Waarde \\
\midrule
Is rechtshandeling & \texttt{true} \\
Datum rechtshandeling & \texttt{2002-04-01} \\
Verricht door de schuldenaar & \texttt{true} \\
Om niet & \texttt{false} \\
Datum faillietverklaring & \texttt{2003-11-19} \\
Onverplicht & \texttt{true} \\
Benadeling schuldeisers & \texttt{false} \\
Wetenschap schuldenaar aangetoond & \texttt{false} \\
Wetenschap wederpartij aangetoond & \texttt{false} \\
Vermoeden wetenschap schuldenaar tegenbewijs & \texttt{null} \\
Vermoeden wetenschap wederpartij tegenbewijs & \texttt{null} \\
Tegen onderwaarde & \texttt{null} \\
Voldoening of zekerheid niet opeisbare schuld & \texttt{null} \\
Gelieerde partij & \texttt{null} \\
Reeds verplicht voor aanvang termijn & \texttt{null} \\
\textbf{Vernietigbaar} & \textbf{\texttt{false}} \\
\bottomrule
\end{tabular}
\end{center}

Opmerking: deze casus wordt in het kader van de faillissementspauliana en art. 42 Fw beperkt tot de vraag of de rechtshandeling omtrent de BMW paulianeus is.

\subsubsection{Uitspraak}\label{uitspraak-6}

(\ldots; red.)

\paragraph{4. De beoordeling}\label{de-beoordeling-2}

In principaal en incidenteel appel
(\ldots; red.)

4.2. Het gaat in dit hoger beroep om het volgende.

4.2.1. Op 28 mei 1999 is de rechtsvoorganger (Advisory Management B.V.) van Brand Station B.V., hierna Brand, opgericht.

Sinds 24 april 2002 is Van Deursen bestuurder en enig aandeelhouder van Brand.

4.2.2. In april 2002 koopt Brand in het kader van een doorstart van de curator in het faillissement van Holland Wear House B.V., verder te noemen HWH, activa uit dit faillissement. HWH was een import- en exportgroothandel in textiel en richtte zich in hoofdzaak op de verkoop van kleding van de merken Canoe en Locai, van welke merken HWH destijds de intellectuele eigendom bezat.

Van Deursen was destijds tot aan het faillissement (uitgesproken op 27 maart 2002) bestuurder van HWH.

4.2.3. Tot de door Brand overgenomen activa behoorde onder meer een auto van het merk BMW en met het kenteken XX-PG-65. Brand heeft deze BMW in april/mei 2002 doorverkocht aan Jon West voor een prijs van € 10.584,09 inclusief BTW (prod. 1 inl. dagv.).

4.2.4. Van Deursen is tevens enig bestuurder/aandeelhouder van John West en van Operations B.V.

4.2.5. John West, toen History Division B.V. genaamd, heeft voor de leasetermijn over april 2002 van de BMW aan Brand in rekening gebracht een bedrag van € 2.023,--.

4.2.6. In de periode van 3 juni 2003 tot en met 17 november 2003 zijn door Brand aan Operations B.V. kasbetalingen verricht voor een totaalbedrag van € 28.608,60 (prod. 5 inl. dagv.) en in de periode van 5 juni 2003 tot en met 13 augustus 2003 girale betalen tot een bedrag van € 52.076,58 (prod. 2 cva).

4.2.7. Ten behoeve van de financiering van de bedrijfsactiviteiten van Brand is door NMB Heller een kredietarrangement beschikbaar gesteld, onder meer inhoudende een bevoorschotting op de facturen.

Bij brief van 5 maart 2003 heeft NMB Heller aan Brand onder meer het volgende mede bericht:

``(\ldots)

Zoals mondeling toegelicht bevestigen wij hierdoor de bevoorschottings- en dienstverleningsovereenkomst met Brand Station B.V. te beëindigen. Onze vordering op u is hierdoor direct opeisbaar.

Teneinde u in staat te stellen maatregelen te treffen die op kort termijn zullen bijdragen aan herstel van continuïteit van uw onderneming zijn wij bereid, tot nader bericht, niet tot het opeisen van onze vordering over te gaan. Wij zullen u, in afwachting van een reorganisatieplan, onder de navolgende voorwaarden de gelegenheid geven de bedrijfsvoering te continueren.

(\ldots)

-- voor 1 april aanstaande zal overschrijding van de kredietfaciliteit, thans € 1.100.000,=, met ten minste € 500.000,= worden teruggebracht

(\ldots)''

In de tweede helft van oktober 2003 heeft NMB Heller geweigerd de financiering van Brand verder te continueren.

4.2.8. Bij vonnis van de rechtbank 's-Hertogenbosch van 19 november 2003 is Brand in staat van faillissement verklaard met aanstelling van mr. R.G. Roeffen tot curator. Operations B.V. is in staat van faillissement verklaard bij vonnis van 1 februari 2006 met aanstelling van mr. P.R. Dekker tot curator. Dit laatste faillissement is opgeheven wegens gebrek aan baten.

4.3. In onderhavige procedure heeft de curator na wijziging van eis in eerste aanleg gevorderd:

\begin{itemize}
\tightlist
\item
  1. Van Deursen en History Division B.V., thans genaamd John West B.V., hoofdelijk te veroordelen, des dat de een betalend de ander zal zijn bevrijd, om aan de curator te voldoen een bedrag van € 22.023,--, te vermeerderen met de wettelijke rente vanaf 10 december 2003, althans een datum in goede justitie te bepalen, althans vanaf de dag der dagvaarding, met veroordeling in de kosten waaronder de kosten van de ten laste van History Division B.V. gelegde beslagen, vermeerderd met de wettelijke rente vanaf de 15e dag na de dag der uitspraak.
\end{itemize}

2. Van Deursen en Operations B.V. hoofdelijk te veroordelen om aan de curator te betalen een bedrag ad € 28.608,60 en een bedrag ad € 52.076,58, te vermeerderen met de wettelijke rente vanaf 1 december 2003, althans een in goede justitie te bepalen datum, althans vanaf de dag der dagvaarding, met veroordeling in de kosten, waaronder begrepen de kosten van de ten laste van Operations B.V. gelegde beslagen, vermeerderd met de wettelijke rente vanaf de 15e dag na de dag der uitspraak.

4.4.1. De curator stelt met betrekking tot de vordering sub 1 jegens History B.V. dat de BMW beneden de reële waarde is verkocht waardoor de schuldeisers van Brand zijn benadeeld, dat hij op grond van artikel 43 jo. 42 Faillissementswet (Fw) de nietigheid heeft ingeroepen van de koopovereenkomst, alsmede dat de leasetermijn van € 2.023,-- onverschuldigd is betaald.

4.4.2. Met betrekking tot de vordering sub 2 jegens Operations B.V. stelt de curator dat Brand in het zicht van haar faillissement betalingen heeft gedaan waardoor Operations B.V. boven andere schuldeisers werd begunstigd en dat deze betalingen op grond van artikel 47 jo. 42 Fw kunnen worden vernietigd nu Van Deursen zowel Brand als Operations B.V. vertegenwoordigde.

4.4.3. Met betrekking tot de vorderingen jegens Van Deursen stelt de curator primair dat Van Deursen onrechtmatig heeft gehandeld door als bestuurder van Brand de in de vorderingen vermelde rechtshandelingen te verrichten en subsidiair dat Van Deursen persoonlijk aansprakelijk is op grond van artikel 2:9 BW wegens kennelijk onbehoorlijke vervulling van zijn taak als bestuurder.

4.5. De rechtbank heeft geoordeeld dat de procedure ten aanzien van Operations B.V. is geschorst. Voorts heeft de rechtbank de vordering ad € 2.023,-- vermeerderd met de wettelijke rente vanaf 1 december 2003 jegens History B.V. (thans: John West B.V.) toegewezen en de overige vorderingen afgewezen. Tenslotte heeft de rechtbank de proceskosten gecompenseerd zodanig dat iedere partij de eigen kosten draagt.

4.6. Naast de hierna te bespreken en duidelijk als grief aangeduide incidentele grief met betrekking tot de compensatie van kosten, ontwaart het hof nog de incidentele grief van Van Deursen en John West dat de curator niet-ontvankelijk is in zijn vordering jegens Van Deursen met betrekking tot de BMW en de girale betalingen. Gelet op hetgeen door de curator bij pleidooi is aangevoerd, heeft de curator een en ander ook als zodanig opgevat.

BMW
4.7. Van Deursen en John West betogen dat indien de BMW in het vermogen van Brand zou zijn gebleven op de BMW een pandrecht ten gunste van NMB Heller zou hebben gerust, en dat NMB Heller haar vordering dan had kunnen verhalen door uitwinning van deze auto. Dit zou hebben geleid tot een verlaging van de schuld van Brand aan NMB Heller, zodat slechts NMB Heller in haar verhaalsrecht zijn benadeeld en niet de overige schuldeisers.

4.7.1. Hoewel Van Deursen en John West die conclusie niet met zoveel woorden trekken, begrijpt het hof dat de strekking van dit betoog is dat de curator niet-ontvankelijk is in zijn vordering jegens Van Deursen op grond van onrechtmatige daad voor wat betreft de BMW.

4.7.2. Het betoog van Van Deursen en John West faalt. Immers, voordat NMB Heller haar schuld als pandhouder zou hebben kunnen verhalen op de BMW, diende aan alle constitutieve vereisten voor pandvestiging te zijn voldaan en wel voor de dag der faillietverklaring. De auto had daarom voor 19 november 2003 in de feitelijke macht moeten zijn genomen door NMB Heller, dan wel door haar zijn opgeëist (artikel 3:236 en 237 BW) om zich als separatist (buiten de faillissementsboedel om) op de auto te kunnen verhalen. Het verwijt van de curator aan Van Deursen is echter niet dat de auto ondanks het daarop rustende pandrecht is verkocht, maar dat deze voor een te lage prijs is verkocht. Door de in de ogen van de curator te lage opbrengst zijn derhalve niet alleen de bank maar alle schuldeisers van Brand in hun verhaalsrechten benadeeld.

Girale betalingen aan Operations
4.8. Van Deursen en John West betogen voorts dat de girale betalingen aan Operations B.V. plaats vonden ten laste van een reeds bestaand debetsaldo bij de bank, zodat het slechts de bank is die schade heeft geleden en niet de overige schuldeisers omdat hun verhaalspositie door de betalingen niet is gewijzigd. De concurrente crediteur Operations B.V. is slechts vervangen door een andere concurrente crediteur, de bank.

4.8.1. De curator is bevoegd voor de belangen van de gezamenlijke schuldeisers op te komen door het geldend maken van een vordering uit hoofde van artikel 6:162 BW bij benadeling van de gezamenlijke schuldeisers in hun verhaalsmogelijkheden (HR 14 januari 1983, NJ 1983, 597 Peeters q.q./Gatzen en HR 24 april 2009, NJ 2009, 416 («JOR» 2010/22, m.nt. NEDF (Dekker/Lutèce); red.)). Een dergelijke benadeling in de verhaalsmogelijkheden kan onder meer bestaan in een doorbreking van de paritas creditorum. Daarvan is in dit geval sprake. Immers, door de girale betalingen is Operations B.V. mogelijk boven andere concurrente schuldeisers begunstigd en dan zijn deze andere concurrente schuldeisers in vergelijking met Operations B.V. benadeeld. De curator is derhalve ontvankelijk in dit deel van zijn vordering.

4.9. Met betrekking tot de afwijzing van de vordering op grond van artikel 42 Fw ten aanzien van de BMW heeft de rechtbank overwogen dat niet vaststaat dat er sprake is van wetenschap van benadeling van schuldeisers ten tijde van de verkoop.

Hiertegen richt zich de principale grief I.

4.9.1. In de toelichting op de principale grief I betoogt de curator dat met de wetenschap (naar het hof begrijpt: bij de handelende partijen) dat de BMW is verkocht voor een prijs beneden de reële waarde de wetenschap van benadeling van de crediteuren is gegeven.

Dit betoog faalt. De juistheid van de stelling dat er sprake was van wetenschap van een verkoop beneden de reële waarde daargelaten -- die stelling is door Van Deursen en John West betwist -- impliceert deze stelling nog niet dat er sprake is van wetenschap van benadeling van crediteuren. Immers, voor die wetenschap van benadeling is meer nodig, namelijk dat ten tijde van de verkoop met een aanzienlijke mate van waarschijnlijkheid was te verwachten dat de schuldeisers niet zouden kunnen worden voldaan. Dit is bijvoorbeeld het geval indien de mogelijkheid van het verwerven van inkomsten niet meer aanwezig is doordat bijvoorbeeld de activiteiten zijn gestaakt ten gevolge waarvan geen inkomsten (meer) worden gegenereerd, of in het geval dat een faillissement onafwendbaar is. Daaromtrent is door de curator echter niets gesteld, terwijl voorts ook niet is gebleken dat al ten tijde van de verkoop van de BMW sprake was van een dergelijke situatie. Het hof onderschrijft op dit punt het oordeel van de rechtbank als vervat in rechtsoverweging 4.3.

De wetenschap omtrent de mogelijkheid van een faillissement en daarmede een mogelijkheid van benadeling is onvoldoende (HR 17 november 2000, NJ 2001, 272 («JOR» 2001/17 (Bakker/Katko); red.)).

Derhalve faalt de principale grief I.

4.10. De principale grief II heeft betrekking op de afwijzing van de vordering tegen Van Deursen wegens onrechtmatig handelen door de betaling van de leasetermijn ad € 2.023,--.

Het hof zal deze grief bespreken tezamen met de bezwaren tegen de afwijzing van alle vorderingen jegens Van Deursen op grond van artikel 6:162 BW (onrechtmatige daad).

4.11. Voor wat betreft de vorderingen van de curator met betrekking tot de girale en kasbetalingen is Operations B.V. wegens de schorsing ex artikel 29 Fw in hoger beroep geen partij meer. Derhalve zijn de vorderingen vermeld onder 4.3. jegens Van Deursen op grond van onrechtmatige daad nog slechts aan de orde de vordering met betrekking tot de girale en kasbetalingen, de leasetermijn ad € 2.023,-- en verder de vordering met betrekking tot de BMW.

4.12. De grieven III en IIII, verder aan te duiden als IV, hebben betrekking op de afwijzing van deze vorderingen jegens Van Deursen.

4.13. De curator heeft gesteld dat door de onder 4.2.6. vermelde betalingen aan Operations B.V. de andere crediteuren van Brand zijn benadeeld ten aanzien van welke betalingen de curator de pauliana heeft ingeroepen en dat deze betalingen het gevolg zijn van overleg tussen Brand en Operations B.V. met als doel Operations B.V. boven andere schuldeisers te begunstigen.

4.13.1. Omtrent de financiële situatie van Brand heeft de curator het volgende aangevoerd.

In het startjaar van de onderneming, het jaar 2002, werd een verlies geleden van ruim 1,7 miljoen euro (mvg onder 4.6. en prod. 18). Tot en met mei 2003 is het verlies opgelopen met een bedrag ad € 591.940,-- (prod. 18). Per eind mei 2003 bedroeg de schuld aan NMB-Heller een bedrag van ruim 2,2 miljoen euro (proces-verbaal cvpp en prod. 11). Per 27 augustus 2003 was er een schuld aan crediteuren van ruim € 399.000,-- waarvan een bedrag ad € 254.408,48 langer dan 90 dagen open stond (proces-verbaal cvpp, mvg onder 4.10. en prod 12). Tegenover deze hoge schuldenlast waren er nauwelijks activa. De voorraden en de intellectuele eigendomsrechten zijn door de curator verkocht en hebben respectievelijk opgebracht € 125.000,-- en € 13.570,--, welke opbrengsten onder inhouding van een boedelbijdrage zijn afgedragen aan de bank (mvg onder 4.7.). Per faillissementsdatum stond er een bedrag ad € 468.876,28 aan debiteuren open, waarvan een bedrag ad € 321.970,21 langer dan 90 dagen (proces-verbaal cvpp en prod. 16).

Na uitwinning van de zekerheden bedraagt de restschuld aan NMB-Heller circa 1,8 miljoen euro.

4.13.2. Bij brief van 5 maart 2003 heeft NMB-Heller aan Brand medegedeeld dat de kredietlimiet per 1 april 2003 werd teruggebracht naar (bedoeld zal zijn met, hof) € 500.000,- (inl. dagv. onder 3.2. en prod. 6).

4.13.3. Voorts heeft de curator gesteld dat vanaf april 2003 de fiscus niet meer werd betaald (proces-verbaal cvpp en prod. 14), terwijl de managementvergoedingen (over 2002 van ruim € 84.000,-- en over 2003 van ruim € 130.000,--) geheel aan Operations B.V. zijn uitbetaald.

4.14. De door de curator aangevoerde feiten en omstandigheden zoals vermeld onder onder 4.13.1. tot en met 4.13.3. zijn door Van Deursen en John West op zich niet bestreden. Zij hebben slechts betwist dat uit deze feiten en omstandigheden volgt dat er sprake was van wetenschap van benadeling als bedoeld in artikel 42 Fw.

4.15. Het hof overweegt als volgt.

Uit de onder 4.13.1. tot en met 4.13.3. aangevoerde feiten en omstandigheden volgt dat een belangrijk deel van het vanaf juni 2003 beschikbare liquide actief dat anders voor de gezamenlijke schuldeisers ter verdeling beschikbaar zou zijn gebleven, is aangewend voor de betaling van de vorderingen van Operations B.V. Door de betalingen aan Operations B.V. zijn de andere concurrente schuldeisers dan Operations B.V. benadeeld.

4.16. Uit voormelde feiten en omstandigheden volgt eveneens dat in de periode vanaf begin juni 2003 tot aan het faillissement de financiële situatie van Brand zeer slecht was, hetgeen aan Van Deursen als bestuurder van Brand kenbaar was. Immers, de verliezen evenals de schulden liepen steeds verder op, terwijl Brand reeds bij brief van 3 maart 2003 was aangezegd dat op korte termijn de schuld met € 500.000,-- zou moeten worden verminderd, terwijl gesteld noch gebleken is dat er in de periode waarin de betalingen hebben plaats gevonden uitzicht bestond op verbetering van de financiële situatie van Brand als ook dat aan deze eis van NMB-Heller zou kunnen worden voldaan.

4.17. De onder 4.13.1. tot en met 4.16. omschreven feiten en omstandigheden in onderlinge samenhang beschouwd, alsmede het feit dat de betalingen hebben plaats gevonden tussen vennootschappen met dezelfde aandeelhouder/bestuurder en de omstandigheid dat gesteld noch gebleken is dat er voor Brand dringende redenen waren om de opeisbare vorderingen van Operations B.V. boven die van andere concurrente schuldeisers te voldoen, leiden tot het oordeel dat het oogmerk tot begunstiging van Operations B.V. boven de andere concurrente schuldeiser zowel bij Brand als bij Operations B.V. in de persoon van Van Deursen bij de betalingen heeft voorgezeten en dat Van Deursen Brand terzake een ernstig verwijt treft.

4.18. De slotsom is dan ook dat de vorderingen jegens Van Deursen op grond van onrechtmatige daad wegens de kas- en girale betalingen toewijsbaar zijn, zodat de grieven III en IV slagen.

4.19. Voor wat betreft de door de curator voor de aansprakelijkheid van Van Deursen met betrekking tot de betaling van de factuur van John West betreffende de leasetermijn over april 2002 overweegt het hof dat er zowel ingeval van aansprakelijkheid ex artikel 2:9 BW (onbehoorlijke taakvervulling als bestuurder) als aansprakelijkheid ex artikel 6:162 BW (onrechtmatige daad) sprake moet zijn van een ernstig verwijt van de bestuurder (HR 10 januari 1997, NJ 1997, 360 («JOR» 1997/29 (Staleman c.s./Van de Ven; red.)). Daarvan is naar het oordeel van het hof echter geen sprake nu de betaling reeds plaats vond in april 2002, de maand waarin de onderneming door Brand na overname werd voortgezet. Het enkele feit dat Brand vanaf 2002 verlies heeft geleden en dat deze verliezen steeds meer opliepen, is daartoe niet voldoende. Immers, het is, zoals Van Deursen en John West terecht aanvoeren, niet ongebruikelijk dat een startende onderneming te maken krijgt met aanvangsverliezen, maar dat men er desondanks in slaagt na verloop van tijd winst te maken.

Grief II faalt derhalve.

4.20. Met betrekking tot vernietiging van de koopovereenkomst van de BMW heeft de curator in hoger beroep de grondslag in hoger beroep vermeerderd in die zin dat hij deze thans tevens grondt op artikel 2:247 BW.

4.21. Van Deursen heeft bij memorie van antwoord geconcludeerd tot niet-ontvankelijkverklaring van de curator in zijn vermeerderde grondslag van de vorderingen, althans tot afwijzing van deze vermeerdering van grondslag. Deze conclusie ontbeert evenwel iedere onderbouwing. De oorspronkelijk eiser heeft de bevoegdheid zijn eis in hoger beroep te wijzigen of te vermeerderen bij memorie van grieven (HR 22 juni 2009, LJN BI8771), tenzij deze wijziging of vermeerdering in strijd is met de goede procesorde. Nu gesteld noch gebleken is dat de eisvermeerdering in strijd is met de goede procesorde, zal het hof op de vermeerderde grondslag beslissen.

4.22. Van Deursen en John West beroepen zich met betrekking tot deze vordering voorts op verjaring ex artikel 3:52 lid 1 sub d.~Nu de curator eerst bij memorie van grieven van 15 juli 2008 de nietigheid heeft ingeroepen, is de vordering verjaard, aldus Van Deursen en John West.

De curator heeft bij pleidooi hiertegen aangevoerd dat de verjaring is gestuit door zijn brief van 10 december 2003.

4.23. De betreffende vordering jegens Van Deursen is een vordering tot schadevergoeding, gebaseerd op een (buitengerechtelijke) vernietiging van een overeenkomst. Dat de grondslag voor deze vernietiging naderhand wordt uitgebreid, maakt dit niet anders (vgl. (HR 11 januari 2002, NJ 2002, 81 («JOR» 2002/94 (Van Dam \& Post/Coöperatieve Land- en Tuinbouwvereniging); red.)).

Ingevolge artikel 3:317 lid 1 BW wordt de verjaringstermijn van een dergelijke vordering gestuit door een schriftelijke aanmaning of mededeling waarbij waarin de schuldeiser zich ondubbelzinnig zijn recht op nakoming voorbehoudt. De brief van 10 december 2003 heeft een dergelijke strekking en heeft derhalve stuitende werking, zodat een nieuwe verjaringstermijn is gaan lopen. Vóór het verstrijken van deze nieuwe termijn is onderhavige procedure aanhangig gemaakt bij dagvaarding d.d. 28 februari 2005. Mitsdien is de vordering niet verjaard.

4.24. De strekking van het vereiste van een schriftelijke vastlegging in artikel 2:247 BW is zo veel mogelijk te waken voor het gevaar dat bij een eenpersoonsvennootschap de enige aandeelhouder handelingen verricht ten eigen bate en ten nadele van de schuldeisers. Naar het oordeel van het hof is aan dit vereiste voldaan, nu van de overeenkomst blijkt middels de factuur (prod. 1 inl. dagv.) die kennelijk in de administratie van Brand was opgenomen, nu de curator de factuur heeft overgelegd.

4.25. Uit het voorgaande volgt dat van de vorderingen van de curator, voor zover in hoger beroep aan de orde, slechts toewijsbaar zijn de vorderingen jegens Van Deursen met betrekking tot de girale en kasbetalingen aan Operations B.V.

4.26. In de zaak van de curator tegen Van Deursen dient laatstgenoemde te worden aangemerkt als de in overwegende mate in het ongelijk gestelde partij en in de zaak van de curator tegen John West, de curator. Nu zowel Van Deursen als John West tezamen verweer hebben gevoerd zowel in eerste aanleg als in hoger beroep, acht het hof het redelijk de kosten aan vast recht en aan salaris procureur van de curator in de eerste zaak te stellen op de helft van de totale kosten van de curator en in de tweede zaak op de helft van de verschotten en de kosten aan salaris procureur van Van Deursen en John West.

4.27. De conclusie is dat de principale grieven III en IV geheel en grief V deels slagen en de overige principale grieven, alsmede de verborgen incidentele grief falen en dat de uitdrukkelijke incidentele grief deels slaagt.

4.28. De slotsom is dat het vonnis waarvan beroep dient te worden vernietigd voorzover de vordering betreffende de girale en de kasbetalingen ad € 80.685,18 jegens Van Deursen is afgewezen en voorzover de proceskosten zijn gecompenseerd en dat dit vonnis voor het overige dient te worden bekrachtigd.

\paragraph{5. De uitspraak}\label{de-uitspraak}

op het principaal en het incidenteel appel

Het hof:

vernietigt het vonnis waarvan beroep voorzover de vordering jegens Van Deursen tot een bedrag ad € 80.685,18 is afgewezen en de proceskosten zijn gecompenseerd;

en opnieuw rechtdoende:

\begin{itemize}
\item
  1. veroordeelt Van Deursen tot betaling aan de curator van een bedrag ad € 80.685,18, vermeerderd met de wettelijke rente ex artikel 6:119 BW vanaf 1 december 2003;
\item
  2. veroordeelt Van Deursen in de proceskosten van de eerste aanleg aan de zijde van de curator (\ldots; red.);
\item
  3. veroordeelt de curator in de proceskosten in eerste aanleg aan de zijde John West (\ldots; red.);
\item
  4. bekrachtigt het vonnis waarvan beroep voor het overige;
\item
  5. veroordeelt Van Deursen in de proceskosten van het principale beroep aan de zijde van de curator (\ldots; red.);
\item
  6. veroordeelt de curator in de proceskosten van het incidentele hoger beroep aan de zijde van John West (\ldots; red.).
\end{itemize}

verklaart dit arrest voor wat betreft de punten 1, 2 en 5 uitvoerbaar bij voorraad.

\newpage

\markboth{JOR 2010/210}{mistral/mistral-medium-3.5-128b:fp8 (k = 0)}

\subsection{k = 0}\label{k-0-6}

\subsubsection{Mistral}\label{mistral-18}

\texttt{mistral/mistral-medium-3.5-128b:fp8}

\begin{center}
\small
\begin{tabular}{@{}p{7.4cm}p{2.6cm}p{2.6cm}@{}}
\toprule
Variabele & Waarde (model) & Waarde (gold) \\
\midrule
Is rechtshandeling & \texttt{true} & \texttt{true} \\
Datum rechtshandeling & \texttt{2002-04-01} & \texttt{2002-04-01} \\
Verricht door de schuldenaar & \texttt{true} & \texttt{true} \\
Om niet & \texttt{false} & \texttt{false} \\
Datum faillietverklaring & \texttt{2003-11-19} & \texttt{2003-11-19} \\
Onverplicht & \texttt{true} & \texttt{true} \\
Benadeling schuldeisers & \textcolor{gray}{\texttt{true}} & \texttt{false} \\
Wetenschap schuldenaar aangetoond & \texttt{false} & \texttt{false} \\
Wetenschap wederpartij aangetoond & \texttt{false} & \texttt{false} \\
Vermoeden wetenschap schuldenaar tegenbewijs & \texttt{null} & \texttt{null} \\
Vermoeden wetenschap wederpartij tegenbewijs & \texttt{null} & \texttt{null} \\
Tegen onderwaarde & \texttt{null} & \texttt{null} \\
Voldoening of zekerheid niet opeisbare schuld & \texttt{null} & \texttt{null} \\
Gelieerde partij & \textcolor{gray}{\texttt{true}} & \texttt{null} \\
Reeds verplicht voor aanvang termijn & \texttt{null} & \texttt{null} \\
\textbf{Vernietigbaar} & \textbf{\texttt{false}} & \textbf{\texttt{false}} \\
\bottomrule
\end{tabular}
\end{center}

\paragraph{Is rechtshandeling - overwegingen}\label{is-rechtshandeling---overwegingen-54}

De doorverkoop van de BMW door Brand aan John West is een koopovereenkomst, wat een rechtshandeling is (art. 3:33 BW).

\paragraph{Verricht door de schuldenaar - overwegingen}\label{verricht-door-de-schuldenaar---overwegingen-54}

Brand (de schuldenaar) heeft de BMW doorverkocht aan John West.

\paragraph{Om niet - overwegingen}\label{om-niet---overwegingen-54}

De BMW is verkocht voor € 10.584,09 inclusief BTW, wat een tegenprestatie inhoudt. Derhalve is de handeling niet om niet.

\paragraph{Onverplicht - overwegingen}\label{onverplicht---overwegingen-54}

De verkoop van de BMW door Brand aan John West (toen nog History Division B.V.) voor € 10.584,09 is een onverplichte rechtshandeling in de zin van art. 42 lid 1 Fw. Er bestond geen op de wet of een overeenkomst berustende rechtsplicht voor Brand om de BMW te verkopen, laat staan tegen een prijs die mogelijk beneden de marktwaarde lag. De curator stelt dat de verkoop beneden de reële waarde heeft plaatsgevonden, wat duidt op een onverplichte handeling. Bovendien was er geen sprake van een verplichting om de auto te verkopen; Brand had de auto kunnen behouden of op een andere (verplichte) wijze kunnen afstoten. De leasetermijn van € 2.023,-- die John West in rekening heeft gebracht aan Brand, is eveneens onverplicht betaald, omdat er geen contractuele verplichting bestond om deze specifieke termijn te betalen in het kader van de doorstart. De stelplicht en bewijslast rusten op de curator, maar op basis van de gegevens in de casus is aannemelijk dat de verkoop en betaling niet op een verplichting berustten.

\paragraph{Benadeling schuldeisers - overwegingen}\label{benadeling-schuldeisers---overwegingen-54}

De verkoop van de BMW door Brand aan John West voor € 10.584,09 inclusief BTW moet worden beoordeeld op benadeling van schuldeisers in de zin van artikel 42 Fw. Uit de casus blijkt dat de BMW afkomstig was uit de doorstart van HWH, waar Van Deursen bestuurder was. De prijs van € 10.584,09 lijkt aanzienlijk lager dan de reële marktwaarde van een BMW (vooral als deze afkomstig is uit een faillissementsboedel, waar activa vaak onder waarde worden verkocht). Bovendien was John West (toen nog History Division B.V.) een onderneming waar Van Deursen ook bestuurder/aandeelhouder van was, wat wijst op een mogelijke verbondenheid en een risico op benadeling. De curator stelt expliciet dat de BMW beneden de reële waarde is verkocht, wat een vermindering van het actief van Brand betekent. Dit leidt tot benadeling van de schuldeisers, omdat het vermogen van Brand (en dus de boedel) kleiner is dan zonder deze transactie het geval zou zijn geweest. Daarnaast is er sprake van een leasetermijn van € 2.023,-- die door John West in rekening is gebracht aan Brand, wat mogelijk een extra vermogensverplaatsing inhouden. Gelet op de verbondenheid tussen Van Deursen, Brand en John West, en de lage verkoopprijs, is er voldoende aanleiding om aan te nemen dat de schuldeisers zijn benadeeld door deze rechtshandeling.

\paragraph{Wetenschap - overwegingen}\label{wetenschap---overwegingen-54}

De verkoop van de BMW door Brand (schuldenaar) aan John West (wederpartij) is verricht binnen een jaar voor de faillietverklaring (19 november 2003). Van Deursen was enig bestuurder/aandeelhouder van zowel Brand als John West, waardoor sprake is van een gelieerde partij in de zin van art. 43 lid 1 sub 4° Fw (bestuurder van de schuldenaar die tevens bestuurder/aandeelhouder is van de wederpartij). Het bewijsvermoeden van art. 43 Fw is daarom van toepassing: de wetenschap van benadeling wordt aan beide zijden (Brand en John West) vermoed. Er is geen sprake van een tegenbewijs dat dit vermoeden ontzenuwt, noch van een eerdere verplichting tot de rechtshandeling. De wetenschap van de schuldenaar en wederpartij hoeft derhalve niet rechtstreeks te worden aangetoond.

\newpage

\markboth{JOR 2010/210}{gpt-5.5 (k = 0)}

\subsubsection{GPT}\label{gpt-18}

\texttt{gpt-5.5}

\begin{center}
\small
\begin{tabular}{@{}p{7.4cm}p{2.6cm}p{2.6cm}@{}}
\toprule
Variabele & Waarde (model) & Waarde (gold) \\
\midrule
Is rechtshandeling & \texttt{true} & \texttt{true} \\
Datum rechtshandeling & \texttt{2002-04-01} & \texttt{2002-04-01} \\
Verricht door de schuldenaar & \texttt{true} & \texttt{true} \\
Om niet & \texttt{false} & \texttt{false} \\
Datum faillietverklaring & \texttt{2003-11-19} & \texttt{2003-11-19} \\
Onverplicht & \texttt{true} & \texttt{true} \\
Benadeling schuldeisers & \textcolor{gray}{\texttt{true}} & \texttt{false} \\
Wetenschap schuldenaar aangetoond & \texttt{false} & \texttt{false} \\
Wetenschap wederpartij aangetoond & \texttt{false} & \texttt{false} \\
Vermoeden wetenschap schuldenaar tegenbewijs & \texttt{null} & \texttt{null} \\
Vermoeden wetenschap wederpartij tegenbewijs & \texttt{null} & \texttt{null} \\
Tegen onderwaarde & \texttt{false} & \texttt{null} \\
Voldoening of zekerheid niet opeisbare schuld & \texttt{null} & \texttt{null} \\
Gelieerde partij & \textcolor{gray}{\texttt{true}} & \texttt{null} \\
Reeds verplicht voor aanvang termijn & \texttt{null} & \texttt{null} \\
\textbf{Vernietigbaar} & \textbf{\texttt{false}} & \textbf{\texttt{false}} \\
\bottomrule
\end{tabular}
\end{center}

\paragraph{Is rechtshandeling - overwegingen}\label{is-rechtshandeling---overwegingen-55}

De gewraakte handeling betreft de verkoop/koopovereenkomst van de BMW door Brand aan John West. Een koopovereenkomst is gericht op rechtsgevolg en kwalificeert als rechtshandeling.

\paragraph{Verricht door de schuldenaar - overwegingen}\label{verricht-door-de-schuldenaar---overwegingen-55}

De BMW behoorde tot de door Brand overgenomen activa en Brand heeft deze BMW in april/mei 2002 doorverkocht aan John West. De rechtshandeling is dus door de schuldenaar Brand verricht.

\paragraph{Om niet - overwegingen}\label{om-niet---overwegingen-55}

Brand heeft de BMW verkocht voor een prijs van € 10.584,09 inclusief btw. Er is dus een tegenprestatie bedongen; dat de curator stelt dat de verkoopprijs beneden de reële waarde lag, maakt de verkoop op basis van de gegeven feiten niet om niet.

\paragraph{Onverplicht - overwegingen}\label{onverplicht---overwegingen-55}

De aangevochten rechtshandeling omtrent de BMW betreft de verkoop/doorverkoop door Brand aan John West. Uit de casus blijkt niet van een op de wet of een eerdere overeenkomst berustende, in rechte afdwingbare verplichting van Brand om de BMW aan John West te verkopen, laat staan tegen de overeengekomen prijs. Dat Van Deursen bij beide vennootschappen betrokken was, dat de transactie mogelijk bedrijfsmatig of praktisch wenselijk werd geacht, of dat John West een leasetermijn in rekening bracht, schept op zichzelf geen rechtsplicht tot verkoop. Nu geen vooraf bestaande afdwingbare verplichting tot het verrichten van deze verkoop blijkt, is de rechtshandeling in de zin van art. 42 Fw onverplicht verricht.

\paragraph{Benadeling schuldeisers - overwegingen}\label{benadeling-schuldeisers---overwegingen-55}

Bij de vermogensvergelijking moet de feitelijke situatie na instandhouding van de verkoop van de BMW worden vergeleken met de hypothetische situatie zonder die verkoop. Uit de casus volgt dat Brand de BMW aan de gelieerde vennootschap John West heeft verkocht voor € 10.584,09 inclusief btw, terwijl de curator stelt dat dit beneden de reële waarde was. Indien de BMW zonder deze rechtshandeling in het vermogen van Brand was gebleven, dan hadden de gezamenlijke schuldeisers zich op de BMW of op een marktconforme verkoopopbrengst daarvan kunnen verhalen. Door verkoop onder de reële waarde is het actief van Brand verminderd met ten minste het verschil tussen de werkelijke waarde en de ontvangen koopprijs. Er zijn geen feiten gesteld waaruit volgt dat de schuldeisers zich ook zonder de verkoop niet op de BMW hadden kunnen verhalen, bijvoorbeeld door zekerheidsrechten van derden, of dat alle schulden ondanks de verkoop volledig voldaan hadden kunnen worden. De transactie leidt daarom tot minder verhaal voor de schuldeisers in het algemeen en levert benadeling op in de zin van artikel 42 Fw.

\paragraph{Wetenschap - overwegingen}\label{wetenschap---overwegingen-55}

De BMW is door Brand verkocht aan John West, terwijl Van Deursen bestuurder/enig aandeelhouder van zowel Brand als John West was; daarmee is sprake van gelieerdheid in de zin van art. 43 lid 1 sub 5 Fw. De verkoop vond echter ruim vóór de faillissementsdatum plaats en uit de casus blijkt niet dat ten tijde van de verkoop het faillissement van Brand en een tekort daarin met een redelijke mate van waarschijnlijkheid voorzienbaar waren. De latere financieringsproblemen bij NMB Heller blijken pas vanaf maart/oktober 2003. Een oogmerk of enkele kans op benadeling volstaat niet. De gestelde verkoop beneden reële waarde is in de gegeven feiten niet concreet genoeg onderbouwd om zelfstandig rechtstreekse wetenschap van benadeling bij Brand of John West aan te nemen. Art. 45 Fw is niet aan de orde, omdat het om een verkoop tegen betaling gaat. Een toepasselijk bewijsvermoeden dat moet worden ontzenuwd speelt hier daarom niet.

\newpage

\markboth{JOR 2010/210}{gemini-3.5-flash (k = 0)}

\subsubsection{Gemini}\label{gemini-18}

\texttt{gemini-3.5-flash}

\begin{center}
\small
\begin{tabular}{@{}p{7.4cm}p{2.6cm}p{2.6cm}@{}}
\toprule
Variabele & Waarde (model) & Waarde (gold) \\
\midrule
Is rechtshandeling & \texttt{true} & \texttt{true} \\
Datum rechtshandeling & \texttt{2002-04-01} & \texttt{2002-04-01} \\
Verricht door de schuldenaar & \texttt{true} & \texttt{true} \\
Om niet & \texttt{false} & \texttt{false} \\
Datum faillietverklaring & \texttt{2003-11-19} & \texttt{2003-11-19} \\
Onverplicht & \texttt{true} & \texttt{true} \\
Benadeling schuldeisers & \textcolor{gray}{\texttt{true}} & \texttt{false} \\
Wetenschap schuldenaar aangetoond & \texttt{false} & \texttt{false} \\
Wetenschap wederpartij aangetoond & \texttt{false} & \texttt{false} \\
Vermoeden wetenschap schuldenaar tegenbewijs & \texttt{null} & \texttt{null} \\
Vermoeden wetenschap wederpartij tegenbewijs & \texttt{null} & \texttt{null} \\
Tegen onderwaarde & \texttt{null} & \texttt{null} \\
Voldoening of zekerheid niet opeisbare schuld & \texttt{null} & \texttt{null} \\
Gelieerde partij & \texttt{null} & \texttt{null} \\
Reeds verplicht voor aanvang termijn & \texttt{null} & \texttt{null} \\
\textbf{Vernietigbaar} & \textbf{\texttt{false}} & \textbf{\texttt{false}} \\
\bottomrule
\end{tabular}
\end{center}

\paragraph{Is rechtshandeling - overwegingen}\label{is-rechtshandeling---overwegingen-56}

De verkoop van de BMW door Brand aan Jon West (John West) betreft een koopovereenkomst. Dit is een meerzijdige rechtshandeling gericht op het scheppen van verbintenissen (levering van de auto tegen betaling van de koopprijs).

\paragraph{Verricht door de schuldenaar - overwegingen}\label{verricht-door-de-schuldenaar---overwegingen-56}

De verkoop van de BMW is verricht door Brand Station B.V. (de latere failliet en dus de schuldenaar) zelf.

\paragraph{Om niet - overwegingen}\label{om-niet---overwegingen-56}

De BMW is verkocht voor een overeengekomen prijs van € 10.584,09 inclusief btw. Er staat dus een tegenprestatie tegenover de levering, waardoor de rechtshandeling om baat (anders dan om niet) is verricht.

\paragraph{Onverplicht - overwegingen}\label{onverplicht---overwegingen-56}

De verkoop van de BMW door Brand aan John West (History Division B.V.) in april/mei 2002 is een onverplichte rechtshandeling. Er bestond voor Brand geen reeds op de wet of een eerdere overeenkomst rustende rechtsplicht om deze specifieke koopovereenkomst aan te gaan en de auto over te dragen. Het sluiten van een dergelijke koopovereenkomst geschiedt in beginsel op basis van contractsvrijheid en is derhalve onverplicht in de zin van artikel 42 Faillissementswet.

\paragraph{Benadeling schuldeisers - overwegingen}\label{benadeling-schuldeisers---overwegingen-56}

De curator stelt dat de BMW door Brand Station B.V. onder de reële waarde is verkocht aan John West B.V. (een gelieerde vennootschap waarvan Van Deursen eveneens bestuurder en aandeelhouder is). Indien een actief onder de marktwaarde wordt overgedragen, treedt er een feitelijke vermindering van het actieve vermogen van de schuldenaar op. In de hypothetische situatie waarin deze verkoop niet (of tegen een reële, marktconforme prijs) had plaatsgevonden, zou er meer verhaalsactief voor de gezamenlijke schuldeisers beschikbaar zijn geweest. De feitelijke situatie na de transactie is derhalve ongunstiger voor de schuldeisers, waardoor sprake is van benadeling in de zin van artikel 42 Fw.

\paragraph{Wetenschap - overwegingen}\label{wetenschap---overwegingen-56}

De verkoop van de BMW vond plaats in april of mei 2002, terwijl het faillissement van Brand pas op 19 november 2003 is uitgesproken. Omdat deze rechtshandeling meer dan een jaar vóór de faillietverklaring is verricht, zijn de bewijsvermoedens van art. 43 en 45 Fw niet van toepassing. De curator draagt derhalve de volledige en rechtstreekse bewijslast van de wetenschap van benadeling aan beide zijden op grond van art. 42 Fw. Ten tijde van de verkoop (vlak na de doorstart in het voorjaar van 2002) waren het faillissement en een tekort daarin nog niet met een redelijke mate van waarschijnlijkheid te voorzien. De financiële problemen en de opzegging van het krediet door NMB Heller vonden pas in de loop van 2003 plaats. Wetenschap van benadeling aan de zijde van de schuldenaar en de wederpartij is dan ook niet aangetoond.

\newpage

\markboth{JOR 2010/210}{mistral/mistral-medium-3.5-128b:fp8 (k = 5)}

\subsection{k = 5}\label{k-5-6}

\subsubsection{Mistral}\label{mistral-19}

\texttt{mistral/mistral-medium-3.5-128b:fp8}

\begin{center}
\small
\begin{tabular}{@{}p{7.4cm}p{2.6cm}p{2.6cm}@{}}
\toprule
Variabele & Waarde (model) & Waarde (gold) \\
\midrule
Is rechtshandeling & \texttt{true} & \texttt{true} \\
Datum rechtshandeling & \texttt{2002-04-01} & \texttt{2002-04-01} \\
Verricht door de schuldenaar & \texttt{true} & \texttt{true} \\
Om niet & \texttt{false} & \texttt{false} \\
Datum faillietverklaring & \texttt{2003-11-19} & \texttt{2003-11-19} \\
Onverplicht & \texttt{true} & \texttt{true} \\
Benadeling schuldeisers & \textcolor{gray}{\texttt{true}} & \texttt{false} \\
Wetenschap schuldenaar aangetoond & \texttt{false} & \texttt{false} \\
Wetenschap wederpartij aangetoond & \texttt{false} & \texttt{false} \\
Vermoeden wetenschap schuldenaar tegenbewijs & \texttt{null} & \texttt{null} \\
Vermoeden wetenschap wederpartij tegenbewijs & \texttt{null} & \texttt{null} \\
Tegen onderwaarde & \texttt{null} & \texttt{null} \\
Voldoening of zekerheid niet opeisbare schuld & \texttt{null} & \texttt{null} \\
Gelieerde partij & \textcolor{gray}{\texttt{true}} & \texttt{null} \\
Reeds verplicht voor aanvang termijn & \texttt{null} & \texttt{null} \\
\textbf{Vernietigbaar} & \textbf{\texttt{false}} & \textbf{\texttt{false}} \\
\bottomrule
\end{tabular}
\end{center}

\paragraph{Is rechtshandeling - overwegingen}\label{is-rechtshandeling---overwegingen-57}

De doorverkoop van de BMW door Brand aan John West is een koopovereenkomst, wat een rechtshandeling is (art. 3:33 BW).

\paragraph{Verricht door de schuldenaar - overwegingen}\label{verricht-door-de-schuldenaar---overwegingen-57}

Brand (de schuldenaar) heeft de BMW doorverkocht aan John West.

\paragraph{Om niet - overwegingen}\label{om-niet---overwegingen-57}

De BMW is verkocht voor € 10.584,09 inclusief BTW, wat een tegenprestatie inhoudt. Derhalve is de handeling niet om niet.

\paragraph{Onverplicht - overwegingen}\label{onverplicht---overwegingen-57}

De verkoop van de BMW door Brand aan John West voor € 10.584,09 is een onverplichte rechtshandeling in de zin van artikel 42 lid 1 Fw. Er bestond geen wettelijke of contractuele verplichting voor Brand om de BMW te verkopen. De BMW was een actief dat Brand had overgenomen uit het faillissement van HWH, en er is geen sprake van een voorafgaande verplichting om dit actief te verkopen. Verkoop gevolgd door verrekening is een onverplichte rechtshandeling van de schuldenaar (Brand), zoals ook blijkt uit de aandachtspunten en de jurisprudentie (bijv. ECLI:NL:RBALM:2011:BU7299, r.o. 3.3). De verrekening zelf is een handeling van de schuldeiser (John West) en valt buiten art. 42 Fw, maar de verkoop als rechtshandeling van Brand is wel onverplicht.

\paragraph{Benadeling schuldeisers - overwegingen}\label{benadeling-schuldeisers---overwegingen-57}

De verkoop van de BMW door Brand aan John West voor € 10.584,09 (inclusief BTW) in april/mei 2002 moet worden beoordeeld op benadeling van schuldeisers in de zin van artikel 42 Fw.

Allereerst is de rechtshandeling (de verkoop) \textbf{onverplicht} in de zin van artikel 42 lid 1 Fw, omdat er geen wettelijke of contractuele verplichting bestond voor Brand om de BMW te verkopen. De verkoop was een vrije keuze van Brand, zonder dat deze daartoe verplicht was.

Vervolgens dient te worden beoordeeld of de schuldeisers door deze verkoop zijn \textbf{benadeeld}. Hiertoe wordt een vermogensvergelijking gemaakt tussen de hypothetische situatie zonder de verkoop en de feitelijke situatie met de verkoop.

\begin{itemize}
\tightlist
\item
  \textbf{Hypothetische situatie zonder verkoop}: De BMW zou deel uitmaken van de boedel van Brand en kon door de curator in het faillissement (uitgesproken op 19 november 2003) worden verhaald ten behoeve van de schuldeisers. De waarde van de BMW had dus ten goede kunnen komen aan de gezamenlijke schuldeisers.
\item
  \textbf{Feitelijke situatie met verkoop}: De BMW is verkocht aan John West voor € 10.584,09. De opbrengst hiervan is niet in de boedel van Brand gekomen, maar ten goede gekomen aan John West. Hierdoor is het actief van Brand (en later de boedel) verminderd met de waarde van de BMW.
\end{itemize}

Uit de casus blijkt dat Brand in financiële moeilijkheden verkeerde. NMB Heller had op 5 maart 2003 al een brief gestuurd waarin werd medegedeeld dat de kredietfaciliteit zou worden beëindigd tenzij Brand de overschrijding met ten minste € 500.000,- zou terugbrengen. In oktober 2003 weigerde NMB Heller de financiering verder te continueren, waarna Brand op 19 november 2003 failliet werd verklaard. Dit wijst erop dat Brand op het moment van de verkoop (april/mei 2002) al in financiële problemen verkeerde en dat de verkoop van de BMW een vermindering van het actief betekende, wat de schuldeisers benadeelde.

Daarnaast is van belang dat Van Deursen zowel bestuurder en enig aandeelhouder van Brand als van John West was. Dit betekent dat hij bij de verkoop van de BMW aan John West \textbf{wist of behoorde te weten} dat deze handeling benadeling van de schuldeisers tot gevolg zou hebben (art. 42 lid 2 Fw). De verkoop was immers een transactie tussen partijen die onder dezelfde invloedssfeer vielen, en de opbrengst kwam niet ten goede aan de boedel, maar aan een met Van Deursen verbonden partij.

Tot slot is er sprake van \textbf{verstoring van de rangorde tussen schuldeisers} (ECLI:NL:RBALM:2011:BU7299, r.o. 3.4). De opbrengst van de verkoop is ten goede gekomen aan John West, terwijl andere schuldeisers van Brand onbetaald bleven. Dit is een vorm van benadeling, ook al is de prijs mogelijk marktconform geweest.

Gelet op het bovenstaande zijn de schuldeisers door de verkoop van de BMW aan John West benadeeld in de zin van artikel 42 Fw.

\paragraph{Wetenschap - overwegingen}\label{wetenschap---overwegingen-57}

De verkoop van de BMW door Brand aan John West (History Division B.V.) voldoet aan de voorwaarden van art. 43 lid 1 sub 3° Fw. Brand (schuldenaar) en John West (wederpartij) zijn gelieerde partijen: Van Deursen is bestuurder en enig aandeelhouder van zowel Brand als John West (art. 43 lid 1 sub 3b Fw). De rechtshandeling is verricht binnen een jaar voor de faillietverklaring (april/mei 2002 vs.~faillissement 19 november 2003) en Brand had zich niet reeds voor de aanvang van die termijn daartoe verplicht. Derhalve wordt de wetenschap van benadeling aan beide zijden vermoed (art. 43 lid 1 Fw). Er is geen sprake van tegenbewijs dat dit vermoeden ontzenuwt, gelet op de financiële problemen van Brand (o.a. beëindiging kredietarrangement door NMB Heller in maart 2003 en weigering verdere financiering in oktober 2003).

\newpage

\markboth{JOR 2010/210}{gpt-5.5 (k = 5)}

\subsubsection{GPT}\label{gpt-19}

\texttt{gpt-5.5}

\begin{center}
\small
\begin{tabular}{@{}p{7.4cm}p{2.6cm}p{2.6cm}@{}}
\toprule
Variabele & Waarde (model) & Waarde (gold) \\
\midrule
Is rechtshandeling & \texttt{true} & \texttt{true} \\
Datum rechtshandeling & \texttt{2002-04-01} & \texttt{2002-04-01} \\
Verricht door de schuldenaar & \texttt{true} & \texttt{true} \\
Om niet & \texttt{false} & \texttt{false} \\
Datum faillietverklaring & \texttt{2003-11-19} & \texttt{2003-11-19} \\
Onverplicht & \texttt{true} & \texttt{true} \\
Benadeling schuldeisers & \texttt{false} & \texttt{false} \\
Wetenschap schuldenaar aangetoond & \texttt{false} & \texttt{false} \\
Wetenschap wederpartij aangetoond & \texttt{false} & \texttt{false} \\
Vermoeden wetenschap schuldenaar tegenbewijs & \texttt{null} & \texttt{null} \\
Vermoeden wetenschap wederpartij tegenbewijs & \texttt{null} & \texttt{null} \\
Tegen onderwaarde & \textcolor{gray}{\texttt{true}} & \texttt{null} \\
Voldoening of zekerheid niet opeisbare schuld & \texttt{null} & \texttt{null} \\
Gelieerde partij & \textcolor{gray}{\texttt{true}} & \texttt{null} \\
Reeds verplicht voor aanvang termijn & \texttt{null} & \texttt{null} \\
\textbf{Vernietigbaar} & \textbf{\texttt{false}} & \textbf{\texttt{false}} \\
\bottomrule
\end{tabular}
\end{center}

\paragraph{Is rechtshandeling - overwegingen}\label{is-rechtshandeling---overwegingen-58}

De te beoordelen BMW-handeling betreft de koopovereenkomst/verkoop van de BMW door Brand aan John West. Een koop/verkoop is gericht op rechtsgevolg en kwalificeert daarom als rechtshandeling. De exacte dag is niet vermeld; de feiten noemen april/mei 2002, daarom is de eerste maand uit die periode als datum opgenomen.

\paragraph{Verricht door de schuldenaar - overwegingen}\label{verricht-door-de-schuldenaar---overwegingen-58}

De failliete schuldenaar is Brand Station B.V. Uit de feiten volgt dat Brand de BMW in april/mei 2002 heeft doorverkocht aan John West.

\paragraph{Om niet - overwegingen}\label{om-niet---overwegingen-58}

Brand heeft de BMW verkocht voor een prijs van € 10.584,09 inclusief btw. Er is dus een wederprestatie bedongen; ook al stelt de curator dat de verkoop beneden de reële waarde plaatsvond, volgt uit de feiten niet dat de handeling om niet was.

\paragraph{Onverplicht - overwegingen}\label{onverplicht---overwegingen-58}

De aangevochten rechtshandeling met betrekking tot de BMW is de verkoop/koopovereenkomst waarbij Brand de BMW in april/mei 2002 aan John West heeft verkocht. Uit de casus blijkt niet van een wettelijke verplichting of een reeds bestaande, in rechte afdwingbare contractuele verplichting van Brand om deze BMW aan John West te verkopen. Dat Brand de BMW kort daarvoor uit de doorstart had verkregen, dat partijen gelieerd waren via Van Deursen, of dat er mogelijk praktische, economische of concernmatige redenen waren om de auto over te dragen, schept geen rechtsplicht. Ook is niet gesteld dat John West vóór de verkoop een afdwingbaar recht op levering van de BMW had. Vergelijkbaar met de lijn in ECLI:NL:RBALM:2011:BU7299 en ECLI:NL:GHDHA:2016:3475 geldt dat een verkoop van activa, ook indien daarmee een bepaalde partij wordt bevoordeeld of een schuldpositie wordt beïnvloed, onverplicht is indien geen vooraf bestaande verplichting tot verkoop bestond. De verkoop van de BMW moet daarom als onverplicht in de zin van artikel 42 Fw worden aangemerkt.

\paragraph{Benadeling schuldeisers - overwegingen}\label{benadeling-schuldeisers---overwegingen-58}

Op basis van de vaststaande feiten kan benadeling door de verkoop van de BMW niet worden vastgesteld. Vast staat slechts dat Brand de BMW in april/mei 2002 aan John West heeft verkocht voor € 10.584,09 inclusief btw. De stelling dat dit beneden de reële waarde was, is alleen als standpunt van de curator weergegeven en wordt niet onderbouwd met een taxatie, marktwaarde, boekwaarde of andere concrete gegevens. Ook blijkt niet dat de koopprijs niet is betaald, dat de opbrengst aan één specifieke schuldeiser is onttrokken of dat de BMW zonder deze verkoop ten tijde van het faillissement nog voor de gezamenlijke schuldeisers beschikbaar zou zijn geweest. Bij de vereiste vermogensvergelijking kan daarom niet worden geconcludeerd dat het actief is verminderd of de verhaalspositie/rangorde van schuldeisers is verslechterd. Als wél zou komen vast te staan dat de BMW materieel meer waard was dan de ontvangen prijs, zou benadeling voor het verschil voor de hand liggen; dat volgt echter niet uit de gegeven casus.

\paragraph{Wetenschap - overwegingen}\label{wetenschap---overwegingen-58}

De verkoop van de BMW door Brand aan John West is een rechtshandeling anders dan om niet, zodat wetenschap van benadeling bij zowel Brand als John West vereist is. Wel is sprake van een gelieerde partij: Van Deursen was bestuurder/enig aandeelhouder van zowel Brand als John West. Ook ligt, uitgaande van de stelling dat de BMW beneden reële waarde is verkocht, aanmerkelijke onderwaarde voor de hand. Deze omstandigheden kunnen in beginsel relevant zijn voor art. 43 Fw, maar uit de feiten volgt niet rechtstreeks dat ten tijde van de verkoop in april/mei 2002 het faillissement van Brand en een tekort daarin met een redelijke mate van waarschijnlijkheid voorzienbaar waren. De latere financieringsproblemen bij NMB Heller blijken pas uit 2003 en zijn onvoldoende om wetenschap bij de verkoop in 2002 aan te tonen. Art. 45 Fw is niet aan de orde, omdat geen sprake is van een rechtshandeling om niet. Er is daarom geen directe wetenschap van benadeling aangetoond bij Brand of John West.

\newpage

\markboth{JOR 2010/210}{gemini-3.5-flash (k = 5)}

\subsubsection{Gemini}\label{gemini-19}

\texttt{gemini-3.5-flash}

\begin{center}
\small
\begin{tabular}{@{}p{7.4cm}p{2.6cm}p{2.6cm}@{}}
\toprule
Variabele & Waarde (model) & Waarde (gold) \\
\midrule
Is rechtshandeling & \texttt{true} & \texttt{true} \\
Datum rechtshandeling & \texttt{2002-04-01} & \texttt{2002-04-01} \\
Verricht door de schuldenaar & \texttt{true} & \texttt{true} \\
Om niet & \texttt{false} & \texttt{false} \\
Datum faillietverklaring & \texttt{2003-11-19} & \texttt{2003-11-19} \\
Onverplicht & \texttt{true} & \texttt{true} \\
Benadeling schuldeisers & \textcolor{gray}{\texttt{true}} & \texttt{false} \\
Wetenschap schuldenaar aangetoond & \texttt{false} & \texttt{false} \\
Wetenschap wederpartij aangetoond & \texttt{false} & \texttt{false} \\
Vermoeden wetenschap schuldenaar tegenbewijs & \texttt{null} & \texttt{null} \\
Vermoeden wetenschap wederpartij tegenbewijs & \texttt{null} & \texttt{null} \\
Tegen onderwaarde & \texttt{null} & \texttt{null} \\
Voldoening of zekerheid niet opeisbare schuld & \texttt{null} & \texttt{null} \\
Gelieerde partij & \texttt{null} & \texttt{null} \\
Reeds verplicht voor aanvang termijn & \texttt{null} & \texttt{null} \\
\textbf{Vernietigbaar} & \textbf{\texttt{false}} & \textbf{\texttt{false}} \\
\bottomrule
\end{tabular}
\end{center}

\paragraph{Is rechtshandeling - overwegingen}\label{is-rechtshandeling---overwegingen-59}

De verkoop van de BMW door Brand aan Jon West betreft een koopovereenkomst. Dit is een meerzijdige rechtshandeling gericht op het scheppen van verbintenissen, te weten de levering van de auto tegen betaling van de koopprijs.

\paragraph{Verricht door de schuldenaar - overwegingen}\label{verricht-door-de-schuldenaar---overwegingen-59}

De koopovereenkomst is aangegaan en verricht door Brand Station B.V. zelf, de schuldenaar die naderhand failliet is verklaard.

\paragraph{Om niet - overwegingen}\label{om-niet---overwegingen-59}

De BMW is verkocht voor een overeengekomen prijs van € 10.584,09 inclusief btw. Omdat er een tegenprestatie is bedongen, is de rechtshandeling anders dan om niet (om baat).

\paragraph{Onverplicht - overwegingen}\label{onverplicht---overwegingen-59}

De aangevochten rechtshandeling betreft de verkoop en levering van de BMW (kenteken XX-PG-65) door Brand Station B.V. aan John West B.V. (destijds History Division B.V.) in april/mei 2002. Een rechtshandeling is onverplicht in de zin van artikel 42 Fw wanneer er geen op de wet of een overeenkomst berustende rechtsplicht bestond om deze te verrichten. Uit de casus blijkt niet dat er voor Brand Station B.V. enige wettelijke verplichting of een voorafgaande contractuele verbintenis bestond om de BMW aan John West B.V. te verkopen of over te dragen. De verkoop en levering van de auto zijn dan ook onverplicht verricht.

\paragraph{Benadeling schuldeisers - overwegingen}\label{benadeling-schuldeisers---overwegingen-59}

In deze casus heeft Brand Station B.V. een BMW verkocht en geleverd aan John West B.V., een gelieerde vennootschap waarvan Van Deursen eveneens de enig bestuurder en aandeelhouder is. De curator stelt dat deze verkoop heeft plaatsgevonden voor een bedrag dat beneden de reële marktwaarde van de auto lag. Door een actiefbestanddeel onder de reële waarde over te dragen, is het vermogen van de schuldenaar verminderd. Indien deze rechtshandeling niet had plaatsgevonden (de hypothetische situatie), zou de BMW nog in het vermogen van Brand Station B.V. aanwezig zijn geweest en in de faillissementsboedel zijn gevallen, dan wel voor een hogere, marktconforme prijs zijn verkocht. Hierdoor is er sprake van een vermindering van het actief, wat leidt tot een feitelijke verslechtering van de verhaalsmogelijkheden van de gezamenlijke schuldeisers ten tijde van het geldend maken van de rechten door de curator. De schuldeisers zijn derhalve benadeeld in de zin van artikel 42 Fw.

\paragraph{Wetenschap - overwegingen}\label{wetenschap---overwegingen-59}

De verkoop van de BMW vond plaats in april/mei 2002, terwijl het faillissement van Brand op 19 november 2003 is uitgesproken. Omdat de transactie meer dan een jaar vóór de faillietverklaring heeft plaatsgevonden, zijn de bewijsvermoedens van art. 43 Fw en art. 45 Fw niet van toepassing. De curator draagt de bewijslast om de wetenschap van benadeling rechtstreeks aan te tonen op grond van art. 42 Fw. Hoewel Van Deursen bestuurder en aandeelhouder was van zowel de schuldenaar (Brand) als de wederpartij (John West), is niet aangetoond dat ten tijde van de transactie (april/mei 2002) het faillissement en een tekort daarin met een redelijke mate van waarschijnlijkheid te voorzien waren. Brand was destijds net gestart met een doorstart en de financiering werd pas in 2003 beëindigd. De wetenschap van benadeling is derhalve niet aangetoond.

\newpage

\markboth{JOR 2010/210}{mistral/mistral-medium-3.5-128b:fp8 (k = 10)}

\subsection{k = 10}\label{k-10-6}

\subsubsection{Mistral}\label{mistral-20}

\texttt{mistral/mistral-medium-3.5-128b:fp8}

\begin{center}
\small
\begin{tabular}{@{}p{7.4cm}p{2.6cm}p{2.6cm}@{}}
\toprule
Variabele & Waarde (model) & Waarde (gold) \\
\midrule
Is rechtshandeling & \texttt{true} & \texttt{true} \\
Datum rechtshandeling & \texttt{2002-04-01} & \texttt{2002-04-01} \\
Verricht door de schuldenaar & \texttt{true} & \texttt{true} \\
Om niet & \texttt{false} & \texttt{false} \\
Datum faillietverklaring & \texttt{2003-11-19} & \texttt{2003-11-19} \\
Onverplicht & \texttt{true} & \texttt{true} \\
Benadeling schuldeisers & \textcolor{gray}{\texttt{true}} & \texttt{false} \\
Wetenschap schuldenaar aangetoond & \texttt{false} & \texttt{false} \\
Wetenschap wederpartij aangetoond & \texttt{false} & \texttt{false} \\
Vermoeden wetenschap schuldenaar tegenbewijs & \texttt{null} & \texttt{null} \\
Vermoeden wetenschap wederpartij tegenbewijs & \texttt{null} & \texttt{null} \\
Tegen onderwaarde & \texttt{null} & \texttt{null} \\
Voldoening of zekerheid niet opeisbare schuld & \texttt{null} & \texttt{null} \\
Gelieerde partij & \textcolor{gray}{\texttt{true}} & \texttt{null} \\
Reeds verplicht voor aanvang termijn & \texttt{null} & \texttt{null} \\
\textbf{Vernietigbaar} & \textbf{\texttt{false}} & \textbf{\texttt{false}} \\
\bottomrule
\end{tabular}
\end{center}

\paragraph{Is rechtshandeling - overwegingen}\label{is-rechtshandeling---overwegingen-60}

De doorverkoop van de BMW door Brand aan John West is een koopovereenkomst, wat een rechtshandeling is (art. 3:33 BW).

\paragraph{Verricht door de schuldenaar - overwegingen}\label{verricht-door-de-schuldenaar---overwegingen-60}

Brand (de schuldenaar) heeft de BMW doorverkocht aan John West.

\paragraph{Om niet - overwegingen}\label{om-niet---overwegingen-60}

De BMW is verkocht voor € 10.584,09 inclusief BTW, wat een tegenprestatie inhoudt. Derhalve is de handeling niet om niet.

\paragraph{Onverplicht - overwegingen}\label{onverplicht---overwegingen-60}

De verkoop van de BMW door Brand aan Jon West is een onverplichte rechtshandeling in de zin van artikel 42 lid 1 Fw. Er bestond geen op de wet of een overeenkomst berustende rechtsplicht voor Brand om de BMW te verkopen. De verkoop was geen verplichting die voortvloeide uit een bestaande overeenkomst of wettelijke verplichting. Morele, feitelijke of economische druk (zoals de doorstart van HWH) schept geen rechtsplicht. De stelplicht en bewijslast van het onverplichte karakter rusten op de curator, en in deze casus is voldaan aan die eis omdat er geen sprake was van een verplichting om de BMW te verkopen.

\paragraph{Benadeling schuldeisers - overwegingen}\label{benadeling-schuldeisers---overwegingen-60}

De curator stelt dat de BMW beneden de reële waarde is verkocht aan John West, waardoor de schuldeisers van Brand zijn benadeeld. Uit de casus blijkt dat Brand de BMW in april/mei 2002 heeft doorverkocht aan John West voor een prijs van € 10.584,09 inclusief BTW. Van Deursen was zowel bestuurder van Brand als van John West, wat een belangenverstrengeling suggereert. Bovendien was Brand op dat moment in financiële moeilijkheden, zoals blijkt uit de latere faillissementsaanvraag en de brief van NMB Heller van 5 maart 2003, waarin wordt meegedeeld dat de kredietfaciliteit wordt beëindigd en dat Brand de faciliteit met ten minste € 500.000,- moet terugbrengen. Dit duidt op liquiditeitsproblemen bij Brand.

De benadeling van de schuldeisers moet worden beoordeeld door een vermogensvergelijking: de hypothetische situatie waarin de schuldeisers zouden hebben verkeerd zonder de gewraakte rechtshandeling (verkoop BMW) wordt vergeleken met de situatie waarin zij feitelijk verkeren als die handeling in stand blijft. In dit geval is de BMW, een actief van Brand, uit de boedel verdwenen zonder dat daar een gelijkwaardige tegenprestatie tegenover staat die ten goede komt aan alle schuldeisers. De opbrengst van € 10.584,09 is naar John West gegaan, een vennootschap die onder controle stond van Van Deursen, wat suggereert dat de opbrengst niet ten goede is gekomen aan de gezamenlijke schuldeisers van Brand.

Daarnaast is er sprake van een verstoring van de rangorde tussen de schuldeisers, omdat de opbrengst van de verkoop van de BMW alleen ten goede is gekomen aan John West, en niet aan alle schuldeisers van Brand. Dit is voldoende om aan te nemen dat de schuldeisers in het algemeen zijn benadeeld, ook als per saldo geen vermogensvermindering optreedt (ECLI:NL:RBALM:2011:BU7299, r.o. 3.4).

Gelet op het bovenstaande is er sprake van benadeling van de schuldeisers in de zin van artikel 42 Fw.

\paragraph{Wetenschap - overwegingen}\label{wetenschap---overwegingen-60}

De verkoop van de BMW door Brand aan John West (History Division B.V.) is een rechtshandeling anders dan om niet, verricht binnen een jaar voor de faillietverklaring van Brand (19 november 2003). Brand en John West zijn gelieerde partijen in de zin van art. 43 lid 1 sub 3b Fw, omdat Van Deursen bestuurder en enig aandeelhouder is van zowel Brand als John West. Hierdoor wordt de wetenschap van benadeling aan beide zijden vermoed (art. 43 lid 1 Fw). Er is geen sprake van tegenbewijs dat dit vermoeden ontzenuwt, noch van een verplichting die voor de jaartermijn bestond. De curator hoeft derhalve de wetenschap niet rechtstreeks aan te tonen.

\newpage

\markboth{JOR 2010/210}{gpt-5.5 (k = 10)}

\subsubsection{GPT}\label{gpt-20}

\texttt{gpt-5.5}

\begin{center}
\small
\begin{tabular}{@{}p{7.4cm}p{2.6cm}p{2.6cm}@{}}
\toprule
Variabele & Waarde (model) & Waarde (gold) \\
\midrule
Is rechtshandeling & \texttt{true} & \texttt{true} \\
Datum rechtshandeling & \texttt{2002-04-01} & \texttt{2002-04-01} \\
Verricht door de schuldenaar & \texttt{true} & \texttt{true} \\
Om niet & \texttt{false} & \texttt{false} \\
Datum faillietverklaring & \texttt{2003-11-19} & \texttt{2003-11-19} \\
Onverplicht & \texttt{true} & \texttt{true} \\
Benadeling schuldeisers & \texttt{false} & \texttt{false} \\
Wetenschap schuldenaar aangetoond & \texttt{false} & \texttt{false} \\
Wetenschap wederpartij aangetoond & \texttt{false} & \texttt{false} \\
Vermoeden wetenschap schuldenaar tegenbewijs & \texttt{null} & \texttt{null} \\
Vermoeden wetenschap wederpartij tegenbewijs & \texttt{null} & \texttt{null} \\
Tegen onderwaarde & \texttt{null} & \texttt{null} \\
Voldoening of zekerheid niet opeisbare schuld & \texttt{null} & \texttt{null} \\
Gelieerde partij & \textcolor{gray}{\texttt{true}} & \texttt{null} \\
Reeds verplicht voor aanvang termijn & \texttt{null} & \texttt{null} \\
\textbf{Vernietigbaar} & \textbf{\texttt{false}} & \textbf{\texttt{false}} \\
\bottomrule
\end{tabular}
\end{center}

\paragraph{Is rechtshandeling - overwegingen}\label{is-rechtshandeling---overwegingen-61}

De gewraakte handeling betreft de verkoop/koopovereenkomst van de BMW door Brand aan John West. Een koopovereenkomst is gericht op rechtsgevolg en kwalificeert als rechtshandeling. De exacte datum is niet vermeld; uit de feiten volgt alleen dat dit in april/mei 2002 plaatsvond.

\paragraph{Verricht door de schuldenaar - overwegingen}\label{verricht-door-de-schuldenaar---overwegingen-61}

De verkoop is verricht door Brand Station B.V., de latere failliet/schuldenaar zelf, aan John West.

\paragraph{Om niet - overwegingen}\label{om-niet---overwegingen-61}

Brand verkocht de BMW voor een prijs van € 10.584,09 inclusief btw. Daarmee is een wederprestatie bedongen, zodat de rechtshandeling anders dan om niet is, ook al stelt de curator dat de prijs beneden de reële waarde lag.

\paragraph{Onverplicht - overwegingen}\label{onverplicht---overwegingen-61}

De aangevochten rechtshandeling betreft de verkoop/doorverkoop door Brand van de BMW aan John West in april/mei 2002. Uit de casus blijkt niet van een wettelijke verplichting of een reeds bestaande overeenkomst op grond waarvan Brand gehouden was deze BMW juist aan John West te verkopen of tegen de gehanteerde prijs over te dragen. Dat de verkoop mogelijk binnen concernverband plaatsvond, praktisch wenselijk was of verband hield met de doorstart, schept geen in rechte afdwingbare rechtsplicht. Beslissend is of John West de verkoop en levering in rechte had kunnen afdwingen vóórdat Brand daartoe besloot; daarvoor ontbreken aanwijzingen. Vergelijkbaar met de lijn in onder meer ECLI:NL:RBALM:2011:BU7299 geldt dat verkoop van activa onverplicht is, tenzij reeds een verplichting tot verkoop bestond. Die verplichting is hier niet gesteld of gebleken. De rechtshandeling omtrent de BMW is daarom onverplicht verricht in de zin van artikel 42 Fw.

\paragraph{Benadeling schuldeisers - overwegingen}\label{benadeling-schuldeisers---overwegingen-61}

Voor benadeling moet de feitelijke situatie met verkoop van de BMW aan John West worden vergeleken met de hypothetische situatie zonder die verkoop. Vast staat slechts dat Brand de BMW in april/mei 2002 aan John West heeft verkocht voor € 10.584,09 inclusief btw. Uit de casus volgt niet concreet wat de reële marktwaarde van de BMW op dat moment was, of wat de boedel in het hypothetische scenario zonder verkoop daadwerkelijk voor de BMW zou hebben kunnen realiseren. De enkele stelling van de curator dat de BMW beneden de reële waarde is verkocht, is daarvoor onvoldoende. Evenmin blijkt dat de koopprijs niet is betaald of dat de opbrengst aan één specifieke schuldeiser is toegekomen op een wijze die de rangorde heeft verstoord. Bij gebrek aan vaststaande feiten waaruit een vermindering van het actief, vermeerdering van passiva of rangordeverstoring volgt, kan benadeling van de schuldeisers door de BMW-transactie niet worden vastgesteld.

\paragraph{Wetenschap - overwegingen}\label{wetenschap---overwegingen-61}

De verkoop van de BMW door Brand aan John West was een meerzijdige rechtshandeling anders dan om niet, zodat wetenschap van benadeling zowel bij Brand als bij John West vereist is. Hoewel sprake is van gelieerdheid: Van Deursen was bestuurder/enig aandeelhouder van Brand en tevens bestuurder/aandeelhouder van John West, volgt daaruit buiten een toepasselijk bewijsvermoeden nog niet rechtstreeks dat ten tijde van de BMW-transactie het faillissement van Brand en een tekort daarin met redelijke mate van waarschijnlijkheid voorzienbaar waren. De BMW is in april/mei 2002 verkocht, terwijl Brand pas op 19 november 2003 failliet is verklaard; de relevante financieringsproblemen met NMB Heller blijken vooral uit 2003 en de beëindiging/weigering van verdere financiering in oktober 2003. Concrete feiten dat Brand en John West al bij de BMW-verkoop een faillissement met tekort redelijkerwijs moesten voorzien, ontbreken. Art. 43 Fw biedt voor de BMW-transactie daarom geen werkzaam vermoeden, ook al is sprake van gelieerdheid; art. 45 Fw is niet aan de orde omdat de verkoop niet om niet was. De door de curator gestelde verkoop beneden reële waarde kan zonder toepasselijk vermoeden of nadere feiten de wetenschap van benadeling niet dragen.

\newpage

\markboth{JOR 2010/210}{gemini-3.5-flash (k = 10)}

\subsubsection{Gemini}\label{gemini-20}

\texttt{gemini-3.5-flash}

\begin{center}
\small
\begin{tabular}{@{}p{7.4cm}p{2.6cm}p{2.6cm}@{}}
\toprule
Variabele & Waarde (model) & Waarde (gold) \\
\midrule
Is rechtshandeling & \texttt{true} & \texttt{true} \\
Datum rechtshandeling & \texttt{2002-04-01} & \texttt{2002-04-01} \\
Verricht door de schuldenaar & \texttt{true} & \texttt{true} \\
Om niet & \texttt{false} & \texttt{false} \\
Datum faillietverklaring & \texttt{2003-11-19} & \texttt{2003-11-19} \\
Onverplicht & \texttt{true} & \texttt{true} \\
Benadeling schuldeisers & \texttt{false} & \texttt{false} \\
Wetenschap schuldenaar aangetoond & \texttt{false} & \texttt{false} \\
Wetenschap wederpartij aangetoond & \texttt{false} & \texttt{false} \\
Vermoeden wetenschap schuldenaar tegenbewijs & \texttt{null} & \texttt{null} \\
Vermoeden wetenschap wederpartij tegenbewijs & \texttt{null} & \texttt{null} \\
Tegen onderwaarde & \texttt{null} & \texttt{null} \\
Voldoening of zekerheid niet opeisbare schuld & \texttt{null} & \texttt{null} \\
Gelieerde partij & \texttt{null} & \texttt{null} \\
Reeds verplicht voor aanvang termijn & \texttt{null} & \texttt{null} \\
\textbf{Vernietigbaar} & \textbf{\texttt{false}} & \textbf{\texttt{false}} \\
\bottomrule
\end{tabular}
\end{center}

\paragraph{Is rechtshandeling - overwegingen}\label{is-rechtshandeling---overwegingen-62}

De verkoop van de BMW door Brand aan Jon West (John West) betreft een koopovereenkomst. Dit is een meerzijdige rechtshandeling gericht op het scheppen van verbintenissen (levering van de auto tegen betaling van de koopprijs).

\paragraph{Verricht door de schuldenaar - overwegingen}\label{verricht-door-de-schuldenaar---overwegingen-62}

De koopovereenkomst is gesloten door Brand Station B.V. (de latere failliet) zelf als verkopende partij.

\paragraph{Om niet - overwegingen}\label{om-niet---overwegingen-62}

De verkoop vond plaats tegen een overeengekomen koopprijs van € 10.584,09 inclusief btw. Er staat dus een tegenprestatie tegenover de levering van de auto, waardoor de rechtshandeling om baat (anders dan om niet) is.

\paragraph{Onverplicht - overwegingen}\label{onverplicht---overwegingen-62}

De aangevochten rechtshandeling betreft de verkoop en levering van de BMW door Brand Station B.V. aan John West (History Division B.V.) in april/mei 2002. Een rechtshandeling is onverplicht in de zin van artikel 42 Fw wanneer daartoe op het moment van het verrichten geen op de wet of een eerdere overeenkomst berustende rechtsplicht bestond die in rechte afdwingbaar was. Uit de casus blijkt niet dat er voor Brand Station B.V. enige wettelijke of contractuele verplichting bestond om de BMW aan John West te verkopen en over te dragen. De koopovereenkomst en de daaruit voortvloeiende levering zijn dan ook onverplicht verricht.

\paragraph{Benadeling schuldeisers - overwegingen}\label{benadeling-schuldeisers---overwegingen-62}

Voor een geslaagd beroep op de faillissementspauliana (artikel 42 Fw) is vereist dat de schuldeisers door de aangevochten rechtshandeling zijn benadeeld. Of er sprake is van benadeling wordt vastgesteld door een vermogensvergelijking: de hypothetische situatie waarin de schuldeisers zonder de rechtshandeling zouden hebben verkeerd, wordt vergeleken met de feitelijke situatie. In deze casus stelt de curator dat de BMW beneden de reële waarde is verkocht aan John West B.V. De BMW is echter verkocht voor een prijs van € 10.584,09 inclusief btw, kort nadat Brand deze auto zelf had overgenomen uit het faillissement van HWH. Indien een goed voor een reële, marktconforme prijs wordt verkocht en de opbrengst in de vennootschap vloeit, worden schuldeisers in beginsel niet benadeeld omdat het actief (de auto) is vervangen door een gelijkwaardig actief (de koopsom). Nu de curator, tegenover de gemotiveerde betwisting, onvoldoende heeft aangetoond dat de reële marktwaarde van de BMW ten tijde van de verkoop wezenlijk hoger lag dan de betaalde koopsom, is er geen sprake van een vermindering van het actief en is de vereiste benadeling van schuldeisers niet komen vast te staan.

\paragraph{Wetenschap - overwegingen}\label{wetenschap---overwegingen-62}

De verkoop van de BMW vond plaats in april/mei 2002, terwijl het faillissement van Brand Station B.V. pas op 19 november 2003 is uitgesproken. Omdat deze rechtshandeling meer dan een jaar vóór de faillietverklaring is verricht, is het bewijsvermoeden van art. 43 Fw wegens de termijnoverschrijding niet van toepassing. De curator draagt derhalve de bewijslast om de wetenschap van benadeling rechtstreeks aan te tonen op grond van art. 42 Fw. Ten tijde van de verkoop (april/mei 2002) was de vennootschap net gestart met een doorstart. De kredietrelatie met de bank werd pas in maart 2003 opgezegd en de financiering werd pas in oktober 2003 definitief gestaakt. Er zijn geen feiten of omstandigheden gesteld waaruit blijkt dat het faillissement en een tekort daarin ten tijde van de verkoop in april/mei 2002 al met een redelijke mate van waarschijnlijkheid te voorzien waren. De wetenschap van benadeling is dan ook niet aangetoond.

\newpage

\markboth{JOR 2011/20}{Gold standard}

\section{JOR 2011/20}\label{jor-201120}

uitspraakdatum 17-06-2010

\subsection{Gold}\label{gold-7}

\begin{center}
\small
\begin{tabular}{@{}p{8cm}p{2.6cm}@{}}
\toprule
Variabele & Waarde \\
\midrule
Is rechtshandeling & \texttt{true} \\
Datum rechtshandeling & \texttt{2009-08-19} \\
Verricht door de schuldenaar & \texttt{true} \\
Om niet & \texttt{false} \\
Datum faillietverklaring & \texttt{2009-11-30} \\
Onverplicht & \texttt{true} \\
Benadeling schuldeisers & \texttt{true} \\
Wetenschap schuldenaar aangetoond & \texttt{false} \\
Wetenschap wederpartij aangetoond & \texttt{false} \\
Vermoeden wetenschap schuldenaar tegenbewijs & \texttt{false} \\
Vermoeden wetenschap wederpartij tegenbewijs & \texttt{false} \\
Tegen onderwaarde & \texttt{true} \\
Voldoening of zekerheid niet opeisbare schuld & \texttt{null} \\
Gelieerde partij & \texttt{null} \\
Reeds verplicht voor aanvang termijn & \texttt{false} \\
\textbf{Vernietigbaar} & \textbf{\texttt{true}} \\
\bottomrule
\end{tabular}
\end{center}

\subsubsection{Uitspraak}\label{uitspraak-7}

(\ldots; red.)

\paragraph{2. De feiten}\label{de-feiten-2}

2.1. Bij vonnis van deze rechtbank van 6 juli 2009 is de besloten vennootschap Drent Goebel B.V. (hierna ook: Drent) in staat van faillissement verklaard en zijn mr. Brandsma en mr. König als curatoren in dit faillissement aangesteld.

Drent hield zich onder meer bezig met het ontwerpen en vervaardigen van machines, in het bijzonder voor de grafische industrie. Zij beschouwde zich als pionier op het gebied van de ontwikkeling van drukmachines in smallere formaten en heeft in de loop van de tijd verschillende nieuwe types ontwikkeld, waaronder de Vision. Daarnaast verrichtte zij garantiewerkzaamheden en servicewerkzaamheden en leverde zij reserveonderdelen.

Door Drent te verrichten service- en garantiewerkzaamheden besteedde Drent met enige regelmaat uit aan (een van) gedaagden.

2.2. Klanten die bij Drent een nieuwe drukmachine kochten, ruilden veelal hun oude machine in bij Drent. Drent verkocht deze ingeruilde machines vervolgens aan haar zustervennootschap Trend Graphic Services B.V., hierna ook: Trend. Trend verkocht deze machines weer door na ze te hebben opgeknapt of aangepast aan de wensen van de klant. Behalve van Drent betrok Trend ook drukmachines rechtstreeks van de markt. In afwachting van doorverkoop werden de machines door Trend opgeslagen.

Trend had geen personeel in dienst. Zij maakte gebruik van werknemers van Drent en huurde incidenteel diensten van (een van) gedaagden in.

2.3. Gedaagden houden zich onder meer bezig met het wereldwijd aan ondernemingen leveren van integrale technische dienstverlening voor industriële apparatuur. Zij zijn vooral actief binnen de grafische- en verpakkingsindustrie. Gedaagden vervaardigen geen machines, maar onderhouden, reviseren en modificeren machines.

2.4. Op 19 augustus 2009 heeft Trend de navolgende orderbevestiging gestuurd aan ``Veldkamp, Kaagstraat 5b, 8102 GZ Raalte'' (productie 1 van de curatoren):

``(\ldots) Geachte heer Veldkamp,

Na een lange discussie, en een aantal bezoeken aan Veldkamp BV in Raalte, zijn we het eens geworden over de verkoop van de volgende machines en machine delen aan Veldkamp BV.

A. Drent Vision US 512, incompleet van bouwjaar 2000, `as it is' en zonder garantie.

Afwikkel unit, 4 druktorens 22'' formaat groepen, outfeed en een wikkel unit.

B. Drent Vision U 995, incompleet van bouwjaar 1998, `as it is' en zonder garantie.

Afwikkel unit, 5 druktorens 22,5'' formaat groepen, flexo unit, UV installatie, luchtdroger en een wikkel unit

C. Drent Vision U 1063, incompleet van bouwjaar 2000, `as it is' en zonder garantie.

Afwikkel unit, 3 druktorens 23,5'' formaat groepen, flexo unit, UV installatie, luchtdroger en een wikkel unit.

D. 4 x Drent Vision druktorens incompleet aan losse onderdelen op pallets. Bouwjaar 1998.

De totale prijs voor de hierboven vermelde machine units en machine delen is euro 45.000,00, af fabriek.

(\ldots)''

2.5. Trend heeft op 19 augustus 2009 een factuur met nummer D 09 040 gestuurd aan ``VELDKAMP Kaagstraat 5b 8102 GZ Raalte (\ldots)'' (productie 2 van de curatoren). Daarin staat voor de Drent Vision 995 een koopprijs vermeld van € 14.800,--, voor de Drent Vision U 1063 een koopprijs van € 14.500,-- en voor de Drent Vision US 512 een koopprijs van € 11.500,--. De factuur vermeldt voorts een totaalbedrag van € 4.200,-- voor vier exemplaren van een ``Druktoren aan onderdelen incompleet''. Alle koopprijzen zijn exclusief BTW.

2.6. Trend heeft op enig ogenblik aan (een van) gedaagden een stuk toegestuurd met als opschrift (productie 3 van de curatoren): ``Aanvulling/omschrijving Factuur D 09-040 Veldkamp''

In dit stuk staan de drie onder 2.4 vermelde Drent Visionmachines vermeld met daarbij per machine een opsomming van de bij die machine behorende onderdelen, zoals afwikkel infeed, 5 x druktoren, 1 x flexo, etc. Daarnaast vermeldt dit stuk als losse componenten:

``6 x Drent Vision Druktorens U 1066 + 2 units 24''

5 x Drent vision drukunits 600 mm. U 1113

1 x Drent vision basis stapellaar U 997.

1 x Drent Vision pakket uitleg U 997

1 x Drent Vision diep stapellaar U 1039

1 x Flexo unit U 1071

1 x 36'' koelcilinder

Diverse kisten, pallets met kabels, losse frames, onderdelen welke bij de hierboven vermelde machines horen. (\ldots)''

2.7. Bij vonnis van deze rechtbank van 30 november 2009 is Trend in staat van faillissement verklaard en is mr. Brandsma als curator in dit faillissement aangesteld. Bij vonnis van 7 december 2009 is mr. König naast mr. Brandsma aangesteld als curator in dit faillissement.

2.8. Bij brief van 4 maart 2010 aan gedaagden hebben de curatoren meegedeeld dat zij zich op het standpunt stellen dat de schuldeisers van Trend door de verkoop van activa van Trend aan (een van) gedaagden op 19 augustus 2009 zijn benadeeld en dat zij de aan die transactie ten grondslag liggende overeenkomst vernietigen. De vernietiging van de verkoop en levering van de activa en de losse componenten heeft tot gevolg dat Trend hiervan altijd eigenaar is gebleven en dat de machines en de componenten aan de curatoren teruggegeven moeten worden, zo schrijven de curatoren aan gedaagden.

2.9. De curatoren hebben op 5 maart 2010 ten laste van gedaagden conservatoir beslag laten leggen op een drukmachine, onderdelen en componenten die zich bevonden in een werkplaats aan de Kaagstraat 5B en in een loods aan de Punterstraat 6 te Raalte (productie 6 van de curatoren).

Bij exploot van 29 maart 2010 hebben de curatoren conservatoir beslag laten leggen op onder meer druktorens, drukmachines, onderdelen en componenten die waren opgeslagen in een door Veldkamp Exploitatiemaatschappij gehuurde loods aan de Tjalkstraat 9-11 te Raalte (productie 7 van de curatoren). De volgende dag hebben de curatoren conservatoir beslag laten leggen op onder meer drukmachines, frames en kasten en kisten met onderdelen en componenten die zich bevonden in een loods aan de Boeierstraat 9 te Raalte (productie 8 van de curatoren).

2.10. Op 23 mei 2010 hebben de curatoren een bod van € 60.000,-- ontvangen van de besloten vennootschap met beperkte aansprakelijkheid DG Press Services B.V. (hierna ook: DG Press) voor de Drent Vision US 512 met een aantal componenten. Het bod is gedaan onder de conditie ``as is where is'', hetgeen betekent dat de machine wordt gekocht in de staat waarin zij zich bevindt en dat het transport van het gekochte voor rekening van de koper komt en onder onder meer de voorwaarde dat de levering uiterlijk 4 juni 2010 diende plaats te vinden (productie 9 van de curatoren).

2.11. Desgevraagd hebben gedaagden aan de curatoren meegedeeld dat zij niet bereid zijn mee te werken aan afgifte van de van Trend gekochte activa aan de curatoren. Nadat de curatoren DG Press hebben meegedeeld dat de beoogde leveringsdatum van 4 juni 2010 niet gehaald zou kunnen worden, heeft DG Press op 27 mei 2010 bericht bereid te zijn de levertermijn te verlengen, met dien verstande dat de levering uiterlijk op 30 juni 2010 voor 17.00 uur dient plaats te vinden (productie 10 van de curatoren).

De curatoren hebben het bod van DG Press geaccepteerd onder de opschortende voorwaarden dat de rechter-commissaris in het faillissement van Trend toestemming voor de verkoop geeft en zij in staat zijn om het gekochte op 30 juni 2010 voor 17.00 uur aan DG Press te leveren.

2.12. Op 2 juni 2010 heeft DG Press onder de conditie ``as is where is'' een bod van € 50.000,-- uitgebracht op alle componenten van de Drent Vision U 1063, waaronder drie druktorens. DG Press heeft voorts een bod uitgebracht op zes druktorens van de Drent Vision U 1066 voor een bedrag van € 7.500,-- per toren. Beide biedingen zijn uitgebracht onder voorbehoud van een inspectie op compleetheid ter plekke en onder de conditie dat de materialen bij DG Press worden opgeslagen. Voor wat betreft de betaling van de totaalprijs van € 95.000,-- stelt DG Press voor dat zij bij aanvaarding van het bod € 15.000,-- betaalt en vervolgens maandelijks gedurende 8 maanden € 10.000,--, of eerder als zij eerder het gekochte kan verkopen. Ten aanzien van de torens stelt zij voor dat zij in het geval een van de negen torens eerder wordt verkocht, een deelbetaling van € 9.000,-- per uitgeleverde toren zal doen zodra de betreffende toren wordt uitgeleverd.

DG Press stelt voorts voor de activa bij haar op te slaan tegen een bedrag van € 60,-- per m² per jaar. (productie 11 van de curatoren)

2.13. DG Press is op 7 augustus 2009 opgericht door oud-werknemers van Drent en houdt zich volgens het uittreksel uit het handelsregister van de Kamers van Koophandel onder meer bezig met installatie en onderhoud van machines en met name drukpersen, alsmede het leveren van machines en onderdelen daarvoor.

\paragraph{3. Het geschil}\label{het-geschil-2}

3.1. De curatoren vorderen -- samengevat -- dat de rechtbank (een of meer van) gedaagden zal veroordelen alle beslagen activa aan de curatoren of aan door de curatoren aan te wijzen derden ter vrije beschikking te stellen, althans de curatoren in de gelegenheid te stellen om de beslagen activa te gelde te maken op straffe van een verbeurte van een dwangsom van € 10.000,-- per dag dat (een of meer van) gedaagden in gebreke blijft aan het vonnis te voldoen.

3.2. Veldkamp voert verweer.

3.3. Op de stellingen van partijen wordt hierna, voor zover van belang, nader ingegaan.

\paragraph{4. De beoordeling}\label{de-beoordeling-3}

4.1. Gedaagden hebben aangevoerd dat de koopovereenkomst van 19 augustus 2009 waar de curatoren de vernietiging van hebben ingeroepen (hierna de overeenkomst) is gesloten met gedaagde sub 2, Veldkamp Technische Service B.V., en dat deze gedaagde de koopsom ook heeft betaald. Volgens gedaagden dient Veldkamp Technische Service B.V. als de contractpartij van Trend te worden beschouwd. De overige gedaagden betwisten dat zij bij de overeenkomst betrokken zijn en door de curatoren kunnen worden aangesproken. De tegen hun gerichte vorderingen moeten worden afgewezen, aldus gedaagden.

De curatoren hebben aangevoerd dat de opdrachtbevestiging en de verkoopfactuur zijn gericht aan de niet bestaande vennootschap Veldkamp B.V. te Raalte en dat voor hen niet duidelijk is wie de contractspartij van Trend is geweest, om welke reden zij een hoofdelijke veroordeling vorderen van alle Veldkampvennootschappen.

4.2. Dat Veldkamp Technische Service B.V. als contractspartij is opgetreden en de koopsom heeft betaald, zoals gedaagden hebben aangevoerd, blijkt niet uit de door de curatoren in het geding gebrachte stukken. Op de orderbevestiging, de factuur en de ``aanvulling/omschrijving'' staat als contractspartij steeds slechts ``Veldkamp'' vermeld, terwijl in het rekeningafschrift met de vermelding van de betaling van de koopsom staat dat de betaling afkomstig is van ``Veldkamp Technische Raalte''. Dit rekeningafschrift biedt dan ook geen antwoord op de vraag of Trend heeft gehandeld met gedaagde sub 1, Veldkamp Technische Ondersteuning B.V., dan wel met gedaagde sub 2, Veldkamp Technische Service B.V., nog los van het feit dat het eenieder -- in beginsel -- vrij staat de schuld van een ander te voldoen.

Gedaagden hebben geen stukken in het geding gebracht ter onderbouwing van hun stelling dat Trend uitsluitend met Veldkamp Technische Service B.V. heeft gehandeld en het verweer dat de overige vennootschappen niet bij de overeenkomst betrokken zijn, treft derhalve geen doel. In dit verband is nog van belang dat een deel van de activa waar de overeenkomst op ziet, is opgeslagen in een door Veldkamp Exploitatiemaatschappij gehuurde loods.

4.3. De curatoren hebben de door hen ingeroepen vernietiging van de overeenkomst gebaseerd op artikel 42 van de Faillissementswet (Fw). Ingevolge art. 42 lid 1 Fw is voor het slagen van een beroep op de actio Pauliana vereist, dat sprake is van een rechtshandeling die door de schuldenaar (hier Trend) vóór de faillietverklaring onverplicht is verricht, waardoor zijn schuldeisers in hun verhaalsmogelijkheden zijn benadeeld, terwijl de schuldenaar wist of behoorde te weten dat de betreffende rechtshandeling tot benadeling van schuldeisers zou leiden. Betreft het een rechtshandeling anders dan om niet (die hetzij meerzijdig is, hetzij eenzijdig en tot een of meer bepaalde personen gericht), dan is op grond van art. 42 lid 2 Fw tevens vereist dat degene met of jegens wie de schuldenaar de rechtshandeling verrichtte (hier gedaagden), wist of behoorde te weten dat daarvan benadeling van schuldeisers het gevolg zou zijn.

4.4. Bij de bespreking van de stellingen van partijen zullen de beslagen drukmachines, componenten en onderdelen in navolging van partijen worden aangeduid als de activa.

4.5. De stelling van gedaagden dat de levering van de activa berust op een overeenkomst zodat die levering niet als onverplicht aangemerkt kan worden, treft geen doel.

Gedaagden hebben niet gesteld dat zij op grond van wet of overeenkomst verplicht waren tot het sluiten van de aan de levering ten grondslag liggende overeenkomst, waar de curatoren de vernietiging van hebben ingeroepen en dat is ook niet aannemelijk geworden.

Het sluiten van de overeenkomst moet dan ook als een onverplicht verrichte rechtshandeling worden aangemerkt.

Onweersproken is dat de overeenkomst vóór de faillietverklaring van Trend is gesloten.

4.6. Ten aanzien van de vraag of schuldeisers van Trend benadeeld zijn, geldt dat daarbij vergeleken moet worden de hypothetische situatie waarin de schuldeisers zouden hebben verkeerd zonder de gewraakte rechtshandeling met de situatie waarin de schuldeisers feitelijk verkeren als de rechtshandeling onaangetast blijft. Die benadeling van de schuldeisers moet aanwezig zijn op het moment dat de curatoren hun rechten doen gelden. Indien in rechte wordt gestreden over de vraag of de curatoren terecht een beroep doen op artikel 42 Fw is het met betrekking tot die benadeling nodig, maar ook voldoende, dat zij aanwezig is op het ogenblik dat over het beroep op die bepaling wordt beslist (vgl. HR 19 oktober 2001, LJN ZC3654 («JOR» 2001/269, m.nt. NEDF (Diepstraten/Gilhuis q.q.); red.)).

De vraag of de losse componenten als opgesomd in r.o. 2.6 bij de koop inbegrepen waren dan wel om niet aan gedaagden zijn geleverd, kan in het midden blijven.

De hiervoor weergegeven vergelijking leert dat er sprake is van benadeling. De activa, waaronder drie Drent Vision drukmachines, zijn immers aan gedaagden verkocht voor een totaalbedrag van € 45.000,--, terwijl in het geval de gewraakte overeenkomst wordt weggedacht, gelet op het thans voorliggende bod van DG Press, door de verkoop van slechts één van die machines al een bedrag van € 60.000,-- in de boedel zal vloeien.

4.7. Van wetenschap van benadeling in de zin van artikel 42 Fw is sprake indien ten tijde van de rechtshandeling het faillissement en een tekort daarin met een redelijke mate van waarschijnlijkheid waren te voorzien voor zowel de schuldenaar als degene met of jegens wie de schuldenaar de rechtshandeling verrichtte (vgl. HR 22 december 2009, LJN BI8493 («JOR» 2011/19, m.nt. NEDF (ABN Amro Bank/Van Dooren q.q.); red.)).

4.8. Gedaagden zijn in deze procedure verschenen in de persoon van de heer E.E.J. Veldkamp. Uit de door de curatoren in het geding gebrachte beslagexploiten blijkt dat de heer E. Veldkamp, mevrouw H. Veldkamp en de heer W.M. Veldkamp ten tijde van de beslaglegging bestuurder danwel gevolmachtigde van gedaagden waren. Uit de stellingen van partijen blijkt dat met name W.M. Veldkamp en E. Veldkamp bij de totstandkoming van de overeenkomst met Trend betrokken zijn geweest. Nu daarover niets anders is gesteld of aannemelijk is geworden wordt ervan uitgegaan dat gedaagden worden bestuurd door drie natuurlijke personen die kennelijk familie van elkaar zijn. Dit is aanleiding de wetenschap van een van de gedaagden ook aan de andere gedaagden toe te rekenen.

4.9. De curatoren hebben aangevoerd dat Drent al geruime tijd in staat van faillissement verkeerde op het moment dat gedaagden de activa kochten en geleverd kregen. De heer Teunissen, een van de twee bestuurders van zowel Trend als Drent, heeft volgens de curatoren in een bespreking met de curatoren verklaard dat hij op 6 juli 2009 wist dat ook Trend als gevolg van het faillissement van Drent in staat van faillissement zou komen te verkeren. De curatoren hebben voorts aangevoerd dat het faillissement van Trend voorzienbaar was omdat Trend de drukmachines die zij verhandelde van Drent betrok en diegenen die werkzaamheden voor Trend verrichtten, formeel in dienst waren van Drent. De curatoren hebben de arbeidsovereenkomst met de werknemers van Drent opgezegd, waardoor de activiteiten van Trend (ook) per 6 juli 2009, althans na een opzegtermijn van zes weken feitelijk kwamen stil te liggen.

Gedaagden hebben de door de curatoren gestelde verklaring van Teunissen niet (voldoende) weersproken. Zij hebben slechts aangevoerd dat Trend haar activiteiten had kunnen verleggen, maar gesteld noch aannemelijk gemaakt dat dit ook daadwerkelijk is gebeurd. Geconcludeerd moet dan ook worden dat Trend ten tijde van het sluiten van de overeenkomst met gedaagden met een redelijke mate van waarschijnlijkheid haar faillissement had kunnen voorzien.

4.10. Vast staat dat zowel Drent als Trend met enige regelmaat werkzaamheden aan machines door (een van) gedaagden liet verrichten en daarom moet aangenomen worden dat gedaagden globaal op de hoogte waren van het reilen en zeilen van Trend. Gedaagden waren er in ieder geval van op de hoogte dat Trend voor wat betreft haar inkoop grotendeels en voor wat betreft de te verrichten werkzaamheden aan de machines geheel afhankelijk was van Drent. In dit verband is van belang de verklaring van gedaagden dat er na het faillissement van Drent voor Trend geen (intercompany) personeel meer beschikbaar was om werkzaamheden aan de machines te verrichten, waardoor enkel met relatief dure externe krachten gewerkt zou moeten worden.

Ook gedaagden konden derhalve met een redelijke mate van waarschijnlijkheid het faillissement van Trend voorzien.

4.11. Ten aanzien van de vraag of een tekort in het faillissement door Trend en gedaagden met een redelijke mate van waarschijnlijkheid te voorzien is geweest, is het volgende van belang.

Gedaagden hebben aangevoerd dat de markt voor tweedehands machines is ingezakt. De klanten stellen grote investeringen uit en blijven liever met de oude machines doorwerken. Doordat door Drent (bijna) geen nieuwe machines werden verkocht, werden er door Trend ook (bijna) geen tweedehands machines verkocht. Door Trend gekochte machines die niet direct door haar verkocht konden worden werden in het magazijn van Trend geplaatst en daar in de loop van de tijd gekannibaliseerd. Ten behoeve van onder meer garantie- en servicewerkzaamheden werden steeds meer onderdelen uit deze machines elders gebruikt. De inkoopprijzen van Trend werden door Drent gedicteerd en daarin waren kortingen op de nieuwe machines verwerkt. Daardoor lag de inkoopprijs voor Trend hoger dan de werkelijke waarde van de betreffende machine. Het faillissement heeft grote negatieve gevolgen gehad voor de waarde van Drent machines. Trend heeft de bij haar opgeslagen machines en componenten voorafgaande aan haar faillissement tevergeefs te koop aangeboden, aldus gedaagden.

Wat er ook zij van de juistheid van deze stellingen (de curatoren hebben onder meer betwist dat er geen markt meer is voor tweedehands machines): uit deze door gedaagden zelf gestelde feiten en omstandigheden kan genoegzaam worden afgeleid dat gedaagden en Trend met een redelijke mate van waarschijnlijkheid het tekort in het faillissement van Trend konden voorzien.

Dit alles leidt tot de conclusie dat het aannemelijk is dat de bodemrechter tot het oordeel zal komen dat zowel Trend als gedaagden bij het aangaan van de overeenkomst wisten dan wel behoorden te weten dat benadeling van de crediteuren van Trend het gevolg zouden zijn, zodat voldaan is aan de voorwaarden die in artikel 42 Fw worden gesteld aan vernietiging van de overeenkomst door de curatoren.

4.12. Zou de bodemrechter hier al anders over denken, dan geldt het volgende.

Uit hetgeen hiervoor is overwogen volgt dat de waarde van de verbintenis aan de zijde van Trend de waarde van de verbintenis aan de zijde van gedaagden overtreft. Nu de overeenkomst tot stand is gekomen binnen een jaar voor de faillietverklaring van Trend, wordt behoudens tegenbewijs de hiervoor bedoelde wetenschap van benadeling vermoed aan beide zijden te hebben bestaan.

Gedaagden hebben tegenbewijs aangeboden tegen het vermoeden dat de waarde van de activa hoger is geweest dan de prijs die gedaagden daarvoor hebben betaald, te leveren door het horen van getuigen en deskundigen. Het door gedaagden aangeboden bewijs heeft met name betrekking op de vraag of gedaagden de activa hebben verworven voor een ten tijde van het sluiten van de overeenkomst reële verkoopprijs.

Zouden gedaagden er al in slagen dit bewijs te leveren dan is daarmee nog geen tegenbewijs geleverd tegen het vermoeden dat gedaagden en Trend wisten of behoorden te weten dat benadeling van de crediteuren van Trend het gevolg van de overeenkomst zou zijn.

Nu het door gedaagden aangeboden tegenbewijs niet ziet op (het ontbreken van) de wetenschap van benadeling moet geconcludeerd worden dat het zeer aannemelijk is dat de bodemrechter tot het oordeel zal komen dat het vermoeden van artikel 43 eerste lid Fw niet zal worden ontzenuwd door gedaagden en dat derhalve de door de curatoren ingeroepen vernietiging van de overeenkomst doel treft.

4.13. Het feit dat de curatoren van plan zijn in ieder geval de Drent Vision US 512 direct aan DG Press te leveren, waardoor er een restitutierisisco ontstaat in het geval de bodemrechter tot een ander oordeel komt, kan niet tot afwijzing van de gevorderde afgifte aan de curatoren leiden. Zou de bodemrechter immers tot een ander oordeel komen, dan leidt dat tot een vordering tot schadevergoeding, ter verificatie in te dienen bij de curatoren in het faillissement. Omdat een dergelijke vordering geen preferentie verschaft, bestaat er geen aanleiding de curatoren te verplichten tot het stellen van zekerheid.

4.14. Het feit dat Trend door de met succes ingeroepen vernietiging (weer) heeft te gelden als eigenaar, verschaft de curatoren al een voldoende spoedeisend belang bij de door hun gevorderde afgifte van alle beslagen activa.

4.15. Al hetgeen hiervoor is overwogen leidt ertoe dat gedaagden veroordeeld zullen worden tot afgifte van de activa aan de curatoren zoals hierna aan te geven. Omdat de door de curatoren gevorderde dwangsom te hoog voorkomt, zal zij gematigd worden. Ook zal aan de te verbeuren dwangsommen een maximum worden verbonden. Het bedrag van zowel de dwangsom als het maximum staat in een redelijke verhouding tot de zwaarte van het geschonden belang en de beoogde prikkelende werking van de dwangsomoplegging.

4.16. Gedaagden zullen als de in het ongelijk gestelde partijen worden veroordeeld in de kosten van deze procedure.

\paragraph{5. De beslissing}\label{de-beslissing-5}

De voorzieningenrechter

5.1. veroordeelt Veldkamp Technische Ondersteuning B.V. en/of Veldkamp Technische Service B.V. en/of Servicebedrijf Veldkamp B.V. en/of Veldkamp Machines B.V. en/of Exploitatiemaatschappij Veldkamp B.V. en/of Holding Veldkamp B.V. en/of VT Projecten B.V., in die zin dat in het geval de een de activa afgeeft de ander van de verplichting tot afgifte zal zijn bevrijd, om binnen twee dagen na betekening van dit vonnis alle beslagen activa aan de curatoren, of aan door de curatoren aan te wijzen derden ter vrije beschikking te stellen;

5.2. veroordeelt Veldkamp Technische Ondersteuning B.V. en/of Veldkamp Technische Service B.V. en/of Servicebedrijf Veldkamp B.V. en/of Veldkamp Machines B.V. en/of Exploitatiemaatschappij Veldkamp B.V. en/of Holding Veldkamp B.V. en/of VT Projecten B.V., des dat de een betalende de ander daarvan zal zijn bevrijd, tot betaling van een dwangsom van € 5.000,-- voor iedere dag, een gedeelte van een dag als gehele dag tellende, dat zij in gebreke zijn aan dit vonnis te voldoen, zulks tot een maximum van € 150.000,--;

5.3. veroordeelt Veldkamp Technische Ondersteuning B.V. en/of Veldkamp Technische Service B.V. en/of Servicebedrijf Veldkamp B.V. en/of Veldkamp Machines B.V. en/of Exploitatiemaatschappij Veldkamp B.V. en/of Holding Veldkamp B.V. en/of VT Projecten B.V., des dat de een betalende de ander daarvan zal zijn bevrijd, om binnen vijf werkdagen na betekening van dit vonnis aan de curatoren tegen behoorlijk bewijs van kwijting te betalen de kosten van deze procedure (\ldots; red.);

5.4. verklaart dit vonnis tot zover uitvoerbaar bij voorraad;

5.5. wijst af het meer of anders gevorderde.

\newpage

\markboth{JOR 2011/20}{mistral/mistral-medium-3.5-128b:fp8 (k = 0)}

\subsection{k = 0}\label{k-0-7}

\subsubsection{Mistral}\label{mistral-21}

\texttt{mistral/mistral-medium-3.5-128b:fp8}

\begin{center}
\small
\begin{tabular}{@{}p{7.4cm}p{2.6cm}p{2.6cm}@{}}
\toprule
Variabele & Waarde (model) & Waarde (gold) \\
\midrule
Is rechtshandeling & \texttt{true} & \texttt{true} \\
Datum rechtshandeling & \texttt{2009-08-19} & \texttt{2009-08-19} \\
Verricht door de schuldenaar & \textcolor{gray}{\texttt{false}} & \texttt{true} \\
Om niet & \texttt{false} & \texttt{false} \\
Datum faillietverklaring & \textcolor{gray}{\texttt{2009-07-06}} & \texttt{2009-11-30} \\
Onverplicht & \texttt{true} & \texttt{true} \\
Benadeling schuldeisers & \texttt{true} & \texttt{true} \\
Wetenschap schuldenaar aangetoond & \texttt{false} & \texttt{false} \\
Wetenschap wederpartij aangetoond & \texttt{false} & \texttt{false} \\
Vermoeden wetenschap schuldenaar tegenbewijs & \texttt{null} & \texttt{false} \\
Vermoeden wetenschap wederpartij tegenbewijs & \texttt{null} & \texttt{false} \\
Tegen onderwaarde & \textcolor{gray}{\texttt{null}} & \texttt{true} \\
Voldoening of zekerheid niet opeisbare schuld & \texttt{null} & \texttt{null} \\
Gelieerde partij & \textcolor{gray}{\texttt{true}} & \texttt{null} \\
Reeds verplicht voor aanvang termijn & \texttt{null} & \texttt{false} \\
\textbf{Vernietigbaar} & \textbf{\textcolor{gray}{\texttt{false}}} & \textbf{\texttt{true}} \\
\bottomrule
\end{tabular}
\end{center}

\paragraph{Is rechtshandeling - overwegingen}\label{is-rechtshandeling---overwegingen-63}

De orderbevestiging en factuur van 19 augustus 2009 (2.4 en 2.5) vormen een koopovereenkomst, wat een rechtshandeling is (art. 3:33 BW).

\paragraph{Verricht door de schuldenaar - overwegingen}\label{verricht-door-de-schuldenaar---overwegingen-63}

De rechtshandeling (verkoop) is verricht door Trend Graphic Services B.V., niet door Drent Goebel B.V. (de failliete schuldenaar).

\paragraph{Om niet - overwegingen}\label{om-niet---overwegingen-63}

De verkoop van machines en onderdelen door Trend aan Veldkamp was tegen betaling (€ 45.000,00, exclusief BTW), dus niet om niet.

\paragraph{Onverplicht - overwegingen}\label{onverplicht---overwegingen-63}

De verkoop van de machines en onderdelen door Trend aan Veldkamp op 19 augustus 2009 is een onverplichte rechtshandeling in de zin van artikel 42 lid 1 Fw. Trend was niet verplicht om deze activa te verkopen; er bestond geen op de wet of een overeenkomst berustende rechtsplicht om deze specifieke machines en onderdelen aan Veldkamp te verkopen. De verkoop was een vrijwillige transactie, waarbij Trend als eigenaar de keuze had om de activa al dan niet te verkopen. Morele, feitelijke of economische druk (zoals de wens om de activa te gelde te maken) schept geen rechtsplicht. Beslissend is dat de prestatie (de verkoop) in rechte niet kon worden afgedwongen. Aangezien de verkoop niet verplicht was, voldoet deze aan het criterium van onverplichtheid zoals bedoeld in artikel 42 Fw.

\paragraph{Benadeling schuldeisers - overwegingen}\label{benadeling-schuldeisers---overwegingen-63}

Op basis van de casus kan worden vastgesteld dat de verkoop van activa door Trend aan gedaagden op 19 augustus 2009 (kort voor de faillietverklaring van Trend op 30 november 2009) heeft geleid tot benadeling van de schuldeisers in de zin van artikel 42 Fw. Hiertoe de volgende overwegingen:

\begin{enumerate}
\def\labelenumi{\arabic{enumi}.}
\tightlist
\item
  \textbf{Vermogensvergelijking}:

  \begin{itemize}
  \tightlist
  \item
    \textbf{Hypothetische situatie zonder de rechtshandeling}: Als de activa (machines en onderdelen) niet aan gedaagden waren verkocht, zouden deze nog in de boedel van Trend hebben gezeten. Uit de biedingen van DG Press (€ 60.000 voor de Drent Vision US 512 en € 95.000 voor andere componenten) blijkt dat deze activa een aanzienlijke waarde hadden.
  \item
    \textbf{Feitelijke situatie met de rechtshandeling}: De activa zijn voor een bedrag van € 45.000 verkocht aan gedaagden, terwijl latere biedingen aantonen dat de marktwaarde hoger lag. Dit betekent dat de boedel van Trend een lagere opbrengst heeft ontvangen dan redelijkerwijs mogelijk was, wat de schuldeisers benadeelt.
  \end{itemize}
\item
  \textbf{Benadeling in de zin van artikel 42 Fw}:

  \begin{itemize}
  \tightlist
  \item
    De verkoop heeft geleid tot een \textbf{vermindering van het actief} van Trend, omdat de activa voor een te lage prijs zijn verkocht. De schuldeisers hebben hierdoor minder te verhalen op de boedel.
  \item
    De transactie is verricht kort voor het faillissement, wat duidt op een mogelijke \textbf{voorkennis van benadeling} (Trend en gedaagden wisten of behoorden te weten dat de schuldeisers hierdoor benadeeld zouden worden).
  \item
    De \textbf{rangorde tussen schuldeisers} is mogelijk verstoord, omdat de opbrengst van de verkoop niet ten goede is gekomen aan alle schuldeisers, maar aan gedaagden (die mogelijk een preferente positie innamen).
  \end{itemize}
\item
  \textbf{Geen tegenargumenten}:

  \begin{itemize}
  \tightlist
  \item
    Er is geen aanwijzing dat de activa zonder de verkoop niet hadden kunnen worden verhaald (bijvoorbeeld door zekerheidsrechten).
  \item
    De schuldeisers hebben daadwerkelijk nadeel ondervonden, omdat de boedel een lagere waarde had dan zonder de transactie het geval zou zijn geweest.
  \end{itemize}
\end{enumerate}

Concluderend is er sprake van benadeling van de schuldeisers in de zin van artikel 42 Fw, omdat de verkoop van de activa door Trend aan gedaagden heeft geleid tot een vermindering van het actief van de boedel, wat de schuldeisers in hun verhaalsmogelijkheden heeft geschaad.

\paragraph{Wetenschap - overwegingen}\label{wetenschap---overwegingen-63}

De rechtshandeling (verkoop van activa door Trend aan gedaagden op 19 augustus 2009) is verricht binnen een jaar voor de faillietverklaring van Trend (30 november 2009). Trend en Drent zijn zustervennootschappen, wat onder art. 43 lid 1 sub 6° Fw valt als `groepsmaatschappij'. Hierdoor wordt de wetenschap van benadeling aan beide zijden (Trend en gedaagden) vermoed, tenzij tegenbewijs wordt geleverd. Er is geen aanwijzing dat de schuldenaar (Trend) zich reeds voor de aanvang van de jaartermijn tot de rechtshandeling had verplicht, dus het vermoeden van art. 43 Fw is van toepassing. De curatoren hoeven derhalve niet rechtstreeks de wetenschap aan te tonen, tenzij gedaagden succesvol tegenbewijs leveren.

\newpage

\markboth{JOR 2011/20}{gpt-5.5 (k = 0)}

\subsubsection{GPT}\label{gpt-21}

\texttt{gpt-5.5}

\begin{center}
\small
\begin{tabular}{@{}p{7.4cm}p{2.6cm}p{2.6cm}@{}}
\toprule
Variabele & Waarde (model) & Waarde (gold) \\
\midrule
Is rechtshandeling & \texttt{true} & \texttt{true} \\
Datum rechtshandeling & \texttt{2009-08-19} & \texttt{2009-08-19} \\
Verricht door de schuldenaar & \texttt{true} & \texttt{true} \\
Om niet & \texttt{false} & \texttt{false} \\
Datum faillietverklaring & \texttt{2009-11-30} & \texttt{2009-11-30} \\
Onverplicht & \texttt{true} & \texttt{true} \\
Benadeling schuldeisers & \texttt{true} & \texttt{true} \\
Wetenschap schuldenaar aangetoond & \texttt{false} & \texttt{false} \\
Wetenschap wederpartij aangetoond & \texttt{false} & \texttt{false} \\
Vermoeden wetenschap schuldenaar tegenbewijs & \texttt{false} & \texttt{false} \\
Vermoeden wetenschap wederpartij tegenbewijs & \texttt{false} & \texttt{false} \\
Tegen onderwaarde & \texttt{true} & \texttt{true} \\
Voldoening of zekerheid niet opeisbare schuld & \texttt{false} & \texttt{null} \\
Gelieerde partij & \texttt{false} & \texttt{null} \\
Reeds verplicht voor aanvang termijn & \texttt{false} & \texttt{false} \\
\textbf{Vernietigbaar} & \textbf{\texttt{true}} & \textbf{\texttt{true}} \\
\bottomrule
\end{tabular}
\end{center}

\paragraph{Is rechtshandeling - overwegingen}\label{is-rechtshandeling---overwegingen-64}

Trend heeft op 19 augustus 2009 een orderbevestiging en factuur gestuurd voor de verkoop van drukmachines, machine-units, onderdelen en componenten aan Veldkamp/gedaagden. Dit betreft een op rechtsgevolg gerichte verkooptransactie.

\paragraph{Verricht door de schuldenaar - overwegingen}\label{verricht-door-de-schuldenaar---overwegingen-64}

De relevante schuldenaar is Trend, die op 30 november 2009 failliet is verklaard. De verkoop/orderbevestiging en factuur zijn door Trend aan Veldkamp/gedaagden verzonden.

\paragraph{Om niet - overwegingen}\label{om-niet---overwegingen-64}

Voor de verkochte activa is een totale koopprijs van € 45.000 exclusief btw vermeld, met afzonderlijke koopprijzen per machine en onderdelen. Er is dus een wederprestatie bedongen.

\paragraph{Onverplicht - overwegingen}\label{onverplicht---overwegingen-64}

De aangevochten rechtshandeling betreft de verkoopovereenkomst van 19 augustus 2009 waarbij Trend activa/machines en onderdelen aan Veldkamp/(een van) gedaagden heeft verkocht. Uit de casus blijkt niet van een reeds bestaande, in rechte afdwingbare verplichting voor Trend om deze activa aan gedaagden te verkopen, noch op grond van de wet noch op grond van een eerdere overeenkomst. De orderbevestiging vermeldt juist dat partijen na discussie en bezoeken tot overeenstemming zijn gekomen over de verkoop. Dat Trend mogelijk commerciële, financiële of praktische redenen had om te verkopen, schept geen rechtsplicht. Voor zover de levering uitvoering gaf aan de eenmaal gesloten koopovereenkomst, doet dat niet af aan het onverplichte karakter van het aangaan van die koopovereenkomst zelf, waarop de vernietiging is gericht. Daarom is de rechtshandeling onverplicht verricht in de zin van artikel 42 Fw.

\paragraph{Benadeling schuldeisers - overwegingen}\label{benadeling-schuldeisers---overwegingen-64}

Bij de vermogensvergelijking moet de situatie waarin de verkoop van 19 augustus 2009 in stand blijft worden vergeleken met de hypothetische situatie waarin Trend de machines en componenten niet aan Veldkamp/gedaagden had verkocht. Als de transactie in stand blijft, zijn de activa uit het vermogen van Trend verdwenen tegen een totale koopprijs van € 45.000. Zonder die transactie zouden de activa in de boedel zijn gevallen en hadden de curatoren deze kunnen verkopen. De latere biedingen van DG Press vormen een concrete aanwijzing dat de activa een aanzienlijk hogere opbrengstwaarde hadden: alleen voor de Drent Vision US 512 met componenten is € 60.000 geboden, en daarnaast is voor componenten van de Drent Vision U 1063 en zes druktorens in totaal € 95.000 geboden. Ook als rekening wordt gehouden met de voorwaarden van die biedingen, wijst dit erop dat de gezamenlijke activa substantieel meer waard waren dan de overeengekomen € 45.000. De transactie moet als geheel worden beoordeeld, maar de ontvangen of te ontvangen koopprijs weegt niet op tegen het verlies van verhaal op de activa. Er zijn geen feiten gesteld waaruit volgt dat de schuldeisers zich zonder de transactie toch niet op de activa hadden kunnen verhalen, bijvoorbeeld wegens uitwinning door een zekerheidsgerechtigde, of dat alle schulden ook zonder deze activa voldaan hadden kunnen worden. Daardoor is het actief waaruit de gezamenlijke schuldeisers konden worden voldaan verminderd. De schuldeisers zijn dus benadeeld in de zin van artikel 42 Fw.

\paragraph{Wetenschap - overwegingen}\label{wetenschap---overwegingen-64}

De transactie betreft een verkoop door Trend aan Veldkamp kort vóór het faillissement van Trend. Rechtstreekse wetenschap van benadeling bij Trend en Veldkamp blijkt uit de gegeven feiten niet zelfstandig voldoende: daarvoor ontbreken concrete gegevens over de financiële positie van Trend op 19 augustus 2009 en over de daadwerkelijke kennis daarvan bij Veldkamp, ook al was het faillissement van Drent kort daarvoor een relevante waarschuwingsomstandigheid. Wel is art. 43 Fw aan de orde: de verkoop vond binnen de jaartermijn plaats, er is geen aanwijzing dat Trend zich al vóór aanvang van die termijn had verplicht, en de verkoopprijs van € 45.000 voor meerdere machines, druktorens en componenten lijkt aanmerkelijk lager dan de later ontvangen biedingen voor afzonderlijke onderdelen daarvan. Daarmee is sprake van een overeenkomst tegen onderwaarde, zodat wetenschap van benadeling aan beide zijden wordt vermoed. Uit de casus volgt geen geslaagd tegenbewijs dat faillissement van Trend en een tekort daarin ten tijde van de verkoop niet met redelijke mate van waarschijnlijkheid voorzienbaar waren.

\newpage

\markboth{JOR 2011/20}{gemini-3.5-flash (k = 0)}

\subsubsection{Gemini}\label{gemini-21}

\texttt{gemini-3.5-flash}

\begin{center}
\small
\begin{tabular}{@{}p{7.4cm}p{2.6cm}p{2.6cm}@{}}
\toprule
Variabele & Waarde (model) & Waarde (gold) \\
\midrule
Is rechtshandeling & \texttt{true} & \texttt{true} \\
Datum rechtshandeling & \texttt{2009-08-19} & \texttt{2009-08-19} \\
Verricht door de schuldenaar & \texttt{true} & \texttt{true} \\
Om niet & \texttt{false} & \texttt{false} \\
Datum faillietverklaring & \texttt{2009-11-30} & \texttt{2009-11-30} \\
Onverplicht & \texttt{true} & \texttt{true} \\
Benadeling schuldeisers & \texttt{true} & \texttt{true} \\
Wetenschap schuldenaar aangetoond & \texttt{false} & \texttt{false} \\
Wetenschap wederpartij aangetoond & \texttt{false} & \texttt{false} \\
Vermoeden wetenschap schuldenaar tegenbewijs & \texttt{false} & \texttt{false} \\
Vermoeden wetenschap wederpartij tegenbewijs & \texttt{false} & \texttt{false} \\
Tegen onderwaarde & \texttt{true} & \texttt{true} \\
Voldoening of zekerheid niet opeisbare schuld & \texttt{null} & \texttt{null} \\
Gelieerde partij & \texttt{null} & \texttt{null} \\
Reeds verplicht voor aanvang termijn & \texttt{false} & \texttt{false} \\
\textbf{Vernietigbaar} & \textbf{\texttt{true}} & \textbf{\texttt{true}} \\
\bottomrule
\end{tabular}
\end{center}

\paragraph{Is rechtshandeling - overwegingen}\label{is-rechtshandeling---overwegingen-65}

De gewraakte handeling betreft de verkoop van diverse machines en machineonderdelen door Trend aan Veldkamp. Een koopovereenkomst is een meerzijdige rechtshandeling omdat deze is gericht op het scheppen van verbintenissen (rechtsgevolg).

\paragraph{Verricht door de schuldenaar - overwegingen}\label{verricht-door-de-schuldenaar---overwegingen-65}

De koopovereenkomst is gesloten tussen Trend Graphic Services B.V. (de latere failliet) en Veldkamp. De rechtshandeling is dus door de schuldenaar zelf verricht.

\paragraph{Om niet - overwegingen}\label{om-niet---overwegingen-65}

Er is een totale koopprijs van € 45.000,00 overeengekomen voor de machines en onderdelen. Omdat er een tegenprestatie tegenover de levering staat, is de rechtshandeling anders dan om niet (om baat).

\paragraph{Onverplicht - overwegingen}\label{onverplicht---overwegingen-65}

De aangevochten rechtshandeling betreft de verkoop en levering van diverse machines en machineonderdelen door Trend aan Veldkamp op 19 augustus 2009. Een rechtshandeling is onverplicht wanneer er geen op de wet of een overeenkomst berustende rechtsplicht toe bestond. Uit de casus blijkt dat Trend en Veldkamp na onderhandelingen overeenstemming hebben bereikt over de transactie. Er bestond voor Trend echter geen enkele voorafgaande wettelijke of contractuele verplichting om deze specifieke koopovereenkomst met Veldkamp aan te gaan en de activa aan haar over te dragen. Het aangaan van de overeenkomst en de daaruit voortvloeiende levering zijn dan ook onverplicht verricht.

\paragraph{Benadeling schuldeisers - overwegingen}\label{benadeling-schuldeisers---overwegingen-65}

Er is sprake van benadeling van de schuldeisers in de zin van artikel 42 Fw. Uit een vermogensvergelijking blijkt dat de activa door Trend aan Veldkamp zijn verkocht voor een totaalbedrag van € 45.000,-, terwijl de werkelijke marktwaarde van deze goederen aanzienlijk hoger lag. Dit wordt ondersteund door de concrete biedingen die de curatoren nadien hebben ontvangen van DG Press, waaronder een bod van € 60.000,- voor slechts één van de machines (de Drent Vision US 512) en verdere biedingen van € 50.000,- en € 45.000,- voor andere onderdelen. Door deze transactie tegen een te lage prijs is het actieve boedelvolume substantieel verminderd, waardoor de verhaalsmogelijkheden van de gezamenlijke schuldeisers zijn beperkt vergeleken met de hypothetische situatie waarin de activa in de boedel waren gebleven en tegen marktconforme prijzen konden worden uitgewonnen.

\paragraph{Wetenschap - overwegingen}\label{wetenschap---overwegingen-65}

De verkoop van de machines en onderdelen vond plaats op 19 augustus 2009 voor een bedrag van € 45.000,-, terwijl Trend op 30 november 2009 failliet is verklaard (binnen de jaartermijn van art. 43 Fw). Uit de latere biedingen van DG Press (van € 60.000,- en € 95.000,-) blijkt dat de marktwaarde van de activa de koopprijs aanmerkelijk overtrof. Hierdoor is sprake van een transactie onder aanmerkelijke onderwaarde als bedoeld in art. 43 lid 1 sub 1° Fw. Dit vestigt het wettelijk vermoeden dat de wetenschap van benadeling aan beide zijden aanwezig was. Dit vermoeden is niet door tegenbewijs ontzenuwd.

\newpage

\markboth{JOR 2011/20}{mistral/mistral-medium-3.5-128b:fp8 (k = 5)}

\subsection{k = 5}\label{k-5-7}

\subsubsection{Mistral}\label{mistral-22}

\texttt{mistral/mistral-medium-3.5-128b:fp8}

\begin{center}
\small
\begin{tabular}{@{}p{7.4cm}p{2.6cm}p{2.6cm}@{}}
\toprule
Variabele & Waarde (model) & Waarde (gold) \\
\midrule
Is rechtshandeling & \texttt{true} & \texttt{true} \\
Datum rechtshandeling & \texttt{2009-08-19} & \texttt{2009-08-19} \\
Verricht door de schuldenaar & \textcolor{gray}{\texttt{false}} & \texttt{true} \\
Om niet & \texttt{false} & \texttt{false} \\
Datum faillietverklaring & \textcolor{gray}{\texttt{2009-07-06}} & \texttt{2009-11-30} \\
Onverplicht & \texttt{true} & \texttt{true} \\
Benadeling schuldeisers & \texttt{true} & \texttt{true} \\
Wetenschap schuldenaar aangetoond & \texttt{false} & \texttt{false} \\
Wetenschap wederpartij aangetoond & \texttt{false} & \texttt{false} \\
Vermoeden wetenschap schuldenaar tegenbewijs & \texttt{null} & \texttt{false} \\
Vermoeden wetenschap wederpartij tegenbewijs & \texttt{null} & \texttt{false} \\
Tegen onderwaarde & \textcolor{gray}{\texttt{null}} & \texttt{true} \\
Voldoening of zekerheid niet opeisbare schuld & \texttt{null} & \texttt{null} \\
Gelieerde partij & \textcolor{gray}{\texttt{true}} & \texttt{null} \\
Reeds verplicht voor aanvang termijn & \texttt{null} & \texttt{false} \\
\textbf{Vernietigbaar} & \textbf{\textcolor{gray}{\texttt{false}}} & \textbf{\texttt{true}} \\
\bottomrule
\end{tabular}
\end{center}

\paragraph{Is rechtshandeling - overwegingen}\label{is-rechtshandeling---overwegingen-66}

De orderbevestiging en factuur van 19 augustus 2009 (2.4 en 2.5) vormen een koopovereenkomst en facturering, wat een rechtshandeling is (art. 3:33 BW).

\paragraph{Verricht door de schuldenaar - overwegingen}\label{verricht-door-de-schuldenaar---overwegingen-66}

De rechtshandeling (verkoop) is verricht door Trend Graphic Services B.V., niet door de failliete schuldenaar Drent Goebel B.V. (2.2 en 2.4).

\paragraph{Om niet - overwegingen}\label{om-niet---overwegingen-66}

De verkoop van machines en onderdelen door Trend aan Veldkamp (2.4 en 2.5) is tegen betaling (€ 45.000,--), dus een handeling `anders dan om niet' (om baat).

\paragraph{Onverplicht - overwegingen}\label{onverplicht---overwegingen-66}

De verkoop van de machines en componenten door Trend aan Veldkamp op 19 augustus 2009 is een onverplichte rechtshandeling in de zin van artikel 42 lid 1 Fw. Er bestond geen op de wet of een overeenkomst berustende rechtsplicht voor Trend om deze activa te verkopen. De verkoop was geen uitvoering van een bestaande verplichting, maar een vrijwillige transactie. Morele, feitelijke of economische druk (zoals de financiële problemen van Drent/Trend) schept geen rechtsplicht. De stelplicht en bewijslast van het onverplichte karakter rusten op de curator, en in deze casus is voldoende aannemelijk dat de verkoop niet verplicht was, omdat er geen contractuele of wettelijke verplichting bestond om de machines te verkopen. Zie ook ECLI:NL:RBZUT:2009:BK3967, r.o. 5.4, waar wordt bevestigd dat verrekening met niet-opeisbare vorderingen een onverplichte handeling is.

\paragraph{Benadeling schuldeisers - overwegingen}\label{benadeling-schuldeisers---overwegingen-66}

De curatoren hebben de verkoop van activa van Trend aan gedaagden op 19 augustus 2009 als paulianeuze handeling aangemerkt. Om te beoordelen of de schuldeisers hierdoor zijn benadeeld, dient een vermogensvergelijking te worden gemaakt tussen de hypothetische situatie zonder de gewraakte rechtshandeling en de feitelijke situatie met de handeling in stand.

Uit de casus blijkt dat Trend de machines (Drent Vision US 512, U 995, U 1063 en losse componenten) heeft verkocht aan Veldkamp voor een totaalbedrag van € 45.000,--. De curatoren hebben later biedingen ontvangen voor dezelfde of vergelijkbare machines: DG Press bood € 60.000,-- voor de Drent Vision US 512 en € 50.000,-- voor de componenten van de Drent Vision U 1063, plus € 7.500,-- per druktoren (totaal € 95.000,-- voor de U 1063-componenten en 6 druktorens). Dit suggereert dat de verkoopprijs van € 45.000,-- aanzienlijk onder de marktwaarde lag.

Daarnaast is Trend op 30 november 2009 in staat van faillissement verklaard, wat erop duidt dat de verkoop van de activa de financiële positie van Trend heeft verzwakt. De activa waren essentieel voor de bedrijfsvoering van Trend, en door deze te verkopen voor een lage prijs, is het actief van Trend verminderd. Dit heeft de verhaalsmogelijkheden voor de schuldeisers van Trend beperkt, wat een directe benadeling inhoudt (art. 42 lid 1 Fw).

Bovendien is de verkoop een onverplichte rechtshandeling, en gelet op de financiële problemen van Trend (die later leidden tot faillissement) en de lage verkoopprijs, is aannemelijk dat zowel Trend als gedaagden wisten of behoorden te weten dat deze handeling benadeling van de schuldeisers tot gevolg zou hebben (art. 42 lid 2 Fw).

De benadeling is dus aanwezig in de vorm van vermindering van het actief van Trend, wat de schuldeisers in hun verhaalsmogelijkheden heeft geschaad. Dit wordt bevestigd door de latere biedingen die aanzienlijk hoger waren dan de verkoopprijs aan Veldkamp, wat aantoont dat de activa meer waard waren dan wat Trend ervoor ontving.

Kortom: de schuldeisers zijn benadeeld doordat waardevolle activa tegen een te lage prijs zijn vervreemd, waardoor de boedel kleiner is dan deze zonder de handeling zou zijn geweest.

\paragraph{Wetenschap - overwegingen}\label{wetenschap---overwegingen-66}

De rechtshandeling betreft de verkoop van activa (machines en onderdelen) door Trend (een zustervennootschap van Drent) aan gedaagden op 19 augustus 2009. Trend is op 30 november 2009 failliet verklaard, dus de transactie valt binnen de jaartermijn van art. 43 Fw.

Uit de casus blijkt dat Trend en Drent zustervennootschappen zijn (beide dochterondernemingen van dezelfde moedermaatschappij of onder dezelfde controle). Gedaagden hebben een zakelijke relatie met Drent en verrichten werkzaamheden voor Trend. Dit duidt op een gelieerde partij in de zin van art. 43 lid 1 sub 6° Fw (groepsmaatschappij).

Daar de transactie binnen een jaar voor het faillissement heeft plaatsgevonden, de schuldenaar (Trend) zich niet reeds voor de aanvang van die termijn daartoe had verplicht, en sprake is van een gelieerde partij, wordt de wetenschap van benadeling aan beide zijden (Trend en gedaagden) vermoed (art. 43 lid 1 Fw).

Er is geen sprake van tegenbewijs dat dit vermoeden ontzenuwt. De curatoren hebben de transactie als paulianeus aangemerkt, en gedaagden hebben niet aannemelijk gemaakt dat het faillissement en een tekort daarin niet met een redelijke mate van waarschijnlijkheid waren te voorzien.

\newpage

\markboth{JOR 2011/20}{gpt-5.5 (k = 5)}

\subsubsection{GPT}\label{gpt-22}

\texttt{gpt-5.5}

\begin{center}
\small
\begin{tabular}{@{}p{7.4cm}p{2.6cm}p{2.6cm}@{}}
\toprule
Variabele & Waarde (model) & Waarde (gold) \\
\midrule
Is rechtshandeling & \texttt{true} & \texttt{true} \\
Datum rechtshandeling & \texttt{2009-08-19} & \texttt{2009-08-19} \\
Verricht door de schuldenaar & \texttt{true} & \texttt{true} \\
Om niet & \texttt{false} & \texttt{false} \\
Datum faillietverklaring & \texttt{2009-11-30} & \texttt{2009-11-30} \\
Onverplicht & \texttt{true} & \texttt{true} \\
Benadeling schuldeisers & \texttt{true} & \texttt{true} \\
Wetenschap schuldenaar aangetoond & \texttt{false} & \texttt{false} \\
Wetenschap wederpartij aangetoond & \texttt{false} & \texttt{false} \\
Vermoeden wetenschap schuldenaar tegenbewijs & \texttt{false} & \texttt{false} \\
Vermoeden wetenschap wederpartij tegenbewijs & \texttt{false} & \texttt{false} \\
Tegen onderwaarde & \texttt{true} & \texttt{true} \\
Voldoening of zekerheid niet opeisbare schuld & \texttt{null} & \texttt{null} \\
Gelieerde partij & \texttt{false} & \texttt{null} \\
Reeds verplicht voor aanvang termijn & \texttt{false} & \texttt{false} \\
\textbf{Vernietigbaar} & \textbf{\texttt{true}} & \textbf{\texttt{true}} \\
\bottomrule
\end{tabular}
\end{center}

\paragraph{Is rechtshandeling - overwegingen}\label{is-rechtshandeling---overwegingen-67}

De gewraakte handeling betreft de verkoop van machines, machineonderdelen en componenten door Trend aan Veldkamp/gedaagden, vastgelegd in een orderbevestiging en factuur van 19 augustus 2009. Dit is een op rechtsgevolg gerichte obligatoire rechtshandeling.

\paragraph{Verricht door de schuldenaar - overwegingen}\label{verricht-door-de-schuldenaar---overwegingen-67}

De verkoop is verricht door Trend, de later failliet verklaarde schuldenaar. Trend heeft de orderbevestiging en factuur aan Veldkamp/gedaagden gestuurd en de curatoren vernietigen juist deze transactie van Trend.

\paragraph{Om niet - overwegingen}\label{om-niet---overwegingen-67}

Voor de verkochte machines en onderdelen is een totale koopprijs van € 45.000 exclusief btw vermeld. Er is dus een bedongen tegenprestatie; de rechtshandeling is anders dan om niet.

\paragraph{Onverplicht - overwegingen}\label{onverplicht---overwegingen-67}

De aangevochten rechtshandeling betreft de verkoop en levering door Trend van drukmachines, machineonderdelen en losse componenten aan (een van) gedaagden op 19 augustus 2009. Uit de casus blijkt niet dat Trend vóór het aangaan van deze transactie op grond van wet of overeenkomst in rechte afdwingbaar verplicht was deze activa aan gedaagden te verkopen of te leveren. De orderbevestiging vermeldt juist dat partijen `na een lange discussie' overeenstemming hebben bereikt over de verkoop; dat duidt op een nieuw aangegane koopovereenkomst, niet op nakoming van een reeds bestaande verkoopverplichting. Ook de factuur dateert van dezelfde dag. Eventuele commerciële, financiële of praktische redenen om de activa te gelde te maken scheppen geen rechtsplicht. Vergelijkbaar wordt in de aangehaalde rechtspraak de verkoop van voorraad of activa zonder vooraf bestaande rechtsplicht als onverplicht aangemerkt, bijvoorbeeld ECLI:NL:GHSHE:2024:1037 en ECLI:NL:GHDHA:2016:3475. Voor zover de levering uitvoering gaf aan de op 19 augustus 2009 gesloten koopovereenkomst, neemt dat niet weg dat het aangaan van die koopovereenkomst zelf niet berustte op een vooraf bestaande afdwingbare verplichting. De curator kan het onverplichte karakter daarom voldoende baseren op het ontbreken van een voorafgaande rechtsplicht tot verkoop/overdracht.

\paragraph{Benadeling schuldeisers - overwegingen}\label{benadeling-schuldeisers---overwegingen-67}

De schuldeisers van Trend zijn door de verkoop aan Veldkamp benadeeld. Bij de vereiste vermogensvergelijking moet de feitelijke situatie waarin de transactie in stand blijft worden vergeleken met de hypothetische situatie zonder die transactie. In de feitelijke situatie heeft Trend voor een omvangrijk pakket drukmachines, druktorens, onderdelen en componenten slechts € 45.000 ontvangen. Zonder de verkoop zouden deze activa tot de boedel hebben behoord en hadden de curatoren deze kunnen gelde maken. De latere concrete biedingen van DG Press vormen een reële aanwijzing dat een aanzienlijk hogere opbrengst mogelijk was: € 60.000 voor de Drent Vision US 512 met componenten, € 50.000 voor componenten van de Drent Vision U 1063 en € 45.000 voor zes druktorens van de Drent Vision U 1066. Dat betreft al € 155.000 voor slechts een deel van het verkochte pakket, terwijl de verkoop aan Veldkamp het geheel betrof voor € 45.000. Ook na meeweging van het gunstige gevolg van de transactie, namelijk de ontvangst van de koopprijs, resteert dus een duidelijke vermindering van het actief/verhaalsvermogen. Er zijn geen aanwijzingen dat schuldeisers zich zonder de verkoop niet op deze activa hadden kunnen verhalen, bijvoorbeeld wegens uitwinning door zekerheidsgerechtigden. De benadeling bestond bovendien op het moment waarop de curatoren hun rechten geldend maakten, omdat de activa toen door de verkoop buiten de boedel waren geraakt en Veldkamp afgifte weigerde.

\paragraph{Wetenschap - overwegingen}\label{wetenschap---overwegingen-67}

De verkoop door Trend aan Veldkamp was een rechtshandeling anders dan om niet, zodat wetenschap van benadeling bij zowel Trend als Veldkamp vereist is. Op basis van de gegeven feiten is die wetenschap niet zelfstandig en rechtstreeks bewezen: er zijn aanwijzingen van financiële problemen door het eerdere faillissement van Drent en de afhankelijkheid van Trend van Drent, maar geen concrete gegevens over de liquiditeits- of schuldenpositie van Trend of over specifieke kennis daarvan bij Veldkamp. Wel is het bewijsvermoeden van art. 43 lid 1 sub 1 Fw van toepassing: de totale koopprijs van € 45.000 voor meerdere Vision-machines en componenten staat, mede gelet op de latere biedingen van DG Press van € 60.000 voor alleen de Vision US 512 en € 95.000 voor de U 1063-componenten en zes druktorens, in een aanmerkelijke wanverhouding tot de waarde van de overgedragen activa. Er is geen aanwijzing dat Trend zich al vóór de jaartermijn tot deze verkoop had verplicht. Er zijn evenmin feiten gesteld waaruit volgt dat het vermoeden aan de zijde van Trend of Veldkamp met succes is ontzenuwd. Art. 45 Fw is niet aan de orde, omdat geen sprake is van een rechtshandeling om niet.

\newpage

\markboth{JOR 2011/20}{gemini-3.5-flash (k = 5)}

\subsubsection{Gemini}\label{gemini-22}

\texttt{gemini-3.5-flash}

\begin{center}
\small
\begin{tabular}{@{}p{7.4cm}p{2.6cm}p{2.6cm}@{}}
\toprule
Variabele & Waarde (model) & Waarde (gold) \\
\midrule
Is rechtshandeling & \texttt{true} & \texttt{true} \\
Datum rechtshandeling & \texttt{2009-08-19} & \texttt{2009-08-19} \\
Verricht door de schuldenaar & \texttt{true} & \texttt{true} \\
Om niet & \texttt{false} & \texttt{false} \\
Datum faillietverklaring & \texttt{2009-11-30} & \texttt{2009-11-30} \\
Onverplicht & \texttt{true} & \texttt{true} \\
Benadeling schuldeisers & \texttt{true} & \texttt{true} \\
Wetenschap schuldenaar aangetoond & \textcolor{gray}{\texttt{true}} & \texttt{false} \\
Wetenschap wederpartij aangetoond & \textcolor{gray}{\texttt{true}} & \texttt{false} \\
Vermoeden wetenschap schuldenaar tegenbewijs & \texttt{false} & \texttt{false} \\
Vermoeden wetenschap wederpartij tegenbewijs & \texttt{false} & \texttt{false} \\
Tegen onderwaarde & \texttt{true} & \texttt{true} \\
Voldoening of zekerheid niet opeisbare schuld & \texttt{null} & \texttt{null} \\
Gelieerde partij & \texttt{null} & \texttt{null} \\
Reeds verplicht voor aanvang termijn & \texttt{false} & \texttt{false} \\
\textbf{Vernietigbaar} & \textbf{\texttt{true}} & \textbf{\texttt{true}} \\
\bottomrule
\end{tabular}
\end{center}

\paragraph{Is rechtshandeling - overwegingen}\label{is-rechtshandeling---overwegingen-68}

De gewraakte handeling betreft de verkoop van diverse machines en machineonderdelen door Trend aan Veldkamp. Een koopovereenkomst is een meerzijdige rechtshandeling omdat deze is gericht op het scheppen van verbintenissen (rechtsgevolg).

\paragraph{Verricht door de schuldenaar - overwegingen}\label{verricht-door-de-schuldenaar---overwegingen-68}

De koopovereenkomst is gesloten tussen Trend Graphic Services B.V. (de latere failliet) en Veldkamp. De rechtshandeling is dus door de schuldenaar zelf verricht.

\paragraph{Om niet - overwegingen}\label{om-niet---overwegingen-68}

Er is een totale koopprijs van € 45.000,00 overeengekomen voor de machines en onderdelen. Omdat er een tegenprestatie tegenover de levering staat, is de rechtshandeling anders dan om niet (om baat).

\paragraph{Onverplicht - overwegingen}\label{onverplicht---overwegingen-68}

De aangevochten rechtshandeling betreft de verkoop en levering van diverse machines en machineonderdelen door Trend aan Veldkamp op of rond 19 augustus 2009. Een rechtshandeling is onverplicht wanneer er ten tijde van het verrichten ervan geen op de wet of een eerdere overeenkomst berustende rechtsplicht bestond die in rechte afdwingbaar was. Uit de casus blijkt dat de koopovereenkomst pas op of kort voor 19 augustus 2009 na onderhandelingen tot stand is gekomen (`Na een lange discussie (\ldots) zijn we het eens geworden'). Er bestond vóór het sluiten van deze overeenkomst geen enkele eerdere verplichting voor Trend om deze specifieke activa aan Veldkamp te verkopen of te leveren. Het aangaan van de koopovereenkomst en de daaruit voortvloeiende levering zijn derhalve onverplichte rechtshandelingen.

\paragraph{Benadeling schuldeisers - overwegingen}\label{benadeling-schuldeisers---overwegingen-68}

Om te bepalen of er sprake is van benadeling van schuldeisers in de zin van artikel 42 Faillissementswet (Fw), dient een vermogensvergelijking te worden gemaakt tussen de hypothetische situatie waarin de schuldeisers zouden hebben verkeerd zonder de gewraakte rechtshandeling en de feitelijke situatie waarin zij verkeren nu de handeling is verricht. In de feitelijke situatie zijn de machines en machineonderdelen door Trend aan Veldkamp verkocht en geleverd voor een totaalbedrag van € 45.000,-. In de hypothetische situatie waarin deze verkoop niet had plaatsgevonden, zouden de activa nog in de boedel van Trend zijn gevallen. Uit de feiten blijkt dat de curatoren na de beslaglegging concrete en reële biedingen hebben ontvangen van DG Press Services B.V. voor deze activa: een bod van € 60.000,- voor de Drent Vision US 512 en een bod van € 95.000,- voor de overige componenten en druktorens (in totaal € 155.000,-). Nu de werkelijke marktwaarde van de activa ten tijde van het geldend maken van de rechten door de curator aanzienlijk hoger bleek te zijn dan de overeengekomen koopprijs van € 45.000,-, heeft de rechtshandeling geleid tot een substantiële vermindering van het actief van de boedel. De schuldeisers zijn hierdoor rechtstreeks in hun verhaalsmogelijkheden benadeeld.

\paragraph{Wetenschap - overwegingen}\label{wetenschap---overwegingen-68}

De verkoop van de machines op 19 augustus 2009 vond plaats binnen een jaar voor het faillissement van Trend (30 november 2009). De overeengekomen koopprijs van € 45.000,- voor alle machines en onderdelen tezamen was aanmerkelijk lager dan de werkelijke waarde, hetgeen blijkt uit de latere concrete biedingen van € 60.000,- en € 50.000,- voor slechts delen van de activa. Hierdoor is het bewijsvermoeden van art. 43 lid 1 sub 1° Fw (onderwaarde) van toepassing. Dit vermoeden is niet ontzenuwd; de wetenschap van benadeling is aan beide zijden ook rechtstreeks aangetoond. Trend wist dat zij na het faillissement van haar zustervennootschap Drent (waarvan zij volledig afhankelijk was voor personeel en machines) in ernstige problemen verkeerde. Veldkamp, als nauw betrokken partner in de grafische industrie, wist of behoorde te weten dat deze transactie onder de marktprijs de schuldeisers van Trend zou benadelen en dat een faillissement met een tekort voorzienbaar was.

\newpage

\markboth{JOR 2011/20}{mistral/mistral-medium-3.5-128b:fp8 (k = 10)}

\subsection{k = 10}\label{k-10-7}

\subsubsection{Mistral}\label{mistral-23}

\texttt{mistral/mistral-medium-3.5-128b:fp8}

\begin{center}
\small
\begin{tabular}{@{}p{7.4cm}p{2.6cm}p{2.6cm}@{}}
\toprule
Variabele & Waarde (model) & Waarde (gold) \\
\midrule
Is rechtshandeling & \texttt{true} & \texttt{true} \\
Datum rechtshandeling & \texttt{2009-08-19} & \texttt{2009-08-19} \\
Verricht door de schuldenaar & \textcolor{gray}{\texttt{false}} & \texttt{true} \\
Om niet & \texttt{false} & \texttt{false} \\
Datum faillietverklaring & \textcolor{gray}{\texttt{2009-07-06}} & \texttt{2009-11-30} \\
Onverplicht & \texttt{true} & \texttt{true} \\
Benadeling schuldeisers & \texttt{true} & \texttt{true} \\
Wetenschap schuldenaar aangetoond & \texttt{false} & \texttt{false} \\
Wetenschap wederpartij aangetoond & \texttt{false} & \texttt{false} \\
Vermoeden wetenschap schuldenaar tegenbewijs & \texttt{null} & \texttt{false} \\
Vermoeden wetenschap wederpartij tegenbewijs & \texttt{null} & \texttt{false} \\
Tegen onderwaarde & \textcolor{gray}{\texttt{null}} & \texttt{true} \\
Voldoening of zekerheid niet opeisbare schuld & \texttt{null} & \texttt{null} \\
Gelieerde partij & \textcolor{gray}{\texttt{true}} & \texttt{null} \\
Reeds verplicht voor aanvang termijn & \texttt{null} & \texttt{false} \\
\textbf{Vernietigbaar} & \textbf{\textcolor{gray}{\texttt{false}}} & \textbf{\texttt{true}} \\
\bottomrule
\end{tabular}
\end{center}

\paragraph{Is rechtshandeling - overwegingen}\label{is-rechtshandeling---overwegingen-69}

De orderbevestiging (2.4) en factuur (2.5) van Trend aan Veldkamp vormen een koopovereenkomst, wat een rechtshandeling is (art. 3:33 BW).

\paragraph{Verricht door de schuldenaar - overwegingen}\label{verricht-door-de-schuldenaar---overwegingen-69}

De rechtshandeling (koopovereenkomst) is verricht door Trend, niet door Drent (de schuldenaar in het faillissement van 6 juli 2009).

\paragraph{Om niet - overwegingen}\label{om-niet---overwegingen-69}

De koopovereenkomst tussen Trend en Veldkamp bevat een koopprijs (€ 45.000,00, exclusief BTW), wat een wederprestatie inhoudt. Derhalve is de handeling `anders dan om niet' (om baat).

\paragraph{Onverplicht - overwegingen}\label{onverplicht---overwegingen-69}

De verkoop van de machines door Trend aan Veldkamp op 19 augustus 2009 is een onverplichte rechtshandeling in de zin van artikel 42 lid 1 Fw. Er bestond geen wettelijke of contractuele verplichting voor Trend om deze machines te verkopen. De machines waren afkomstig van inruil door klanten van Drent en werden door Trend opgeslagen in afwachting van doorverkoop. Er is geen sprake van een voorafgaande verplichting tot verkoop, noch van een wettelijke plicht om de machines te vervreemden. Morele, feitelijke of economische druk (zoals de financiële situatie van Trend of Drent) schept geen rechtsplicht. Beslissend is dat de verkoop niet kon worden afgedwongen in rechte. De stelplicht en bewijslast van het onverplichte karakter rusten op de curator, en in deze casus is voldoende aannemelijk dat de verkoop vrijwillig en zonder rechtsplicht heeft plaatsgevonden.

\paragraph{Benadeling schuldeisers - overwegingen}\label{benadeling-schuldeisers---overwegingen-69}

Op basis van de casus en de toepasselijke jurisprudentie (met name ECLI:NL:RBROT:2014:4807 en ECLI:NL:RBALM:2011:BU7299) kan worden geconcludeerd dat de schuldeisers van Trend door de verkoop van activa aan gedaagden op 19 augustus 2009 zijn benadeeld in de zin van artikel 42 Fw.

Allereerst is sprake van een onverplichte rechtshandeling: Trend was niet verplicht om de machines en componenten aan gedaagden te verkopen. Vervolgens dient te worden beoordeeld of de schuldeisers zijn benadeeld.

De benadeling moet worden vastgesteld door een vermogensvergelijking: de hypothetische situatie waarin de schuldeisers zouden hebben verkeerd zonder de gewraakte rechtshandeling wordt vergeleken met de situatie waarin zij feitelijk verkeren als die handeling in stand blijft.

Uit de casus blijkt dat de machines en componenten die door Trend aan gedaagden zijn verkocht, activa van Trend betroffen die in de boedel hadden moeten komen. Door de verkoop zijn deze activa aan de boedel onttrokken, waardoor de schuldeisers in hun verhaalsmogelijkheden zijn beperkt. Dit is een klassiek geval van vermindering van het actief, wat een vorm van benadeling is (vgl. ECLI:NL:RBROT:2014:4807, r.o. 4.4).

Daarnaast is er sprake van een verstoring van de rangorde tussen de schuldeisers, omdat de opbrengst van de verkoop aan één specifieke schuldeiser (gedaagden) ten goede is gekomen, in plaats van dat deze opbrengst ten behoeve van alle schuldeisers in de boedel zou zijn gekomen (vgl. ECLI:NL:RBALM:2011:BU7299, r.o. 3.4).

De omstandigheid dat de curatoren later biedingen hebben ontvangen voor de machines (bijv. € 60.000 voor de Drent Vision US 512 en € 95.000 voor andere componenten) toont aan dat de activa waarde hadden en dus een bate voor de boedel hadden kunnen vormen. Het feit dat deze activa zijn verkocht voor een lagere prijs (€ 45.000) en dat de opbrengst niet in de boedel is gekomen, bevestigt de benadeling.

Tot slot is er geen aanwijzing dat de schuldeisers ook zonder de handeling volledig zouden zijn voldaan, of dat de activa niet voor verhaalsdoeleinden hadden kunnen worden gebruikt. Derhalve is er sprake van benadeling in de zin van artikel 42 Fw.

\paragraph{Wetenschap - overwegingen}\label{wetenschap---overwegingen-69}

De transactie betreft de verkoop van activa door Trend (een zustervennootschap van Drent) aan Veldkamp op 19 augustus 2009, binnen een jaar voor de faillietverklaring van Drent op 6 juli 2009. Trend en Drent zijn gelieerde partijen in de zin van art. 43 lid 1 sub 3° Fw, omdat Drent een rechtspersoon is en Trend een zustervennootschap (sub 6°). Hierdoor wordt de wetenschap van benadeling aan beide zijden (Trend en Veldkamp) vermoed, tenzij tegenbewijs wordt geleverd. Er is geen sprake van een verplichting die voor de jaartermijn bestond, noch van een onderwaarde of voldoening van een niet-opeisbare schuld. De curatoren zullen moeten aantonen dat de transactie onverplicht was en dat benadeling van schuldeisers heeft plaatsgevonden. Het bewijsvermoeden van art. 43 Fw is van toepassing, maar kan worden weerlegd door Veldkamp.

\newpage

\markboth{JOR 2011/20}{gpt-5.5 (k = 10)}

\subsubsection{GPT}\label{gpt-23}

\texttt{gpt-5.5}

\begin{center}
\small
\begin{tabular}{@{}p{7.4cm}p{2.6cm}p{2.6cm}@{}}
\toprule
Variabele & Waarde (model) & Waarde (gold) \\
\midrule
Is rechtshandeling & \texttt{true} & \texttt{true} \\
Datum rechtshandeling & \texttt{2009-08-19} & \texttt{2009-08-19} \\
Verricht door de schuldenaar & \texttt{true} & \texttt{true} \\
Om niet & \texttt{false} & \texttt{false} \\
Datum faillietverklaring & \texttt{2009-11-30} & \texttt{2009-11-30} \\
Onverplicht & \texttt{true} & \texttt{true} \\
Benadeling schuldeisers & \texttt{true} & \texttt{true} \\
Wetenschap schuldenaar aangetoond & \texttt{false} & \texttt{false} \\
Wetenschap wederpartij aangetoond & \texttt{false} & \texttt{false} \\
Vermoeden wetenschap schuldenaar tegenbewijs & \texttt{false} & \texttt{false} \\
Vermoeden wetenschap wederpartij tegenbewijs & \texttt{false} & \texttt{false} \\
Tegen onderwaarde & \texttt{true} & \texttt{true} \\
Voldoening of zekerheid niet opeisbare schuld & \texttt{false} & \texttt{null} \\
Gelieerde partij & \texttt{false} & \texttt{null} \\
Reeds verplicht voor aanvang termijn & \texttt{false} & \texttt{false} \\
\textbf{Vernietigbaar} & \textbf{\texttt{true}} & \textbf{\texttt{true}} \\
\bottomrule
\end{tabular}
\end{center}

\paragraph{Is rechtshandeling - overwegingen}\label{is-rechtshandeling---overwegingen-70}

De gewraakte handeling betreft de verkoopovereenkomst, en kennelijk ook levering, van activa/machines en componenten van Trend aan Veldkamp/gedaagden. Een verkoop en levering zijn op rechtsgevolg gerichte handelingen.

\paragraph{Verricht door de schuldenaar - overwegingen}\label{verricht-door-de-schuldenaar---overwegingen-70}

De verkoop is verricht door Trend, de later failliet verklaarde schuldenaar: Trend stuurde op 19 augustus 2009 de orderbevestiging en factuur aan Veldkamp.

\paragraph{Om niet - overwegingen}\label{om-niet---overwegingen-70}

De activa zijn verkocht tegen een koopprijs: volgens de orderbevestiging bedroeg de totale prijs € 45.000 exclusief btw, met specificatie in de factuur. Er was dus een bedongen tegenprestatie.

\paragraph{Onverplicht - overwegingen}\label{onverplicht---overwegingen-70}

De aangevochten rechtshandeling betreft de verkoop en levering door Trend van de Drent Vision-machines en losse componenten aan (een van) gedaagden/Veldkamp op of omstreeks 19 augustus 2009. Uit de casus blijkt niet van een reeds vóór deze transactie bestaande, in rechte afdwingbare verplichting van Trend om juist deze activa aan gedaagden te verkopen of te leveren. De orderbevestiging en factuur van 19 augustus 2009 vormen de totstandkoming/uitvoering van de bestreden verkoop zelf en leveren geen vooraf bestaande rechtsplicht op. Dat Trend handelde in gebruikte machines en verkoop van voorraad tot haar normale bedrijfsactiviteiten behoorde, maakt het sluiten van deze concrete koopovereenkomst niet verplicht in de zin van art. 42 Fw. Evenmin blijkt van een wettelijke verplichting of van een eerdere overeenkomst die Trend tot deze verkoop noopte. Voor zover de levering uitvoering gaf aan de koopovereenkomst, neemt dat niet weg dat de koopovereenkomst als dragende rechtshandeling onverplicht is aangegaan; vergelijkbaar wordt verkoop gevolgd door verrekening in de rechtspraak als onverplichte rechtshandeling van de schuldenaar aangemerkt, indien geen voorafgaande verkoopplicht bestond (vgl. o.a. ECLI:NL:RBALM:2011:BU7299 en ECLI:NL:GHDHA:2016:3475). De rechtshandeling is daarom onverplicht verricht.

\paragraph{Benadeling schuldeisers - overwegingen}\label{benadeling-schuldeisers---overwegingen-70}

Bij de vermogensvergelijking moet de feitelijke situatie waarin de verkoop aan Veldkamp/gedaagden in stand blijft worden vergeleken met de hypothetische situatie zonder die verkoop. In de feitelijke situatie is een omvangrijk pakket machines, druktorens en componenten uit het vermogen van Trend verdwenen tegen een totale koopprijs van € 45.000 exclusief btw. Uit de latere concrete biedingen van DG Press blijkt dat ten minste een deel van diezelfde activa aanzienlijk meer kon opbrengen: alleen al voor de Drent Vision US 512 met componenten is € 60.000 geboden, terwijl daarnaast nog biedingen zijn gedaan van € 50.000 voor componenten van de Drent Vision U 1063 en € 45.000 voor zes druktorens. Ook als rekening wordt gehouden met de voorwaarden bij die biedingen, volgt hieruit dat de verkochte activa een substantieel hogere realiseerbare waarde hadden dan de overeengekomen totale prijs. Zonder de gewraakte verkoop zouden deze activa in de boedel zijn gebleven en hadden de curatoren deze ten behoeve van de gezamenlijke schuldeisers kunnen gelde maken. Door de verkoop is het actief van Trend dus verminderd althans is verhaal voor schuldeisers beperkt. Er zijn geen aanwijzingen dat schuldeisers zich ook zonder de verkoop niet op deze activa hadden kunnen verhalen of dat de transactie per saldo een gelijkwaardig voordeel voor de boedel opleverde. De schuldeisers zijn daarom benadeeld in de zin van artikel 42 Fw.

\paragraph{Wetenschap - overwegingen}\label{wetenschap---overwegingen-70}

De verkoop was een rechtshandeling anders dan om niet. Rechtstreekse wetenschap van benadeling is op basis van de aangeleverde feiten niet zelfstandig bewezen: er zijn wel sterke aanwijzingen door het eerdere faillissement van Drent, de afhankelijkheid van Trend van Drent en de korte periode tot het faillissement van Trend, maar zonder het bewijsvermoeden is met name aan de zijde van Veldkamp/gedaagden niet concreet aangetoond dat het faillissement van Trend en een tekort daarin met redelijke mate van waarschijnlijkheid voorzienbaar waren. Wel is het bewijsvermoeden van art. 43 lid 1 sub 1 Fw toepasselijk: de verkochte machines en componenten zijn voor € 45.000 verkocht, terwijl latere concrete biedingen voor slechts delen daarvan (€ 60.000 voor de US 512 en € 95.000 voor onderdelen/torens) op een aanmerkelijk hogere waarde wijzen. Er is geen aanwijzing dat Trend zich al vóór de jaartermijn tot deze verkoop had verplicht. Van gelieerdheid of voldoening/zekerheid voor een niet-opeisbare schuld blijkt niet. Nu geen feiten zijn gesteld die het vermoeden ontzenuwen, wordt de wetenschap van benadeling aan beide zijden vermoed en is tegenbewijs niet geslaagd. Art. 45 Fw is niet aan de orde omdat geen rechtshandeling om niet voorligt.

\newpage

\markboth{JOR 2011/20}{gemini-3.5-flash (k = 10)}

\subsubsection{Gemini}\label{gemini-23}

\texttt{gemini-3.5-flash}

\begin{center}
\small
\begin{tabular}{@{}p{7.4cm}p{2.6cm}p{2.6cm}@{}}
\toprule
Variabele & Waarde (model) & Waarde (gold) \\
\midrule
Is rechtshandeling & \texttt{true} & \texttt{true} \\
Datum rechtshandeling & \texttt{2009-08-19} & \texttt{2009-08-19} \\
Verricht door de schuldenaar & \texttt{true} & \texttt{true} \\
Om niet & \texttt{false} & \texttt{false} \\
Datum faillietverklaring & \texttt{2009-11-30} & \texttt{2009-11-30} \\
Onverplicht & \texttt{true} & \texttt{true} \\
Benadeling schuldeisers & \texttt{true} & \texttt{true} \\
Wetenschap schuldenaar aangetoond & \texttt{false} & \texttt{false} \\
Wetenschap wederpartij aangetoond & \texttt{false} & \texttt{false} \\
Vermoeden wetenschap schuldenaar tegenbewijs & \texttt{false} & \texttt{false} \\
Vermoeden wetenschap wederpartij tegenbewijs & \texttt{false} & \texttt{false} \\
Tegen onderwaarde & \texttt{true} & \texttt{true} \\
Voldoening of zekerheid niet opeisbare schuld & \texttt{null} & \texttt{null} \\
Gelieerde partij & \texttt{null} & \texttt{null} \\
Reeds verplicht voor aanvang termijn & \texttt{false} & \texttt{false} \\
\textbf{Vernietigbaar} & \textbf{\texttt{true}} & \textbf{\texttt{true}} \\
\bottomrule
\end{tabular}
\end{center}

\paragraph{Is rechtshandeling - overwegingen}\label{is-rechtshandeling---overwegingen-71}

De gewraakte handeling betreft de verkoop van diverse machines en machineonderdelen door Trend aan Veldkamp. Een koopovereenkomst is een meerzijdige rechtshandeling omdat deze is gericht op het scheppen van verbintenissen (rechtsgevolg).

\paragraph{Verricht door de schuldenaar - overwegingen}\label{verricht-door-de-schuldenaar---overwegingen-71}

De koopovereenkomst is gesloten tussen Trend Graphic Services B.V. (de latere failliet) en Veldkamp. De rechtshandeling is dus door de schuldenaar zelf verricht.

\paragraph{Om niet - overwegingen}\label{om-niet---overwegingen-71}

Er is een totale koopprijs van € 45.000,00 overeengekomen voor de machines en onderdelen. Omdat er een tegenprestatie tegenover de levering staat, is de rechtshandeling anders dan om niet (om baat).

\paragraph{Onverplicht - overwegingen}\label{onverplicht---overwegingen-71}

De aangevochten rechtshandeling betreft de verkoop en levering van diverse machines en machineonderdelen door Trend aan Veldkamp op 19 augustus 2009. Een rechtshandeling is onverplicht wanneer daartoe geen op de wet of een voorafgaande overeenkomst berustende rechtsplicht bestond die in rechte kon worden afgedwongen. Het aangaan van een nieuwe koopovereenkomst en de daaropvolgende levering zijn in beginsel onverplichte rechtshandelingen. Uit de casus blijkt niet van een reeds bestaande wettelijke of contractuele verplichting voor Trend om deze specifieke activa aan Veldkamp te verkopen en te leveren; de transactie is het resultaat van vrije onderhandelingen die op 19 augustus 2009 zijn geformaliseerd. De rechtshandeling is dan ook onverplicht verricht.

\paragraph{Benadeling schuldeisers - overwegingen}\label{benadeling-schuldeisers---overwegingen-71}

De beoordeling van benadeling van schuldeisers in de zin van artikel 42 Fw vereist een vergelijking tussen de hypothetische situatie waarin de schuldeisers zonder de rechtshandeling zouden hebben verkeerd en de feitelijke situatie waarin zij verkeren als de handeling in stand blijft. In deze casus heeft Trend op 19 augustus 2009 diverse machines en onderdelen verkocht aan Veldkamp voor een totaalbedrag van € 45.000,-. Kort na het faillissement hebben de curatoren echter concrete en reële biedingen ontvangen van DG Press voor dezelfde activa, waaronder een bod van € 60.000,- voor de Drent Vision US 512 en een bod van € 50.000,- voor de componenten van de Drent Vision U 1063. De totale marktwaarde van de activa lag derhalve aanzienlijk hoger dan de overeengekomen koopprijs van € 45.000,-. Zonder de gewraakte rechtshandeling hadden de curatoren deze activa voor een veel hogere opbrengst kunnen verkopen ten behoeve van de boedel. Door de verkoop tegen een te lage prijs zijn de activa aan het verhaal van de gezamenlijke schuldeisers onttrokken, wat heeft geleid tot een rechtstreekse vermindering van het actieve boedelvermogen. De schuldeisers zijn hierdoor in hun verhaalsmogelijkheden benadeeld.

\paragraph{Wetenschap - overwegingen}\label{wetenschap---overwegingen-71}

De verkoop van de machines en onderdelen door Trend aan Veldkamp vond plaats op 19 augustus 2009, binnen een jaar voor de faillietverklaring van Trend op 30 november 2009. De transactie was onverplicht en geschiedde tegen een koopprijs van € 45.000,-, terwijl de werkelijke waarde (blijkens latere biedingen van € 60.000,- en € 95.000,-) aanmerkelijk hoger lag. Hierdoor is sprake van een transactie onder de marktwaarde, wat het bewijsvermoeden van art. 43 lid 1 sub 1° Fw activeert. De wetenschap van benadeling wordt derhalve aan beide zijden vermoed. Er is geen sprake van een succesvol geleverd tegenbewijs om dit vermoeden te ontzenuwen, noch is de wetenschap los van dit vermoeden rechtstreeks aangetoond.

\newpage

\markboth{JOR 2012/196}{Gold standard}

\section{JOR 2012/196}\label{jor-2012196}

uitspraakdatum 15-03-2011

\subsection{Gold}\label{gold-8}

\begin{center}
\small
\begin{tabular}{@{}p{8cm}p{2.6cm}@{}}
\toprule
Variabele & Waarde \\
\midrule
Is rechtshandeling & \texttt{true} \\
Datum rechtshandeling & \texttt{2006-09-05} \\
Verricht door de schuldenaar & \texttt{true} \\
Om niet & \texttt{false} \\
Datum faillietverklaring & \texttt{2006-12-13} \\
Onverplicht & \texttt{false} \\
Benadeling schuldeisers & \texttt{false} \\
Wetenschap schuldenaar aangetoond & \texttt{false} \\
Wetenschap wederpartij aangetoond & \texttt{null} \\
Vermoeden wetenschap schuldenaar tegenbewijs & \texttt{null} \\
Vermoeden wetenschap wederpartij tegenbewijs & \texttt{null} \\
Tegen onderwaarde & \texttt{null} \\
Voldoening of zekerheid niet opeisbare schuld & \texttt{null} \\
Gelieerde partij & \texttt{null} \\
Reeds verplicht voor aanvang termijn & \texttt{null} \\
\textbf{Vernietigbaar} & \textbf{\texttt{false}} \\
\bottomrule
\end{tabular}
\end{center}

\subsubsection{Uitspraak}\label{uitspraak-8}

(\ldots; red.)

\paragraph{4. De beoordeling}\label{de-beoordeling-4}

4.1. Het gaat in dit hoger beroep om het volgende.

\begin{enumerate}
\def\labelenumi{\alph{enumi}.}
\item
  Op 13 december 2006 is Nutriscience B.V. (hierna: de failliet) op eigen verzoek in staat van faillissement verklaard, nadat de Algemene Vergadering van Aandeelhouders daartoe op 11 december 2006 had besloten. De curator is als zodanig aangesteld. Tot het faillissement was Thomassen de enige statutaire bestuurder van de failliet. Universiteit Maastricht Holding B.V. (hierna: UMH) was eigenaar van 100\% van de aandelen in de failliet.
\item
  Op 15 mei 2006 heeft Thomassen met de failliet een schriftelijk vastgelegde overeenkomst van geldlening gesloten welke, voor zover thans van belang, het volgende inhoudt:
\end{enumerate}

``IN OVERWEGING NEMENDE:

-- De schuldenaar (de failliet; toev. hof) heeft behoefte aan een kortlopend overbruggingskrediet van € 70.000,- voor de periode 25 april tot en met uiterlijk 31 december;

-- De Schuldeiser (Thomassen; toev. hof) bereid is om dit (\ldots) tijdelijk te verstrekken.

-- De Schuldeiser tevens directeur is van NutriScience (de Schuldernaar). Daarom wordt de overeenkomst mede ondertekend door de aandeelhouder.

Artikel 1

De schuldeiser leent per 16 mei 2006 aan de schuldenaar € 50.000,00 (\ldots) en per 21 juni 2006 (\ldots) € 20.000,00 (\ldots) in totaal € 70.000,- (\ldots).

Artikel 2

De op de geldlening verschuldigde rente bedraagt 6\% per jaar (\ldots)

Artikel 3

Deze geldlening wordt aangegaan voor de periode van 16 mei 2006 tot uiterlijk 31 december 2006. Voor of uiterlijk op die datum zal de integrale hoofdsom van de lening door de schuldenaar zijn/worden afgelost.

(\ldots)

Artikel 5

Vervroegde, volledige of gedeeltelijke terugbetaling is steeds toegestaan.(\ldots)

Artikel 9

De geldlening is van de zijde van de schuldeiser opeisbaar (\ldots) in de navolgende gevallen: indien en zodra de schuldenaar overlijdt, in staat van faillissement wordt verklaard, hem surséance van betaling of schuldsanering (\ldots) wordt verleend, hij onder curatele of onder bewind wordt gesteld, hij zijn bedrijf staakt of (\ldots) overdraagt (\ldots)''

\begin{enumerate}
\def\labelenumi{\alph{enumi}.}
\setcounter{enumi}{2}
\item
  In augustus 2006 is tussen Thomassen en de failliet wederom een overeenkomst van geldlening gesloten, thans voor een bedrag van € 20.000,--. Deze overeenkomst is niet schriftelijk vastgelegd.
\item
  Op 18 mei 2006 heeft de failliet van Thomassen ontvangen een bedrag van € 50.000,--, op 23 juni 2006 een bedrag van € 20.000,-- en op 30 augustus 2006 wederom een bedrag van € 20.000,--.
\item
  Op 5 september 2006 heeft de failliet aan Thomassen terugbetaald het onder c) genoemde bedrag van € 20.000,--. Op 14 oktober 2006 respectievelijk 24 oktober 2006 zijn aan Thomassen bedragen terugbetaald van € 51.250,-- respectievelijk € 20.400,--. In totaal is aan Thomassen derhalve terugbetaald € 91.650,--.
\item
  Bij brieven van 30 november 2006 aan de fiscus, het UWV en het ABP heeft Thomassen namens de failliet betalingsonmacht gemeld ter zake van door de failliet verschuldigde loonbelasting en pensioenpremies over de maand oktober 2006.
\item
  De failliet had tien werknemers. Deze werknemers hebben hun loon over de maand november 2006 (in totaal ongeveer € 30.000,--) niet ontvangen.
\item
  De notulen van de Raad van Commissarissen (hierna: RvC) van de failliet van 9 november 2006 behelzen onder meer het volgende:
\end{enumerate}

``(\ldots) De realisatie voor 2006 geeft een negatieve winst voor belastingen van € 50.000,-- Het worst case scenario geeft -/- € 125.000,--. (\ldots) Drie klanten worden nu afgewerkt. Lopende projecten kunnen hopelijk wel gewoon afgerond worden. Geprobeerd moet worden om nieuwe klanten binnen te halen. Verder moeten er nieuwe investeerders gezocht worden. (\ldots) Er komt nog wel geld binnen in 2007. Als er geen projecten meer binnenkomen en lopende projecten worden afgewerkt, dan wordt in 2007 het laagste resultaat verwacht, namelijk -/- € 350.000,--.

(\ldots)

UMH wil zeker meewerken aan een reddingsscenario, maar stelt als voorwaarde dat er een substantiële financiering uit de markt gehaald wordt. (\ldots) UMH kan bijvoorbeeld een debiteurenfinanciering bieden als nieuwe grote klanten binnengehaald worden die pas over een aantal maanden gaan betalen. Ook wil men meewerken aan het uitgeven van nieuwe aandelen aan nieuwe investeerders.(\ldots)

Het gaat nu om de korte termijn. (\ldots) Als de salarissen niet betaald kunnen worden, moet dat gemeld worden bij de fiscus. Dat moet nu gebeuren. (\ldots) Externe financiering en nieuwe manier van werken moeten tegelijk gebeuren. (\ldots) Met het verhaal van vandaag moeten investeerders worden benaderd. (\ldots)''

4.2. De curator heeft Thomassen in rechte betrokken en gevorderd 1) voor recht te verklaren dat bij buitengerechtelijke verklaring van 8 maart 2007 de nietigheid van de rechtshandelingen verricht op 5 september 2006, 4 oktober 2006 en 14 oktober 2006, door de curator terecht is ingeroepen en 2) Thomassen te veroordelen tot betaling van € 95.455,--, althans € 93.438,-- te vermeerderen met rente en kosten.

De curator heeft zich in eerste aanleg beroepen op art. 42 Fw. en aan de vorderingen ten grondslag gelegd dat de terugbetalingen door de failliet van de van Thomassen geleende gelden op voormelde data onverplicht waren en dat van die terugbetalingen, naar Thomassen en de failliet wisten, benadeling van de overige crediteuren van de failliet het gevolg was. De terugbetalingen waren paulianeus en derhalve vernietigbaar, zodat de betaalde bedragen als onverschuldigd betaald door de curator kunnen worden teruggevorderd. Voorts heeft de curator zich op het standpunt gesteld dat de terugbetalingen door de failliet aan Thomassen selectieve betalingen aan een crediteur waren, zodat Thomassen, als bestuurder van de failliet, onrechtmatig heeft gehandeld door die betalingen door de failliet te doen verrichten.

4.3. Na door Thomassen gevoerd verweer heeft de rechtbank de vorderingen afgewezen. De rechtbank overweegt daartoe dat de terugbetalingen door de failliet niet onverplicht waren omdat (r.o. 3.2.) uit de aanhef van de overeenkomst van geldlening blijkt dat het doel van de lening was een overbruggingskrediet aan de failliet te verschaffen tot uiterlijk 31 december 2006 en dat een redelijke uitleg van die overeenkomst meebrengt dat de termijn van de lening gekoppeld was aan de periode waarvoor het krediet nodig was ter overbrugging van een gebrek aan liquiditeiten. Volgens de rechtbank ontstond daarom voor de failliet de rechtsplicht tot terugbetaling zodra de behoefte aan krediet ter overbrugging van liquiditeitskrapte niet langer aanwezig was.

Voorts overweegt de rechtbank (r.o. 3.3.) dat voormelde uitleg van de overeenkomst wordt ondersteund door eerdere overbruggingskredieten die Thomassen aan de failliet heeft verstrekt. Het gaat daarbij om een bedrag van € 80.000,-- op 25 april 2005 (geleend tot uiterlijk 31 december 2005) en een bedrag van € 45.000,-- op 19 september 2005 (geleend tot uiterlijk 30 september 2005). Volgens de rechtbank blijkt uit productie 6 bij conclusie van antwoord, kort gezegd, dat deze geleende bedragen, kort nadat (het hof begrijpt:) een debiteur van de failliet aan de failliet een betaling had verricht, aan Thomassen zijn terugbetaald. Ook blijkt uit die overzichten dat de terugbetalingen waarom het in deze procedure gaat steeds nagenoeg zijn samengevallen met betalingen aan de failliet door haar debiteuren, aldus de rechtbank.

4.4. Met de grieven 6 en 7 (zo stelt Thomassen naar het hof begrijpt) wijzigt de curator de grondslag van de vorderingen. Thomassen maakt daartegen bezwaar omdat hem daardoor een instantie wordt ontnomen. Het hof overweegt als volgt.

De curator heeft in eerste aanleg (inl. dagv. sub 28) de vorderingen mede gegrond op een onrechtmatige daad van Thomassen als bestuurder van de failliet. In zoverre is dus geen sprake van een wijziging van (de grondslag van) de eis.

Voor het overige heeft te gelden dat het verlies van een instantie inherent is aan het feit dat de wet toestaat dat een eis ook in hoger beroep kan worden gewijzigd. Slechts onder bijkomende omstandigheden kan dit feit het oordeel rechtvaardigen dat sprake is van strijd met de eisen van een goede procesorde. Dergelijke omstandigheden zijn echter niet gesteld of gebleken, zodat het bezwaar van Thomassen ongegrond is.

4.5. Na wijziging van de eis legt de curator primair art. 42 Fw. (juncto art. 43 Fw), subsidiair art. 47 Fw en meer subsidiair een door Thomassen gepleegde onrechtmatige daad (art. 6:162 BW) aan de vorderingen ten grondslag.

4.6. Op grond van art. 42 Fw. kan de curator ten behoeve van de boedel iedere rechtshandeling van de failliet die deze vóór de faillietverklaring onverplicht heeft verricht en waarvan deze bij dit verrichten wist of behoorde te weten dat daarvan benadeling van de schuldeisers het gevolg zou zijn, door een buitengerechtelijke verklaring vernietigen.

De vraag moet derhalve beantwoord worden of de terugbetalingen door de failliet op 5 september 2006, 14 oktober 2006 en 24 oktober 2006 onverplicht verrichte rechtshandelingen zijn. Dit is het geval als er niet een op wet of overeenkomst berustende verplichting tot terugbetaling bestaat, waarbij heeft te gelden dat betaling van een niet-opeisbare schuld steeds onverplicht is. Een verbintenis is voorts op grond van art. 6:38 BW terstond opeisbaar indien geen tijd voor de nakoming is bepaald.

4.6.1. De overeenkomst op grond waarvan terugbetaling van € 20.000,-- op 5 september 2006 heeft plaats gevonden is niet op schrift gesteld. Uit enig schriftelijk stuk is dus niet af te leiden dat een tijd voor de nakoming van de verbintenis tot terugbetaling is bepaald. Feiten die dit anders zouden maken heeft de curator niet gesteld. Dit betekent dat de terugbetaling door de failliet van 5 september 2006 gelet op art. 6:38 BW niet onverplicht geschiedde.

4.6.2. De terugbetalingen van (afgerond) € 50.000,-- en € 20.000,-- op respectievelijk 14 en 24 oktober 2006 zijn verricht op grond van de overeenkomst van geldlening die de onder 4.1. b) weergegeven bepalingen behelst. In geschil is of partijen bij die overeenkomst met die bepalingen bedoeld hebben een tijd voor de nakoming van de verbintenis tot terugbetaling te bepalen als bedoeld in art. 6:38 BW. De curator stelt dat de vordering tot terugbetaling op grond van artikel 3 van de overeenkomst pas op 31 december 2006 opeisbaar was, zodat de eerdere betalingen door de failliet aan Thomassen onverplicht waren. Volgens Thomassen was sprake van een overbruggingskrediet dat onmiddellijk moest worden afgelost zodra de failliet daartoe de middelen had: Thomassen schoot bedragen voor die later door (het hof begrijpt:) debiteuren aan de failliet zouden worden voldaan. Aldus kon de failliet eerder over die bedragen beschikken dan wanneer zij op betaling door de debiteuren zou moeten wachten. Zodra de debiteur betaalde werd de vordering van Thomassen op de failliet opeisbaar, zo begrijpt het hof het standpunt van Thomassen.

4.6.3. Het hof overweegt als volgt. De tekst van artikel 3 van de overeenkomst luidt dat ``voor of uiterlijk op'' 31 december 2006 de hoofdsom moet worden afgelost.

De betekenis van een omstreden beding in een schriftelijke overeenkomst moet door de rechter worden vastgesteld aan de hand van hetgeen partijen over en weer hebben verklaard en uit elkaars verklaringen en gedragingen overeenkomstig de zin die zij daaraan in de gegeven omstandigheden redelijkerwijs mochten toekennen, hebben afgeleid en van hetgeen zij te dien aanzien redelijkerwijs van elkaar mochten verwachten. Uit een en ander volgt dat redelijkheid en billijkheid hierbij een rol spelen (vgl. HR 12 januari 2001, NJ 2001,199).

4.6.4. Het enkele feit dat partijen van mening verschillen over de betekenis van artikel 3 van de overeenkomst brengt met zich dat die overeenkomst conform voormeld criterium dient te worden uitgelegd. Grief 2, waarmee de curator betoogt dat de rechtbank de schriftelijke overeenkomst ten onrechte heeft uitgelegd, nu de inhoud daarvan duidelijk is, faalt om die reden.

4.6.5. Het hof is met de rechtbank van oordeel dat de overeenkomst aldus moet worden uitgelegd dat de vorderingen van Thomassen tot terugbetaling opeisbaar waren zodra de failliet de middelen daartoe voorhanden had omdat zij betaling door haar debiteuren had verkregen. Weliswaar overweegt de rechtbank (r.o.3.2.) dat de rechtsplicht tot terugbetaling ontstond ``zodra de behoefte aan krediet ter overbrugging van liquiditeitskrapte niet langer aanwezig was'', maar deze overweging dient begrepen te worden zoals hiervoor weergegeven, gelet op het door Thomassen gevoerde verweer (cva, blz 6, vijfde alinea) en de verklaring van Thomassen ter comparitie in eerste aanleg. Deze houdt het volgende in: ``In de geldleningovereenkomsten zijn geen specifieke data opgenomen waarop de lening terugbetaald diende te worden omdat ik niet wist wanneer het geld van debiteuren binnen zou komen waarmee ik dan weer terugbetaald zou worden.''

In hoger beroep (mva blz 6, eerste alinea en blz 12, laatste alinea) heeft Thomassen dit standpunt herhaald.

4.6.6. Het hof baseert zijn uitleg van het omstreden beding in de overeenkomst evenals de rechtbank op de gang van zaken rond de eerdere overeenkomsten van geldlening die Thomassen en de failliet in 2005 sloten en op de door de curator ook in hoger beroep niet weersproken inhoud van productie 6 bij conclusie van antwoord. Uit die productie blijkt dat de terugbetalingen door de failliet steeds (nagenoeg) samenvielen met betalingen door debiteuren aan de failliet. Dat zulks anders dan op grond van een verplichting daartoe zou zijn gebeurd, is gesteld noch gebleken. Het feit dat in 2005 bij de failliet van een geheel andere liquiditeitspositie sprake was dan in 2006 (zoals de curator in de toelichting op grief 5 II stelt) doet, indien al juist, er niet aan af dat de gang van zaken rond de geldleningen in 2005 naar het oordeel van het hof medebepalend is voor de uitleg van de leningsovereenkomst van 15 mei 2006.Uit die gang van zaken leidt het hof een aanwijzing af voor de opeisbaarheid van de schuld op het door Thomassen gestelde moment.

4.6.7. Bij zijn uitleg betrekt het hof voorts dat de redactie van de overeenkomst van 15 mei 2006, die is opgesteld door de UMH, de aandeelhouder van de failliet, onweersproken gebrekkig is geweest. Het beroep dat de curator op de artikelen 3 en 9 van de overeenkomst doet, ter ondersteuning van zijn betoog dat de vordering van Thomassen op de failliet slechts in bepaalde situaties vóór 31 december 2006 opeisbaar was, heeft om die reden geen doorslaggevende betekenis.

4.6.8. Uit het voorgaande vloeit voort dat de terugbetalingen door de failliet aan Thomassen niet onverplicht geschiedden en dat de grieven 1, 3, 5 I en 5 II falen. Ook leidt het voorgaande tot de conclusie dat het oordeel van de rechtbank, zoals door het hof in r.o. 4.6.5. hiervoor begrepen, voldoende gemotiveerd is en voldoende inzicht in haar gedachtegang geeft, zodat ook grief 4 faalt. Voor zover dat laatste al anders zou zijn baat die grief de curator, gelet op de door het hof aan de overeenkomst gegeven uitleg, niet.

4.7. Subsidiair legt de curator artikel 47 Fw. aan de vordering ten grondslag. Op grond van dit artikel kan, voor zover thans relevant, de voldoening door de schuldenaar van een opeisbare schuld worden vernietigd wanneer de betaling het gevolg was van overleg tussen de schuldenaar en de schuldeiser dat ten doel had laatstgenoemde door die betaling boven andere schuldeisers te begunstigen. Voor de aanwezigheid van overleg als in het artikel bedoeld, dat op grond van de wetsgeschiedenis niet spoedig mag worden aangenomen, is vereist dat sprake was van samenspanning, dat wil zeggen dat niet alleen bij de schuldeiser (Thomassen) maar ook bij de schuldenaar (de failliet) het oogmerk heeft voorgezeten door de gewraakte betaling deze schuldeiser boven andere te begunstigen.

4.7.1. Het hof is van oordeel dat de curator het door voormeld artikel vereiste overleg en het eveneens vereiste oogmerk onvoldoende heeft onderbouwd. Ook in 2005 zijn door de failliet aan Thomassen terugbetalingen van geleende bedragen gedaan zodra een debiteur van de failliet aan laatstgenoemde een betaling had verricht. Bovendien dateren de terugbetalingen van 5 september, 14 en 24 oktober 2006, terwijl op (en gedurende enige tijd na) 6 november 2006 blijkens de notulen van de vergadering van de RvC en de (onvoldoende) weersproken stellingen van Thomassen, er nog zicht was op de ontvangst van een drietal opdrachten door de failliet en er voorts getracht is het voortbestaan van de failliet te verzekeren door middel van een Management Buy-In proposal (door het personeel van de failliet of Thomassen zelf) gedateerd 21 november 2006 (prod. 8 cva). Toen de opdrachten niet kwamen en dit voorstel niet door de aandeelhouder werd geaccepteerd, is op 11 december 2006 het besluit genomen het faillissement aan te vragen. Dit feitencomplex duidt niet op, kort gezegd, samenspanning vóór of op de data van terugbetaling van de geleende gelden aan Thomassen, tussen de failliet en Thomassen met het oogmerk Thomassen als crediteur boven andere crediteuren te begunstigen. Op de verschillende momenten van terugbetaling was immers nog geen sprake van beëindiging van de onderneming van de failliet, ondanks de slechte financiële situatie van laatstgenoemde. Hieraan doet onder die omstandigheden niet af het enkele feit dat Thomassen als bestuurder van de failliet besloot tot terugbetaling aan zichzelf als geldverstrekker. Grief 6 faalt daarom.

4.8. Met grief 7 legt de curator, meer subsidiair, aan de vorderingen ten grondslag dat Thomassen als bestuurder van de failliet onrechtmatig heeft gehandeld door op de genoemde data aan zichzelf de aan de failliet geleende gelden terug te doen betalen, in plaats van dat bedrag aan te wenden ter gelijke verdeling onder alle andere crediteuren. Deze crediteuren hadden in totaal op 30 september 2006 een bedrag van € 328.346,-- van de failliet te vorderen, aldus de curator. Volgens de curator heeft Thomassen als bestuurder van de failliet selectief aan zichzelf betaald, waardoor de overige crediteuren schade hebben geleden.

4.8.1. Het beroep dat de curator doet op de uitspraak van de HR van 12 juni 1998 (JOR 1998, 107 (m.nt. Van den Ingh (Coral/Salt); red.) faalt reeds omdat in het onderhavige geval, anders dan in het geval dat in voormelde uitspraak aan de orde was, geenszins vast staat dat de failliet ten tijde van de gewraakte terugbetalingen aan Thomassen, besloten had haar activiteiten te beëindigen. Het hof verwijst naar hetgeen in r.o. 4.7.1. is overwogen. Onder de daar genoemde omstandigheden vormt de door de curator gestelde selectieve betaling geen onrechtmatige daad van Thomassen als bestuurder van de failliet. De enkele selectieve betaling door een vennootschap aan een crediteur is immers niet onrechtmatig jegens de overige crediteuren, ook al zou de vennootschap op het moment van die betaling niet tot betaling van al haar crediteuren in staat zijn. Feiten die dit oordeel anders zouden doen uitvallen heeft de curator niet gesteld. Ook grief 7 faalt.

4.9. Het door de curator gedane bewijsaanbod is onvoldoende gespecificeerd. Het hof passeert daarom dat aanbod.

4.10. De grieven 8 en 9 hebben, anders dan de curator in de toelichting op grief 9 lijkt te veronderstellen, geen zelfstandige betekenis. Deze grieven delen het lot van de overige grieven.

4.11. De slotsom luidt dat, nu alle grieven falen, het vonnis waarvan beroep zal worden bekrachtigd. De curator zal in de proceskosten van het hoger beroep veroordeeld worden. De proceskosten veroordeling zal op verzoek van Thomassen uitvoerbaar bij voorraad worden verklaard.

\paragraph{5. De uitspraak}\label{de-uitspraak-1}

Het hof:

bekrachtigt het vonnis waarvan beroep;

veroordeelt de curator in de proceskosten van het hoger beroep (\ldots; red.);

verklaart de proceskostenveroordeling uitvoerbaar bij voorraad;

wijst af het meer of anders gevorderde.

\newpage

\markboth{JOR 2012/196}{mistral/mistral-medium-3.5-128b:fp8 (k = 0)}

\subsection{k = 0}\label{k-0-8}

\subsubsection{Mistral}\label{mistral-24}

\texttt{mistral/mistral-medium-3.5-128b:fp8}

\begin{center}
\small
\begin{tabular}{@{}p{7.4cm}p{2.6cm}p{2.6cm}@{}}
\toprule
Variabele & Waarde (model) & Waarde (gold) \\
\midrule
Is rechtshandeling & \texttt{true} & \texttt{true} \\
Datum rechtshandeling & \texttt{2006-09-05} & \texttt{2006-09-05} \\
Verricht door de schuldenaar & \texttt{true} & \texttt{true} \\
Om niet & \texttt{false} & \texttt{false} \\
Datum faillietverklaring & \texttt{2006-12-13} & \texttt{2006-12-13} \\
Onverplicht & \textcolor{gray}{\texttt{true}} & \texttt{false} \\
Benadeling schuldeisers & \textcolor{gray}{\texttt{true}} & \texttt{false} \\
Wetenschap schuldenaar aangetoond & \textcolor{gray}{\texttt{true}} & \texttt{false} \\
Wetenschap wederpartij aangetoond & \textcolor{gray}{\texttt{true}} & \texttt{null} \\
Vermoeden wetenschap schuldenaar tegenbewijs & \texttt{null} & \texttt{null} \\
Vermoeden wetenschap wederpartij tegenbewijs & \texttt{null} & \texttt{null} \\
Tegen onderwaarde & \texttt{null} & \texttt{null} \\
Voldoening of zekerheid niet opeisbare schuld & \texttt{null} & \texttt{null} \\
Gelieerde partij & \textcolor{gray}{\texttt{true}} & \texttt{null} \\
Reeds verplicht voor aanvang termijn & \textcolor{gray}{\texttt{true}} & \texttt{null} \\
\textbf{Vernietigbaar} & \textbf{\textcolor{gray}{\texttt{true}}} & \textbf{\texttt{false}} \\
\bottomrule
\end{tabular}
\end{center}

\paragraph{Is rechtshandeling - overwegingen}\label{is-rechtshandeling---overwegingen-72}

De terugbetaling van € 20.000,- op 5 september 2006 is een rechtshandeling, omdat dit een op een rechtsgevolg gerichte wil (betaling) is die zich door een verklaring (de daadwerkelijke betaling) heeft geopenbaard.

\paragraph{Verricht door de schuldenaar - overwegingen}\label{verricht-door-de-schuldenaar---overwegingen-72}

De terugbetaling is verricht door de failliet (de schuldenaar), zoals blijkt uit de feiten (onder e).

\paragraph{Om niet - overwegingen}\label{om-niet---overwegingen-72}

De terugbetaling van de lening is een rechtshandeling `anders dan om niet', omdat deze voortvloeit uit een geldleningsovereenkomst (art. 7A:1571 BW) waarbij een wederprestatie (terugbetaling met rente) was bedongen.

\paragraph{Onverplicht - overwegingen}\label{onverplicht---overwegingen-72}

De rechtshandelingen betreffen de terugbetaling van leningen die Thomassen aan de failliet had verstrekt. Uit de casus blijkt dat de leningen waren afgesloten met een opeisbaarheidsclausule (art. 9 van de leningsovereenkomst van 15 mei 2006), maar dat de leningen op het moment van terugbetaling (september en oktober 2006) nog niet opeisbaar waren, omdat de failliet nog niet in faillissement was verklaard en de afspraak was dat de lening uiterlijk op 31 december 2006 zou worden afgelost. De terugbetalingen vonden dus plaats voordat de schuld opeisbaar was. Betaling van een niet-opeisbare schuld is onverplicht in de zin van art. 42 Fw. Bovendien was er geen sprake van een wettelijke of contractuele verplichting om de leningen vervroegd af te lossen, ook al was vervroegde aflossing toegestaan (art. 5). De stelplicht en bewijslast van het onverplichte karakter rusten op de curator, en uit de feiten blijkt dat de terugbetalingen niet afgedwongen konden worden op het moment van betaling. Derhalve zijn de rechtshandelingen onverplicht verricht.

\paragraph{Benadeling schuldeisers - overwegingen}\label{benadeling-schuldeisers---overwegingen-72}

Op basis van de casus kan worden vastgesteld dat de schuldeisers zijn benadeeld in de zin van artikel 42 Fw. Hiertoe wordt het volgende overwogen:

\begin{enumerate}
\def\labelenumi{\arabic{enumi}.}
\tightlist
\item
  \textbf{Vermogensvergelijking}:

  \begin{itemize}
  \tightlist
  \item
    Zonder de gewraakte terugbetalingen (totaal € 91.650) aan Thomassen zou de failliet over deze bedragen hebben kunnen beschikken om haar schulden aan andere schuldeisers (zoals de fiscus, UWV, ABP en werknemers) te voldoen.
  \item
    De failliet had op het moment van de terugbetalingen (september-oktober 2006) al financiële problemen: in november 2006 werd betalingsonmacht gemeld voor loonbelasting en pensioenpremies, en werknemers ontvingen geen loon voor november 2006. De RvC-notulen van 9 november 2006 tonen aan dat de failliet een negatieve winst had en worst-case scenario's voorzag met aanzienlijke verliezen.
  \item
    De terugbetalingen aan Thomassen hebben het actief van de failliet verminderd, terwijl de passiva (schulden aan andere crediteuren) ongewijzigd bleven. Dit leidt tot een vermindering van het voor de schuldeisers beschikbare vermogen.
  \end{itemize}
\item
  \textbf{Benadeling in de zin van artikel 42 Fw}:

  \begin{itemize}
  \tightlist
  \item
    De terugbetalingen waren \textbf{onverplicht} in de zin van artikel 42 lid 1 Fw, omdat de leningen kortlopend waren (tot uiterlijk 31 december 2006) en vervroegde terugbetaling was toegestaan (artikel 5 van de leningsovereenkomst). Er was derhalve geen dringende verplichting om de leningen in september-oktober 2006 af te lossen.
  \item
    Thomassen, als bestuurder en lener, \textbf{wist of behoorde te weten} dat de failliet in financiële moeilijkheden verkeerde. De RvC-notulen en de betalingsonmachtmeldingen in november 2006 (maar met terugwerkende kracht voor oktober 2006) duiden op een situatie waarin de failliet haar verplichtingen niet kon nakomen. Thomassen was als bestuurder op de hoogte van deze financiële toestand.
  \item
    De benadeling is \textbf{middellijk}: op het moment dat de curator zijn rechten doet gelden (na faillietverklaring), is sprake van een vermindering van het actief ten opzichte van de hypothetische situatie waarin de terugbetalingen niet hadden plaatsgevonden.
  \end{itemize}
\item
  \textbf{Geen tegenwerping van reële verhaalsmogelijkheden}:

  \begin{itemize}
  \tightlist
  \item
    Er is geen aanwijzing dat de andere schuldeisers (zoals de fiscus, UWV, ABP of werknemers) zich niet op de terugbetaalde bedragen hadden kunnen verhalen. De bedragen waren immers liquide middelen die voor alle schuldeisers beschikbaar hadden moeten zijn.
  \end{itemize}
\item
  \textbf{Conclusie}:
  De terugbetalingen aan Thomassen hebben het vermogen van de failliet ten nadele van de overige schuldeisers verminderd, terwijl Thomassen wist of behoorde te weten dat dit tot benadeling zou leiden. Derhalve is sprake van benadeling in de zin van artikel 42 Fw.
\end{enumerate}

\paragraph{Wetenschap - overwegingen}\label{wetenschap---overwegingen-72}

De leningen zijn verricht binnen een jaar voor faillissement (art. 43 lid 1 Fw), maar de schuldenaar had zich reeds vóór de jaartermijn (13 december 2005) tot de lening van 15 mei 2006 verplicht (art. 43 lid 1 Fw: `niet reeds voor de aanvang van die termijn daartoe had verplicht'). Hierdoor is het bewijsvermoeden van art. 43 Fw niet van toepassing. Wel is Thomassen als bestuurder een gelieerde partij (art. 43 lid 1 sub 4° Fw), maar dit vermoeden geldt niet omdat de verplichting reeds bestond. De curator moet derhalve de wetenschap van benadeling (art. 42 Fw) rechtstreeks aantonen. Uit de notulen van 9 november 2006 blijkt dat de failliet in financiële moeilijkheden verkeerde (negatieve winst, worst-case scenario, noodzaak tot externe financiering). De betalingsonmachtmelding van 30 november 2006 en het niet betalen van lonen in november 2006 onderstrepen dit. Gelet hierop was het faillissement en een tekort daarin met een redelijke mate van waarschijnlijkheid te voorzien, zodat de wetenschap van de schuldenaar (failliet) en Thomassen (als wederpartij bij een niet-om niet handeling) kan worden aangenomen.

\newpage

\markboth{JOR 2012/196}{gpt-5.5 (k = 0)}

\subsubsection{GPT}\label{gpt-24}

\texttt{gpt-5.5}

\begin{center}
\small
\begin{tabular}{@{}p{7.4cm}p{2.6cm}p{2.6cm}@{}}
\toprule
Variabele & Waarde (model) & Waarde (gold) \\
\midrule
Is rechtshandeling & \texttt{true} & \texttt{true} \\
Datum rechtshandeling & \texttt{2006-09-05} & \texttt{2006-09-05} \\
Verricht door de schuldenaar & \texttt{true} & \texttt{true} \\
Om niet & \texttt{false} & \texttt{false} \\
Datum faillietverklaring & \texttt{2006-12-13} & \texttt{2006-12-13} \\
Onverplicht & \textcolor{gray}{\texttt{true}} & \texttt{false} \\
Benadeling schuldeisers & \textcolor{gray}{\texttt{true}} & \texttt{false} \\
Wetenschap schuldenaar aangetoond & \texttt{false} & \texttt{false} \\
Wetenschap wederpartij aangetoond & \texttt{false} & \texttt{null} \\
Vermoeden wetenschap schuldenaar tegenbewijs & \texttt{false} & \texttt{null} \\
Vermoeden wetenschap wederpartij tegenbewijs & \texttt{false} & \texttt{null} \\
Tegen onderwaarde & \texttt{false} & \texttt{null} \\
Voldoening of zekerheid niet opeisbare schuld & \textcolor{gray}{\texttt{true}} & \texttt{null} \\
Gelieerde partij & \textcolor{gray}{\texttt{true}} & \texttt{null} \\
Reeds verplicht voor aanvang termijn & \texttt{false} & \texttt{null} \\
\textbf{Vernietigbaar} & \textbf{\textcolor{gray}{\texttt{true}}} & \textbf{\texttt{false}} \\
\bottomrule
\end{tabular}
\end{center}

\paragraph{Is rechtshandeling - overwegingen}\label{is-rechtshandeling---overwegingen-73}

De gewraakte handelingen betreffen terugbetalingen door de failliet aan Thomassen van geleende bedragen. Een terugbetaling strekt tot rechtsgevolg, namelijk voldoening/vermindering van de schuld uit geldlening, en kwalificeert daarom als rechtshandeling.

\paragraph{Verricht door de schuldenaar - overwegingen}\label{verricht-door-de-schuldenaar---overwegingen-73}

Uit de feiten volgt dat de failliet op 5 september 2006, 14 oktober 2006 en 24 oktober 2006 bedragen aan Thomassen heeft terugbetaald. De rechtshandelingen zijn dus door de schuldenaar zelf verricht.

\paragraph{Om niet - overwegingen}\label{om-niet---overwegingen-73}

De terugbetalingen vonden plaats ter aflossing van eerder door Thomassen aan de failliet verstrekte geldleningen. Daartegenover stond de bevrijding van de failliet van haar terugbetalingsverplichting; de handelingen zijn daarom anders dan om niet.

\paragraph{Onverplicht - overwegingen}\label{onverplicht---overwegingen-73}

Voor de terugbetalingen op de schriftelijke geldlening van € 70.000 bestond op 14 en 24 oktober 2006 nog geen in rechte afdwingbare betalingsplicht. In de overeenkomst is bepaald dat de lening liep tot uiterlijk 31 december 2006 en dat vervroegde terugbetaling steeds was toegestaan. Die bepaling gaf de failliet slechts de bevoegdheid om eerder af te lossen; zij maakte de schuld niet reeds vóór 31 december 2006 opeisbaar en verplichtte de failliet niet tot eerdere betaling. De in artikel 9 genoemde opeisbaarheidsgronden deden zich ten tijde van de betalingen nog niet voor. Dat de failliet praktisch of economisch aanleiding had om Thomassen terug te betalen, schept geen rechtsplicht. Voor zover het gaat om de terugbetaling van de in augustus 2006 mondeling verstrekte lening van € 20.000 is minder duidelijk of een termijn was afgesproken; uit de casus volgt geen schriftelijke niet-opeisbaarheid. De aangevochten aflossingen op de schriftelijke lening zijn echter vóór opeisbaarheid gedaan en daarmee onverplicht in de zin van artikel 42 Fw.

\paragraph{Benadeling schuldeisers - overwegingen}\label{benadeling-schuldeisers---overwegingen-73}

De aangevochten rechtshandelingen betreffen de terugbetalingen aan Thomassen van in totaal € 91.650 kort vóór het faillissement. Bij de vermogensvergelijking moet de feitelijke situatie waarin deze betalingen in stand blijven worden vergeleken met de hypothetische situatie zonder deze betalingen. In de feitelijke situatie is Thomassen vrijwel volledig voldaan, terwijl andere schuldeisers, waaronder werknemers, fiscus, UWV/ABP en mogelijk overige crediteuren, onbetaald zijn gebleven. Zonder de betalingen aan Thomassen zou dit bedrag in beginsel beschikbaar zijn gebleven voor de gezamenlijke schuldeisers, althans zou het niet uitsluitend aan één schuldeiser ten goede zijn gekomen. Dat door de terugbetaling ook de schuld aan Thomassen is verminderd, neemt de benadeling niet weg: bij betaling van een bestaande schuld kan benadeling bestaan in verstoring van de rangorde/paritas creditorum, doordat één concurrente schuldeiser volledig wordt voldaan ten koste van de verhaalsmogelijkheden van de overige schuldeisers. Uit de casus blijkt niet dat alle overige schulden ook zonder deze betalingen volledig voldaan hadden kunnen worden of dat de gezamenlijke schuldeisers zich hoe dan ook niet op de betaalde gelden hadden kunnen verhalen. De schuldeisers zijn daarom door de terugbetalingen benadeeld in de zin van artikel 42 Fw.

\paragraph{Wetenschap - overwegingen}\label{wetenschap---overwegingen-73}

De terugbetalingen aan Thomassen zijn niet om niet verricht, zodat ook wetenschap aan de zijde van Thomassen als wederpartij vereist is en art. 45 Fw niet aan de orde is. Rechtstreekse wetenschap van benadeling is op basis van de gegeven feiten niet zelfstandig voor alle betalingen aangetoond. Wel geldt het bewijsvermoeden van art. 43 Fw: de betalingen vonden binnen de jaartermijn plaats, de verplichting tot deze rechtshandelingen bestond niet al vóór aanvang van die termijn, en Thomassen was de enige statutair bestuurder van de failliet, zodat sprake is van een gelieerde partij in de zin van art. 43 lid 1 sub 4 Fw. Daarnaast betroffen in elk geval de terugbetalingen van de schriftelijke lening voldoening vóór de uiterste vervaldatum van 31 december 2006; de bepaling dat vervroegde terugbetaling was toegestaan maakte de schuld nog niet opeisbaar voor Thomassen. Het vermoeden van wetenschap aan beide zijden is niet ontzenuwd: de failliet verkeerde kort daarna in betalingsonmacht, lonen bleven onbetaald, een faillissement op eigen verzoek volgde spoedig en uit de stukken blijkt dat voortzetting afhankelijk was van onzekere externe financiering en nieuwe klanten. Een enkele hoop op redding of mogelijke steun van UMH is onvoldoende tegenbewijs.

\newpage

\markboth{JOR 2012/196}{gemini-3.5-flash (k = 0)}

\subsubsection{Gemini}\label{gemini-24}

\texttt{gemini-3.5-flash}

\begin{center}
\small
\begin{tabular}{@{}p{7.4cm}p{2.6cm}p{2.6cm}@{}}
\toprule
Variabele & Waarde (model) & Waarde (gold) \\
\midrule
Is rechtshandeling & \texttt{true} & \texttt{true} \\
Datum rechtshandeling & \texttt{2006-09-05} & \texttt{2006-09-05} \\
Verricht door de schuldenaar & \texttt{true} & \texttt{true} \\
Om niet & \texttt{false} & \texttt{false} \\
Datum faillietverklaring & \texttt{2006-12-13} & \texttt{2006-12-13} \\
Onverplicht & \textcolor{gray}{\texttt{true}} & \texttt{false} \\
Benadeling schuldeisers & \textcolor{gray}{\texttt{true}} & \texttt{false} \\
Wetenschap schuldenaar aangetoond & \textcolor{gray}{\texttt{true}} & \texttt{false} \\
Wetenschap wederpartij aangetoond & \textcolor{gray}{\texttt{true}} & \texttt{null} \\
Vermoeden wetenschap schuldenaar tegenbewijs & \texttt{false} & \texttt{null} \\
Vermoeden wetenschap wederpartij tegenbewijs & \texttt{false} & \texttt{null} \\
Tegen onderwaarde & \texttt{null} & \texttt{null} \\
Voldoening of zekerheid niet opeisbare schuld & \textcolor{gray}{\texttt{true}} & \texttt{null} \\
Gelieerde partij & \textcolor{gray}{\texttt{true}} & \texttt{null} \\
Reeds verplicht voor aanvang termijn & \texttt{false} & \texttt{null} \\
\textbf{Vernietigbaar} & \textbf{\textcolor{gray}{\texttt{true}}} & \textbf{\texttt{false}} \\
\bottomrule
\end{tabular}
\end{center}

\paragraph{Is rechtshandeling - overwegingen}\label{is-rechtshandeling---overwegingen-74}

De terugbetalingen van de leningen door de failliet aan Thomassen zijn rechtshandelingen, aangezien deze gericht zijn op het teweegbrengen van rechtsgevolg, te weten het voldoen en daarmee tenietgaan van de verbintenissen uit de geldleningsovereenkomsten en de overdracht van de geldsommen.

\paragraph{Verricht door de schuldenaar - overwegingen}\label{verricht-door-de-schuldenaar---overwegingen-74}

De terugbetalingen zijn rechtstreeks verricht door de schuldenaar (Nutriscience B.V.) aan de schuldeiser (Thomassen).

\paragraph{Om niet - overwegingen}\label{om-niet---overwegingen-74}

De betalingen zijn verricht ter aflossing van de leningen die Thomassen eerder aan de failliet had verstrekt. Omdat hier de eerdere ontvangst van de leensommen tegenover staat, zijn de rechtshandelingen om baat (anders dan om niet) verricht.

\paragraph{Onverplicht - overwegingen}\label{onverplicht---overwegingen-74}

De terugbetalingen van de geldleningen vonden plaats op 5 september 2006, 14 oktober 2006 en 24 oktober 2006. De schriftelijke lening van € 70.000,- had een looptijd tot uiterlijk 31 december 2006. Hoewel artikel 5 van de overeenkomst bepaalde dat vervroegde terugbetaling was toegestaan, betrof dit een bevoegdheid voor de schuldenaar en geen verplichting. De vordering van Thomassen was derhalve pas op 31 december 2006 opeisbaar. De vervroegde terugbetalingen in oktober 2006 betreffen dus de voldoening van een nog niet opeisbare schuld, wat als een onverplichte rechtshandeling kwalificeert. Ten aanzien van de mondelinge lening van € 20.000,- die in augustus 2006 is verstrekt en op 5 september 2006 is terugbetaald, is evenmin gesteld of gebleken dat er op dat moment al een contractuele of wettelijke verplichting tot terugbetaling bestond. Nu voor geen van de betalingen een in rechte afdwingbare rechtsplicht aanwezig was, zijn de rechtshandelingen onverplicht verricht.

\paragraph{Benadeling schuldeisers - overwegingen}\label{benadeling-schuldeisers---overwegingen-74}

Door de terugbetalingen van de geldleningen aan Thomassen in de periode van september tot en met oktober 2006 voor een totaalbedrag van € 91.650,- is er een aanzienlijk bedrag aan liquide middelen aan het vermogen van Nutriscience B.V. onttrokken. Indien deze betalingen niet hadden plaatsgevonden (de hypothetische situatie), had dit actieve vermogen ter beschikking gestaan van de gezamenlijke schuldeisers in het faillissement. In de feitelijke situatie is het actief verminderd en zijn de verhaalsmogelijkheden van de overige schuldeisers, waaronder preferente schuldeisers zoals de Belastingdienst, het UWV, het ABP en de werknemers (die hun loon over november 2006 niet hebben ontvangen), ernstig beperkt. Daarnaast is de rangorde tussen de schuldeisers verstoord doordat Thomassen, als concurrente schuldeiser en tevens bestuurder, volledig is voldaan ten nadele van de overige (deels preferente) schuldeisers. Er is derhalve sprake van benadeling van de schuldeisers in de zin van artikel 42 Fw.

\paragraph{Wetenschap - overwegingen}\label{wetenschap---overwegingen-74}

De terugbetalingen door Nutriscience B.V. aan haar statutair bestuurder Thomassen vonden plaats binnen een jaar voor de faillietverklaring. Omdat Thomassen de bestuurder (natuurlijk persoon) is van de schuldenaar (rechtspersoon), is het bewijsvermoeden van art. 43 lid 1 sub 4° onder a Fw (gelieerde partij) van toepassing. Daarnaast waren de terugbetalingen in oktober 2006 ter voldoening van een niet-opeisbare schuld (de lening was immers pas uiterlijk 31 december 2006 opeisbaar), waardoor ook het vermoeden van art. 43 lid 1 sub 2° Fw geldt. De leningen zijn pas in mei en augustus 2006 aangegaan, dus er was geen eerdere verplichting van voor de jaartermijn. Het tegenbewijs slaagt niet: gezien de penibele financiële situatie (negatieve resultaten, betalingsonmacht fiscus/UWV/ABP en onbetaalde salarissen) was het faillissement met een redelijke mate van waarschijnlijkheid voorzienbaar. Bovendien is de wetenschap van benadeling aan beide zijden ook rechtstreeks aangetoond, aangezien Thomassen als bestuurder en als wederpartij volledig op de hoogte was van deze financiële situatie.

\newpage

\markboth{JOR 2012/196}{mistral/mistral-medium-3.5-128b:fp8 (k = 5)}

\subsection{k = 5}\label{k-5-8}

\subsubsection{Mistral}\label{mistral-25}

\texttt{mistral/mistral-medium-3.5-128b:fp8}

\begin{center}
\small
\begin{tabular}{@{}p{7.4cm}p{2.6cm}p{2.6cm}@{}}
\toprule
Variabele & Waarde (model) & Waarde (gold) \\
\midrule
Is rechtshandeling & \texttt{true} & \texttt{true} \\
Datum rechtshandeling & \texttt{2006-09-05} & \texttt{2006-09-05} \\
Verricht door de schuldenaar & \texttt{true} & \texttt{true} \\
Om niet & \texttt{false} & \texttt{false} \\
Datum faillietverklaring & \texttt{2006-12-13} & \texttt{2006-12-13} \\
Onverplicht & \textcolor{gray}{\texttt{true}} & \texttt{false} \\
Benadeling schuldeisers & \textcolor{gray}{\texttt{true}} & \texttt{false} \\
Wetenschap schuldenaar aangetoond & \texttt{false} & \texttt{false} \\
Wetenschap wederpartij aangetoond & \texttt{false} & \texttt{null} \\
Vermoeden wetenschap schuldenaar tegenbewijs & \texttt{false} & \texttt{null} \\
Vermoeden wetenschap wederpartij tegenbewijs & \texttt{false} & \texttt{null} \\
Tegen onderwaarde & \texttt{null} & \texttt{null} \\
Voldoening of zekerheid niet opeisbare schuld & \texttt{null} & \texttt{null} \\
Gelieerde partij & \textcolor{gray}{\texttt{true}} & \texttt{null} \\
Reeds verplicht voor aanvang termijn & \texttt{false} & \texttt{null} \\
\textbf{Vernietigbaar} & \textbf{\textcolor{gray}{\texttt{true}}} & \textbf{\texttt{false}} \\
\bottomrule
\end{tabular}
\end{center}

\paragraph{Is rechtshandeling - overwegingen}\label{is-rechtshandeling---overwegingen-75}

De terugbetaling van € 20.000,- op 5 september 2006 is een rechtshandeling, omdat dit een op een rechtsgevolg gerichte wil (betaling) is die zich door een verklaring (de daadwerkelijke betaling) heeft geopenbaard.

\paragraph{Verricht door de schuldenaar - overwegingen}\label{verricht-door-de-schuldenaar---overwegingen-75}

De terugbetaling is verricht door de failliet (de schuldenaar), zoals blijkt uit de feiten onder e.

\paragraph{Om niet - overwegingen}\label{om-niet---overwegingen-75}

De terugbetaling is een verplichting voortvloeiend uit de leningsovereenkomst (art. 1 en 3 van de overeenkomst van 15 mei 2006 en de mondelinge overeenkomst van augustus 2006). Er was sprake van een wederprestatie (de lening zelf), dus de handeling is niet `om niet'.

\paragraph{Onverplicht - overwegingen}\label{onverplicht---overwegingen-75}

De betalingen van de failliet aan Thomassen (€ 50.000 op 18 mei 2006, € 20.000 op 23 juni 2006, en € 20.000 op 30 augustus 2006) en de terugbetalingen (€ 20.000 op 5 september 2006, € 51.250 op 14 oktober 2006, en € 20.400 op 24 oktober 2006) zijn onverplicht verricht in de zin van artikel 42 lid 1 Fw.

De leningsovereenkomst van 15 mei 2006 (art. 3) voorzag in een looptijd tot uiterlijk 31 december 2006, met de mogelijkheid tot vervroegde aflossing (art. 5). De terugbetalingen vonden echter plaats \emph{voor} 31 december 2006, maar dit betekent niet dat de betalingen verplicht waren. De schuldenaar (failliet) was namelijk niet verplicht om de lening vervroegd af te lossen; de verplichting bestond slechts tot terugbetaling \emph{uiterlijk} op 31 december 2006. De vervroegde aflossing was derhalve een onverplichte rechtshandeling, omdat er geen rechtsplicht bestond om de lening eerder af te lossen dan de overeengekomen uiterste datum.

Daarnaast is de leningsovereenkomst van augustus 2006 (€ 20.000) niet schriftelijk vastgelegd. Hoewel de curator de stelplicht en bewijslast draagt voor het onverplichte karakter, kan uit de omstandigheden worden afgeleid dat deze lening eveneens niet opeisbaar was op het moment van terugbetaling (5 september 2006). Er was geen sprake van een afdwingbare verplichting tot terugbetaling op dat moment, zeker niet gelet op de financiële toestand van de failliet (betalingsonmacht gemeld in november 2006 voor loonbelasting en pensioenpremies, en onbetaalde lonen in november 2006).

De notulen van de RvC van 9 november 2006 bevestigen dat de failliet al in een financieel moeilijke positie verkeerde, met een negatieve winst en een worst-case scenario van -€ 125.000. Dit onderstreept dat de terugbetalingen aan Thomassen niet verplicht waren, omdat de failliet in feite al in een toestand van betalingsonmacht verkeerde.

Kortom: de betalingen en terugbetalingen waren onverplicht, omdat er geen afdwingbare verplichting bestond om deze op het moment van betaling te verrichten. De stelplicht en bewijslast voor dit onverplichte karakter rusten op de curator, en de omstandigheden van de casus (met name de financiële toestand van de failliet en de afwezigheid van een opeisbare verplichting tot vervroegde aflossing) rechtvaardigen de conclusie dat de rechtshandelingen onverplicht zijn verricht.

\paragraph{Benadeling schuldeisers - overwegingen}\label{benadeling-schuldeisers---overwegingen-75}

Op basis van de casus kan worden geconcludeerd dat de schuldeisers door de terugbetalingen aan Thomassen zijn benadeeld in de zin van artikel 42 Fw. Hiertoe wordt als volgt gemotiveerd:

\begin{enumerate}
\def\labelenumi{\arabic{enumi}.}
\item
  \textbf{Vermogensvergelijking}: De hypothetische situatie zonder de gewraakte rechtshandelingen (terugbetalingen aan Thomassen) moet worden vergeleken met de feitelijke situatie. In de hypothetische situatie zouden de bedragen van € 91.650 (terugbetaald aan Thomassen) in de boedel van Nutriscience B.V. zijn gebleven en beschikbaar zijn geweest voor verdeling onder alle schuldeisers. In de feitelijke situatie zijn deze bedragen aan Thomassen betaald, waardoor het actief van de boedel is verminderd.
\item
  \textbf{Benadeling in de zin van artikel 42 Fw}: De vermindering van het actief van de boedel door de terugbetalingen aan Thomassen leidt tot benadeling van de schuldeisers, omdat deze bedragen niet meer ter beschikking staan voor verdeling onder de gezamenlijke schuldeisers. Dit is een klassiek geval van vermindering van het actief (i), wat voldoende is om benadeling aan te nemen.
\item
  \textbf{Tijdstip van benadeling}: De benadeling moet aanwezig zijn op het moment waarop de curator zijn rechten doet gelden (leer van de middellijke benadeling). Aangezien de terugbetalingen zijn gedaan in september en oktober 2006, en het faillissement is uitgesproken op 13 december 2006, zijn de bedragen op dat moment niet meer in de boedel aanwezig. De schuldeisers zijn dus daadwerkelijk benadeeld doordat deze middelen niet meer voor verdeling beschikbaar zijn.
\item
  \textbf{Geen tegenargumenten}: Er zijn geen aanwijzingen dat de schuldeisers ook zonder de terugbetalingen volledig zouden zijn voldaan (er was sprake van betalingsonmacht, zie punt f en g van de casus). Ook is niet gebleken dat de schuldeisers zich toch niet op deze bedragen hadden kunnen verhalen (bijvoorbeeld door zekerheidsrechten).
\item
  \textbf{Relevante uitspraken}: In de uitspraak ECLI:NL:RBOVE:2014:5599 (rov. 9) wordt benadrukt dat benadeling optreedt wanneer geld uit het vermogen van de schuldenaar verdwijnt dat anders voor de schuldeisers beschikbaar zou zijn geweest. Dit is in deze casus het geval. Ook in ECLI:NL:GHARL:2017:5673 (rov. 2.11) wordt benadrukt dat benadeling aanwezig is wanneer bedragen die anders voor de schuldeisers beschikbaar zouden zijn, aan een derde zijn betaald.
\end{enumerate}

Conclusie: De schuldeisers zijn benadeeld doordat de boedel is verminderd met € 91.650, wat anders voor verdeling onder de schuldeisers beschikbaar zou zijn geweest.

\paragraph{Wetenschap - overwegingen}\label{wetenschap---overwegingen-75}

De rechtshandelingen (terugbetalingen van de lening door de failliet aan Thomassen) zijn verricht binnen een jaar voor de faillietverklaring (13 december 2006). Thomassen was als bestuurder van de failliet een gelieerde partij in de zin van art. 43 lid 1 sub 4° Fw (bestuurder van de rechtspersoon-schuldenaar). Het bewijsvermoeden van art. 43 lid 1 Fw is daarom van toepassing: de wetenschap van benadeling wordt aan beide zijden (schuldenaar en wederpartij) vermoed. Er is geen sprake van een reeds bestaande verplichting vóór de jaartermijn (de lening was pas in mei/juni 2006 gesloten). Tegenbewijs is niet geslaagd: uit de notulen van de RvC (9 november 2006) blijkt dat de financiële situatie van de failliet zeer slecht was (negatieve winst, worst-case scenario van -€125.000,-, en verwachte resultaten in 2007 van -€350.000,-). Bovendien was er al betalingsonmacht gemeld (fiscus, UWV, ABP) en waren salarissen niet betaald. Dit maakt aannemelijk dat het faillissement en een tekort daarin met een redelijke mate van waarschijnlijkheid waren te voorzien. Het vermoeden van art. 43 Fw blijft dus staan.

\newpage

\markboth{JOR 2012/196}{gpt-5.5 (k = 5)}

\subsubsection{GPT}\label{gpt-25}

\texttt{gpt-5.5}

\begin{center}
\small
\begin{tabular}{@{}p{7.4cm}p{2.6cm}p{2.6cm}@{}}
\toprule
Variabele & Waarde (model) & Waarde (gold) \\
\midrule
Is rechtshandeling & \texttt{true} & \texttt{true} \\
Datum rechtshandeling & \texttt{2006-09-05} & \texttt{2006-09-05} \\
Verricht door de schuldenaar & \texttt{true} & \texttt{true} \\
Om niet & \texttt{false} & \texttt{false} \\
Datum faillietverklaring & \texttt{2006-12-13} & \texttt{2006-12-13} \\
Onverplicht & \textcolor{gray}{\texttt{true}} & \texttt{false} \\
Benadeling schuldeisers & \textcolor{gray}{\texttt{true}} & \texttt{false} \\
Wetenschap schuldenaar aangetoond & \texttt{false} & \texttt{false} \\
Wetenschap wederpartij aangetoond & \texttt{false} & \texttt{null} \\
Vermoeden wetenschap schuldenaar tegenbewijs & \texttt{false} & \texttt{null} \\
Vermoeden wetenschap wederpartij tegenbewijs & \texttt{false} & \texttt{null} \\
Tegen onderwaarde & \texttt{false} & \texttt{null} \\
Voldoening of zekerheid niet opeisbare schuld & \textcolor{gray}{\texttt{true}} & \texttt{null} \\
Gelieerde partij & \textcolor{gray}{\texttt{true}} & \texttt{null} \\
Reeds verplicht voor aanvang termijn & \texttt{false} & \texttt{null} \\
\textbf{Vernietigbaar} & \textbf{\textcolor{gray}{\texttt{true}}} & \textbf{\texttt{false}} \\
\bottomrule
\end{tabular}
\end{center}

\paragraph{Is rechtshandeling - overwegingen}\label{is-rechtshandeling---overwegingen-76}

De gewraakte handelingen betreffen terugbetalingen door de failliet aan Thomassen ter voldoening van geldleningen. Daarmee zijn zij gericht op een rechtsgevolg, namelijk betaling/kwijting van een schuld.

\paragraph{Verricht door de schuldenaar - overwegingen}\label{verricht-door-de-schuldenaar---overwegingen-76}

Uit de feiten volgt dat de failliet de bedragen aan Thomassen heeft terugbetaald. De rechtshandelingen zijn dus door de schuldenaar zelf verricht.

\paragraph{Om niet - overwegingen}\label{om-niet---overwegingen-76}

De terugbetalingen vonden plaats ter voldoening van eerder door Thomassen verstrekte geldleningen. Daartegenover stond dus de vermindering of delging van een bestaande schuld; de handelingen zijn daarom anders dan om niet.

\paragraph{Onverplicht - overwegingen}\label{onverplicht---overwegingen-76}

Voor de terugbetalingen op de schriftelijke lening van € 70.000 bestond op de betaaldata in oktober 2006 geen in rechte afdwingbare betalingsplicht. De lening liep tot uiterlijk 31 december 2006; vóór die datum was de hoofdsom niet opeisbaar, behoudens de in artikel 9 genoemde gevallen, die zich op dat moment niet hadden voorgedaan. Artikel 5 bepaalt slechts dat vervroegde terugbetaling was toegestaan, maar een bevoegdheid tot vervroegde aflossing schept geen afdwingbare verplichting om vóór opeisbaarheid te betalen. Het gaat dus om betaling van een niet-opeisbare schuld, hetgeen in de zin van artikel 42 Fw onverplicht is. Voor de mondelinge lening van € 20.000 zijn geen concrete afspraken gesteld waaruit een opeisbare rechtsplicht tot terugbetaling op 5 september 2006 volgt; voor zover die betaling eveneens als vervroegde aflossing moet worden gezien, geldt hetzelfde. De aangevochten terugbetalingen zijn daarom onverplicht verricht.

\paragraph{Benadeling schuldeisers - overwegingen}\label{benadeling-schuldeisers---overwegingen-76}

De benadeling moet worden beoordeeld door vergelijking van de feitelijke situatie waarin de terugbetalingen aan Thomassen in stand blijven met de hypothetische situatie zonder die terugbetalingen. In de feitelijke situatie is in september/oktober 2006 in totaal € 91.650 uit het vermogen van de failliet aan Thomassen betaald. Daarmee is Thomassen, als schuldeiser/bestuurder, volledig of nagenoeg volledig voldaan, terwijl kort daarna betalingsonmacht is gemeld bij onder meer fiscus, UWV en ABP en de werknemers hun loon over november 2006 niet meer hebben ontvangen. Het faillissement volgde op 13 december 2006. Zonder de terugbetalingen zouden de betaalde bedragen in beginsel beschikbaar zijn gebleven voor de gezamenlijke schuldeisers, terwijl Thomassen slechts als schuldeiser in het faillissement had kunnen meedelen. Dat de schuld aan Thomassen door de betalingen is verminderd, neemt de benadeling niet weg: bij een tekort in de boedel leidt betaling aan één concurrente schuldeiser tot verstoring van de paritas creditorum en vermindering van de verhaalsmogelijkheden van de overige schuldeisers. Er zijn geen aanwijzingen dat alle schuldeisers ook zonder vernietiging volledig zouden zijn voldaan, of dat de overige schuldeisers zich niet op deze gelden hadden kunnen verhalen. De schuldeisers zijn daarom door de terugbetalingen benadeeld in de zin van artikel 42 Fw.

\paragraph{Wetenschap - overwegingen}\label{wetenschap---overwegingen-76}

De terugbetalingen aan Thomassen vonden plaats vóór faillissement en betroffen geen rechtshandelingen om niet, maar voldoening van geldleningen. Rechtstreekse wetenschap dat ten tijde van alle terugbetalingen het faillissement en een tekort daarin met redelijke mate van waarschijnlijkheid waren te voorzien, is op basis van de gegeven feiten niet zelfstandig overtuigend aangetoond; de zwaarste aanwijzingen (RvC-notulen 9 november, melding betalingsonmacht 30 november, onbetaald loon november en eigen faillissementsaanvraag in december) liggen grotendeels ná de betalingen. Wel is art. 43 Fw aan de orde: Thomassen was tot faillissement enig statutair bestuurder van de failliet, zodat sprake is van een gelieerde partij ex art. 43 lid 1 sub 4° Fw. Daarnaast zijn in elk geval de terugbetalingen op de schriftelijke lening vóór de uiterste aflossingsdatum van 31 december 2006 aan te merken als voldoening van een nog niet opeisbare schuld ex art. 43 lid 1 sub 2° Fw. De rechtshandelingen vonden binnen de relevante jaarperiode plaats en er blijkt niet van een verplichting die al vóór aanvang van die termijn bestond. Daarmee wordt de wetenschap van benadeling aan beide zijden vermoed. Uit de casus volgt geen voldoende tegenbewijs dat faillissement en tekort toen niet voorzienbaar waren; Thomassen had als bestuurder juist inzicht in de financiële positie.

\newpage

\markboth{JOR 2012/196}{gemini-3.5-flash (k = 5)}

\subsubsection{Gemini}\label{gemini-25}

\texttt{gemini-3.5-flash}

\begin{center}
\small
\begin{tabular}{@{}p{7.4cm}p{2.6cm}p{2.6cm}@{}}
\toprule
Variabele & Waarde (model) & Waarde (gold) \\
\midrule
Is rechtshandeling & \texttt{true} & \texttt{true} \\
Datum rechtshandeling & \texttt{2006-09-05} & \texttt{2006-09-05} \\
Verricht door de schuldenaar & \texttt{true} & \texttt{true} \\
Om niet & \texttt{false} & \texttt{false} \\
Datum faillietverklaring & \texttt{2006-12-13} & \texttt{2006-12-13} \\
Onverplicht & \textcolor{gray}{\texttt{true}} & \texttt{false} \\
Benadeling schuldeisers & \textcolor{gray}{\texttt{true}} & \texttt{false} \\
Wetenschap schuldenaar aangetoond & \textcolor{gray}{\texttt{true}} & \texttt{false} \\
Wetenschap wederpartij aangetoond & \textcolor{gray}{\texttt{true}} & \texttt{null} \\
Vermoeden wetenschap schuldenaar tegenbewijs & \texttt{false} & \texttt{null} \\
Vermoeden wetenschap wederpartij tegenbewijs & \texttt{false} & \texttt{null} \\
Tegen onderwaarde & \texttt{null} & \texttt{null} \\
Voldoening of zekerheid niet opeisbare schuld & \textcolor{gray}{\texttt{true}} & \texttt{null} \\
Gelieerde partij & \textcolor{gray}{\texttt{true}} & \texttt{null} \\
Reeds verplicht voor aanvang termijn & \texttt{false} & \texttt{null} \\
\textbf{Vernietigbaar} & \textbf{\textcolor{gray}{\texttt{true}}} & \textbf{\texttt{false}} \\
\bottomrule
\end{tabular}
\end{center}

\paragraph{Is rechtshandeling - overwegingen}\label{is-rechtshandeling---overwegingen-77}

De terugbetalingen van de leningen door de failliet aan Thomassen zijn rechtshandelingen, aangezien deze gericht zijn op het teweegbrengen van rechtsgevolg, te weten het voldoen en daarmee tenietgaan van de verbintenissen uit de geldleningsovereenkomsten en de overdracht van de geldsommen.

\paragraph{Verricht door de schuldenaar - overwegingen}\label{verricht-door-de-schuldenaar---overwegingen-77}

De terugbetalingen zijn rechtstreeks verricht door de schuldenaar (Nutriscience B.V.) aan de schuldeiser (Thomassen).

\paragraph{Om niet - overwegingen}\label{om-niet---overwegingen-77}

De betalingen zijn verricht ter aflossing van de leningen die Thomassen eerder aan de failliet had verstrekt. Omdat hier de eerdere ontvangst van de leensommen tegenover staat, zijn de rechtshandelingen om baat (anders dan om niet) verricht.

\paragraph{Onverplicht - overwegingen}\label{onverplicht---overwegingen-77}

De schriftelijke geldlening van € 70.000,- had een looptijd tot uiterlijk 31 december 2006. De terugbetalingen van € 51.250,- en € 20.400,- vonden plaats in oktober 2006, derhalve vóór de overeengekomen uiterste aflosdatum. Hoewel artikel 5 van de leningsovereenkomst vervroegde terugbetaling toestond, creëert dit slechts een bevoegdheid voor de schuldenaar en geen in rechte afdwingbare rechtsplicht. Omdat de schuldeiser (Thomassen) de terugbetaling op dat moment nog niet rechtens kon afdwingen, zijn deze vervroegde betalingen onverplicht verricht in de zin van artikel 42 Fw.

\paragraph{Benadeling schuldeisers - overwegingen}\label{benadeling-schuldeisers---overwegingen-77}

In deze casus heeft Nutriscience B.V. (de failliet) in de maanden september en oktober 2006 in totaal € 91.650,- terugbetaald aan haar bestuurder Thomassen ter aflossing van leningen. Om te bepalen of er sprake is van benadeling van de schuldeisers in de zin van artikel 42 Fw, moet een vergelijking worden gemaakt tussen de feitelijke situatie waarin de schuldeisers verkeren na de betalingen en de hypothetische situatie waarin zij zouden hebben verkeerd als de betalingen achterwege waren gebleven. Zonder deze terugbetalingen zou het bedrag van € 91.650,- in het vermogen van de vennootschap zijn gebleven en na het faillissement beschikbaar zijn geweest voor de gezamenlijke schuldeisers. Door de betalingen is het actief van de vennootschap met dit bedrag verminderd. Bovendien is de paritas creditorum (de gelijkheid van schuldeisers) verstoord, aangezien Thomassen (als bestuurder/schuldeiser) volledig is voldaan, terwijl andere schuldeisers, zoals de fiscus, het UWV, het ABP en de werknemers (die hun loon over november 2006 niet hebben ontvangen), onbetaald zijn gebleven. Er rustten geen zekerheidsrechten op de goederen van de failliet ten behoeve van Thomassen die deze benadeling zouden wegnemen. De schuldeisers zijn derhalve daadwerkelijk in hun verhaalsmogelijkheden benadeeld door de aangevochten rechtshandelingen.

\paragraph{Wetenschap - overwegingen}\label{wetenschap---overwegingen-77}

De terugbetalingen door Nutriscience B.V. aan haar bestuurder Thomassen zijn onverplicht verricht binnen een jaar voor het faillissement. Omdat Thomassen zowel de enig statutair bestuurder van de schuldenaar als de wederpartij in privé was, is de wetenschap van benadeling aan beide zijden rechtstreeks aangetoond. Bovendien gelden de bewijsvermoedens van art. 43 lid 1 sub 2 Fw (voldoening van een niet-opeisbare schuld, aangezien de lening pas op 31 december 2006 opeisbaar was) en sub 4 Fw (transactie met een bestuurder). Er is geen succesvol tegenbewijs geleverd om deze vermoedens te ontzenuwen.

\newpage

\markboth{JOR 2012/196}{mistral/mistral-medium-3.5-128b:fp8 (k = 10)}

\subsection{k = 10}\label{k-10-8}

\subsubsection{Mistral}\label{mistral-26}

\texttt{mistral/mistral-medium-3.5-128b:fp8}

\begin{center}
\small
\begin{tabular}{@{}p{7.4cm}p{2.6cm}p{2.6cm}@{}}
\toprule
Variabele & Waarde (model) & Waarde (gold) \\
\midrule
Is rechtshandeling & \texttt{true} & \texttt{true} \\
Datum rechtshandeling & \texttt{2006-09-05} & \texttt{2006-09-05} \\
Verricht door de schuldenaar & \texttt{true} & \texttt{true} \\
Om niet & \texttt{false} & \texttt{false} \\
Datum faillietverklaring & \texttt{2006-12-13} & \texttt{2006-12-13} \\
Onverplicht & \textcolor{gray}{\texttt{true}} & \texttt{false} \\
Benadeling schuldeisers & \textcolor{gray}{\texttt{true}} & \texttt{false} \\
Wetenschap schuldenaar aangetoond & \texttt{false} & \texttt{false} \\
Wetenschap wederpartij aangetoond & \texttt{false} & \texttt{null} \\
Vermoeden wetenschap schuldenaar tegenbewijs & \texttt{false} & \texttt{null} \\
Vermoeden wetenschap wederpartij tegenbewijs & \texttt{false} & \texttt{null} \\
Tegen onderwaarde & \texttt{null} & \texttt{null} \\
Voldoening of zekerheid niet opeisbare schuld & \texttt{null} & \texttt{null} \\
Gelieerde partij & \textcolor{gray}{\texttt{true}} & \texttt{null} \\
Reeds verplicht voor aanvang termijn & \texttt{false} & \texttt{null} \\
\textbf{Vernietigbaar} & \textbf{\textcolor{gray}{\texttt{true}}} & \textbf{\texttt{false}} \\
\bottomrule
\end{tabular}
\end{center}

\paragraph{Is rechtshandeling - overwegingen}\label{is-rechtshandeling---overwegingen-78}

De terugbetaling van € 20.000,- op 5 september 2006 is een rechtshandeling, omdat dit een op een rechtsgevolg gerichte wil (betaling) is die zich door een verklaring (de daadwerkelijke betaling) heeft geopenbaard.

\paragraph{Verricht door de schuldenaar - overwegingen}\label{verricht-door-de-schuldenaar---overwegingen-78}

De terugbetaling is verricht door de failliet (de schuldenaar), zoals blijkt uit de feiten onder e.

\paragraph{Om niet - overwegingen}\label{om-niet---overwegingen-78}

De terugbetaling is een verplichting voortvloeiend uit de leningsovereenkomst (art. 1 en 3 van de overeenkomst van 15 mei 2006 en de mondelinge overeenkomst van augustus 2006). Er was sprake van een wederprestatie (de lening zelf), dus de handeling is niet `om niet'.

\paragraph{Onverplicht - overwegingen}\label{onverplicht---overwegingen-78}

De betalingen door Nutriscience B.V. (de failliet) aan Thomassen ter aflossing van de geldleningen van € 70.000 (mei 2006) en € 20.000 (augustus 2006) zijn onverplicht verricht in de zin van artikel 42 lid 1 Fw.

De leningsovereenkomsten voorzagen in een looptijd tot uiterlijk 31 december 2006 (art. 3 van de overeenkomst van 15 mei 2006) en, voor de lening van augustus 2006, een vergelijkbare termijn. De betalingen vonden plaats op 5 september 2006, 14 oktober 2006 en 24 oktober 2006, dus \emph{na} de vervaldatum van 31 december 2006.

Uit de leningsovereenkomsten blijkt dat de leningen \emph{opeisbaar} waren op 31 december 2006 (art. 3 van de eerste overeenkomst en de impliciete termijn voor de tweede lening). Echter, de betalingen vonden plaats \emph{na} deze datum, wat betekent dat de schuldenaar (Nutriscience) op dat moment \emph{geen verplichting} had om de leningen af te lossen, tenzij Thomassen de leningen daadwerkelijk opeiste.

Uit de casus blijkt niet dat Thomassen de leningen opeiste vóór de betalingen. Artikel 5 van de eerste leningsovereenkomst stelt weliswaar dat vervroegde terugbetaling is toegestaan, maar dit betekent niet dat de schuldenaar verplicht was om vervroegd af te lossen. De betalingen waren dus \emph{niet verplicht} op het moment dat ze werden verricht, omdat de vervaldatum nog niet was bereikt of de leningen niet waren opgevraagd.

Daarnaast blijkt uit de notulen van de RvC van 9 november 2006 (punt h) dat de financiële situatie van Nutriscience B.V. verslechterde en dat er sprake was van een negatieve winst en een worst-case scenario van -€ 125.000. Dit duidt erop dat de betalingen aan Thomassen mogelijk ten koste gingen van andere schuldeisers, wat relevant is voor de benadelingseis in art. 42 Fw.

Kortom: de betalingen waren onverplicht omdat er op het moment van betaling geen opeisbare verplichting bestond (de leningen waren nog niet vervallen of opgevraagd) en de schuldenaar had derhalve geen rechtsplicht om op dat moment te betalen. Morele, feitelijke of economische druk (zoals de relatie tussen Thomassen als bestuurder en de failliet) schept geen rechtsplicht (vgl. ECLI:NL:RBOVE:2014:5599, r.o. 8).

\paragraph{Benadeling schuldeisers - overwegingen}\label{benadeling-schuldeisers---overwegingen-78}

Op basis van de casus kan worden geconcludeerd dat de schuldeisers door de aangevochten rechtshandelingen zijn benadeeld in de zin van artikel 42 Faillissementswet (Fw).

Allereerst is van belang dat de failliet (Nutriscience B.V.) in de periode voor het faillissement leningen heeft afgesloten met Thomassen, haar statutair bestuurder. De leningen bedroegen in totaal € 90.000 (€ 50.000 op 16 mei 2006, € 20.000 op 21 juni 2006 en € 20.000 in augustus 2006). Deze bedragen zijn vervolgens in de periode september tot oktober 2006 terugbetaald, namelijk € 20.000 op 5 september 2006, € 51.250 op 14 oktober 2006 en € 20.400 op 24 oktober 2006, totaal € 91.650.

Uit de notulen van de Raad van Commissarissen van 9 november 2006 blijkt dat de financiële situatie van de failliet slecht was. Er was sprake van een negatieve winst voor belastingen van € 50.000 in 2006, met een worst-case scenario van -€ 125.000. Voor 2007 werd een laagste resultaat van -€ 350.000 verwacht als er geen nieuwe projecten zouden binnenkomen. Bovendien had de failliet betalingsonmacht gemeld aan de fiscus, het UWV en het ABP voor loonbelasting en pensioenpremies over oktober 2006, en hadden de werknemers hun loon over november 2006 niet ontvangen.

De terugbetalingen aan Thomassen hebben geleid tot een vermindering van het actief van de failliet. Zonder deze terugbetalingen zouden de bedragen van € 91.650 beschikbaar zijn geweest voor de schuldeisers van de failliet. Dit betekent dat de schuldeisers in hun verhaalsmogelijkheden zijn benadeeld, omdat het vermogen dat anders voor hen beschikbaar was, is weggestroomd naar Thomassen.

Daarnaast is er sprake van een verstoring van de rangorde tussen de schuldeisers. Thomassen, als bestuurder, heeft een preferente behandeling gekregen ten opzichte van andere schuldeisers, zoals de fiscus, het UWV, het ABP en de werknemers, die hun vorderingen niet of niet tijdig hebben ontvangen. Dit is een vorm van benadeling in de zin van artikel 42 Fw, ook al is er per saldo mogelijk geen vermogensvermindering opgetreden (ECLI:NL:GHARL:2017:5673, r.o. 2.11).

De benadeling moet worden beoordeeld op het moment waarop de curator zijn rechten doet gelden (leer van de middellijke benadeling). Op dat moment is de benadeling aanwezig, omdat de schuldeisers in de hypothetische situatie zonder de gewraakte rechtshandelingen een groter vermogen tot hun beschikking zouden hebben gehad.

Tot slot is er geen sprake van een situatie waarin de schuldeisers toch niet op de betrokken goederen hadden kunnen verhalen, zoals wanneer er zekerheidsrechten op rustten. De benadeling is dus daadwerkelijk aanwezig.

\paragraph{Wetenschap - overwegingen}\label{wetenschap---overwegingen-78}

De rechtshandelingen (terugbetalingen van de lening door de failliet aan Thomassen) zijn verricht binnen een jaar voor de faillietverklaring (13 december 2006). Thomassen was als bestuurder van de failliet een gelieerde partij in de zin van art. 43 lid 1 sub 4° Fw (bestuurder van de rechtspersoon-schuldenaar). Het bewijsvermoeden van art. 43 lid 1 Fw is daarom van toepassing, tenzij de schuldenaar zich reeds voor de aanvang van de jaartermijn (13 december 2005) tot de rechtshandelingen had verplicht. Uit de casus blijkt dat de leningovereenkomst op 15 mei 2006 is gesloten, dus binnen de jaartermijn. Er is geen sprake van een eerdere verplichting (zoals een eerdere lening of afspraak). Het vermoeden van wetenschap van benadeling aan beide zijden (art. 43 lid 1 Fw) is dus van toepassing. Er is geen succesvol tegenbewijs geleverd: de notulen van de RvC van 9 november 2006 (h) tonen aan dat de financiële situatie van de failliet zeer slecht was (negatieve winst, worst-case scenario van -€125.000, en verwachting van -€350.000 in 2007). Bovendien was er sprake van betalingsonmacht (f) en niet-betaalde lonen (g). Dit maakt aannemelijk dat het faillissement en een tekort daarin met een redelijke mate van waarschijnlijkheid waren te voorzien. Het tegenbewijs van Thomassen (dat de failliet `midden in het leven stond') is niet geloofwaardig gelet op deze feiten (vgl. ECLI:NL:GHARL:2017:5658, r.o. 4.4).

\newpage

\markboth{JOR 2012/196}{gpt-5.5 (k = 10)}

\subsubsection{GPT}\label{gpt-26}

\texttt{gpt-5.5}

\begin{center}
\small
\begin{tabular}{@{}p{7.4cm}p{2.6cm}p{2.6cm}@{}}
\toprule
Variabele & Waarde (model) & Waarde (gold) \\
\midrule
Is rechtshandeling & \texttt{true} & \texttt{true} \\
Datum rechtshandeling & \texttt{2006-09-05} & \texttt{2006-09-05} \\
Verricht door de schuldenaar & \texttt{true} & \texttt{true} \\
Om niet & \texttt{false} & \texttt{false} \\
Datum faillietverklaring & \texttt{2006-12-13} & \texttt{2006-12-13} \\
Onverplicht & \textcolor{gray}{\texttt{true}} & \texttt{false} \\
Benadeling schuldeisers & \textcolor{gray}{\texttt{true}} & \texttt{false} \\
Wetenschap schuldenaar aangetoond & \textcolor{gray}{\texttt{true}} & \texttt{false} \\
Wetenschap wederpartij aangetoond & \textcolor{gray}{\texttt{true}} & \texttt{null} \\
Vermoeden wetenschap schuldenaar tegenbewijs & \texttt{false} & \texttt{null} \\
Vermoeden wetenschap wederpartij tegenbewijs & \texttt{false} & \texttt{null} \\
Tegen onderwaarde & \texttt{false} & \texttt{null} \\
Voldoening of zekerheid niet opeisbare schuld & \textcolor{gray}{\texttt{true}} & \texttt{null} \\
Gelieerde partij & \textcolor{gray}{\texttt{true}} & \texttt{null} \\
Reeds verplicht voor aanvang termijn & \texttt{false} & \texttt{null} \\
\textbf{Vernietigbaar} & \textbf{\textcolor{gray}{\texttt{true}}} & \textbf{\texttt{false}} \\
\bottomrule
\end{tabular}
\end{center}

\paragraph{Is rechtshandeling - overwegingen}\label{is-rechtshandeling---overwegingen-79}

De gewraakte handelingen betreffen terugbetalingen door de failliet aan Thomassen ter aflossing van geldleningen. Daarmee zijn zij gericht op rechtsgevolg, namelijk betaling/aflossing en vermindering of tenietgaan van de schuld.

\paragraph{Verricht door de schuldenaar - overwegingen}\label{verricht-door-de-schuldenaar---overwegingen-79}

Uit de feiten volgt dat de failliet op 5 september 2006, 14 oktober 2006 en 24 oktober 2006 bedragen aan Thomassen heeft terugbetaald. De handelingen zijn dus door de schuldenaar zelf verricht.

\paragraph{Om niet - overwegingen}\label{om-niet---overwegingen-79}

De terugbetalingen stonden tegenover eerder door Thomassen aan de failliet verstrekte geldleningen. De failliet betaalde daarmee een eigen schuld af; dit is een rechtshandeling anders dan om niet.

\paragraph{Onverplicht - overwegingen}\label{onverplicht---overwegingen-79}

Voor de schriftelijke geldlening van 15 mei 2006 bestond wel een terugbetalingsverplichting, maar die was pas uiterlijk op 31 december 2006 in rechte afdwingbaar, behoudens de in art. 9 genoemde opeisbaarheidsgronden. Die gronden deden zich ten tijde van de betalingen in oktober 2006 niet voor. Art. 5 bepaalt slechts dat vervroegde terugbetaling was toegestaan; dat schept geen rechtsplicht tot vervroegde betaling. De terugbetalingen vóór 31 december 2006 betroffen daarom betaling van een nog niet-opeisbare schuld en zijn in de zin van art. 42 Fw onverplicht verricht. Voor de mondelinge lening van augustus 2006 zijn geen concrete afwijkende opeisbaarheidsafspraken vermeld; voor zover die betaling op dezelfde tijdelijke overbruggingsfinanciering moet worden betrokken, geldt eveneens dat geen afdwingbare verplichting tot terugbetaling op 5 september 2006 blijkt. Economische of feitelijke druk om Thomassen terug te betalen levert geen rechtsplicht op.

\paragraph{Benadeling schuldeisers - overwegingen}\label{benadeling-schuldeisers---overwegingen-79}

De terugbetalingen aan Thomassen hebben de schuldeisers benadeeld. Bij de vereiste vermogensvergelijking moet de feitelijke situatie waarin Thomassen in september/oktober 2006 in totaal € 91.650 heeft ontvangen, worden vergeleken met de hypothetische situatie zonder die betalingen. Zonder de betalingen zou dit bedrag in beginsel in het vermogen van de failliet zijn gebleven en beschikbaar zijn geweest voor verhaal door de gezamenlijke schuldeisers; Thomassen zou dan voor zijn lening slechts als schuldeiser in het faillissement zijn opgekomen. In de feitelijke situatie is Thomassen daarentegen vóór faillissement volledig voldaan, terwijl andere schuldeisers, waaronder werknemers voor ongeveer € 30.000 aan novemberloon en onder meer fiscus, UWV en ABP, onbetaald zijn gebleven. Ook indien tegenover de betaling een vermindering van de schuld aan Thomassen stond, is sprake van benadeling doordat de rangorde tussen schuldeisers is verstoord en verhaalsactief aan de gezamenlijke schuldeisers is onttrokken. Er zijn geen aanwijzingen dat ook zonder de betalingen alle schulden voldaan hadden kunnen worden, of dat andere schuldeisers zich niet op het betaalde bedrag hadden kunnen verhalen. Dit sluit aan bij de benaderingswijze dat betaling aan één schuldeiser in het zicht van insolventie benadeling kan opleveren doordat actief verdwijnt dat anders voor de boedel beschikbaar was, vgl. bijvoorbeeld ECLI:NL:GHARL:2017:5673 en ECLI:NL:RBMNE:2016:3159.

\paragraph{Wetenschap - overwegingen}\label{wetenschap---overwegingen-79}

De terugbetalingen vonden kort vóór het faillissement plaats aan Thomassen, de enige statutair bestuurder van de failliet. Daarmee is sprake van een gelieerde partij in de zin van art. 43 lid 1 sub 4° Fw. Voor zover het de schriftelijke lening van 15 mei 2006 betreft, was aflossing uiterlijk 31 december 2006 overeengekomen; vervroegde betaling was wel toegestaan, maar de schuld was voor Thomassen nog niet opeisbaar, zodat ook art. 43 lid 1 sub 2° Fw in beeld komt. De schuldenaar had zich niet vóór de jaartermijn tot deze terugbetalingen verplicht. De betalingen waren niet om niet, zodat art. 45 Fw niet aan de orde is, en er is geen duidelijke onderwaarde. Gelet op de ernstige liquiditeitsproblemen, de noodzaak van overbruggingskrediet, de latere meldingen betalingsonmacht, onbetaalde lonen, de RvC-notulen over acute financieringsnood en de eigen faillissementsaanvraag kort daarna, is wetenschap van benadeling bij de failliet en bij Thomassen als bestuurder/crediteur rechtstreeks voldoende aannemelijk. In ieder geval wordt die wetenschap op grond van art. 43 Fw aan beide zijden vermoed en uit de casus blijkt niet van geslaagd tegenbewijs dat faillissement en tekort ten tijde van de betalingen niet met een redelijke mate van waarschijnlijkheid voorzienbaar waren.

\newpage

\markboth{JOR 2012/196}{gemini-3.5-flash (k = 10)}

\subsubsection{Gemini}\label{gemini-26}

\texttt{gemini-3.5-flash}

\begin{center}
\small
\begin{tabular}{@{}p{7.4cm}p{2.6cm}p{2.6cm}@{}}
\toprule
Variabele & Waarde (model) & Waarde (gold) \\
\midrule
Is rechtshandeling & \texttt{true} & \texttt{true} \\
Datum rechtshandeling & \texttt{2006-09-05} & \texttt{2006-09-05} \\
Verricht door de schuldenaar & \texttt{true} & \texttt{true} \\
Om niet & \texttt{false} & \texttt{false} \\
Datum faillietverklaring & \texttt{2006-12-13} & \texttt{2006-12-13} \\
Onverplicht & \textcolor{gray}{\texttt{true}} & \texttt{false} \\
Benadeling schuldeisers & \textcolor{gray}{\texttt{true}} & \texttt{false} \\
Wetenschap schuldenaar aangetoond & \textcolor{gray}{\texttt{true}} & \texttt{false} \\
Wetenschap wederpartij aangetoond & \textcolor{gray}{\texttt{true}} & \texttt{null} \\
Vermoeden wetenschap schuldenaar tegenbewijs & \texttt{false} & \texttt{null} \\
Vermoeden wetenschap wederpartij tegenbewijs & \texttt{false} & \texttt{null} \\
Tegen onderwaarde & \texttt{null} & \texttt{null} \\
Voldoening of zekerheid niet opeisbare schuld & \textcolor{gray}{\texttt{true}} & \texttt{null} \\
Gelieerde partij & \textcolor{gray}{\texttt{true}} & \texttt{null} \\
Reeds verplicht voor aanvang termijn & \texttt{false} & \texttt{null} \\
\textbf{Vernietigbaar} & \textbf{\textcolor{gray}{\texttt{true}}} & \textbf{\texttt{false}} \\
\bottomrule
\end{tabular}
\end{center}

\paragraph{Is rechtshandeling - overwegingen}\label{is-rechtshandeling---overwegingen-80}

De terugbetalingen van de leningen door de failliet aan Thomassen zijn rechtshandelingen, aangezien deze gericht zijn op het teweegbrengen van rechtsgevolg, te weten het voldoen en daarmee tenietgaan van de verbintenissen uit de geldleningsovereenkomsten en de overdracht van de geldsommen.

\paragraph{Verricht door de schuldenaar - overwegingen}\label{verricht-door-de-schuldenaar---overwegingen-80}

De terugbetalingen zijn rechtstreeks verricht door de schuldenaar (Nutriscience B.V.) aan de schuldeiser (Thomassen).

\paragraph{Om niet - overwegingen}\label{om-niet---overwegingen-80}

De betalingen zijn verricht ter aflossing van de leningen die Thomassen eerder aan de failliet had verstrekt. Omdat hier de eerdere ontvangst van de leensommen tegenover staat, zijn de rechtshandelingen om baat (anders dan om niet) verricht.

\paragraph{Onverplicht - overwegingen}\label{onverplicht---overwegingen-80}

De schriftelijke geldleningsovereenkomst van 15 mei 2006 bepaalde in artikel 3 dat de lening liep tot uiterlijk 31 december 2006. De lening was derhalve op de data van de terugbetalingen (september en oktober 2006) nog niet opeisbaar door de schuldeiser, nu er geen opeisbaarheidsgrond als bedoeld in artikel 9 aan de orde was. Hoewel artikel 5 van de overeenkomst vervroegde terugbetaling toestond, schiep dit slechts een bevoegdheid voor de schuldenaar en geen in rechte afdwingbare rechtsplicht. Het betalen van een niet-opeisbare schuld is in het Nederlandse insolventierecht aan te merken als een onverplichte rechtshandeling. Ook ten aanzien van de mondelinge lening van augustus 2006 is niet gebleken van een rechtsplicht tot terugbetaling op 5 september 2006. De aangevochten terugbetalingen zijn dan ook onverplicht verricht.

\paragraph{Benadeling schuldeisers - overwegingen}\label{benadeling-schuldeisers---overwegingen-80}

De aangevochten rechtshandelingen betreffen de terugbetalingen door Nutriscience B.V. aan haar bestuurder/schuldeiser Thomassen van in totaal € 91.650,- in september en oktober 2006. Om te bepalen of er sprake is van benadeling van de schuldeisers in de zin van artikel 42 Fw, moet de feitelijke situatie van de schuldeisers worden vergeleken met de hypothetische situatie waarin zij zouden hebben verkeerd zonder deze terugbetalingen. Indien de terugbetalingen niet hadden plaatsgevonden, zou het bedrag van € 91.650,- in de boedel zijn gevloeid en beschikbaar zijn geweest voor het verhaal van de gezamenlijke schuldeisers. Door de betalingen is het actief van de vennootschap echter substantieel verminderd, terwijl er ten tijde van het faillissement (en reeds kort na de betalingen) aanzienlijke onvoldane schulden waren, waaronder achterstallige belastingen, pensioenpremies en circa € 30.000,- aan onbetaalde werknemerssalariaten over november 2006. De schuldeisers zijn hierdoor feitelijk in hun verhaalsmogelijkheden benadeeld (vermindering van het actief). Er is derhalve sprake van benadeling van de schuldeisers.

\paragraph{Wetenschap - overwegingen}\label{wetenschap---overwegingen-80}

De terugbetalingen in september en oktober 2006 vonden plaats binnen een jaar voor de faillietverklaring (13 december 2006). Omdat de wederpartij (Thomassen) de enige statutair bestuurder van de schuldenaar was, is sprake van een gelieerde partij als bedoeld in art. 43 lid 1 sub 4 Fw. Daarnaast betrof het de voldoening van een niet-opeisbare schuld (art. 43 lid 1 sub 2 Fw), aangezien de lening pas uiterlijk 31 december 2006 opeisbaar was. Het wettelijk bewijsvermoeden van wetenschap van benadeling aan beide zijden is derhalve van toepassing. Dit vermoeden is niet ontzenuwd. Bovendien is de wetenschap van benadeling ook rechtstreeks aangetoond: de vennootschap verkeerde in ernstige liquiditeitsproblemen (getuige de noodzaak van de overbruggingskredieten in mei en augustus 2006) en Thomassen was als enig bestuurder volledig op de hoogte van deze situatie, waardoor het faillissement en een tekort daarin met een redelijke mate van waarschijnlijkheid voorzienbaar waren.

\newpage

\markboth{JOR 2014/76}{Gold standard}

\section{JOR 2014/76}\label{jor-201476}

uitspraakdatum 23-01-2013

\subsection{Gold}\label{gold-9}

\begin{center}
\small
\begin{tabular}{@{}p{8cm}p{2.6cm}@{}}
\toprule
Variabele & Waarde \\
\midrule
Is rechtshandeling & \texttt{true} \\
Datum rechtshandeling & \texttt{2009-02-19} \\
Verricht door de schuldenaar & \texttt{true} \\
Om niet & \texttt{false} \\
Datum faillietverklaring & \texttt{2009-03-09} \\
Onverplicht & \texttt{true} \\
Benadeling schuldeisers & \texttt{true} \\
Wetenschap schuldenaar aangetoond & \texttt{false} \\
Wetenschap wederpartij aangetoond & \texttt{null} \\
Vermoeden wetenschap schuldenaar tegenbewijs & \texttt{false} \\
Vermoeden wetenschap wederpartij tegenbewijs & \texttt{false} \\
Tegen onderwaarde & \texttt{null} \\
Voldoening of zekerheid niet opeisbare schuld & \texttt{true} \\
Gelieerde partij & \texttt{null} \\
Reeds verplicht voor aanvang termijn & \texttt{false} \\
\textbf{Vernietigbaar} & \textbf{\texttt{true}} \\
\bottomrule
\end{tabular}
\end{center}

\subsubsection{Uitspraak}\label{uitspraak-9}

(\ldots; red.)

\paragraph{2. De feiten}\label{de-feiten-3}

2.1. De stichting Stichting Thuiszorg Groningen (verder te noemen: TZG) behoorde tot een groep stichtingen geleid door de stichting Stichting Meavita Nederland (verder: Meavita) die actief was op het gebied van thuiszorgverlening.

2.2. Menzis stelt zich ten doel het bevorderen van zorg voor chronisch zieken en gehandicapten door het uitvoeren van de Algemene wet bijzondere ziektekosten (AWBZ). In dat kader geeft Menzis beschikkingen op basis waarvan gelden worden verstrekt aan zorginstellingen, waaronder TZG. De betalingen geschieden door CAK.

2.3. Tussen de zorgverzekeraars en de zorgaanbieders bestaat een raamovereenkomst, waarin de vergoeding van de voor zorg gemaakte kosten door de zorgverzekeraars en de voorwaarden waaronder die vergoeding wordt verstrekt worden geregeld. TZG is een zorgaanbieder als bedoeld in die raamovereenkomst; Menzis wordt daarin aangeduid als ``het zorgkantoor''. In die raamovereenkomst is onder meer het volgende bepaald:

``Artikel 5 Declaratie en betaling (extramurale zorg)

(\ldots)

\begin{enumerate}
\def\labelenumi{\arabic{enumi}.}
\setcounter{enumi}{3}
\tightlist
\item
  Het zorgkantoor draagt er zorg voor dat maandelijks of per vier weken een voorschot aan de zorgaanbieder wordt overgemaakt. Het voorschot wordt in eerste instantie gebaseerd op de gemaakte productieafspraak per 15 oktober 2006 en vervolgens vanaf de derde maand/periode aangepast aan de gerealiseerde volumes tot maximaal een volume waarbij verwacht kan worden dat de regionale contracteerruimte door de gezamenlijk zorgaanbieders in de loop van het jaar zal worden bereikt. Dit zal worden bepaald volgens een systeem van lineaire extrapolatie tot en met de laatst bekende maand of periode De realisatie in de meest recent gedeclareerde periode wordt daarbij als verwachting voor de resterende periode gehanteerd. Partijen kunnen, indien wenselijk, inzake de bevoorschotting andere afspraken maken.
\end{enumerate}

(\ldots)

Artikel 9 Materiële controle

(\ldots)

\begin{enumerate}
\def\labelenumi{\arabic{enumi}.}
\setcounter{enumi}{1}
\tightlist
\item
  Indien het zorgkantoor heeft vastgesteld dat ten onrechte gedane bedragen zijn gedeclareerd en betaald, behoudt het zorgkantoor zich het recht voor deze terug te vorderen, dan wel te verrekenen met de opvolgende declaraties.''
\end{enumerate}

2.4. Menzis schrijft op 18 december 2008 aan Meavita naar aanleiding van eerdere correspondentie over de situatie van Meavita en TZG onder meer:

``We hebben onder meer het volgende vastgesteld;

-- Het herstelplan, zoals het er ligt, blijkt onhaalbaar.

-- Werkmaatschappij Thuiszorg Groningen is in feite technisch failliet, en heeft na periode 3 geen middelen meer. De basis onder Sensire is minder acuut, maar ook wankel, met name door het risico op vorderingen tussen de werkmaatschappijen onderling.

-- De bestuurder kan de continuïteitsparagraaf met betrekking tot de jaarrekening 2007 niet ondertekenen.

-- De accountant kan geen verklaring afgeven.

-- En er zijn substantiële schuldeisers.''

2.5. Bij brief van 17 februari 2008 heeft Menzis aan Meavita meegedeeld (de rechtbank gaat ervan uit dat de datering op een vergissing berust en dat dit moet zijn: 17 februari 2009):

``Gedurende enige maanden heeft de NZa (de Nederlandse Zorgautoriteit, rechtbank) erin toegestemd dat wij aan u periodiek voorschotten vóóraf voldeden voor nog te verlenen AWBZ-zorg en wel op basis van een relatief hoog productieniveau. Laatstelijk hebben wij u aldus bevoorschot voor de periode tot en met 22 februari 2009.

Gelet op de ontwikkelingen in de afgelopen periode is bij ons ernstige twijfel ontstaan over de vraag of deze wijze van bevoorschotting uit een oogpunt van rechtmatige en doelmatige uitvoering van de AWBZ nog langer te verantwoorden zou zijn. Wij willen daarom niet zonder gerichte toestemming van de NZa het verstrekken van voorschotten op de wijze waarop dit de afgelopen periode is geschied, blijven continueren.

Desgevraagd heeft de NZa ons bij brief van vandaag laten weten dat voor zover Meavita, ter borging van de continuïteit van zorg, nog zorg levert, financiering van deze productie alleen op rechtmatige en doelmatige wijze kan plaatsvinden als dit gebeurt op basis van aantoonbaar gerealiseerde productie.

Op grond hiervan hebben wij besloten ten behoeve van Thuiszorg Groningen na 22 februari 2009 geen voorschotten vooraf meer te verlenen. Er zal vanaf die datum uitsluitend maandelijks achteraf aan u worden betaald voor de in die maand geleverde productie, een en ander zoals in de AWBZ-sector algemeen gebruikelijk is.''

2.6. TZG heeft op 18 februari 2009 een e-mailbericht aan Menzis gestuurd met als bijlage een overzicht van de kosten van de in januari en februari 2009 verleende zorg en het daarvoor ontvangen voorschotbedrag. Volgens dat overzicht beloopt het verschil tussen het voorschotbedrag en de kosten van de verleende zorg € 1.191.126,-.

2.7. Menzis heeft bij brief van 18 februari 2009 aan TZG meegedeeld dat er over januari en februari 2009 een bedrag van € 1.191.000,- te veel is bevoorschot ten behoeve van TZG en gevraagd om terugbetaling van dat bedrag door overschrijving op een bankrekening van CAK.

2.8. Op 19 februari 2009 is van een bankrekening van TZG een bedrag van € 1.191.000,- afgeschreven ten gunste van CAK met vermelding van ``Te veel voorschot jaar 2009 19150'' (verder te noemen: de betaling).

2.9. Op 24 februari 2009 heeft de rechtbank te Groningen aan TZG surseance van betaling verleend en de curatoren tot bewindvoerders benoemd.

2.10. Bij vonnis van diezelfde rechtbank van 9 maart 2009 is TZG failliet verklaard en zijn de curatoren benoemd tot curatoren.

2.11. De curatoren hebben bij brief van 1 oktober 2009 aan Menzis op grond van het bepaalde, primair in artikel 42 Fw en subsidiair in artikel 47 Fw de nietigheid van de betaling ingeroepen.

2.12. Bij brief van 18 november 2009 aan CAK hebben de curatoren aangekondigd (ook) CAK in rechte te zullen betrekken als de betaling niet ongedaan gemaakt wordt.

\paragraph{3. Het geschil}\label{het-geschil-3}

3.1. De curatoren vorderen dat de rechtbank bij vonnis, voor zoveel mogelijk uitvoerbaar bij voorraad,

primair en subsidiair,

\begin{enumerate}
\def\labelenumi{(\roman{enumi})}
\item
  zal verklaren voor recht, dat de betaling door TZG aan Menzis via CAK, althans door TZG aan CAK, ten bedrage van € 1.191.000,-, nietig is;
\item
  Menzis althans CAK zal veroordelen tot betaling van € 1.191.000,- aan de curatoren, te vermeerderen met de wettelijke handelsrente vanaf 19 februari 2009 tot aan de dag van de algehele voldoening;
\end{enumerate}

meer subsidiair en meest subsidiair,

\begin{enumerate}
\def\labelenumi{(\roman{enumi})}
\item
  zal verklaren voor recht, dat Menzis althans CAK onrechtmatig jegens de schuldeisers van TZG heeft gehandeld door de betaling door TZG aan Menzis via CAK, althans door TZG aan CAK, te bewerkstelligen, althans toe te staan, dan wel te verrekenen;
\item
  Menzis althans CAK zal veroordelen tot betaling van € 1.191.000,- aan de curatoren, te vermeerderen met de wettelijke rente vanaf 19 februari 2009 tot aan de dag der algehele voldoening;
\end{enumerate}

primair, subsidiair, meer subsidiair en meest subsidiair,

Menzis althans CAK zal veroordelen in de kosten van deze procedure.

3.2. De vordering is -- tegen de achtergrond van de hiervoor weergegeven feiten -- gegrond op de volgende stellingen.

Primair, gelet op artikel 42 Fw: de betaling geschiedde onverplicht, door de betaling zijn de schuldeisers van TZG benadeeld, doordat het bedrag van de betaling niet aan hen ten goede gekomen is, terwijl dat wel het geval zou zijn geweest als betaling niet was verricht, en Menzis en CAK wisten dat de betaling benadeling van de schuldeisers tot gevolg zou hebben. Dit laatste volgt uit een brief van 18 december 2008 waarin Menzis heeft vastgesteld dat TZG technisch failliet was en na periode 3 van 2009 geen middelen meer had, en uit de brief van 30 december 2009 aan het ministerie van Volksgezondheid, Welzijn en Sport (VWS) en de Nederlandse Zorgautoriteit (NZa).

Subsidiair, gelet op artikel 47 Fw: er heeft overleg plaatsgevonden tussen TZG enerzijds en Menzis en CAK anderzijds. Menzis heeft zich gemengd in de bedrijfsvoering van Meavita en (daarmee) TZG en voor zichzelf bestuurlijke invloed bedongen op het beleid van Meavita. Menzis kreeg gelegenheid om twee leden van de raad van commissarissen (RvC) van Meavita te benoemen, waarbij de besluitvorming in de RvC plaats zou vinden op basis van een ``blocking minority''. Menzis kreeg op die manier doorslaggevende invloed op het beleid van Meavita.

Meer subsidiair: Menzis, althans CAK, heeft onrechtmatig gehandeld jegens de schuldeisers van TZG door anders dan op grond van door de wet erkende redenen van voorrang laatstgenoemde de betaling met voorrang aan CAK te doen verrichten, althans die toe te staan, en dusdoende zichzelf te bevoordelen boven andere schuldeisers. De curatoren zijn onder de gegeven omstandigheden bevoegd voor de gezamenlijke schuldeisers op te treden en de vordering op deze grond in te stellen.

Meest subsidiair, gelet op artikel 54 Fw: Menzis en/of CAK kunnen zich niet te goeder trouw beroepen op een rechtsgeldige verrekening, omdat zij de vordering hebben verkregen in de wetenschap dat de financiële positie van TZG desastreus was en dat TZG afstevende op een faillissement.

3.3. Menzis en CAK voeren gemotiveerd verweer.

3.4. Op de stellingen en weren van partijen wordt, voor zover hier van belang, hieronder ingegaan.

\paragraph{4. De beoordeling}\label{de-beoordeling-5}

4.1. Het verweer van Menzis houdt, in het kort weergegeven, het volgende in. De betaling is bijgeschreven op de bankrekening van CAK, zodat het niet Menzis is die haar heeft ontvangen. De publiekrechtelijke regeling van de geldstromen in de AWBZ-zorg brengt mee, zowel dat de uitkering van de voorschotbedragen niet uit het vermogen van Menzis geschiedde als dat de betaling niet in het vermogen van Menzis is gevloeid, zodat -- wat er verder zij van de standpunten van de curatoren -- Menzis niet gehouden kan zijn tot voldoening van de vordering van de curatoren.

4.2. Verder voert Menzis aan dat het systeem van bevoorschotting meebrengt dat een te hoog voorschot heeft te gelden als onverschuldigd betaald. Menzis licht toe dat dit gedurende enige tijd welbewust gebeurde om de zorgverlening door TZG zeker te stellen, maar dat daarmee niet werd beoogd de verhaalsmogelijkheden voor de commune crediteuren van TZG te verruimen.

4.3. De curatoren weerspreken niet, dat de betaling heeft plaatsgevonden door overschrijving van het bedrag van € 1.191,000,- op een bankrekening van CAK. Zij stellen echter, dat de betaling heeft plaatsgevonden op instigatie van Menzis. Menzis maakte immers deel uit van de begeleidingscommissie die was ingesteld met het oog op de dreigende deconfiture van TZG. Deze commissie stelde regels voor de bevoorschotting. Menzis heeft een vordering ter grootte van het bedrag van de betaling ter verificatie ingediend in het faillissement van TZG. Daaruit volgt, zo stellen de curatoren verder, dat Menzis ervan uitgaat dat het haar vordering is.

4.4. De rechtbank overweegt als volgt. In deze zaak is sprake van een driepartijenverhouding, waarin Menzis op grond van de hiervoor onder 2.3 aangehaalde raamovereenkomst voorschotten verstrekt aan zorgaanbieders, onder wie TZG, en CAK als betaalkantoor fungeert. Menzis is blijkens de hiervoor onder 2.7 weergegeven brief ook de opdrachtgever voor de terugbetaling van het te veel verstrekte voorschot.

4.5. Noch de verstrekking van voorschotten noch de terugvordering van te veel verstrekte voorschotten is een rechtshandeling om niet. De voorschotten worden immers verstrekt ter voldoening aan een wettelijke verplichting, die is uitgewerkt in de meergenoemde raamovereenkomst. Deze raamovereenkomst voorziet in afrekening van de verstrekte voorschotten op basis van de werkelijk verleende zorg, waarbij te veel betaalde voorschotten moeten worden terugbetaald dan wel verrekend met voorschotten over latere perioden.

4.6. Uit dit voorschotstelsel volgt dat, zolang geen eindafrekening heeft plaatsgevonden, een te veel betaald voorschot niet opeisbaar is. Dat de voorschotregeling op enig moment opnieuw is vastgesteld in afwijking van wat oorspronkelijk op basis van de raamovereenkomst gold maakt dit niet anders. Dat op het tijdstip van het verzoek van Menzis aan TZG om terugbetaling van het te veel verstrekte voorschot nog geen eindafrekening over de desbetreffende periode(n) was opgemaakt, staat vast, reeds omdat de periode waarvoor het voorschot was verstrekt nog niet was verstreken en de zorgverlening, naar kan worden aangenomen, nog niet was gestaakt. Volgens vaste rechtspraak geschiedt betaling van een niet opeisbare schuld steeds onverplicht. Nu Menzis het te veel verstrekte voorschot heeft teruggevorderd en de betaling daadwerkelijk heeft plaatsgevonden vóór de eindafrekening, geschiedde de betaling onverplicht als bedoeld in artikel 42 Fw.

4.7. Het bedrag van de betaling is aan één enkele schuldeiser ten goede gekomen met voorbijgaan aan de andere schuldeisers, waarvan een overzicht is gehecht aan het proces-verbaal van de verificatievergadering (productie 10 bij de conclusie van repliek). Die andere schuldeisers zijn daardoor benadeeld.

4.8. Vervolgens is de vraag aan de orde of TZG en Menzis wisten of behoorden te weten dat benadeling van haar schuldeisers het gevolg zou zijn van de betaling. Nu het ging om de betaling van een niet opeisbare schuld, wordt die wetenschap op grond van artikel 43 aanhef en onder 2° Fw behoudens tegenbewijs vermoed aan beide zijden te bestaan.

4.9. De rechtbank verwerpt het verweer van Menzis in dit verband, dat de bevoorschotting gebeurde met het oog op de continuïteit van de zorgverlening door TZG en niet ter voldoening van de handelscrediteuren. Menzis heeft niet toegelicht welke verplichtingen TZG aanging met het oog op de continuïteit van de zorg en welke verplichtingen zij was aangegaan met de handelscrediteuren, en ook niet wat dit onderscheid rechtvaardigt, ervan uitgaande dat TZG laatstbedoelde verplichtingen aanging in het kader van het realiseren van haar bedrijfsdoelstelling, te weten het verlenen van zorg.

4.10. Voorts overweegt de rechtbank dat Menzis op de hoogte was van de precaire toestand waarin TZG verkeerde. Dit volgt genoegzaam uit haar hiervoor onder 2.4 aangehaalde brief aan Meavita, de houdstermaatschappij van TZG. Menzis heeft derhalve onvoldoende argumenten en feiten naar voren gebracht om toegelaten te kunnen worden tot het tegenbewijs tegen het wettelijk vermoeden dat de aan het slot van artikel 42, eerste lid, eerste zin Faillissementswet bedoelde wetenschap aan beide zijden bestaat bij rechtshandelingen ter voldoening van of zekerheidstelling voor een niet opeisbare schuld

4.11. Een en ander leidt tot de conclusie dat de gevorderde verklaring voor recht toewijsbaar is tegen Menzis.

4.12. Ten aanzien van Menzis treft echter het verweer dat de betaling niet door haar is ontvangen doel. De curatoren hebben ter zitting erkend dat het geld (lees: de betaling, rechtbank) niet in het vermogen van Menzis is gevloeid en niet op een bankrekening van Menzis is bijgeschreven, zodat zij ook niet gehouden kan zijn tot ongedaanmaking. De vordering is in zoverre dan ook niet toewijsbaar tegen Menzis.

4.13. Het verweer van CAK houdt primair in, dat zij krachtens de geldende regelgeving slechts fungeert als ``doorgeefluik'' voor betalingen aan verleners van AWBZ-zorg en dat zij dus tussen Menzis en TZG geen zelfstandige positie inneemt. Die betalingen geschieden door CAK namens de zorgverzekeraars met geld uit het Algemeen Fonds Bijzondere Ziektekosten (AFBZ). CAK is dan ook van mening dat zij niets voor zichzelf heeft ontvangen. CAK betwist overigens niet, dat de betaling onverplicht was.

4.14. Dit verweer kan niet slagen, omdat vast staat dat CAK de betaling heeft ontvangen en zij niet vermeldt aan wie zij deze heeft doorgeleid. Zij heeft ook niet betwist dat zij de voorschotten namens Menzis aan TZG uitkeerde (artikel 4 lid 2 Administratiebesluit Bijzondere Ziektekostenverzekering) en bijgevolg ook de betaling wegens te veel uitgekeerd voorschot namens Menzis heeft ontvangen. Daaruit volgt dat de wetenschap van Menzis dat benadeling van de schuldeisers van TZG het gevolg zou zijn van de betaling, aan CAK moet worden toegerekend. De betaling is dan ook op grond van het bepaalde in artikel 42 Fw vernietigbaar. Uit de vernietiging vloeit een ongedaanmakingsverplichting voor CAK voort. De vordering is dus tegen CAK toewijsbaar, zowel wat de verklaring voor recht als wat de ongedaanmaking van de betaling betreft, dit laatste vermeerderd met de verder niet betwiste wettelijke handelsrente.

4.15. Menzis en CAK zullen als de in het ongelijk gestelde partijen worden veroordeeld in de kosten van het geding. (\ldots; red.)

\paragraph{Beslissing}\label{beslissing-1}

De rechtbank

-- verklaart voor recht, dat de betaling door TZG aan Menzis via CAK ten bedrage van € 1.191.000,- nietig is;

-- veroordeelt CAK om tegen behoorlijk bewijs van kwijting aan de curatoren te betalen € 1.191,000,-, vermeerderd met de wettelijke handelsrente als bedoeld in artikel 6:119a BW vanaf 19 februari 2009 tot aan de dag van de algehele voldoening;

-- veroordeelt Menzis en CAK ieder voor haar aandeel (\ldots; red.) in de kosten van dit geding, (\ldots; red.);

-- verklaart dit vonnis tot zover uitvoerbaar bij voorraad;

-- wijst het meer of anders gevorderde af.

\newpage

\markboth{JOR 2014/76}{mistral/mistral-medium-3.5-128b:fp8 (k = 0)}

\subsection{k = 0}\label{k-0-9}

\subsubsection{Mistral}\label{mistral-27}

\texttt{mistral/mistral-medium-3.5-128b:fp8}

\begin{center}
\small
\begin{tabular}{@{}p{7.4cm}p{2.6cm}p{2.6cm}@{}}
\toprule
Variabele & Waarde (model) & Waarde (gold) \\
\midrule
Is rechtshandeling & \texttt{true} & \texttt{true} \\
Datum rechtshandeling & \texttt{2009-02-19} & \texttt{2009-02-19} \\
Verricht door de schuldenaar & \texttt{true} & \texttt{true} \\
Om niet & \texttt{false} & \texttt{false} \\
Datum faillietverklaring & \texttt{2009-03-09} & \texttt{2009-03-09} \\
Onverplicht & \texttt{true} & \texttt{true} \\
Benadeling schuldeisers & \texttt{true} & \texttt{true} \\
Wetenschap schuldenaar aangetoond & \texttt{false} & \texttt{false} \\
Wetenschap wederpartij aangetoond & \texttt{false} & \texttt{null} \\
Vermoeden wetenschap schuldenaar tegenbewijs & \texttt{null} & \texttt{false} \\
Vermoeden wetenschap wederpartij tegenbewijs & \texttt{null} & \texttt{false} \\
Tegen onderwaarde & \texttt{null} & \texttt{null} \\
Voldoening of zekerheid niet opeisbare schuld & \textcolor{gray}{\texttt{null}} & \texttt{true} \\
Gelieerde partij & \textcolor{gray}{\texttt{true}} & \texttt{null} \\
Reeds verplicht voor aanvang termijn & \texttt{false} & \texttt{false} \\
\textbf{Vernietigbaar} & \textbf{\texttt{true}} & \textbf{\texttt{true}} \\
\bottomrule
\end{tabular}
\end{center}

\paragraph{Is rechtshandeling - overwegingen}\label{is-rechtshandeling---overwegingen-81}

De betaling van € 1.191.000,- door TZG aan CAK (op verzoek van Menzis) is een op een rechtsgevolg gerichte wil die zich door een verklaring (overschrijving) heeft geopenbaard. Dit kwalificeert als een rechtshandeling in de zin van art. 3:33 BW.

\paragraph{Verricht door de schuldenaar - overwegingen}\label{verricht-door-de-schuldenaar---overwegingen-81}

De betaling is verricht door TZG (de schuldenaar), zoals blijkt uit de afschrijving van de bankrekening van TZG (2.8).

\paragraph{Om niet - overwegingen}\label{om-niet---overwegingen-81}

De betaling geschiedde ter terugbetaling van een te hoog voorschot (art. 9 lid 2 raamovereenkomst). Dit is een verplichting die voortvloeit uit de contractuele afspraken, waardoor er sprake is van een wederprestatie (de verstrekking van het voorschot). Derhalve is de handeling niet `om niet'.

\paragraph{Onverplicht - overwegingen}\label{onverplicht---overwegingen-81}

De betaling van € 1.191.000,- door TZG aan CAK (op instigatie van Menzis) is onverplicht in de zin van artikel 42 lid 1 Fw. Uit de raamovereenkomst (art. 5 lid 4) blijkt dat voorschotten worden verstrekt op basis van verwachte productie, maar dat deze achteraf worden gecorrigeerd op basis van de werkelijk gerealiseerde productie. Artikel 9 lid 2 van de raamovereenkomst biedt Menzis het recht om ten onrechte betaalde bedragen terug te vorderen of te verrekenen met toekomstige declaraties. Echter, op het moment van de betaling (19 februari 2009) was er nog geen sprake van een opeisbare vordering van Menzis/CAK op TZG. De betaling was derhalve niet gebaseerd op een bestaande, afdwingbare rechtsplicht. De brief van Menzis d.d. 18 februari 2009, waarin om terugbetaling wordt verzocht, toont aan dat Menzis de betaling als een onverschuldigde betaling beschouwde, maar dit creëert nog geen rechtsplicht voor TZG om direct te betalen. Pas na de betaling zou Menzis/CAK een vordering kunnen doen gelden, maar op het moment van betalen bestond deze vordering nog niet als opeisbare schuld. Derhalve was de betaling onverplicht.

\paragraph{Benadeling schuldeisers - overwegingen}\label{benadeling-schuldeisers---overwegingen-81}

Op basis van de casus kan worden vastgesteld dat de betaling van € 1.191.000,- door TZG aan CAK (op instigatie van Menzis) heeft geleid tot benadeling van de schuldeisers in de zin van artikel 42 Fw. Hiertoe wordt als volgt gemotiveerd: (1) \textbf{Vermogensvergelijking}: Zonder de betaling zou het bedrag van € 1.191.000,- in de boedel van TZG zijn gebleven en ten goede zijn gekomen aan de gezamenlijke schuldeisers. In de feitelijke situatie is dit bedrag afgevloeid naar CAK (en indirect Menzis), waardoor de boedel met dit bedrag is verminderd. (2) \textbf{Onverplicht karakter}: De betaling was onverplicht, omdat Menzis zelf had besloten om de bevoorschotting stop te zetten en achteraf te betalen (brief van 17 februari 2009). De terugvordering van het te veel betaalde voorschot was derhalve geen contractuele verplichting, maar een eenzijdige maatregel van Menzis. (3) \textbf{Kennis van benadeling}: Menzis wist of behoorde te weten dat TZG technisch failliet was en na periode 3 van 2009 geen middelen meer had (brief van 18 december 2008). De betaling heeft derhalve geleid tot een vermindering van het actief van TZG, terwijl Menzis en CAK wisten dat dit ten koste zou gaan van de andere schuldeisers. (4) \textbf{Geen compenserend voordeel}: Er is geen sprake van een compenserend voordeel voor de schuldeisers, omdat het bedrag niet is gebruikt om andere schulden af te lossen of de continuïteit van TZG te waarborgen. De betaling heeft derhalve geleid tot een directe benadeling van de schuldeisers door vermindering van het actief.

\paragraph{Wetenschap - overwegingen}\label{wetenschap---overwegingen-81}

De betaling van € 1.191.000,- door TZG aan CAK (op instigatie van Menzis) vond plaats op 19 februari 2009, binnen één jaar voor de faillietverklaring (9 maart 2009). TZG en Menzis/CAK behoren tot een gelieerde structuur (art. 43 lid 1 sub 6° Fw: groepsmaatschappij). Het bewijsvermoeden van art. 43 lid 1 Fw is dus van toepassing: de wetenschap van benadeling wordt aan beide zijden (TZG en Menzis/CAK) vermoed. Er is geen sprake van een reeds bestaande verplichting vóór de jaartermijn (de betaling was een terugvordering van een te hoog voorschot, geen contractuele verplichting). Tegenbewijs is niet aannemelijk: Menzis had in december 2008 al vastgesteld dat TZG technisch failliet was en geen middelen meer had, wat de wetenschap van benadeling bevestigt. De curator hoeft derhalve de wetenschap niet rechtstreeks aan te tonen.

\newpage

\markboth{JOR 2014/76}{gpt-5.5 (k = 0)}

\subsubsection{GPT}\label{gpt-27}

\texttt{gpt-5.5}

\begin{center}
\small
\begin{tabular}{@{}p{7.4cm}p{2.6cm}p{2.6cm}@{}}
\toprule
Variabele & Waarde (model) & Waarde (gold) \\
\midrule
Is rechtshandeling & \texttt{true} & \texttt{true} \\
Datum rechtshandeling & \texttt{2009-02-19} & \texttt{2009-02-19} \\
Verricht door de schuldenaar & \texttt{true} & \texttt{true} \\
Om niet & \texttt{false} & \texttt{false} \\
Datum faillietverklaring & \texttt{2009-03-09} & \texttt{2009-03-09} \\
Onverplicht & \textcolor{gray}{\texttt{false}} & \texttt{true} \\
Benadeling schuldeisers & \texttt{true} & \texttt{true} \\
Wetenschap schuldenaar aangetoond & \textcolor{gray}{\texttt{true}} & \texttt{false} \\
Wetenschap wederpartij aangetoond & \textcolor{gray}{\texttt{true}} & \texttt{null} \\
Vermoeden wetenschap schuldenaar tegenbewijs & \texttt{null} & \texttt{false} \\
Vermoeden wetenschap wederpartij tegenbewijs & \texttt{null} & \texttt{false} \\
Tegen onderwaarde & \texttt{false} & \texttt{null} \\
Voldoening of zekerheid niet opeisbare schuld & \textcolor{gray}{\texttt{false}} & \texttt{true} \\
Gelieerde partij & \texttt{false} & \texttt{null} \\
Reeds verplicht voor aanvang termijn & \texttt{null} & \texttt{false} \\
\textbf{Vernietigbaar} & \textbf{\textcolor{gray}{\texttt{false}}} & \textbf{\texttt{true}} \\
\bottomrule
\end{tabular}
\end{center}

\paragraph{Is rechtshandeling - overwegingen}\label{is-rechtshandeling---overwegingen-82}

De gewraakte handeling is de betaling/overschrijving van € 1.191.000 vanaf een bankrekening van TZG ten gunste van CAK met de vermelding dat het een te veel ontvangen voorschot betrof. Een betaling is gericht op rechtsgevolg, namelijk vermogensverschuiving en, volgens de strekking daarvan, voldoening/terugbetaling van een schuld.

\paragraph{Verricht door de schuldenaar - overwegingen}\label{verricht-door-de-schuldenaar---overwegingen-82}

Het bedrag is op 19 februari 2009 afgeschreven van een bankrekening van TZG. De curatoren duiden dit ook aan als de betaling door TZG aan Menzis via CAK, althans door TZG aan CAK.

\paragraph{Om niet - overwegingen}\label{om-niet---overwegingen-82}

De betaling werd gedaan als terugbetaling van een volgens Menzis te veel verstrekt voorschot over januari en februari 2009. Daarmee stond tegenover de betaling de beoogde vermindering of voldoening van een terugbetalingsverplichting; uit de feiten blijkt niet dat sprake was van een gift of dividenduitkering zonder wederprestatie.

\paragraph{Onverplicht - overwegingen}\label{onverplicht---overwegingen-82}

Voor onverplichtheid is beslissend of TZG tot de betaling van € 1.191.000 in rechte niet gehouden was. Uit de casus volgt dat TZG op grond van de AWBZ-/raamovereenkomst werd bevoorschot voor te leveren zorg en dat zij zelf op 18 februari 2009 aan Menzis een overzicht stuurde waaruit bleek dat over januari/februari 2009 € 1.191.126 meer was bevoorschot dan door geleverde zorg werd gedekt. Menzis heeft daarop terugbetaling van het te veel bevoorschotte bedrag gevraagd. Een te hoog voorschot waarvoor geen grondslag in gerealiseerde productie bestaat, levert een terugbetalingsverplichting op, althans een rechtens afdwingbare verplichting tot restitutie van het onverschuldigd of te veel betaalde. Dat betaling plaatsvond via CAK maakt dit niet anders; dat betreft de wijze van uitvoering binnen de AWBZ-betaalstructuur. De casus biedt onvoldoende aanknopingspunten dat de terugbetalingsschuld nog niet opeisbaar was of dat slechts later via eindafrekening/verrekening mocht worden gecorrigeerd. Daarom was de betaling niet onverplicht in de zin van artikel 42 Fw.

\paragraph{Benadeling schuldeisers - overwegingen}\label{benadeling-schuldeisers---overwegingen-82}

De betaling van € 1.191.000 kort vóór de surseance en het faillissement heeft het actief van TZG verminderd met dat bedrag. Daar staat wel tegenover dat een eventuele terugbetalingsvordering wegens te hoge bevoorschotting met hetzelfde bedrag is verminderd. Per saldo hoeft daarom geen zuivere vermogensvermindering te zijn ontstaan. Dat neemt de benadeling echter niet weg: door de betaling is één schuldeiser, Menzis/CAK, volledig voldaan, terwijl dit bedrag zonder betaling in de boedel zou zijn gebleven en beschikbaar zou zijn geweest voor gezamenlijke schuldeisers. In faillissement zou Menzis/CAK hoogstens als schuldeiser hebben meegedeeld naar rang. Daarmee is de paritas creditorum verstoord en zijn de overige schuldeisers in hun verhaalsmogelijkheden benadeeld. Dat Menzis stelt dat de voorschotten niet waren bedoeld om de verhaalsmogelijkheden van concurrente schuldeisers te vergroten, maakt dit niet anders; beslissend is de feitelijke vermogensvergelijking op het moment dat de curator de vernietiging inroept. Uit de casus volgt onvoldoende concreet dat Menzis/CAK het bedrag zonder de betaling toch volledig en rechtmatig zou hebben kunnen verrekenen of verhalen, zodat dit niet aan benadeling in de weg staat.

\paragraph{Wetenschap - overwegingen}\label{wetenschap---overwegingen-82}

De betaling vond plaats zeer kort voor de surseance en het faillissement. Uit de feiten volgt rechtstreeks dat TZG haar financiële noodsituatie kende. Ook Menzis beschikte vóór de betaling over concrete informatie waaruit bleek dat TZG technisch failliet was, na korte tijd geen middelen meer zou hebben, geen continuïteitsverklaring kon krijgen en substantiële schuldeisers had; bovendien heeft Menzis kort voor de betaling de bevoorschotting beëindigd wegens ernstige twijfels. Daarmee was voor TZG en in elk geval voor Menzis ten tijde van de betaling het faillissement en een tekort daarin met een redelijke mate van waarschijnlijkheid voorzienbaar. Voor de bewijsvermoedens van art. 43 Fw biedt de casus geen toereikende grond: de betaling strekte tot terugbetaling van een gesteld te veel ontvangen voorschot en niet tot een prestatie tegen aanmerkelijke onderwaarde, niet tot voldoening van een niet-opeisbare schuld, en er is geen gelieerdheid in de zin van art. 43 lid 1 sub 3°-6° Fw. Art. 45 Fw is niet relevant omdat de betaling niet om niet is verricht.

\newpage

\markboth{JOR 2014/76}{gemini-3.5-flash (k = 0)}

\subsubsection{Gemini}\label{gemini-27}

\texttt{gemini-3.5-flash}

\begin{center}
\small
\begin{tabular}{@{}p{7.4cm}p{2.6cm}p{2.6cm}@{}}
\toprule
Variabele & Waarde (model) & Waarde (gold) \\
\midrule
Is rechtshandeling & \texttt{true} & \texttt{true} \\
Datum rechtshandeling & \texttt{2009-02-19} & \texttt{2009-02-19} \\
Verricht door de schuldenaar & \texttt{true} & \texttt{true} \\
Om niet & \texttt{false} & \texttt{false} \\
Datum faillietverklaring & \texttt{2009-03-09} & \texttt{2009-03-09} \\
Onverplicht & \textcolor{gray}{\texttt{false}} & \texttt{true} \\
Benadeling schuldeisers & \texttt{true} & \texttt{true} \\
Wetenschap schuldenaar aangetoond & \textcolor{gray}{\texttt{true}} & \texttt{false} \\
Wetenschap wederpartij aangetoond & \textcolor{gray}{\texttt{true}} & \texttt{null} \\
Vermoeden wetenschap schuldenaar tegenbewijs & \texttt{null} & \texttt{false} \\
Vermoeden wetenschap wederpartij tegenbewijs & \texttt{null} & \texttt{false} \\
Tegen onderwaarde & \texttt{null} & \texttt{null} \\
Voldoening of zekerheid niet opeisbare schuld & \textcolor{gray}{\texttt{null}} & \texttt{true} \\
Gelieerde partij & \texttt{null} & \texttt{null} \\
Reeds verplicht voor aanvang termijn & \texttt{null} & \texttt{false} \\
\textbf{Vernietigbaar} & \textbf{\textcolor{gray}{\texttt{false}}} & \textbf{\texttt{true}} \\
\bottomrule
\end{tabular}
\end{center}

\paragraph{Is rechtshandeling - overwegingen}\label{is-rechtshandeling---overwegingen-83}

De betaling van een bedrag van € 1.191.000,- door TZG aan CAK is een rechtshandeling, aangezien deze handeling is gericht op het rechtsgevolg van het tenietgaan van een verbintenis (de terugbetaling van te veel ontvangen voorschotten).

\paragraph{Verricht door de schuldenaar - overwegingen}\label{verricht-door-de-schuldenaar---overwegingen-83}

De betaling is verricht vanaf een bankrekening van TZG, de schuldenaar zelf.

\paragraph{Om niet - overwegingen}\label{om-niet---overwegingen-83}

De betaling is verricht ter terugbetaling van te veel ontvangen voorschotten (ter delging van een schuld), waardoor er sprake is van een rechtshandeling onder bezwarende titel (om baat) en dus anders dan om niet.

\paragraph{Onverplicht - overwegingen}\label{onverplicht---overwegingen-83}

De betaling van € 1.191.000,- door TZG op 19 februari 2009 betrof de terugbetaling van een te veel ontvangen voorschot voor de maanden januari en februari 2009. Op grond van de wet (onverschuldigde betaling ex artikel 6:203 BW) en de tussen partijen geldende raamovereenkomst (waarin is bepaald dat te veel betaalde bedragen kunnen worden teruggevorderd) bestond er voor TZG een rechtsplicht om dit te veel ontvangen bedrag aan voorschotten terug te betalen. Nadat Menzis bij brief van 18 februari 2009 expliciet om terugbetaling had verzocht, was deze vordering opeisbaar en in rechte afdwingbaar. De nakoming van een dergelijke opeisbare verbintenis is een verplichte rechtshandeling. Er bestond derhalve een op de wet of overeenkomst berustende rechtsplicht tot de betaling, waardoor de rechtshandeling niet onverplicht was in de zin van artikel 42 Fw.

\paragraph{Benadeling schuldeisers - overwegingen}\label{benadeling-schuldeisers---overwegingen-83}

De betaling van € 1.191.000,- door TZG aan CAK/Menzis heeft geleid tot een directe vermindering van de liquide middelen van TZG vlak voor het faillissement. Indien deze betaling achterwege was gebleven, zou dit bedrag in de faillissementsboedel zijn gevallen en beschikbaar zijn geweest voor de gezamenlijke schuldeisers. Hoewel Menzis/CAK een vordering had wegens te veel betaalde voorschotten, betrof dit een concurrente vordering. Door de betaling is deze vordering volledig voldaan, waardoor Menzis/CAK is bevoordeeld boven andere schuldeisers en de paritas creditorum is verstoord. Er is derhalve sprake van benadeling van de schuldeisers in de zin van artikel 42 Fw.

\paragraph{Wetenschap - overwegingen}\label{wetenschap---overwegingen-83}

Zowel de schuldenaar (TZG) als de wederpartij (Menzis/CAK) beschikten over directe, geobjectiveerde wetenschap van de benadeling. Uit de brief van Menzis van 18 december 2008 blijkt dat zij expliciet wist dat TZG `technisch failliet' was en spoedig geen middelen meer zou hebben. TZG wist dit zelf ook, aangezien de continuïteitsparagraaf niet kon worden ondertekend en de surseance van betaling reeds vijf dagen na de betaling werd verleend. Het faillissement en een tekort daarin waren ten tijde van de betaling op 19 februari 2009 met een redelijke mate van waarschijnlijkheid te voorzien. De bewijsvermoedens van artikel 43 Fw zijn niet van toepassing, nu er geen sprake is van onderwaarde, een niet-opeisbare schuld of een gelieerde relatie in de zin van dat artikel.

\newpage

\markboth{JOR 2014/76}{mistral/mistral-medium-3.5-128b:fp8 (k = 5)}

\subsection{k = 5}\label{k-5-9}

\subsubsection{Mistral}\label{mistral-28}

\texttt{mistral/mistral-medium-3.5-128b:fp8}

\begin{center}
\small
\begin{tabular}{@{}p{7.4cm}p{2.6cm}p{2.6cm}@{}}
\toprule
Variabele & Waarde (model) & Waarde (gold) \\
\midrule
Is rechtshandeling & \texttt{true} & \texttt{true} \\
Datum rechtshandeling & \texttt{2009-02-19} & \texttt{2009-02-19} \\
Verricht door de schuldenaar & \texttt{true} & \texttt{true} \\
Om niet & \texttt{false} & \texttt{false} \\
Datum faillietverklaring & \texttt{2009-03-09} & \texttt{2009-03-09} \\
Onverplicht & \texttt{true} & \texttt{true} \\
Benadeling schuldeisers & \texttt{true} & \texttt{true} \\
Wetenschap schuldenaar aangetoond & \texttt{false} & \texttt{false} \\
Wetenschap wederpartij aangetoond & \texttt{false} & \texttt{null} \\
Vermoeden wetenschap schuldenaar tegenbewijs & \texttt{null} & \texttt{false} \\
Vermoeden wetenschap wederpartij tegenbewijs & \texttt{null} & \texttt{false} \\
Tegen onderwaarde & \texttt{null} & \texttt{null} \\
Voldoening of zekerheid niet opeisbare schuld & \textcolor{gray}{\texttt{null}} & \texttt{true} \\
Gelieerde partij & \textcolor{gray}{\texttt{true}} & \texttt{null} \\
Reeds verplicht voor aanvang termijn & \texttt{null} & \texttt{false} \\
\textbf{Vernietigbaar} & \textbf{\texttt{true}} & \textbf{\texttt{true}} \\
\bottomrule
\end{tabular}
\end{center}

\paragraph{Is rechtshandeling - overwegingen}\label{is-rechtshandeling---overwegingen-84}

De betaling van € 1.191.000,- door TZG aan CAK (op verzoek van Menzis) is een op een rechtsgevolg gerichte wil die zich door een verklaring (overschrijving) heeft geopenbaard. Dit kwalificeert als een rechtshandeling in de zin van art. 3:33 BW.

\paragraph{Verricht door de schuldenaar - overwegingen}\label{verricht-door-de-schuldenaar---overwegingen-84}

De betaling is verricht door TZG (de schuldenaar), zoals blijkt uit de afschrijving van de bankrekening van TZG (2.8).

\paragraph{Om niet - overwegingen}\label{om-niet---overwegingen-84}

De betaling geschiedde ter terugbetaling van een te hoog voorschot (art. 9 lid 2 raamovereenkomst). Dit is een verplichting die voortvloeit uit de contractuele afspraken, waardoor er sprake is van een wederprestatie (de verstrekking van het voorschot). Derhalve is de handeling niet `om niet'.

\paragraph{Onverplicht - overwegingen}\label{onverplicht---overwegingen-84}

De betaling van € 1.191.000,- door TZG aan CAK (in opdracht van Menzis) op 19 februari 2009 is onverplicht verricht in de zin van artikel 42 lid 1 Fw.

Allereerst is vastgesteld dat TZG op 18 december 2008 door Menzis al als `technisch failliet' was bestempeld en dat TZG na periode 3 van 2009 geen middelen meer had (2.4). Menzis heeft op 17 februari 2009 (vermoedelijk een typfout, bedoeld wordt 2009) aan Meavita meegedeeld dat voorschotten voortaan alleen nog achteraf zouden worden betaald op basis van aantoonbaar gerealiseerde productie, omdat de wijze van bevoorschotting niet langer rechtmatig en doelmatig was (2.5).

De betaling van € 1.191.000,- betrof de terugbetaling van een te hoog voorschotbedrag voor januari en februari 2009 (2.6 en 2.7). Uit de raamovereenkomst tussen TZG en Menzis (2.3) blijkt dat voorschotten werden verstrekt op basis van productieafspraken en later werden aangepast aan de gerealiseerde volumes. Artikel 9 lid 2 van de raamovereenkomst bepaalt dat Menzis ten onrechte betaalde bedragen kan terugvorderen of verrekenen met opvolgende declaraties.

Hoewel Menzis dus het recht had om het te veel betaalde voorschot terug te vorderen, was TZG \textbf{niet verplicht} om dit bedrag onmiddellijk en in één keer terug te betalen. De raamovereenkomst voorziet in de mogelijkheid van verrekening met toekomstige declaraties (2.3, artikel 9 lid 2). Dit betekent dat TZG de keuze had om het bedrag in termijnen terug te betalen of te verrekenen met toekomstige vorderingen. Er bestond derhalve \textbf{geen afdwingbare rechtsplicht} voor TZG om het bedrag van € 1.191.000,- op 19 februari 2009 in één keer af te lossen.

Daarnaast was TZG op dat moment al in een financieel desastreuse situatie (2.4 en 2.5), en de surseance van betaling werd slechts vijf dagen later, op 24 februari 2009, verleend (2.9). De betaling was dus onverplicht, omdat TZG niet gedwongen kon worden om het bedrag op die wijze en op dat moment te voldoen. Morele, feitelijke of economische druk (zoals de druk van Menzis om de voorschotten aan te passen) schept geen rechtsplicht (vgl. ECLI:NL:GHSHE:2015:1681, rov. 4.5).

Tot slot kan worden opgemerkt dat de betaling niet voortvloeide uit een bestaande, afdwingbare verplichting tot onmiddellijke terugbetaling. De raamovereenkomst bood ruimte voor verrekening, wat aantoont dat de betaling niet verplicht was in de zin van artikel 42 Fw.

\paragraph{Benadeling schuldeisers - overwegingen}\label{benadeling-schuldeisers---overwegingen-84}

Op 19 februari 2009 heeft TZG een bedrag van € 1.191.000,- overgemaakt aan CAK ter terugbetaling van een te hoog voorschot voor AWBZ-zorg. Deze betaling vond plaats in een periode waarin TZG financieel in zware moeilijkheden verkeerde: Menzis had op 18 december 2008 al vastgesteld dat TZG technisch failliet was en na periode 3 van 2009 geen middelen meer had (2.4). Bovendien had Menzis op 17 februari 2009 (waarschijnlijk een typefout voor 2009) meegedeeld dat voorschotten voortaan alleen nog achteraf zouden worden betaald op basis van aantoonbaar gerealiseerde productie, omdat de wijze van bevoorschotting niet langer rechtmatig en doelmatig was (2.5).

De betaling van € 1.191.000,- was een onverplichte rechtshandeling in de zin van artikel 42 lid 1 Fw. TZG was niet verplicht om het te veel ontvangen voorschot onmiddellijk terug te betalen; de raamovereenkomst tussen Menzis en TZG voorzag in de mogelijkheid om te veel betaalde bedragen te verrekenen met opvolgende declaraties (artikel 9 lid 2 raamovereenkomst, 2.3). Door de betaling direct te verrichten in plaats van te verrekenen, heeft TZG een bedrag van € 1.191.000,- uit haar vermogen laten vloeien dat anders in de boedel zou zijn gebleven en ten goede zou zijn gekomen aan alle schuldeisers.

De benadeling van de schuldeisers is aanwezig: zonder deze betaling zou het bedrag van € 1.191.000,- deel hebben uitgemaakt van het actief van TZG op het moment van de faillietverklaring (2.10), en derhalve beschikbaar zijn geweest voor verdeling onder alle schuldeisers. De hypothetische situatie (zonder de betaling) is dus gunstiger voor de schuldeisers dan de feitelijke situatie (met de betaling).

Daarnaast was TZG zich bewust van de benadeling: uit de brief van Menzis van 18 december 2008 (2.4) blijkt dat TZG wist dat zij technisch failliet was en geen middelen meer had. Ook Menzis en CAK wisten of behoorden te weten dat de betaling benadeling van de schuldeisers tot gevolg zou hebben, gelet op de kennis van de financiële situatie van TZG (2.4, 2.5 en 3.2).

\paragraph{Wetenschap - overwegingen}\label{wetenschap---overwegingen-84}

De betaling van € 1.191.000,- door TZG aan CAK (op instigatie van Menzis) op 19 februari 2009 valt binnen de termijn van één jaar voor de faillietverklaring (24 februari 2009 surseance, 9 maart 2009 faillissement). TZG had zich niet reeds voor de aanvang van die termijn tot deze betaling verplicht (de brief van Menzis van 17 februari 2009 maakt duidelijk dat de betaling pas na die datum werd geëist). De betaling is verricht met CAK, maar Menzis was als zorgkantoor betrokken bij de bevoorschotting en heeft de terugvordering geïnitieerd. TZG en Menzis/CAK behoren tot een groepsstructuur (TZG valt onder Meavita, en Menzis is een zorgverzekeraar die in de AWBZ-context nauw betrokken is bij de financiering). Op grond van art. 43 lid 1 sub 6 Fw wordt de wetenschap van benadeling aan beide zijden (TZG en Menzis/CAK) vermoed, omdat de betaling is verricht met een groepsmaatschappij (Menzis/CAK kan in dit verband als gelieerde partij worden beschouwd, gelet op de nauwe samenwerking in de AWBZ-structuur). Er is geen sprake van tegenbewijs: Menzis had in de brief van 18 december 2008 al vastgesteld dat TZG technisch failliet was en geen middelen meer had, wat aantoont dat de benadeling van schuldeisers voorzienbaar was. De betaling was onverplicht, omdat TZG niet contractueel verplicht was om het te veel betaalde voorschot direct terug te betalen (art. 9 lid 2 van de raamovereenkomst geeft Menzis slechts het recht om te verrekenen of terug te vorderen, maar niet de plicht voor TZG om onmiddellijk te betalen).

\newpage

\markboth{JOR 2014/76}{gpt-5.5 (k = 5)}

\subsubsection{GPT}\label{gpt-28}

\texttt{gpt-5.5}

\begin{center}
\small
\begin{tabular}{@{}p{7.4cm}p{2.6cm}p{2.6cm}@{}}
\toprule
Variabele & Waarde (model) & Waarde (gold) \\
\midrule
Is rechtshandeling & \texttt{true} & \texttt{true} \\
Datum rechtshandeling & \texttt{2009-02-19} & \texttt{2009-02-19} \\
Verricht door de schuldenaar & \texttt{true} & \texttt{true} \\
Om niet & \texttt{false} & \texttt{false} \\
Datum faillietverklaring & \texttt{2009-03-09} & \texttt{2009-03-09} \\
Onverplicht & \textcolor{gray}{\texttt{false}} & \texttt{true} \\
Benadeling schuldeisers & \texttt{true} & \texttt{true} \\
Wetenschap schuldenaar aangetoond & \textcolor{gray}{\texttt{true}} & \texttt{false} \\
Wetenschap wederpartij aangetoond & \textcolor{gray}{\texttt{true}} & \texttt{null} \\
Vermoeden wetenschap schuldenaar tegenbewijs & \texttt{null} & \texttt{false} \\
Vermoeden wetenschap wederpartij tegenbewijs & \texttt{null} & \texttt{false} \\
Tegen onderwaarde & \texttt{false} & \texttt{null} \\
Voldoening of zekerheid niet opeisbare schuld & \textcolor{gray}{\texttt{false}} & \texttt{true} \\
Gelieerde partij & \texttt{false} & \texttt{null} \\
Reeds verplicht voor aanvang termijn & \texttt{null} & \texttt{false} \\
\textbf{Vernietigbaar} & \textbf{\textcolor{gray}{\texttt{false}}} & \textbf{\texttt{true}} \\
\bottomrule
\end{tabular}
\end{center}

\paragraph{Is rechtshandeling - overwegingen}\label{is-rechtshandeling---overwegingen-85}

De gewraakte handeling is de betaling/overschrijving van € 1.191.000 vanaf een bankrekening van TZG ten gunste van CAK met als omschrijving dat het een te veel ontvangen voorschot betrof. Een betaling is gericht op rechtsgevolg, namelijk vermogensoverdracht en voldoening van een gestelde terugbetalingsverplichting.

\paragraph{Verricht door de schuldenaar - overwegingen}\label{verricht-door-de-schuldenaar---overwegingen-85}

Uit de feiten volgt dat op 19 februari 2009 een bedrag van € 1.191.000 is afgeschreven van een bankrekening van TZG. De betaling is daarmee door de schuldenaar TZG zelf verricht.

\paragraph{Om niet - overwegingen}\label{om-niet---overwegingen-85}

De betaling vond plaats naar aanleiding van de mededeling van Menzis dat over januari en februari 2009 te veel was bevoorschot en het verzoek dit bedrag terug te betalen. De betaling strekte dus tot terugbetaling/voldoening van een gestelde schuld en was daarmee niet zonder wederprestatie in de relevante zin.

\paragraph{Onverplicht - overwegingen}\label{onverplicht---overwegingen-85}

De aangevochten betaling betrof de terugbetaling van een door TZG zelf berekend te veel ontvangen voorschot voor AWBZ-zorg over januari/februari 2009. Op grond van de raamovereenkomst was de bevoorschotting gekoppeld aan gerealiseerde productie en kon het zorgkantoor ten onrechte gedeclareerde en betaalde bedragen terugvorderen of verrekenen. Voor zover het voorschot hoger was dan de gerealiseerde zorgproductie, had TZG geen aanspraak op dat meerdere; terugbetaling daarvan kon in rechte worden afgedwongen, hetzij contractueel, hetzij als onverschuldigd betaalde bevoorschotting. Dat Menzis betaling op een rekening van CAK verlangde, maakt dit niet anders: betaling aan een door de schuldeiser aangewezen uitvoerings-/betaaladres is slechts de wijze van nakoming. De situatie lijkt daarmee eerder op betaling ter uitvoering van een bestaande doorbetalings- of terugbetalingsafspraak dan op een onverplichte vermogensverschuiving; vgl. in algemene zin ECLI:NL:GHLEE:2007:BB3627. Er zijn onvoldoende aanknopingspunten dat TZG vóór de betaling geen juridisch afdwingbare terugbetalingsplicht had. De betaling is daarom niet onverplicht verricht in de zin van art. 42 Fw.

\paragraph{Benadeling schuldeisers - overwegingen}\label{benadeling-schuldeisers---overwegingen-85}

De betaling van € 1.191.000 heeft het voor de gezamenlijke schuldeisers beschikbare actief van TZG kort vóór surseance en faillissement verminderd. In de hypothetische situatie zonder de betaling zou dit bedrag op de bankrekening van TZG zijn gebleven en in beginsel in de boedel zijn gevallen. Daartegenover zou Menzis/CAK hoogstens een terugvorderingsvordering wegens te veel betaald voorschot hebben gehad, maar dan als schuldeiser in de verdeling. In de feitelijke situatie is Menzis/CAK door de betaling volledig voldaan, terwijl andere schuldeisers onbetaald bleven. Ook als de passiva door de betaling met hetzelfde bedrag zijn verminderd, is sprake van benadeling doordat de paritas creditorum/rangorde tussen schuldeisers is doorbroken. Uit de casus volgt niet dat zonder de betaling alle schulden hadden kunnen worden voldaan, of dat Menzis/CAK over een zekerheidsrecht of een zonder meer daadwerkelijk en rechtsgeldig uit te oefenen verrekeningspositie beschikte waardoor de boedel het bedrag hoe dan ook niet had kunnen behouden. De omstandigheid dat het bedrag zag op te veel ontvangen AWBZ-voorschotten en niet bedoeld zou zijn geweest om het verhaal van concurrente schuldeisers te verruimen, neemt niet weg dat het bedrag feitelijk tot het vermogen van TZG behoorde en bij uitblijven van betaling voor verhaal beschikbaar zou zijn geweest. Daarom zijn de schuldeisers door de betaling benadeeld in de zin van artikel 42 Fw.

\paragraph{Wetenschap - overwegingen}\label{wetenschap---overwegingen-85}

De wetenschap van benadeling is rechtstreeks aannemelijk/aangetoond. Ten tijde van de betaling op 19 februari 2009 was voor TZG het faillissement met een tekort met een redelijke mate van waarschijnlijkheid voorzienbaar: zij verkeerde in acute financiële nood, werd kort daarna op 24 februari 2009 in surseance geplaatst en op 9 maart 2009 failliet verklaard. Ook Menzis had die wetenschap rechtstreeks: zij had op 18 december 2008 zelf vastgesteld dat TZG feitelijk technisch failliet was, na periode 3 geen middelen meer had en substantiële schuldeisers had, en staakte op 17 februari 2009 de bevoorschotting wegens ernstige twijfel over voortzetting daarvan. Voor CAK als louter betaaladres blijken minder zelfstandige feiten, maar voor Menzis als materiële wederpartij/instigator is de wetenschap aanwezig. Art. 43 Fw levert geen aanvullend bewijsvermoeden op: geen aanmerkelijke onderwaarde, geen voldoening van een niet-opeisbare schuld en geen gelieerde partij in de zin van art. 43 lid 1 sub 3°-6° Fw. Art. 45 Fw is niet van toepassing, omdat de betaling strekte tot terugbetaling/verrekening van een gesteld te hoog voorschot en dus niet als rechtshandeling om niet is te kwalificeren. Tegenbewijs tegen een wettelijk vermoeden speelt daarom niet.

\newpage

\markboth{JOR 2014/76}{gemini-3.5-flash (k = 5)}

\subsubsection{Gemini}\label{gemini-28}

\texttt{gemini-3.5-flash}

\begin{center}
\small
\begin{tabular}{@{}p{7.4cm}p{2.6cm}p{2.6cm}@{}}
\toprule
Variabele & Waarde (model) & Waarde (gold) \\
\midrule
Is rechtshandeling & \texttt{true} & \texttt{true} \\
Datum rechtshandeling & \texttt{2009-02-19} & \texttt{2009-02-19} \\
Verricht door de schuldenaar & \texttt{true} & \texttt{true} \\
Om niet & \texttt{false} & \texttt{false} \\
Datum faillietverklaring & \texttt{2009-03-09} & \texttt{2009-03-09} \\
Onverplicht & \textcolor{gray}{\texttt{false}} & \texttt{true} \\
Benadeling schuldeisers & \texttt{true} & \texttt{true} \\
Wetenschap schuldenaar aangetoond & \textcolor{gray}{\texttt{true}} & \texttt{false} \\
Wetenschap wederpartij aangetoond & \textcolor{gray}{\texttt{true}} & \texttt{null} \\
Vermoeden wetenschap schuldenaar tegenbewijs & \texttt{null} & \texttt{false} \\
Vermoeden wetenschap wederpartij tegenbewijs & \texttt{null} & \texttt{false} \\
Tegen onderwaarde & \texttt{null} & \texttt{null} \\
Voldoening of zekerheid niet opeisbare schuld & \textcolor{gray}{\texttt{null}} & \texttt{true} \\
Gelieerde partij & \texttt{null} & \texttt{null} \\
Reeds verplicht voor aanvang termijn & \texttt{null} & \texttt{false} \\
\textbf{Vernietigbaar} & \textbf{\textcolor{gray}{\texttt{false}}} & \textbf{\texttt{true}} \\
\bottomrule
\end{tabular}
\end{center}

\paragraph{Is rechtshandeling - overwegingen}\label{is-rechtshandeling---overwegingen-86}

De betaling van een bedrag van € 1.191.000,- door TZG aan CAK is een rechtshandeling, aangezien deze handeling is gericht op het rechtsgevolg van het tenietgaan van een verbintenis (de terugbetaling van te veel ontvangen voorschotten).

\paragraph{Verricht door de schuldenaar - overwegingen}\label{verricht-door-de-schuldenaar---overwegingen-86}

De betaling is verricht vanaf een bankrekening van TZG, de schuldenaar zelf.

\paragraph{Om niet - overwegingen}\label{om-niet---overwegingen-86}

De betaling is verricht ter terugbetaling van te veel ontvangen voorschotten (ter delging van een schuld), waardoor er sprake is van een rechtshandeling onder bezwarende titel (om baat) en dus anders dan om niet.

\paragraph{Onverplicht - overwegingen}\label{onverplicht---overwegingen-86}

De aangevochten rechtshandeling betreft de betaling van € 1.191.000,- door TZG aan CAK op 19 februari 2009. Deze betaling werd verricht ter terugbetaling van een te hoog voorschot over de maanden januari en februari 2009. Op grond van de systematiek van de raamovereenkomst en de wet (artikel 6:203 BW wegens onverschuldigde betaling) bestond er voor TZG een rechtsplicht om het te veel ontvangen voorschot terug te betalen. Nadat Menzis bij brief van 18 februari 2009 de omvang van het te veel bevoorschotte bedrag had vastgesteld en om terugbetaling had verzocht, was deze vordering opeisbaar en in rechte afdwingbaar. De voldoening aan een opeisbare verplichting is een verplichte rechtshandeling in de zin van artikel 42 Fw. Er bestond derhalve een op de wet en overeenkomst berustende rechtsplicht tot betaling.

\paragraph{Benadeling schuldeisers - overwegingen}\label{benadeling-schuldeisers---overwegingen-86}

De benadeling van de schuldeisers in de zin van artikel 42 Fw moet worden beoordeeld door de feitelijke situatie na de rechtshandeling te vergelijken met de hypothetische situatie waarin de schuldeisers zouden hebben verkeerd zonder die handeling. In de feitelijke situatie heeft TZG een bedrag van € 1.191.000,- overgemaakt aan CAK/Menzis ter terugbetaling van te veel ontvangen voorschotten. In de hypothetische situatie waarin deze betaling achterwege was gebleven, zou dit bedrag op de bankrekening van TZG zijn blijven staan en na de faillietverklaring (die kort daarna volgde) in de boedel zijn gevallen. Dit actieve vermogen zou dan beschikbaar zijn geweest voor de gezamenlijke schuldeisers. Door de betaling is het boedelactief met dit substantiële bedrag verminderd, waardoor de verhaalsmogelijkheden van de schuldeisers rechtstreeks zijn beperkt. Het verweer van Menzis dat het voorschot onverschuldigd was en niet bedoeld was om de verhaalsruimte van commune crediteuren te vergroten, doet hieraan niet af. De vordering tot terugbetaling van de voorschotten betreft in beginsel een concurrente vordering. Door deze vordering vlak voor het faillissement volledig te voldoen, is Menzis bevoordeeld en is de rangorde tussen schuldeisers verstoord (vergelijkbaar met de situatie in ECLI:NL:GHSHE:2009:BK7363, waarin onverplichte terugbetalingen van voorschotten eveneens tot benadeling leidden). Er is derhalve sprake van benadeling van de schuldeisers.

\paragraph{Wetenschap - overwegingen}\label{wetenschap---overwegingen-86}

In deze casus is de wetenschap van benadeling aan beide zijden rechtstreeks aangetoond op grond van artikel 42 Fw. Uit de brief van Menzis van 18 december 2008 blijkt dat Menzis wist dat TZG technisch failliet was, na periode 3 geen middelen meer zou hebben en substantiële schuldeisers had. TZG was hiervan eveneens op de hoogte en vroeg enkele dagen na de betaling (op 24 februari 2009) surseance van betaling aan. De bewijsvermoedens van artikel 43 Fw en artikel 45 Fw zijn niet van toepassing, nu er geen sprake is van een rechtshandeling om niet, onderwaarde, een niet-opeisbare schuld of een transactie met een gelieerde partij in de zin van deze artikelen.

\newpage

\markboth{JOR 2014/76}{mistral/mistral-medium-3.5-128b:fp8 (k = 10)}

\subsection{k = 10}\label{k-10-9}

\subsubsection{Mistral}\label{mistral-29}

\texttt{mistral/mistral-medium-3.5-128b:fp8}

\begin{center}
\small
\begin{tabular}{@{}p{7.4cm}p{2.6cm}p{2.6cm}@{}}
\toprule
Variabele & Waarde (model) & Waarde (gold) \\
\midrule
Is rechtshandeling & \texttt{true} & \texttt{true} \\
Datum rechtshandeling & \texttt{2009-02-19} & \texttt{2009-02-19} \\
Verricht door de schuldenaar & \texttt{true} & \texttt{true} \\
Om niet & \texttt{false} & \texttt{false} \\
Datum faillietverklaring & \texttt{2009-03-09} & \texttt{2009-03-09} \\
Onverplicht & \texttt{true} & \texttt{true} \\
Benadeling schuldeisers & \texttt{true} & \texttt{true} \\
Wetenschap schuldenaar aangetoond & \texttt{false} & \texttt{false} \\
Wetenschap wederpartij aangetoond & \texttt{false} & \texttt{null} \\
Vermoeden wetenschap schuldenaar tegenbewijs & \texttt{null} & \texttt{false} \\
Vermoeden wetenschap wederpartij tegenbewijs & \texttt{null} & \texttt{false} \\
Tegen onderwaarde & \texttt{null} & \texttt{null} \\
Voldoening of zekerheid niet opeisbare schuld & \textcolor{gray}{\texttt{null}} & \texttt{true} \\
Gelieerde partij & \textcolor{gray}{\texttt{true}} & \texttt{null} \\
Reeds verplicht voor aanvang termijn & \texttt{false} & \texttt{false} \\
\textbf{Vernietigbaar} & \textbf{\texttt{true}} & \textbf{\texttt{true}} \\
\bottomrule
\end{tabular}
\end{center}

\paragraph{Is rechtshandeling - overwegingen}\label{is-rechtshandeling---overwegingen-87}

De betaling van € 1.191.000,- door TZG aan CAK (op verzoek van Menzis) is een op een rechtsgevolg gerichte wil die zich door een verklaring (overschrijving) heeft geopenbaard. Dit kwalificeert als een goederenrechtelijke handeling (betaling) en valt derhalve onder het begrip rechtshandeling.

\paragraph{Verricht door de schuldenaar - overwegingen}\label{verricht-door-de-schuldenaar---overwegingen-87}

De betaling is verricht door TZG (de schuldenaar), zoals blijkt uit de afschrijving van haar bankrekening (2.8).

\paragraph{Om niet - overwegingen}\label{om-niet---overwegingen-87}

De betaling geschiedde ter terugbetaling van een te hoog voorschot (art. 9 lid 2 raamovereenkomst). Dit impliceert een verplichting tot terugbetaling, zodat er sprake is van een tegenprestatie (het niet hoeven terugvorderen via andere middelen). Derhalve is de handeling niet `om niet'.

\paragraph{Onverplicht - overwegingen}\label{onverplicht---overwegingen-87}

De betaling van € 1.191.000,- door TZG aan CAK op 19 februari 2009 is onverplicht verricht in de zin van artikel 42 lid 1 Fw.

Uit de raamovereenkomst tussen Menzis (als zorgkantoor) en TZG (als zorgaanbieder) blijkt dat Menzis voorschotten verstrekt op basis van een lineaire extrapolatie van de gerealiseerde productie (artikel 5, lid 4). Artikel 9, lid 2 van de raamovereenkomst bepaalt dat Menzis het recht heeft om ten onrechte betaalde bedragen terug te vorderen of te verrekenen met opvolgende declaraties.

De brief van Menzis van 17 februari 2009 (waarschijnlijk 2009, gelet op de context) maakt duidelijk dat Menzis vanaf 22 februari 2009 geen voorschotten meer vooraf zou verlenen, maar alleen nog achteraf zou betalen op basis van aantoonbaar gerealiseerde productie. Dit betekent dat het voorschotbedrag voor januari en februari 2009, dat € 1.191.000,- te hoog bleek te zijn, niet direct opeisbaar was. Er bestond geen rechtsplicht voor TZG om dit bedrag onmiddellijk terug te betalen; Menzis had slechts het recht om dit bedrag terug te vorderen of te verrekenen met toekomstige declaraties.

Daarom was de betaling van € 1.191.000,- door TZG aan CAK een onverplichte rechtshandeling, omdat er geen op de wet of een overeenkomst berustende verplichting bestond om dit bedrag op dat moment te voldoen. De stelplicht en bewijslast van het onverplichte karakter rusten op de curator, en deze heeft voldoende onderbouwd dat de betaling niet verplicht was.

\paragraph{Benadeling schuldeisers - overwegingen}\label{benadeling-schuldeisers---overwegingen-87}

Op basis van de casus en de toepasselijke jurisprudentie kan worden geconcludeerd dat de schuldeisers door de betaling van € 1.191.000,- door TZG aan CAK (op instigatie van Menzis) zijn benadeeld in de zin van artikel 42 Faillissementswet (Fw).

Allereerst is de betaling onverplicht in de zin van artikel 42 lid 1 Fw. Uit de raamovereenkomst (artikel 5, lid 4) blijkt dat Menzis voorschotten verstrekt op basis van een lineaire extrapolatie van de gerealiseerde volumes, met de mogelijkheid om deze aan te passen. Artikel 9, lid 2 van de raamovereenkomst bepaalt dat Menzis het recht heeft om ten onrechte betaalde bedragen terug te vorderen of te verrekenen met opvolgende declaraties. Dit betekent dat TZG niet verplicht was om het te veel ontvangen voorschotbedrag direct terug te betalen aan CAK. De betaling was derhalve onverplicht.

Vervolgens is er sprake van benadeling van de schuldeisers. De benadeling moet worden beoordeeld door een vermogensvergelijking: de hypothetische situatie waarin de schuldeisers zouden hebben verkeerd zonder de gewraakte rechtshandeling wordt vergeleken met de situatie waarin zij feitelijk verkeren als die handeling in stand blijft. In dit geval had TZG, zonder de betaling van € 1.191.000,- aan CAK, dit bedrag kunnen aanwenden om haar schuldeisers te voldoen. Omdat TZG op 24 februari 2009 surseance van betaling heeft gekregen en op 9 maart 2009 failliet is verklaard, is het aannemelijk dat de schuldeisers in de hypothetische situatie een grotere verhaalsmogelijkheid zouden hebben gehad. De betaling heeft dus geleid tot een vermindering van het actief van TZG, wat de schuldeisers benadeelt.

Ten slotte is er sprake van wetenschap van benadeling. Uit de brief van Menzis aan Meavita van 18 december 2008 (2.4) blijkt dat Menzis wist dat TZG technisch failliet was en na periode 3 van 2009 geen middelen meer had. Daarnaast blijkt uit de brief van 17 februari 2009 (2.5) dat Menzis besloten had om na 22 februari 2009 geen voorschotten meer vooraf te verlenen aan TZG, omdat de financiële situatie van TZG dermate slecht was dat bevoorschotting niet langer verantwoord was. Dit betekent dat Menzis wist of behoorde te weten dat de betaling van € 1.191.000,- door TZG aan CAK benadeling van de schuldeisers tot gevolg zou hebben. Omdat de betaling een rechtshandeling anders dan om niet betreft (er was een tegenprestatie in de vorm van de verplichting tot terugbetaling van het te veel ontvangen voorschot, maar deze verplichting was niet onmiddellijk opeisbaar), is ook vereist dat CAK wist of behoorde te weten dat de betaling benadeling van de schuldeisers tot gevolg zou hebben (artikel 42 lid 2 Fw). Gelet op de nauwe samenwerking tussen Menzis en CAK in het kader van de AWBZ-regeling en de kennis van Menzis over de financiële situatie van TZG, kan worden aangenomen dat CAK eveneens wist of behoorde te weten dat de betaling benadelend was voor de schuldeisers.

Gelet op het voorgaande zijn alle voorwaarden voor de toepassing van artikel 42 Fw vervuld: de betaling was onverplicht, heeft geleid tot benadeling van de schuldeisers, en zowel TZG als CAK wisten of behoorden te weten dat de betaling benadeling van de schuldeisers tot gevolg zou hebben.

\paragraph{Wetenschap - overwegingen}\label{wetenschap---overwegingen-87}

De betaling van € 1.191.000,- door TZG aan CAK (op instigatie van Menzis) op 19 februari 2009 valt binnen de termijn van één jaar voor de faillietverklaring van TZG (24 februari 2009 surseance, 9 maart 2009 faillissement). TZG en Menzis/CAK zijn gelieerde partijen in de zin van art. 43 lid 1 sub 6 Fw, omdat TZG en Menzis deel uitmaken van een groep stichtingen geleid door Meavita (art. 43 lid 1 sub 6 Fw). Hierdoor wordt de wetenschap van benadeling aan beide zijden vermoed (art. 43 lid 1 Fw). TZG had zich niet reeds voor de aanvang van de jaartermijn (19 februari 2008) tot de betaling verplicht, zodat het vermoeden van toepassing is. Er is geen sprake van tegenbewijs dat het vermoeden ontkracht, gelet op de brief van Menzis aan Meavita van 18 december 2008, waarin Menzis vaststelde dat TZG technisch failliet was en na periode 3 van 2009 geen middelen meer had. Dit betekent dat de wetenschap van benadeling aan beide zijden wordt vermoed en niet is weerlegd.

\newpage

\markboth{JOR 2014/76}{gpt-5.5 (k = 10)}

\subsubsection{GPT}\label{gpt-29}

\texttt{gpt-5.5}

\begin{center}
\small
\begin{tabular}{@{}p{7.4cm}p{2.6cm}p{2.6cm}@{}}
\toprule
Variabele & Waarde (model) & Waarde (gold) \\
\midrule
Is rechtshandeling & \texttt{true} & \texttt{true} \\
Datum rechtshandeling & \texttt{2009-02-19} & \texttt{2009-02-19} \\
Verricht door de schuldenaar & \texttt{true} & \texttt{true} \\
Om niet & \texttt{false} & \texttt{false} \\
Datum faillietverklaring & \texttt{2009-03-09} & \texttt{2009-03-09} \\
Onverplicht & \textcolor{gray}{\texttt{false}} & \texttt{true} \\
Benadeling schuldeisers & \texttt{true} & \texttt{true} \\
Wetenschap schuldenaar aangetoond & \textcolor{gray}{\texttt{true}} & \texttt{false} \\
Wetenschap wederpartij aangetoond & \textcolor{gray}{\texttt{true}} & \texttt{null} \\
Vermoeden wetenschap schuldenaar tegenbewijs & \texttt{null} & \texttt{false} \\
Vermoeden wetenschap wederpartij tegenbewijs & \texttt{null} & \texttt{false} \\
Tegen onderwaarde & \texttt{false} & \texttt{null} \\
Voldoening of zekerheid niet opeisbare schuld & \textcolor{gray}{\texttt{false}} & \texttt{true} \\
Gelieerde partij & \texttt{false} & \texttt{null} \\
Reeds verplicht voor aanvang termijn & \texttt{null} & \texttt{false} \\
\textbf{Vernietigbaar} & \textbf{\textcolor{gray}{\texttt{false}}} & \textbf{\texttt{true}} \\
\bottomrule
\end{tabular}
\end{center}

\paragraph{Is rechtshandeling - overwegingen}\label{is-rechtshandeling---overwegingen-88}

De gewraakte handeling is de betaling/overschrijving van € 1.191.000 vanaf een bankrekening van TZG ten gunste van CAK, met vermelding dat het ging om een te veel betaald voorschot. Een betaling door overschrijving is gericht op rechtsgevolg, namelijk voldoening of terugbetaling van een gestelde schuld.

\paragraph{Verricht door de schuldenaar - overwegingen}\label{verricht-door-de-schuldenaar---overwegingen-88}

Het bedrag is op 19 februari 2009 afgeschreven van een bankrekening van TZG. De betaling is daarmee door de schuldenaar zelf verricht.

\paragraph{Om niet - overwegingen}\label{om-niet---overwegingen-88}

De betaling vond plaats als terugbetaling van een volgens Menzis te veel verstrekt voorschot over januari en februari 2009. Daarmee stond tegenover de betaling de beoogde voldoening of vermindering van een gestelde terugbetalingsvordering; de handeling is daarom anders dan om niet.

\paragraph{Onverplicht - overwegingen}\label{onverplicht---overwegingen-88}

De aangevochten rechtshandeling is de betaling op 19 februari 2009 van € 1.191.000 aan CAK, met als omschrijving dat sprake was van een te veel ontvangen voorschot. Uit de raamovereenkomst en het systeem van AWBZ-bevoorschotting volgt dat voorschotten voorlopig waren en moesten worden afgestemd op de daadwerkelijk gerealiseerde productie. Artikel 9 van de raamovereenkomst geeft het zorgkantoor bovendien het recht om ten onrechte betaalde bedragen terug te vorderen dan wel te verrekenen. TZG heeft zelf op 18 februari 2009 een overzicht verstrekt waaruit bleek dat over januari/februari 2009 € 1.191.126 te veel was bevoorschot, waarna Menzis terugbetaling heeft gevraagd. Daarmee bestond voor TZG een in rechte afdwingbare verplichting tot terugbetaling van het te veel ontvangen voorschot, hetzij contractueel, hetzij uit onverschuldigde betaling. Dat de betaling via CAK liep, sluit aan bij de AWBZ-betaalstructuur en maakt de voldoening niet onverplicht. De curator heeft onvoldoende aanknopingspunten om aan te nemen dat slechts sprake was van morele, feitelijke of economische druk of van betaling van een nog niet opeisbare schuld. De betaling was daarom niet onverplicht in de zin van artikel 42 Fw.

\paragraph{Benadeling schuldeisers - overwegingen}\label{benadeling-schuldeisers---overwegingen-88}

Bij de vereiste vermogensvergelijking moet worden vergeleken de feitelijke situatie waarin TZG op 19 februari 2009 € 1.191.000 aan CAK/Menzis heeft betaald met de hypothetische situatie waarin die betaling achterwege was gebleven. In die hypothetische situatie zou dit bedrag op de bankrekening van TZG zijn gebleven en bij de kort daarop gevolgde surseance en het faillissement in beginsel beschikbaar zijn geweest voor verhaal door de gezamenlijke schuldeisers. Daartegenover zou CAK/Menzis hoogstens een vordering tot terugbetaling van het te veel bevoorschotte bedrag hebben gehad. In de feitelijke situatie is het actief van TZG met € 1.191.000 verminderd en is CAK/Menzis kort vóór faillissement volledig voldaan, terwijl andere substantiële schuldeisers onbetaald bleven. Dat het ging om een te hoog voorschot of mogelijk onverschuldigd betaalde AWBZ-gelden neemt de benadeling niet weg: ook betaling van een bestaande schuld kan in insolventie de paritas creditorum verstoren doordat één schuldeiser volledig wordt voldaan ten koste van de overige schuldeisers. Er zijn onvoldoende concrete aanwijzingen dat CAK/Menzis zich ook zonder deze betaling volledig en onvermijdelijk buiten faillissement had kunnen verhalen, bijvoorbeeld door een zekerheidsrecht of volledige verrekening, of dat de overige schuldeisers zich niet op het bedrag hadden kunnen verhalen. De aangevochten betaling leidt daarom tot vermindering van het voor de schuldeisers beschikbare actief en tot rangordeverstoring.

\paragraph{Wetenschap - overwegingen}\label{wetenschap---overwegingen-88}

De wetenschap van benadeling is hier rechtstreeks aannemelijk gemaakt. Ten tijde van de betaling op 19 februari 2009 waren faillissement en een tekort daarin met een redelijke mate van waarschijnlijkheid te voorzien: Menzis had al op 18 december 2008 zelf vastgesteld dat TZG ``in feite technisch failliet'' was, na periode 3 geen middelen meer had, dat het herstelplan onhaalbaar was, de accountant geen verklaring kon afgeven en er substantiële schuldeisers waren. Op 17 februari 2009 heeft Menzis bovendien wegens ernstige twijfel de bevoorschotting vooraf stopgezet. Voor TZG gold, gelet op haar eigen financiële situatie en de kort daarop gevolgde surseance en faillietverklaring, eveneens dat zij dit wist of behoorde te weten. Daarmee is voldaan aan de strenge maatstaf dat niet slechts een kans op benadeling, maar faillissement en tekort redelijkerwijs voorzienbaar moeten zijn (vgl. de maatstaf zoals ook benadrukt in ECLI:NL:GHSHE:2014:4460 en HR 2016:1088). Voor zover Menzis als materiële wederpartij geldt, is haar wetenschap rechtstreeks aangetoond; voor CAK als louter betalingskanaal blijkt uit de feiten geen zelfstandige wetenschap. De bewijsvermoedens van art. 43 Fw lijken niet toepasselijk: geen gelieerde partij, geen duidelijke onderwaarde-overeenkomst en geen voldoening van een niet-opeisbare schuld. Tegenbewijs tegen een vermoeden speelt daarom niet.

\newpage

\markboth{JOR 2014/76}{gemini-3.5-flash (k = 10)}

\subsubsection{Gemini}\label{gemini-29}

\texttt{gemini-3.5-flash}

\begin{center}
\small
\begin{tabular}{@{}p{7.4cm}p{2.6cm}p{2.6cm}@{}}
\toprule
Variabele & Waarde (model) & Waarde (gold) \\
\midrule
Is rechtshandeling & \texttt{true} & \texttt{true} \\
Datum rechtshandeling & \texttt{2009-02-19} & \texttt{2009-02-19} \\
Verricht door de schuldenaar & \texttt{true} & \texttt{true} \\
Om niet & \texttt{false} & \texttt{false} \\
Datum faillietverklaring & \texttt{2009-03-09} & \texttt{2009-03-09} \\
Onverplicht & \textcolor{gray}{\texttt{false}} & \texttt{true} \\
Benadeling schuldeisers & \texttt{true} & \texttt{true} \\
Wetenschap schuldenaar aangetoond & \textcolor{gray}{\texttt{true}} & \texttt{false} \\
Wetenschap wederpartij aangetoond & \textcolor{gray}{\texttt{true}} & \texttt{null} \\
Vermoeden wetenschap schuldenaar tegenbewijs & \texttt{null} & \texttt{false} \\
Vermoeden wetenschap wederpartij tegenbewijs & \texttt{null} & \texttt{false} \\
Tegen onderwaarde & \texttt{null} & \texttt{null} \\
Voldoening of zekerheid niet opeisbare schuld & \textcolor{gray}{\texttt{null}} & \texttt{true} \\
Gelieerde partij & \texttt{null} & \texttt{null} \\
Reeds verplicht voor aanvang termijn & \texttt{null} & \texttt{false} \\
\textbf{Vernietigbaar} & \textbf{\textcolor{gray}{\texttt{false}}} & \textbf{\texttt{true}} \\
\bottomrule
\end{tabular}
\end{center}

\paragraph{Is rechtshandeling - overwegingen}\label{is-rechtshandeling---overwegingen-89}

De betaling van een bedrag van € 1.191.000,- door TZG aan CAK is een rechtshandeling, aangezien deze handeling is gericht op het rechtsgevolg van het tenietgaan van een verbintenis (de terugbetaling van te veel ontvangen voorschotten).

\paragraph{Verricht door de schuldenaar - overwegingen}\label{verricht-door-de-schuldenaar---overwegingen-89}

De betaling is verricht vanaf een bankrekening van TZG, de schuldenaar zelf.

\paragraph{Om niet - overwegingen}\label{om-niet---overwegingen-89}

De betaling is verricht ter terugbetaling van te veel ontvangen voorschotten (ter delging van een schuld), waardoor er sprake is van een rechtshandeling onder bezwarende titel (om baat) en dus anders dan om niet.

\paragraph{Onverplicht - overwegingen}\label{onverplicht---overwegingen-89}

De betaling van € 1.191.000,- door TZG aan CAK/Menzis op 19 februari 2009 vond plaats nadat was vastgesteld dat er over de maanden januari en februari 2009 een bedrag van € 1.191.126,- te veel aan voorschotten was verstrekt. Een te veel ontvangen voorschot leidt tot een rechtsplicht tot terugbetaling, hetzij op grond van de overeenkomst, hetzij op grond van de wet (onverschuldigde betaling, art. 6:203 BW). Menzis heeft dit bedrag op 18 februari 2009 expliciet teruggevorderd, waardoor de vordering opeisbaar werd. De betaling door TZG strekte derhalve tot voldoening van een opeisbare schuld. Aangezien er een in rechte afdwingbare rechtsplicht bestond tot terugbetaling van het te veel ontvangen voorschot, is de betaling aan te merken als een verplichte rechtshandeling. Het feit dat Menzis op grond van de raamovereenkomst ook had kunnen verrekenen met toekomstige declaraties, doet niet af aan het bestaan van de zelfstandige terugbetalingsverplichting na terugvordering.

\paragraph{Benadeling schuldeisers - overwegingen}\label{benadeling-schuldeisers---overwegingen-89}

Om te bepalen of er sprake is van benadeling van schuldeisers in de zin van artikel 42 Faillissementswet (Fw), dient een vermogensvergelijking te worden gemaakt tussen de feitelijke situatie na de betaling en de hypothetische situatie waarin de betaling achterwege zou zijn gebleven. In de feitelijke situatie is een bedrag van € 1.191.000,- aan het vermogen van TZG onttrokken ten gunste van CAK/Menzis. In de hypothetische situatie waarin deze betaling niet was verricht, zou dit bedrag van € 1.191.000,- in de boedel van TZG zijn gevallen. Menzis zou in dat geval voor het te veel ontvangen voorschot slechts een concurrente vordering hebben gehad in het faillissement. Aangezien TZG ten tijde van de betaling reeds in ernstige financiële problemen verkeerde (`technisch failliet' met substantiële schuldeisers) en het faillissement zeer kort daarna is uitgesproken, is het onvermijdelijk dat de concurrente schuldeisers in het faillissement niet volledig voldaan kunnen worden. Door de betaling is Menzis voor 100\% voldaan, waardoor het actieve boedelvolume met € 1.191.000,- is verminderd en de verhaalsmogelijkheden van de overige schuldeisers zijn beperkt. Dit levert een duidelijke verstoring van de rangorde (paritas creditorum) en een vermindering van het actief op. De stelling dat het hier gaat om de terugbetaling van onverschuldigd betaalde voorschotten of publieke middelen doet hieraan niet af; een terugbetaling van een vordering (zelfs uit onverschuldigde betaling) vóór faillissement leidt tot benadeling van de gezamenlijke schuldeisers, omdat de desbetreffende crediteur anders slechts als concurrent had moeten meedelen (overeenkomstig ECLI:NL:GHSHE:2014:4568). Er is derhalve sprake van benadeling van schuldeisers.

\paragraph{Wetenschap - overwegingen}\label{wetenschap---overwegingen-89}

De wetenschap van benadeling is aan beide zijden rechtstreeks aangetoond op grond van artikel 42 Fw. Uit de brief van Menzis van 18 december 2008 blijkt dat Menzis wist dat TZG technisch failliet was en na periode 3 geen middelen meer zou hebben. TZG was uiteraard ook volledig op de hoogte van haar eigen desastreuze financiële situatie. Ten tijde van de betaling op 19 februari 2009, slechts vijf dagen voor de surseanceaanvraag, was het faillissement en een tekort daarin met een redelijke mate van waarschijnlijkheid te voorzien. Bewijsvermoedens op grond van artikel 43 of 45 Fw zijn niet van toepassing.

\end{document}
